\bookpart{知识的获取}{Knowledge Acquisition}
\label{part:efficient-knowledge-use}

\BookPartGuide{%
训练需要可用的答案、关系、约束和目标。人工、规则、模型及样本内部的结构可以提供不同依据；取得这些依据后，融合既组织多来源样本标注，也说明已有模型和模块怎样共同使用。质量评估与修正再检查学习依据是否可靠且有用。}

\chapter{知识的获取概述}
\label{chap:knowledge-acquisition-overview}

假设我们要训练一个识别当地花卉物种的模型。使用时，把一张花卉照片交给模型，希望它返回“鸢尾”或“玫瑰”等物种名称。照片就是这项任务的预测输入，物种名称就是希望得到的预测输出。现在我们已经拍下了一批照片，但这些文件只有像素，没有附上物种答案。如果模型把一张鸢尾照片判成玫瑰，学习算法还缺少一项依据，无法按物种是否正确来检查这次预测。

一个直接的办法是请植物学家看照片，把物种名称写在每张照片旁边。这样，照片甲就由单独的图像变成了“照片甲及其物种答案”。训练时，学习算法把模型的预测与这个答案比较，再据此调整模型内部的参数。参数是决定模型怎样从输入计算输出的数值；训练就是利用样本调整这些数值的过程。这里已经有两个不同的问题：怎样根据答案训练模型，以及怎样取得值得使用的答案。本章讨论后一个问题。

任务要求改变时，缺少的依据也会改变。整张照片的物种答案不能告诉模型花朵位于哪个区域；若希望模型从不同角度认出同一株花，还要说明哪些照片拍的是同一对象。有时没有人能逐张给出物种名称，但我们可以遮住照片的一部分，让模型根据其余部分恢复原内容。类别答案、对象关系、区域位置和恢复要求都可以影响学习，不过它们提供的信息不同，必须按各自用途取得和检查。

本章把为学习补入这些依据的过程称为知识获取。这里的“知识”指与任务相联系、可以用于训练的答案、关系、约束或目标，并不意味着凡是写入训练文件的内容都已经正确。我们需要说明补入了什么、依据来自哪里、在哪些对象上适用，以及经过了什么检查。沿着这批花卉照片，先弄清需要保存和区分的内容，再讨论如何选择来源、组织获取和验收；最后用代表性研究说明这些办法怎样扩展、为什么仍有使用条件。阅读本章只需理解函数输入输出、向量、概率和按条件选择方案的基本思想，各类方法的完整算法由后续章节展开。

\Needspace{6\baselineskip}
\section{知识获取的基本概念}
\label{sec:knowledge-acquisition-problem}

拍摄花卉得到的颜色、纹理和形状是我们实际取得的观测。花卉的物种、所在地区和光照条件描述被研究的环境；镜头取景、曝光和传感器成像决定我们能观察到什么；选择哪些花、在哪些时刻拍照，则决定这批样本怎样取得。例如，只在晴天拍路边常见花，即使照片很多，也没有观察到雨天或稀有物种。把环境、观测方式和选样方式分开，才能判断缺少的是照片，还是照片所对应的任务依据。

\subsection{从观测到训练依据}
\label{sec:knowledge-acquisition-observation-supervision}

先看最直接的类别答案。给照片甲写入“鸢尾”，得到的是\term{标签}（label），即这张照片在指定分类任务中的目标类别。为了利用它，学习算法需要给预测的偏差计分，这个计分函数称为损失函数（loss function）。一个用于说明思想的计分办法是：预测类别与标签相同计0，不同计1；多张照片上的平均分就表示错分比例。实际训练可以采用其他便于调整参数的损失函数，但都要先规定预测与目标怎样比较。标签提供比较依据，损失函数规定比较方法，两者承担不同职责。

也可能有两张照片乙、丙，我们不知道花的物种，却知道它们拍的是同一株花。希望模型认出同一对象时，可以要求它对两张照片产生相近的描述。模型用一组数值描述图像，这组数值称为表示（representation），可写成特征向量；\term{表示学习}（representation learning）就是通过训练取得这种描述方式，而非完全由人预先规定每个特征。设两张照片的表示分别为$\bm{h}_{\text{乙}},\bm{h}_{\text{丙}}\in\mathbb{R}^{d}$，其中$d$是特征数。为了惩罚同一对象的表示差异，可以在训练中加入
\[
  \ell_{\mathrm{pair}}
  =\lVert\bm{h}_{\text{乙}}-\bm{h}_{\text{丙}}\rVert_2^2.
\]
这里$\lVert\cdot\rVert_2$计算向量的欧氏长度，作用于两组表示之差，就得到它们的欧氏距离；$\ell_{\mathrm{pair}}$是这对照片的距离平方。例如，两组二维表示为$(1,2)$与$(2,2)$时，距离平方为1；调整为相同向量时，这项损失为0。配对关系告诉算法哪些表示应当接近，并没有提供物种名称。还须确认配对确实指向同一对象，而且不能只最小化这一项：把所有照片都映射到同一个向量也能使它为0，却失去了区分不同花卉的能力。完整训练还需保留区分其他对象的信息。

另一类依据规定哪些预测允许同时成立。如果任务已经规定一株花只能属于一个物种，“玫瑰”和“鸢尾”就不能同时成为该对象的最终类别。这种限制是预测约束（prediction constraint），它排除不合法的答案组合，却没有直接告诉我们该选哪一个类别。约束针对的是一株花；一张同时拍到玫瑰和鸢尾的照片可以包含两种花，不能把对象层面的互斥条件误用到整张照片上。采用不同任务或不同对象单位时，约束也要重新说明。

没有物种标签时，照片原有内容还可以提供另一种比较依据。我们先保留完整照片，再遮住其中一块花瓣，把剩余图像交给模型，要求它预测被遮住的内容。原来保留的图像块就是恢复目标（reconstruction target），预测图像块与它的差异可以用于计算损失。只有遮挡后的图像进入预测输入，原图像块另作目标，才能形成需要根据上下文完成的任务。这属于\term{自监督学习}（self-supervised learning）：从观测自身构造学习目标，减少对逐例人工答案的依赖。恢复颜色和纹理可能帮助学习图像表示，但恢复得好还不能证明模型已经会辨认物种。

这些内容进入训练、规定模型应当回答什么或满足什么要求时，都承担\term{监督信号}（supervision signal）的作用。这里的“监督”不限于人工提供物种答案；对象关系、预测约束和自动构造的恢复目标也能指导参数调整。尚未确认适用性与正确性的监督信号称为\term{候选监督信号}（candidate supervision signal）。能够据某个标签计算损失，只说明训练可以使用这个答案，并不说明答案可靠；错误标签也可以让算法执行一次参数更新。标签怎样进入学习的通用分工可参考\BookChapterRef{chap:machine-learning-paradigms}{卷3《基础、理论与可信性》的“范式简介”章}。

\subsection{训练依据的保存与状态}
\label{sec:knowledge-acquisition-record-status}

取得答案以后，我们还要让另一位研究者能够正确读取和复核它。只把“鸢尾”写进文件，会丢掉几个必要信息：它属于照片甲还是照片乙，依据哪份物种名录，由谁在什么材料下作答，适用于哪个地区。如果名录版本改变，同一个名称也可能对应不同的分类范围。因此，训练依据需要连同这些说明一起保存；本部分把这样的记录称为\term{样本标注}（sample annotation）。这个用法比单个类别标签宽：一条记录也可以保存两张照片之间的关系或一个恢复目标。

以人工给照片甲作答为例，一条可追溯的记录可以写成：对象为照片甲，内容为“鸢尾”，来源为标注者甲及其所见照片和所用名录，适用范围为本地区、名录版本$v$下的物种分类。作答者已提交答案，说明获取操作完成、内容已经给出；在独立复核之前，这个答案仍是候选。对象、内容和范围供训练解释答案，来源供复核追查依据，状态则说明当前取得和检查到了哪一步。下面的记号只是把这几项合写成一条记录。

\begin{definition}[样本标注]
  \label{def:knowledge-acquisition-record}
  本部分的样本标注保存一项训练依据及其来源、范围和状态，记为
  \[
    r=(o,s,c,\Omega,z_{\mathrm{sem}},z_{\mathrm{ver}}).
  \]
  其中$o$是所针对的对象，$c$是答案、关系、约束或目标等内容，$\Omega$规定内容适用的样本范围与任务范围。

  来源字段$s=(s_{\mathrm{src}},s_{\mathrm{cond}},z_{\mathrm{act}},z_{\mathrm{exec}})$依次保存来源标识、取得条件、来源动作和执行状态。

  语义状态$z_{\mathrm{sem}}$说明任务内容是否确定、对象是否适用，取值为已给出、未知或不适用；验证状态$z_{\mathrm{ver}}$说明指定检查的进展，取值为候选、已核验或已拒绝。
\end{definition}

在刚才的例子中，$o$对应照片甲，$c$对应“鸢尾”，$\Omega$对应地区、名录版本和任务范围。$s_{\mathrm{src}}$保存标注者身份，$s_{\mathrm{cond}}$保存他看到的材料与作答条件；$z_{\mathrm{act}}$为“作答”，$z_{\mathrm{exec}}$为“成功”。最后两个字段分别为$z_{\mathrm{sem}}=\text{已给出}$与$z_{\mathrm{ver}}=\text{候选}$。这里的“语义”指内容在任务中的含义，“验证”指按指定依据进行的检查。它们分开以后，系统就能保存一个已经提交、但尚未核验的答案。

为了看清为什么需要多个状态，假设照片甲的关键花瓣被遮住。植物学家认真检查了照片，仍不能辨认物种，于是主动不提交类别判断。他完成了请求，因此执行状态是“成功”；他没有给出类别，因此内容是\term{未知}（unknown），即任务答案尚未确定；主动不作判断的动作则称为\term{弃权}（abstention）。执行成功与物种未知可以同时成立。如果程序根本没有打开照片就失败，虽然也没有答案，却没有发生检查后的主动弃权，后续应先修程序或重试。

再假设误把一张石头照片送进只标注花卉的任务。石头不在任务对象范围内，应记为\term{不适用}（not applicable）；继续寻找它的花卉物种没有意义。未知对象可能需要补拍局部或转交专家，不适用对象需要排除或转交另一任务，程序失败需要恢复执行。把它们都存成一个空值，会让后续处理无法选择行动。表~\ref{tab:knowledge-acquisition-record-states}先列各层状态，再用照片甲的几种情形说明它们怎样组合，后续章节沿用这套约定。

\begin{table*}[htbp]
  \centering
  \small
  \setlength{\tabcolsep}{4pt}
  \begin{tabularx}{\linewidth}{@{}>{\RaggedRight\setlength{\parindent}{0pt}\arraybackslash}p{25mm}>{\RaggedRight\setlength{\parindent}{0pt}\arraybackslash}p{55mm}>{\RaggedRight\setlength{\parindent}{0pt}\arraybackslash}X@{}}
    \toprule
    \textbf{样本标注层} & \textbf{允许取值} & \textbf{回答的问题} \\
    \midrule
    执行状态 & 未请求、未处理、处理中、成功、失败 & 请求是否发起；发起后是否已开始、正在处理、完成或失败 \\
    来源动作 & 作答、弃权；无动作时为空 & 来源提交了判断，还是检查后主动不提交判断 \\
    语义状态 & 已给出、未知、不适用 & 内容是否已确定，或者对象本就不属于任务范围 \\
    验证状态 & 候选、已核验、已拒绝 & 针对指定属性的检查是否完成，结果是否允许使用 \\
    规则触发结果 & 触发、未触发、无法确定 & 规则条件是否成立；以具名结果保存在$s_{\mathrm{cond}}$中 \\
    \bottomrule
  \end{tabularx}

  \medskip

  \begin{tabularx}{\linewidth}{@{}>{\RaggedRight\setlength{\parindent}{0pt}\arraybackslash}p{30mm}>{\RaggedRight\setlength{\parindent}{0pt}\arraybackslash}p{19mm}>{\RaggedRight\setlength{\parindent}{0pt}\arraybackslash}p{17mm}>{\RaggedRight\setlength{\parindent}{0pt}\arraybackslash}p{30mm}>{\RaggedRight\setlength{\parindent}{0pt}\arraybackslash}p{18mm}>{\RaggedRight\setlength{\parindent}{0pt}\arraybackslash}X@{}}
    \toprule
    \textbf{样本标注情形} & \textbf{执行} & \textbf{动作} & \textbf{内容／语义} & \textbf{验证} & \textbf{条件或核验依据} \\
    \midrule
    尚未向人工发起请求 & 未请求 & — & 空／未知 & 候选 & 尚未发起任务 \\
    人工提交物种判断 & 成功 & 作答 & 鸢尾／已给出 & 候选 & 依据照片与名录作答 \\
    人工因遮挡而弃权 & 成功 & 弃权 & 空／未知 & 候选 & 辨认物种所需的花瓣局部不可见 \\
    规则运行但未触发 & 成功 & — & 空／未知 & 候选 & 条件检查结果：未触发 \\
    复核确认第二条正确 & 成功 & 作答 & 鸢尾／已给出 & 已核验 & 独立证据支持第二行的物种判断 \\
    \bottomrule
  \end{tabularx}
  \caption{分层保存获取与核验状态。下表都以照片甲为对象，采用同一名录和地区范围；前四行是独立情形，末行表示第二行经过核验后的状态。执行成功可以与内容未知同时成立。破折号表示动作字段为空，末行保留原作答动作，复核依据另行保存。}
  \label{tab:knowledge-acquisition-record-states}
\end{table*}

表中的第四行说明，完成一次条件检查与提交一个物种判断不是同一动作。规则条件不成立时记“未触发”，无需把空答案伪装成弃权；对象若在适用范围外，则根据任务前提记为不适用。核验同样要注明所检查的属性。复核者确认植物学家的弃权符合规范，只能确认动作合规，不能把未知的物种改成已知，也不能把这项检查当成物种正确性的核验。

多份来源作答合在一起后，如果还不能形成可用结论，我们称结果为“未决”，即综合判断尚未确定；这不是单个来源的一次弃权。审计是否开始、是否完成，与“正确、错误、无法裁决”等质量结论也要另行保存，具体处理见第~\ref{chap:knowledge-quality-assessment-correction}章。类似地，“观测不足”“规范歧义”说明没有答案的原因，“不是玫瑰”却已经给出了一项否定判断。如果把这些内容与未知都写成类别0，学习算法就无法分辨它读到的是负标签，还是根本没有任务答案。

例如，把“是否玫瑰”作为一个二分类问题时，标签0可以明确表示“不是玫瑰”，仍然是一项可用于训练的答案。对物种无法判断的照片，应保留未知并说明缺少什么依据，不能仅为方便存储也编码成0。否则，学习算法会把证据不足的照片当成已经确认的负例，原本没有答案的情形就被变成了错误监督。

有了样本标注与状态的区分，再看图~\ref{fig:knowledge-acquisition-observation-to-knowledge}中的三层对象：上层像素是观测，中层答案、关系、约束与目标是候选样本标注，下层才是训练对样本标注的使用。沿四条分支向下看，同一批照片可以支持不同训练要求，样本标注形式不能由文件位置来判断。中层虚线框表示仍待核验：能够自动构造恢复目标，只说明有了一项可训练的要求，不能据此认定它正确或有助于物种识别。

\begin{figure*}[htbp]
  \centering
  \begin{tikzpicture}[
  x=1mm,y=1mm,
  every node/.style={font=\figurefont\footnotesize,text=BookInk},
  layer label/.style={anchor=west,font=\figurefont\bfseries\footnotesize,text=BookTeal},
  observation/.style={rectangle,draw=FigureInput,line width=0.9pt,fill=FigureInput!9,minimum width=52mm,minimum height=12mm,align=center},
  candidate/.style={rectangle,rounded corners=1.2mm,draw=BookAmber,dashed,line width=0.8pt,fill=BookAmber!7,text width=30mm,minimum height=20mm,align=center,inner sep=2mm},
  use/.style={rectangle,draw=FigureOutput,line width=0.8pt,fill=FigureOutput!6,text width=30mm,minimum height=16mm,align=center,inner sep=2mm},
  flow/.style={knowledge arrow,draw=BookMuted,line width=0.8pt,rounded corners=1.2mm}
]
  \node[layer label] at (0,7) {原始观测};
  \node[observation] (photos) at (88,3) {一组花卉照片\\像素与拍摄条件};

  \draw[BookRule,line width=0.6pt] (0,-9) -- (164,-9);
  \node[layer label] at (0,-18) {候选监督信号};
  \node[candidate] (answer) at (30,-27) {类别答案\\“照片甲是鸢尾”\\{\scriptsize 语义：已给出；验证：候选}};
  \node[candidate] (relation) at (70,-27) {对象关系\\“乙、丙来自同一株”\\{\scriptsize 语义：已给出；验证：候选}};
  \node[candidate] (constraint) at (110,-27) {类别约束\\“玫瑰与鸢尾互斥”\\{\scriptsize 语义：已给出；验证：候选}};
  \node[candidate] (restore) at (150,-27) {恢复目标\\“预测被遮蔽花瓣”\\{\scriptsize 语义：已给出；验证：候选}};

  \draw[flow] (photos.south) -- ++(0,-5) -| (answer.north);
  \draw[flow] (photos.south) -- ++(0,-5) -| (relation.north);
  \draw[flow] (photos.south) -- ++(0,-5) -| (constraint.north);
  \draw[flow] (photos.south) -- ++(0,-5) -| (restore.north);

  \draw[BookRule,line width=0.6pt] (0,-40) -- (164,-40);
  \node[layer label] at (0,-49) {训练使用};
  \node[use] (class loss) at (30,-57) {监督分类损失\\学习物种答案};
  \node[use] (relation loss) at (70,-57) {关系一致性目标\\拉近同一对象表示};
  \node[use] (constrained prediction) at (110,-57) {受约束预测\\排除不合法组合};
  \node[use] (reconstruction) at (150,-57) {内容恢复目标\\学习图像表示};

  \draw[flow,draw=BookTeal] (answer.south) -- (class loss.north);
  \draw[flow,draw=BookTeal] (relation.south) -- (relation loss.north);
  \draw[flow,draw=BookTeal] (constraint.south) -- (constrained prediction.north);
  \draw[flow,draw=BookTeal] (restore.south) -- (reconstruction.north);
\end{tikzpicture}

  \caption{观测、候选样本标注与训练使用的三个层次。由同一批花卉照片得到的四类样本标注，分别连接到分类、关系一致性、受约束预测和内容恢复；中层分开标明语义状态与验证状态，虚线框表示样本标注仍需核验。向下的连接说明使用方式，不保证识别收益。内容均为教学构造。}
  \label{fig:knowledge-acquisition-observation-to-knowledge}
\end{figure*}

\subsection{从任务要求确定获取缺口}
\label{sec:knowledge-acquisition-gap-diagnosis}

现在可以把“还不够用”落实到具体缺口。如果目标地区连照片都没有，缺少的是原始观测，需要先采集花卉照片；本章重点讨论已有观测缺少什么解释依据、关系、约束或目标，以及现有标注样本的覆盖和质量哪里不足。一个稀有物种没有任何已确认样本标注，是内容缺失；有了样本标注却只有一次拍摄、没有覆盖特殊光照，是覆盖不足；有整图类别却没有花朵位置，是粒度不足。这三种情况要求补充的内容不同，不能靠重复拍摄同一盆常见花一并解决。

这里的覆盖（coverage）考察目标范围中的哪些情形已经有依据，例如哪些物种、光照和拍摄角度已有标注。粒度（granularity）考察一条依据细到什么对象层次，例如整张照片、一朵花的区域，还是某片花瓣的位置。即使100张照片已经覆盖所有目标物种，若每张只有整图类别，仍不能作为花朵位置的直接答案。反过来，给同一盆玫瑰精确标出100个区域，也没有补入稀有鸢尾的依据。覆盖与粒度是两个不同的检查方向。

有些缺口发生在已经取得的答案中：相似物种被混淆，是内容错误；分类名录调整以后仍沿用旧名称，则是时效失配。拍摄设备或目标地区改变，也可能使旧样本标注的适用条件需要重查，不能只复制旧标签便认为更新完成。表~\ref{tab:knowledge-acquisition-gap-types}把五类缺口接到相应获取需求；验收一列说明每种修补需要观察什么，而不是预先指定某类来源最优。这些缺口可以同时出现，例如稀有物种既缺少多种光照下的样本标注，现有的唯一答案又可能错误。

\begin{table*}[htbp]
  \centering
  \small
  \setlength{\tabcolsep}{3pt}
  \begin{tabularx}{\linewidth}{@{}>{\RaggedRight\setlength{\parindent}{0pt}\arraybackslash}p{25mm}>{\RaggedRight\setlength{\parindent}{0pt}\arraybackslash}p{43mm}>{\RaggedRight\setlength{\parindent}{0pt}\arraybackslash}p{45mm}>{\RaggedRight\setlength{\parindent}{0pt}\arraybackslash}X@{}}
    \toprule
    \textbf{不足类型} & \textbf{花卉识别中的表现} & \textbf{需要补充或改变的内容} & \textbf{可观察的验收依据} \\
    \midrule
    内容缺失 & 目标地区的某个稀有物种没有任何已确认样本标注 & 取得该物种的类别依据、适用范围和来源身份与作答条件 & 目标类别出现已核验样本标注，且对象与范围字段完整 \\
    覆盖不足 & 常见物种有大量近重复照片，稀有物种或特殊光照几乎未覆盖 & 补入困难物种、环境和外观变化的样本标注，而非重复已有高频情形 & 预先划定的物种与拍摄条件各组达到约定的已标注样本数量要求 \\
    粒度不足 & 只有整图类别，无法回答花朵位置或同一对象关系 & 补入与目标问题匹配的区域、实例或成对关系 & 留出的区域或关系样本能按预先规定的对象与语义复核 \\
    内容错误 & 相似物种被混淆，或自动答案把背景叶片当成花朵 & 复核并修正内容，保留原来源、错误类型和修订轨迹 & 在相同独立参考样本上核对修正前后答案，确认已知错误得到处理 \\
    时效失配 & 分类体系已调整，旧样本标注仍用废弃名称；部署地区的物种分布也已变化 & 更新类别映射、适用范围，或重新获取受影响样本标注 & 旧类别均映射到新体系或明确废弃；目标时间窗口的抽查满足预设标准 \\
    \bottomrule
  \end{tabularx}
  \caption{从监督信号缺口确定获取需求。每行由花卉任务中的不足，连接到应补入或改变的内容及验收依据；缺少稀有物种样本标注与缺少位置不能用同一种补充代替。已标注样本数、复核语义和抽查标准须随目标预先确定，各类缺口并非互斥。}
  \label{tab:knowledge-acquisition-gap-types}
\end{table*}

\subsection{从可读取到可靠、有用}
\label{sec:knowledge-acquisition-acceptance}

补入样本标注后，可以按三个不同的问题验收。先检查记录能否读取：照片标识是否存在，标签是否属于约定的类别，名录版本和来源是否齐全，状态组合是否合法。这称为\term{形式完整性}（formal completeness），即记录满足预先规定的字段与格式要求。一条各字段都填齐的“照片甲是玫瑰”记录仍可能把鸢尾认错，因此格式检查之后还需要内容检查。

检查类别内容时，可以请未参与原标注的专家，依据另行提供的清晰局部照片和同一名录重新判断，再比较两份答案。这称为\term{参考一致性}（reference consistency），即候选内容与指定参考证据是否相符。可靠性（reliability）则要求说明这种检查支持了哪些质量判断：例如独立抽查中有多少类别答案满足约定的正确性标准。参考证据的选择也有限制；若两位专家都只看到同一张模糊照片，或者都沿用了同一个错误名录，答案一致仍可能是共享错误。因而“已核验”总要注明检查的属性、依据和标准，不能解释为所有方面永远正确。

即使一批标签按参考证据正确，也还要检查它们是否帮助当前学习任务。假设已有许多清晰玫瑰照片，新标签全部来自相似的玫瑰照片，可能只增加了重复材料，对稀有鸢尾识别的帮助很小。把交付标注放入既定学习算法，再检查训练后模型在目标任务中的结果，评价的是\term{下游效用}（downstream utility）。这里“下游”指取得训练依据之后的模型学习和实际预测；效用衡量这些依据在该任务中带来的收益。

评价这种收益时，要区分训练样本、验证样本和测试样本。训练样本直接参与参数调整；验证样本用于选择获取来源、筛选方案或训练设置；测试样本在方案确定以后，用于检查面对新对象的结果。三者应按用途隔离，并说明如何抽取。随机拆开照片文件仍可能把同一株花的近重复照片分到两侧，使评价高估对新对象的识别能力；因此需要检查对象是否重叠、拍摄条件是否符合未来任务。这些名字本身不保证样本彼此独立或具有相同分布。

形式完整性、参考一致性和下游效用因而逐步回答“能否读取”“指定证据是否支持”“是否帮助任务”。它们检查的对象不同：第一项检查记录，第二项检查内容与依据的关系，第三项检查使用这些记录训练后的模型。训练损失下降也不能代替后两项，因为模型可能只是更贴近已经写入的错误答案。

这几类检查也适用于类别以外的监督。对象关系要核对两张照片是否确实属于同一株花；恢复目标要核对遮挡位置与原内容是否正确对应，再检查这种训练是否服务于识别。模型给出的\term{置信度}（confidence）是它对某个候选答案的确信程度或所赋的分数，具体含义取决于模型及评分方式。高分可以用于决定优先检查哪些对象，却不能自行证明答案正确；模型也可能对训练中反复出现的错误给出高分。

\subsection{在资源限制下选择获取方案}
\label{sec:knowledge-acquisition-plan-selection}

现在假设人工和已有模型都能为稀有物种补入候选标签，但经费、工时和计算资源有限。选择方案不能只比较生成速度，还要比较最后留下的标注是否可靠、是否覆盖缺失情形，以及用于训练有什么收益。一个获取方案应当写明询问哪些对象、调用哪些来源、怎样取得答案、怎样复核，以及什么内容允许交付。仅写“采用模型标注”还不足以确定完整方案。

对每种方案，先区分产生的全部候选与验收后保留的标注。例如，模型提出1000条标签，其中600条因不适用或证据不足而没有通过交付要求，交付集合只有剩下的400条。获取1000条候选和检查被排除的600条都耗费了资源。若只把调用费用除以1000，就无法说明取得一条可交付标注的完整成本。

为了得到可比较的效用，需要固定学习算法、评价样本和比较基线。这里的基线是约定的参照方案，例如使用已有标签训练的模型。对花卉分类，可用留出准确率相对基线增加了多少来计效用；从80\%提高到85\%，增加的是5个百分点。也可以使用平均损失的相反数，让损失降低对应效用升高。可靠性要明确检查类别正确性、关系正确性还是目标构造是否合规，覆盖则要预先划定物种与拍摄条件组，以及每组需要多少标注。不同任务可以采用不同度量，但比较同一任务的方案时应保持口径一致。

一个教学比较可以说明为什么要把可靠性和覆盖单列。假设两个方案均满足预算，交付标准要求独立抽查的标签正确率至少为95\%，四个预设物种组每组至少有20条合格标注。方案甲的抽查正确率为98\%，训练后的留出准确率为88\%，却遗漏了一个稀有物种组；方案乙的抽查正确率为96\%，留出准确率为85\%，四组均达到数量要求。即使甲的总体准确率更高，它仍未完成指定覆盖要求，应先排除，再比较达标方案。这些数值只演示选择规则；有限抽查存在统计不确定性，实际方案还需规定抽查方式及接受标准。

这些量确定后，我们选择预算以内、可靠性和覆盖均达标的方案，再按任务效用从高到低比较，因此目标使用最大化。可靠性与覆盖单独设为约束，是为了避免总体效用较高的方案仍遗漏稀有物种，或交付不可接受的错误；它们不能仅靠效用平均值来代替。下面用统一符号表示这一选择问题。可行方案（feasible plan）在这里指人员、模型、规则和材料允许执行的完整安排，预算和交付门槛再进一步限制其中可以选用哪些。

\begin{definition}[监督信号获取目标]
  \label{def:knowledge-acquisition-objective}
  给定目标任务$T$、可行获取方案集合$\mathcal{A}$和资源上限$\bm{B}$。每个$a\in\mathcal{A}$指定来源、对象、获取动作与复核方式，产生候选样本标注集合$\mathcal{K}_{\mathrm{cand}}(a)$及按交付要求选出的子集$\mathcal{K}(a)\subseteq\mathcal{K}_{\mathrm{cand}}(a)$。

  设$U_T(\mathcal{K}(a))$为交付样本标注的任务效用，$R(\mathcal{K}(a))$为其可靠性度量，$G(\mathcal{K}(a))$为其目标范围覆盖度量，三者均按数值越大越好的方向定义；$r_0$与$g_0$分别为最低可靠性和覆盖要求。获取目标写为
  \begin{equation}
    \begin{aligned}
      \max_{a\in\mathcal{A}}\quad
      & U_T\bigl(\mathcal{K}(a)\bigr) \\
      \text{s.t.}\quad
      & \bm{C}(a)\preceq\bm{B}, \\
      & R\bigl(\mathcal{K}(a)\bigr)\geq r_0,\qquad
        G\bigl(\mathcal{K}(a)\bigr)\geq g_0,
    \end{aligned}
    \label{eq:knowledge-acquisition-objective}
  \end{equation}
  其中$\bm{C}(a)=(C_{\mathrm{dev}}(a),C_{\mathrm{acq}}(a),C_{\mathrm{review}}(a),C_{\mathrm{maint}}(a))$依次计量准备开发、逐项获取、复核修正和持续维护的成本，$\preceq$表示各项都不超过$\bm{B}$中的对应上限。成本按完整方案核算，包括后来被拒绝的候选所耗费的获取与复核投入。人工工时、调用费用与计算资源应分别报告；同一成本项涉及不同单位时也须列明，除非已有可审计的换算规则，否则不能直接相加。
\end{definition}

式~\eqref{eq:knowledge-acquisition-objective}因而把“获得多少条”改为“交付样本标注在任务中有什么收益、是否达标、耗费多少”。这个目标说明如何比较方案，不保证存在达标方案或最大值能够取到，也不提供寻找最优方案的算法；实际效用和可靠性还需要通过试点、抽查与留出评价估计。固定预算时，可以比较达标方案取得的效用，同时报告覆盖和可靠性；固定目标效果时，则反过来比较达到相同要求需要多少分项投入。

读这个式子时，可以先检查三条限制，再看第一行的最大化。$a$代表完整方案，$\mathcal{A}$是当前人员、模型、规则和材料条件下允许考虑的方案集合；$\mathcal{K}_{\mathrm{cand}}(a)$保存全部候选，$\mathcal{K}(a)$保存按要求保留的部分。$R$与$G$分别检查可靠性和覆盖是否达到$r_0$与$g_0$；只有通过这些检查并满足预算的方案，才进入效用$U_T$的比较。式中的$\text{s.t.}$表示“满足以下约束”。这个选择顺序也说明，效用再高都不能补偿一个必须满足却未满足的条件。

成本向量用来保留不同投入，而不是把不同单位强行合成一个数。例如，开发规则需要20工时、调用模型需要100元，20与100没有可以直接相加的共同单位；复核还可能同时耗费工时和调用费用。$\bm{C}(a)\preceq\bm{B}$表示每一项都分别不超过相应上限，工时有剩余不能自动抵消调用费用超限。准备开发、逐项获取、复核修正和持续维护的分工，则帮助我们看出成本发生在一次性投入还是每轮重复操作中。

批量获取之前的小规模试点，用来确认上述比较能否实际执行：目标样本范围是否清楚，哪些未知和错误可以接受，困难物种或条件是否优先覆盖，以及达到什么条件应当停止。当前批次验收后可以交付，但这只完成了当前范围内的任务。分类体系和样本分布继续变化时，长期方案还要规定何时重新抽查、更新或废弃旧样本标注，并把这部分投入计入持续维护成本。

\section{知识获取的研究内容}
\label{sec:knowledge-acquisition-research-scope}

假设试点已经确认稀有物种缺少类别依据，也缺少不同角度下的对象关系，下一步就要决定从哪里补入这些内容。一条样本标注的判断依据来自哪里，称为它的\term{监督来源}（supervision source）。我们可以询问人工、执行规则、让已有模型预测，或利用照片本身构造目标；取得标注以后，再判断哪些能共同使用、哪些需要修正。来源负责提出内容，融合负责组织多份内容或已有模型，质量检查负责判断是否满足要求。把这几种职责分开，就能看清一个完整方案究竟改变了哪一步。

\subsection{按判断依据选择监督来源}
\label{sec:knowledge-acquisition-source-selection}

人工来源让人根据材料和任务规范作出判断。例如，植物学家观察照片中的花瓣、叶片和生长位置，查阅名录后提交“鸢尾”及依据；必要时也可以要求补拍局部，或者因证据不足而弃权。取得答案依赖人的观察与专业判断，因而需要安排可回答的问题、所见材料和人员条件。人工有机会处理规则尚未写明的边界情形，但不会因为作答者是人就自动正确，逐例判断与复核也持续耗费时间。

规则来源把已经规定的条件反复应用到对象。假设照片附有采集编号，另有一份经确认的“采集编号到物种名称”名录，一条规则可以按编号查表，匹配成功时给出候选类别，找不到时不给类别判断。它重复执行的是已写明的关联条件，而非每次重新作植物学判断。是否误配编号、名录是否过期、条件是否适用于当前照片，都影响结果。规则还可以给出前节的类别互斥约束；它提供哪种内容，取决于所编码的条件和结论。

模型来源把已经学得的预测能力用于新对象。例如，已有视觉模型读取照片，计算各物种的分数，再提出分数最高的类别作为候选标签；它也可能给出候选花朵区域。这里的依据是已有模型从过去训练中学得的规律，不是这张照片上已经核验的物种答案。把这类预测转成新训练样本的目标标签，称为\term{伪标签}（pseudo-label），即由模型提出、用作标签的预测结果。它可以减少逐例人工作答，但若过去的训练只覆盖清晰正面照片，新照片中的遮挡仍可能导致高置信错误。

样本本身也可以提供依据。前节把原图像块保留为恢复目标，不需要先取得物种标签；若拍摄日志明确记录照片乙、丙来自同一株花，这个关系也能用于要求两组表示接近。使用的是原内容或已有关系，而非另一个模型对物种的猜测。不过，构造所依赖的关系仍要成立。随意把外观相似的两朵花配成同一对象，可能把不同物种拉近；为构造不同视图而裁掉全部花朵，也可能只留下背景。目标容易构造与目标适合识别，是两项不同判断。

四类来源按内容主要通过什么机制取得来区分。同一个“鸢尾”标签可以由人工、查表规则或视觉模型提供；一段程序既可能执行规则，也可能遮挡图片来构造自监督目标，不能只凭“用了程序”判定来源。构造过程调用模型时，还应说明模型是提出任务答案，还是仅帮助准备材料。表~\ref{tab:knowledge-acquisition-chapter-map}把这些判断接到后续章节，表~\ref{tab:knowledge-acquisition-source-conditions}再把准备投入、易错条件与检查办法放到同一任务下比较。

\begin{table*}[htbp]
  \centering
  \small
  \setlength{\tabcolsep}{4pt}
  \begin{tabularx}{\linewidth}{@{}>{\RaggedRight\setlength{\parindent}{0pt}\arraybackslash}p{35mm}>{\RaggedRight\setlength{\parindent}{0pt}\arraybackslash}X>{\RaggedRight\setlength{\parindent}{0pt}\arraybackslash}p{22mm}@{}}
    \toprule
    \textbf{研究内容} & \textbf{核心问题与主要内容} & \textbf{后续章节} \\
    \midrule
    从人工获取监督信号 & 怎样把需求变成可回答的任务，并选择交互方式、样本、标注者与投入 & 第~\ref{chap:knowledge-acquisition-human}章 \\
    从规则获取监督信号 & 怎样把模式、名录关联和约束写成可重复执行的条件，并检查候选判断 & 第~\ref{chap:rule-based-knowledge-acquisition}章 \\
    从模型获取监督信号 & 怎样准备已有模型、组织输入和请求、取得候选答案，再根据检查结果调整获取过程 & 第~\ref{chap:knowledge-acquisition-model-based}章 \\
    从样本本身获取监督信号 & 怎样用原内容、已有关系或分组依据构造训练目标，并检查构造条件 & 第~\ref{chap:data-derived-knowledge-acquisition}章 \\
    融合多来源标注与模型 & 怎样使多份标注指向相同对象和含义，再综合使用；已有模型又怎样共同预测、组合计算或提供新模型的训练依据 & 第~\ref{chap:knowledge-fusion}章 \\
    评估与修正监督信号 & 怎样建立判断依据、估计质量、定位问题，并验证修正是否有效 & 第~\ref{chap:knowledge-quality-assessment-correction}章 \\
    \bottomrule
  \end{tabularx}
  \caption{由获取需求查找后续章节。前四行回答不同来源怎样产生候选监督信号，后两行回答样本标注怎样共同使用、检查和修正。来源与后续处理承担不同职责，不是六种互斥算法，也不要求每个方案全部采用。}
  \label{tab:knowledge-acquisition-chapter-map}
\end{table*}

\begin{table*}[htbp]
  \centering
  \small
  \setlength{\tabcolsep}{3pt}
  \begin{tabularx}{\linewidth}{@{}>{\RaggedRight\setlength{\parindent}{0pt}\arraybackslash}p{15mm}>{\RaggedRight\setlength{\parindent}{0pt}\arraybackslash}p{37mm}>{\RaggedRight\setlength{\parindent}{0pt}\arraybackslash}p{28mm}>{\RaggedRight\setlength{\parindent}{0pt}\arraybackslash}p{43mm}>{\RaggedRight\setlength{\parindent}{0pt}\arraybackslash}X@{}}
    \toprule
    \textbf{来源} & \textbf{准备与投入} & \textbf{典型标注结果} & \textbf{限制与易错处} & \textbf{核验入口} \\
    \midrule
    人工 & 规范、界面与培训；逐项判断和复核 & 类别、区域、解释或弃权 & 受能力与可见证据限制；歧义、疲劳和偏好可致错 & 独立专家、预先确认的参考答案与留出抽查 \\
    规则 & 名录、编号映射与适用条件；开发和版本维护 & 类别、关联与约束 & 编号误配、未收录对象和名录过期 & 错配反例、名录版本核对与覆盖抽查 \\
    模型 & 模型与预测请求的组织；调用、筛选和更新 & 类别、区域、关系或生成材料 & 依赖能力与分布；高置信错误和来源偏差可被复制 & 独立标注、可验证规则与任务留出评价样本集 \\
    样本本身 & 观测与构造依据；样本处理、构造变换和训练 & 恢复目标、视图关系与分组 & 依赖保留的结构；错误配对、改变物种含义的变换，或依赖背景完成训练 & 检查构造、移除相应目标作对照、独立下游评价 \\
    \bottomrule
  \end{tabularx}
  \caption{在同一花卉任务下选择监督来源。各行并列列出准备与逐项投入、典型标注结果、容易出错的条件和核验办法；依据这些条件选择来源，而不能仅凭标注结果形式或自动化程度判断。弃权是人工来源的动作，不是额外的类别答案。}
  \label{tab:knowledge-acquisition-source-conditions}
\end{table*}

\subsection{多来源结果的对齐与融合}
\label{sec:knowledge-acquisition-alignment-fusion}

单独比较来源以后，可以把互补的办法放进同一个稀有花卉方案。植物学家先为少量代表照片作答并写出规范，规则按采集名录取得更多候选类别，视觉模型为尚无名录答案的照片提出预测，同一对象的多角度照片则构造关系目标。它们分别补充已有缺口，因而可以同时采用，而不是必须挑一个来源处理全部问题。

共同使用以前，先确认不同结果回答的是同一个问题。人工对照片甲中的左侧花朵作答，模型对整张照片作答，两份结果的对象不同；人工用“鸢尾”指一个具体物种，规则用它指一组近似类别，两份结果的含义也不同。对象对齐（object alignment）把结果连到同一对象，语义对齐（semantic alignment）确认类别名称、任务和适用范围采用同一含义。两种对齐完成之后，才能把不同来源的“鸢尾”作为可比较的票。

多份可比较判断仍可能冲突，因此需要融合（fusion），即按照明确办法把多来源信息共同用于形成结果。一个最小办法是对同一对象的类别票按多数决定类别；更复杂的办法还会估计来源在不同条件下的可靠性。多数票也有前提：如果三条规则都从同一份错误名录复制答案，它们的三票只是同一错误的重复，不能当作三份独立支持。恢复目标和关系目标则没有与物种类别相同的回答空间，可以分别进入训练，无须融合成一个物种名称。

已有模型还可以作为整体共同使用，例如让几个模型各自预测后再组合输出，或者由一个模型向另一个模型提供训练目标。这与综合同一照片的标签票是不同层面的操作。第~\ref{chap:knowledge-fusion}章会分别处理多来源标注与已有模型的联合使用；在当前获取方案中，需要先说明交付的是标注内容还是模型，以及后续怎样使用它。

\subsection{获取、检查与再获取}
\label{sec:knowledge-acquisition-feedback}

把这些操作组织起来，一次性方案可以从试点进入批量获取，再按需要对齐融合、抽查修正并验收交付。如果抽查发现稀有物种在遮挡下经常被误判，持续方案还要把这个结果送回下一轮，改选样本、追加人工判断或调整来源与预算。每轮维护日志应说明改变了什么，哪些旧结果的条件仍成立、可以保留，哪些对象必须重新获取，以及何时停止。只是再次调用同一个模型生成同类答案，没有补入独立标注依据时，增加轮次并不能作为可用监督信号增加的证据。

例如，试点发现大多数错误来自看不见花瓣的照片，下一轮可以只对这些对象补拍局部，再交给专家判断；若错误来自旧名录中的名称映射，则应更新映射并重查受影响的记录。前一种处理增加可见证据，后一种处理改变答案解释依据。检查结果之所以形成反馈（feedback），是因为它决定了下一轮改什么，而不仅是把同样的获取步骤再执行一次。

图~\ref{fig:knowledge-acquisition-workflow}把样本标注的去向与方案的调整分开画出。四类来源都能送入候选集合，实线监督信号路径随后分叉：多份可比较结果需要对齐融合，单一来源或独立训练目标可以直接检查。虚线反馈则从质量判断返回来源选择或获取机制，表示检查结果改变下一轮怎样做。沿这条反馈看，发现遮挡下的错误可能只要求重做特定对象，而不是把全部步骤无条件重跑。

\begin{figure*}[htbp]
  \centering
  \begin{tikzpicture}[
  x=1mm,y=1mm,
  every node/.style={font=\figurefont\footnotesize,text=BookInk},
  stage/.style={rectangle,draw=BookRule,line width=0.8pt,fill=BookPaper,text width=24mm,minimum height=13mm,align=center,inner sep=2mm},
  source/.style={rectangle,rounded corners=1.2mm,draw=FigureInput,line width=0.8pt,fill=FigureInput!8,text width=14mm,minimum height=10mm,align=center,inner sep=1.5mm},
  fusion/.style={stage,draw=FigureModel,fill=FigureModel!8,text width=20mm},
  check/.style={stage,draw=BookAmber,fill=BookAmber!8,text width=21mm},
  result/.style={stage,draw=FigureOutput,fill=FigureOutput!6,text width=15mm},
  flow/.style={knowledge arrow,draw=BookTeal,line width=0.9pt,rounded corners=1.2mm},
  control/.style={knowledge arrow,draw=BookMuted,line width=0.8pt,rounded corners=1.2mm},
  bus/.style={draw=BookMuted,line width=0.8pt},
  knowledge bus/.style={draw=BookTeal,line width=0.9pt},
  feedback bus/.style={draw=BookAmber,dashed,line width=0.8pt,rounded corners=1.2mm},
  adjust/.style={knowledge arrow,draw=BookAmber,dashed,line width=0.8pt,rounded corners=1.2mm},
  edge label/.style={font=\figurefont\scriptsize,inner sep=1pt,text=BookMuted}
]
  \node[stage] (need) at (15,12) {需求与资源\\范围、效果、预算};
  \node[stage] (select) at (47,12) {来源选择\\对象、机制、投入};

  \node[source] (human) at (83,36) {人工};
  \node[source] (rule) at (83,20) {规则};
  \node[source] (model) at (83,4) {模型};
  \node[source] (data) at (83,-12) {样本构造};
  \node[font=\figurefont\scriptsize\bfseries,text=BookTeal] at (83,46) {执行选定的获取机制};

  \node[stage,draw=FigureModel,fill=FigureModel!8] (candidate) at (126,12) {候选监督信号\\含来源与状态};
  \node[fusion,text width=24mm] (align) at (83,-46) {对象与语义对齐\\需要时再融合};
  \node[check] (quality) at (126,-46) {质量判断\\覆盖与可靠性\\下游效用};
  \node[result] (usable) at (159,-46) {可用结果\\或停止};

  \draw[control] (need.east) -- (select.west);
  \draw[control] (select.east) -- (67,12);
  \draw[bus] (67,-12) -- (67,36);
  \draw[control] (67,36) -- (human.west);
  \draw[control] (67,20) -- (rule.west);
  \draw[control] (67,4) -- (model.west);
  \draw[control] (67,-12) -- (data.west);

  \draw[knowledge bus] (100,-12) -- (100,36);
  \draw[flow] (human.east) -- (100,36);
  \draw[flow] (rule.east) -- (100,20);
  \draw[flow] (model.east) -- (100,4);
  \draw[flow] (data.east) -- (100,-12);
  \draw[flow] (100,12) -- (candidate.west);

  \draw[flow] ([xshift=-8mm]candidate.south) -- (118,-25) -| (align.north);
  \node[edge label,anchor=south] at (99,-24) {多份可比结果};
  \draw[flow] (align.east) -- (quality.west);
  \draw[flow] (candidate.south) -- (quality.north);
  \node[edge label,anchor=west,align=left] at (129,-16) {单一来源\\或独立目标};
  \draw[flow] (quality.east) -- (usable.west);

  \coordinate (feedback split) at (100,-64);
  \draw[feedback bus] (quality.south) |- (feedback split);
  \draw[adjust] (feedback split) -- (100,-72) -| (select.south);
  \node[edge label,text=BookAmber,anchor=south] at (70,-71) {调整来源选择};
  \draw[adjust] (feedback split) -- (64,-64) -- (64,-23) -- (67,-23) -- (67,-12);
  \node[edge label,text=BookAmber,anchor=south] at (82,-63) {调整获取机制};

  \draw[flow] (4,-82) -- (16,-82);
  \node[anchor=west,font=\figurefont\scriptsize] at (18,-82) {监督信号流动};
  \draw[adjust] (57,-82) -- (69,-82);
  \node[anchor=west,font=\figurefont\scriptsize] at (71,-82) {质量反馈};
\end{tikzpicture}

  \caption{获取、融合与质量改进的一种组织方式。实线监督信号路径由四类来源汇入候选集合，按是否需要综合判断选择融合或直接检查；虚线表示依据覆盖、可靠性与效用调整来源选择和获取机制。需求与来源选择之间的连接组织执行，不传递新的标签或关系。图不要求每次启用全部来源或全部环节。}
  \label{fig:knowledge-acquisition-workflow}
\end{figure*}

比较两个完整方案时，固定的不只是花卉分类这个任务名称，还包括目标地区与时间范围、评价样本、训练设置和成本核算周期。覆盖、可靠性、下游指标与分项投入需要同时报告，准备和持续成本还要按来源列明。如果根据同一评价样本集反复选样、调整规则或改写模型提示，这组评价样本就已经参与方案选择，不能再当作独立的最终效果证据，否则评价可能偏向被选中的方案。跨来源预算协调见第~\ref{chap:knowledge-fusion}章，人工内部的选样、选人与重复标注分配见第~\ref{chap:knowledge-acquisition-human}章；日志、存储与调度的系统实现不在本部分展开。

\section{知识获取的研究历史}
\label{sec:knowledge-acquisition-history}

前节的方案既可以逐例取得答案，也可以把专家判断写成规则反复执行。历史中同样能看到这两种安排。Rosenblatt在1958年的感知机研究中讨论用带答案样本调整连接；感知机根据输入特征的加权结果给出类别，样本答案用于判断是否需要调整权重。1984年出版的MYCIN项目回顾则说明怎样把领域专家的判断依据编码为条件规则，并检查、维护规则库。前者让答案随样本进入学习，后者把判断依据写成可重复执行的规则。两者为理解获取方式提供早期例证，并不是监督学习或专家系统的唯一起源\scite{rosenblatt1958,buchanan1984rulebased}

逐例提供答案需要持续投入人工，规则库需要维护覆盖范围与例外。后续研究从这些不同限制出发：有限答案应当用在哪些对象上，多份不一致的答案怎样使用，没有任务标签时还能从哪里取得学习目标，以及模型生成内容后怎样检查和保留它们。本节按这些问题组织代表工作，说明每项工作改变的机制。自动化让人参与的环节发生变化，但界定任务、处理难例和核验结果仍可能需要人工，这些路线也继续并存。

\subsection{有限标签的选择与利用}
\label{sec:knowledge-acquisition-history-label-budget}

当未标注照片很多、预算只允许追加少量人工答案时，我们需要选择询问对象。设现有已标注照片恰好都是“紫色鸢尾”和“红色玫瑰”，两个候选预测办法可能都答对：一个主要依赖颜色，另一个主要依赖花瓣形态。它们都符合现有答案，却可能在一张粉色鸢尾照片上给出不同预测。询问这种照片，有机会知道哪种判断办法需要修正；继续询问两者都答对的近重复照片，未必能提供同样的区分信息。这是一个教学情形，用来说明有限标签为什么可能留下多个不同候选。

Seung、Opper与Sompolinsky在1992年的委员会查询（query by committee）中，利用这种预测分歧选择对象。委员会是一组当前允许考虑的候选假设，每个假设是一种从输入得到类别的预测办法；让它们对未标注样本分别作答，再优先询问分歧大的对象。取得人工答案以后，就能据此更新候选及其判断。查询价值取决于允许的假设范围、委员会能否反映当前不确定性，以及人能否给出有依据的答案。若分歧只是因为照片严重模糊、任何人都无法辨认，花费同样预算未必取得有效区分\scite{seung1992querycommittee}

另一条路线利用已有少量标签和大量未标注样本，减少新答案必须逐个来自人工的程度。Blum与Mitchell在1998年的协同训练（co-training）中，让两个分类器相互提供候选标签。它们观察同一对象的两种特征视图（feature view），即两组不同的可用特征。原工作中的网页例子有助于理解：网页正文是一组描述，指向该网页的链接文字是另一组描述，两者都可能提供网页类别的线索。

先用少量已标注网页分别训练两个分类器，再让它们给未标注网页预测。一个分类器选出的高置信候选标签，可以成为另一个分类器后续训练的依据；两者使用的特征不同，因而有机会互相补入信息。原工作的理论设定要求每个视图都足以预测类别，并要求给定类别后两视图条件独立：在已知类别的条件下，再知道正文视图，不会改变链接视图的条件分布，反之亦然。仅仅把特征分成两组并不能满足这些条件。两视图共享偏差时，可能反复把同一种错误传给对方。委员会查询主要改变向人询问的对象，协同训练改变标签怎样在模型之间扩展；扩展出来的仍是候选答案\scite{blum1998cotraining}

选样还有一个起步问题：标签极少时，新训练的任务模型尚未形成可靠判断，用它的预测分数来选择待标注样本也可能不可靠。这称为冷启动（cold start）困难，即缺少足够初始依据，后续选择机制难以开始有效工作。一个思路是借助已在其他材料上训练的大语言模型，也就是能处理和生成文本的已有大规模模型。2026年发表于TACL的ActiveLLM在文本少样本场景中利用大语言模型选择待人工标注的实例。它提供询问优先级，最终答案仍需人工取得；效果证据限于论文的模型、样本集与选择协议，不能直接推广到所有任务\scite{bayer2026activellm}

\subsection{多来源答案的质量推断}
\label{sec:knowledge-acquisition-history-source-quality}

一张照片交给三个人，两人判断玫瑰、一人判断鸢尾，多数票会选择玫瑰，但票数本身没有说明三人的专业能力和照片难度。某人平时可靠，也可能在遮挡照片上出错；另一人的少数意见则可能来自更充分的专业知识。Whitehill等人在2009年的模型中，联合推断图像标签、标注者能力和样本难度：用同一个概率模型描述不同能力的人在不同难度样本上怎样产生答案，再依据多次作答估计这些未知量。这样，来源与样本条件进入了综合判断，票的作用不必始终相同。推断仍依赖模型设定和作答信息，不能仅凭某人总与多数一致就证明他正确\scite{whitehill2009crowdlabels}

专家还可以把经验写成程序，让它反复提交候选票。Ratner等人在2017年发布的Snorkel用标注函数（labeling function）实现这一安排：程序检查指定条件，满足时返回候选类别，没有类别判断时返回不投票编码。例如，前节按采集编号查名录的规则就可以写成这样的函数。多条函数可能覆盖不同照片，也可能在同一照片上冲突，因而还需要估计哪些函数较可靠、哪些票共享依据。

Snorkel使用生成式来源模型（generative source model）描述这些票怎样产生。这里建模的对象是来源输出：假定存在尚未确定的真实类别，再用来源准确性和相关结构解释为什么会看到这一组类别票。观察多条函数在许多对象上的同意与分歧，帮助估计来源性质，随后形成用于训练的综合标签。它减少了逐例人工作答，却把一部分工作转为编写函数、规定来源模型和检查其条件。若多条函数复制同一错误，或者观测不足以区分来源质量与真实类别，综合结果仍可能不可靠。未触发与主动弃权都可能被程序编码为“不投票”，原始触发结果与来源动作仍需按前节约定分别保存\scite{ratner2017}

\subsection{从原内容构造学习目标}
\label{sec:knowledge-acquisition-history-internal-target}

如果观测很多，却没有目标任务的类别答案，仍可以利用原内容构造学习要求。前节以图像遮挡说明了这一点，文本也有相应机制。设原句是“花瓣呈紫色”，我们遮住“紫色”，让模型根据其余内容预测原来被遮住的单位；在这个教学例子中，原句自身就提供了一个可检查的答案，无需另外询问句子的分类标签。文本先切分成词元（token），即模型读取和预测的基本单位，实际词元可以是词，也可以是词的一部分。

BERT采用的遮蔽语言模型（masked language model）目标，要求根据上下文预测选中位置的原词元。它通过大量这样的预测学习文本表示。先用一般学习目标训练模型，作为后续具体任务的起点，称为\term{预训练}（pre-training）。取得的直接依据是原词元恢复目标，训练的结果是模型及其表示；原文本并没有因此自动得到情感类别或其他下游标签。BERT于2018年发布预印本、2019年发表于NAACL，在这里代表内容恢复路线，并不是自监督学习的唯一起点\scite{devlin2019}

从一般目标学得的表示到具体任务输出，还缺少一层联系。一个模型能根据上下文恢复词元，不等于它已经知道当前情感任务有哪些类别；同样，图像模型能恢复花瓣，也不等于它已经会输出本地区的物种名称。任务适配（task adaptation）就是把已有模型接到具体任务的输入与输出要求上，例如利用已有表示，再用物种标签训练分类输出。若继续调整已有模型的参数，这种后续训练称为\term{微调}（fine-tuning）。预训练提供起点，适配确定怎样用于当前任务，任务评价再检查结果。

Bommasani等人在2021年的\term{基础模型}（foundation model）报告中，把在大规模且来源多样的训练样本上训练、随后可适配多种下游任务的模型作为共同研究对象。一个已训练模型可以服务多种任务，准备投入因而有机会被复用；报告也强调下游系统会继承上游模型的缺陷。回到花卉例子，构造恢复目标、学得适合识别的表示、取得正确物种答案是三种不同结果，前一种成立不能代替后两种的检查\scite{bommasani2021foundation}

\subsection{模型生成材料与递归使用}
\label{sec:knowledge-acquisition-history-generated-material}

模型来源还可以从“给已有样本作答”扩展到“连问题和材料一起生成”。Self-Instruct让语言模型生成指令、输入材料和示范答案，再过滤无效或过于相似的项目，用于后续微调。指令（instruction）说明要执行什么任务，输入材料说明这一次对什么执行，示范答案则给出训练要模仿的输出。例如，“提取句中的花卉名称”是任务指令，“花坛里种着鸢尾”是这次输入，“鸢尾”是示范答案。三者共同组成一项训练实例，任何一项不匹配都可能产生错误监督。这个例子只说明组成方式，不代表论文使用的具体实例。

相比预测一张已有照片的物种，这种安排扩大了可以构造的训练材料，也扩大了需要检查的对象：指令是否明确、输入是否包含足够依据、答案是否符合指令和材料，以及新任务是否只是在重复少数模式。过滤能排除指定的不合格项，却不能保证留下的答案事实正确或任务分布完整。Self-Instruct提供的是生成与筛选训练实例的办法，生成数量不能单独作为可靠知识增加的证据\scite{wang2023selfinstruct}

生成候选以后，还可以按明确标准筛选。DeepSeek-R1的2025年原始报告（arXiv v1）组合了模型生成、拒绝采样、可验证任务的规则奖励与一般场景的奖励模型。这里的拒绝采样（rejection sampling）指报告所用的候选筛选环节：先生成，再拒绝不合格项、保留合格项。奖励（reward）则为训练结果提供评分反馈；规则奖励按可执行标准评分，奖励模型（reward model）通过另一个模型计算分数。以可核对最终答案的数学题为例，检查规则能判断最终答案是否符合标准，却不自动证明每一步解释都正确。一般任务中，模型评分也可能继承评分模型的偏差。获取、筛选和反馈可以相连，但每种检查都要说明它能覆盖哪些属性；本章不展开强化学习的训练算法\scite{deepseek2025r1}

模型提供标签时，还会遇到来源能力与学习者能力不一致的问题。弱到强泛化（weak-to-strong generalization）研究把弱模型生成的标签用于训练更强模型，考察学习者能否取得超过标签来源的任务结果。这里的强弱指同一任务设置下相对的模型能力。强学习者已有的训练基础可能比弱监督者更丰富，弱标签只是后续指导它的一部分信息；它并不是从零开始、只能照抄这批标签。不过，弱标签中的系统错误仍可能被学习者继承。该工作于2023年公开预印本、2024年发表于ICML，在若干任务上观察到学习者超过弱监督者，但仍未恢复强模型的全部能力。结果支持特定设置下弱监督可能有价值，不能作为弱来源普遍足以监督更强系统的保证\scite{burns2023weaktostrong}

使用一次模型生成的训练实例，与反复依赖前代模型生成材料，是不同的训练安排。合成训练样本（synthetic training sample）指由模型等生成机制构造、用于训练的实例。递归使用时，第1代模型产生第2代的训练材料，第2代再产生第3代的材料，后续观测越来越依赖前代模型保留了什么。若生成器很少产生稀有鸢尾，下一代就可能更少见到这类样本；每轮总条数不变，也可能逐渐只剩常见玫瑰。这是用于说明覆盖退化的花卉情形，不是论文中的花卉实验。

Shumailov等人在递归合成训练的理论与实验设置中研究\term{模型坍塌}（model collapse）：原分布中少见的情形逐代丢失，后代模型越来越难保留这些尾部内容。结论依赖具体的递归训练机制，不能推出任意一次使用合成样本都会退化。它要求我们在方案中记录来源谱系（source lineage），即样本由哪个模型、依据哪一代材料生成，以及独立于这条生成链的原始观测是否仍被保留。只统计每代条数，无法检查稀有内容是否已经丢失\scite{shumailov2024collapse}

递归训练还要区分替换与累计。替换会丢弃旧训练集、只保留新合成样本；累计则保留最初从真实对象取得的样本，并继续加入历代合成样本。以教学计数为例，最初有1000张真实照片，每代生成200张，第5轮累计方案保留$1000+5\times200=2000$张，替换方案只保留最近200张。累计至少没有按这项规则主动丢弃最初的照片，同时增加了存储与重训投入；这个条数计算没有证明识别能力更好。

Gerstgrasser等比较这两种安排，其线性回归分析和有限实验表明，累计可以避免所考察替换机制中的误差持续累积，但结论依赖抽样、噪声与模型条件。筛选办法、旧样本权重或目标分布改变以后，还须重新检查稀有内容的覆盖。保留真实样本是方案的一个机制条件，不能被简化为“只要累计就能保证任意递归训练有效”\scite{ka-gerstgrasser2024-accumulation}

这些问题可以再放回年份顺序中对照。表~\ref{tab:knowledge-acquisition-history-works}按所引版本的年份排列代表工作，用相同栏目比较研究问题、获取机制、新增依据与验证边界。它还列出后续主讲章节，便于从当前问题继续阅读。各行的任务和证据不同，不能合并性能数字作统一排名；发表较晚也不说明在当前花卉任务上更适合。

\begin{table*}[p]
  \centering
  \footnotesize
  \setlength{\tabcolsep}{2.5pt}
  \begin{tabularx}{\linewidth}{@{}>{\RaggedRight\setlength{\parindent}{0pt}\arraybackslash}p{28mm}>{\RaggedRight\setlength{\parindent}{0pt}\arraybackslash}p{29mm}>{\RaggedRight\setlength{\parindent}{0pt}\arraybackslash}p{34mm}>{\RaggedRight\setlength{\parindent}{0pt}\arraybackslash}p{34mm}>{\RaggedRight\setlength{\parindent}{0pt}\arraybackslash}X@{}}
    \toprule
    \textbf{代表工作} & \textbf{研究问题} & \textbf{来源或获取机制} & \textbf{新增标签、关系或训练目标} & \textbf{主要假设、验证边界与后续入口} \\
    \midrule
    Rosenblatt感知机（1958） & 怎样依据带答案样本调整模型参数 & 人工答案随样本进入模型参数更新 & 把观测与类别答案组成可学习样本 & 只代表早期监督学习实例；模型准备见第~\ref{chap:knowledge-acquisition-model-based}章 \\
    MYCIN回顾（1984） & 怎样获取并维护专家判断依据 & 访谈专家并把专家判断编码为规则库 & 可重复执行、检查和修订的规则说明 & 规则覆盖与例外依赖领域维护；见第~\ref{chap:rule-based-knowledge-acquisition}章 \\
    委员会查询（1992） & 哪些未标注样本值得询问 & 候选假设的预测分歧指导人工查询 & 把有限人工答案集中到当前模型难以确定的样本 & 依赖假设空间与委员会；人工选择见第~\ref{chap:knowledge-acquisition-human}章 \\
    协同训练（1998） & 怎样利用少量标签和大量未标注样本 & 两种特征视图的分类器相互提供候选标签 & 扩展可训练的样本标签 & 依赖视图充分性与给定类别后的条件独立性；样本关系见第~\ref{chap:data-derived-knowledge-acquisition}章 \\
    Whitehill等（2009） & 怎样融合能力不同的标注者 & 联合建模标签、标注者能力和样本难度 & 估计综合标签及来源可靠性 & 观测能否区分真值与来源质量、模型设定是否合适，均需检验；人员与融合见第~\ref{chap:knowledge-acquisition-human}章和第~\ref{chap:knowledge-fusion}章 \\
    Snorkel（2017） & 怎样反复执行并组合专家启发式 & 返回类别票或不投票编码的标注函数与来源模型 & 批量候选标签及来源准确性、依赖估计 & 规则覆盖与相关结构可能设错；规则与融合见第~\ref{chap:rule-based-knowledge-acquisition}章和第~\ref{chap:knowledge-fusion}章 \\
    BERT（2018/2019） & 怎样让无下游标签文本参与预训练 & 遮蔽观测并预测原词元 & 内容恢复目标和预训练表示 & 内部目标不等于任务答案；样本构造见第~\ref{chap:data-derived-knowledge-acquisition}章 \\
    基础模型报告（2021） & 怎样理解大规模预训练与任务适配 & 在多种来源的训练样本上训练后适配多种任务 & 对可复用表示与共同问题的研究框架 & 下游会继承上游缺陷；模型来源见第~\ref{chap:knowledge-acquisition-model-based}章 \\
    Self-Instruct（2023） & 怎样扩大指令训练材料 & 模型生成指令、预测请求和示范答案并过滤 & 新的任务描述与训练实例 & 过滤不保证事实正确或分布完整；模型来源见第~\ref{chap:knowledge-acquisition-model-based}章 \\
    弱到强泛化（2023预印本／2024发表） & 弱来源能否支持更强学习者 & 弱模型标签监督强模型 & 特定实验中的弱监督训练信号 & 未恢复全部强模型能力；质量边界见第~\ref{chap:knowledge-quality-assessment-correction}章 \\
    递归合成训练样本研究（2024） & 模型生成的训练样本反复回流会发生什么 & 前代模型生成后代训练样本 & 对覆盖退化条件的证据 & 结论依赖递归合成样本的训练机制；质量与修正见第~\ref{chap:knowledge-quality-assessment-correction}章 \\
    DeepSeek-R1原始报告（2025） & 怎样组合生成样本与筛选反馈 & 模型生成、拒绝采样、规则及奖励模型 & 推理与非推理训练材料 & 验证器和生成模型都有覆盖边界；见第~\ref{chap:knowledge-acquisition-model-based}章和第~\ref{chap:knowledge-quality-assessment-correction}章 \\
    ActiveLLM（2026） & 少样本冷启动时怎样选择待标注实例 & 大语言模型根据任务描述选择样本 & 面向人工查询的样本优先级 & 依赖模型和选择协议；人工选样见第~\ref{chap:knowledge-acquisition-human}章 \\
    \bottomrule
  \end{tabularx}
  \caption{按获取问题横向比较代表研究。各行分别说明机制、新增标签、关系或训练目标及其使用边界，并给出后续入口，不比较跨任务性能。BERT和弱到强泛化同时列预印本与正式发表年份，其余采用所引版本年份；代表工作不等于相应方向的唯一起点。}
  \label{tab:knowledge-acquisition-history-works}
\end{table*}

横向比较以后，图~\ref{fig:knowledge-acquisition-history}再显示这些问题在哪些年份同时发展。四条泳道按选择人工判断、组织多来源与规则、利用样本内在结构、模型生成与核验来归位，同一工作也可能服务于不止一条路线。横向位置保留年份先后，但间距不按时间长度缩放；泳道与连接说明本章怎样组织问题，不证明直接继承或因果关系。人工标注和专家规则放在持续背景中，不为这些实践强行指定一个统一起点。

\begin{figure*}[!t]
  \centering
  \begin{tikzpicture}[
  x=1mm,y=1mm,
  every node/.style={font=\figurefont\scriptsize,text=BookInk},
  lane label/.style={anchor=west,font=\figurefont\bfseries\scriptsize,text width=20mm,align=left},
  timeline/.style={draw=BookRule,line width=0.8pt},
  event/.style={rectangle,rounded corners=1mm,line width=0.8pt,minimum height=9mm,align=center,inner sep=1.2mm},
  selection event/.style={event,draw=FigureInput,fill=FigureInput!8,text width=20mm},
  source event/.style={event,draw=FigureModel,fill=FigureModel!8,text width=20mm},
  data event/.style={event,draw=BookTeal,fill=BookTeal!7,text width=20mm},
  generation event/.style={event,draw=BookAmber,fill=BookAmber!8,text width=20mm},
  stem/.style={draw=BookMuted,line width=0.6pt},
  time arrow/.style={knowledge arrow,draw=BookMuted,line width=0.8pt}
]
  \fill[BookPaper] (0,-82) rectangle (22,43);
  \draw[BookRule,line width=0.6pt] (22,-82) -- (22,43);

  \node[lane label,text=FigureInput] at (1,28) {选择人工判断};
  \node[lane label,text=FigureModel] at (1,9) {组织多来源与规则};
  \node[lane label,text=BookTeal] at (1,-10) {利用样本内在结构};
  \node[lane label,text=BookAmber] at (1,-31) {模型生成与核验};

  \node[anchor=south west,font=\figurefont\scriptsize\bfseries,text=BookTeal] at (24,46) {持续背景：人工逐例标注、专家规则、任务设计与核验};
  \draw[time arrow] (24,42) -- (166,42) node[above left,font=\figurefont\scriptsize] {年份：位置保持先后，间距不按比例};
  \draw[timeline] (24,28) -- (166,28);
  \draw[timeline] (24,9) -- (166,9);
  \draw[timeline] (24,-10) -- (166,-10);
  \draw[timeline] (24,-31) -- (166,-31);

  \node[selection event] (qbc) at (32,28) {1992\\委员会查询};
  \node[selection event,text width=17mm] (activellm) at (159,28) {2026\\ActiveLLM};

  \node[source event,text width=22mm] (whitehill) at (64,9) {2009\\Whitehill等};
  \node[source event] (snorkel) at (88,9) {2017\\Snorkel};

  \node[data event] (cotraining) at (49,-10) {1998\\协同训练};
  \node[data event,text width=22mm] (bert) at (95,-10) {2018/2019\\BERT};
  \node[data event,text width=22mm] (foundation) at (120,-10) {2021\\基础模型报告};

  \foreach \yearx/\eventyear in {131/2023,142/2024,153/2025}{
    \fill[BookAmber] (\yearx,-31) circle (1.2pt);
    \node[anchor=south,font=\figurefont\scriptsize] at (\yearx,-29) {\eventyear};
  }
  \node[generation event] (self instruct) at (102,-43) {2023\\Self-Instruct};
  \node[generation event,text width=28mm] (weak strong) at (102,-66) {2023预印本／2024发表\\弱到强泛化};
  \node[generation event,text width=22mm] (collapse) at (151,-49) {2024\\递归生成退化};
  \node[generation event,text width=22mm] (deepseek) at (151,-75) {2025\\DeepSeek-R1};
  \draw[stem] (131,-31) -- (131,-43) -- (self instruct.east);
  \draw[stem] (131,-43) -- (131,-66) -- (weak strong.east);
  \draw[stem] (142,-31) -- (142,-39) -| (collapse.north);
  \draw[stem] (153,-31) -- (166,-31) |- (deepseek.east);

\end{tikzpicture}

  \caption{获取问题随时间扩展而并存。四条泳道按研究问题归位，横向位置保留年份先后，间距不按比例；下方引线把生成相关工作接到年份点，不表示方法继承。图不表达因果或后来路线取代早期路线，“2018/2019”分别指BERT预印本与正式发表。}
  \label{fig:knowledge-acquisition-history}
\end{figure*}

\FloatBarrier
\section{本章小结}
\label{sec:knowledge-acquisition-summary}

回到开篇的花卉模型，照片提供观测，物种标签提供检查类别预测的答案，损失函数再把预测与答案的差异变成训练依据。没有物种名称时，对象关系仍能要求表示接近，类别约束能排除不合法判断，恢复目标也能从观测中构造学习要求。它们都可以指导训练，但给出的信息不同，恢复得好或满足一个约束，都不能直接代替正确的物种答案。

获取这些内容以后，还要保存对象、来源、内容、适用范围和状态。执行成功只说明操作完成，内容已给出只说明取得判断，按指定依据核验后才有相应的质量结论。未知、弃权和不适用之所以分开，是因为它们会带来补证据、恢复执行、转交或排除等不同处理。这样保存，后续训练就能分辨一条负标签与一条尚无答案的记录。

选择获取方案时，需要从实际缺口出发，再决定来源和处理方式。人工、规则、模型与样本本身提供的依据和易错条件不同，可以互补；可比较答案要先对齐对象与含义，不同训练目标则可以分别使用。记录能否读取、参考证据是否支持、用于训练是否帮助目标任务，应分别检查。在预算、可靠性和覆盖满足要求的方案中，再比较任务效用；候选生成数、置信度或训练损失都不能单独替代这些判断。

代表性研究表明，获取已经从逐例答案和专家规则扩展到查询选择、来源建模、内部目标以及模型生成与筛选，但这些路线并未相互取代。批量处理带来新的获取机会，也需要应对来源共享错误、候选回流和覆盖退化；生成或筛选能力扩大，仍不保证新增样本标注可靠、有用。后续各章据此展开四类来源的具体机制，再处理融合与质量改进。本章确定了需求、样本标注和评价口径，具体算法怎样实现、在什么设置下值得采用，还须沿这些入口继续检验。

\chapter{基于人工的知识获取}
\label{chap:knowledge-acquisition-human}

监督学习要调整模型，先要知道对每一份输入期望得到什么答案。第~\ref{chap:knowledge-acquisition-overview}章把这种指导训练的依据统称为监督信号，其中最直接的一类是人给出的判断：人读懂材料、作出判断，再把判断连同来源和状态保存成样本标注。本章只讨论这类由人提供的依据，研究它怎样取得才可靠、可用。机器也能给出候选答案，但那属于第~\ref{chap:knowledge-acquisition-model-based}章的内容；本章关心的是，当我们决定“请人来判断”以后，怎样让有限的人力用在刀刃上。

先把任务具体化。假设我们已经收集了一批产品使用文档，希望训练一个模型，自动从中找出故障报告，并取得支持判断的原句。文档数量很大，而人工预算只够阅读其中一部分，因此不能把每篇文档都交给专家逐项检查。这时真正稀缺的不是文档，而是“读懂并判断一篇文档”所需要的人工时间和专业能力。本章反复出现的“预算”，指的就是这种可以折算成费用或工时的有限人力。

即使已经决定“请人来标”，问题也远没有解决。我们至少还要回答五件事：问整篇文档的类别还是问某句的含义，要求直接作答还是核对系统给出的建议，先处理哪些文档，交给谁判断，以及对有争议的样本应追加几次检查。这五件事各管一个对象，混在一起只凭感觉决定，很容易花钱买到不能用的结果：只追求答案条数，可能留下一堆没有依据、无法复核的标签；只挑模型自己最困惑的材料去问，又可能反复撞上缺页或仅凭可见内容无法判定的文档，预算耗尽却没得到多少确定答案。

在展开方法之前，先说明本章材料里有什么，以及判断的边界在哪里。这些文档中既有表述明确的常见故障，也有稀有的故障类型、生僻的专业术语，还有仅凭可见上下文无法判断的情形。我们要辨认的是“文档是否报告了故障”，而不是仅凭作者的话就断定设备确实存在故障；作者可能只是转述、抱怨或猜测。若任务需要确认设备真的出了故障，还须另外提供检测等客观证据。本章以类别判断和证据片段定位作为主要的标注结果；偶尔出现的图像关键点、对象边界等例子，只用来解释“局部标注”到底补充了哪些标注字段，提醒读者不要把描绘几何边界的成本直接套到阅读文本上。

这五件事之间存在先后依赖，因此本章按一个自然的决策顺序展开。首先要做的是任务设计：答案的含义不固定，后面的一切都无从谈起。含义确定以后，标注方式再决定人通过什么操作把答案交出来。在任务可执行、交互也确定的前提下，样本选择决定下一批询问哪些文档，标注者选择决定这些文档交给谁，预算分配则决定各类行动各执行多少次。这五项相互制约，却不能用一个“总分”把它们合并成单次选择——失败时我们需要知道究竟是问题问错了、人没选对，还是投入不够。

本章还沿用第~\ref{chap:knowledge-acquisition-overview}章建立的样本标注接口：执行状态说明一项操作进展到哪一步，来源动作区分人是作答还是弃权，语义状态说明内容是否已经确定，验证状态说明答案是否经过核验。这套字段让“没给出答案”的不同情形能够分开处理。本章在它之上补充任务依据、具体操作和成本；至于多份答案怎样融合成一个结论、又怎样做独立的质量审计，则留到第~\ref{chap:knowledge-fusion}章和第~\ref{chap:knowledge-quality-assessment-correction}章。

\section{标注任务的设计}
\label{sec:knowledge-acquisition-human-task-design}

把一篇文档交给一位标注者，只说一句“帮助提高故障识别能力”，他其实无法开始工作。这句话说明了大方向，却没有告诉他检查哪个对象、可以依据哪些材料，以及什么情形允许给出确定答案。不同人会按各自的理解补全这些空缺，于是同一篇文档可能被标成不同结果，事后也无法判断分歧出在哪里。任务设计要解决的，正是把下游对模型的真实需求翻译成一套可以照着执行的判断，并预先规定材料或规则不足时应当怎样停止，而不是强迫每个人都交出一个看似确定的标签。任务定义、书面规范和操作界面共同决定标注者实际执行的内容，因此都要在小规模试行中接受检验；试行若暴露缺口，应改动相应环节，再回头检查原有答案是否仍然成立。

\subsection{标注目标的确定}
\label{sec:knowledge-acquisition-human-task-objective}

要把任务讲清楚，最直接的办法是先写一句可以交付、可以照做的完整请求。在文档任务里，这句话可以是：检查一篇文档，只使用正文及给定附件，回答“作者是否报告了产品故障”，并保存支持答案的原句；若无法判断，则保存无法判断的原因。

把这句话拆开看，它其实规定了五件事，称为一个标注任务的五个要素：\emph{待判断对象}是这一篇文档而非单个句子，\emph{允许依据}是正文和给定附件而不能引用外部传闻，\emph{提问}是“作者是否报告了产品故障”，\emph{回答空间}是一组事先约定的可选答案，\emph{补充字段}则要求附上支持原句或无法判断的原因。下游的每一种用途，最终都要落到这五要素中的某一项上；要素不齐，执行者就只能靠猜测补全。

先看回答空间。界面上可以提供“是、否、未知、不适用”四个选项，但要注意它们的地位并不相同。真正用于训练模型的类别只有“是、否”两个；“未知”说明这篇文档的内容尚未确定，“不适用”则说明被送来的对象根本不在任务范围内（例如把一封无关邮件误放进文档队列）。把这两者单独列出，是为了不让执行者在材料不足时被迫二选一。当一个人认真完成检查、却因为依据不足而不提交类别判断时，按第~\ref{chap:knowledge-acquisition-overview}章的接口，他的来源动作记为“弃权”，语义状态记为“未知”，而不是硬给一个“否”。

为什么连证据范围都要随答案一起保存？因为单独一个“是”字含义太宽。它可能表示作者亲历了故障，可能表示作者转述了他人的故障，也可能表示文档附上了经检测确认的故障——这三者对训练的价值和可信程度都不同。只保存“是”而丢掉任务版本和证据范围，事后就无法分辨这个标签到底回答了哪一种情形。

一个请求往往不只产生一种结果，因此还要分清不同标注字段各自回答什么。在这个请求中，文档类别和证据片段承担不同作用。类别概括整篇材料，标注者必须阅读足以判断事件内容和说话者身份的上下文；片段则进一步指出“是哪些文字支持了这个类别”，需要明确起止位置。

这两项的确定性并不同步。几句文字可能共同支持结论，此时类别已经清楚，却找不出唯一的一句证据；反过来，单独一句“有点发热”只提到异常现象，并不足以判断整篇文档是不是故障报告。因此，若训练确实同时需要两种结果，合理的顺序是先判类别，再只对判为正类的文档定位证据，并用稳定的对象标识把两项标注关联起来。这里要特别提醒：没有选中片段，只代表证据定位尚未完成或没找到合适片段，\emph{不能}据此反推出类别为负——“没找到证据句”和“确认没有故障”是两回事。

有时我们还希望标注者额外写一段理由。这时先要问清楚：这段理由写出来是给谁用的？读者不同，写法要求也不同。供争议复核使用时，理由应包含引用位置、采用了规范的哪一条、还缺什么证据，以便另一个人能据此重做一遍判断。若是用来训练文本生成模型，则要限定事实范围、长度和可接受的表述，否则模型会把标注者的自由发挥也当成目标答案去模仿。

还要警惕一个常见误解：作答之后补写的理由，可以帮助别人复核，却未必忠实重现了这个人当初形成判断的真实心理过程——人常常是先有直觉判断、再补出一套说法。因此在没有额外验证的情况下，不能把这种事后说明直接当作可靠的因果解释。

一个任务里的多个字段，难度往往不同。一篇文档可能足以判断“是否故障”和“故障部件”，却完全没有判断“严重程度”的依据。若把这些内容捆成一个必须一次填完的整体，标注者就只能在留空和猜测之间二选一。把它拆成“是否故障、故障部件、严重程度、证据片段”四项后，他可以先交出已经确定的两项，再把需要专业判断能力的严重程度单独转交给合适的人。

不过拆分不是没有代价：同一篇材料可能被反复打开，分散得到的各项结果还要重新关联回同一个文档和对象。是否值得拆，关键看时间主要花在哪一步。如果阅读全文最费时，拆成多项会让人重复阅读，反而更慢；如果最费时的是某种精细操作或专业裁决，拆分才能让大多数样本不必都经历那个最昂贵的步骤。

更重要的是，同样表现为“没有给出确定答案”，背后的原因可能完全不同，相应的补救动作也不同。附件根本缺失，属于“证据不足”，正确做法是补材料；看到了专业检测报告却读不懂，属于“缺少专业判断能力”，应转给具备相应能力的人；而文档落在规则没有写明的边界情形上，则属于“规范歧义”，应先反馈给规范负责人去澄清规则，而不是让标注者各自发挥。

这三种情况下，内容的语义状态都先保留为“未知”，完成检查却不作判断的来源动作都是“弃权”，区别在于弃权时必须记下原因。分开记录原因，才能避免两类浪费：对只是缺材料的文档不断换人重看，或对规则本身不清的任务不断追加投票——这两种做法都解决不了真正的问题。

% 小节入口连同首段长引文预留，避免边注伸入页底留白。
\Needspace{95mm}
\subsection{标注规范的制定}
\label{sec:knowledge-acquisition-human-guideline}

目标定下来以后，下一个问题是：即使两个人都认识“故障”这个中文词，他们心中所指的范围也未必相同。一位标注者可能把用户自己操作不当造成的问题算进故障，另一位可能把论坛里转述的投诉、或作者尚未确认的怀疑也算进去。两个人都很认真，答案却对不上，分歧不来自能力，而来自他们判断时遵循的不是同一条规则。

要让不同的人在不同文档上回答的是同一个问题，就需要一份\term{标注规范}（annotation guideline），也就是一套可以反复应用、写在纸面上的判断规则。一份最小可用的规范应说清：目标对象是什么、每个类别如何定义、哪些情形应纳入或排除、对象之间的边界在哪里、多个条件冲突时按什么顺序判断，以及证据不足或规则冲突时应交给谁处理。类别名称本身（比如一个叫“故障”的按钮）只能指向这些规则，不能代替规则——看到同一个词，不同人仍可能有不同理解。

除了字面规则，人的判断还会受到工作状态和界面安排的影响：长时间标注会疲劳，选项的展示位置、系统预先给出的建议都会形成先入为主的“锚”，把判断拉向某个方向。因此，为了能解释结果差异，规范不应只记录文字条款，还应连同界面顺序、每批处理多长、休息条件等一起记录\scite{knowledge-acquisition-human-thomas2022crowd}。这些条件不同，同一套规则也可能产生不同的标注分布。

规则是否清楚，可以用具体句子来检验。考虑“论坛里有人说电池坏了”。句子里确实出现了“电池坏了”，但提出故障主张的并不是作者本人，而是被转述的某个论坛用户。如果任务限定为“作者亲历的故障”，标注者就要检查说话者与事件的归属，并把这句排除；如果任务收集的是“所有被提及的故障主张”，同一句又应该纳入。同一句、同一个关键词，在两种规则下得到相反答案，差别在于对象和纳入条件改变了，而不在于标注者能不能识别“电池坏了”这几个字。这个例子说明，规则不能靠关键词匹配来执行，必须落到归属和条件上。

撰写规范时通常配三类样例：正例演示怎样才算命中条件，反例演示依据什么排除，边界例则专门用来补清证据缺失或归属冲突时怎么处理。三类样例各有分工，只给正例最容易让人过度纳入。

还要区分两类性质不同的判断，它们需要的规范不一样。一类是判断文档中是否包含某项可以核查的事实，例如“文档是否写明了更换电池的日期”；这类问题原则上有客观对错，规范应给出核查依据，并说明允许查阅哪些外部资料。另一类是判断某种回答是否可接受、质量如何，例如“这段摘要是否简洁、语气是否恰当”；这类问题带有主观性，规范无法指定唯一正确答案，只能固定评价维度、观察条件和量表每一档的含义，让不同人的看法至少能按同一口径记录下来。

对后一类主观判断，随机化展示顺序、用流程约束标注者的操作步骤，有助于控制已知的位置和顺序影响，但这些手段不能自动消除人与人之间固有的主观差异，更不能消除那些没有被观察到的偏差。承认这一点，才不会把“流程做了随机化”误当成“答案已经客观”。

表~\ref{tab:knowledge-acquisition-human-guideline-cases}把任务固定为作者亲历的故障报告，让我们逐行追踪“对象归属和证据是否满足规则”。第二行的“否”是否定作者亲历，不是否定电池可能损坏；后两行保留未知，分别送去补材料或澄清规范。表中文档都是教学构造，只说明规则怎样执行，不表示四行已经覆盖全部语言边界。

\begin{table*}[htbp]
  \centering
  \small
  \begingroup
  \setlength{\tabcolsep}{4pt}
  \renewcommand{\arraystretch}{1.15}
  \begin{tabularx}{\linewidth}{@{}>{\RaggedRight\setlength{\parindent}{0pt}\arraybackslash}p{29mm}>{\RaggedRight\setlength{\parindent}{0pt}\arraybackslash}p{27mm}>{\RaggedRight\setlength{\parindent}{0pt}\arraybackslash}p{34mm}>{\RaggedRight\setlength{\parindent}{0pt}\arraybackslash}p{14mm}>{\RaggedRight\setlength{\parindent}{0pt}\arraybackslash}X@{}}
    \toprule
    待标注文档 & 判断对象 & 适用条款 & 标注结果 & 证据与去向 \\
    \midrule
    “我升级后无法开机。” & 作者本人的设备经历 & 故障明确且归属于作者 & 是 & 保留整句及文档标识 \\
    “论坛里有人说电池坏了。” & 作者本人的设备经历 & 转述不计入作者亲历 & 否 & 保留说话者线索 \\
    “升级后情况更糟，详见附件。”但附件缺失 & 当前文档的故障判断 & 所需证据不在可见材料中 & 未知 & 证据不足：记录缺少附件，转入补材料队列 \\
    设备与外接配件并述，代词指向不明且规范未给判定顺序 & 哪个对象发生故障 & 对象归属须唯一或可核查；判定顺序缺失 & 未知 & 规范歧义：反馈规范负责人 \\
    \bottomrule
  \end{tabularx}
  \endgroup
  \caption{标注规范怎样落实为可执行判断。四行都判断作者亲历的故障报告：“是、否”针对这一任务，“未知”保留尚未确定的内容。规则同时决定保留哪项证据和交给谁继续处理；教学样例不表示已覆盖全部边界。}
  \label{tab:knowledge-acquisition-human-guideline-cases}
\end{table*}

书面规则写清之后，还要保证它真正落实到标注者接触的另外两个环节：操作界面和培训材料。三者必须使用同一套规则、共享同一个版本标识：界面上的选项要与规范规定的回答空间一一对应，培训中的反馈要明确指出标注者具体违反了哪一条，而不是笼统地说“标错了”。

规范会随认识深入而更新，这里有一个容易被忽视的陷阱。当规范从版本$v$更新到$v+1$时，应列出哪些条款改变了答案的含义、哪些历史样本会受影响，再分别决定这些样本需要重标、复核还是仍按旧版本保留。如果规则的实际语义已经改变，却只是替换了说明文字、然后把新旧标签直接合并，那么同一个类别名称下就会悄悄混着两套不同标准产生的答案，模型无从区分。

\subsection{标注任务的试标与调整}
\label{sec:knowledge-acquisition-human-pilot}

规范在纸面上完整，并不等于任务在真实操作中行得通。正式标注者可能根本找不到界面里的附件、误解某个选项，或在某类文档上耗费远超预期的时间。这些问题不经过实际操作很难发现。因此，在扩大规模之前，应让真正将要参与工作的人员，按完整流程先处理一小批材料，这就是\term{人工试标}（human pilot annotation），即一次由人执行的小规模试运行。

试标样本要刻意搭配：既要有典型情形以估计正常速度，也要有事先预判的难例和已知的边界例，专门用来检验规则覆盖得够不够。过程中要完整保存答案、证据、弃权原因、阅读与操作时间以及标注者提出的疑问。试标要回答的是三个操作性问题——任务能不能执行、人员在哪里受阻、单条成本大概是多少；至于样本集整体的标注质量和其中的统计不确定性，属于第~\ref{chap:knowledge-quality-assessment-correction}章的内容。需要说明，“试标”这一动作的执行者也可以是模型，那是另一种执行主体，在第~\ref{sec:knowledge-acquisition-model-pilot-validation}节讨论，本节专指由人执行的试标。

试标中出现不同答案时，不要立刻归因于“某人能力差”，而应按顺序排查三个可能的环节：先确认大家读到的材料是否相同，再确认他们采用的是不是同一条规则。如果材料一致、却是多人对同一条款作出不同解释，问题出在规范本身有歧义，可能需要改规则；如果规则清楚、但证据在材料里根本没有出现，那么加强培训也补不上，要改的是材料供给；只有在规则明确、证据也充分的前提下仍然做错，才直接指向培训、界面设计或人员准入的问题。按这个顺序排查，才能对症下药。

对难例可以组织讨论，但务必先把每个人各自的独立初始答案保存下来，再坐下来讨论。若先讨论、后记录，最终写下的是讨论后的共识，分歧最初发生在哪一个判断步骤就被掩盖掉了，而那恰恰是最有价值的诊断信息。

用一个教学式的试行过程说明这种连锁调整。第一轮刻意集中难例，结果发现许多文档的结论依赖附件中的检测报告，而界面当时只显示正文，标注者无从判断。团队据此做了两处修改：在界面中补充附件摘要，并把“附件缺失”明确列为一种证据不足的弃权原因。到第二轮，判断变得可执行了，但因为要多读附件，每篇材料变长，单例的中位处理时间也随之上升。这个例子说明，补证据和估成本必须一起更新：让更多判断可执行的同一改动，也改变了每条判断的成本，不能只沿用第一轮测得的时间。

还要注意试标样本的构成会影响成本外推。第一轮是为了暴露问题而刻意集中难例，这批样本本来就偏难、偏慢，因此不能把它的平均耗时直接乘以生产环境的文档总数来估算总投入。正式扩大规模前，应另按预计的类别、长度和难度构成抽取一批有代表性的样本重测，同时始终保留一个反馈入口，让上线后出现的新型问题还能回流到任务设计。

\section{标注方式的选择}
\label{sec:knowledge-acquisition-human-interaction}

任务和规范确定之后，还要决定标注者通过什么具体操作把答案交出来。面对同一篇文档，他可以自己从零选出类别，也可以核对系统预先填好的类别；可以从头逐句寻找证据，也可以只修改系统推荐的片段。任务要求（判断是否为故障报告）并没有变，变的是提交标注这个动作本身。\term{人工交互方式}（human annotation interaction）指的就是这样一组安排：标注者在界面上看见什么、需要做什么操作、系统最终保存什么。

交互方式与标注粒度可以自由组合：选上位大类时不必从零开始，划定细片段时也不必全部由人亲手圈出。但正因为操作形态差别很大，比较不同方案时必须把阅读、判断、录入、模型推理和返工的成本一起算进去，不能把“点一下按钮”“在两个答案间二选一”和“写出一条完整答案”当作等量的监督信号——它们消耗的人力和携带的信息量都不同。

\subsection{直接给出答案}
\label{sec:knowledge-acquisition-human-direct-answer}

最基础的交互方式，是系统只提供文档、问题和规范，由标注者自己寻找证据并产生答案，不提供任何预填结果，这称为\term{直接作答}（direct annotation）。它对应“人独立完成”的情形。

答案的形态应跟随问题的类型来选。类别集合明确、选项有限时，用点选最合适；带单位和量程的测量值（如温度、时长）应直接填数值；证据片段或图像区域需要用边界工具划定起止；摘要、改写这类没有固定选项的结果则要用开放文本框。无论哪种形态都要注意两点：选项必须覆盖所有合理答案，并保留“未知”和“不适用”的出口；开放回答则要事先限定事实范围、长度和可接受的同义变体，否则后续处理时很容易把“措辞不同”误当成“标注结论不同”。

仍用同一篇故障文档，可以分别安排三种任务，看到它们的完成条件并不相同。类别任务要求提交“是、否、未知、不适用”之一，选未知时还要附上原因码；片段任务要求标出支持判断的最小连续文本，当证据分散时允许保留多个不相邻片段；一句话摘要则只允许复述文档内可以核查的事实，不能加入外部信息。

三种任务用的是同一份材料，验收标准却不能统一成一句“都已填写”——类别任务看分类是否正确，片段任务看位置是否准确，摘要看事实是否忠实。涉及边界的任务还要预先规定片段能否重叠、能否嵌套、最小单位是什么，例如标注型号名称时是否包含前面的品牌名。把这些约定写明，学习算法才不必去猜测人圈出的范围到底是什么意思。

一次直接作答可以拆成四个耗时不同的环节：理解问题、寻找证据、形成判断、录入结果。不同任务的瓶颈所在完全不同：长文档可能主要耗在全文里查找证据，图像多边形标注可能主要耗在沿轮廓精细描绘。因此，“把点击步骤从三次减到一次”到底能省多少总时间，取决于录入环节原本占多大比例；若时间几乎都花在阅读判断上，优化按钮并无意义。快捷键和批量录入只能减少重复的机械操作，替代不了逐例的理解与判断。

是否强制要求填写额外理由，也要按真实用途决定。只有当理由确实会用于训练、审计或争议处理时，才值得设为必填；若只是“多写一点总没错”，它既增加每份的工作量，还可能产生反效果——人为了填满必填框，事后补写出自己当初并无充分依据的解释。

\subsection{提供粗粒度或局部标签}
\label{sec:knowledge-acquisition-human-partial-label}

当逐篇确定具体故障部件、再定位全部证据的做法超出预算时，可以退一步，先只问较少的标注字段。这里有两种不同的“少做一点”，要注意区分。\term{粗粒度标签}（coarse-grained label）是在答案的类别层次上让步：只缩小答案范围，例如确认“这是硬件问题”，但暂时不区分到底是电池、电源还是主板。\term{局部标签}（partial annotation）则是在对象范围上让步：只判断对象的一部分，例如只标出一处证据片段、或只标一个图像关键点，而不处理其余部分。

这两种标签的共同点是：它们都对最终想要的完整答案构成了某种约束，却没有把完整答案补齐——它们是“关于答案的信息”，而不是答案本身。图像任务尤其能说明标注类型如何影响成本：对象越复杂、要求的标注越精细，人工工作量越大；因此在选择标注方式时，需要同时估计一条候选标签能带来多少价值、又要投入多少标注成本，这正是成本感知选样的出发点\scite{knowledge-acquisition-human-vijayanarasimhan2009cost}。

要让学习算法只使用人真正确认过的部分、而不把空缺误读成答案，就必须把两类信息分开记录：一类是“哪些细类别还没有被排除”，另一类是“哪些位置已经被检查过”。

先形式化粗标签。设完整类别$y$属于细类别集合$Y$。对象$x$的粗标签依据预先建立的父子类别映射，返回一个集合$S(x)\subseteq Y$，它确认真值落在这个集合内，即$y\in S(x)$，但并不指定集合中究竟哪一个元素才是真值。再形式化局部标签。把对象可划分的位置编号为$j=1,\ldots,\ell$，用$m_j\in\{0,1\}$标记该位置是否经过检查；只有当$m_j=1$时，才记录该位置上的局部答案$a_j$。于是这一阶段已知的约束可以写成
\begin{equation}
  y\in S(x),
  \qquad
  \mathcal{A}_{\mathrm{local}}(x)
  =\{(j,a_j):m_j=1\}.
  \label{eq:knowledge-acquisition-human-partial-label}
\end{equation}
式~\eqref{eq:knowledge-acquisition-human-partial-label}让“类别被限制在某个范围”和“某些位置已查、答案是什么”各自独立地进入训练，凡是没有确认的部分一律留空。两个方向的空缺都不能误读：$m_j=0$只表示这个位置没检查过，绝不表示它是负类；集合$S(x)$也不能随便挑其中一个元素当作真值，因为人只确认了“真值在集合里”，没有确认是哪一个。保存父子类别映射、对象标识和已检查范围，正是为了在把不同标注统一成同一种存储格式时，不丢掉这些细微但关键的区别。

这样的约束虽然不完整，仍然可以用于训练。既然人已经确认真值属于$S(x)$，训练时就可以提高模型分配给这个集合的总概率，而不必强迫集合内各个细类别均分概率——没有任何依据要求它们等概率。第~\ref{chap:knowledge-fusion}章的式~\eqref{eq:knowledge-fusion-set-loss}给出了对应的集合损失。

使用时要守住两条边界。其一，这个约束用到的是“真值确实在集合中”这一已被确认的事实，但仅凭集合约束，模型依然无法分辨集合内部的细类别，不能假装它能。其二，只有经过人确认的集合才有此资格；如果$S(x)$只是模型自己给出的一串候选、尚未经过核验，就不能直接把它当作已确认约束塞进训练。

图~\ref{fig:knowledge-acquisition-human-granularity}改用同一个几何场景，直观展示不同标注字段到底补充了什么信息：整图类别只说明这张图里出现了目标；关键点进一步给出目标上的具体位置；包围框再进一步给出一个能包住已检查对象的矩形范围。三者的信息量依次增加，但要注意包围框内仍可能含有背景，所以它并不等于沿对象轮廓勾出的分割边界。图中最后一个面板只为“已检查的对象”画了框，这不代表整张图都已核查过，其余区域仍是未知。

减少标注字段是否真能省钱，仍要回到时间结构上判断。如果最费时的是在图中寻找对象，那么把包围框简化成一个点，可能只省下少量录入操作；如果最费时的是框的定位与录入，简化成点则可能省得较多。单凭“标了多少条标签”这个数字，无法表达这种成本差别。

\begin{figure*}[htbp]
  \centering
  % 同一教学场景中的图像级、关键点与包围框标注；斜线区表示未检查区域。
\begin{tikzpicture}[x={\dimexpr\linewidth/182\relax},y=1mm,scale=.98,transform shape,font=\figurefont\footnotesize,text=BookInk]
  \tikzset{
    scene/.style={draw=BookRule,fill=BookPaper,line width=.8pt},
    object/.style={line width=1pt},
    selected boundary/.style={draw=FigureModel,line width=1.3pt},
    panel title/.style={font=\figurefont\bfseries\footnotesize,anchor=south},
    output/.style={draw=BookRule,fill=white,minimum width=48mm,minimum height=8mm,align=center}
  }
  \foreach \shift/\panel/\title in {0/a/图像级类别,61/b/对象关键点,122/c/对象包围框}{
    \begin{scope}[xshift=\shift mm]
      \node[panel title] at (26,49) {(\panel)\quad\title};
      \path[scene] (1,8) rectangle (51,43);
      \draw[object,draw=BookBlue,fill=FigureInput!18] (14,27) circle (6);
      \draw[object,draw=BookAmber,fill=FigureLoss!18] (25,13)--(37,13)--(31,27)--cycle;
      \draw[object,draw=BookTeal,fill=FigureModel!18] (39,30) rectangle (48,38);
    \end{scope}
  }
  \begin{scope}
    \node[output] at (26,3) {标注：场景含目标对象};
  \end{scope}
  \begin{scope}[xshift=61mm]
    \begin{scope}
      \clip (25,8) rectangle (51,43);
      \foreach \d in {-20,-15,...,70}{\draw[BookRule,line width=.55pt] (\d,8)--({\d+35},43);}
    \end{scope}
    \draw[BookMuted,dashed,line width=.8pt] (25,8) rectangle (51,43);
    \fill[FigureModel] (14,27) circle (1.5);
    \draw[white,line width=.7pt] (12.6,27)--(15.4,27) (14,25.6)--(14,28.4);
    \node[output] at (26,3) {标注：对象甲的关键点};
    \node[anchor=north east,text=BookMuted] at (49.5,41.5) {未检查};
  \end{scope}
  \begin{scope}[xshift=122mm]
    \begin{scope}
      \clip (38,8) rectangle (51,43);
      \foreach \d in {-20,-15,...,70}{\draw[BookRule,line width=.55pt] (\d,8)--({\d+35},43);}
    \end{scope}
    \draw[BookMuted,dashed,line width=.8pt] (38,8) rectangle (51,43);
    \draw[selected boundary] (7,20) rectangle (21,34);
    \draw[selected boundary] (23,11) rectangle (38,29);
    \node[output] at (26,3) {标注：已检查对象甲、乙的框};
    \node[anchor=north east,text=BookMuted] at (49.5,41.5) {未检查};
  \end{scope}
\end{tikzpicture}

  \caption{同一场景的图像级、关键点与包围框标注。（a）给整图类别，（b）给对象甲的关键点，（c）给已检查对象甲、乙的包围框；三个面板共享同一教学场景。矩形框不提供分割轮廓；框已给出不表示全图已检查，斜线区域仍未知，不能作负类。几何形状不是真实识别结果，也不表示成本比例。}
  \label{fig:knowledge-acquisition-human-granularity}
\end{figure*}

\subsection{比较候选结果}
\label{sec:knowledge-acquisition-human-comparison}

有时让人“写出一个理想答案”反而比“判断两个现成答案哪个更好”更困难。例如让人凭空概括故障原因，需要自己组织语言、回忆要点；而给他两份现成概括，他往往能较快看出哪份更准确。当生成比判断更费力时，可以改用\term{候选比较}（candidate comparison）：不要求人给出标准答案，只要求他对系统提供的若干候选结果记录偏好、排序或按维度打分。

比较要先固定一个评价维度$d$，否则“哪个更好”无从谈起。记录时用三个符号区分不同结果：$a\succ_d b$表示在维度$d$上更偏好$a$而不是$b$，$a\sim_d b$表示两者打平，$a\parallel_d b$表示依据不足、或两者根本不可比较。必须强调这类标注回答的问题很有限：它只说明给定的这几个候选在该维度上的相对关系，既不告诉我们候选集之外是否还有更好的答案，也不给出任何候选的绝对质量分数。

举个例子，候选甲完整引用了诊断证据但篇幅较长，候选乙很简短却遗漏了关键的限定条件。若按“事实准确”这一维度比较，人可能选甲；若改按“简洁”比较，则可能选乙。同两份候选、同一个人，只因为评价维度不同就给出相反偏好，这正说明比较结果必须连同维度一起记录。这里新增的信息是“在某个具体维度上的取舍”，而不是对两份回答给出一个笼统的总体排名。

因此每条比较标注都应关联问题、候选的具体内容、它们的展示位置、评价维度和作答者。若只在数据库里留下“甲更好”三个字，将来训练偏好模型时就无从知道人当时看重的到底是事实准确还是文字简洁，这个标签也就失去了意义。

如果不用两两比较、而采用按维度打分，关键是让量表的每一档都有可以重复理解的具体含义，而不是只写“1到5分，越高越好”。以“事实准确”的五档序数量表为例，可以这样锚定：“1”表示核心事实与材料直接冲突，“3”表示核心结论成立但遗漏了重要限定，“5”表示关键事实及其限定都有材料支持；“2”和“4”则位于相邻两个锚点之间。有了这些描述，不同的人才大致按同一把尺子打分。

同时要记住，不同维度的分数不能互相替代或简单加总。同一份回答在“事实准确”上可以得5分，在“简洁”上却只有2分，因为这两个分数约束的是完全不同的属性。此外，序数量表只提供顺序信息，它并不保证“从5到3”与“从3到1”代表相等的质量差距；在没有测量模型和校准证据时，也不能直接比较两个打分习惯未经对齐的人给出的分数。

候选增加到三个以上时，还会出现新的问题。成对比较的结果未必能拼成一个唯一的总排序：除了平局，还可能出现循环偏好——甲优于乙、乙优于丙、丙又优于甲，三者无法排出先后。这时需要另外规定排序规则或引入偏好模型来汇总，不能指望从一次比较直接推出全局次序。还有一种容易被忽略的情形：即使某个候选在比较中“胜出”，也可能是“矮子里拔将军”，候选集里本来就没有一个真正合格的答案。因此界面应允许“两者都不合格”，并补充追问，而不是强迫人在烂答案里选一个。

把候选展示给人时，随机或均衡它们的出现位置，可以检查并部分抵消位置带来的影响，但无法消除措辞、长度和标注者专业背景等造成的全部差异。最后还要算清比较次数：$J$个候选若每一对都比较一次，共需$J(J-1)/2$次判断，随候选数快速增长；为控制成本，可以只查询其中一部分成对关系，但必须明确哪些顺序仍未确定，不能把没比过当成已经相等。

\subsection{确认或修改模型建议}
\label{sec:knowledge-acquisition-human-model-suggestion}

前面几种方式都由人主导，现在引入模型的参与。如果模型已经在文档中找出了可能的证据句，人就不必再从整篇文档里重新搜索一遍，而可以直接核对模型给的结果。\term{模型辅助标注}（model-assisted annotation）指的就是这样一种分工：先由模型提出类别、边界或草稿，再由人来确认、修改或拒绝。

它的收益与代价是同时出现的：模型的预填可能替人省下搜索和录入的时间，但也新增了核对建议、定位模型错误、以及撤销错误建议的工作，一旦人偷懒直接采纳，模型的错误就会被带进标签。因此界面必须把人的不同操作区分开：确认、局部编辑、完全重写、拒绝后另给答案、弃权，各自单独记录。尤其要规定，只有人明确点下“提交且不作改动”时，才能算作确认；超时未处理或随手跳过都不代表认可，否则系统会把大量无人过问的建议误记成人已同意。

这些操作上的区别要落实到第~\ref{chap:knowledge-acquisition-overview}章的状态字段里，否则事后无法还原发生了什么。任务已经发出、但人还没开始，执行状态是“未处理”；已经开始却尚未提交，是“处理中”，此时来源动作仍为空。人正常提交判断，来源动作记为作答；人完成检查却决定不判断，记为弃权——这两种情况下执行状态都可以是“成功”。

这里要再次强调“成功”的准确含义：它只表示这一步操作顺利完成，既不表示内容正确，也不表示已经过核验。反过来，系统根本没有向人发出请求、或程序在中途执行失败，也绝不能记成人的主动弃权——前者没有动作发生，后者是技术故障，二者与人在检查后选择不答是完全不同的情况。

为了在出错时能追查原因，系统需要同时保存三样东西：模型当初给出的建议及其版本、人实际做了什么操作、最终提交的答案。有了这些记录，才能判断一个错误标签究竟是人盲目沿用了模型的错误建议，还是他无视建议、另行作答时自己产生的。

要研究预填到底怎样影响人，可以对同一篇文档比较三种界面：只给材料、让人从零作答；把模型的多个答案列成候选清单；以及直接默认填入一个答案。三种界面下任务完全相同，唯一改变的是人首先看到什么。一个规律是：预填越具体，越可能替人省下操作，但也越容易让他的核对依附于那个默认答案——人倾向于接受摆在眼前的结果。这种便利与盲从之间的张力，需要靠实验而不是直觉来权衡。

自动建议究竟有没有好处，不能一概而论，必须结合具体任务检验。Schulz等人在医学和教学两个领域的诊断推理文本上，分别比较了“有建议”和“无建议”两种条件下的标注耗时，以及标注者之间的一致性；后者用Krippendorff一致性系数$\alpha_U$来衡量，它刻画不同人标出的片段边界和类别彼此吻合的程度。他们的研究观察到建议确实节省了时间，但一致性上的收益随领域不同而有差异，并且在所研究的任务、模型和界面下，没有发现值得注意的“建议偏差”，即人被建议系统性带偏的现象\scite{knowledge-acquisition-human-schulz2019suggestions}。

解读这类结果时要保留两层谨慎。其一，标注者之间一致性提高，并不能直接证明标签更接近客观真值——大家可能被同一个建议带向同一个错误。其二，平均值可能掩盖局部问题：即使总体上平均节省了时间，若模型对某一类文档集中给出错误建议，那一类文档会出现系统性的大量返工，平均数字看不出这个隐患。

因此，比较不同建议界面时要控制变量：固定任务、材料和最终的质量检查口径，再把模型推理成本、人的核对时间和修改时间一并记录。一个有用的做法是留出一部分随机样本，故意不展示建议、让人独立作答，用这部分对照来检查默认答案是否改变了错误的分布；对重要的争议样本，也应始终保留一次不受建议影响的独立判断。

这里有一个评价上的逻辑陷阱必须避开：人在模型建议辅助下共同产出的答案，已经受到该模型的影响，因此不能反过来又把它当成评价“同一个模型”的完全独立参照——用被模型影响过的结果去证明模型正确，等于自我循环论证。

表~\ref{tab:knowledge-acquisition-human-interactions}把上面四种交互方式放在一起，对照每种方式新增了什么（标签、约束或偏好）、又留下了哪些未知；通用的来源和状态字段仍按章首的接口统一保存，不因方式不同而改变。

需要强调，这四行并不是互斥的四选一方案，完全可以组合使用，例如先确认模型给的类别建议、再由人从零定位证据。但组合时成本不能简单相加：两步可能共享同一次阅读而节省时间，也可能因为多出一次核对而增加返工。因此组合方案的成本要重新实测，不能直接把两种方式各自的平均时间加起来当作组合后的成本。

\begin{table*}[htbp]
  \centering
  \small
  \begingroup
  \setlength{\tabcolsep}{3pt}
  \renewcommand{\arraystretch}{1.15}
  \begin{tabularx}{\linewidth}{@{}>{\RaggedRight\setlength{\parindent}{0pt}\arraybackslash}p{18mm}>{\RaggedRight\setlength{\parindent}{0pt}\arraybackslash}p{24mm}>{\RaggedRight\setlength{\parindent}{0pt}\arraybackslash}p{25mm}>{\RaggedRight\setlength{\parindent}{0pt}\arraybackslash}p{29mm}>{\RaggedRight\setlength{\parindent}{0pt}\arraybackslash}p{24mm}>{\RaggedRight\setlength{\parindent}{0pt}\arraybackslash}X@{}}
    \toprule
    方式 & 可见文档及规范 & 人工操作 & 所得标签、约束或偏好 & 仍未知 & 成本与风险 \\
    \midrule
    直接作答 & 文档、问题、规范 & 选类别、定位或撰写 & 答案及支持证据 & 未问属性、未查区域 & 阅读与生成；可能遗漏限定 \\
    粗粒度或局部判断 & 文档、类别层级或局部范围 & 选上位类、集合或局部标记 & 完整目标的约束、已查范围 & 细类别、未查位置 & 查找对象；可能把遗漏当负类 \\
    候选比较 & 问题、候选、评价维度 & 偏好、平局、不可比或评分 & 相对关系或维度评分 & 候选集外答案、未比较关系 & 比较对数增加；位置与措辞偏差、循环偏好 \\
    确认或修改建议 & 文档、建议、上下文 & 确认、编辑、重写、另答或弃权 & 原建议及版本、人工操作、最终答案 & 未独立检查的替代答案 & 生成、核对与返工；可能沿用系统性错误 \\
    \bottomrule
  \end{tabularx}
  \endgroup
  \caption{人工交互方式与所得标签、约束或偏好。四行固定文档任务，比较可见文档及规范、人工动作、所得标签、约束或偏好和仍未知的部分；来源与状态字段另按章首接口保存。成本列指出需要测量的环节，不表示任何方式无条件更快或质量更高。}
  \label{tab:knowledge-acquisition-human-interactions}
\end{table*}

\section{标注样本的选择}
\label{sec:knowledge-acquisition-human-sample-selection}

前面解决了“问什么、怎么问”，现在任务和界面都固定下来，回到一个现实约束：候选文档仍然远远多于预算能支付的人次，不可能每篇都问。于是问题变成：下一批有限的人工，应该先花在哪些文档上最值得？这一节不再改动答案的含义，而是研究怎样挑选待标注样本。

最朴素的做法是随机抽取，但直觉上不同文档的价值并不相同：在模型已经很有把握的常见文档上再问一次，学到的东西很少；在它拿不准的文档上问，可能更有帮助。让学习系统自己参与这个挑选过程，根据已有的训练样本和当前模型的状态来决定下一个去问谁，就称为\term{主动学习}（active learning）。这里“主动”的含义是：新增哪些训练样本不再完全随机，而是由当前的学习状态来决定。

本章采用最常见的池式设定：把所有手头的未标注文档想象成一个静止的“候选池”。在第$t$轮，记已有标签的样本集合为$L_t$，尚未标注的候选池为$U_t$，并用$L_t$训练出当前模型，其参数记为$\bm{\theta}^{(t)}$。一轮的流程是：先固定当前模型和现有证据，给池中的候选打分、选出一批去问人；拿到人的反馈之后，才更新模型，进入下一轮。要注意打分依据的是“更新前”的旧模型，不能边收答案边改变同一轮的选择标准。若文档是随时间陆续到达、而非一开始就全在池中，那只是流式到达，它只改变某一时刻能看到哪些候选，并不构成一种新的价值准则。

主动选择是否真的比随机好，需要通过比较来判断，而不是默认它一定更优。比较的基线就是随机抽样；如果我们事先知道某些类别、来源或时间段特别重要，还可以把分层抽样作为更强的基线——先按这些维度分层，再在每层内抽取。

主动选样有一个必须记住的副作用：它改变了各篇文档进入已标集合的概率，被选到的不再是总体的一个随机缩影。因此，这样积累起来的标签不能不加校正就当成对总体做质量审计的样本来用，否则统计结论会有偏。图~\ref{fig:knowledge-acquisition-human-query-loop}画出了各种策略共享的查询循环；后面要讲的不确定性、代表性和预期收益，差别主要就在“如何打分或如何选一批”这一步，循环骨架是相同的。循环里还要区分两种结果：成功提交的可用答案进入已标集合并更新模型；而来源弃权只增加一条弃权原因和已付出的成本，内容仍未知，不增加确定标签、也不更新模型。无论哪种结果，候选池都要据此更新，尤其要避免把根本无法作答的文档原样反复送回来空耗预算。

\begin{figure}[htbp]
  \centering
  % 池式主动选择流程；虚线路径保留来源弃权及原因，不把查询数等同于有效标签数。
\begin{tikzpicture}[x=1mm,y=1mm,font=\figurefont\footnotesize,text=BookInk,
  box/.style={draw=BookRule,fill=BookPaper,line width=.8pt,rounded corners=1.2mm,minimum width=42mm,minimum height=11mm,align=center,inner sep=2.5pt},
  decision/.style={box,draw=FigureModel,fill=FigureModel!10},
  flow/.style={knowledge arrow,draw=BookTeal,line width=1.1pt,rounded corners=1.2mm},
  secondary/.style={knowledge arrow,draw=BookAmber,line width=.85pt,dashed,rounded corners=1.2mm},
  edge label/.style={font=\figurefont\scriptsize,fill=white,inner sep=1.5pt}]
  \node[box,draw=FigureInput,fill=FigureInput!10,minimum width=31mm] (state) at (16,0) {已有标注$L_t$\\与候选池$U_t$};
  \node[box,minimum width=31mm] (criterion) at (51,0) {按当前依据\\计算查询分数};
  \node[box,minimum width=31mm] (select) at (86,0) {选择样本$x_t$\\并发出任务};
  \node[decision,minimum width=42mm] (judge) at (86,-18) {人工能否依据\\现有观测作答？};
  \node[box,draw=FigureOutput,fill=FigureOutput!7,minimum width=44mm] (usable) at (28,-36) {可用答案\\加入$L_{t+1}$并更新模型};
  \node[box,draw=FigureLoss,fill=FigureLoss!10,minimum width=45mm] (pending) at (80,-36) {来源弃权\\记录原因，$L_t$与模型不变};
  \node[decision,minimum width=54mm] (finish) at (48,-54) {更新$U_{t+1}$与累计成本后\\停止条件是否满足？};
  \node[box,draw=FigureOutput,fill=FigureOutput!7,minimum width=31mm] (end) at (95,-54) {停止并保留\\未知及其原因};
  \draw[flow] (state)--(criterion);
  \draw[flow] (criterion)--(select);
  \draw[flow] (select)--(judge);
  \draw[flow] (judge.south west)--++(0,-4)-|node[edge label,pos=.72,above] {可用}(usable.north);
  \draw[secondary] (judge.south east)--++(0,-4)-|node[edge label,pos=.72,above] {弃权}(pending.north);
  \draw[flow] (usable.south)--++(0,-4)-|(finish.north west);
  \draw[secondary] (pending.south)--++(0,-4)-|(finish.north east);
  \draw[flow] (finish)--node[edge label,above] {是}(end);
  \draw[flow] (finish.west)--node[edge label,above,pos=.25] {否}(3,-54)--(3,-12)--(51,-12)--(criterion.south);
\end{tikzpicture}

  \caption{主动选择样本的查询循环。实线路径把成功提交的可用答案加入已标集合并更新模型；虚线路径记录成功提交的来源弃权及原因，内容仍未知，已标集合和模型不变。两条路径都更新候选池与累计成本，再检查预算、覆盖或独立验证。返回箭头只表示重算依据，不保证质量提高。}
  \label{fig:knowledge-acquisition-human-query-loop}
\end{figure}

\subsection{基于不确定性的样本选择}
\label{sec:knowledge-acquisition-human-uncertainty-selection}

第一种自然的想法，是把宝贵的人工答案用在模型自己最拿不准的地方。如果模型对大部分常见故障已经能稳定判断，再去问这些文档意义不大；而它犹豫不决的文档，往往正是当前决策边界附近、最需要一个标签来厘清的地方。\term{不确定性采样}（uncertainty sampling）就是按当前模型的犹豫程度给候选排序、越犹豫越优先询问。

“犹豫”可以从两个角度衡量，要注意它们并不等价：一是看\emph{单个}模型如何在各个类别之间分配概率，二是看\emph{多个}候选模型是否对同一文档给出不同意见。前者反映一个模型内部的不确定，后者反映不同合理解释之间的分歧。还必须提前说明一个贯穿本节的警告：不确定性分数高，并不一定意味着这篇文档“值得且可以”标注。文档损坏、规范有歧义、或出现了模型从未见过的新类别，都可能打出很高的不确定分，但这些情形靠再问一次人未必解决得了。因此分数只提供“优先问谁”的依据，仍要另外检查可答性和覆盖，绝不能把高分直接等同于“将得到一条可用新标签”。

\subsubsection{基于预测不确定性的样本选择}
\label{sec:knowledge-acquisition-human-predictive-uncertainty}

对同一篇候选文档，模型可能几乎把概率全压在一个类别上，也可能在几个类别之间难分高下。要把这种差别用于排序，需要一个能从模型输出直接算出来的分数。固定当前模型和同一个$K$类空间，记模型认为文档$x$属于类别$c$的概率为
\[
  p_c(x)=p_{\bm{\theta}^{(t)}}(y=c\mid x),
\]
再把这$K$个概率从大到小排列为$p_{(1)}(x)\geq p_{(2)}(x)\geq\cdots\geq p_{(K)}(x)$，其中$K\geq2$，括号下标表示名次而非类别编号。

“犹豫”至少有三种合理的度量，它们各自盯着概率分布的不同特征：最低置信度（least confidence）只看第一名是否足够突出，类别间隔（margin）看前两名咬得有多紧，预测熵（predictive entropy）则看整个概率分布是否普遍分散。为了让“越犹豫越优先”方向一致，把三者都写成分数越大、越应优先查询的形式：
\begin{equation}
  \begin{aligned}
    u_{\mathrm{LC}}(x)
      &=1-p_{(1)}(x),\\
    u_{\mathrm{M}}(x)
      &=1-\bigl[p_{(1)}(x)-p_{(2)}(x)\bigr],\\
    u_{\mathrm{H}}(x)
      &=-\sum_{c=1}^{K}p_c(x)\log p_c(x),
  \end{aligned}
  \label{eq:knowledge-acquisition-human-predictive-uncertainty}
\end{equation}
其中约定$0\log0=0$，以免某个类别概率为零时对数无定义；本章统一用自然对数，此时熵以nat为单位。理解熵时要注意它加权求和的对象：它平均的是“这\emph{一篇}文档可能属于各个类别”所携带的不确定性，既不是不同候选文档之间的差异，也不是历史上已观测错误的平均。

式~\eqref{eq:knowledge-acquisition-human-predictive-uncertainty}把三种角度不同的犹豫程度统一成了“越大越该问”的分数。实际使用时，每轮在$U_t$中按选定的一种分数从高到低排序，查询预算允许的前若干篇；若出现分数完全相同、无法排序的情况，可按事先约定随机打破平局。这种顺序式的不确定性采样最早由Lewis与Gale在文本分类中系统使用，这里保留其“每轮挑最不确定样本”的思路，并给出三种分数方向统一的多分类写法\scite{knowledge-acquisition-human-lewis1994uncertainty}。

\begin{example}[三种不确定性分数可能给出不同排序]
  \label{ex:knowledge-acquisition-human-uncertainty-scores}
  为了只比较分数而不混入模型差异，设同一模型对两个候选给出三分类教学分布$\bm{p}^{A}=(0.50,0.49,0.01)$与$\bm{p}^{B}=(0.45,0.30,0.25)$，两者都已归一化。$A$的最低置信度、间隔分数和自然对数熵依次是$0.50$、$1-(0.50-0.49)=0.99$和$0.742$；$B$依次是$0.55$、$1-(0.45-0.30)=0.85$和$1.067$。$A$的前两类几乎平分，所以间隔准则选$A$；$B$的第一名更弱、整体更分散，所以最低置信度和熵选$B$。二分类中，给定最大类别概率就确定另一类概率，三种分数都随它单调变化，除平局外排序相同；多分类增加了剩余概率怎样分配的自由，等价关系不再成立。
\end{example}

前面的分数针对的是“整篇文档给一个类别”的任务。如果任务变成逐位置打标签（例如给文档中每个句子、或图像中每个像素分类），单个位置的犹豫还不能直接决定整篇文档值不值得打开——人是按整篇付出阅读成本的。这时可以先对$j=1,\ldots,\ell$各位置用同一口径算出局部分数$h_j$，再用三种方式汇总成文档级分数：取总和$\sum_jh_j$、取平均$\ell^{-1}\sum_jh_j$，或取最大值$\max_jh_j$。

三种汇总各有偏差，不能随意互换：总和随位置数增长，会系统性偏向长文档；平均消除了长度影响，却可能把单个关键难点稀释掉；最大值只盯住最难的一点，又可能被一个异常位置主导整篇的排序。选择哪一种，取决于我们关心的是“整篇总体多难”“平均多难”还是“有没有特别难的点”，这是评价单位的改变。当标注成本随文档长度变化时，还可以进一步比较性价比$u(x)/\widehat c(x)$，即用价值除以预计成本；但$\widehat c(x)$必须是另行用数据验证过的人工成本估计，字符数之类的代理量不能直接当作耗时。

上述三种分数都建立在模型给出的概率之上，因此它们的可靠性受限于概率本身是否可信。当概率失准时，$p_c(x)$只是模型的一厢情愿，不能解释成它真实的正确率。这会造成两个方向的误判：高熵既可能来自正当的类别边界（值得问），也可能来自缺页、跨语言或规范缺口这类根本无法作答的情形（问了也白问）；低熵同样不可靠，模型可能对一个从未见过的新类别非常自信地给出错误答案，看上去反而“很确定”。

因此实践中不能全押在不确定性上，应保留随机或分层抽样作为兜底，保证各类情形仍有覆盖；同时把弃权的具体原因反馈回任务设计环节。若不对可答性加以把关，系统可能对同一类“无法作答”的文档反复发出查询，预算被持续消耗，而底层缺失的证据或规则却始终没有被补上。

\subsubsection{基于模型分歧的样本选择}
\label{sec:knowledge-acquisition-human-model-disagreement}

单个模型的概率分散，刻画的是它自己拿不准；但换一个同样合理、训练得同样好的模型，会不会也拿不准，则是另一回事。设想已有的故障标签主要覆盖常见表述，那么可能存在两套不同的预测规律：它们在所有已标注的常见文档上都符合当前的拟合准则，却在某种少见的技术写法上给出截然不同的预测。有限的训练样本只约束了“已经观察到的情形”，并没有排除这些尚未被观测区分开的规律。

既然如此，武断地把某一次训练得到的单一模型当成唯一可能的解释，就会丢掉这种“其实还有别的合理解释”的信息。更稳妥的做法是同时保留一组合格的候选模型，把它们看成一个“委员会”。\term{委员会查询}（query by committee, QBC）让委员会中的每个模型都对同一篇文档作出判断，并优先询问那些委员们意见不一致的文档——分歧正说明现有样本还不足以在这些合理解释之间作出裁决。

这里要严格限定“委员”的资格。所谓合格，指的是每个模型都与训练样本相容，即没有被所用的训练样本和拟合准则排除；这\emph{不}意味着可以随便抓几个训练失败、连已有样本都拟合不好的模型来凑数。一组本身就不合格的模型，它们的分歧毫无信息量。

委员会的成员从哪里来，决定了“分歧”到底在表达什么，因此构造方式不是无关紧要的细节。常见有三种来源：对训练集做重采样、用不同子集训练，考察的是结论对训练样本扰动的敏感程度；用不同随机初始化训练同一模型，探索的是不同优化路径会落到哪些解；贝叶斯式的后验采样则在给定概率模型下，直接表达模型参数本身的不确定性。要清醒地认识到，实践中常用的前两种构造只是对“模型可能有哪些差异”的近似，并不自动等价于完整的贝叶斯后验。

QBC的原始工作以“委员分歧最大”为标准选择查询，并在一个特定的理论模型下分析了这样选样能在多大程度上减少标签的不确定性\scite{seung1992querycommittee}。那个收益结论依赖其理论假设，不能无条件地搬到任意构造出的委员会上。

先考虑信息利用得最少的一种委员会：每个成员只报告它最终预测的类别，不报概率。为了把“意见分成了几派”转成一个可排序的分数，固定同一候选$x$、同一个$K$类空间和$M\geq2$个成员。第$m$个成员本来能给出完整概率$p_m(c\mid x)$，这里先只取它的硬预测$\widehat y_m(x)$，即概率最大的那个类别；若某个成员内部出现平票，应按事先固定的规则打破，不能临场随意决定。

统计各类别得了几票：令$v_c(x)=\sum_{m=1}^{M}\mathbb{1}\{\widehat y_m(x)=c\}$为投给类别$c$的成员数，其中$\mathbb{1}\{\cdot\}$在括号条件成立时取1、否则取0；$v_c(x)/M$就是投给类别$c$的成员比例。这组比例本身构成一个分布，对它取熵，就得到投票熵（vote entropy），衡量委员们硬预测的类别有多分散：
\begin{equation}
  u_{\mathrm{VE}}(x)
  =-\sum_{c=1}^{K}\frac{v_c(x)}{M}
    \log\frac{v_c(x)}{M}.
  \label{eq:knowledge-acquisition-human-vote-entropy}
\end{equation}
投票熵越大，说明成员们的硬类别越四分五裂，按这个准则就越该询问这篇文档。要注意它平均的对象是“成员的投票”，描述的是委员之间的分裂程度，而\emph{不是}在对“真实标签出错的概率”作平均——委员一致也可能一致地错。

只用硬票还有一个明显的信息损失：两个模型即使都投了同一类，把握程度也可能天差地别，一个概率只有$0.51$勉强胜出，另一个却是$0.99$高度确信，硬票把这种差别完全抹掉了。要利用完整概率，得先给所有成员设一个共同的比较基准，最自然的是把$M$个成员的概率等权平均：
\[
  \bar p(c\mid x)
  =\frac{1}{M}\sum_{m=1}^{M}p_m(c\mid x),
\]
这个平均分布代表委员会的“共识”，但它既不是某个新成员的预测，也不直接证明哪个成员更接近正确。有了共同基准，就可以衡量每个成员的概率与它相差多远：用成员自己的概率对“成员概率与共识之比的对数”加权，再在成员之间取平均，得到下式。
\begin{equation}
  u_{\mathrm{KL}}(x)
  =\frac{1}{M}\sum_{m=1}^{M}
    \sum_{c=1}^{K}p_m(c\mid x)
    \log\frac{p_m(c\mid x)}{\bar p(c\mid x)}.
  \label{eq:knowledge-acquisition-human-kl-disagreement}
\end{equation}
式~\eqref{eq:knowledge-acquisition-human-kl-disagreement}就是各成员到共识分布的平均Kullback--Leibler散度（KL散度），可读作“成员们的概率分配平均偏离共同基准有多远”：数值越大，委员会内部分歧越大。KL散度是衡量两个概率分布差异的标准工具，这里把每个成员与共识逐一比较后再平均。

使用这两个量时要统一口径：熵和散度都用自然对数，并且必须在同一个委员会、同一个类别空间下比较不同候选，否则分数不可比。计算中还有一个细节：当某一项概率为零时，按极限约定$0\log0=0$。正常情况下，只要有一个成员给某类别正概率，等权平均就必然也为正，因此“成员给正概率、平均却为零”在数学上不应发生；若真的出现，多半是数值计算出错或各成员的类别编号没有对齐，应当回头检查。

用三个成员、三个类别的教学例子，可以清楚区分“单个模型分散”和“委员之间分歧”这两个对象。

第一种情形，三个成员的概率都在$(0.40,0.35,0.25)$附近且都是第一类最大。每个成员自己的分布都相当分散（单模型熵较高），但它们的硬预测全部指向第一类，于是投票熵为零——委员们其实意见一致，只是每个委员都不太自信。第二种情形，三个成员分别给出$(0.70,0.20,0.10)$、$(0.20,0.70,0.10)$和$(0.20,0.10,0.70)$，每个成员自己都很自信，但它们各自押注了不同类别，每类恰好一票，投票熵达到这个“三成员、三类别”设定下的最大值。第三种情形，三者都给$(0.90,0.06,0.04)$，既高度自信又完全一致，单模型熵和委员分歧都很低。

对比前两种情形可见关键结论：单个模型的概率分散程度，与多个模型之间的意见分歧，是两个独立的判断——可能人人都犹豫却方向一致，也可能人人都自信却各执一词。因此不能只凭其中一个量去反推另一个。

委员会方法的代价也很实在：要利用这种分歧，必须训练并保存$M$个模型、对每个候选都预测$M$次，并在每轮新标签到来后重新训练或刷新这些成员，计算和存储开销大致随$M$增长。更要紧的是，分歧大并不保证问得准。如果成员们因为共享了相同的训练样本或相同的模型架构而带有共同的偏差，它们可能在某些文档上整齐地一起犯错，此时投票熵很低，看起来“很一致”，却是一致地错。同样，用不同初始化制造的差异只覆盖了优化层面的一部分不确定性，远不等于完整的模型参数后验，不能赋予它超出其实际含义的解释。

至于各种策略孰优孰劣，Settles与Craven在序列标注任务上比较过委员会、预期梯度长度等多种方法，发现它们的具体排序取决于任务、所用模型和成本的计算口径\scite{knowledge-acquisition-human-settles2008analysis}。因此不存在一个对所有情形普遍最优的准则，选择应结合自己的任务实测，这一点在后面几节还会反复出现。

\subsection{基于代表性与多样性的样本选择}
\label{sec:knowledge-acquisition-human-representative-diverse-selection}

不确定性采样有一个盲区：它只问“模型哪里不会”，却不管这些问题是否重复。假设候选池里有大量彼此近似的常见投诉，外加少量涉及不同部件的报告，那么只按不确定性挑，很可能一口气选回许多几乎雷同的文档——第一个标签学到的规律，后面几十个标签都在重复，人工被浪费在冗余上。

因此还要补上两个角度。其一是\term{代表性}（representativeness）：一个候选是否处在候选文档分布中“常见、典型”的高概率区域，优先它有助于用一个标签代表一大片相似文档。其二是\term{多样性}（diversity）：当我们一次选一批时，这批样本是否覆盖了彼此不同的情形，而不是挤在同一簇里。代表性看的是单个候选“离大众有多近”，多样性看的是一批候选“彼此有多不同”。

要判断相似，先要把文档变成可以计算距离的形式：用一个表示函数把文档$x$映射成向量$\bm{\phi}(x)$，再在向量空间上定义距离或相似度。这一步隐含一个前提风险：如果学出的表示主要反映的是文档长度、套用的模板这类表面特征，而不是故障的语义，那么算法所谓的“覆盖”只覆盖了长度和模板，并没有覆盖我们真正关心的故障类型。代表性和多样性的质量，首先取决于表示是否抓住了任务相关的结构。

先用代表性刻画“哪些区域最常见”。给定一个固定的非负相似度函数$K(\bm{\phi}(x),\bm{\phi}(z))$，数值越大表示两篇文档越相似；注意这里的$K$是相似度函数，与前面表示类别数的$K$只是符号相同、含义完全不同。对非空候选池$U_t$，把某篇文档$x$与池中每一篇$z$的相似度取平均，就得到它的局部密度分数：
\begin{equation}
  \rho(x)=\frac{1}{|U_t|}
  \sum_{z\in U_t}K\bigl(\bm{\phi}(x),\bm{\phi}(z)\bigr).
  \label{eq:knowledge-acquisition-human-density}
\end{equation}
$\rho(x)$高，说明池子里与$x$相似的候选很多，$x$处在一个密集区域，优先标注它比较容易覆盖一大片常见情形。但这也是一把双刃剑：一味追求高密度，会系统性地忽略那些数量稀少的类别。

密度衡量单个候选，可批量选择关心的是另一个问题：把“已有的标签”和“本批准备选的标签”合在一起后，池中的每一篇文档是否都能在附近找到一个相似的、已被标注的代表？为此记$X_t=\{x:(x,y)\in L_t\}$为已标注样本的预测输入，设本批$B$是$U_t$中一个大小为$b$的子集，$1\leq b\leq|U_t|$。衡量覆盖时，对池中每篇$z$，只取它在$X_t\cup B$中能找到的最高相似度（一个最近的代表就够了，多个相似代表不重复计分），再对全池求和。我们要选的，就是使这个总覆盖最大的那一批：
\begin{equation}
  B^\ast\in\argmax_{\substack{B\subseteq U_t\\|B|=b}}
  \sum_{z\in U_t}
  \max_{x\in X_t\cup B}
  K\bigl(\bm{\phi}(x),\bm{\phi}(z)\bigr).
  \label{eq:knowledge-acquisition-human-coverage}
\end{equation}
内层最大值让每个$z$只由最相似代表贡献一次覆盖，避免把多个代表与同一文档的相似度重复相加；外层求和衡量池内总体覆盖，而非最坏文档的覆盖保证。在教学设定中，某文档当前最大相似度为$0.8$，新样本与它相似度为$0.9$，该项增益是$\max(0.8,0.9)-0.8=0.1$；已有覆盖若为$0.95$，增益就变成$\max(0.95,0.9)-0.95=0$。随着已选集合增大，同一新样本的覆盖增益不会增加，这种边际收益递减称为\term{次模性}（submodularity）。所以每选一个样本都应重算它对当前$X_t\cup B$的新增覆盖；直接取密度最高的$b$个，仍可能集中在同一簇。

式~\eqref{eq:knowledge-acquisition-human-coverage}称为facility-location型覆盖目标，它用已选代表与其余文档的相似度评价一组样本。目标定义了怎样比较批次，还没有给出怎样找到最优批次。这类次模目标已用于训练样本子集选择和主动学习\scite{knowledge-acquisition-human-wei2015submodularity}

为计算每次加入样本的增益，我们把任意代表集合$S$的总覆盖记为$F(S)$，再减去固定已标集合的覆盖，得到批次新增覆盖$G(B)$：
\[
  F(S)=\sum_{z\in U_t}\max_{x\in S}
  K\bigl(\bm{\phi}(x),\bm{\phi}(z)\bigr),
  \qquad
  G(B)=F(X_t\cup B)-F(X_t).
\]
约定空代表集合时最大相似度为0，便有$G(\varnothing)=0$，即一个样本都没加时新增覆盖为零，这称为归一化。再看两条性质：增加代表绝不会让任何文档的最近相似度下降，所以$G$单调不减；结合前面观察到的边际收益递减，$G$就是一个归一化、单调的次模函数。

那么怎样求$B$？这个优化问题在一般形式下包含最大覆盖问题，而后者的精确求解是NP-hard的，即不存在已知的、对任意规模都高效的通用精确算法。实践中改用贪心近似：从$B=\varnothing$开始，每一轮在尚未选中的候选里，挑出使边际增益$G(B\cup\{x\})-G(B)$最大的那一个加入，直到选满$b$个。在相似度固定非负、且约束只限制批次大小（基数约束）的条件下，这个简单贪心有一个经典保证：它得到的覆盖增益至少不低于最优解的$1-1/e$（约$63\%$）\scite{knowledge-acquisition-human-nemhauser1978submodular}。但要注意，这个保证的对象仅仅是“相似度覆盖”这个代理目标，并不等于保证分类准确率提高——覆盖得更均匀，最终任务未必一定更好。

图~\ref{fig:knowledge-acquisition-human-selection-criteria}用同一组二维教学点，把前面三种选样依据并排放在一起，看清它们各自会挑中哪些点。靠近给定分类边界的点体现预测不确定性，落在高密度区域内部的点体现代表性，而分散在不同区域的一批点体现多样性。

图中右上角那个远离所有簇的孤立点值得专门讨论：不确定性和多样性都可能给它高分，但它既可能是一个珍贵的稀有类别样本，也可能只是无意义的异常点或噪声，单凭“距离远”无法分辨这两种身份，仍需结合可答性判断。另外要尊重池式查询的约束：我们只能选择池子里真实存在的候选。聚类算出的几何中心若并不对应任何一篇真实文档，它只能用来指引我们去附近找真实样本，而不能把这个虚拟中心点直接当成一篇文档交给人标注。

\begin{figure*}[htbp]
  \centering
  % 三种选样准则共享完全相同的教学点集、坐标范围和孤立点。
\begin{tikzpicture}[x={\dimexpr\linewidth/174\relax},y=1mm,font=\figurefont\footnotesize,text=BookInk]
  \tikzset{
    candidate/.style={circle,draw=BookMuted,fill=white,line width=.7pt,minimum size=3.2mm,inner sep=0pt},
    chosen/.style={circle,draw=FigureModel,fill=FigureModel,line width=1pt,minimum size=3.8mm,inner sep=0pt},
    panel title/.style={font=\figurefont\bfseries\footnotesize,anchor=south},
    criterion/.style={font=\figurefont\scriptsize,text=BookMuted,align=center},
    leader/.style={draw=BookMuted,line width=.65pt}
  }
  \def\candidatepoints{5/9,8/13,10/7,12/17,15/11,20/27,24/10,24/31,27/15,29/29,30/8,33/19,36/13,45/34}
  \foreach \shift/\panel/\title in {0/a/预测不确定性,61/b/代表性,122/c/批内多样性}{
    \begin{scope}[xshift=\shift mm]
      \node[panel title] at (26,46) {(\panel)\quad\title};
      \draw[BookRule,line width=.7pt] (2,5) rectangle (50,40);
      \foreach \x/\y in \candidatepoints{\node[candidate] at (\x,\y) {};}
      \node[criterion] at (26,1.5) {\strut 候选为空心，选中样本为实心};
    \end{scope}
  }
  \begin{scope}
    \draw[FigureLoss,dashed,line width=1pt] (22,5)--(22,40);
    \node[chosen] at (20,27) {};
    \node[chosen] at (24,10) {};
    \node[anchor=south west,text=FigureLoss] at (22.5,35.5) {给定边界};
    \node[anchor=north east,text=BookMuted] at (47,33) {孤立点};
  \end{scope}
  \begin{scope}[xshift=61mm]
    \draw[FigureInput!35,line width=.7pt] (10,11) ellipse (8 and 6);
    \draw[FigureInput!35,line width=.7pt] (31,13) ellipse (8 and 7);
    \draw[FigureInput!35,line width=.7pt] (25,30) circle (7);
    \node[chosen] at (8,13) {};
    \node[chosen] at (27,15) {};
    \node[chosen] at (24,31) {};
    \node[anchor=north east,text=BookMuted] at (47,33) {孤立点};
  \end{scope}
  \begin{scope}[xshift=122mm]
    \node[chosen] at (5,9) {};
    \node[chosen] at (29,29) {};
    \node[chosen] at (36,13) {};
    \draw[FigureLoss,line width=.8pt] (43,32)--(47,36) (43,36)--(47,32);
    \draw[leader] (38.5,35.5)--(43.2,34.3);
    \node[anchor=north east,text=BookMuted,align=right] at (38,39) {孤立点\\未必有用};
  \end{scope}
\end{tikzpicture}

  \caption{不确定性、代表性与多样性选择的差别。三面板共享教学点集与坐标范围，空心圆为候选，实心圆为池内选中样本。（a）按距给定虚线边界的接近程度选择；（b）从三个高密度区域各选一个实际候选，浅轮廓只标区域；（c）选覆盖不同区域的一批，并用孤立点说明距离远未必有用。坐标不是真实类别或性能证据。}
  \label{fig:knowledge-acquisition-human-selection-criteria}
\end{figure*}

这些依据并非只能单独使用，实践中常常组合，但组合的先后顺序会改变结果，不能随意。若先用可答性和密度把关、排除损坏文档与极端孤立点，再在剩下的合格候选里按不确定性排序，那么是“可答性和代表性”在决定谁有资格入围，不确定性只在入围者之间排序。反过来，若先保留高不确定候选、再在其中最大化批内覆盖，则变成由不确定性决定资格、多样性负责在其中挑一批，这时仍要对选出的结果独立检查可答性，因为不确定性本身不保证可答。

另一种常见做法是把几个分数线性加权成一个总分，但这要求先解决不同分数的量纲尺度问题，并用实验验证权重是否合理；随手给出两个权重，既不能证明普遍最优，也不能据此断言被排除的孤立点一定没有价值——它可能恰恰承载着稀有类别。

覆盖方法好看，但计算开销要如实核算，分两步看。在表示固定、且单次相似度计算成本给定的口径下，若显式算出所有候选两两之间的相似度，需要$O(|U_t|^2)$的时间和存储来构造一个$|U_t|\times|U_t|$的矩阵。在此基础上跑朴素贪心：每一轮都要扫描所有剩余候选、并对池中全部$z\in U_t$重算边际增益，最坏情况下总代价约为$O(b|U_t|^2)$，池子一大就很重。

工程上有若干加速手段：缓存每篇文档当前的最近相似度以避免重算、用“懒惰贪心”推迟那些增益明显不可能最大的候选的评估、或先粗筛掉一批候选再精选。这些方法能降低实际运行时间，却不改变问题本身的组合优化性质和贪心的近似保证。同时要警惕加速带来的偏差：近似最近邻或激进预筛可能悄悄漏掉位于稀有区域的样本——而那往往正是多样性想保住的部分，因此需要专门检查加速后损失了多少覆盖。

\subsection{基于预期收益的样本选择}
\label{sec:knowledge-acquisition-human-expected-utility-selection}

假设两篇文档同样不确定，却只有一篇的标签会影响大量待处理报告，单看当前概率无法分辨它们的后续价值。预期收益方法因而问：得到标签并更新模型后，目标误差或模型参数预计改变多少？标签现在未知，不能只挑最有利的答案来计算；必须逐个假设可能标签，再按查询前的模型概率取条件期望。预期误差降低比较目标分布上的损失，预期模型变化比较模型参数或梯度幅度。前者直接面对任务误差却计算昂贵，后者省去部分计算，却只是一种变化代理。以下先按返回可用且标签准确的理想条件计算，弃权、噪声和人工费用仍要另行纳入实际决策。

\subsubsection{基于预期误差降低的样本选择}
\label{sec:knowledge-acquisition-human-expected-error-reduction}

要比较“标注这篇之后模型会不会少出错”，思路上得先做一系列“假想训练”：假设这篇文档的标签是某个类别，用现有数据加上这条假想标签重新训练，看看得到的模型表现如何，再对每种可能标签都做一遍。这里回到分类任务，符号$K$重新表示类别数，不再是上一节那个相似度函数。

先固定学习算法和“如何从数据读出模型参数”的规则，把“喂给它一个数据集、它返回一组参数”这个复合映射记为$\mathcal{A}_{\theta}$。于是当前模型参数就是$\bm{\theta}^{(t)}=\mathcal{A}_{\theta}(L_t)$。对候选$x\in U_t$和每一个允许的类别$c$，只在一份\emph{临时}训练集里加入假想样本$(x,c)$，用完全相同的训练过程更新一次：
\begin{equation}
  \bm{\theta}^{x,c}
  =\mathcal{A}_{\theta}\bigl(L_t\cup\{(x,c)\}\bigr).
  \label{eq:knowledge-acquisition-human-hypothetical-update}
\end{equation}
每个$\bm{\theta}^{x,c}$回答的是一个“假如”问题：假如$x$的真实标签是$c$，会训练出哪一个模型。要强调这些只是假想分支，系统并不会真的把所有$(x,c)$同时塞进训练集——真实标签只有一个。

有了各个假想模型，还要在\emph{同一份材料}上比较它们的错误，而不是比较各自在训练集上的损失（在训练集上表现好不代表对未来有用）。取一份与实际使用场景相符、含$m>0$个样本的评价集$E=\{(\bm{z}_j,r_j)\}_{j=1}^{m}$，其中$\bm{z}_j$是预测输入、$r_j$是参考答案，$f_{\bm{\theta}}$是参数为$\bm{\theta}$的模型的预测函数。用损失函数$\ell$为每次预测与参考答案的偏差计分、越小越好，在同一份$E$上取平均，得到经验风险：
\begin{equation}
  \widehat R_E(\bm{\theta})
  =\frac{1}{m}\sum_{j=1}^{m}
  \ell\bigl(f_{\bm{\theta}}(\bm{z}_j),r_j\bigr).
  \label{eq:knowledge-acquisition-human-evaluation-risk}
\end{equation}
式~\eqref{eq:knowledge-acquisition-human-evaluation-risk}给每个假想模型一把相同的尺子。损失的选择决定了这把尺子量的是什么：取零一损失（预测错计1、对计0）时，平均损失就是错误率；取对数损失时，评价的则是概率预测的质量。两种尺度不能混用。一个现实困难是，$E$往往没有真实标签$r_j$可用，这时可以像Roy与McCallum的池式方法那样，改用当前模型或假想模型自己的预测分布来估计对数损失或零一损失\scite{knowledge-acquisition-human-roy2001error}。但这样做又额外引入了一个假设：模型给出的概率、以及候选池的构成，足以近似真实的评估分布；这个假设一旦不成立，估计就会系统性偏差。

单个分支的风险还不能决定该不该查，因为人工最终只会返回\emph{一个}真实标签，我们事先不知道是哪一个。合理的做法是把所有分支按其发生的可能性加权：用查询\emph{前}模型给出的$p_{\bm{\theta}^{(t)}}(c\mid x)$作为“标签恰为$c$”的估计概率，对各分支的风险求加权平均，得到查询后的预期风险；再用当前风险减去它，就是预计能降低的误差。\term{预期误差降低}（expected error reduction, EER）正是这项“预计减少的平均损失”，其分支汇总和查询分数写成
\begin{align}
  \overline R_E(x)
    &=\sum_{c=1}^{K}
      p_{\bm{\theta}^{(t)}}(c\mid x)
      \widehat R_E(\bm{\theta}^{x,c}),
  \label{eq:knowledge-acquisition-human-expected-future-risk}\\
  u_{\mathrm{EER}}(x)
    &=\widehat R_E(\bm{\theta}^{(t)})
      -\overline R_E(x).
  \label{eq:knowledge-acquisition-human-expected-error-reduction}
\end{align}
这一轮在$U_t$中挑选使$u_{\mathrm{EER}}$最大的候选。由于当前风险$\widehat R_E(\bm{\theta}^{(t)})$对所有候选都是同一个数，最大化$u_{\mathrm{EER}}$与最小化预期未来风险$\overline R_E$完全等价。保留“当前风险减去预期风险”这个形式仍有意义：它能区分“更新后预计风险依然很高”和“相对当前确实有下降”两种情况。分数还可能为负，含义是模型预计学到这个标签后损失反而上升——这通常提示该标签可能误导，或当前概率估计不可靠。最后务必记住，权重$p_{\bm{\theta}^{(t)}}(c\mid x)$只是模型在查询前对未知标签的猜测，并非真实观测到的频率；因此纸面算出的下降只是一种预测，不保证真实查询后一定兑现。

\begin{example}[可复算的预期误差降低]
  \label{ex:knowledge-acquisition-human-expected-error}
  设同一评价样本集和损失下，当前风险为$0.30$，以下均为二分类教学数值。$x_A$的标签概率是$0.6$和$0.4$，对应临时更新后的风险是$0.18$和$0.27$，所以$\overline R_E(x_A)=0.6\times0.18+0.4\times0.27=0.216$，降低量为$0.30-0.216=0.084$。$x_B$两类各占$0.5$，分支风险是$0.12$和$0.35$，得到$\overline R_E(x_B)=0.235$、降低量$0.065$。按EER应选$x_A$，虽然$x_B$有一个更有利的分支。决定排序的是所有未知标签分支的加权结果，不是最有利标签下的最低风险。
\end{example}

图~\ref{fig:knowledge-acquisition-human-expected-utility}把例~\ref{ex:knowledge-acquisition-human-expected-error}中$x_A$的两个互斥假设分开画出：每个假设各自临时训练、各自估计风险，最后才按概率加权。图上出现两份计算，纯粹是为了在决策前权衡两种可能，绝不意味着要向人索取两份互相矛盾的真值。

这个方法的主要代价是计算量：朴素实现下，$N$个候选、$K$个类别，需要做$NK$次完整的临时训练和评估，候选稍多就难以承受。预筛候选、只对部分标签抽样、用增量更新近似重训等手段都能降低开销，但每一种近似都在改变分数本身，因此最终仍应用“真实查询之后、在独立任务上观察到的效果”来检验，而不能只看预期值是否好看。

\begin{figure*}[htbp]
  \centering
  % 二分类候选的互斥标签分支，以及不做完整重训的期望梯度长度替代路径。
\begin{tikzpicture}[x={\dimexpr\linewidth/174\relax},y=1mm,font=\figurefont\footnotesize,text=BookInk,
  box/.style={draw=BookRule,fill=BookPaper,line width=.8pt,rounded corners=1.2mm,minimum height=11mm,align=center,inner sep=2.5pt},
  flow/.style={knowledge arrow,draw=BookTeal,line width=1.1pt,rounded corners=1.2mm},
  alternative/.style={knowledge arrow,draw=BookAmber,line width=.9pt,dashed,rounded corners=1.2mm},
  edge label/.style={font=\figurefont\scriptsize,fill=white,inner sep=1.3pt}]
  \node[box,draw=FigureInput,fill=FigureInput!10,minimum width=25mm] (candidate) at (15,25) {候选样本$x$};
  \node[box,minimum width=34mm] (update1) at (58,37) {假设$y=1$\\临时更新$\bm{\theta}^{x,1}$};
  \node[box,minimum width=28mm] (risk1) at (103,37) {估计风险\\$0.18$};
  \node[box,minimum width=34mm] (update0) at (58,14) {假设$y=0$\\临时更新$\bm{\theta}^{x,0}$};
  \node[box,minimum width=28mm] (risk0) at (103,14) {估计风险\\$0.27$};
  \node[box,draw=FigureOutput,fill=FigureOutput!7,minimum width=39mm] (aggregate) at (151,25) {期望风险$0.216$\\降低量$0.084$};
  \node[font=\figurefont\scriptsize,text=BookMuted] (current) at (151,43) {当前风险$0.30$};
  \draw[flow] (candidate.east)--node[edge label,above,pos=.46] {$p(y=1\mid x)=0.6$}(update1.west);
  \draw[flow] (candidate.east)--node[edge label,below,pos=.46] {$p(y=0\mid x)=0.4$}(update0.west);
  \draw[flow] (update1)--(risk1);
  \draw[flow] (update0)--(risk0);
  \draw[flow] (risk1.east)--(aggregate.west);
  \draw[flow] (risk0.east)--(aggregate.west);
  \draw[flow] (current)--node[edge label,right] {比较}(aggregate);
  \node[box,draw=FigureLoss,fill=FigureLoss!9,minimum width=45mm] (gradients) at (66,-8) {分别计算$\lVert\bm{g}_1(x)\rVert_2$\\与$\lVert\bm{g}_0(x)\rVert_2$};
  \node[box,draw=FigureLoss,fill=FigureLoss!9,minimum width=47mm] (egl) at (128,-8) {$0.6\lVert\bm{g}_1\rVert_2+0.4\lVert\bm{g}_0\rVert_2$};
  \draw[alternative] (candidate.south)|-node[edge label,below,pos=.4] {不作完整重训}(gradients.west);
  \draw[alternative] (gradients.east)--(egl.west);
\end{tikzpicture}

  \caption{标签未知时的预期查询收益。实线使用例~\ref{ex:knowledge-acquisition-human-expected-error}的二分类教学数值：互斥标签各自临时更新与评估，按$0.6$和$0.4$汇总成风险$0.216$，比当前$0.30$减少$0.084$。下方虚线是下一方法的替代计算，先取各假设标签的梯度范数，再按概率平均，不先平均可能相反的梯度方向。}
  \label{fig:knowledge-acquisition-human-expected-utility}
\end{figure*}

使用EER还要避开几个评价陷阱。第一，如果带标签的$E$被反复用于每一轮选样，它就成了被这个选择策略“用过”的验证证据，模型会逐渐朝它优化；因此最终验收必须另留一批完全独立的材料，否则会把对$E$的过拟合虚报成在新样本上的收益。第二，若用无标签池加模型概率来近似风险，估计的可靠性依赖模型概率和池子本身；池子的分布若与真实目标分布不同，风险估计可能系统性失真。第三，所有计算都只基于当前的$L_t$逐点独立进行，这会导致重复计数：两篇相似文档各自算出来都很有价值，不代表把它们一起加入时价值还能相加——第一篇的标签很可能已经厘清了第二篇原本待区分的大部分规律，第二篇的边际收益所剩无几。

\subsubsection{基于预期模型变化的样本选择}
\label{sec:knowledge-acquisition-human-expected-model-change}

若为每个候选和标签重训、再遍历评价样本集的成本过高，可以先估计这个标签会让当前模型参数移动多少。\term{预期模型变化}（expected model change）用预计更新幅度作为查询依据，它不是直接估计未来错误率。为在不完整重训的条件下取得可计算代理，固定当前模型参数与可微的单样本损失$\ell(f_{\bm{\theta}}(x),c)$，对每个假设类别$c$求梯度：
\begin{equation}
  \bm{g}_c(x)
  =\left.
    \nabla_{\bm{\theta}}
    \ell(f_{\bm{\theta}}(x),c)
    \right|_{\bm{\theta}=\bm{\theta}^{(t)}}.
  \label{eq:knowledge-acquisition-human-hypothetical-gradient}
\end{equation}
梯度$\bm{g}_c(x)$指向的是：在当前参数处，若把$x$的标签当作$c$，损失会沿哪个方向上升；普通梯度下降沿它的反方向更新参数。若用学习率$\eta>0$只做一步普通梯度更新，就可以用
\[
  \bm{\theta}^{x,c}\approx\bm{\theta}^{(t)}-\eta\bm{g}_c(x)
\]
来近似前面那个完整重训得到的$\bm{\theta}^{x,c}$。这一步移动的位移长度，由梯度的范数$\eta\lVert\bm{g}_c(x)\rVert_2$近似。

由于标签仍然未知，和EER一样要对所有假设类别取期望。关键次序是：\emph{先对每个分支分别取梯度范数（长度），再按概率平均}，由此得到\term{预期梯度长度}（expected gradient length, EGL），即预计这一条样本的单步梯度有多大：
\begin{equation}
  u_{\mathrm{EGL}}(x)
  =\sum_{c=1}^{K}
    p_{\bm{\theta}^{(t)}}(c\mid x)
    \lVert\bm{g}_c(x)\rVert_2.
  \label{eq:knowledge-acquisition-human-expected-gradient-length}
\end{equation}
在模型参数化方式、当前状态和学习率都固定时，$\eta u_{\mathrm{EGL}}(x)$就近似“参数一步位移的预期长度”，因此可在候选池中优先询问EGL高的样本。再次强调，它回答的是“预计更新幅度有多大”，而不是“这次更新是否降低错误”。Settles与Craven把这一思路用于序列标注，并讨论了对未知标签求期望的计算方式\scite{knowledge-acquisition-human-settles2008analysis}。

为什么“先取范数、再平均”的次序如此重要，不能反过来？如果先把各分支梯度按概率平均、最后才取范数，即计算$\lVert\sum_c p(c\mid x)\bm{g}_c(x)\rVert_2$，那么方向相反的梯度会在相加时互相抵消，最后衡量的是“平均\emph{方向}的长度”，而不是“各个分支平均会移动多远”。这两者差别很大：一个无论答案是什么都会让模型明显改变的样本，其两个相反方向的梯度平均后可能恰好为零，被错误地判定为“没有价值”。

用采用对数损失的Logistic二分类可以把这一点看得很具体。把文档特征向量末尾添一个常数1以容纳偏置项，记为增广特征$\widetilde{\bm{x}}$；其单样本梯度形如“模型预测概率与假设标签之差”再乘以该特征（推导可参考\BookChapterRef{chap:classification}{卷4《模型：经典模型》的“分类”章}）。当模型对该文档的预测恰为$p(y=1\mid x)=0.5$时，两个假设标签对应的梯度分别是$-0.5\widetilde{\bm{x}}$和$0.5\widetilde{\bm{x}}$，方向相反、长度相同。若先平均，两者抵消为零向量；但EGL先各自取长度再平均，得到$0.5\lVert\widetilde{\bm{x}}\rVert_2>0$——这正确反映了“无论真实标签是哪一个，模型都会发生一次非零更新”。

把EGL在Logistic下的表达式化简，还能暴露出“分数高未必更值得问”的另一个原因。令$p=p(y=1\mid x)$，两个分支的梯度范数分别是$(1-p)\lVert\widetilde{\bm{x}}\rVert_2$和$p\lVert\widetilde{\bm{x}}\rVert_2$，按概率$p$与$1-p$加权后可得
\[
  u_{\mathrm{EGL}}(x)=2p(1-p)\lVert\widetilde{\bm{x}}\rVert_2.
\]
这个乘积由两部分组成：$2p(1-p)$在$p=0.5$（恰在决策边界）时最大，刻画“标签的不确定性”；$\lVert\widetilde{\bm{x}}\rVert_2$则是样本特征自身的尺度。看一个教学对比：一个$p=0.5$、特征范数为1的边界样本，EGL为$0.5$；而一个$p=0.9$（模型其实很有把握）、但特征范数高达10的异常样本，EGL反而有$1.8$。后者得高分主要靠特征尺度大，并不代表它信息更丰富或更值得标注。这提醒我们：对特征做标准化、做梯度裁剪、或只选取部分模型参数来算范数，都会改变幅度的可比性，因此解释EGL分数前必须固定模型的参数化方式和当前优化状态，否则分数差异可能只是尺度差异。

EGL确实免去了为每个分支完整重训、以及反复扫描评价集的开销，但它仍需对每个候选、每个可能标签各做一次反向传播来求梯度，候选和类别很多时同样不便宜。更根本的局限在于近似本身：真实优化器常带有动量、自适应预条件（如Adam），甚至使用二阶信息，它们会利用优化器的内部状态改变实际步长和方向，此时简单的欧氏梯度范数只能粗略代理参数位移，并不能等同于这些优化器下模型真正移动的距离。因此和所有代理指标一样，最终是否成功仍要回到查询之后的独立任务损失或事先预定的业务指标上检验，而不能把“模型参数变化很大”本身当成目标达成。

\section{标注者的选择}
\label{sec:knowledge-acquisition-human-annotator-selection}

选出了值得标注的文档，还要决定把它交给谁。一篇直白写着“无法开机”的中文文档，和一篇需要核对专业检测报告的外文文档，显然未必适合交给同一个人：后者要求外语能力、相关领域知识，可能还要求有查阅检测系统的权限。标注者的领域专长、语言能力、对标注工具的熟练度、能否访问所需证据，以及他当前的工作负载，都会实实在在地影响每一篇文档能不能被正确作答。

因此“选谁”要分三步、由粗到细地推进：先用任务的硬性要求确定哪些人\emph{有资格}承接，再用他们的历史表现补充能力证据，最后把具体样本与合适的人员逐一匹配。本节解决的是“谁来做什么”；至于每个样本值得投入多少人次，是第~\ref{sec:knowledge-acquisition-human-budget}节的预算问题，而多个人的答案如何融合成一个结论，则在第~\ref{chap:knowledge-fusion}章处理，三者不要混为一谈。

\subsection{基于专业能力的标注者选择}
\label{sec:knowledge-acquisition-human-expertise-selection}

判断一个人能否承接某个任务，不能只看他是不是“某领域专家”这样的笼统身份。对一篇需要核对检测报告的文档，真正要逐项检查的是四种具体能力：能不能识别出目标对象，能不能理解任务用到的专业概念，能不能查到任务允许使用的证据，以及能不能按规范把结果正确录入。这四项任何一项不满足，标注都可能失败。

还要避免两个相反的想当然。一方面，领域资历只说明他具备进入这类任务的基础，并不自动证明他熟悉\emph{当前这个任务}划定的类别边界——一位资深工程师未必清楚本书任务如何区分“转述”与“亲历”。另一方面，经过培训的普通人员完全可能正确执行那些规则明确的常见判断，不必凡事故意上交专家；真正关键的是，当任务超出他的专业判断能力时，他应当能如实暴露这个缺口、选择弃权或转交，而不是硬猜一个结论。

这些能力不应靠印象判断，而要用培训例题和试标逐项验证。分界大致是：表述明确的普通故障，理解规范的受训人员就能处理；而需要判读电池热失控这类机理、并核查检测报告的文档，才真正要求相应的专业判断能力和访问权限。

准入的依据应当随着“辅助是否到位”而调整，而非固定在身份上。如果补上术语表、把检测报告改写成通俗的证据摘要之后，普通人员已经能依据规则正确判断，就不必仅凭“这是专业领域”而全部转给专家。但要承认辅助的边界：若任务本身要求在相互冲突的医学证据之间作出专业裁决，这种判断依赖长期专业训练，短期的操作培训无法替代。（这里临时换成医学例子，只是为了说明准入依据始终由“具体要求人做什么判断”决定，并不改变前面的故障任务设定。）

“非专家到底能不能用”没有一概而论的答案，必须放在具体任务和具体的答案聚合方式下验证。Snow等人在若干自然语言处理任务上，把非专家（如众包标注者）的结果与专家标准比较，结果显示在他们研究的那些任务里，汇集多份非专家答案可以达到可用水平\scite{knowledge-acquisition-human-snow2008nonexpert}。但这一结论有明确边界：它并不证明任意复杂的任务都能用非专家替代专家，也不意味着“多数人一致”就天然正确——多数人可能共享同一个误解。

还有一种情况会彻底改变“该找谁”：当任务目标从客观判断转为主观偏好测量时，所需的人群也随之改变。要了解用户喜欢哪种回答，就应该去问预定的服务对象本人、记录他们判断的分布，而不是请专家替目标人群表达偏好——专家的审美不等于用户的审美。

最后，人员准入本身也是一项要花钱、要持续维护的投入，培训、资格试题、以及规则更新后的重新确认，都应计入任务的准备成本，而不是默认免费。一个人通过了旧版本的资格考核，只代表他胜任过去那个版本的任务；当类别边界、标注工具或证据来源发生实质变化时，应当用受影响的例题重新考查，而不是让旧资格无限期有效。与此配套，界面必须给人一条正当的出口，允许他明确报告“这项判断超出我的专业能力”，系统据此触发转交。若界面强制作答、不允许弃权，人就会被迫猜测，结果是把一个已知的能力缺口伪装成一个看似确定的答案埋进数据，危害比留空更大。

\subsection{基于历史表现的标注者选择}
\label{sec:knowledge-acquisition-human-history-selection}

资格解决“能不能做”，历史表现则进一步回答“以往做得怎么样”。但直接拿两个人的总体符合率排序，很容易被任务难度污染。假设甲主要处理表述直白的普通故障，乙主要啃晦涩的技术报告，那么即使乙能力更强，他的总体符合率也可能更低——拉低他的是任务更难，而不是他水平更差。

为了让历史记录能支撑公平比较，除了第~\ref{chap:knowledge-acquisition-overview}章要求的通用来源和状态字段，还要额外关联三类信息：这条标注属于什么任务类型、是否有独立参考或后续裁决结论、以及花了多长时间。选人时，应优先比较两个人在\emph{相互重叠的同类任务}上的表现，或者按可以解释的类别、难度分组后再比。一个只做简单样本、总体比例很高的人，并不能据此被认定适合困难样本——他可能从未被难题考验过。

把这种分组比较量化。固定在同一个任务组$g$、同一版规范下，设标注者$a$完成了$n_{a,g}>0$项，其中有$k_{a,g}$项与独立参考一致，用比值
\[
  \widehat q_{a,g}=\frac{k_{a,g}}{n_{a,g}}
\]
描述他在这\emph{一组已完成任务}上与参考相符的比例。措辞上要克制：它只是对已有记录的描述，不能直接当成他在所有未来任务上的正确率——未来任务的构成可能不同。

同样重要的是比值背后的样本量。$9/10$和$90/100$的点估计都是$0.9$，但后者建立在十倍的证据上，可信程度明显更高；前者只是好运气就能达到。因此读历史表现时，还要把弃权率、超时情况和对应的规范版本并列来看，不能让一个孤零零的综合高分掩盖“某一组其实样本极少、证据不足”的事实。对这种比例如何做可靠性估计、如何量化其统计不确定性，第~\ref{chap:knowledge-quality-assessment-correction}章会专门展开。

这套方法对两类人会失灵，需要专门处理：一类是新加入者，根本没有历史；另一类是长期未参与、或任务已经变更的老成员，他们的旧结论不能直接沿用。对这两类人，可先分配少量\emph{校准材料}——这些材料有已知参考答案，且难度构成尽量贴近真实任务——用一小批结果建立初步能力估计；在积累到足够证据之前，限制他们承接高风险任务的数量。这样既给了机会，又避免凭偶然的几次表现就把人永久定级。

同时不能盲信“参考”本身。用于对照的参考答案也需要检查：如果多数标注者共享同一个对规范的误解，那么与多数人一致并不等于正确；如果参考集已经过时、或只包含常见类别，那么在它上面表现再好，也无法证明一个人具备处理困难区域的能力。参考的质量决定了历史表现评价的天花板。

\subsection{基于样本与标注者匹配的任务分配}
\label{sec:knowledge-acquisition-human-matching}

若本批包含不同语言和专业方向的文档，拥有合格人员还不等于每篇都有可用承接人。需要把逐样本资格与本轮工时写成可检查约束。对本轮要求分配的样本$i$和人员$a$，令$e_{i,a}\in\{0,1\}$表示语言、领域、权限和任务资格是否满足，$z_{i,a}\in\{0,1\}$表示是否实际分配。预计工时$t_{i,a}$与可用总工时$H_a$采用同一单位，例如小时；基本约束为
\begin{equation}
  z_{i,a}\leq e_{i,a},
  \qquad
  \sum_a z_{i,a}\geq1,
  \qquad
  \sum_i t_{i,a}z_{i,a}\leq H_a,
  \label{eq:knowledge-acquisition-human-assignment}
\end{equation}
第一项排除无资格分配，第二项使每个本轮样本至少有一名承接人，第三项避免预计工时超过个人容量。它们检验可行性，并不优化人员排名或保证得到答案；没有合格承接人或容量不足时，这组约束可能无解。允许弃权时，分配至少一次仍不等于取得可用标签，因此还要指定转交人员及等待上限，到期未解决则保留未知与原因。

图~\ref{fig:knowledge-acquisition-human-matching}用混合领域文档显示这项逐例匹配：普通中文材料交受训通用人员，医疗证据交医疗专家，英文合同交双语法律人员。这里更换领域是为了显出语言与专业资格的对应，不继承故障任务的类别规则。细线仅表示可承接，粗线才是实际分配，虚线箭头是通用人员遇专业难例后弃权转交；图中“余1项”等名额是按该教学任务预计耗时折算的负载提示。总体历史均分既不能表达这些资格，也不能表达专家只剩一个名额。

\begin{figure}[htbp]
  \centering
  % 样本需求与标注者能力的候选关系、实际分配和弃权转交。
\begin{tikzpicture}[x=1mm,y=1mm,font=\figurefont\footnotesize,text=BookInk,
  sample/.style={draw=FigureInput,fill=FigureInput!10,line width=.9pt,rounded corners=1.2mm,minimum width=38mm,minimum height=12mm,align=center,inner sep=2.5pt},
  worker/.style={draw=FigureModel,fill=FigureModel!10,line width=.9pt,rounded corners=1.2mm,minimum width=39mm,minimum height=12mm,align=center,inner sep=2.5pt},
  candidate/.style={draw=BookMuted!65,line width=.65pt},
  assigned/.style={draw=BookTeal,line width=1.5pt},
  transfer/.style={knowledge arrow,draw=BookAmber,line width=.9pt,dashed,rounded corners=1.2mm}]
  \node[font=\figurefont\bfseries\small] at (20,51) {样本需求};
  \node[font=\figurefont\bfseries\small] at (89,51) {标注者能力与负载};
  \node[sample] (s1) at (20,36) {普通中文文档\\低难度};
  \node[sample] (s2) at (20,16) {医疗证据片段\\需查阅术语};
  \node[sample] (s3) at (20,-4) {英文合同条款\\高难度};
  \node[worker] (w1) at (89,36) {受训通用标注者\\中文，余6项};
  \node[worker] (w2) at (89,16) {医疗领域专家\\中文，余1项};
  \node[worker] (w3) at (89,-4) {双语法律人员\\中英，余2项};
  \draw[candidate] (s1.east)--(w1.west);
  \draw[candidate] (s2.east)--(w1.west);
  \draw[candidate] (s2.east)--(w2.west);
  \draw[candidate] (s3.east)--(w3.west);
  \draw[assigned] (s1.east)--(w1.west);
  \draw[assigned] (s2.east)--(w2.west);
  \draw[assigned] (s3.east)--(w3.west);
  \draw[transfer] (w1.east)--(112,36)--node[left,align=right] {难例弃权\\转交}(112,16)--(w2.east);
  \draw[candidate] (8,-18)--(27,-18);
  \node[anchor=west] at (29,-18) {可承接};
  \draw[assigned] (58,-18)--(77,-18);
  \node[anchor=west] at (79,-18) {本轮分配};
  \draw[transfer] (8,-27)--(27,-27);
  \node[anchor=west] at (29,-27) {来源弃权后的转交};
\end{tikzpicture}

  \caption{样本需求与标注者能力的匹配。细实线表示具备承接条件，粗实线表示本轮分配，带箭头虚线表示来源弃权后的转交。领域、语言、难度和剩余名额均为教学假设；名额需按预计耗时换算为工时，连线数量不表示能力排名。}
  \label{fig:knowledge-acquisition-human-matching}
\end{figure}

常见样本构成稳定时，可以预设通用与专业队列；任务变化后，则应结合新答案、弃权和负载调整分工。处理较少的人员也需适量校准任务，以更新能力证据。若只把容易任务交给当前高分人员，并持续推迟难例，困难区域会长期缺少标签，历史比例也会越来越受分配方式影响。记录每次调整的依据和版本，才能区分人员能力变化与任务构成变化。

\section{标注预算的分配}
\label{sec:knowledge-acquisition-human-budget}

前面几节分别决定了问什么、怎么问、问谁，但还有一个更高层的选择：手头这笔预算，究竟应该花在哪几类行动上。可选的行动不止一种：可以询问全新的文档，也可以给已有文档追加一份独立判断，可以调用更贵的专家，可以先粗标一批、再挑其中一部分细化，还可以干脆先启动一轮模型更新、用新模型指导后续选择。剩余预算在这些行动之间如何分配、各类各做多少次，无法只凭某一个样本的打分来回答——这是一个对“整组行动配置”的决策。

为了能在完整配置之间比较，先把它形式化。设$z_k$为第$k$类行动执行的次数，它必须是非负整数（行动按次发生）；$c_k$为该行动单次的费用；$F(\bm{z})$为采用某个配置时需要一次性投入的固定费用（例如搭建某个标注流程的成本），所有费用统一用同一种货币或事先约定的费用单位计量。再考虑不止资金一类约束：资源$r$可以是某位专家、一组通用标注人员、或一个受限的计算队列，记行动$k$每执行一次占用资源$r$的小时数为$t_{k,r}$，该资源总容量为$H_r$。人工资源按人工工时计，不要把计算队列的机时与人工工时直接相加，二者性质不同。最后，用$\widehat V(\bm{z})$表示依据任务证据估计的、该配置带来的总价值，越大越好。于是预算分配写成一个规划问题：
\begin{equation}
  \begin{aligned}
    \max_{\bm{z}}\quad
      &\widehat V(\bm{z})\\
    \text{s.t.}\quad
      &F(\bm{z})+\sum_k c_k z_k\leq B,\\
      &\sum_k t_{k,r}z_k\leq H_r
      \quad\text{对每类受限资源$r$}.
  \end{aligned}
  \label{eq:knowledge-acquisition-human-budget-allocation}
\end{equation}
式~\eqref{eq:knowledge-acquisition-human-budget-allocation}说的是：在资金预算$B$和每类资源总工时的双重限制内，挑选总价值尽可能高的行动组合。价值$\widehat V$可以具体定义为“预计减少的标签不确定性”或某项预定的任务收益，但一旦选定口径就必须全程固定，否则不同配置无法公平比较。

读这个模型时要厘清几个容易混淆的量。固定费用只对真正启动了的配置收取，与按次计费的$c_kz_k$分开；实际投入的人次由$z_k$和每项行动各需几名人员共同决定，不能把“人次”“费用”“工时”揉成同一个数字。尤其要区分\emph{总工时}与\emph{交付期限}：两个相互独立的任务各需1小时，派两人并行，1小时即可交付，但总共消耗了2工时；而普通标注与专家复核各需1小时且必须先后进行，就算有两人也至少要2小时才能交付。可见同样的总工时，并行性和依赖关系会显著改变实际工期；要判断能否赶上期限，还必须另外考察并行安排、任务间依赖和等待时间，无法只凭总工时推出。

这个规划还提醒我们，各类行动的价值一般不能简单逐项相加。一次到位的培训可能让一大批原本无法处理的标注变得可执行，它的价值摊薄在后续许多样本上；而对同一份相似答案重复标注两次，第二次提供的往往只是冗余的同一判断依据，边际价值很低。固定启动费用的存在也会改变“是否值得开一个新批次”的决定：量太小，固定费摊不薄，反而不划算。单位费用换来的价值这类指标可以帮助做局部比较，却不保证把各局部最优拼起来就是完整组合的最优。下面分别剖析三对最常见的取舍：新增覆盖与重复核验、普通标注与专家标注、不同粒度与不同轮次。

\subsection{新增标注与重复标注的预算分配}
\label{sec:knowledge-acquisition-human-new-versus-repeat}

第一对取舍，是把下一个人力花在“新对象”还是“旧对象”上。假设手头已有一些答案存在争议的故障文档，同时还有一些部件类别完全没有标签。把下一人次用于一篇全新文档，增加的是\emph{对象覆盖}——让更多种类的情况进入已标集合；把它用于让另一个人重新检查一篇旧文档，增加的则是\emph{同一对象的独立证据}——用第二份互不影响的判断来核验或暴露分歧。两者买到的东西不同。

把这笔账记清楚。设样本$i$一共接受了$r_i$次判断，第$j$次的成本为$c_{i,j}$，则在它身上花的逐例费用是$\sum_{j=1}^{r_i}c_{i,j}$；整个项目的总费用还要再加上准备、培训、协调等与单篇无关的固定项。在总人次相同的前提下，让更多样本满足$r_i=1$（每篇只看一次）能扩大覆盖面；而让少数样本的$r_i>1$（重复看多次）则提供了观察分歧的机会，并在错误彼此独立等适当条件下，能通过综合多份判断降低随机错误。究竟偏向哪一端，应由真正关心的因素决定——是否存在大量尚未覆盖的遗漏类别、争议有多频繁、人员能力如何、以及错误标签对下游造成多大损失——而\emph{不}应由“标签条数”这种表面指标决定：重复标注并不增加对象数，却同样消耗人次。

\Needspace{65mm}
重复一份答案究竟能不能增加证据，关键取决于两个人的错误是否相互独立，这一点常被想当然。如果多人共享同一套错误的规范、或都带着相同偏见，第二份答案只会把同一个偏差再重复一遍，谈不上核验；如果后答者在作答前已经看到了前答者的结果，他的判断会被前者锚定，也不再是一份独立证据。Sheng等人对含噪标注的研究表明，在某些标注质量与成本条件下，重复确实能改善最终标签和训练出的模型，但这种收益并非总会出现；有选择地、针对高风险或高分歧样本追加判断，通常比对所有样本无差别重复更划算\scite{knowledge-acquisition-human-sheng2008another}。因此“一个样本看几次”本身应当成为一项预算决策，依据证据动态决定，而不是不假思索地套用“每例固定三人”这类普遍规则。

与“固定次数”相比，更合理的是自适应安排：固定次数执行起来简单，却对清晰样本和争议样本一视同仁地投入相同人力，造成浪费。自适应做法是先取得一份或两份相互独立的答案，只有当触发了事先规定的条件——如两人答案分歧、证据被标记为缺失、或该样本属于高风险类别——才追加下一份判断，同时规定好最高追加人次和停止后样本的去向。

追加之前务必先诊断“为什么需要追加”，因为不是所有问题都能靠多投一次票解决：若根因是证据缺失，正确动作是补材料；若根因是规范本身有缺陷，则应先修正规则、再按新规则重查。指望靠不断增加投票次数来消除这两类问题，只会空耗预算。最后要在预算中为下游两件事单独留出资源：多份答案如何融合成结论属于第~\ref{chap:knowledge-fusion}章，独立的质量审计属于第~\ref{chap:knowledge-quality-assessment-correction}章，它们各自都要消耗人力，不能默认免费。

\subsection{普通标注与专家标注的预算分配}
\label{sec:knowledge-acquisition-human-generalist-versus-expert}

第二对取舍，是用昂贵的专家还是便宜的普通标注者。面对“大部分是规则清楚的故障报告、夹杂少量专业难例”的典型构成，至少有三种组织方式：全部交给专家；先由普通人员处理、遇到困难再转交专家；或先请专家制定规范、培训普通人员，使更多样本能被普通人按规则解决，只把残余难例上交。

要在费用上公平比较，先要固定“$n$个样本都要达到同一完成要求”，并用统一的费用单位记三种单价：普通人员单次$c_{\mathrm{O}}$、专家单次$c_{\mathrm{E}}$、每次转交额外的协调成本$c_{\mathrm{X}}$。第一种方案全程用专家，固定启动费$F_{\mathrm{E}}$，总费用很直接：
\[
  C_{\mathrm{all}}=F_{\mathrm{E}}+nc_{\mathrm{E}}.
\]
第二种方案让普通人员先处理全部$n$篇，再把其中比例$q$转交给专家。这里有一笔容易漏掉的钱：普通人员处理过转交样本的费用\emph{不会}因为转交而退还，这篇等于被处理了两次。再加上路由流程的固定费用$F_{\mathrm{R}}$，总费用为
\begin{equation}
  C_{\mathrm{route}}
  =F_{\mathrm{R}}+nc_{\mathrm{O}}
   +qn(c_{\mathrm{E}}+c_{\mathrm{X}}).
  \label{eq:knowledge-acquisition-human-expert-routing-cost}
\end{equation}
式~\eqref{eq:knowledge-acquisition-human-expert-routing-cost}把普通判断、专家判断、转交协调三笔钱分开计入，避免把一个被转交的样本误算成只做过一次判断。第三种方案在前端投入专家精力：让专家先开发规范、再培训普通人员，期待培训后需要转交的比例下降。记开发与培训的固定费为$F_{\mathrm{D}}$、$F_{\mathrm{T}}$，培训后仍需转交的比例为$q'$（理想情况下$q'<q$），则
\begin{equation}
  C_{\mathrm{train}}
  =F_{\mathrm{D}}+F_{\mathrm{T}}+nc_{\mathrm{O}}
   +q'n(c_{\mathrm{E}}+c_{\mathrm{X}}).
  \label{eq:knowledge-acquisition-human-expert-training-cost}
\end{equation}
培训方案在前面增加了一大笔固定费，它只有在“转交减少、专家单价节省”等后续收益足以抵补这笔投入时，才可能真的更省钱。下面用一组教学数值具体核算，请注意$q'$能降到多少是假设、不是实测结论。设三种方案都让$n=100$篇文档在同一规范下得到可用的最终判断，普通人员无法完成的必须转交；取$c_{\mathrm{O}}=1$、$c_{\mathrm{E}}=4$、$c_{\mathrm{X}}=0.5$，全程专家$F_{\mathrm{E}}=10$，普通后转交方案$F_{\mathrm{R}}=20$、$q=0.20$，培训方案$F_{\mathrm{D}}+F_{\mathrm{T}}=70$、$q'=0.05$。代入上面各式，三个方案的费用依次是
\[
  C_{\mathrm{all}}=410,\qquad
  C_{\mathrm{route}}=210,\qquad
  C_{\mathrm{train}}=192.5.
\]
光比较钱还不够，专家的\emph{时间}往往是比预算更紧的约束，需要单独核算。仍设专家处理一例要12分钟。全程专家处理100例需要$100\times12/60=20$小时；转交方案只有$20$例到专家，用4小时；培训方案除了残余的$5$例（1小时）外，开发和培训还另占专家6小时，合计$6+5\times12/60=7$小时。于是出现一个重要结论：费用最低的培训方案（$192.5$）并不是占用专家时间最少的方案——转交方案只花专家4小时。钱和稀缺资源的最优配置可能指向不同选择。

所有这些数字都是教学假设，不能当成真实报价。尤其“培训能否真把转交比例从$20\%$压到$5\%$”必须靠试标验证，三种方案在同一完成要求下的标签质量也须另行测量；绝不能因为某个方案账面费用低，就直接推出它质量可接受或一定该选。

把专家时间约束一般化。记专家处理一例的预计工时为$t_{\mathrm{E}}$，其可用总容量为$H_{\mathrm{E}}$（两者同单位），则路由方案要可行必须满足$qnt_{\mathrm{E}}\leq H_{\mathrm{E}}$；培训方案还要把开发、培训占用的$h_{\mathrm{D}}$、$h_{\mathrm{T}}$算进去，满足$h_{\mathrm{D}}+h_{\mathrm{T}}+q'nt_{\mathrm{E}}\leq H_{\mathrm{E}}$。除了总量，转交\emph{门槛}的高低也是个操作痛点：门槛过低、什么都上交，专家队列会积压；门槛过高、要求普通人员尽量自己解决，又可能迫使他们在能力不足时硬猜。

还要看到专家两种工作的影响范围不同：专家亲自解决一个疑难样本，只影响这一个样本；而他写出一条可复用的判断规则，可能让一大批同类样本今后都能由普通人员解决，潜在价值大得多。但规则究竟能覆盖多少文档，同样要靠试标确认，不能在规划时就乐观地把开发费摊到全部样本、假定人人都会受益。

最后破除一个常见迷信：用了专家不等于标签就一定可靠。普通人员和专家可能共享同一个错误的任务定义——如果问题本身就问错了、或规范有系统性偏差，专家一样会在错误的框架里作答。因此专家身份不能替代独立证据，付费更高也不直接等于新增的可用标签更多。是否调用专家，应看专家在\emph{当前这个任务}上能带来多少增量价值，而非“贵的总没错”。当交付期限紧张时，决策还要加入时间维度：比较“干等专家”会耽误什么，与“先用普通人员处理其他可答文档”能换来多少进展，不能只盯着账面费用最小。

\subsection{不同标注粒度与轮次的预算分配}
\label{sec:knowledge-acquisition-human-granularity-round-budget}

第三对取舍，是“一次标到位”还是“先粗后细、分两轮”，以及“大批量少轮次”还是“小批量多轮次”。如果一个便宜的粗判断就能先排除大量明显无关的文档，那么昂贵的精细标注只需施加在剩下的少数报告上，直觉上能省钱。两阶段方案因此先为所有对象取得粗标签，再根据粗标签划定的范围，挑选其中值得细化的对象做精细标注。

但它在省下“未进入第二轮那部分细标费用”的同时，会引入一项新成本：第二轮重新打开材料、回忆并恢复上次判断上下文的开销，这在一次到位的方案里不存在。把两类成本都记全：设第一轮处理$n$个对象，粗标单价$c_{\mathrm{coarse}}$；第二轮只细化其中$m\leq n$个，每个要付重开费$c_{\mathrm{reopen}}$和细标费$c_{\mathrm{fine}}$；两轮各自有固定费$F_1,F_2$，涵盖相应的流程配置、调度和所需的模型更新。两阶段总费用为
\begin{equation}
  C_{\mathrm{two}}
  =F_1+F_2+nc_{\mathrm{coarse}}
   +m(c_{\mathrm{reopen}}+c_{\mathrm{fine}}).
  \label{eq:knowledge-acquisition-human-two-stage-cost}
\end{equation}
式~\eqref{eq:knowledge-acquisition-human-two-stage-cost}把粗标、细标、重开三笔钱都记上，避免把第二轮误以为“只需补填几个新字段”而漏算重开成本。作为对照，一次性完成全部细标签的费用是$C_{\mathrm{full}}=F_{\mathrm{full}}+nc_{\mathrm{full}}$。两者必须在\emph{同一个目标标注要求}下比较才公平。两阶段真正更便宜的条件是：第一轮确实排除了大量无须细化的对象（$m$显著小于$n$），而且重复阅读、对象重新对齐和模型更新的额外开销没有把省下的细标费用吞掉。

即使确定要分轮迭代，每轮选多少文档（批大小）仍会同时影响成本和反馈速度，是另一组权衡。大批量的好处是能摊薄每轮的启动、调度和模型重训成本，规模效应好；但它必须等整批全部标完，才能用上这批新标签更新模型，反馈周期长，期间用的一直是旧模型。小批量则能更早地把新标签反馈进模型、用更新后的依据指导后续选择，但批次数变多会重复支付启动与计算开销。

无论批大批次，每一轮结束都应记录同一套台账：投入了多少、新增的对象覆盖如何、有多少弃权及弃权原因、独立的质量证据、以及观测到的下游效果。停止项目也要有明确且不同的理由，至少能区分四种：预算耗尽、覆盖条件已满足、预定验证指标达标、以及再标下去的新增收益已经不足。这四种理由所需的证据完全不同，不能混用。停止时，对仍未解决的重要区域要保留其状态和后续去向，不能无声丢弃。特别要警惕一种似是而非的停止理由：模型自己变得“自信”了。自信只说明它的概率分布更集中，并不意味着它更正确——模型可能在错误上愈发笃定，因此单凭自信不能支持停止。

表~\ref{tab:knowledge-acquisition-human-fixed-budget}改用固定120个费用单位比较配置，区别于前面要求100篇都取得最终判断的比较。设普通完整判断1单位、专家4单位、粗判断0.5单位、局部细化2单位，专家每例12分钟；前三方案准备费20，两阶段两轮配置费30。为演示核算，教学场景设$c_{\mathrm{X}}=c_{\mathrm{reopen}}=0$，转交行总费$20+80\times1+5\times(4+0)=120$，两阶段行总费$30+100\times0.5+20\times(0+2)=120$。表中对象覆盖与所得粒度不同，同额费用不表示同等质量或标注价值。

\begin{table*}[htbp]
  \centering
  \small
  \begingroup
  \setlength{\tabcolsep}{3pt}
  \renewcommand{\arraystretch}{1.15}
  \begin{tabularx}{\linewidth}{@{}>{\RaggedRight\setlength{\parindent}{0pt}\arraybackslash}p{22mm}>{\RaggedRight\setlength{\parindent}{0pt}\arraybackslash}p{14mm}>{\RaggedRight\setlength{\parindent}{0pt}\arraybackslash}p{28mm}>{\RaggedRight\setlength{\parindent}{0pt}\arraybackslash}p{16mm}>{\RaggedRight\setlength{\parindent}{0pt}\arraybackslash}p{14mm}>{\RaggedRight\setlength{\parindent}{0pt}\arraybackslash}p{32mm}>{\RaggedRight\setlength{\parindent}{0pt}\arraybackslash}X@{}}
    \toprule
    方案 & 对象数 & 判断配置 & 专家工时 & 准备费 & 逐例费用 & 待核验 \\
    \midrule
    扩大新增覆盖 & 100 & 每例1次普通判断 & 0小时 & 20 & $100\times1=100$ & 歧义与稀有类别审计 \\
    增加重复判断 & 50 & 每例2次独立普通判断 & 0小时 & 20 & $50\times2\times1=100$ & 覆盖较窄；检查共性规范错误 \\
    普通后转专家 & 80 & 每例1次普通判断；其中5例追加专家 & 1小时 & 20 & $80\times1+5\times4=100$ & 转交门槛、专家增量价值 \\
    粗标后细化 & 100 & 每例1次粗标；其中20例再细化 & 0小时 & 30 & $100\times0.5+20\times2=90$ & 80例无细标签；对象对应与重开成本 \\
    \bottomrule
  \end{tabularx}
  \endgroup
  \caption{固定资源下的人工标注方案比较。对象数为新增覆盖数，不是判断人次；准备费加逐例费均为120个费用单位，专家工时另计。四行所得覆盖、重复证据和细粒度标注不同；单价及数值为教学假设，末列保留待核验需求，不预填准确率。}
  \label{tab:knowledge-acquisition-human-fixed-budget}
\end{table*}

比较完这些配置，可以把全章的三项核心决策重新连起来看。一篇文档值不值得去获取标签，由第~\ref{sec:knowledge-acquisition-human-sample-selection}节的选样依据回答；它是否有人具备作答条件、该交给谁，由第~\ref{sec:knowledge-acquisition-human-annotator-selection}节决定；而“是请一位专家、请两位普通人员、还是用同样的钱多覆盖几个对象”这类投入强度问题，则是本节预算分配的选择。

三者实际上相互牵连、需要联合决定，但在记录时应当把各自的依据分开保存。这样一旦效果不理想，才能回溯失败究竟出在哪个环节：是对象选错了（问了不该问的文档）、人员不匹配（交给了做不了的人）、还是投入不足（该重复或该请专家时舍不得花钱）。分开留痕也让未来的重新配置更容易——当费用或某类资源容量变化时，可以只调整受影响的那一环，而不必推倒重来。

\section{本章小结}
\label{sec:knowledge-acquisition-human-summary}

回顾全章，当人工资源有限时，我们做的并不是“尽可能多买标签”这一件事，而是在配置一套\emph{每个环节含义都明确、可追溯}的获取过程。任务定义和标注规范保证不同的人回答的是同一个问题；交互方式决定系统如实保存人实际取得的答案、约束或偏好，包括弃权；样本选择与人员匹配让那些真正值得询问的材料，有机会落到有能力作答的人手上；预算分配则在覆盖新对象、重复核验、调用专家和分轮迭代之间做平衡。贯穿始终的一条方法论是：比较任何两个方案，都要先固定相同的完成条件，并同时满足费用与工时两类约束；而能否赶上交付期限，还必须额外核查任务依赖、可并行程度和等待时间，无法由总工时直接推算。

把这套方法落回开篇的故障报告任务，可以拼出一个具体方案。依据试标结果，把规则清楚的常见文档交给经过培训的普通人员，用分层随机样本守住整体覆盖，再叠加不确定性与代表性准则挑选需要补充的材料；需要判读专业检测报告的文档转交给相应专家；对高风险或出现分歧的样本，按预先规定的人次上限追加独立判断；同时始终另留一批既不参与选样、也不参与调参的审计样本，供最终独立评估。遇到附件缺失的文档，正确动作是补齐材料，而不是简单地多派一个人去猜。若先做便宜的粗类别判断就能排除大量无关对象，那么只对剩下的报告定位精细证据可能确实省钱；但如果主要成本本来就在“读完全文”，分两轮反复打开文档反而更贵——是否分阶段，仍取决于成本实际发生在哪一步。

最后必须重申本章所有具体安排的条件性：它们是否成立，都要用当前任务自己的证据来支持，不能把这里的取舍当成任何场景都适用的通用配方。选样的高分可能只是在追逐噪声或无法作答的文档，委员会可能因为共享偏差而一致地犯错，模型建议可能悄悄影响人的独立判断，重复标注也可能只是把同一套规范误解复制多遍；而“操作成功完成”同样不等于答案已经核验。因此，交付每一条候选监督信号时，都应能回答五个可追溯的问题：问的是什么、谁依据什么材料作答、通过什么操作完成、哪些内容仍然未知、以及为此付出了多少成本。后续两章将在这个基础上，把多份人工答案与其他来源按共同语义融合（第~\ref{chap:knowledge-fusion}章），并用独立证据评估和修正质量（第~\ref{chap:knowledge-quality-assessment-correction}章）；在那之前，获取了多少标签、模型对自己有多自信，都不能被当成质量已经达标的结论。

\chapter{基于规则的知识获取}
\label{chap:rule-based-knowledge-acquisition}

第~\ref{chap:knowledge-acquisition-human}章讨论的依据来自人的逐篇判断，可靠但昂贵。是否存在一类可以重复执行、几乎不增加边际人力的判断依据？一个自然的想法是：把人判断时反复用到的条件直接写成程序规则，让机器自动执行。本章研究这种\emph{由规则提供依据}的知识获取方式，先从一个能暴露其全部难点的例子开始。

假设我们要从中文新闻中自动找出任职变动，并记录人员与组织之间的任职关系。最容易想到的规则是：句子里只要出现“任命”这个词，就判定它报告了任职变动。但拿三条新闻一测，这条规则会同时命中三种截然不同的情形：“海岬科技任命林舟为首席架构师”直接陈述了一次任职；“海岬科技未任命林舟”虽然出现了“任命”，却说的是\emph{没有}这回事；“关于任命林舟的传言已被否认”中，任命只是被转述的传言，而且随后被否认。三句话都含“任命”，却只有第一句真正支持目标事件。这个对照说明一个关键区分：程序在字符串里找到词语，只完成了\emph{字面检查}；要得到可用的标签，还必须再判断两件事——这条规则在当前语境下究竟\emph{应不应当触发}，以及一旦触发，推出的结论\emph{是否符合任务定义}。本章的大部分内容，都是在回答怎样把这两层判断也写进规则、并留下可供复核的依据。

要让上述判断稳定、可重复地执行，而不是每次都靠人临场决定，我们把判断所依据的东西显式写成\emph{规则}：有的规则规定“检查样本是否满足某些条件”，有的规定“去查询已有的实体、概念或关系事实”，还有的规定“最终答案必须满足哪些约束”。按这些事先写明的明确依据，由程序为样本产生任务标签、候选范围或关系标签，就是\term{规则标注}（rule-based annotation）。

要特别注意规则标注的产物地位：它产生的仍是\emph{候选}标签或约束，并不会因为出自程序就自动正确，因此每条结果都必须能说明“凭什么产生”，以便复核。理解规则时，最好把两个容易混淆的问题分开。一个是\emph{规则怎样形成}：规则文本由谁给出——可以是领域专家直接编写，可以是从大量样本的重复结构中归纳，也可以由模型先生成候选再由人确认。另一个是\emph{规则执行时依据什么产生标签}：是检查样本自身的模式，是关联外部已有的实体与事实，还是求解答案之间的约束。这两个维度相互独立：模型完全可以提出一条“去查询实体关系库事实”的规则（来源是模型、依据是外部事实），专家也可以亲手写一条“只比较数值阈值”的规则（来源是专家、依据是样本模式）。正因为来源和依据可以自由组合，二者必须分开记录，出错时才能定位到底是规则写错了，还是规则依赖的事实出了问题。

要把规则接进第~\ref{chap:knowledge-acquisition-overview}章那套含六个字段的样本标注，先要把规则执行的结果分得足够细，因为“没有得到标签”背后有三种完全不同的情况：句子里根本没有出现任命线索；否定语境过于复杂、程序无法可靠分析；以及这篇材料压根不是任务要处理的中文新闻。把三者都存成同一个空值，后续就无法决定是该补规则、弃权还是直接排除。

为此，一条规则在逻辑上拆成两个判断部分。其一是\emph{适用前提}，回答“这条规则是否适用于当前样本”，例如样本必须是中文新闻语体、事件时间在任务口径内，记为$A_v(x)$，其中$x$是样本、$v$是规则版本。其二是\emph{触发条件}，回答“在适用的前提下，样本是否命中了规则关心的模式”，记为$C_v(x)$；条件满足后如何拼出标签则由结果构造$O_v(x)$负责。

这里立刻遇到普通“是/否”布尔逻辑不够用的问题：除了条件成立、不成立，还常有\emph{暂时无法确定}，比如某个实体识别器没有输出，程序无从判断角色条件。因此本章用三值$\{\mathsf{T},\mathsf{F},\mathsf{U}\}$分别表示条件成立、不成立和无法确定。多个条件组合时，需要一种与三值配套的合取规则，称为\term{强Kleene合取}（strong Kleene conjunction）：只要有任一项为$\mathsf{F}$，合取就为$\mathsf{F}$（一个已确认不成立的条件足以否决）；只有所有项都为$\mathsf{T}$，合取才为$\mathsf{T}$；其余情形（没有明确假值、但至少一项未知）为$\mathsf{U}$。这样规定的好处是：一个已确认的假值能直接决定“不触发”，不会因为另一个条件恰好未知而被掩盖成“无法确定”。

令$D_v(x)=A_v(x)\land C_v(x)$表示适用前提与触发条件的合取。在一次条件检查\emph{成功执行完毕}之后，规则落入哪个分支由下式给出：
\begin{equation}
  b_v(x)=
  \begin{cases}
    \text{范围外}, & A_v(x)=\mathsf{F},\\
    \text{触发}, & D_v(x)=\mathsf{T},\\
    \text{未触发}, & A_v(x)\ne\mathsf{F}\ \text{且}\ D_v(x)=\mathsf{F},\\
    \text{无法确定}, & D_v(x)=\mathsf{U}.
  \end{cases}
  \label{eq:rule-based-knowledge-acquisition-interface}
\end{equation}
式中的$b_v(x)$只回答“这次检查走到了四个分支中的哪一个”，它本身并不代替那条完整的六字段样本标注，各分支还要翻译成相应的状态。

若适用前提$A_v(x)=\mathsf{F}$，说明对象根本不在规则的处理范围内，记为“范围外”：语义状态是“不适用”，来源动作为空——规则并没有尝试回答它。在适用前提未被否决的前提下，其余三种结果都作为\emph{具名触发结果}保存在取得条件字段$s_{\mathrm{cond}}$里，以便区分：$D_v(x)=\mathsf{T}$为“触发”，$D_v(x)=\mathsf{F}$为“未触发”，$D_v(x)=\mathsf{U}$为“无法确定”。对于那些“条件一满足就能直接拼出任务内容”的规则，三种状态对应不同动作：触发时来源动作为“作答”，给出$O_v(x)$，语义状态为已给出；未触发时规则只是说“我没看到目标模式”，因此没有来源动作、任务内容仍为未知（这一点对后文区分负类至关重要）；合取无法确定时，以“弃权”动作把内容保留为未知。

还要区分一类更复杂的规则：有些规则触发后并不立刻产出答案，而是进入一段后续求解（比如本章第三节要讲的约束推理）。对这类规则，“触发”只代表\emph{获准进入求解}，求解是否成功、最终来源动作是什么，必须另外记录，绝不能把“开始求解”误当成“已经作答”。

读者可能会把几种“没得到答案”的情形混为一谈，这里专门厘清执行状态。上述四个分支都建立在“条件检查这一\emph{动作本身}顺利跑完”的前提上，所以这一步的执行状态都是“成功”；但执行成功只表示程序正常工作，并不表示答案已被核验为真。如果后面还有求解步骤，求解有没有完成要单独保存其执行状态；而程序中途崩溃、网络中断属于技术故障，应记为“执行失败”，这不是规则在检查后主动选择弃权——二者绝不能混淆。

另一个贯穿全章的要点是负类。范围外、未触发、弃权这三种情况\emph{都没有证明样本属于负类}：规则没看到目标模式，不等于目标模式不存在。只有当一条规则明确检查负类定义、并据此返回一个负标签时，得到的样本标注才能当作负例。记住这一点，能避免把海量“规则没命中”的样本错误地当成反例喂给模型。

解决了“何时作答”，还要约定“规则返回的东西长什么样”，因为不同形态的结果需要的复核依据完全不同。规则结果$O_v(x)$至少有三种形态：它可以是一个\emph{候选标签}，比如“任职变动”；可以是一个\emph{候选范围}，只把答案缩小到$\{\text{管理岗位},\text{技术岗位}\}$这样的集合而不指定其一；还可以是一条\emph{关系约束}，比如“两个实体不得取同一身份”。

形态不同，复核时必须留下的证据也不同：靠样本自身匹配的规则，要留下命中的具体样本位置（哪个词、哪段字符）；靠查询外部事实的规则，要留下查到了哪条外部事实、来自哪个版本的库；靠推理的规则，则要留下完整的推理步骤。判断不出来时还要保存原因。表~\ref{tab:rule-based-knowledge-acquisition-record}沿用任命例子，把一条规则为了可复核而应记录的字段集中列出，读者可在后续三节中反复回到这张表。

\begin{table*}[htbp]
  \centering
  \begingroup
  \setlength{\tabcolsep}{4pt}
  \setlength{\parindent}{0pt}
  \setlength{\RaggedRightParindent}{0pt}
  \renewcommand{\arraystretch}{1.12}
  \begin{tabularx}{\linewidth}{>{\raggedright\arraybackslash}p{19mm}>{\raggedright\arraybackslash}p{32mm}>{\raggedright\arraybackslash}p{49mm}>{\raggedright\arraybackslash}X}
    \toprule
    字段 & 作用 & 教学取值 & 缺少时无法判断的问题 \\
    \midrule
    待检查对象 & 确定规则检查什么 & 一句中文新闻及其中的实体提及 & 规则究竟作用于整篇文档、句子还是实体对 \\
    适用前提 & 限定允许使用规则的范围 & 中文新闻语体，事件时间符合任务口径 & 对象在范围外，还是缺少判断范围所需的证据 \\
    触发条件 & 说明何时产生规则结果 & “任命”连接组织与人员，且不在否定或被否认的引语中 & 命中由哪个条件造成，新增例外改变了什么 \\
    规则结果 & 说明产生何种候选标签或约束 & 候选关系$(\text{林舟},\text{任职于},\text{海岬科技})$ & 返回的是类别、候选集合还是关系约束 \\
    无规则结果时的状态 & 区分范围、触发结果与动作 & 条件检查：范围外记不适用；未触发动作为\textemdash；合取无法确定时弃权 & 无任务答案是因为范围外、未触发，还是依据不足 \\
    命中证据 & 保留从待检查样本到规则结果的路径 & 关键词跨度、两个实体标识及局部依存路径 & 程序成功是否等于样本语义支持结论 \\
    适用范围及版本 & 支持更新与追溯 & 新闻语体；规则\texttt{appointment-v1}；实体关系库2026-09版 & 词表、解析器或实体属性或关系事实更新后哪些旧结果需要重算 \\
    \bottomrule
  \end{tabularx}
  \endgroup
  \caption{一条标注规则需要记录的内容。教学规则可以返回关系候选，也可以按任务改为类别或候选集合；字段的作用是把语义状态、来源动作、执行状态、规则依据和版本范围分开。}
  \label{tab:rule-based-knowledge-acquisition-record}
\end{table*}

沿任命例子继续想，规则为了得到结论实际上在求助三类不同的依据：它可能直接检查句子里的词语和角色（这是\emph{样本自身}就有的信息），可能去查明“林舟”对应现实中的哪个实体、岗位该归到哪个类别（这要借助\emph{样本之外已有的事实}），也可能不直接指定答案，而是规定几个候选答案之间必须满足什么关系（这是在\emph{答案空间}上做推理）。下面三节就按这三类依据展开：基于模式匹配的规则检查样本中的可观察条件，基于知识关联的规则使用已有的实体、概念与关系事实，基于约束推理的规则限制候选答案之间的合法组合。需要强调，一条真实规则常常同时使用这三类办法；但无论怎样组合，都必须为每一步单独留下依据，出错时才能判断问题究竟发生在匹配、关联还是求解环节。

% 标题、入口段与完整Snorkel边注共同检查剩余空间。
\Needspace{85mm}
\section{基于模式匹配的规则利用}
\label{sec:rule-based-knowledge-acquisition-pattern}

第一类依据完全来自样本自身。对一句待标注新闻，规则首先能利用的就是句中出现了哪些词语、它们在什么位置、数值是多少——这些都不需要求助样本之外的信息。规则可以检查这些可观察材料是否满足预先写好的条件，条件成立时给出候选结果；这种\emph{直接检查样本自身特征}的办法，称为\term{模式匹配}（pattern matching）。它是最基础、也最容易实现的一类规则依据。

模式匹配有一个被广泛使用的工程载体：Snorkel框架中的“标注函数”（labeling function）。一个标注函数就是一段可重复执行的程序，它依据字符串模式、启发式规则或外部资源，为送进来的样本给出一个类别标签，也可以在依据不足时明确放弃投票\scite{ratner2017}。把规则写成函数而非散落的脚本，好处是每条规则的触发条件、返回值和弃权行为都显式、可测试，也方便后续把成百上千条规则的结果统一融合（融合方法见第~\ref{chap:knowledge-fusion}章）。

这里要澄清一个术语差别，以免把Snorkel的约定和本章的状态字段错误对应。Snorkel原始工作把标注函数“不返回类别票”的所有情况统称为abstain（弃权），其中既包括“仔细看过但无法判断”，也包括“模式压根没匹配上”。但本章坚持把这两者分开：在条件检查阶段，\emph{明确未触发}（模式没命中、合取为假）时来源动作留空，只有\emph{无法确定而主动停止判断}时才记为弃权。第~\ref{chap:knowledge-acquisition-overview}章已经说明，“没看到目标”和“看到了但判断不了”是两件事，后续处理方式不同。

不过在送入标签融合模型的那一步，为了计算方便，可以把“未触发”和“主动弃权”这两种都没有类别票的结果，\emph{在融合层}映射成同一个“不投票”编码——前提是规则已成功执行、且对象确实处于共同适用范围内。必须强调这只是为融合产生的临时编码，它\emph{不覆盖}原始记录里的触发结果、来源动作、适用范围和弃权原因；这些字段仍要原样保留。而且，未请求、执行失败、适用范围外、或适用范围尚未对齐的结果，不属于“正常检查完但不投票”，必须另行处理，不能混进这个编码。如果将来还想分别为“未触发”和“主动弃权”建模、估计它们各自的发生规律，就需要在来源模型中保留条件状态并做相应扩展，此时融合层的单一“不投票”编码就不够用了。

在真正计算模式之前，通常还要做一道\emph{前处理}：把原始样本转换成可以直接检查的表示，例如把句子切分成词、把日志解析成字段值，等等；检查过程中命中的字符跨度、提取出的数值和各个中间条件的结果也都要保存下来。

之所以强调保存中间结果，是因为“规则没有给出可用答案”至少可能由三种不同障碍造成，必须能彼此区分。第一种是\emph{字段缺失}：前处理没能提供所需字段，条件根本无法求值，结果是无法确定。第二种是\emph{条件算出假值}：字段齐全、计算正常，但样本确实不满足条件，规则正常地未触发。第三种最隐蔽：\emph{条件成立却不支持目标语义}——字符串模式命中了，但命中处的真实含义并不是目标事件（正如“未任命”也含“任命”），这时产生的候选缺少语义依据，需要回头检查候选本身。前两种影响的是“触发/未触发/无法确定”的分支判断，第三种则发生在触发之后、针对已产生的候选；把它们一概记成“执行失败”会完全误导排查方向。下面先看符号条件和数值条件这两种最基本的条件，再看多个条件如何组合。它们全都依赖样本格式和前处理，一旦分词器或字段格式变化，原有规则必须重新验证。

\subsection{基于符号模式的规则利用}
\label{sec:rule-based-knowledge-acquisition-symbolic-pattern}

先看最直接的符号条件。一句新闻中的“任命”是一个清晰可观察的符号线索，但正如章首三例所示，只找到这个词还无法确定“谁任命了谁、任命是否真的发生”。要把判断做扎实，可以让符号条件由浅入深地分三层。

最浅的一层是\emph{关键词}：只检查某个词是否出现。它实现最简单，但表达力也最弱，无法约束词的位置和上下文。进一层是\emph{正则表达式}：它除了要求词出现，还能限制字符的相对位置和相邻形式，比如要求“任命”后面紧跟一个人名、前面是一个机构名。最深的一层是\emph{句式模板或局部结构模式}：它借助句法分析，检查若干成分（任命者、被任命者、谓词）是否以规定的语法关系同时出现，而不只是位置邻近。这三层用到的前处理越来越多，能区分的情形也越来越细。需要说明，这套思路并不限于文本：当检查对象从新闻换成日志或其他结构化观测时，被检查的“符号”可以变成某个字段的取值或某个事件类型标记，但判断依据仍然是\emph{当前样本内部的符号结构}，没有借助样本之外的事实——这正是模式匹配与下一节“知识关联”的分界。

把任命规则具体化。待检查样本限定为一个已经切分好的句子后，规则要识别句中的谓词及其参与者：谁是任命者（应当是一个组织）、谁是被任命者（应当是一个人员）。光有角色还不够，触发条件还要排除两种“假任命”：谓词不能受否定词支配（否则是“未任命”），任命也不能只出现在一句随后被报道者否认的引语里（否则是“传言被否认”）。在样本处于适用范围的前提下，只有所有条件都成立，规则才返回任职变动候选以及相应的关系三元组；只要有一个条件明确为假就不触发；没有明确假值、却仍有条件无法确定时才弃权。

以“海岬科技任命林舟为首席架构师”为例，一条可复核的规则不应只吐出一个“是”，而应同时记录：“任命”一词的字符跨度、“海岬科技”作为组织的跨度、“林舟”作为人员的跨度，以及三者在句法上的局部结构。有了这些记录，复核者才能回到原句逐一核对——这两个实体确实充当了规则声称的角色吗？谓词与它们的关系确实支持“任职”吗？如果规则只留下最终标签、不留命中位置，出错时就无从复查。

现在可以看清章首三例在多层条件下如何被区分。三条文本都能通过最浅层的关键词检查（都含“任命”），真正的差别在于任命是否被\emph{肯定地陈述}。加入“任命谓词不受否定词支配”这一条件后，直接否定句“海岬科技未任命林舟”被排除，因为否定词明确支配了谓词，该条件为假。但“关于任命林舟的传言已被否认”更麻烦：在引语\emph{内部}，词语可能仍是肯定形式，否定发生在引语\emph{之外}，由报道者对整段引述作出否认。因此规则不能只看句内否定词，还要识别引语的边界、并判断报道者是否否认了引语内容；当句法分析无法可靠确定这些时，对应条件应记为未知而非简单地判真或判假。

图~\ref{fig:rule-based-knowledge-acquisition-pattern-rule}把适用前提、谓词条件、角色条件和非否定条件的结果分别保存、再统一合取，读者应重点观察“未知项”和“假值”如何共同决定最终分支：一个明确假值足以使规则未触发，哪怕另有条件未知；只有在不存在假值时，未知才导致弃权。最后要破除一个想当然：增加条件确实会\emph{收窄}触发范围、减少明显误命中，但新条件本身（比如否定检测、引语识别）也可能判断错误，因此不能用“条件写得多”来论证“规则更准确”——条件数量与正确性之间没有等号。

\begin{figure*}[htbp]
  \centering
  \begin{tikzpicture}[
  x=1cm,
  y=1cm,
  font=\figurefont\footnotesize,
  panel/.style={draw=BookRule,line width=0.8pt,fill=white},
  text box/.style={draw=BookMuted,line width=0.8pt,fill=BookPaper,minimum height=10mm,text width=43mm,align=left,inner sep=2mm},
  condition/.style={draw=FigureModel,line width=1pt,fill=FigureModel!10,minimum height=10mm,text width=38mm,align=center,inner sep=1.5mm},
  output/.style={draw=FigureOutput,line width=1pt,fill=FigureOutput!7,minimum height=12mm,text width=39mm,align=center,inner sep=1.5mm},
  combine/.style={condition,draw=BookInk,fill=BookPaper},
  main/.style={knowledge arrow,draw=BookInk,line width=1.1pt,rounded corners=1mm}
]
  \fill[white] (0,0) rectangle (16.6,8);
  \draw[panel] (0,0) rectangle (5.15,8);
  \draw[panel] (5.4,0) rectangle (11.05,8);
  \draw[panel] (11.3,0) rectangle (16.6,8);

  \node[anchor=north west,font=\figurefont\bfseries] at (0.2,7.8) {(a) 相同字面命中};
  \node[text box,anchor=north west] at (0.25,7.1) {肯定：海岬科技\textbf{任命}林舟。\\\textcolor{BookTeal}{字面命中}};
  \node[text box,anchor=north west] at (0.25,4.95) {否定：海岬科技\textbf{未任命}林舟。\\\textcolor{BookTeal}{字面命中}};
  \node[text box,anchor=north west] at (0.25,2.8) {引用后否认：“\textbf{任命}林舟”的传言被否认。\\\textcolor{BookTeal}{字面命中}};
  \node[anchor=south,text=BookMuted,font=\figurefont\scriptsize,align=center] at (2.575,0.25) {匹配成功只说明找到了跨度};

  \node[anchor=north west,font=\figurefont\bfseries] at (5.6,7.8) {(b) 条件记录后统一合取};
  \node[condition] (scope) at (8.2,6.65) {$A_v$：适用前提\\记录$\mathsf{T}/\mathsf{F}/\mathsf{U}$};
  \node[condition] (keyword) at (8.2,5.35) {$K$：出现目标谓词\\记录$\mathsf{T}/\mathsf{F}/\mathsf{U}$};
  \node[condition] (roles) at (8.2,4.05) {$E$：组织与人员角色匹配\\记录$\mathsf{T}/\mathsf{F}/\mathsf{U}$};
  \node[condition] (context) at (8.2,2.75) {$N$：无否定或引语否认\\记录$\mathsf{T}/\mathsf{F}/\mathsf{U}$};
  \node[combine] (conjunction) at (8.2,1.0) {保留$A_v$并计算\\$D_v=A_v\land K\land E\land N$};
  \foreach \condition in {scope,keyword,roles,context} {
    \draw[main,-] (\condition.east) -- (10.7,0 |- \condition.east);
    \fill[BookInk] (10.7,0 |- \condition.east) circle (1pt);
  }
  \draw[main] (10.7,6.65) -- (10.7,1.95) -| (conjunction.north);

  \node[anchor=north west,font=\figurefont\bfseries] at (11.5,7.8) {(c) 分支与来源动作};
  \node[output] (triggered) at (14.1,6.65) {$D_v=\mathsf{T}$：触发、作答\\给出类别、范围或关系候选};
  \node[output,draw=BookMuted,fill=BookPaper] (untriggered) at (14.1,4.9) {$A_v\ne\mathsf{F},\ D_v=\mathsf{F}$\\未触发；动作：\textemdash\\任务内容未知};
  \node[output,draw=BookAmber,fill=BookAmber!8] (abstain) at (14.1,3.15) {$D_v=\mathsf{U}$：无法确定\\动作：弃权；任务内容未知};
  \node[output,draw=BookMuted,fill=white] (outside) at (14.1,1.4) {$A_v=\mathsf{F}$：范围外\\语义：不适用；动作：\textemdash};
  \coordinate (output-split) at (11.55,1.0);
  \draw[main,-] (conjunction.east) -- (output-split);
  \draw[main,-] (output-split) -- (11.55,6.65);
  \fill[BookInk] (output-split) circle (1pt);
  \foreach \result in {triggered,untriggered,abstain,outside} {
    \draw[main] (11.55,0 |- \result.west) -- (\result.west);
  }
\end{tikzpicture}

  \caption{从字面命中到带例外的组合规则。左面板中相同谓词出现在肯定、否定和被否认的引语里；中面板分别记录适用前提$A_v$、谓词条件$K$、角色条件$E$和非否定条件$N$，再计算$D_v=A_v\land K\land E\land N$，并列记录不要求并行执行。右面板按式~\eqref{eq:rule-based-knowledge-acquisition-interface}区分范围外、触发、未触发与无法确定；\textemdash{}表示来源动作为空。所有分支均以执行成功为前提。遇到未知仍须检查可独立求值的条件，例如$A_v=\mathsf{T}$且$(K,E,N)=(\mathsf{T},\mathsf{U},\mathsf{F})$时应未触发，而非弃权。}
  \label{fig:rule-based-knowledge-acquisition-pattern-rule}
\end{figure*}

前面逐层加码的条件，代价是对前处理的依赖也层层加深，这一点在维护时必须算清楚。字面关键词看似最简单，但它依赖文本规范化和分词是否正确，而且极易受多义词影响。带位置约束的正则表达式能规定局部顺序，却仍然难以处理否定成分与谓词分居两句的跨句否定。句式模板能检查实体角色和依存路径，可它又把可靠性押在了实体识别器和句法分析器的质量上。

由此得到一条规律：规则使用的语言结构越丰富，需要核验的中间结果就越多——关键词命中、分词边界、实体类型、依存弧，每一环都可能出错。这意味着维护一条复杂规则，绝不只是偶尔改改规则文本：词形变化、新的拼写方式、分词器升级、领域特有的句式、乃至解析器版本更新，都可能改变中间结果，从而让原本正确的规则静默失效。这些前处理的版本应像规则本身一样被记录和回归测试。

这些符号条件\emph{从哪里来}，是与“怎么执行”相连但不同的另一个问题，呼应章首对“规则来源”和“规则依据”的区分。专家可以凭借经验直接写出条件；也可以从语料中反复出现的结构归纳出模板；还可以让模型先生成一批候选条件。无论候选来自哪里，都不能直接上线，而要回到原始样本上同时检验支持情形和反例：它在哪些句子上正确触发、在哪些句子上误触发？误触发究竟是因为条件写得太宽，还是适用前提划得不对？

排查时尤其要利用保存下来的触发位置和解析结果，因为这能区分两类性质不同的错误：一类是\emph{规则解释错了}——前处理给出的结构是对的，但规则把这个结构错误地对应到了目标语义；另一类是\emph{前处理就错了}——句法分析或实体识别根本没给出正确结构，规则在错误的结构上判断，自然不会对。前者要改规则，后者要改或更换前处理，混为一谈就会修错地方。

最后要承认模式匹配的两类固有遗漏，即使规则本身执行得完全正确也无法避免。第一类是\emph{语义/时效上的过度触发}：规则准确匹配到了一则任命公告，但任务只收录“已经生效”的任职关系，而这则公告可能是尚未生效的任命预告；模式命中了，目标事实却还不成立，此时规则结果不符合任务定义，不能直接当成事实。第二类方向相反，是\emph{覆盖不足}：某句确实属于目标类，但因为句子写法特殊、解析依据不足，规则无法确定而弃权，于是这个本该作答的样本被漏掉了。

这两类问题提醒我们：规则执行正确（程序没报错、条件求值无误）与规则结果在任务意义上正确，是两个层面的事。前者属于本章的工程正确性，后者要靠第~\ref{chap:knowledge-quality-assessment-correction}章的统一质量检查来评估和修正，不能在写规则阶段就默认两者一致。

\subsection{基于数值条件的规则利用}
\label{sec:rule-based-knowledge-acquisition-numeric-pattern}

符号模式处理的是“出现什么词、结构如何”，但很多判断对象根本不是词语，而是一个有大小、可比较的量。要判断设备是否过热，检查新闻里有没有“烫”“热”这些字是靠不住的，真正该看的是传感器的温度读数。对于测量值、计数、比例或距离，规则先取得可比较的数值，再判断它是否落在指定阈值或区间内；这种办法称为\term{数值条件规则}（numeric-condition rule）。

必须从一开始就划清一条界限：数值比较是两个层面的事。“比较计算本身有没有算错”是一回事——82 是否大于 80，可以机械地、确定地验证；“这个阈值 80 对当前设备、当前任务是否合理”则是另一回事，它依赖物理特性和任务口径，无法靠比较运算自证。计算正确只能说明规则\emph{忠实于}给定的观测值和公式，完全不说明阈值选得对。这一节后半会专门处理阈值从哪来、何时需要怀疑它。

用一个教学规则看清数值判断里的各种边界。假设任务规定：设备读数不低于$80\,^\circ\mathrm{C}$时返回“高温候选”；落在$[78,80)\,^\circ\mathrm{C}$这个紧贴阈值的区间内时，因为测量误差可能让真值跨过临界线，规则选择弃权；低于$78\,^\circ\mathrm{C}$时这条规则不触发。每个观测值还应附带单位、采集时间和传感器状态，否则数值的含义不完整。

在其余前提都满足时：读数$82\,^\circ\mathrm{C}$明确触发；$79\,^\circ\mathrm{C}$落在缓冲区里，规则弃权；若连单位都缺失、无法完成有意义的比较，同样弃权而不是乱猜。单位换算是一个特别容易出错、又特别能说明“忠实观测值”重要性的环节：$176\,^\circ\mathrm{F}$换算成摄氏恰好是$80\,^\circ\mathrm{C}$，达到了\emph{包含端点}的下界$80$，所以应当触发。但如果程序偷懒、把数字$176$直接当成$176\,^\circ\mathrm{C}$去比较，比较运算照样会顺利执行（176当然大于80）、甚至也触发，可这个候选依据的是一个被错误理解的观测值——程序没报错，结论却不可信。这类“能跑但语义错”的故障，正是数值规则必须保存单位和换算过程的原因。

这个温度例子说明，一条数值规则远不止“写一个大于号”那么简单，它实际是一条处理链：先从原始观测\emph{计算出}可比较的特征，再\emph{统一单位或基准}，然后根据区间端点是包含还是排除，执行严格（$>$）或非严格（$\geq$）比较，最后把比较结果翻译成触发、未触发或无法确定。

当规则使用多个区间（如分档报警）时，还要明确两件事：每个端点归哪个区间（$80$ 到底算上一档还是下一档），以及区间之间是否允许重叠；重叠或留下空隙都会让某些数值对应多个或零个结果。对缺失值、传感器溢出（如读数恒为最大值）、以及根本不可比较的值，也要预先规定处理方式，不能让它们静默地落入某一档。例子中$[78,80)$这个“不确定缓冲区”不是随便画的：它应当来自测量误差范围、数值舍入范围或其他有证据支持的限制。如果误差来源未知，规则可以保守地弃权，但不能为了显得严谨，就把任意划定的一条缓冲带美其名曰“经过验证的置信区间”——没有误差依据的缓冲区只是主观裕量。

阈值的\emph{来源}决定了它的身份，也决定了当情况变化时该怎么对待它，这里有三种典型来源，不能混用。其一，阈值由任务定义直接规定，比如教学任务硬性规定$80\,\mathrm{^\circ C}$为高温线。这时阈值其实在\emph{定义标签口径}：什么叫“高温”完全由它说了算；一旦口径改变，按旧口径标注的样本就需要重新标注。其二，阈值来自某个领域标准（如行业安全规范）。这时要保存它适用的环境条件和标准版本，因为标准会修订、也只在规定环境下成立。其三，阈值是用少量样本“调”出来的经验值。这时必须保存用来调参的验证样本和整个调整过程，并且用\emph{没有参与调参}的新样本做独立检验，否则很容易把“在开发样本上恰好吻合”误当成“对新样本也成立”。

无论哪种来源，阈值都不是一劳永逸的。当监测对象的群体构成、采集设备型号或时间分布发生变化时，原本合适的阈值可能整体偏移，需要重新评估它对\emph{当前}对象是否仍然适用。

即使阈值数值已经确定，要拿到一个候选仍需完成一系列前置步骤：计算特征、查询参照基准、换算单位、处理缺失值，每一步都是潜在的错误来源。当不同地区、设备或人群使用不同基准时（比如不同的温度标准、不同的正常值范围），比较之前还必须先选对\emph{与观测设备和单位相匹配}的那个基准；拿 A 设备的读数去套 B 人群的基准，结论没有意义。

正因为错误可能来自链条上的不同环节——基准选错、特征算错、或把范围外的样本误交给规则——规则应在每一步都记录观测值、计算结果和所用基准版本。这样维护者面对一个错误候选时，才能沿链条定位是哪一环出了问题，而不是笼统地说“规则不准”。最后要把一项\emph{持续成本}计入：传感器会漂移、观测分布会随时间变化，因此对输入分布的监测不是一次性工作，而是规则上线后长期要做的维护。

\subsection{基于组合条件的规则利用}
\label{sec:rule-based-knowledge-acquisition-composite-pattern}

实际规则很少只检查单个条件，任命规则就是典型：它必须\emph{同时}确认谓词出现、实体角色正确、语境没有否定，才有充分依据返回候选。把这三项分别记为$K$（出现任命谓词）、$E$（实体角色匹配）、$N$（没有否定或被否认的引语），触发条件就是它们的合取
\[
  C_v(x)=K\land E\land N.
\]
各项仍然取章首规定的三值$\{\mathsf{T},\mathsf{F},\mathsf{U}\}$，并按强Kleene合取计算。这里值得重申三值在组合时的现实意义：如果实体识别器对某个片段根本没有预测输出，角色条件$E$是无法确定（$\mathsf{U}$），而\emph{不是}“实体类型不符合”（$\mathsf{F}$）——把“没看出来”误记成“看出来不是”，会错误地让规则从不触发，漏掉本应弃权待查的样本。

当任务里既有特殊情形又有一般情形时，多条规则之间还需要一套\emph{控制策略}：先试哪条？某条给出结果后还要不要继续看别的？这与单条规则内部的合取是两个层面的问题。

第一种组织是\term{规则列表}（rule list）：规则按优先级从上到下依次检查。本节采用一种保守策略——只有当某条规则被确定地判为“未触发”或“范围外”时，才继续往下试；一旦触发就作答并停止；一旦无法确定就弃权并停止。这种“遇未知即停”的保守性体现在：假设排在前面的特殊规则给出$D_v(x)=\mathsf{U}$，那么即使后面的一般规则本来能够触发，列表也\emph{不}采用一般规则的答案，因为特殊情形尚未被排除，贸然采用一般答案可能正好踩中特例。列表末尾若有一个“默认分支”想返回负类，条件必须非常严格：只有在列表的适用前提已确认、前面各项都\emph{确定地}未触发或范围外、且任务明定“这些都不命中即构成负类”时，默认分支才能给负标签；否则只能保留为未触发、动作为空。

第二种组织是\term{规则树}（rule tree）：沿不同条件分支向下走。本节同样规定，遇到未知时沿一个专门的弃权出口离开，而不是把未知默认成“否”。需要强调，列表和树在未知时“弃权并停止/走弃权出口”只是\emph{本节选择的控制策略}，并不是强Kleene合取本身强制规定的；换一种更激进的策略，未知处理可以不同。第三种组织则是互不设优先级的\term{规则集合}（rule set）：各条规则独立执行、互不短路，系统保留每一条的原始结果（包括各自的未知），把冲突与融合留给后面的多来源融合阶段。

逐项检查并不意味着每次都要把所有昂贵计算跑完。合取有一个可以利用的性质：一旦某一项已经确定为假，整个合取的结果就已被锁定为假，其余项不再影响结果，因此可以\emph{短路}——直接跳过它们。短路既有计算意义，也有正确性要求。

举例来说，若适用前提已确认成立、但最廉价的关键词检查$K$已经为假，那就不必再运行昂贵的句法解析去求$E$和$N$了，因为合取必为假，省下计算。但短路不能用错方向：当$K$为真、$E$未知时，规则\emph{还不能}弃权，因为$N$仍可独立求值，应当继续检查否定条件。具体地，若$(K,E,N)=(\mathsf{T},\mathsf{U},\mathsf{F})$，由于$N=\mathsf{F}$这个明确假值已经存在，合取为假，标注状态应是“未触发、动作为空”，而\emph{不是}弃权。请记住这个次序：\textbf{假值优先于未知}——只有在“不存在任何假值”的前提下，未知项才会让最终结果变成无法确定。短路时还要记录哪些条件是真正求过值、哪些只是被跳过，以便复核。

同一条“假值优先”的规则也适用于适用前提与触发条件之间，要注意两种情形的处理次序不同。如果适用前提未知、但触发条件已经确定为假，那么$D_v(x)$仍为假（前提未知不足以挽救一个已被否决的触发条件），可以确认“未触发”；但如果适用前提\emph{本身}已为假，则应优先按“范围外”处理，而不是“未触发”——前者说对象不归这条规则管，后者说对象归你管但没命中，语义不同。标注状态里应能区分一个条件究竟是：已经求出真值、仍然未知、还是因短路被跳过。

再看否定例外与正向模式冲突的情形：一旦否定条件$N$被确认不成立（即确实存在否定或否认），那么即使一般的词语模式命中了“任命”，规则也应记为未触发，同时把\emph{两处证据}都保留下来——既要保留“任命”的命中位置，也要保留否定依据，因为复核者需要看到这两者如何相互否决。

搭好一个多条件规则后，不能只看“它在开发样本上是否全对”就收工，还要逐项理解每个条件\emph{实际起了什么作用}。有两个简单有效的诊断动作。一是\emph{逐项删除}：每次去掉一个条件，看哪些新样本会因此进入触发范围，从而知道这个条件究竟在拦截什么。二是检查\emph{条件共现}：若两个条件几乎总是同时成立或同时不成立，其中一个可能是冗余的，白白增加维护成本和前处理依赖。

还要避免一个推理错误：给合取规则增加条件，通常确实会\emph{缩小}触发覆盖，但覆盖变窄绝不等于可靠性提高——新增条件自身也可能判断错误，它在挡住一部分误触发的同时，可能引入新的漏触发。对从样本归纳或由模型生成的条件组合，这些诊断同样必不可少；而且条件组合得越深、前处理链条越长，需要调试的路径和依赖就越多，复杂度本身就是一种成本。

最后把三个层次的“决策”理清，避免混用：规则列表/规则树的控制策略决定\emph{多条规则或多个分支之间}如何选择结果；合取短路决定\emph{同一条规则内部}哪些计算可以安全跳过；而规则集合中相互独立的规则则是另一回事——它们应当各自保留结果、不做选择。

这一区分在多规则冲突时尤其重要。如果多条独立规则对同一样本给出不同标签，而实现时简单地“以后执行的覆盖先执行的”或随手丢弃某些结果，冲突就被悄悄隐藏了，事后无从判断哪个标签其实更有依据。正确做法是：在规则层老老实实保留每一条规则的版本、证据和原始结果（包括分歧），不提前裁决；如何把这些可能冲突的标签综合成最终结论，是第~\ref{chap:knowledge-fusion}章多来源融合要解决的问题，规则层不应越俎代庖。

\section{基于知识关联的规则利用}
\label{sec:rule-based-knowledge-acquisition-association}

模式匹配的全部依据都在样本内部，但有些判断单凭样本里的文字怎么都做不出来。在任命新闻里，就算程序准确地圈出了“林舟”这两个字，也无法知道他指的是现实中哪一位同名的人——叫这个名字的作者可能有好几位；就算确认了头衔是“首席架构师”，也还要知道任务采用的岗位体系把这个头衔归到哪一类。这里缺的不是“再写一个更精巧的字面模式”，因为无论模式怎么写，样本里根本没有区分两位“林舟”的信息。缺的是\emph{样本与外部已有知识之间的一座桥}。

这就引出第二类依据：先确定样本中的提及对应到知识库中的哪个已有实体、概念或事实，再利用这层对应去产生候选标签或约束，称为\term{外部事实关联}（external fact association）。它与模式匹配的根本区别在于：模式匹配回答“样本自身满足什么可观察条件”，而外部事实关联回答“样本里的东西对应到了外部世界的哪个对象，以及那个对象有哪些已知属性”。

使用外部事实时，初学者容易以为“查到了就可以直接用”，其实查到之后还有两个必须分开回答的问题。第一个问题是这项事实\emph{本身是否成立、是否仍然有效}：实体的属性可能已经过时（人已离职、机构已更名），知识库中的关系可能有时间有效期。第二个问题是这项事实\emph{能否用于当前这个样本}：即便事实在知识库中为真，它也未必在当前这句话里被表达——两个人真实存在雇佣关系，不等于每一句同时提到他们的话都在谈论雇佣。把“事实为真”错误地等同于“本句表达了该事实”，正是后面远程监督要警惕的问题。

外部事实常组织成\term{本体}（ontology）：一个明确记录概念及其相互关系（如上下位、从属）的体系。查询本体时必须核对它采用的分类口径——同一个岗位词，在不同岗位体系里可能落到不同位置。表~\ref{tab:rule-based-knowledge-acquisition-association}按\emph{实体、概念、关系}三种查询对象，分别比较需要建立什么对应、以及错误容易发生在哪一步。需要说明，一条规则可以同时使用多种关联，但它们之间没有固定的先后顺序，不必机械地“先实体、再概念、后关系”，应按样本和任务的实际需要组织。

\begin{table*}[htbp]
  \centering
  \begingroup
  \setlength{\tabcolsep}{2.5pt}
  \setlength{\parindent}{0pt}
  \setlength{\RaggedRightParindent}{0pt}
  \renewcommand{\arraystretch}{1.12}
  \begin{tabularx}{\linewidth}{>{\raggedright\arraybackslash}p{17mm}>{\raggedright\arraybackslash}p{25mm}>{\raggedright\arraybackslash}p{26mm}>{\raggedright\arraybackslash}p{30mm}>{\raggedright\arraybackslash}p{29mm}>{\raggedright\arraybackslash}X}
    \toprule
    关联类型 & 样本线索 & 查询对象 & 对应条件 & 候选标签或约束 & 主要歧义 \\
    \midrule
    实体匹配 & 名称、别名、标识及上下文 & 实体标识及属性 & 提及指向该实体；属性适合当前时间与任务 & 类别、属性或查询入口 & 同名未必同实体；无候选与多候选有别 \\
    概念关联 & 词语、短语及局部用法 & 词典、本体或分类概念 & 表达指向该概念；目标属性允许沿层级继承 & 概念标签、上位类或候选范围 & 相似未必同义；上位类不能确定唯一细类 \\
    关系匹配 & 已消歧实体对及证据范围 & 有方向和时间的关系事实 & 实体与关系定义对齐；句级需表达证据，组级需另立假设 & 句级关系候选或提及组监督 & 共现未必表达；缺失未必为假 \\
    \bottomrule
  \end{tabularx}
  \endgroup
  \caption{三类外部事实关联的待检查样本、依据与错误位置。表中每一行都要求同时保存样本侧证据和实体、概念与关系库版本，避免只留下最终标签而无法判断错误发生在哪一步。}
  \label{tab:rule-based-knowledge-acquisition-association}
\end{table*}

\subsection{基于实体匹配的规则利用}
\label{sec:rule-based-knowledge-acquisition-entity-matching}

要查询新闻中“林舟”的属性，前提是先确定这个名字到底指谁——这是一切后续查询的入口。本章用\term{实体匹配}（entity matching）泛指把样本中的对象与已有实体对应起来的过程；在文本场景下，它的具体操作称为\term{实体链接}（entity linking），即把文中的一个提及（mention）对应到知识库里那个有唯一标识的实体对象。

实体链接带来一个关键转变：匹配成功后，规则处理的对象从一段有歧义的文字跨度，变成了一个可以直接查询的实体标识。因此它绝不能退化成单纯的字符串比较——“林舟”这两个字相同，不代表指向同一人。实体链接可利用的样本线索包括提及本身的文字跨度、名称、别名、显式标识符（如工号、网址）以及周围上下文；它的结果有四种可能，且含义各不相同：唯一确定一个实体、得到多个暂时无法区分的候选、一个候选都没有、或证据不足而弃权。

用一个同名的教学例子说明消歧如何进行。设实体库里有两位都叫“林舟”的作者：\texttt{E-017}的属性与海岬科技相关，\texttt{E-204}的属性与东岭大学相关。对“林舟介绍了海岬科技的新系统”，规则分两步走：先按名字把两个实体都取为候选，再借助额外证据\emph{消歧}——句中出现了海岬科技、或文档带有指向某作者主页的标识、或来源本身就与某机构相关，据此判断这个提及应当对应哪位作者。要注意，句中出现海岬科技只是一条\emph{待检查的线索}，“名字相同”这一点本身完全没有确定身份，不能跳过证据直接拍板。

如果上下文极少，只有一句“林舟发表了报告”，没有任何能区分两位的线索，规则就应当老实地保留候选集合$\{\texttt{E-017},\texttt{E-204}\}$，而不是任意挑一个——猜错的代价远大于保留未知。还要区分两种“拿不到唯一实体”的不同情况：库里\emph{一个候选都没有}，说明知识库可能缺这个实体，需要补充实体或检索线索；库里\emph{有多个候选却无法区分}，说明实体存在但证据不足，需要补充能辨别身份的上下文。两种情况的补救方向不同，应分别保存。

除了用上下文消歧，还有几种常见的链接手段，各自的可靠性边界不同。如果提及本身带有序号、主页链接这类\emph{唯一标识}，按标识查询几乎可以直接确定实体，是最可靠的情形。别名表把简称、旧名都作为查询入口，能召回更多写法，但别名往往有时间性（机构曾用名、人物婚前姓），需要记录名称在什么时间段有效。模糊字符串匹配放宽了对拼写完全一致的要求，能容忍错别字和写法变体，但放得越宽，越可能把“名字相近、其实不同”的两个实体混为一谈。模型可以利用更广的上下文辅助判断该选哪个候选，但它承担的仍然是“在候选中选择关联对象”这一步，不能替代对最终链接的核验。

把实体链接拆成两个阶段在排错时特别有用，这也是实体链接研究中的标准划分\scite{knowledge-acquisition-milne2008}：先做\emph{候选生成}（按名字、别名等从库里检索出可能的实体集合），再做\emph{基于上下文的候选选择}（从中挑一个或保留多个）。这样拆开后，错误就能对号入座：最终链接错了，可能是因为正确实体压根没进入候选集（召回失败），也可能是它在候选集里、却被选错了（排序失败），两者的改进方法完全不同。

链接到实体只是手段，查到属性之后还要再做一次“属性—任务”对照，否则容易把看似相关实则不对口的属性用错。例如实体\texttt{E-017}的职业属性写着“软件工程师”，但任务要回答的是“文档所述那个时间点上他担任什么职务”；“软件工程师”是一个相对静态的职业身份，并不等同于他在特定时间、特定机构的具体任职关系，二者不能互相替代。

这一步要能区分三类不同来源的错误：其一是\emph{事实本身有误}（库里的属性就是错的）；其二是\emph{链接错误}（属性是对的，但提及被接到了错误的实体上）；其三是\emph{属性不适用}（实体对、属性也对，但这条属性不是当前任务或当前时点所问的那个）。为了事后能区分三者，规则应保存实体库版本、属性的生效时间、当初的候选列表、消歧所用证据以及最终选择，形成一条完整可回溯的链接链。

实体词表也需要持续维护。新别名可以由专家补充、从历史匹配样本中归纳，或采纳模型的建议，但每加入一个别名都要检查两件事：它会不会其实是另一个实体的名字（导致两个实体被错误合并），以及它是否只在特定时间段有效。改动别名后，还应回看旧的标注结果，确认这次改动会影响哪些历史链接、是否需要重算。

评估实体链接质量时，要按实体分组、并分别盯住前述两个阶段，因为不同实体的“知名度”差异很大。对每个目标实体组，分别检查\emph{候选召回}（正确实体到底有没有进入候选集）和\emph{选择错误}（在候选里是否选错了身份）。一个重要的反直觉点是：关于某个实体的材料多，既不保证它的召回率更高，也不保证消歧更准——材料多有时反而带来更多歧义别名。因此，在资料充分的热门实体上测得的好结果，不能直接外推到资料稀少的冷门实体，后者往往才是错误集中的地方。

\subsection{基于概念关联的规则利用}
\label{sec:rule-based-knowledge-acquisition-concept-association}

实体匹配解决“这个名字指谁”，但新闻里的“首席架构师”还面临另一个独立问题：它在任务规定的岗位分类里属于哪一类？这与确定“林舟”是哪位作者是\emph{不同}的问题——判断头衔的岗位类别，并不需要先确定戴这个头衔的人的唯一身份，哪怕有两位同名的林舟，“首席架构师”这个头衔的归类方式也是一样的。

\term{概念关联}（concept association）处理的就是这种“词义到类别”的对应：把文本中的词语或短语对应到词典、本体或某个分类体系中的概念标识，然后依据该体系明确规定的上下位（从属）关系，推出类别或缩小候选范围。这里有一个清晰的分工：\emph{文字}只负责在映射表里找到对应概念，真正的\emph{推导}则依据概念标识与概念之间在体系中的关系完成，而不是再靠字面猜测。

沿概念层级推导时，\textbf{方向}至关重要，最容易出错的就是把上下位关系反着用。假设教学岗位体系把“首席架构师”对应到细概念\texttt{role:\allowbreak{}technical-\allowbreak{}lead}，又规定该概念从属于上位类\texttt{role:\allowbreak{}management-\allowbreak{}position}。那么当文本确认了“首席架构师”这个细类后，规则可以\emph{沿从属关系向上}推出“管理岗位”这一上位类别，这个推导是成立的。

但反过来不成立：只知道某人属于“管理岗位”，绝不能反推出他就是“首席架构师”，因为管理岗位这个上位类之下还包含技术总监、产品负责人等许多其他岗位——上位类信息不足以唯一确定下位类。这个方向差别直接决定了规则能输出什么：做\emph{粗类别}判断的任务，上位概念就够用，可以直接给类；做\emph{细类别}判断的任务，若缺乏进一步证据，只能保留下位候选集合、标为无法唯一确定，不能硬选一个。还要注意，这里的归类结论依据的是\emph{指定的这套教学体系}，换成另一套岗位划分方式结果可能不同，不能把它当成对所有岗位体系的普遍断言。

建立“词语到概念”映射时，不能把词与词之间的各种关系混为一谈，因为它们允许的操作不同。在当前用法下确属同义的词，可以指向同一个概念；缩写词需要结合领域和上下文才能确定展开成哪个概念（同一个缩写在不同领域可能对应完全不同的词组）。尤其要防止两种过度合并：\emph{近义词}只是含义相近，在任务要求精确区分时不足以认定为同一概念；\emph{上下位词}（如“工程师”与“首席架构师”）之间存在层级方向，更不能合并成一个概念而丢掉从属关系。

WordNet是这类资源的经典范例：它先把严格意义上的同义词组织成一个个“同义词集”（synset），每个同义词集对应一个概念，再用上位、下位等语义关系把这些概念连接起来\scite{knowledge-acquisition-miller1995}。这种“词面集合”与“概念关系”分开记录的做法值得借鉴——它迫使使用者区分“两个词能不能互换”和“两个概念谁从属于谁”这两类问题。

为了让归类结果可追查，三类知识应当\emph{分开维护}，因为它们回答的是链路上不同的问题：概念映射表确定“这段文本指向哪个概念”，本体关系确定“概念与概念之间是什么关系”，推导规则确定“哪些关系可以、以及如何转成任务标签”。混在一张表里，出错时就无法判断是映射错了、关系定义错了、还是推导规则用错了。映射候选同样可以由人工整理、从词语使用模式归纳或由模型建议，但无论来源如何，都要分别核验“文本表达的含义是否真对应该概念”和“所用的体系关系是否正确”。

使用层级时还要警惕\emph{属性继承}被滥用。沿上下位关系能继承\emph{类别成员资格}（首席架构师是管理岗位），但并非该概念的一切属性都能无条件向下或向上传递。某个具体岗位的薪酬区间、任期长度、资格要求，未必适用于上位类下的所有岗位；反过来，上位类的某些泛化属性，下位的特殊岗位也可能有例外。能不能继承某条属性，要回到该属性的语义本身去核对，不能因为“是子类”就默认一切属性都继承。

概念关联还有一个不可回避的多体系问题：同一个岗位词，在“按职责划分”“按组织级别划分”“按行业标准划分”的不同体系里，可能落在完全不同的分类位置。因此规则必须记录它当前采用的是\emph{哪一套体系、哪个版本}；当体系发生概念合并、拆分或层级调整时，才能沿“映射到推导”的链条，反查出哪些旧结果需要重算。

最后同样要守住“未关联并非负类”的边界：当前体系里找不到某个词的概念映射，只说明\emph{规则没能建立关联}（可能是词典没收这个新词），并不能证明这个样本就不属于目标类别。这与模式匹配中“未触发不等于负例”是同一个道理。至于使用了多套体系、需要把不同来源的概念彼此对齐时该怎么处理，则留到第~\ref{chap:knowledge-fusion}章讨论。

\subsection{基于关系匹配的规则利用}
\label{sec:rule-based-knowledge-acquisition-relation-matching}

实体匹配确定了“谁是谁”、概念关联确定了“属于哪类”，但若任务要标注的是“林舟\emph{任职于}海岬科技”，只查单个实体的属性或岗位类别仍然不够——这里要判断的是\emph{两个实体之间}的关系。\term{关系匹配}（relation matching）把样本中的一对实体与知识库中的关系事实对应起来，再用关联到的事实产生关系标签，或成组地提供监督信号。

从单实体转到关系，检查对象发生了实质变化：现在的基本单位是一个\emph{有次序的实体对}（谁在前、谁在后不能颠倒），而且必须同时确定关系的定义和它的有效时间。例如三元组$(\text{林舟},\text{任职于},\text{海岬科技})$是有方向的，不能换成$(\text{海岬科技},\text{任职于},\text{林舟})$（方向反了，关系就不成立）；它也有时间性，不能不加区分地套用到该事实生效之前或失效之后的所有文本上。

关系标注需要大量“这句话是否表达了某关系”的逐句标签，全靠人工标注极其昂贵。一个能批量、自动产生训练候选的想法应运而生：既然知识库已经记录了“林舟任职于海岬科技”，那么凡是\emph{同时提到}这两个实体的句子，很可能都在表达这项关系——于是直接用知识库事实给这些共现句子打上关系标签。这种\emph{借助已有知识库事实、自动为提及它的文本提供标签}的办法，称为\term{远程监督}（distant supervision）。Mintz等人的经典设定正是把知识库关系与文本中共同出现的实体对对齐，据此自动构造训练实例\scite{knowledge-acquisition-mintz2009}。

必须立刻点明这个假设的强度：远程监督依据的是一条启发式——“若两个实体存在某种关系，则同时提到它们的句子倾向于表达该关系”。它\emph{没有逐句证明}每一个共现句子真的在表达那项事实。共现只是关系表达的必要线索，远非充分条件，下一段的对照就会暴露由此产生的错误标签。理解这一点，是正确使用远程监督的前提：它能极大降低标注成本，但代价是标签里混入噪声。

用一个严格控制的教学对照看清噪声从何而来。继续使用构造事实“林舟任职于海岬科技”，并把目标关系限定为“文本所述时间已经生效的任职状态”；假设句甲中的任命已经生效、所述时点仍在该事实的有效期内。在这些明确前提下比较两句话：句甲“海岬科技任命林舟为首席架构师”确实支持目标任职关系；句乙“林舟在论坛上批评了海岬科技的发布流程”虽然也同时提到了林舟和海岬科技，却只表达了“批评”这一动作，与任职毫无关系。

如果机械地执行远程监督、把事实复制给所有共现句子，句乙就会被错误地贴上一个“任职关系”的正标签，而这个标签在文本中毫无依据——这正是远程监督噪声的典型来源。需要说明，这里“句甲支持、句乙不支持”是教学中我们\emph{已知}的语义对照，用来帮助读者理解问题，并不是机器仅凭共现就能作出的判断；真实系统面对成百万条句子时并不知道哪一句是“句乙”。图~\ref{fig:rule-based-knowledge-acquisition-relation-path}中，两句经过实体消歧和关系对齐后会被归入同一个“提及组”；图中表示语义对照的点线只是教学标注，不属于自动判定步骤，构造提及组时也不要求预先知道组内哪一句真正表达了关系。

\begin{figure*}[htbp]
  \centering
  \begin{tikzpicture}[
  x=1cm,
  y=1cm,
  font=\figurefont\footnotesize,
  source/.style={draw=FigureInput,line width=0.9pt,fill=FigureInput!9,minimum height=10mm,text width=32mm,align=left,inner sep=1.5mm},
  fact/.style={draw=FigureModel,line width=1pt,fill=FigureModel!10,minimum height=9mm,text width=53mm,align=center,inner sep=1.5mm},
  process/.style={draw=BookMuted,line width=0.8pt,fill=BookPaper,minimum height=10mm,text width=13mm,align=center,inner sep=1.5mm},
  comparison/.style={draw=BookMuted,line width=0.8pt,dotted,fill=white,minimum height=15mm,text width=36mm,align=center,inner sep=1.5mm},
  group/.style={draw=FigureOutput,line width=1.1pt,fill=FigureOutput!7,minimum height=37mm,text width=36mm,align=center,inner sep=1.5mm},
  flow/.style={knowledge arrow,draw=BookInk,line width=1.1pt,rounded corners=1mm},
  evidence/.style={knowledge arrow,draw=FigureModel,line width=0.9pt,dashed,rounded corners=1mm},
  contrast/.style={draw=BookMuted,line width=0.8pt,dotted},
  edge label/.style={text=BookMuted,inner sep=1pt,font=\figurefont\scriptsize}
]
  \fill[white] (0,0) rectangle (16.7,7);

  \node[source] (sent-a) at (1.8,5.35) {句甲：海岬科技任命林舟为首席架构师。};
  \node[source] (sent-b) at (1.8,2.25) {句乙：林舟在论坛上批评了海岬科技的发布流程。};

  \node[process] (entity-a) at (4.85,5.35) {实体消歧};
  \node[process] (entity-b) at (4.85,2.25) {实体消歧};
  \node[process] (mapping-a) at (7.1,5.35) {关系对齐};
  \node[process] (mapping-b) at (7.1,2.25) {关系对齐};

  \node[fact] (kb) at (7.35,3.8) {已有构造事实\\$(\text{林舟/E-017},\text{任职于},\text{海岬科技/E-003})$};
  \node[comparison] (positive) at (10.3,5.35) {教学语义对照\\句甲表达关系\\可作句子级候选};
  \node[comparison] (cooccur) at (10.3,2.25) {教学语义对照\\句乙仅共同提及\\不应给句子正标签};
  \node[group] (bag) at (14.55,3.8) {提及组级候选\\句甲、句乙组成一组\\由实体关系库事实关联关系\\假设至少一条表达关系};

  \draw[flow] (sent-a) -- (entity-a);
  \draw[flow] (entity-a) -- (mapping-a);
  \draw[contrast] (mapping-a.east) -- (positive.west);
  \draw[flow] (mapping-a.north) -- (7.1,6.55)
    -- node[edge label,above] {按实体对归组} (14.55,6.55) -- (bag.north);

  \draw[flow] (sent-b) -- (entity-b);
  \draw[flow] (entity-b) -- (mapping-b);
  \draw[contrast] (mapping-b.east) -- (cooccur.west);
  \draw[flow] (mapping-b.south) -- (7.1,1.05)
    -- (14.55,1.05) -- (bag.south);

  \draw[evidence] (kb.north -| mapping-a.south) -- (mapping-a.south);
  \draw[evidence] (kb.south -| mapping-b.north) -- (mapping-b.north);
  \draw[evidence] (kb.east) -- (bag.west);

  \draw[flow] (0.35,0.35) -- (1.25,0.35);
  \node[anchor=west,text=BookMuted,font=\figurefont\scriptsize] at (1.35,0.35) {处理步骤};
  \draw[evidence] (4.1,0.35) -- (5.0,0.35);
  \node[anchor=west,text=BookMuted,font=\figurefont\scriptsize] at (5.1,0.35) {事实关联};
  \draw[contrast] (7.85,0.35) -- (8.75,0.35);
  \node[anchor=west,text=BookMuted,font=\figurefont\scriptsize] at (8.85,0.35) {已知语义对照，非自动判定};
\end{tikzpicture}

  \caption{从已有事实到文本候选关系的关联路径。实体和关系均为构造例子，并假设句甲的任命已生效、所述时点在事实有效期内；实线箭头表示处理步骤，虚线箭头表示实体关系库事实的证据关联，点线引出教学中已知的语义对照，不表示自动判定。实体消歧确认两句中的提及指向同一实体对；关系对齐检查定义、方向与时间。在上述前提下，句甲支持任职关系，句乙仅共同提及，但两句都可按实体对进入同一提及组。组级候选依赖至少一次表达的假设，不要求预先知道哪句支持关系，也不能据此给每句确定的正标签。}
  \label{fig:rule-based-knowledge-acquisition-relation-path}
\end{figure*}

要缓解上述噪声，需要先把整条关联路径上的几个环节走扎实：确认实体身份与关系方向，核对知识库关系与任务关系的定义是否对齐，最后还要明确一个关键问题——\emph{证据应当出现在什么范围}：是要求就在当前这一句里表达，还是放宽到整段话、乃至同一实体对的多条提及中只要有一次表达即可？证据范围一旦改变，被监督的对象（单句 vs. 实体对）也随之改变。

远程监督的原始设定近似要求“每个共现提及都表达该关系”，过强而不符合实际。Riedel等人把它放松为\emph{至少一次表达}的假设：分别建模两个层次——“这个实体对\emph{总体上}是否存在该关系”和“\emph{某一条具体提及}是否表达该关系”，只假设一个存在关系的实体对，其多条提及中\emph{至少有一条}真正表达了关系，而不要求条条都表达\scite{knowledge-acquisition-riedel2010}。这样一来，像句乙那样的噪声提及就可以被容忍：它不必为正，只要同组里存在像句甲那样真正表达关系的提及即可。

在“至少一次”假设下，数据的组织方式随之改变：把同一实体对的多条提及打包成一个\emph{提及组}（bag），赋予组一个\emph{组级候选}标签，表示“组内至少有一条提及可能支持该关系”，而不再给组内每一句话都打正标签。模型学习时可以先对每一句“是否表达关系”分别打分，再按某种方式（如取最大分）聚合成组级预测，训练时只把组级预测与组级标签比较。第~\ref{chap:knowledge-acquisition-model-based}章的式~\eqref{eq:knowledge-acquisition-model-bag-prediction}给出了这种多实例预测的具体计算和所依赖的假设。

要厘清本章规则在其中的角色：规则负责产出提及组、实体标识和关联到的知识库事实，即把原料和组结构准备好；但“组内究竟哪一句构成证据”仍是一个需要\emph{提及级推断}的学习问题，规则并不回答。因此不能用教学里我们人为已知的“句甲支持、句乙不支持”去代替模型本应输出的预测——那是讲解用的答案，不是自动流程的产物。

把整条路径（实体消歧、关系对齐、证据范围判断）的每一步都记录下来，出错时才能区分性质不同的故障：消歧环节出错会把提及关联到无关实体；使用了过时事实会把关系套到错误的时间；而文本本身没表达关系，则是证据范围条件未被满足。

还有一种常见情况：知识库里\emph{根本没有}查到该实体对的关系。这时应采取什么态度至关重要。本章采用\term{开放世界假设}（open-world assumption）：承认知识库可能是不完整的，“没有查到这条关系”只说明库中未记录，\emph{不等于}“已证明该关系不存在”——也许只是还没收录。因此规则在查不到时应保留未知或候选关系，而不是草率地返回负标签；只有当存在明确的排除性证据、且它正好符合任务对负类的定义时，才能给出负例。这与全章反复强调的“未触发/未查到不等于负类”一脉相承。

开放世界假设带来大量“只有正例已知、其余未知”的实体对，自然引出一个问题：训练关系分类器时，是否\emph{必须}先把这些未知对象猜成负例？把未知一律当负，显然会把其中真正的正关系误标，扭曲训练。

其实不必如此。在一定抽样条件下，可以直接用“一批确认的正例”加“一批未标注样本（其中正负混杂）”来学习，这就是正例—未标注学习（PU learning）。其条件是：正例样本能够代表目标正类的分布、未标注样本能够代表目标总体分布、并且正类比例已知或可估计；满足这些条件时，可以从两批样本的平均损失\emph{重建}出完整的分类风险，而无须逐条把未知改写成负类。具体构造见第~\ref{sec:ka-pu-learning}节。需要强调这些条件并不免费：如果知识库偏向于只收录显著、知名的关系，那么抽中的正例往往是“显著关系”的代表、而非全部正关系的代表，代表性假设就不能默认成立；开放世界假设本身只说明“未查到并不等于不存在”，它并不自动提供PU学习所需的抽样保证。

即使上述各步都能正确执行，最终候选的\emph{分布}仍可能失真，需要单独检查两类量：每个关系类别的候选总量，以及每个实体对的提及组大小。读这些数字时要避免错误归因：两个关系类别的候选总数不同，可能只是因为它们涉及的实体对数量本来就不同，不能据此推断“某类中单个实体对被提及得更多”——总量与组内密度是两回事。

还要警惕重复带来的虚高：如果同一实体对在新闻里被反复共现，或相关文本被多家媒体转载，在没有去重的情况下，这些近乎重复的提及都会被计入该实体对的提及数，使组显得很大、证据很“充分”，但信息其实是重复的。这类分布差异仅靠“规则执行成功”无法发现，因此还要保存知识库版本、事实来源、实体对、组内提及列表以及去重标识，供后续的融合与可靠性检查追溯。

\section{基于约束推理的规则利用}
\label{sec:rule-based-knowledge-acquisition-constraint}

前两类规则都是在\emph{直接产生}答案：模式匹配说“这个样本像正类”，知识关联说“该实体有某属性”。但还有一类常见情形，单个答案看起来都合理，放在一起却自相矛盾。比如对同一岗位，细类判断说“他是首席架构师（为真）”，上位类判断却说“他不属于管理岗位（为假）”，而按岗位体系首席架构师必属管理岗——两个结论不可能同时成立。

处理这种矛盾，最自然的办法不是再分别重判一次，而是退一步规定：\emph{这些答案之间必须满足什么关系}。这就引出第三类依据。\term{约束推理}（constraint-based inference）先把每个未知标签所有可能的取值列为候选，然后用逻辑、结构或领域知识给出的条件，\emph{排除掉那些不允许的组合}；它通常不直接指定唯一答案，而是缩小可行答案的范围。排除之后若只剩一个候选，答案随之确定；若仍剩多个，可以把这一整个范围作为结果保留，也可以再引入一类“允许违反、但违反要付出代价”的软偏好，在多个可行候选中选取代价最小的。

要准确说出“约束到底排除了什么”，先要划定一个出发点：在不考虑任何约束时，任务\emph{理论上允许}哪些完整答案？把所有这些完整答案收集成集合$\bm{Y}$，其中每个元素$\bm{y}$是对一组未知标签或数值的一次\emph{完整赋值}（即给每个未知量都指定一个取值）；$\bm{Y}$由任务规定的取值范围确定，比如每个二值标签可取0或1，则$n$个标签共有$2^n$种赋值。

然后再施加必须无条件满足的条件。设有$m$条这样的硬约束$c_1,\ldots,c_m$，每条约束对完整赋值$\bm{y}$返回真或假。\emph{同时满足全部硬约束}的那些赋值构成可行集合：
\begin{equation}
  \bm{F}
  =
  \{\bm{y}\in\bm{Y}:c_k(\bm{y})=\mathsf{T},\ k=1,\ldots,m\}.
  \label{eq:rule-based-knowledge-acquisition-feasible-set}
\end{equation}
式中要求每一条硬约束同时为真，因此$\bm{F}$回答的问题是“哪些\emph{完整答案仍然被允许}”，请特别注意它\emph{不}回答“哪个答案更好、更受偏好”——可行与否是一个是非判断，偏好排序则是另一回事（将由软约束处理）。

在约束确实能对所有候选求值的前提下，可行集合$\bm{F}$会出现三种情况，必须区别对待，因为它们传递的信息完全不同。其一，$\bm{F}$只有唯一元素：当前前提恰好把答案完全确定，这个唯一赋值可作为结果。其二，$\bm{F}$有多个元素：依据只足以把范围缩小，但不足以唯一确定，此时应保留整个可行范围而不是任取一个冒充确定答案。其三，$\bm{F}$为空集：没有任何赋值能同时满足所有条件，这是一个有价值的信号——说明已知事实、约束规则或观测值之间\emph{存在冲突}（例如数据违反了物理守恒），需要排查，而不是随便给个答案。还要区分“无解”与“无法判断”：若只是缺少某些前提、导致约束暂时无法求值，应保留为无法确定，而不能草率宣告无解。

把这些映射回样本标注状态。当任务\emph{允许返回范围}时，一个非空的可行集合本身就可以作为内容提交：来源动作为作答、语义状态为已给出，但要意识到其中单个值可能仍未知。而当无法判断（缺前提）或已证明无解（$\bm{F}$为空）时，本节都规定\emph{不}提交任务赋值：来源动作为弃权、内容为空、语义为未知；但两者要分别留痕——缺前提时保存“缺的是哪个前提”，无解时另存无解诊断和相互冲突的依据，不能把“算不出来”和“算出来发现矛盾”混成一类。与前面一样，上述执行状态都以“相应的检查或求解确实跑完”为成功前提；程序中断记为执行失败。最后，求解顺利完成只在逻辑上得到了可行集，验证状态并不会因此自动变成“已核验”——满足约束不等于符合事实。还要记住触发与求解的区别：条件检查触发了求解，绝不代表求解已经给出任务内容。

\subsection{基于逻辑约束的规则利用}
\label{sec:rule-based-knowledge-acquisition-logical-constraint}

最基本的约束作用于若干真假命题之间的逻辑关系。要排除“细类成立、上位类却不成立”这种矛盾，可以把两项类别判断分别写成命题，再规定它们之间的\emph{蕴含}关系。一般地，\term{逻辑约束}（logical constraint）用蕴含（“若$p$则$q$”）、互斥（“$p$与$q$不能同时成立”）、等价（“$p$与$q$同真同假”）等关系，限制若干命题允许怎样\emph{共同}取值。

逻辑约束有两种强弱不同的用法。当作\emph{硬约束}时，它代表必须遵守的规则，任何违反该关系的候选组合都直接从可行集合中剔除，没有商量余地。但有时我们想表达的并不是“绝对禁止”，而是“最好如此、违反了会差一些”的\emph{偏好}，这时可改用\term{软约束}（soft constraint）：违反它的候选仍然允许保留，但要为此付出一个事先指定的代价。由于违反软约束不会使候选不可行，最终选谁就必须结合代价的权重或其他比较依据，而不能只看是否违反。

用最简单的两个命题把逻辑约束算一遍。设二值命题$p$表示“属于技术管理岗位”，$q$表示“属于管理岗位”，教学体系要求$p\Rightarrow q$（技术管理岗位必然是管理岗位）。两个二值命题共有$2^2=4$种赋值组合：$(0,0),(0,1),(1,0),(1,1)$。逐一检验，只有$(p,q)=(1,0)$违反蕴含（是技术管理岗却不是管理岗），被硬约束排除；其余三种仍在可行集合中。这一步只是在做逻辑筛选，并没有判断现实中$p,q$到底取什么值。

如果还想在三个可行候选中进一步挑一个，就需要引入偏好：人为给每个候选设置一个软约束代价$L(\bm{y})$，然后按代价从小到大比较。图~\ref{fig:rule-based-knowledge-acquisition-feasible-set}给了一组教学代价，其中三个可行项里最小者恰好唯一。但这里有一个必须向读者讲透的要点：\textbf{代价最小只代表它最符合“我们当前设定的偏好”，完全不证明它就是真实答案}。图的下方面板则演示了另一种情形：当观测之间本身矛盾（比如同时观测到$p=1$、$q=0$，又坚持$p\Rightarrow q$），没有任何赋值能满足全部条件，可行集合变为空——这提示数据或约束冲突，而不是某个候选“代价高”。

\begin{figure*}[htbp]
  \centering
  \begin{tikzpicture}[
  x=1cm,
  y=1cm,
  font=\figurefont\footnotesize,
  candidate/.style={draw=BookMuted,line width=0.8pt,fill=BookPaper,minimum width=30mm,minimum height=13mm,align=center,inner sep=1.5mm},
  selected/.style={candidate,draw=FigureModel,line width=1.2pt,fill=FigureModel!10},
  premise/.style={draw=FigureInput,line width=0.9pt,fill=FigureInput!9,minimum width=28mm,minimum height=10mm,align=center,inner sep=1.5mm},
  constraint/.style={draw=BookAmber,line width=1pt,fill=BookAmber!9,minimum width=31mm,minimum height=10mm,align=center,inner sep=1.5mm},
  combine/.style={circle,draw=BookInk,line width=0.9pt,fill=white,minimum size=7mm,inner sep=0pt},
  empty/.style={draw=BookAmber,line width=1.2pt,fill=BookAmber!9,minimum width=32mm,minimum height=11mm,align=center},
  flow/.style={knowledge arrow,draw=BookInk,line width=1.1pt,rounded corners=1mm},
  soft/.style={text=BookMuted,font=\figurefont\scriptsize}
]
  \fill[white] (0,-1.15) rectangle (16.6,6.5);
  \node[anchor=north west,font=\figurefont\bfseries] at (0,6.25) {(a) 硬约束筛选与软偏好};
  \node[constraint,minimum width=44mm] (hard) at (8.3,5.75) {硬约束：$p\Rightarrow q$};

  \node[candidate] (c00) at (1.8,4.25) {$(p,q)=(0,0)$\\{\scriptsize 可行；$L=1.0$}};
  \node[candidate] (c01) at (5.65,4.25) {$(p,q)=(0,1)$\\{\scriptsize 可行；$L=0.6$}};
  \node[candidate,draw=BookAmber,fill=BookAmber!7] (c10) at (9.5,4.25) {$(p,q)=(1,0)$\\{\scriptsize 违反硬约束}};
  \node[selected] (c11) at (13.35,4.25) {$(p,q)=(1,1)$\\{\scriptsize 可行；$L=0.2$（最小）}};
  \draw[BookAmber,line width=1.2pt] ([xshift=-10mm,yshift=-4mm]c10.center) -- ([xshift=10mm,yshift=4mm]c10.center);
  \draw[BookAmber,line width=1.2pt] ([xshift=-10mm,yshift=4mm]c10.center) -- ([xshift=10mm,yshift=-4mm]c10.center);
  \node[soft,anchor=north] at (8.3,3.42) {硬约束留下三个答案；软代价只负责排序};

  \draw[BookRule,line width=0.8pt] (0,2.95) -- (16.6,2.95);
  \node[anchor=north west,font=\figurefont\bfseries] at (0,2.75) {(b) 新增互相矛盾的已知条件};
  \node[premise] (p-one) at (2.2,1.65) {观测：$p=1$};
  \node[premise] (q-zero) at (7.2,1.65) {观测：$q=0$};
  \node[constraint] (imply) at (12.2,1.65) {硬约束：$p\Rightarrow q$};
  \node[combine] (and) at (7.2,-0.35) {$\land$};
  \node[empty] (none) at (12.4,-0.35) {无可行解\\报告冲突};
  \draw[flow] (p-one.south) |- (and.west);
  \draw[flow] (q-zero.south) -- (and.north);
  \draw[flow] (imply.south) -- ++(0,-0.5) -| (and.45);
  \draw[flow] (and.east) -- (none.west);
\end{tikzpicture}

  \caption{约束缩小可行答案并产生冲突的教学示意。上面板从两个二值命题的四种赋值出发，硬约束$p\Rightarrow q$排除$(1,0)$，软约束代价再在三个可行候选中给出唯一最小值；这些代价是人为设置，不是统计估计。下面板同时加入观测$p=1$、观测$q=0$和同一蕴含后，可行集合变为空集。逻辑可行、偏好最小和真实正确是三种不同判断。}
  \label{fig:rule-based-knowledge-acquisition-feasible-set}
\end{figure*}

实际偏好往往不止一条，需要一个可计算的量来汇总多条软约束的共同作用。令第$j$条软约束对候选$\bm{y}$的违反代价为$\ell_j(\bm{y})\geq0$：取零表示完全不违反，取正值表示有违反，值越大违反越严重；再给每条软约束配一个权重$w_j\geq0$，规定这一项偏好在总比较中占多大分量。

一个关键安排是：\textbf{软约束只在硬约束筛过的可行集合$\bm{F}$内部起作用}。也就是说，先无条件满足全部硬约束，再对选定的有限条软约束求加权代价之和，总数越小表示越受偏好。绝不允许为了降低软代价而选择违反硬约束的候选——硬约束是入场资格，软约束只是入场后的排序依据。当可行集非空、且最小值确实能取到时，选择总加权代价最小的候选，得到
\begin{equation}
  \widehat{\bm{y}}
  \in
  \argmin_{\bm{y}\in\bm{F}}
  \sum_j w_j\ell_j(\bm{y}).
  \label{eq:rule-based-knowledge-acquisition-soft-constraint}
\end{equation}
式中$\bm{y}\in\bm{F}$这一限制保证了被比较的候选都没有违反硬约束；目标值更小，才表示它在“当前给定的权重”下更受偏好——换一套权重，最优候选可能改变。$\argmin$会保留\emph{所有}达到最小值的候选而不强行挑一个，因此若最小值由多个候选并列取得，答案实际上仍未唯一确定，此时应如实报告多解。存在性也要分清：在有限的非空候选集合上、代价又有限，最小值一定存在；但若是连续取值的优化问题，最小值未必能取到（可能只是下确界），需要另外确认前述存在条件。

这种“硬约束定边界、软代价做排序”的框架，可以通过整数规划或线性规划来求解，从而把多个局部模型的预测与任务层面的逻辑条件融合起来；Roth和Yih在自然语言任务中给出了有代表性的形式化\scite{knowledge-acquisition-roth2004}。还可以再往前走一步：若希望把这套加权逻辑\emph{解释成一个概率模型}（让权重对应“某条约束成立的可能性”），就必须明确权重的统计含义，并确认公式代入具体对象后有足够的实例支撑；Markov逻辑网络为此提供了一种把加权逻辑规则概率化的实现方式\scite{pgm-energy-models-richardson2006}。但要注意，随手设定的权重并不天然具备概率解释。

把知识写成逻辑约束时，最容易犯的错误是丢掉关系的\emph{方向}和\emph{适用对象}，这与概念关联中上下位不可逆是同一个警示。由“首席架构师属于技术负责人”只能得到有方向的蕴含$\text{首席架构师}\Rightarrow\text{技术负责人}$，绝不能反过来说所有技术负责人都是首席架构师。此外，每条约束往往只在特定对象范围内成立：针对“正式岗位”设定的互斥条件（一人不能同时任一岗）未必适用于允许兼任的临时角色；把约束的适用范围悄悄扩大，会错误地砍掉合法候选。

因此，无论约束来自专家编写、从样本归纳还是模型建议，都要逐项核对四件事：蕴含方向有没有写反、量词（对“所有”还是“某些”对象成立）是否正确、作用对象是否恰当、以及有没有被忽略的例外情况。

为了让约束结果可维护、可在规则更新后复查，一次推理应留下完整的“决策说明”：用了哪些前提、据此排除了哪些候选、每条软约束分别贡献了多少代价。这样当某条约束或权重调整后，才能沿说明重新计算并看出结果为何改变。成本也要如实计入：约束求解本身的计算量、软权重如何设定、出现冲突时如何解释，都不是免费的。

最后划清本章与下一章的边界：这里讨论的约束，作用是\emph{从既定前提出发去产生或限制新的候选}；如果约束主要被用来\emph{综合多个已经存在的来源}给出的答案（比如调和多个模型、多种标注的分歧），那属于第~\ref{chap:knowledge-fusion}章的融合方法，不应在规则推理层提前完成。

\subsection{基于结构约束的规则利用}
\label{sec:rule-based-knowledge-acquisition-structural-constraint}

逻辑约束处理的是抽象命题之间的关系，但很多任务的标签带有\emph{空间或序列结构}：一个实体名可能横跨好几个词，每个词各自选了标签之后，还必须保证这些标签能拼出一个合法的片段。

\term{结构约束}（structural constraint）利用样本在序列、树、图或成组形式中的组织关系，来限制标签\emph{之间}如何组合。它与一般逻辑约束的区别在于：除了标签的取值关系，还必须先明确“哪些位置（或对象）彼此相邻、谁包含谁、哪些属于同一组”——也就是约束作用的那张“结构”从哪来。这张结构并非凭空假设，它可以直接来自任务格式（如固定的字段顺序）、标注规范（如片段如何界定）或样本字段，也可以由实体匹配等前置步骤构造出来；来源不同，可靠性也不同。

要让一串逐词标签表达“片段从哪开始、到哪结束”，需要一套边界编码约定。本节采用BIO体系中的IOB2约定，三个标记分别是beginning（开始）、inside（内部）和outside（外部）的缩写：对任意实体类型X，一个X类型片段的第一个词标为\texttt{B-X}，片段内部后续的词标为\texttt{I-X}，不属于任何片段的词标为\texttt{O}。

这套约定立刻导出一条结构约束：\texttt{I-X}只能出现在\texttt{B-X}或另一个\texttt{I-X}之后（即延续一个已经开始的同类型片段），它\emph{不能凭空出现}在句首或一个\texttt{O}之后——否则就成了一个“没有开头的片段”，是非法组合。这种把片段结构编码进逐词标签、再检查局部转移合法性的做法，源自早期的文本分块工作；需要提醒的是，历史上还有一种IOB1约定，它只在两个同类片段紧临时才用\texttt{B}标记后一个片段的开头，其余片段内部用\texttt{I}开头，与本节“每个片段都以\texttt{B}开始”的IOB2约定不同，阅读文献时不要混淆\scite{knowledge-acquisition-ramshaw1995}。

用一个四词教学序列把约束的作用和局限都看清楚。对“林舟~加入~海岬~科技”，先由局部规则给每个位置提供候选标签，其中\texttt{PER}、\texttt{ORG}分别表示人员和组织。假设逐点独立选择后得到一串标签
\[
  \texttt{B-PER},\ \texttt{I-ORG},\ \texttt{B-ORG},\ \texttt{I-ORG}.
\]
检查相邻转移就会发现非法之处：第二个位置的\texttt{I-ORG}想延续一个组织片段，但它前面的位置是\texttt{B-PER}（人员片段的开头），并不存在一个已开始的组织片段可供延续，违反了“I必须承接同类型B/I”的规则，因此整串组合非法，必须修改。图~\ref{fig:rule-based-knowledge-acquisition-sequence-constraint}给出的另一串\texttt{B-PER, O, B-ORG, I-ORG}则满足所有转移条件，是一个合法序列。

但这个例子同时暴露了结构约束的\emph{局限}：合法不等于正确。把“海岬科技”两个字全都标成\texttt{O}（序列变成\texttt{B-PER, O, O, O}）其实也完全满足转移规则、照样合法，可它显然漏掉了真正的组织提及。结构约束只能排除“拼不成合法片段”的组合，却无法判断众多合法组合中哪一个在语义上正确；要选出真正的任务答案，还必须借助词语、上下文或其他来源提供的\emph{语义}证据。约束负责“剔除不可能”，语义负责“在可能中选对”。

\begin{figure*}[htbp]
  \centering
  \begin{tikzpicture}[
  x=1cm,
  y=1cm,
  font=\figurefont\footnotesize,
  token/.style={draw=FigureInput,line width=0.9pt,fill=FigureInput!9,minimum width=30mm,minimum height=8mm,align=center},
  tag/.style={draw=FigureModel,line width=0.9pt,fill=FigureModel!9,minimum width=30mm,minimum height=8mm,align=center},
  bad tag/.style={tag,draw=BookAmber,fill=BookAmber!9},
  choices/.style={draw=BookMuted,line width=0.7pt,dashed,fill=white,minimum width=34mm,minimum height=9mm,align=center,font=\figurefont\scriptsize},
  allowed/.style={knowledge arrow,draw=BookInk,line width=1pt},
  invalid/.style={knowledge arrow,draw=BookAmber,line width=1pt,dashed},
  edge label/.style={fill=white,text=BookMuted,inner sep=1pt,font=\figurefont\scriptsize}
]
  \fill[white] (0,0) rectangle (16.85,6.3);
  \node[anchor=east,font=\figurefont\bfseries] at (1.25,5.85) {词元};
  \foreach \x/\word in {3.1/林舟,7.3/加入,11.5/海岬,15.1/科技} {
    \node[token] at (\x,5.85) {\word};
  }

  \node[anchor=east,font=\figurefont\bfseries,text=BookAmber] at (1.25,4.45) {局部};
  \node[tag] (bad1) at (3.1,4.45) {\texttt{B-PER}};
  \node[bad tag] (bad2) at (7.3,4.45) {\texttt{I-ORG}};
  \node[tag] (bad3) at (11.5,4.45) {\texttt{B-ORG}};
  \node[tag] (bad4) at (15.1,4.45) {\texttt{I-ORG}};
  \draw[invalid] (bad1) -- (bad2);
  \draw[allowed] (bad2) -- (bad3);
  \draw[allowed] (bad3) -- (bad4);
  \node[circle,draw=BookAmber,fill=white,text=BookAmber,line width=0.8pt,minimum size=4mm,inner sep=0pt,font=\figurefont\bfseries\scriptsize] at (5.2,4.45) {$\times$};

  \node[anchor=east,font=\figurefont\bfseries,text=FigureModel] at (1.25,3.0) {合法};
  \node[tag] (good1) at (3.1,3.0) {\texttt{B-PER}};
  \node[tag] (good2) at (7.3,3.0) {\texttt{O}};
  \node[tag] (good3) at (11.5,3.0) {\texttt{B-ORG}};
  \node[tag] (good4) at (15.1,3.0) {\texttt{I-ORG}};
  \draw[allowed] (good1) -- node[edge label,above] {允许} (good2);
  \draw[allowed] (good2) -- node[edge label,above] {允许} (good3);
  \draw[allowed] (good3) -- node[edge label,above] {允许} (good4);

  \node[anchor=east,font=\figurefont\bfseries,text=BookMuted] at (1.25,1.45) {候选};
  \node[choices] at (3.1,1.45) {\{\texttt{B-PER}, \texttt{O}\}};
  \node[choices] at (7.3,1.45) {\{\texttt{O}, \texttt{I-ORG}\}};
  \node[choices] at (11.5,1.45) {\{\texttt{B-ORG}, \texttt{O}\}};
  \node[choices] at (15.1,1.45) {\{\texttt{I-ORG}, \texttt{O}\}};
  \node[anchor=south,text=BookMuted,font=\figurefont\scriptsize,align=center] at (8.4,0.2) {结构约束排除非法组合，但合法序列仍可能不止一个};
\end{tikzpicture}

  \caption{局部标签与整体结构的约束关系。采用本节的IOB2边界约定：\texttt{B-X}开始类型X片段，\texttt{I-X}只能延续同类型的\texttt{B-X}或\texttt{I-X}，\texttt{O}位于片段外。上行的\texttt{B-PER}\,$\to$\,\texttt{I-ORG}违反转移条件；下行是一种合法候选。每个位置下方保留局部候选集合，联合约束筛选后仍可能留下多个合法序列。}
  \label{fig:rule-based-knowledge-acquisition-sequence-constraint}
\end{figure*}

结构约束不仅能事后验收，还能在求解过程中主动缩小候选，这就是\emph{局部传播}。其做法是：反复扫描相邻位置，凡是某个位置的候选标签在相邻位置中\emph{找不到任何合法搭配}，就把它删掉；不断重复，直到再也删不动为止。

在上面的序列里，第二位置的\texttt{I-ORG}在第一位置的候选集$\{\texttt{B-PER},\texttt{O}\}$中找不到任何一个能充当其合法前驱的标签（两者后面都不能直接跟\texttt{I-ORG}），因此它在第一轮就被删除。传播会持续进行：第四位置的\texttt{I-ORG}暂时不能删，因为第三位置还保留着\texttt{B-ORG}可以作它的前驱；但要注意这种支持是\emph{有条件}的——一旦第三位置最终选了\texttt{O}，第四位置的\texttt{I-ORG}就失去依托、不能再选。

这说明一个重要事实：传播结束后每个位置\emph{各自}仍留有候选，并不代表把这些候选\emph{任意}拼起来都是合法序列，跨位置的搭配约束依然存在。局部传播只是剪掉了“局部就注定不可能”的标签；要得到一个完整、全局合法的赋值，或要在所有合法序列中按语义评分挑出最优的一个，还必须进行\emph{联合选择}，根据结构形态选用动态规划（如序列）、搜索或通用约束求解器。

局部传播在链状序列上很有效，但把结构从序列换成一般的图时，“每处局部都可行”并不能保证“整体存在解”，一个经典例子是三节点的着色式矛盾。设三个变量$a,b,c$都只能取$0$或$1$，并要求\emph{每一对}变量取值都不同（即$a\ne b$、$b\ne c$、$c\ne a$）。

逐个局部检查会发现一切正常：对任意一对变量、任意一个取值，都能在另一个变量上找到相反值来满足“二者不同”，因此局部传播删不掉任何候选。但全局地看，两种赋值都会走进死胡同：若令$a=0$，为与$a$不同只能有$b=1$、$c=1$，可这样$b$与$c$相等，违反$b\ne c$；令$a=1$同理导致$b=c=0$，仍矛盾。结论是：三条两两约束\emph{局部都有支持}，合在一起却\emph{没有任何共同可行的赋值}（这正是奇环不能二染色的现象）。这提醒我们，局部一致性只是全局有解的必要条件而非充分条件，某些冲突只有在联合求解时才暴露。

结构约束的思想还可以扩展到更多关系：片段之间的\emph{边界包含}（大片段能否套小片段）、标签的\emph{父子层级}（子标签出现时父标签是否必须出现）、以及\emph{跨样本的身份一致性}（同一个实体在不同样本中应保持同一标识）。使用时都要把具体要检查的结构关系讲清楚，不能只说“加结构约束”。特别要注意身份一致与属性可变的区分：跨样本要求一致的是实体的\emph{标识}（这个人始终是同一个人），并\emph{不}要求他的岗位、位置或状态也一直不变——人会换工作、会移动，把身份一致错误地扩大成属性不变，会产生大量伪冲突。

所有这些结构推理都建立在一个前提上：那张“谁挨着谁、谁包含谁”的结构本身是有依据的。而结构可能来自可靠程度很不相同的渠道：文档层级可能直接来自文件格式（较可靠），跨样本的同一对象关系要靠实体匹配确定（可能出错），图中的邻接边则可能由距离阈值或某条关系规则构造出来（带有主观性）。

结构一旦有错，约束推理会“忠实地”沿着错误结构得出错误结论，且形式上毫无破绽：缺了一条本应存在的边，会阻断本该发生的传播，使两个本应联动的对象失去约束；多出一条错误的边，会把毫不相干的对象强行联结，逼它们满足本不该有的关系；边界定错，则可能把本来合法的标签组合误判成冲突。因此结构的来源与版本必须像标签证据一样保存；而且要清醒，约束求解\emph{成功}（得到了一个干净的合法解）绝不能反过来证明所用结构正确——结构错了，求解照样能成功。

结构约束的求解成本高度依赖结构的规模和连接方式，选择算法时要据此判断。对只含相邻转移条件的长序列，存在高效的动态规划解法（如Viterbi式递推），可以在与序列长度成线性或多项式的时间内求出全局最优合法序列；但当结构是连接稠密的图、约束又是全局性的（如三节点那种相互耦合），问题可能要求指数级的搜索空间，未必存在高效精确算法。

若因此改用近似方法，求解说明中必须如实写明两件事：还有哪些冲突没有真正解决，以及返回的结果到底是“满足全部约束的完整解”，还是“只满足了部分局部条件的近似结果”。把近似解包装成完整解，会让下游误以为所有结构约束都已满足。

\subsection{基于领域约束的规则利用}
\label{sec:rule-based-knowledge-acquisition-domain-constraint}

逻辑约束和结构约束具有较强的通用性，与具体行业无关；最后一类约束则直接来自某个\emph{具体领域的客观规律或业务规则}。一个直观例子是盘点物料：如果知道期初库存量、这段时间的流入量和期末库存量，那么那段没有直接测量的流出量其实被“物料守恒”限制住了——东西不会凭空产生或消失。

\term{领域约束}（domain constraint）把这种守恒关系、数量平衡、取值范围限制或业务一致性规则，写成未知标签或未知数值必须满足的条件，再用已有的观测去缩小未知量的候选范围。引入它并不意味着要换一个全新的求解器（很多情况下仍是在解同一类约束满足问题），真正改变的是\emph{约束的依据}：它来自特定领域的知识，因此必须额外讲清每个变量代表什么物理量、在什么环境与前提下成立、用什么单位、允许多大容差。脱离这些限定，守恒方程可能根本不适用。

先在理想条件下算一个能得到唯一值的例子。考虑同一个核算时段内的一组教学观测：初始库存$s_0=100\,\mathrm{kg}$，期间流入$q_{\mathrm{in}}=30\,\mathrm{kg}$，期末盘点得到结束量$s_1=108\,\mathrm{kg}$，未知的是期间流出量$q_{\mathrm{out}}$。

在一个\emph{封闭}前提下——除了已记录的流入、流出外，没有任何其他来源的进料，也没有损耗、失窃等未计量损失——物料在账面上应当严格平衡：期末量等于期初量加流入、再减流出。于是有精确的平衡关系
\begin{equation}
  s_1=s_0+q_{\mathrm{in}}-q_{\mathrm{out}},
  \label{eq:rule-based-knowledge-acquisition-balance}
\end{equation}
把三个已知量代入，$108=100+30-q_{\mathrm{out}}$，解得$q_{\mathrm{out}}=22\,\mathrm{kg}$。这里能得到唯一、确定的值，靠的是\emph{两个条件同时成立}：三个观测完整准确，以及“没有其他进出与损失”的封闭前提。这个前提一旦不成立（比如有未登记的损耗），就不能再沿用同一条等式作同样的推断，否则算出来的$22$只是“账面差额”，不等于真实流出。

现实中测量必有误差，要求等式严格成立往往不现实，此时问题应从“求唯一值”改成“判断哪些取值仍然与观测相容”。引入\emph{整体平衡残差}来量化偏离程度：它等于“观测到的期末量”与“按平衡关系算出的期末量”之差；对残差取绝对值，是为了同时限制“盘多了”和“盘少了”\emph{两个方向}的偏离，而不只是单侧。

假设各项观测的误差综合起来，使整体残差的绝对值不超过$2\,\mathrm{kg}$，那么不再要求等式精确成立，而只要求
\[
  \left|s_1-(s_0+q_{\mathrm{in}}-q_{\mathrm{out}})\right|
  \leq 2\,\mathrm{kg},
\]
仍代入同一组数值，即$|108-(100+30-q_{\mathrm{out}})|=|q_{\mathrm{out}}-22|\leq2$，规则只能推出一个可行区间$q_{\mathrm{out}}\in[20,24]\,\mathrm{kg}$。请注意：这时$22$虽然位于区间中央，但它\emph{不再是被确定的答案}，真实流出量可能是区间内任意值，不能把区间中点误报成精确结论。还要厘清$2\,\mathrm{kg}$这个数的口径：它约束的是\emph{合成后的整体残差}，并不是对三个观测各自误差界的分别声明，二者不能混用。

边界情形也要能识别：如果连流入量都缺失，已知量不足，原观测就不再能确定上述区间，应保留无法确定；若把期末量改成$s_1=140\,\mathrm{kg}$、同时仍要求流出非负$q_{\mathrm{out}}\geq0$，那么在“无其他来源”的前提下，账怎么算都平不了（期初加流入只有130，不可能有140），这是\emph{无解}，提示观测或前提冲突。因此求解器应能分别返回三种结果：唯一值、候选范围、观测冲突，并保存单位、容差与求解精度。

一条领域关系能不能当作“必须精确成立的硬等式”，取决于它的\emph{依据有多强}，可以区分三个层次。第一，若关系是按任务定义\emph{必然成立}的业务恒等式（例如复式记账里借贷必须相等、各互斥账项之和必须等于总量），它可以直接作为硬约束，违反即说明记账出错。第二，若是一条\emph{领域规律}（如物料守恒），但它依赖系统封闭、过程稳态或其他理想化条件，就必须把这些条件明确写进适用前提；前提不满足时，规律不能硬套。第三，若关系只是从有限样本中\emph{观察、拟合}出来的经验相关性，它通常更适合写成带误差的软约束，并保留估计误差和“条件变化需重新估计”的更新条款，而不能当成恒等式。

表~\ref{tab:rule-based-knowledge-acquisition-domain-constraint}固定同一个物料平衡例子，并排比较这三种依据各自允许推出多强的结论：精确等式给唯一值，带容差的理想规律给可行区间，经验关系只给偏好或异常候选。三者的数学表达式可能看起来很像，但证据强度和可用结论完全不同，绝不能因为“形式相似”就互换使用。

\begin{table*}[htbp]
  \centering
  \begingroup
  \setlength{\tabcolsep}{2.5pt}
  \setlength{\parindent}{0pt}
  \setlength{\RaggedRightParindent}{0pt}
  \renewcommand{\arraystretch}{1.12}
  \begin{tabularx}{\linewidth}{>{\raggedright\arraybackslash}p{23mm}>{\raggedright\arraybackslash}p{32mm}>{\raggedright\arraybackslash}p{25mm}>{\raggedright\arraybackslash}p{26mm}>{\raggedright\arraybackslash}p{32mm}>{\raggedright\arraybackslash}X}
    \toprule
    依据类型 & 前提 & 观测 & 误差或容差 & 可得结论 & 失效处理 \\
    \midrule
    按定义成立的关系 & 各量是完整且互斥的记账项 & $s_0=100$、$q_{\mathrm{in}}=30$、$s_1=108$，单位为$\mathrm{kg}$ & 无额外容差 & 精确值$q_{\mathrm{out}}=22\,\mathrm{kg}$ & 报告定义冲突或无解 \\
    理想条件下的领域规律 & 除已列流入、流出外，无其他物质通量或损失；观测误差已合成 & 同上；另核验核算边界及传感器版本 & 整体平衡残差绝对值不超过$2\,\mathrm{kg}$ & 可行范围$q_{\mathrm{out}}\in[20,24]\,\mathrm{kg}$ & 范围外停用；无法确定时弃权 \\
    样本归纳的经验关系 & 关系仅在指定设备、时间与样本范围内观察到 & 同类测量及规则拟合样本的版本 & 软代价与经核验的残差范围 & 偏好排序或异常候选 & 条件变化后重新估计 \\
    \bottomrule
  \end{tabularx}
  \endgroup
  \caption{领域约束的前提决定结论强度。表中数值均为单位一致的假设教学观测值；精确等式、带容差范围和经验软关系分别支持唯一值、可行区间和偏好候选，不能相互替代。}
  \label{tab:rule-based-knowledge-acquisition-domain-constraint}
\end{table*}

拿到一条候选的领域约束后，落到数据上还要先核对“公式里的变量”与“样本里的观测字段”是否真的表示同一个量——库存模型里的“流入”，是否对应数据表中那个字段的口径？字段名相近不代表物理含义相同。对由模型建议出来的关系尤其要警惕：模型很容易把“两个量高度相关”误写成“两者恒等”，而相关关系既无方向保证、也允许偏差，根本不具备守恒式的强度。

领域约束同样有时效性和环境依赖：当生产工艺、计量设备或业务流程改变后，原来“无额外物料通量、无损失”的封闭前提、误差范围甚至业务规则都可能不再成立，需要重新检查前提是否仍有效。因此开发和维护一套领域约束，工作内容远不止“把公式写出来”：还包括观测对齐、单位换算、误差建模、无解情形的解释以及版本管理，这些都是保证约束长期可用的必要成本。

最后把领域约束的几种结果按章首的状态约定逐一落成样本标注，并再次强调：平衡求解\emph{完成}这件事，本身并没有独立验证测量观测是否可靠——数学上解得出值，不代表传感器读数就是对的。

在“任务允许返回可行范围”的设定下：精确例提交唯一值$22\,\mathrm{kg}$，带容差例提交区间$[20,24]\,\mathrm{kg}$；两者的来源动作都是作答、执行成功、语义状态为已给出、验证状态为候选（尚未独立核验），区别在于后者给出的是范围、其中单个值仍未确定。再看两种“答不出”的不同情形：如果无法确认是否还存在其他物料通量（封闭前提是否成立都不知道），则检查虽完成但应弃权，内容为空、语义为未知，并记录“缺少核算边界证据”这一具体原因；如果进一步\emph{确认}存在未计入的通量，则是规则的适用前提不成立，该规则对此样本不适用，内容与动作都为空。最后是无解情形：当完整观测为$s_1=140\,\mathrm{kg}$、且非负流出要求不变时，求解成功地\emph{证明了无解}，但仍不提交流出量——动作是弃权、内容为空、语义为未知，同时把无解诊断和相互冲突的依据另行存档（这些诊断绝不能反过来被当成一个流出量标签）。

请特别区分最后这一对：“因前提不足而无法判断”与“已严格证明条件冲突而无解”，二者都导致没有任务赋值，但留下的原因和求解结果截然不同，必须分别记录；而“任务允许的多解范围”本身则是一项\emph{可以正式提交}的内容。规则产出的标签或约束质量究竟如何、错误怎样修正，将由第~\ref{chap:knowledge-quality-assessment-correction}章继续检查和处理。

\section{本章小结}
\label{sec:rule-based-knowledge-acquisition-summary}

回到章首的任命新闻可以看到，基于规则的知识获取绝不能止步于“程序在句子里找到了‘任命’这个词”。从一次字面命中到一个可用的候选标签，中间要走一条可追溯的链：先确认样本是否满足条件、再把提及关联到正确的对象、最后用答案之间的约束排除矛盾组合，而且每走一步都要说清依据是什么。

贯穿全章有两条彼此独立的主线，读者应当把它们分清。一条关于\emph{规则从哪来}：专家编写、从样本归纳、由模型生成，决定的是候选规则\emph{怎样被得到}。另一条关于\emph{规则靠什么下判断}：模式匹配、外部事实关联、约束推理，决定的是候选标签或约束\emph{凭什么产生}。这两条线可以自由交叉，一条规则也完全可以同时调用多种机制；但无论怎样组合，每一步都要分别保存待检查样本、规则结果、支持证据和版本信息，否则出错时无法沿链定位。

三类依据各有其能力边界，选用时应对各自的失效方式保持清醒。模式匹配不依赖外部资源、可以直接返回标签或关系候选，但它受符号变化、数值基准和条件组合的牵制，需要持续维护分词、阈值与分支逻辑。外部事实关联能够复用知识库里已有的实体、概念和关系，省去重复判断，但必须逐一检查实体消歧是否正确、事实版本与时间是否适用、以及证据范围是否真的覆盖当前句子——知识库为真不等于本句表达了它。约束推理不直接产生答案，而是排除不合法的组合、缩小候选范围或在无解时暴露冲突；它的一切结论都建立在前提和结构本身正确之上，结构错了，求解再干净也无济于事。尤其要记住，按软约束选出的“最小代价”候选，只代表它最符合\emph{当前人为设定的偏好}，这个选择动作本身并没有把它变成被证实的真实答案。

把任命例子中的各步串成一条完整链路，可以复习全章的方法：符号模式先在句中找到任命谓词及其参与者；实体链接接着把“林舟”“海岬科技”这两个提及对应到知识库中的唯一标识；关系匹配再按关系定义、方向、有效时间和证据范围，检查这个任职候选是否成立；当多个标签需要联合决定时，还可以用前面讲过的逻辑约束或结构约束筛掉矛盾组合。

这条链上所有的条件检查都遵循同一套三值约定，值得再集中复述一次：适用前提不成立时记为范围外、语义为不适用；在适用的前提下，前提与触发条件的合取为真才触发，为假则未触发且来源动作为空，无法确定时才弃权。合取遵循“假值优先于未知”——一个已确认的假值不应被另一个未知条件遮蔽。但强Kleene合取本身并不规定规则列表或规则树遇到未知时是否继续，那属于需要另行制定的\emph{控制策略}；触发之后若还有求解，求解结果也要单独记录。到了产出阶段：任务允许返回范围时，多个可行解可以作为候选内容正式作答；可行集为空说明无解，应如实报告冲突，而不是伪造一个标签。最后再强调全章反复出现的那条底线：\textbf{规则没有结果（范围外、未触发、弃权、未查到）都不等于负例}，一个真正的负标签仍然需要规则去明确检查负类的定义。

规则之所以值得使用，很大程度上是因为它的判断过程是\emph{显式}的：从待检查样本，到适用前提、模式匹配、事实关联，再到约束求解，每一步都留有记录，使错误可以沿链条被追查和定位——这是它相对“黑箱”判断的独特优势。但显式和可追溯\emph{并不}保证所有样本都能得到正确答案，规则会过时、会有覆盖不到的写法、也会依赖出错的前处理和外部事实。

因此是否继续使用一条规则，应当在三类因素之间权衡：它的依据是否仍可核验、它的适用范围是否还覆盖当前数据、以及维护它所需的前处理与更新成本是否划算。当缺少可用的规则条件或外部事实时，工作方向是补充覆盖；当多个来源（多条规则、规则与模型、规则与人）给出冲突结果时，不应在规则层强行裁决，而应进入第~\ref{chap:knowledge-fusion}章做多来源融合；而这些规则产出的标签和约束究竟可不可靠、发现错误后如何修正，则由第~\ref{chap:knowledge-quality-assessment-correction}章用统一的质量评估与修正流程收口。

\chapter{基于模型的知识获取}
\label{chap:knowledge-acquisition-model-based}

前面两章分别用了两类获取依据：第~\ref{chap:knowledge-acquisition-human}章靠人逐篇判断，准确但人力有限；第~\ref{chap:rule-based-knowledge-acquisition}章靠显式规则，可重复但只能覆盖人能事先写清的条件。现在考虑第三种常见情形：一批花卉照片已经采集完成，却没有逐株对应的物种答案；照片数量庞大、每张可见的部位各不相同，而“山茶与茶梅如何区分”这类鉴别依据又极其繁多，很难穷尽成一条条规则。人工看不过来、规则写不全，这时还能借助什么？

一个自然的答案是\emph{模型}：让一个从大量过去样本中学到规律的模型来参与作答。模型的判断依据既可能来自训练时沉淀在参数里的关系（见过成千上万张花后，它隐式学会了花瓣形态与物种的对应），也可能来自调用现场才提供的材料——比如随请求附上带答案的参考照片，或样本之间预先建好的连接。它据此产出的答案形态也很多样：一个类别、一个数值、一对对象之间的关系，甚至一段说明文本。

但要立刻重申第~\ref{chap:knowledge-acquisition-overview}章定下的立场：模型吐出的内容并不会因为“算得快、看起来自信”就成为可靠依据。这些预测输出连同其来源、适用范围和核验状态一起保存后，仍称为\term{候选监督信号}——它们已经被取得，但是否能作为当前任务的判断依据、能否进一步当作训练样本，都还需要检查。本章研究的正是怎样\emph{规范地取得和管理}这些由模型产生的候选，而不是默认它们正确。

一次模型调用成功，只回答了最表层的问题——“模型返回了什么字符串”。要敢用这个答案，还必须回答两个更深的问题：这条输出\emph{依据什么}算出、又\emph{适用于哪里}。为此至少要沿三条依赖逐一确认：先确认模型本身能接收当前这类照片、并按任务指定的物种表给出输出（能力与接口对得上）；再确认这次调用\emph{实际}给它看了哪些观测、示范和外部材料（不能假设它知道没提供的信息）；最后辨认答案主要是通过哪条机制算出的——是模型参数里的映射、是现场检索到的参考样本，还是图结构上的传播。

把这些依赖记录清楚，最大的价值在出错时体现：复核发现一个物种标错了，只有沿依赖链往回查，才能判断问题出在模型没学会、输入给错了材料、还是图连错了边，进而决定该\emph{改模型、改输入、还是改下一轮处理哪些对象}。因此本章要解决的核心问题是：怎样把“准备模型、组织调用、产生预测、根据反馈改进”连成一条\emph{可追溯}的链路来获取任务候选，同时始终保留一批不参与方案选择、专供独立核验的依据——否则我们将无法判断模型究竟是真的学会了，还是只是在反复迎合用来调它的那批样本。

本章用一个贯穿始终的花卉识别任务把各环节串起来：目标是为一张张照片补上物种名称；当可见证据不足以辨认时，允许模型弃权，但必须像第~\ref{chap:knowledge-acquisition-human}章的人工标注那样保存弃权依据，而不是硬猜。值得注意的是，现实中可用的起点资源可能差别很大——也许只有少量人工标注的照片，也许已有一个现成的植物识别模型，也许只有一个能接收图文指令的通用大模型。起点不同，“准备一个可用模型”所需的工作就完全不同。

但无论起点如何，对\emph{记录}的要求始终一致：每一条预测输出都必须能对应到具体的模型版本、具体的预测输入版本、以及具体的预测机制。有了这三层版本信息，复核者面对一条错误标签时，才能一步步追查错误究竟是从哪一个环节进入的；缺少任何一层，错误都会变成无法复现的“黑箱偶发”。

本章的结构顺着这条依赖链自然展开，共四步。第一步（4.1节）先得到一个\emph{可调用}的模型：评估、训练或适配它，并用小批量试运行确认能力边界。第二步（4.2节）为模型\emph{组织这一次的输入}：同一张照片，给不给他类别表、参考照片、外部名录，会显著改变它能给出的判断。第三步（4.3节）由\emph{预测机制}真正产出候选样本标注：直接用参数映射、参考相似样本、或在图上传播标签。第四步（4.4节）根据复核结果\emph{选择下一轮如何修改}：更新模型、更新输入，还是调整下一批处理对象。

整条链路中，样本标注仍严格按定义~\ref{def:knowledge-acquisition-record}分开保存内容、来源动作、执行状态与验证状态。请读者带着两条贯穿全章的提醒阅读：一次\emph{成功的弃权}（模型正常运行后判断证据不足）不等于调用失败（程序出错）；一条已经给出的物种名也不等于已经核验为真——“有输出”和“输出正确”永远是两回事。

\section{标注模型的准备}
\label{sec:knowledge-acquisition-model-preparation}

面对一堆尚无物种答案的照片，第一步不是立刻调用，而是先\emph{找出一个能按本任务要求作答的模型}。本章把承担这一职责的模型称为\term{标注模型}（labeling model）——它泛指任何“为待处理对象提供候选监督信号”的模型，并不特指某种网络结构：它可以是一个专用图像分类器，也可以是一个能读图的通用大模型。

准备一个标注模型，要依次回答三个问题，对应本节的三个小节：\emph{有哪些可行方案、各自要做什么准备}（方案选择）；\emph{需要的话，怎样训练或适配它的参数}（训练与适配）；\emph{得到可调用版本后，它在当前任务上到底行不行}（用少量照片试标确认能力）。这三步分别交付三样东西：一批候选方案、一个确定版本的模型、以及“采用它还是继续修改”的决策依据。

这里最容易有的误解是“是不是一定要训练一个模型”。是否需要训练，取决于\emph{现有模型与当前任务之间到底差了什么}，而不取决于这个模型来自哪里、是专用还是通用、参数量多大。一个专用识别模型若用的是另一套物种表，类别对不上，可能仍需适配；一个通用模型若读完任务说明和几个示范就能按要求作答，可能完全不必更新参数；而想从头训练一个模型，若手头标签太少，也可能根本训不起来、只能暂缓。

图~\ref{fig:knowledge-acquisition-model-preparation}把这个判断画成分支：不要求改参数的方案（直接复用、靠示范作答）可以直接进入试标；要求改参数的方案则先训练或适配再试标。试标之后，反馈还要能沿分支\emph{退回}到对应的准备环节——如果问题出在接口，回去改接口；出在标签口径，回去改标签或类别体系；出在任务范围本身，回去重新界定任务，而不是一律归咎于“模型不够强”。

\begin{figure}[htbp]
  \centering
  % 图4-1：标注模型的准备与试标决策，按112 mm正文宽度设计。
\begin{tikzpicture}[
  x=1mm,
  y=1mm,
  font=\figurefont\footnotesize,
  text=BookInk,
  prep input/.style={book node,draw=FigureInput,fill=FigureInput!7,line width=0.8pt,text width=25mm,minimum height=11mm,inner sep=1.5mm},
  prep process/.style={book node,draw=FigureModel,fill=FigureModel!7,line width=0.8pt,text width=28mm,minimum height=11mm,inner sep=1.5mm},
  prep decision/.style={book node,draw=BookAmber,fill=BookAmber!7,line width=0.9pt,text width=24mm,minimum height=11mm,inner sep=1.5mm},
  prep result/.style={book node,draw=BookMuted,fill=BookPaper,line width=0.7pt,text width=25mm,minimum height=11mm,inner sep=1.5mm},
  prep flow/.style={knowledge arrow,draw=BookMuted,line width=0.8pt,rounded corners=1mm},
  prep feedback/.style={knowledge arrow,draw=BookAmber,line width=0.8pt,dashed,rounded corners=1mm},
  edge text/.style={fill=white,inner sep=0.8pt,font=\figurefont\scriptsize,text=BookMuted}
]
  \node[prep input] (requirements) at (16,0) {任务要求\\类别、边界、弃权};
  \node[prep input] (models) at (55,0) {可用模型\\专用或通用};
  \node[prep input] (labels) at (94,0) {初始标签\\样本与覆盖};

  \node[prep process] (select) at (55,-19) {选择候选方案\\并记录资源约束};
  \node[prep decision] (need) at (55,-37) {是否需要\\模型参数训练？};
  \node[prep process] (adapt) at (87,-51) {训练或适配\\形成模型版本};
  \node[prep process] (pilot) at (55,-66) {试标确认\\能力范围与失败类型};

  \node[prep result,draw=BookAmber] (source) at (17,-86) {补充其他来源\\或重定任务};
  \node[prep result,draw=FigureOutput] (accept) at (55,-86) {采用当前方案\\进入批量获取};
  \node[prep result,draw=FigureModel] (continue) at (93,-86) {继续适配\\再做试标};

  \draw[prep flow] (requirements.south) -- (16,-10) -| (select.west);
  \draw[prep flow] (models.south) -- (select.north);
  \draw[prep flow] (labels.south) -- (94,-10) -| (select.east);
  \draw[prep flow] (select.south) -- (need.north);
  \draw[prep flow] (need.south) -- node[edge text,right] {否} (pilot.north);
  \draw[prep flow] (need.east) -- node[edge text,above] {是} (72,-37) |- (adapt.north);
  \draw[prep flow] (adapt.south) |- (pilot.east);
  \draw[prep flow] (pilot.south) -- (accept.north);
  \draw[prep flow] (pilot.west) -- (36,-66) |- (source.north);
  \draw[prep flow] (pilot.east) -- (75,-66) |- (continue.north);

  \draw[prep feedback] (source.west) -- (1,-86) |- (select.west);
  \draw[prep feedback] (continue.east) -- (110,-86) |- (adapt.east);
  \node[edge text,anchor=west] at (3,-29) {资源或要求变化};
  \node[edge text,anchor=east] at (108,-42) {新训练版本};
\end{tikzpicture}

  \caption{标注模型的准备与试标决策。上方三类资源共同限定候选方案；“是否需要模型参数训练”分支决定直接试标还是先适配。下方反馈返回需要改变的步骤。试标用于准备阶段的方案选择，不能代替方案冻结后的独立验收。}
  \label{fig:knowledge-acquisition-model-preparation}
\end{figure}

\subsection{模型方案的选择}
\label{sec:knowledge-acquisition-model-scheme-selection}

要在多个模型之间做选择，不能拿“谁准确率高”这种笼统印象比较，而应先把\emph{本任务对模型的具体要求}写清楚，再逐个核对模型是否满足。需要明确三件事。

其一是\emph{预测输入}：模型会收到什么？可能只给一张照片，也可能同时给地点、季节、拍到了植株哪个部位等元信息；输入越丰富，对模型接口的要求越高。其二是\emph{预测输出}：希望它回什么？可以要求只给物种表中的一个名称，也可以允许给多个候选并附置信度；同时必须约定证据不足时怎样返回“未知”、怎样弃权，而不是逼它在没把握时也硬给一个名字。其三是\emph{范围}：任务覆盖的是本地那几十种常见花卉，还是不预先限定物种的所有植物照片？二者要求的类别范围天差地别，前者一个小分类器也许就够，后者则接近开放式识别。

这三件事确定以后，才有客观依据判断一个现有模型是否\emph{支持}本任务：它能不能接收所需的观测模态（图像、文本、元信息），它的类别含义与我们的物种表对不对得上，它的输出字段能不能承载我们要的候选与弃权标记。范围、输入、输出没定清就选模型，等于没说清需求就采购工具。

在统一的要求下，可以按“需要做多少准备动作”把候选方案归为三类。

第一类是\emph{从头构建}：当没有任何可复用的现成模型、但手里有任务样本和标签时，可以自己训练一个模型。它的好处是可以完全按任务定制，代价是需要足够的标签和训练投入。第二类是\emph{直接复用}：如果已有模型的任务语义、输入条件与当前任务相容（比如已有一个用同一套物种表训练的花卉识别器），可以先不改参数、直接拿来试运行。第三类是\emph{适配}（迁移学习的典型用法）：手头有一个在大量数据上预训练过的模型，它学到的视觉表示是有用的，但类别集合、样本分布或接口与当前任务不完全一致，于是在它的基础上做针对性更新，让它按当前任务工作。

要避免两个想当然。其一，适配\emph{不是万能胶}：它不保证只更新少量参数就能弥合任意大的差异——如果预训练任务与花卉识别相去太远，或缺口在“根本看不见所需部位”，再怎么微调也补不上。其二，“专用/通用”和“模型大小”都不能直接决定该用哪类方案：专用模型和通用模型都可能可以直接复用、也可能都需要适配；参数量大也不代表它天然适用于你的物种表。为减少这种判断的主观性，负责任的模型通常附有\emph{模型卡}（model card），把预期用途、评价条件和已知限制写成一份模型报告，供使用者核对“这个模型是否适合我的场景”\scite{mitchell2019modelcards}。

确认方案在技术上相容之后，还要算清两类资源账，因为“能做”不等于“划算、可持续”。表~\ref{tab:knowledge-acquisition-model-preparation}把成本分成两段。\emph{一次性投入}发生在准备阶段：建设标签、训练或适配模型、改造接口、以及试标试验所花的人力与算力。\emph{持续条件}发生在长期使用中：每次推断占用的资源、调用服务的配额或限额、以及服务版本不断更新带来的维护。

三种方案在这两本账上各有取舍。直接调用一个现成服务可能几乎省掉初始训练，但它把成本转移到了每一轮：照片或特征要持续传输、受调用限额约束、服务方一旦升级版本返回结果可能悄然改变，你对模型本身几乎没有控制权。从头构建需要一笔前期训练投入，却通常换来对模型参数和部署范围的完全掌控。直接复用主要只需调整调用配置和输出约束；适配则要额外决定“允许改动哪些参数”（只改最后一层，还是连骨干一起微调），改动范围越大，需要验证的东西越多。

需要强调，表中列的是\emph{成本构成和控制范围的差异}，并不是一张未经测量的“谁优谁劣”排名表——真实的费用和效果必须结合具体任务实测，不能凭方案名称断定。

\begin{table*}[htbp]
  \centering
  \small
  \setlength{\RaggedRightParindent}{0pt}
  \renewcommand{\arraystretch}{1.15}
  \begin{tabularx}{\linewidth}{@{}>{\RaggedRight\arraybackslash}p{20mm}*{4}{>{\RaggedRight\arraybackslash}X}@{}}
    \toprule
    准备方案 & 起点与准备 & 一次性投入 & 持续条件 & 重点试验 \\
    \midrule
    从头构建 & 任务标签；训练并接入模型 & 标签、训练与调参 & 算力、维护与重训 & 标签覆盖、过拟合 \\
    直接复用 & 已有模型；对齐接口 & 接入、试标与合规检查 & 服务、限额与版本 & 语义错配、接口漂移 \\
    适配已有模型 & 预训练模型参数；确定更新范围 & 适配训练与能力检查 & 推断资源与再适配 & 负迁移、原能力下降 \\
    \bottomrule
  \end{tabularx}
  \caption{模型准备方案的资源与适用条件。按行比较起点、一次性投入、持续条件与待检验风险；三行描述准备动作，同一已有模型可由直接复用转入适配。接口漂移指服务版本或返回约定改变，负迁移指已有表示用于当前任务后反而降低任务表现；它与原模型能力下降是不同的检查对象。各成本仍需实际测量，表中不预设价格或效果名次。}
  \label{tab:knowledge-acquisition-model-preparation}
\end{table*}

把三种方案落回花卉任务，可以看到不同起点对应不同选择。如果只有少量物种标签，既可以试一个轻量分类器，也可以复用一个已有的视觉编码器：先把照片编码成特征向量，再只训练或微调上面的分类部分，从而用很少标签借力预训练表示。如果已有一个恰好使用同一物种体系的现成识别模型，最省事的做法是先不加训练、直接拿它试标一批，看效果再说。如果手头只有一个通用的图文多模态模型，则要先检查它能否读懂任务说明、仿照示范、并严格在当前类别表内作答——通用能力不等于开箱即用地契合你的具体输出约束。

无论选哪种方案，试标时都要\emph{刻意挑难的来考}：容易混淆的近缘物种、拍糊或缺部位的低质量照片、以及根本不在类别表里的对象。要专门检查三类问题：有没有未覆盖的类别、输入稍有变化结果是否稳定、以及有没有“高置信错误”。这里的高置信错误特别危险也特别值得查——它指模型给出了很高的接收得分、看起来十分确定，答案却与参照判断不符；这种错误若不被发现，最容易被当成可靠标签直接采用。最后要记住，方案选择阶段得到的只是一个\emph{附带资源限制和假定适用范围的候选}，它能不能正式采用，仍要经过下一小节的试标确认，不能在选择时就拍板。

\subsection{标注模型的训练与适配}
\label{sec:knowledge-acquisition-model-training-adaptation}

假设已经有一个能把照片编码成有用视觉表示的预训练模型，但它的输出层还认不全当前这套物种表。这时缺的不是再写一段调用说明（说明无法改变模型的判断），而是一个\emph{按任务答案去调整模型参数的准则}。

这个准则的基本思想很直接：用每一张训练照片，比较模型当前预测与它的真实物种标签差多少，这个差距用\emph{损失}来衡量；然后在“允许更新的参数范围”内，寻找一组让平均损失尽量小的参数。先把记号列全：训练样本集为$\mathcal{D}_{\mathrm{tr}}=\{(\bm{x}_i,y_i)\}_{i=1}^{n}$（$n>0$），其中$\bm{x}_i$是第$i$张照片的特征、$y_i$是它的任务答案；模型参数记为$\bm{\theta}$，开始训练前的预训练参数为$\bm{\theta}_0$。

适配时还有一个常见诉求：我们既想让模型学会新任务，又怕改动太大、把它原来具备的能力（比如通用的视觉辨识力）破坏掉。为此可以在平均损失之外再加一项\emph{惩罚}，约束新参数不要离预训练参数$\bm{\theta}_0$太远，用以表达“尽量保留已有能力”的偏好。于是训练目标写成
\begin{equation}
  \widehat{\bm{\theta}}
  \in
  \argmin_{\bm{\theta}\in\Theta_{\mathrm{upd}}}
  \frac{1}{n}\sum_{i=1}^{n}\ell\!\left(f_{\bm{\theta}}(\bm{x}_i),y_i\right)
  +\lambda\Omega(\bm{\theta},\bm{\theta}_0),
  \label{eq:knowledge-acquisition-adaptation-objective}
\end{equation}
逐项解释这个目标。$\Theta_{\mathrm{upd}}$是\emph{允许选择的参数配置集合}：哪些参数被冻结、哪些允许更新，都由它规定；被冻结的坐标在整个优化中保持原值。$f_{\bm{\theta}}$是参数为$\bm{\theta}$时模型作出的预测，$\ell$衡量一次预测与答案之间的差异；$\Omega(\bm{\theta},\bm{\theta}_0)$表达“离起点多远”的保留偏好，系数$\lambda\geq0$控制这项惩罚的强度——$\lambda$越大，优化越保守、越不敢偏离预训练模型。损失取训练样本的\emph{平均}而非总和，是为了衡量模型在这批任务样本上的\emph{总体}拟合程度、且不随样本数量放缩；但平均值低只代表整体拟合好，绝不保证每张照片都预测正确。

不同训练策略其实只是这个式子的不同取法。\emph{从头训练}时，可以让全部参数都可变、并去掉相对$\bm{\theta}_0$的惩罚（允许任意$\bm{\theta}$、$\lambda=0$）。只训练\emph{输出层}时，表示部分参数全部冻结、只让最后的分类层变化。\emph{微调}（fine-tuning）则介于两者之间，继续更新部分或全部预训练参数。

使用这个目标还要有三点数学与工程上的清醒认识。第一，只有当允许集合非空、且目标确实能取到最小值时，最小点集合$\widehat{\bm{\theta}}$才有意义；实际训练几乎总是用梯度下降\emph{近似}求解，必须另行检查收敛与停止条件，不能因为目标写成了最小化形式，就保证到达了全局最小点。第二，加了保留惩罚只代表“\emph{倾向于}不破坏原能力”，并不\emph{保证}原能力一定保住；有没有发生“灾难性遗忘”，要另找一批对应原能力的样本去实测。第三，这个优化能学到什么，完全由喂进来的训练对象和标签决定——下一段就讨论数据本身的陷阱。

训练目标刻画了“怎样优化”，但\emph{能拟合出什么}，根本上取决于放进训练集的对象和标签。这里有两个初学者常踩的坑。

第一个是数据划分造成的“假性泛化”。如果同一株植物的几张连拍照片被随机打散、分别落进训练集、方案选择集和最终测试集，那么模型在训练时其实已经见过几乎一模一样的照片；这时测试分数再高，也只说明它记住了这几株植物，\emph{不足以}证明它能识别没见过的新植株。正确做法是按\emph{植株或采集批次}成组划分——同一株的所有照片必须待在同一组，再核对各组在物种、拍摄条件和“观测不足”情形上的覆盖是否均衡。

第二个是标签的质量与覆盖。错误标签会把损失直接拉向错误答案，优化只会更努力地去拟合这些错误，训练再久也不会自动纠正；而若某个物种在训练集里压根没有标签，模型就没有收到任何“学习这个类别”的任务信号，延长训练时间也变不出它来。这两类缺失都不能靠“多训几轮”解决，只能回到数据建设。

数据按用途分成训练、验证、测试三组后，还要严格分开\emph{三种不同的判断}，混用会让评价失真。训练损失只用于观察“优化有没有在拟合当前目标”，它回答的是学得动不动，不是泛化好不好。验证集用于在开发过程中作\emph{选择}：冻结哪些参数、$\lambda$取多大、何时停止训练，都靠比较候选在验证集上的表现决定。测试集则必须被“封存”，只有在上述所有决定都冻结之后才用一次，用来估计最终方案对新对象的表现。

评价一个候选时也不能只看单一数字。假设某个方案在验证集上把常见花卉的正确率提高了，但同时让“低清照片上本应弃权却强行作答”的比例也上升了，那么这是一好一坏两种变化，必须同时记录——只报平均正确率会把后者的退化掩盖掉。最后，一旦选定版本，就应把它与训练样本版本、类别体系、参数更新范围和推断配置\emph{绑定}记录，这样将来产生的每条样本标注才能注明“它究竟由哪个版本、在什么配置下产出”，错误才可追溯。

也要避免“为训练而训练”。如果一个已有模型在试标中已经满足接口和任务要求，那就直接冻结它的版本和调用配置使用即可，不必画蛇添足地再做参数训练——多余训练反而可能破坏原本可用的状态。反过来，有些缺口\emph{不是微调能解决的}：如果模型根本无法接收所需类型的图像输入，或其结构承载不下任务所需的类别规模，那么再多样本、再多轮微调也补不上。这类问题应当诚实地退回到表~\ref{tab:knowledge-acquisition-model-preparation}重新选择方案（换一个支持该模态/规模的模型），而不能因为“训练损失确实在下降”就自欺欺人地继续使用一个本质上不胜任的模型。

\subsection{标注能力的试验与确认}
\label{sec:knowledge-acquisition-model-pilot-validation}

训练好的模型有了一个可调用版本，但“在开发集上损失低”并不等于“在真实标注流程里可靠”。在扩大使用之前，应先让它\emph{严格按正式流程}处理一小批照片做演练，这就是\term{模型试标}（model pilot annotation）。它与第~\ref{sec:knowledge-acquisition-human-pilot}节的人工试标目的相同——都是在投入大批资源前小规模跑通整个任务、暴露问题；区别在于作答主体是模型而非人，因此故障的表现和修正方式也不同（模型会出现接口报错、输出无法解析等人工不会有的问题）。

试标样本要刻意覆盖五类情形：典型的正常照片、易混物种的边界案例、低质量输入（失焦、遮挡、缺部位）、类别表之外的对象、以及按规范应当弃权的情形。为了让问题能分门别类，每条样本标注都要保存：样本标识、模型版本与输入版本、模型返回的\emph{原始输出}、输出能否被成功解析、解析出的候选、以及对应的参照答案。

把三类失败\emph{分开计数}尤其关键：调用根本没完成（接口故障/超时）、返回的格式不合法（比如不是约定的JSON，属于解析问题）、格式正常但内容判错（属于判断问题）。这三类问题的责任环节和解法完全不同——接口问题修系统、解析问题修输出约定、判断问题才需要改模型或补数据；若只统计一个笼统的“失败率”，就无法决定该修哪里。

根据这些分项结果，可以对下一步做分流决策。如果在预定的适用范围内，内容错误率、弃权情况和单次调用成本都满足事先设定的“批量准入”门槛，就采用当前方案、准备扩大使用。如果问题明确可以通过补标签、调整模型参数或修复接口解决，则回到适配环节继续改进。如果根因是“关键观测根本不可见”（照片里就没拍到鉴别部位）或“任务超出了模型可确认的范围”，那不是再训练能解决的，应改为补充其他来源——转人工、加规则、或提供外部文档。

这里有一条与第2章人工试标一致的纪律：试标集一旦被反复用来挑选模型、调阈值、改提示，它就\emph{变成了开发依据}，不能再充当“方案冻结后”的独立证明。用同一批样本既做选择又做验收，等于既当运动员又当裁判。具体的准入指标如何设定、怎样抽样、分数的区间估计如何计算，将在第~\ref{chap:knowledge-quality-assessment-correction}章展开，本节只负责确认“能力范围到哪里”。

% 长作者引文从段中开始，需同时预留段首至标记和完整边注的高度。
\Needspace{100mm}
一种常见的后续用法，会在“产生候选”和“训练新模型”之间再增加一层依赖关系：让一个模型先为未标注样本提供答案，再用这些答案去训练另一个模型。提供目标答案的那个模型在这里称为\term{任务教师}（task teacher），因为它给“学习者”提供了具有既定任务含义的目标；被训练的模型则称为学生。

这两个名称只描述信息流向，\emph{不}暗示谁的模型更大、谁更准确：教师可以是一个小而专的模型，学生反而可能是一个已经具备强大预训练能力的大模型，学生最终的行为也不完全由教师标签决定。教师弱、学生强时会发生什么？弱到强泛化研究在语言、棋类和奖励建模等实验中观察到一个值得注意的现象：一个强预训练学生，用弱模型给出的标签微调后，有时反而能\emph{超过}那个弱监督者自身的表现；但朴素的微调并不能恢复一个真正强模型的全部能力\scite{burns2023weaktostrong}。这个结果说明学生\emph{有时}能纠正教师错误，但它要求我们在自己的任务上分别、独立地检查教师和学生的表现，绝不能据此推出“学生总会自动纠正教师”的乐观结论。

用花卉任务把试标的检查点具体化：可以同时放入山茶与茶梅这对易混物种（检查物种混淆）、只有花瓣的局部照片（检查观测不足时是否弃权）、以及类别表之外的植物（检查范围识别）。如果模型只在“典型正面清晰照片”上表现可用，那就只能确认它在\emph{这一范围}内可用，不能把结论外推到所有照片；当本应弃权的照片被强行赋予类别时，要补上弃权接口和弃权原因，或为这些对象增加其他来源。

如果准入决策依赖模型的置信度分数，还必须检查分数与真实正确率之间的关系，这就是\term{概率校准}（probability calibration）要回答的问题。它考察的是：当我们把某个接收分数\emph{解释成“正确的概率”}时，落在该分数段的那些样本，其实际答对比例是否真的与平均预测概率相符；校准关心的是这种“分数—频率”的一致程度，而不是分数本身谁大谁小。例如，所有给出0.9分的预测，理应在长期里约90%正确，若实际只有60%，分数就是被严重高估、未校准的。

做这项检查时，分数段怎么划、参照答案是什么、允许多大偏差，都要随评价范围事先确定；而且有限样本算出的正确比例本身带有估计误差。还须记住校准关系的边界：即使整体校准良好，也不保证某一条高分答案一定正确；校准关系也不会自动迁移到换了地区、换了模型版本之后。因此最终的能力确认报告应保存“可用范围、各类失败类型、以及什么情况下必须转人工接管”，而不能只留一个笼统的总体正确率。

\section{标注输入的组织与构造}
\label{sec:knowledge-acquisition-input-organization}

模型选定之后，下一个同样关键、却常被忽视的问题是：\emph{这一次作答，究竟允许模型看到什么？}要意识到，模型的答案不是它孤立地“凭空知道”的，而是\emph{它这一次输入里有什么}的函数。同一张花瓣局部照片，附不附类别表（告诉它只能在山茶、茶梅、月季里选）、附不附同一植株的叶片照片（补上鉴别所必需的部位），会让它可能作出的判断截然不同。把输入准备好，往往比反复更换模型更能决定成败。

为便于组织和排查，我们按\emph{作用}把一次输入中的内容分成四类。\textbf{任务要求}规定“要回答什么问题、答案的含义和格式是什么”；\textbf{待处理样本及上下文}提供“当前这一个对象的观测”；\textbf{示范案例}通过别的已完成例子，展示“规范在具体实例上是怎么用的”；\textbf{辅助证据}则从样本之外补入外部事实或对象关系。这四类回答的问题不同——前两类主要说明“判断谁、答成什么样”，后两类主要提供“可以参考什么”。

需要说明两点。其一，四类功能可以自由组合，但并不要求每个模型都接收全部四类：一个固定的分类器可能只吃样本、用配置文件表达任务要求，示范和证据则由外围检索组件提供——功能是按作用划分的，与具体用什么接口实现无关。其二，正因为输入内容如此影响结果，后面会看到：必须记录每个片段起了什么作用，否则结果变化时无法归因。

图~\ref{fig:knowledge-acquisition-input-composition}用两种线型把输入分成两组：虚线传入的是“作答约定与示范”（教模型怎么答），实线传入的是“当前观测与外部证据”（供模型据此作答）。

这里藏着一个容易混淆、却对归因很重要的情况：\emph{同一个片段可能同时承担两种作用}。一张带答案的参考照片，既可能在示范“输出字段该怎么填”（格式引导），又可能真的为当前这株植物提供了“它和谁长得像”的鉴别证据。如果它确实一身二任，就必须在记录里同时注明这两种用途。否则一旦预测结果改变，我们将无法辨认变化究竟来自“格式示范让输出更规范了”，还是“新参考提供了实质的物种证据”——这两种原因的改进方向完全不同。

为了让四类输入的作用具体可感，下面围绕同一个对象S042，\emph{依次}构造一段完整请求：先给任务要求，再加入当前观测，再加入一个示范，最后加入辅助证据；把这四段按顺序拼起来，就是本例模型实际收到的完整预测输入。

为方便只凭纸面文字也能核对“哪些判断有依据、哪些没有”，我们把照片里\emph{可见的内容}用一段观测文字描述出来（而不是真的嵌入图片）。需要事先声明：本章的对象编号、植物名录、日期和所有模型输出都是\emph{教学构造}，没有真正调用任何模型，这些名称也不用于现实中的植物鉴定；它们的唯一目的是把“模型能看到什么、据此能/不能推出什么”演示清楚。

\begin{figure}[htbp]
  \centering
  % 图4-2：一次模型标注中的四类预测输入，按112 mm正文宽度设计。
\begin{tikzpicture}[
  x=1mm,
  y=1mm,
  font=\figurefont\footnotesize,
  text=BookInk,
  input block/.style={book node,draw=FigureInput,fill=FigureInput!6,line width=0.8pt,text width=37mm,minimum height=12mm,inner sep=1.4mm},
  model block/.style={book node,draw=FigureModel,fill=FigureModel!8,line width=0.9pt,text width=25mm,minimum height=18mm,inner sep=1.5mm},
  output block/.style={book node,draw=FigureOutput,fill=FigureOutput!6,line width=0.8pt,text width=17mm,minimum height=18mm,inner sep=1.2mm},
  convention/.style={knowledge arrow,draw=BookAmber,line width=0.8pt,dashed},
  evidence/.style={knowledge arrow,draw=BookBlue,line width=0.9pt},
  edge text/.style={fill=white,inner sep=0.7pt,font=\figurefont\scriptsize,text=BookMuted}
]
  \node[input block,draw=BookAmber] (requirement) at (21,0) {任务要求与预测输出约定\\Q1：类别C1、弃权、字段};
  \node[input block,draw=BookAmber] (demonstration) at (21,-18) {示范案例与参照答案\\S017：观测不足，弃权};
  \node[input block] (sample) at (21,-36) {待处理样本与必要上下文\\S042：花瓣局部，无叶片};
  \node[input block] (support) at (21,-54) {辅助证据与关联事实\\E03：名录R3的地块名录};

  \node[model block] (model) at (70,-27) {标注模型\\版本$M_5$};
  \node[output block] (output) at (101,-27) {S042\\样本标注\\未知、弃权\\证据与版本};

  \draw[convention] (requirement.east) -- node[edge text,above] {约定} ([yshift=6mm]model.west);
  \draw[convention] (demonstration.east) -- node[edge text,above] {示范} ([yshift=2mm]model.west);
  \draw[evidence] (sample.east) -- node[edge text,below] {当前观测} ([yshift=-2mm]model.west);
  \draw[evidence] (support.east) -- node[edge text,below] {外部证据} ([yshift=-6mm]model.west);
  \draw[knowledge arrow,draw=FigureModel,line width=1.0pt] (model.east) -- (output.west);
\end{tikzpicture}

  \caption{一次模型标注中的四类预测输入。虚线传入作答约定和示范，实线传入当前观测与外部证据；S042、O042与E03的标识限定各片段的对象和用途。四类内容按功能区分，不要求所有模型使用同一种提示接口。}
  \label{fig:knowledge-acquisition-input-composition}
\end{figure}

在把四段拼成最终请求之前，还要对每段内容做四项检查，它们分别防一种常见错误。一是\emph{对象归属}：这段材料说的是不是当前对象？把别株植物的信息塞进来会直接导致张冠李戴。二是\emph{空间预算}：模型一次能读入的长度有限，一段冗长材料可能把更关键的观测挤掉。三是\emph{内容一致}：不同片段之间是否互相矛盾（如名录说一种、示范暗示另一种）。四是\emph{防泄漏}：有没有不小心把“实际作答时模型不可能拿到”的参照答案写进输入——用未来才知道的答案提示现在的判断，等于作弊。

还要在流程上区分两个环节：\emph{检索}负责“把可能相关的候选材料找出来”，\emph{组织}负责“确定其中哪一段对应当前对象、以及怎样放进输入”。把这两步的责任分清，出错时才能定位：模型用错了证据，到底是因为检索阶段就找回了错误材料，还是材料本身正确、却在组织阶段被关联到了错误的对象？混为一谈就无法对症修复。

\subsection{任务要求与输出约定}
\label{sec:knowledge-acquisition-task-output-contract}

先看四类输入中的第一类。对S042，我们要的是“在一张\emph{指定物种表}内作出的判断”，而不是让模型随意报出任何植物名称——没有边界的答案无法用于训练，也无法核验。因此任务要求必须先讲清几件事：使用哪一份物种名录、杂交种和别名如何处理（要不要折算回规范名）、多张照片如何对应到同一对象，以及在什么条件下应返回未知（证据不足）或不适用（对象出界）。

任务要求和\emph{输出格式约定}要分开看：前者决定模型“判断什么”，后者决定程序“怎样把结果读回来”。输出约定要规定如何保存对象标识、内容、来源动作、弃权原因和证据引用；模型若还给出能解释含义的置信度分数，则作为附加字段另存，不和内容混淆。下面是给一个能接受自然语言指令的模型所写的说明，读者可对照体会“任务要求”与“输出格式”是如何同时交代清楚的：

\begin{quote}
  请按规范Q1，仅根据本请求提供的观测和证据，判断S042所指植株的物种。任务范围限于类别表C1中的山茶、茶梅、月季。若作答，只返回其中一个规范名称，不自行补入别名、杂交种或其他类别。

  只有证据足以支持唯一类别时才作答，此时内容为物种名称，语义状态为“已给出”，动作为“作答”。观测不足、材料冲突或只能缩小候选范围时，内容留空，语义状态为“未知”，动作为“弃权”。只有已确认对象不属于本任务范围时，才记为“不适用”，内容和动作均留空。不得臆测不可见部位，也不得把某地出现过的物种直接当作该地每株植物的物种。

  只返回一个JSON对象，字段固定为object、content、semantic、action、reason、evidence，依次表示对象标识、内容、语义状态、来源动作、原因和证据编号列表。空值写为null，reason简述判断依据，不虚构概率分数。后附S017是完整示范，不是本次待作答对象，也不是S042的证据。
\end{quote}

读这段说明要注意它在职责上的克制：它只让模型交付\emph{内容}和\emph{来源动作}（以及原因、证据编号），并没有让模型自己声明“这次调用成没成功”“答案对不对”。执行是否完成，应由调用程序写入运行日志；指定属性是否核验通过，应由后续独立检查记录。模型在reason里自称“判断正确”“高可信”都不作数——自报成功不能代替运行日志，自报正确也不能代替核验。这种“生成者不兼任验收者”的分工，是保证结果可追溯的前提。

上面那段自然语言说明，只能发给能读指令的模型；一个结构固定的分类接口读不懂散文，需要把同样的要求\emph{落实成配置}。用一个具体的数值例子说明配置如何工作。假设类别表C1把输出索引与类别对应为$1\mapsto$山茶、$2\mapsto$茶梅、$3\mapsto$月季；模型对S042输出一个得分向量$(0.55,0.40,0.05)$，其中最高分0.55对应山茶。但“取最大项”这一步只产生了一个\emph{类别候选}，还没有决定要不要接收它。

为此在模型外面再包一层接收规则（同样是配置）：教学上门槛设为$0.70$，只有最高分不低于门槛才作答，否则返回未知并弃权。本例最高分0.55低于门槛，因此正确结果是弃权，而\emph{不是}退而求其次地把0.55对应的山茶、或干脆把“未知”编成第四个类别索引——后者会把“没把握”伪装成一个物种，污染训练。要特别强调门槛的身份：它在这里只是一条\emph{用来区分“接收”还是“弃权”的构造配置}，选0.70并不保证准确率达到70\%。照片和特征仍按分类器自身的接口传入；示范、外部证据这类它读不了的内容，则由支持这些功能的外围组件单独处理。

任务规范日后会变化，要区分两种变化的代价。如果只是输出字段改了个名字、语义不变，旧的样本标注或许可以随字段重命名迁移；但如果某个类别被拆分（山茶拆成数个品种）或含义改变，旧标签所指就不再明确，通常必须重新预测或人工复核。

文字规范是否被一致理解，光看定义不够，还要配上\emph{边界实例}去校准。山茶和茶梅应各有清晰正例、互相混淆的反例、以及无法区分时的未知例，用它们具体演示“叶缘、花期这类线索，只有在对应部位可见、或拍摄时间已知时才允许使用”，而不能凭名字联想。

当任务从“图像分类”换成别的类型，输出空间的整个形态都要重新约定。例如实体抽取任务，输出不再是一个照片类别，而是“原文中的片段＋实体类型”，这时还必须规定抽出的片段如何对应到规范化名称（多个别名怎么归一）。按答案空间的开放程度，大致有三种设定，要在规范里讲清：\emph{封闭分类}只能在给定的有限类别集合中选择；\emph{开放文本生成}允许出现集合之外的内容；\emph{多解任务}则还要规定多个候选之间是否互斥、能否并列。任务对象或输出空间一旦改变，记录里那个“答案”字段的含义也必须\emph{随之重新说明}——同一个字段名，在分类、抽取、生成任务里装的东西完全不同。

\subsection{待处理样本与必要上下文}
\label{sec:knowledge-acquisition-sample-context}

任务要求讲清了“怎么答”，但模型还必须知道“现在判断的到底是哪一个对象”。这一类输入由两部分组成：\emph{待处理样本}本身用来锁定当前对象（是S042这株植物，而不是旁边那株），\emph{必要上下文}则补入解释这个对象所需要、且在真实使用时确实能够获得的观测。

不同模态需要绑定的定位信息不同。照片要对应到原图或一个稳定可查的引用，并注明判断区域、拍摄时间和对象标识；文本要对应到原文中的具体位置和必要段落；语音或时序数据还要限定时间区间。操作上应当\emph{先绑定对象、再据此裁取上下文}：先确认“说的是这一株”，再围绕它取周边信息。这样模型的输出才能准确还原到正确位置，避免把对邻近对象（旁边的植物、上一句话的实体）的判断错误地写回当前样本。

现在向S042请求加入当前观测。用O042标识这一片段，是为了以后追查模型看过什么，不是提前提供物种标签：

\begin{quote}
  当前对象：S042，对应植株P042。观测O042：2026年4月12日在教学园地块P03拍摄的花朵局部，仅见粉红色花瓣；叶片、叶缘、完整花形与株形均不可见。地点和日期来自该照片的采集日志。本请求没有其他照片，也没有这株植物的已知物种标签。
\end{quote}

对照这段构造观测可以看清“提供信息”和“空喊注意”的区别。O042确实给了地点和日期，这两条可以用于匹配花期、地块名录等上下文；但它\emph{明确没有}提供叶缘——叶片不可见。此时若真的补入一张同一植株的叶片照片，那是\emph{实实在在地增加了观测}，模型因此获得了判断叶缘的依据；但若只是把提示措辞改成“请仔细看叶缘”，模型并不会因此获得任何新信息——它无法从一张只有花瓣的照片里看出根本没拍到的叶缘。语言上的强调不能创造视觉证据，这个区分在后文判断“模型该不该弃权”时至关重要。

怎么从原始材料里\emph{裁取}喂给模型的那部分，直接决定哪些线索被保留、哪些被丢弃。花瓣局部特写能把花瓣形态看得很细，却可能同时丢掉叶片和整株形态这些鉴别近缘物种的关键线索；改用整图固然补回了这些信息，却也可能把花盆、背景里的标牌等干扰物一起带入，诱使模型“抄近路”（比如靠背景标牌而非花形判断）。这是一种需要权衡的信息取舍。

文本任务中对应的问题是\emph{切分边界}：把长文档切成片段时，用重叠窗口可以保留跨句、跨段的语义线索（一个实体的描述可能延续到下一段），但同一段被重复读入时，重复出现的实体要按原文位置合并，不能当成两个对象。时序任务还多一层“判断时点”的硬约束：在某一时刻做在线预测，只能使用\emph{当时已经允许获得}的历史信息；若把后来才发生的观测也喂进去，任务就从“当时能否预测”偷换成了“事后回溯”，评价的不再是同一种能力。

当一次输入要合并照片、文本、元信息等多种观测时，\emph{对象、位置、时间三者必须对齐}，否则会产生隐蔽却严重的错误。把植株A的照片配上植株B的拍摄地点，模型给出的答案即使措辞看起来合理，也很可能指向了错误的对象；若某个缺失观测被悄悄填上一个默认值（比如没拍到时间就写一个常见日期），样本标注就会\emph{虚报}模型实际看到的内容，让后续复核误以为它用了某条其实不存在的证据。正因如此，必须如实保存这次预测输入里实际包含的每一项观测、示范和证据。

最后是一条防止评价造假的纪律：\emph{上下文如何整理的规则，必须在查看参照答案之前就冻结}。如果先看到了正确答案、再回头挑选“哪些背景是必要的”，就很可能不自觉地专挑那些能导向正确答案的上下文，相当于把评价对象的答案直接泄漏进输入，测出来的“高准确率”毫无意义。

% 为首段及GPT-3完整作者边注共同预留空间。
\Needspace{125mm}
\subsection{示范案例与参照答案}
\label{sec:knowledge-acquisition-demonstrations}

到这一步，模型已经有了作答规范Q1和当前观测O042，但规范是抽象的文字，模型未必能准确体会“证据不足时到底该怎么填那几个字段”。与其反复解释，不如\emph{直接给它看一个已经做好的例子}。做法是加入\emph{另一个对象}的完整输入和它的参照答案，作为\term{示范案例}（demonstration example）：让模型通过具体实例看懂规则的用法，而这个例子绝不直接给出S042的物种答案。

这种用法对应一个重要概念。某些语言模型能够只依据输入文本里给出的几个例子就照着完成新任务、\emph{完全不更新自身参数}，这称为\term{上下文学习}（in-context learning, ICL）。GPT-3相关工作中的“少样本”设置——把任务说明和若干示范一起放进文本对话、模型据此作答——是这种用法的代表性实例\scite{brown2020}。要抓住它与训练/微调的本质区别：这里改变的只是\emph{这一次调用的输入内容}，模型权重一个字节都没变；因此“加示范有没有用、输出稳不稳定”不能凭模型名气假设，必须在\emph{当前这个具体模型}上用试验确认。

选择示范时要有明确的覆盖意图：要展示的情况应当与规范里的各个\emph{判断分支}一一对应，示范里的字段也必须和正式输入完全一致，模型才会照着同一格式输出。如果所有示范都只给清晰的正面照片，模型就只学会了“有把握”这一种情形，无从知道观测不足时该如何填写；加入一个“未知/弃权”示范，能让它看到弃权分支长什么样——但要清醒，\emph{看到弃权范例并不保证模型在该弃权时就一定正确弃权}，示范只是引导而非强制。下面就为S042的请求加入一个完整的教学示范，其答案是严格按Q1的弃权条件构造的：

\begin{quote}
  示范预测输入：对象S017，植株P017，类别表C1。观测O017：2026年4月10日在教学园地块P08拍摄，仅见失焦的粉红色花瓣局部，叶片和株形不可见；没有辅助证据。按同一规范Q1判断物种。示范预测输出如下：
\end{quote}
\begin{lstlisting}[style=book code,numbers=none]
{
  "object": "S017", "content": null,
  "semantic": "未知", "action": "弃权",
  "reason": "花瓣局部失焦，缺少区分物种的观测",
  "evidence": ["O017"]
}
\end{lstlisting}

这个示范的作用边界要说清楚。它展示的是“证据不足时，content、action、reason这些字段分别怎么写”，是一种\emph{格式与决策风格}的范例。它\emph{没有}声称S017和S042是同一株植物，也没有提供“这两株物种相同”的任何证据——模型不能因为例子里S017弃权，就推断S042也该弃权，或反之。此外，示范里的reason也必须经过检查：如果示范本身给出的理由是错的（比如用了不相关的判断依据），模型会把这套\emph{错误的规范解释}当成模板复制到后续大量调用中，错误被成倍放大。

加示范不是“多多益善”，它会同时改变输入的长度、示范的排列顺序和类别构成，因此不能只按“给了几条”来评估效果。一个已被反复观察到的现象是\emph{顺序敏感}：关于少样本提示的研究在分类实验中发现，仅仅是把同一组示范换一种排列顺序，就可能改变模型的结果；而且在一个模型上试验出的“最佳顺序”，换到另一个模型上未必仍然最好\scite{knowledge-acquisition-prompt-order2022}。

因此对待示范也要像对待超参数一样严谨：保持任务说明和评价口径一致，在\emph{开发样本}上比较不同的示范集合与排列、并保存版本，绝不能拿最终评价样本反复试到满意为止——那又是用测试集做选择。还要算一笔空间账：示范越长越多，占用的输入预算越大，比较方案时必须检查当前真正关键的观测或证据有没有被大段示范\emph{挤出}输入；示范反而盖过了待判断对象，是典型的本末倒置。

最后澄清一个极易混淆的点：同一个参考样本库，可以服务于两种目的完全不同的机制。如果某个固定案例的主要作用是“展示作答规则和字段怎么填”，它是本节的\emph{示范}；如果系统是\emph{根据当前查询临时检索}出带标签的相似照片、并让这些照片的标签真正参与当前答案的计算，那它就是下一小节的\emph{辅助证据}，进一步还会构成第~\ref{sec:knowledge-acquisition-neighbor-prediction}节的近邻预测。二者选择材料的依据不同：示范按“要教会模型哪些规范分支”来选，证据按“哪些参考在内容上支持当前对象”来选。共享同一个底层样本库，并不使这两种用途可以等同对待——它们影响模型的方式和需要核验的环节都不一样。

\subsection{辅助证据与关联信息}
\label{sec:knowledge-acquisition-auxiliary-evidence}

示范教会了模型“怎么答”，但回到S042，它手里仍缺“凭什么把这株判成某个物种”的当前依据——O042只有花瓣，S017只教了弃权格式。这时可以从样本之外补入\term{辅助证据}（auxiliary evidence）：来源可以是文档或文档库、实体关系库、其他带标签的参考样本、或对象之间的关系；它们都属于“当前观测之外、但可能支持本次判断”的材料。

组织辅助证据要回答四个问题：要找什么（检索字段）、怎样确认它与当前对象对应（对象匹配）、取用其中哪一段（范围截取）、以及模型输出怎样回指到原始来源（证据编号，便于复核）。这里要切断一个常见的等号：\emph{检索命中}只说明查询词与某段材料相似、或某个字段匹配上了；这段材料\emph{是否真的支持}一个具体物种答案，是检索之后\emph{另一项独立的判断}。“查到了相关文档”不等于“文档证明了答案”。

在花卉任务里，辅助证据大致有两条检索路线。一是按\emph{照片特征}找视觉上相似、且带物种标签的参考照片；二是按\emph{地点与花期}之类的元信息去查植物名录，看某地块在某季节可能出现哪些物种。

无论哪条路线，取回的每条材料都必须保留四项背景信息：它来自哪里、适用的时间范围、对应的对象标识、以及其上的标签是否经过核验。原因在于这类外部材料有多种不可靠方式：同名的植物可能根本不是同一对象（需消歧）；一份过时的名录可能已不再适用；同一条信息被多处转载，并不增加\emph{独立}依据（看起来有三个来源，其实出自一处）。对文本中的实体，同样要靠上下文消歧，确定同一个名称在当前材料里到底指谁，而不能只因字符串相同就把事实合并——这与第~\ref{chap:rule-based-knowledge-acquisition}章实体链接的告诫一致。

对S042，我们固定检索字段为地块P03和日期2026年4月12日，取覆盖这两项条件的名录片段。把它接在示范后，就完成了四类预测输入的组合：

\begin{quote}
  辅助证据E03：教学园维护名录R3，第8条，名录编写时间为2026年4月1日。适用对象是地块P03，时间范围是2026年3月1日至6月30日。原文：“本期地块名录中有山茶和茶梅；此表不是完整物种清单，未建立逐株标识与物种的对应。”该片段已与R3核对，但没有核验P042的物种。请在这一范围内使用E03，只为当前对象S042返回一次结果。
\end{quote}

逐条核对E03的证据力。P042的拍摄地点和日期确实落在E03声明的适用范围内（地块P03、2026年3—6月），所以这段名录\emph{可以合法地关联}到当前请求，对象与时间这两道关都过了。但它的证据力到此为止，有两个关键限制：其一，名录只说到“地块中有山茶和茶梅”，\emph{没有}建立“具体哪一株是哪个物种”的逐株对应；其二，它\emph{明确声明}自己不是完整物种清单。因此E03既不能确定P042是山茶还是茶梅，也不能排除清单之外的类别。

把O042（只有不可靠的花瓣）和E03（只有地块级、非穷尽的名录）合起来，证据仍不足以支持任何一个唯一物种。按Q1“只有证据足以支持唯一类别时才作答”的规定，正确的构造输出应当是弃权：

\begin{lstlisting}[style=book code,numbers=none]
{
  "object": "S042", "content": null,
  "semantic": "未知", "action": "弃权",
  "reason": "O042缺少鉴别特征，E03没有逐株对应",
  "evidence": ["O042", "E03"]
}
\end{lstlisting}

若实际调用返回这条预测输出，程序将执行状态记为成功，来源记为模型$M_5$，取得条件保存Q1、C1、O042、S017和E03的引用，验证状态先记为候选。复核若确认弃权符合Q1，确认的属性是动作合理，物种内容仍为未知。再看另一条构造预测输出：它的JSON可以解析，却把证据说成了不同范围的事实：

\begin{lstlisting}[style=book code,numbers=none]
{
  "object": "S042", "content": "山茶",
  "semantic": "已给出", "action": "作答",
  "reason": "R3已经证明P042是山茶",
  "evidence": ["E03"]
}
\end{lstlisting}

对比之下，第二条输出的问题就清楚了：它的JSON格式完全合法、字段也填得很满，但它把一条“地块级名录”悄悄\emph{拔高}成了“P042这一株就是山茶”的逐株证明，而E03原文恰恰说没有这种对应。这是一个典型的“证据被过度解读”的错误，应在\emph{内容核验}环节记为已拒绝，同时保留模型的原始输出留证。

这里再次体现状态分离：这次调用本身已经正常完成，不能因为内容被拒绝就把\emph{执行状态}改成失败——程序没出错，是模型的结论超出了证据。两条输出放在一起，决定命运的不是“字段填得齐不齐”（第二条反而更齐全），而是\emph{证据能支持到哪一层}：第一条诚实地停在证据允许的范围内（弃权、合规），第二条越过证据编造了一个确定标签（内容被拒、不能作训练用）。教会模型（和读者）判断“证据支持到什么强度为止”，正是这个例子的用意。

“检索外部材料、再让模型据此作答”这套组合，在大模型场景下有一个通用名称：检索增强生成（retrieval-augmented generation, RAG）。它由两个角色配合：检索器先从外部索引中取回相关文档，生成模型再\emph{依据}这些检索结果（而非仅凭参数记忆）产生文本答案。RAG的价值在于把一部分知识放在模型参数之外、可随时更新的文档里，从而把“外部可检索的记忆”与“参数刻画的预测能力”结合起来\scite{knowledge-acquisition-rag2020}。

但RAG只是给出了一种检索与作答的\emph{组合架构}，它并不替你完成、也不保证辅助证据流程里的任何一项核验。仍要分别检查：材料是怎么找到的（检索质量）、所需证据有没有找全（召回是否充分）、模型是否\emph{严格在材料支持的范围内}作答（有没有脱离材料编造）。模型能流畅地生成一段答案，绝不等于这三件事都已做到——“会说话”与“说得有依据”是两回事。

最后谈谈检索的取舍与边界。扩大检索范围（多取回一些材料）是一把双刃剑：它可能找回原先漏掉的关键证据，但同时会带入无关内容、重复材料，并增加额外的检索与调用成本。因此要事先约定取回多少候选、给每份证据标明它对应哪个样本；当不同材料\emph{互相冲突}时，正确做法是把冲突\emph{保留}下来留待判断，而不是让后取回的材料随意覆盖先取回的。

到这里，本节产出的只是一批\emph{可追溯的材料与关系}，尚未回答“如何综合它们得出一个结论”。这些材料若按相似度直接参与数值计算，就是下一节4.3中的近邻预测；若用它们在对象间连边、联立求解，则是图结构预测；而当多个\emph{来源}（人、规则、不同模型）给出需要调和的结论时，则交给第~\ref{chap:knowledge-fusion}章的融合方法处理。

\section{候选知识的预测}
\label{sec:knowledge-acquisition-candidate-prediction}

预测输入已经绑定到当前对象、材料也备齐了，还差最后一步：这些输入\emph{究竟通过什么计算变成一个答案}？对同一张没有标签的照片，至少有三种自然的算法路径，它们“向谁要答案”的来源根本不同。

第一种是\emph{问模型参数}：直接用训练好的分类器做一次前向计算，答案来自参数里已经学好的输入—类别映射。第二种是\emph{问相似的旧案例}：在带标签的照片库里找与它最像的几张，看那些照片是什么物种、投票决定。第三种是\emph{问它在关系网络中的位置}：把已知与未知的照片一起放进一张图，用“相连的照片倾向于同类”这一约束，联立所有未知照片的类别得分。图~\ref{fig:knowledge-acquisition-prediction-mechanisms}把三条路径画在一起，分别对应：读取模型参数中的已学关系、读取近邻的参考答案、利用图上连接传播标签。

需要先说明分类标准，以免读者误以为这是“神经网络 vs 传统方法”之别。这里区分机制的唯一依据是\emph{答案主要靠什么算出来}：用参数映射、查近邻标签、还是在图上传播。三者完全可以串在同一个方案里（比如先用神经网络学表示、再用这个表示去检索近邻），所以它们不是互斥的模型家族。

\begin{figure*}[htbp]
  \centering
  % 图4-3：三种预测机制对照，按169 mm全版心宽度设计。
\begin{tikzpicture}[
  x=1mm,
  y=1mm,
  font=\figurefont\footnotesize,
  text=BookInk,
  unknown/.style={circle,draw=FigureInput,fill=white,line width=0.9pt,minimum size=8mm,inner sep=0pt},
  known/.style={rectangle,draw=FigureInput,fill=FigureInput!10,line width=0.9pt,minimum size=8mm,inner sep=0pt},
  compute/.style={book node,draw=FigureModel,fill=FigureModel!8,line width=0.8pt,text width=16mm,minimum height=11mm,inner sep=1mm},
  answer/.style={book node,draw=FigureOutput,fill=FigureOutput!6,line width=0.8pt,text width=13mm,minimum height=11mm,inner sep=1mm},
  relation/.style={draw=BookMuted,line width=0.7pt},
  error relation/.style={draw=BookAmber,line width=0.8pt,dashed,rounded corners=1mm},
  calculation/.style={knowledge arrow,draw=FigureModel,line width=0.9pt},
  annotation/.style={font=\figurefont\scriptsize,text=BookMuted,align=center},
  edge label/.style={font=\figurefont\scriptsize,text=BookMuted,fill=white,inner sep=.5pt}
]
  \draw[BookRule,line width=0.7pt] (55,2) -- (55,-52);
  \draw[BookRule,line width=0.7pt] (110,2) -- (110,-52);

  \node[anchor=north west,font=\figurefont\bfseries] at (1,1) {A\quad 已学映射};
  \node[unknown] (direct-query) at (9,-25) {S};
  \node[annotation] at (9,-34) {查询S042};
  \node[compute] (direct-model) at (29,-25) {预测模型\\$f_{\bm{\theta}}$};
  \node[answer] (direct-answer) at (48,-25) {候选\\山茶};
  \draw[calculation] (direct-query.east) -- (direct-model.west);
  \draw[calculation] (direct-model.east) -- (direct-answer.west);
  \node[annotation] at (29,-44) {预测时主要读取模型参数};

  \node[anchor=north west,font=\figurefont\bfseries] at (57,1) {B\quad 近邻参考};
  \node[unknown] (neighbor-query) at (64,-27) {S};
  \node[annotation] at (64,-37) {查询S042};
  \node[known] (neighbor-a1) at (79,-13) {A};
  \node[known] (neighbor-a2) at (79,-27) {A};
  \node[known] (neighbor-b) at (79,-41) {B};
  \node[compute,text width=13mm] (vote) at (98,-27) {加权\\投票};
  \node[answer,text width=13mm] (neighbor-answer) at (98,-46) {候选\\山茶};
  \draw[relation] (neighbor-query) -- (neighbor-a1);
  \draw[relation] (neighbor-query) -- (neighbor-a2);
  \draw[relation] (neighbor-query) -- (neighbor-b);
  \draw[calculation] (neighbor-a1.east) -- (vote.west);
  \draw[calculation] (neighbor-a2.east) -- (vote.west);
  \draw[calculation] (neighbor-b.east) -- (vote.west);
  \draw[calculation] (vote.south) -- (neighbor-answer.north);
  \node[annotation] at (79,-51) {$A=$山茶，$B=$茶梅};

  \node[anchor=north west,font=\figurefont\bfseries] at (112,1) {C\quad 图结构};
  \node[known] (graph-a) at (119,-11) {$A:1$};
  \node[known] (graph-b) at (119,-42) {$B:0$};
  \node[unknown] (graph-u) at (137,-27) {$u$};
  \node[unknown] (graph-v) at (151,-11) {$v$};
  \node[answer,text width=14mm] (graph-answer) at (160,-33) {$u$的A分数\\$\frac23\to\frac{11}{15}$\\错边后$\frac{17}{27}$};
  \draw[relation] (graph-a) -- node[edge label,above] {$2$} (graph-u);
  \draw[relation] (graph-b) -- node[edge label,below] {$1$} (graph-u);
  \draw[relation] (graph-u) -- node[edge label,above] {$1$} (graph-v);
  \draw[relation] (graph-a) -- node[edge label,above] {$3$} (graph-v);
  \draw[error relation] (graph-b.east) -- (graph-v.south |- graph-b.east) -- (graph-v.south);
  \node[annotation,text=BookAmber] at (137,-45) {错误边$+3$};
  \draw[calculation] (graph-u.east) -- (graph-answer.west);
  \node[annotation] at (136,-51) {方形为标签锚点；边权可按例~\ref{ex:knowledge-acquisition-graph-propagation}复算};
\end{tikzpicture}

  \caption{三种预测机制怎样使用已有模型参数、参考标签与图连接。圆形是未知答案的查询或节点，方形是带答案参考或标签锚点；无箭头边表示相似关系，箭头表示计算路径。面板A、B示意候选产生，面板C对照直接邻居、多步传播和错误连边的得分。类别、二维位置、距离与边权均为教学构造，未应用Q1接收约定，不改写S042的弃权判断。}
  \label{fig:knowledge-acquisition-prediction-mechanisms}
\end{figure*}

正因为答案来源不同，三条路径\emph{出错时需要保留和追查的依据}也不同，表~\ref{tab:knowledge-acquisition-prediction-mechanisms}把每种机制依赖的东西（模型参数/参考样本/图结构）、核心计算、处理新样本的方式和主要风险一一对应。例如：直接模型的风险在训练覆盖和接口变化，近邻的风险在距离是否真有语义、参考标签质量和检索遗漏，图方法的风险则在错误连边、错误锚点和孤立无援的连通分量。

在真实的组合方案里，往往是直接模型先把照片编码成表示、近邻方法再用这个表示找参考、图推断最后协调一批节点。采用这类组合时，中间每一步——参数映射、参考选择、图连接——都应当\emph{可单独检查}。最后强调一个会反复出现的核查原则：两条路径即使给出了\emph{相同的最终类别}，也不能说明它们的中间依赖相同或同样可靠；只对最终答案、不看过程，就无法在出错时定位是哪一环的问题。

\begin{table*}[htbp]
  \centering
  \small
  \setlength{\RaggedRightParindent}{0pt}
  \renewcommand{\arraystretch}{1.15}
  \begin{tabularx}{\linewidth}{@{}>{\RaggedRight\arraybackslash}p{20mm}*{4}{>{\RaggedRight\arraybackslash}X}@{}}
    \toprule
    预测机制 & 读取的模型参数、参考标签或图结构 & 核心计算 & 新样本处理 & 主要依赖与风险 \\
    \midrule
    已学映射 & 模型参数、接口与当前预测输入 & 前向计算或解码 & 按接口推断 & 训练覆盖、接口与分布变化 \\
    近邻参考 & 参考答案、表示与索引 & 距离排序与加权 & 检索并聚合 & 距离含义、标签质量与检索遗漏 \\
    图结构 & 边权、标签锚点与节点 & 联立求解或传播 & 接入图并更新解 & 错边、错标签与无锚点分量 \\
    \bottomrule
  \end{tabularx}
  \caption{三种预测机制的使用条件。各行按当前答案的主要计算路径，列出读取的模型参数、参考标签或图结构、新样本的处理及依赖风险；学习表示、检索与图推断可在同一方案中组合，不是模型结构的永久分类。}
  \label{tab:knowledge-acquisition-prediction-mechanisms}
\end{table*}

\subsection{基于已学映射的直接预测}
\label{sec:knowledge-acquisition-direct-prediction}

先看最直接的路径。一个已经训练好的花卉分类器，在每次调用时\emph{不需要}再去翻找带标签照片——它把过去大量样本中学到的规律“压缩并固定”在模型参数里，作答时只做一次前向计算，把当前输入通过这些参数\emph{映射}成每个类别的得分。这种答案完全由当前输入和固定参数算出、不查阅外部样本的路径，本章称为\term{直接预测}（direct prediction）。

把记号列清楚。$\mathsf{u}_i$表示样本$i$这一次的完整输入，包括接口允许传入的任务要求、上下文和辅助内容（而不仅是照片本身）；$\bm{\theta}$是这次调用版本的模型参数。类别表规定候选集合为$Y=\{1,\ldots,K\}$，模型对每个类别$c$算出一个得分$s_{\bm{\theta},c}(\mathsf{u}_i)$，约定分数越高表示模型越支持该类别。要从这$K$个得分里得到一个具体的类别候选，最常见的做法就是选取得分最大的类别：
\begin{equation}
  \widehat{y}_i
  \in
  \argmax_{c\in Y}s_{\bm{\theta},c}(\mathsf{u}_i)
  \label{eq:knowledge-acquisition-direct-classification}
\end{equation}
取最大值这个动作要正确理解其能力边界：它只完成了“在\emph{同一张给定类别表}内部给类别排序、挑出最靠前的一个”，仅此而已。若最高分出现并列，不能随便挑一个，而要按事先预定的规则处理，或保留多个候选。若希望模型在没把握时弃权，光有$\argmax$不够，还要另行规定接收条件，比如最高分要达到某门槛、或最高分与次高分的差距要足够大（差距过小说明在两个物种间犹豫）。

还有两个常见误解必须澄清。其一，模型输出的得分\emph{不是天然就是概率}：只有当模型的定义方式和相应的校准评价都支持时，才能把某个分数解释成“正确的可能性”；未经校准的分数只是内部排序用的相对量。其二，即使分数已经校准良好，用$\argmax$选出最高分也\emph{不构成}“这张照片答案正确”的证明——概率再高也可能落在那小部分错误里。

直接预测这套“输入—得分—取最大类”的映射，还可以用在另一种\emph{监督单位}下，理解它需要先回顾第~\ref{sec:rule-based-knowledge-acquisition-relation-matching}节的远程监督。远程监督只给得到“组”这一级的标签：对同一个实体对的一组句子，我们只知道“其中至少有一句表达了某关系”，却不知道\emph{具体是哪一句}。沿用那里的两句话：句乙“林舟在论坛上批评了海岬科技的发布流程”只是同时提及两者，句甲“海岬科技任命林舟为首席架构师”才真正表达任职。

这里要严格区分\emph{教学讲解}与\emph{训练时实际拥有的信息}。我们作为讲解者当然知道甲真乙假，但在真实的远程监督训练里，只有“组内至少一句表达任职”这一个组级标签，\emph{没有}甲、乙各自的逐句答案，模型并不知道哪句是证据。于是产生一个技术问题：模型是逐句给出预测的，监督却只在组一级，怎样把逐句预测\emph{聚合成}一个能与组级标签比较的量？这正是接下来要解决的。

先把组形式化。令$B$为某个非空提及组内句子的索引集合，$Z_j\in\{0,1\}$表示第$j$句是否表达目标关系（1为表达）。模型对每一句，根据由该句文本和实体标识组成的输入$\mathsf{u}_j$，算出一个“这句表达关系”的概率
\[
  p_j=p_{\bm{\theta}}(Z_j=1\mid\mathsf{u}_j).
\]
我们想要的是“组内\emph{至少有一句}表达关系”的概率。直接处理“至少一句”不方便，概率论里常用的技巧是转而去算它的\emph{补事件}“\emph{一句也不}表达”，再用1减去它。

\term{噪声或}（noisy-or）就是完成这一聚合的经典模型，它基于两个明确假设：给定整组输入后，各句的$Z_j$相互\emph{条件独立}；每句是否表达关系只取决于它自己的$\mathsf{u}_j$。在这两个假设下，“所有句子都不表达”这个联合事件的概率，等于各句“不表达”概率$(1-p_j)$的乘积。于是“至少一句表达”的概率就是1减去这个乘积。得到组级概率后，再用组标签$y_B\in\{0,1\}$与它比较，用二元交叉熵损失去惩罚“模型给真实事件赋了过低概率”；采用自然对数时，损失写成
\begin{equation}
  \begin{aligned}
    \widehat p_B
    &=1-\prod_{j\in B}(1-p_j),\\
    \ell_B
    &=-y_B\log\widehat p_B
      -(1-y_B)\log(1-\widehat p_B).
  \end{aligned}
  \label{eq:knowledge-acquisition-model-bag-prediction}
\end{equation}
例如构造一个正组，按句乙、句甲的顺序设$p_1=0.1$、$p_2=0.8$、$y_B=1$。两句都不表达的概率为$0.9\times0.2$，因此
\[
  \widehat p_B=1-(1-0.1)(1-0.8)=0.82,
  \qquad
  \ell_B=-\log0.82\approx0.19845.
\]
这个计算直观地说明了聚合在做什么：$(1-0.1)(1-0.8)=0.18$是“两句都不表达”的概率（句乙虽有0.9的可能不表达，但句甲只有0.2的可能不表达，两者同时成立的概率相乘后很小），因此至少一句表达的概率高达$0.82$；$-\\log0.82\\approx0.19845$是当组标签为正时，这个尚不算高的组概率所受的损失。

训练时，梯度会\emph{穿过这个聚合式}回传给逐句模型，但它优化的目标始终是整组的“至少一句”事件，\emph{从不要求组内每一句都为正}——这正是它能容忍句乙这类噪声提及的原因。预测时可以返回组分数$0.82$，并把句内得分更高的句甲作为“可能的证据句”列出供人参考；但务必记住$0.82$是\emph{组}的概率，不是句甲、也不是句乙各自表达关系的概率。而且，多组不同的$(p_1,p_2)$搭配可以凑出同一个$0.82$（这是一个方程、两个未知数），所以仅凭组标签，数学上根本无法唯一确定究竟哪一句是真正的支撑句。

noisy-or的结论完全建立在“各句条件独立”这个假设上，使用时要警惕假设失效。如果组内的句子其实是同一篇报道的重复转载、或高度相似，它们共享同一处证据、并不独立，却被当成多个独立证据把$(1-p_j)$反复相乘，就可能\emph{高估}整组的支持强度——一条信息被算了好几遍。

另一个风险在监督的源头：知识库里存在某条关系，并不保证当前这个组里真有句子在表达它（回顾第3章远程监督的噪声）。因此当$y_B=1$只是一个远程监督候选时，应先决定如何对待它——直接接收、降权、还是送复核；上面的聚合公式只演示了“\emph{一旦接受}这条监督，具体怎么训练”，它本身并没有、也不能把那条远程标签核验成事实。

前面的输出要么是固定类别表中的一个类别，要么是组级的0/1事件；一旦任务改变，这套参数映射输出的\emph{含义}也随之改变。回归任务输出的是连续数值；结构化预测（如分词、句法）输出的标签还要满足序列、树或集合上的组合约束。而当我们让模型\emph{生成}物种名称和一段说明文字时，候选不再是$Y$中的一个索引，而是由词元表中一个个词元（token）按顺序拼成的序列。

生成一个长度为$m$的序列，概率应如何计算？核心工具是概率论的\emph{链式法则}：把“整个序列出现”的概率，分解成“在给定输入、以及此前已经生成的前缀条件下，下一个词元出现”的一系列逐步条件概率之积。设输入仍为$\mathsf{u}_i$，第$r$个词元为$y_r$，则完整序列$y_{1:m}$的条件概率为
\begin{equation}
  p_{\bm{\theta}}(y_{1:m}\mid\mathsf{u}_i)
  =
  \prod_{r=1}^{m}
  p_{\bm{\theta}}(y_r\mid y_{1:r-1},\mathsf{u}_i).
  \label{eq:knowledge-acquisition-direct-generation}
\end{equation}
这个乘积清楚地显示了生成过程的\emph{自回归}特性：每一步的概率都依赖此前已生成的所有词元$y_{1:r-1}$，因此生成是一个逐词推进、前面影响后面的过程。但要注意，链式法则只给出了“\emph{如何给任意一个完整序列算概率}”，它本身\emph{还没有替我们选定}最终要输出哪一段文本——可能的序列数量极其巨大，不可能逐一比较。

真正\emph{从这些条件概率里选出或采样出一个序列}的过程，称为解码（decoding）。怎么解码对结果影响很大：是每步都选概率最大的词（贪心）、还是在若干候选序列中综合择优（集束搜索）、还是按概率随机采样；保留多少候选、何时停止、允许生成多长，都会改变最终文本。得到文本之后还有任务层面的收尾：生成的物种名要映射回类别表、处理各种别名写法。最后必须守住一条底线：一段说明文字\emph{通顺、像那么回事}，只证明模型完成了语言生成；它陈述的证据是否真实、是否真的支持所给物种，仍需另行核查——流畅不等于有据。

无论分类还是生成，产出原始输出后，批量流程还要做一道“解析并把结果还原到对象”的工序，这一步有自己特有的错误。用重叠窗口切分长文本时，同一个实体在不同窗口里可能得到不一致的答案，需要消解；损坏的图像可能让模型返回空结果；更隐蔽的是，一个字段可能\emph{语法上完全合法、语义上却填错了}（JSON能解析，但物种名是编的）。

为分别定位这三类问题，应同时保存\emph{模型的原始输出}和\emph{解析后}得到的候选，这样才能把“模型本身预测错”“解析程序读错/对不上字段”“结果写回了错误对象”三件事分开检查。此外，让模型一次采样多个候选虽可能提升命中机会，却也会扩大待处理集合、增加后续融合与复核的工作量；这些多份候选最终如何综合取舍，属于第~\ref{chap:knowledge-fusion}章的内容。

前面都是“输入对象固定、让模型给答案”，还可以更进一步，连“要做什么任务、给什么输入”都让模型参与生成。\term{Self-Instruct}就是这种材料自动构造方式：从少量人工写的种子任务出发，让模型自己生成新的指令、配套的输入、以及示范答案，再过滤掉无效或高度重复的材料，用保留下来的实例继续微调自身\scite{wang2023selfinstruct}。

这一步显著扩大了需要检查的范围，因为模型不再只改变标签，而是\emph{同时改变了训练任务和输入的分布}。原来针对“固定样本、固定标签”设计的质检项已不够用，检查必须覆盖四件事：生成的任务本身是否成立（是不是一个有意义、可学习的问题）、输入是否真的符合该指令、答案是否可用、以及材料之间是否大量重复。单纯生成更多实例只增加了候选\emph{数量}；它们究竟有没有扩大任务覆盖、有没有真正带来训练效用，需要另外的证据，不能用数量代替质量。

\subsection{基于近邻参考的预测}
\label{sec:knowledge-acquisition-neighbor-prediction}

设想直接预测这条路暂时走不通（分类器在目标物种上还不适用），但我们手里保存着一批\emph{带物种标签}的照片。一个朴素而有效的想法是：要判断一张新照片是什么，就去看“和它长得最像的那些已知照片”分别是什么。这种答案\emph{在调用现场从参考样本取得}、而不是从事先固定的模型参数读取的路径，称为\term{近邻参考预测}（neighbor-reference prediction）：先选出与查询最相近的若干带答案样本，再把它们的答案聚合起来。

把它写成计算。设参考集为$\mathcal{R}=\{(\bm{x}_j,y_j)\}_{j=1}^{n}$，其中$y_j$是参考照片$j$的类别标签、属于任务类别表$Y$；用一个映射$\phi$把照片变到一个可以比较的特征空间，$d$是该空间上的距离。对新照片$\bm{x}$，按距离从小到大选出最近的$k$个参考，其索引集合记为$N_k(\bm{x})$（$1\leq k\leq n$，这就是“$k$近邻”里的$k$）。为了把邻居们的标签汇总成可比较的类别支持，可以先给每个邻居一个非负权重并归一化（越近通常权重越大），再累加每个类别得到的总支持、选支持最大的类别：
\begin{equation}
  \begin{aligned}
    \widehat{y}(\bm{x})
    &\in
    \argmax_{c\in Y}
    \sum_{j\in N_k(\bm{x})}
    \alpha_j(\bm{x})\mathbf{1}\{y_j=c\}, \\
    \alpha_j(\bm{x})
    &=
    \frac{w\!\left(d(\phi(\bm{x}),\phi(\bm{x}_j))\right)}
    {\sum_{r\in N_k(\bm{x})}w\!\left(d(\phi(\bm{x}),\phi(\bm{x}_r))\right)}.
  \end{aligned}
  \label{eq:knowledge-acquisition-neighbor-prediction}
\end{equation}
逐项看懂这个式子。函数$w$把“距离”换算成非负权重，通常让距离越小、权重越大（越像的参考越有发言权），也可以简单取等权（$k$个邻居一人一票）；分母把所选邻居的权重归一化到总和为1，因此前提是至少有一个邻居具有正权重，否则无法归一化。归一化后，某个类别的得分恰好等于“投给该类的邻居权重之和”。

读这个得分时要保持克制：它是一个\emph{加权投票比例}，\emph{不是}自动校准过的正确概率——邻居们可能一起错。同一套框架稍加改动就能用于回归：把式中那个“标签是否等于类别$c$”的指示函数$\mathbf{1}\{\cdot\}$换成参考数值$y_j$，输出就变成邻居数值的加权平均。还有若干边界情况必须在实现前\emph{预先规定}如何处理：距离并列时怎么选邻居、类别票数相同时取谁、一个有效邻居都没有时怎么办，以及查询照片本身就和某张参考完全重复时如何处理。最后一种尤其要警惕：如果评价对象的照片（带着已知标签）也混在参考集里，最近邻就会直接把它\emph{自己的标签}查回来，得到虚高的满分；这不是泛化，必须把这类重复参考排除，或明确采用一种不同的评估口径。

这种“依据局部参考实例下判断”是经典近邻分类的基本思想\scite{coverhart1967}。这里我们只保留理解当前机制所必需的检索与聚合两步；完整的算法细节（$k$如何选择、距离度量、收敛性质等）与统计保证，可参考\BookChapterRef{chap:classification}{卷4《模型：经典模型》的“分类”章}。

值得点出近邻法与直接预测在\emph{维护方式}上的一个重要差别。在近邻法里，答案是调用时实时读取参考标签得到的，因此一旦某个参考标签被更正或撤回，它\emph{附近}查询的结果可能立刻改变，\emph{不需要}先重新训练表示模型——知识更新直接改数据即可。而直接预测把知识固化在参数里，要改变行为通常得重新训练。这既是近邻法易于维护的优点，也意味着一条错误参考会长期、直接地影响周围一片查询。

近邻法能不能给出\emph{有物种意义}的依据，完全取决于一个前提：特征空间里的“距离近”是否真的对应“任务相关的相似”。这个前提可能不成立。如果学到的表示$\phi$主要编码的是背景（草地、蓝天）、拍摄设备或光线，而不是花的鉴别特征，那么“近邻”可能只是“在同一地点、用同款手机拍”的照片——公园里的一张山茶可能因此检索到同地点的一张月季，形态上却毫无关系。只有当表示保留了足够的物种鉴别特征，形态近邻才有资格支持物种判断。

此外还有两类偏差。其一，多张邻居标签一致，可能不是因为它们都正确，而是因为它们\emph{来自同一个错误来源}（同一批误标数据），错误被反复引用。其二，类别失衡会让多数类天然挤占邻域：若参考集里月季占绝大多数，那么无论查询是什么，最近的$k$个邻居多半都是月季，少数类几乎没有发言权。还要注意参考稀疏的情形：当某区域参考照片很少时，所谓“最近邻”只是“矮子里拔将军”，“这批里最近”绝不等于“足够相似”；此时应按事先定好的距离与覆盖门槛，降低对该投票的接收程度乃至直接弃权。

近邻法把训练成本转移到了\emph{推断与存储}：保留的参考越多，判断时要搜索的库越大，存储、索引构建和每次查询的成本都随之上升。为应对大规模库，工程上常用近似近邻索引（approximate nearest-neighbor index），它\emph{用可能漏掉真正最近邻为代价}换取更快的搜索；这个代价必须被意识到——漏掉的若是关键参考，投票结果就会改变。此外，新增或撤回参考照片时，索引和缓存可能需要重建或更新，否则会用到过期数据。

为保证可追溯与可复现，每条近邻产生的候选样本标注都应保存：实际采用了哪些参考的标识、它们各自的距离与权重、以及所用索引的版本。这样日后才能\emph{重放}式~\eqref{eq:knowledge-acquisition-neighbor-prediction}的计算，并把“参考答案本身被更正导致结果变化”与“检索实现/索引版本变化导致找到不同邻居”这两种原因区分开。

\subsection{基于图结构的预测}
\label{sec:knowledge-acquisition-graph-prediction}

若一批未标注照片互相相似，只汇总每张查询附近的已知标签，会忽略未知照片之间的连接。我们可以将已知与未知样本共同放入图，让连接约束多个节点的答案，这就是\term{图结构预测}（graph-based prediction）。这里采用的假设是：任务相关的强连接应对应相近的类别得分，并由已知标签确定类别含义。\term{标签传播}（label propagation）沿这些连接传递类别得分，不把相似边解释成因果关系。设无向图$\mathcal{G}=(V,E)$有非负对称边权$w_{i,j}=w_{j,i}$，无边处权重为0；$V$分为已标节点$L$与未知节点$U$。已标节点称为标签锚点，因为它们提供类别参照。用$\bm{F}_i\in\mathbb{R}^{K}$保存节点$i$的$K$类得分，所有行组成$\bm{F}\in\mathbb{R}^{|V|\times K}$，已知标签$\bm{Y}_i$为独热向量，即对应类别为1、其余为0。我们需要一组未知得分，既保持$\bm{F}_L=\bm{Y}_L$，又让强连接两端差异小，因此在满足边界的得分矩阵中最小化加权平方差：
\begin{equation}
  \min_{\bm{F}:\,\bm{F}_L=\bm{Y}_L}
  \frac{1}{2}\sum_{i,j\in V}w_{i,j}\lVert\bm{F}_i-\bm{F}_j\rVert_2^2
  =
  \min_{\bm{F}:\,\bm{F}_L=\bm{Y}_L}
  \operatorname{tr}\!\left(\bm{F}^{\top}\bm{L}_{\mathcal{G}}\bm{F}\right),
  \label{eq:knowledge-acquisition-graph-energy}
\end{equation}
其中$\bm{W}$是边权矩阵，$\bm{D}$是对角矩阵，$\bm{D}_{i,i}=\sum_jw_{i,j}$；图Laplacian矩阵$\bm{L}_{\mathcal{G}}=\bm{D}-\bm{W}$把同一平方差写成矩阵形式，$\operatorname{tr}$表示对角元素之和。对称权重使$(i,j)$和$(j,i)$重复计数，$1/2$抵消重复；平方差展开后的两项节点平方范数合并为度数项，内积项由$\bm{W}$给出，便得到右侧。最小化的是图上加权不一致，不是未知样本的真实分类错误；其作用依赖连边符合刚才的任务相似假设。

要从目标得到实际未知得分，按$L,U$排列行列并固定$\bm{F}_L=\bm{Y}_L$。未知部分的梯度为$2(\bm{L}_{U,U}\bm{F}_U+\bm{L}_{U,L}\bm{Y}_L)$，令其为0，便得到线性方程；可逆时再写成逆矩阵形式：
\begin{equation}
  \bm{L}_{U,U}\bm{F}_U
  =
  -\bm{L}_{U,L}\bm{Y}_L,
  \qquad
  \bm{F}_U
  =
  -\bm{L}_{U,U}^{-1}\bm{L}_{U,L}\bm{Y}_L.
  \label{eq:knowledge-acquisition-graph-solution}
\end{equation}
这里有一个自然的疑问：矩阵$\bm{L}_{U,U}$一定可逆、解一定唯一吗？答案是\emph{不一定}，取决于图的连通结构。对非负、无向的边权，有一个清晰的充分条件：如果\emph{每一个}包含未知节点的、沿正权重连通的分量，都能沿边\emph{到达至少一个标签锚点}，那么固定锚点后，未知部分就有唯一解。

可以这样理解其依据：假设在锚点不变的前提下，未知得分发生了某个扰动却\emph{不增加任何能量}（零能量扰动）。由于任何非零的相邻差异都会带来正的平方差，零能量要求扰动在每个连通分量内\emph{处处相等}（是一个常数）；而这个分量连着锚点、锚点处扰动必须为0，所以这个常数只能是0。唯一的零扰动就是“不变”，说明解唯一，对应二次型正定。反过来，若某个未知分量里\emph{一个锚点都没有}，那么把所有未知节点都加上同一个常数，相邻差异仍是0、能量仍为0，却不违反任何锚点约束——于是绝对类别得分无法确定（只能定到相差一个常数），对应的矩阵就是奇异的、不可逆的。Gaussian fields与调和函数方法系统地给出了这种构造、它的矩阵解以及一个随机游走视角的解释\scite{zhu2003}。

同一个线性方程还能导出一种\emph{不需要显式求逆}的局部算法。把未知节点$i$的驻点条件移项整理，会得到一个很直观的结论：它的得分应当等于\emph{所有邻居得分的加权平均}，这称为调和条件。于是可以不直接解方程组，而是\emph{逐轮迭代}更新
\[
  \bm{F}_i
  \leftarrow
  \frac{\sum_j w_{i,j}\bm{F}_j}{\sum_j w_{i,j}},
\]
其中分母要求节点$i$至少有一个正权重邻居（否则无法做平均）。迭代规则是：每一轮用上式更新所有未知节点，但\emph{每轮结束都把已知锚点强行重置回}$\bm{Y}_i$，防止锚点值在传播中被稀释；在前面的“每分量可达锚点”条件下，这种边界迭代会收敛到与线性求解相同的结果。

无论用哪种方式得到得分矩阵，最后都还要从每行的$K$个类别得分里选出最高项，并应用事先定好的平票处理与弃权规则——\emph{得分矩阵本身只是连续的软得分，还不是已核验的类别标签}。相比近邻法“每张查询只看自己附近的已知标签”，图传播多了一个能力：未知节点也能作为\emph{中继}，把类别得分继续传给更远的节点（信息可以跨过没有直接锚点的中间节点）。但能力扩大也意味着风险扩大：一个错误的锚点、或一条连接了两个不同类别的错误边，其影响可能沿着图传播得\emph{更远}、污染一整片。还要清醒：迭代在数值上\emph{收敛}，只说明找到了图能量目标的那个最优解，完全不说明底层“强边即同类”的图假设是正确的。

前面的得分还可以换一个非常直观的角度理解——\emph{在图上随机行走}。把边权归一化后，设想从一个未知照片节点出发，每一步都按“各条边占该节点总边权的比例”随机选择一个邻居走过去，一旦\emph{第一次走到某个标签锚点}就停止。在有限图且每个未知分量都能到达锚点的条件下，可以证明：某节点在类别$c$上的得分，恰好等于“从该节点出发的随机行走，\emph{最先}抵达的锚点是类别$c$”这一事件的概率。

这个随机游走解释让“标签传播”变得可触摸：强连接多的路径被走中的概率大，得分本质上是在问“顺着相似度走，最容易先撞上哪一类的已知照片”。但必须守住解释的边界：这是一个\emph{由我们构造的图和行走规则决定的到达概率}，是模型内部的量，\emph{不是}通过独立测试得到的“分类正确率”。如果背景相似造成了错误连边，随机行走就会被引向错误类别的锚点，最先到达的终点随之改变，得分也就被带偏了。

下面用一个只算类别A一维得分的小图，亲手算出三种情形的差别：只看查询的直接邻居、在完整图上多步传播、以及人为加入一条错误连边后结果如何恶化。建议读者跟着把每一步的加权平均算一遍，这比抽象推导更能建立直觉。

\begin{example}[直接近邻、多步传播与错误连边]
  \label{ex:knowledge-acquisition-graph-propagation}
  用一个构造图检查未知节点的连接改变了什么，只计算类别A的一维得分。锚点A、B分别取$f_A=1$、$f_B=0$；查询$u$以权重2、1连接A、B，以权重1连接未知$v$；$v$再以权重3连接A。若只读取$u$直接相连的已标答案，忽略$v$并归一化，得到
  \[
    f_u^{\mathrm{direct}}=\frac{2f_A+f_B}{2+1}=\frac{2}{3}.
  \]
  完整图则保留$u$与$v$的连接。两节点的总边权都为4，调和条件给出
  \[
    4f_u=2+f_v,\qquad 4f_v=3+f_u,
  \]
  联立得$f_u=11/15$、$f_v=14/15$，说明A锚点经$v$影响了$u$。再人为加入错误边：让$v$以权重3连接B，第二式变为$7f_v=3+f_u$，解变为$f_u=17/27$、$f_v=14/27$。假设教学规则仅当A得分至少$0.7$时接收A候选，原图中的$11/15$达到阈值，错边后的$17/27$没有达到。阈值在此展示得分变化如何影响接收，不是实际质量标准。图~\ref{fig:knowledge-acquisition-prediction-mechanisms}的面板C使用同一组节点、权重和结果。
\end{example}

回到花卉任务，图的节点是一张张照片，连边的\emph{线索}可以来自形态表示相似、拍摄地点相同、或“同属一株”的元信息。但每一种线索都要单独检查它\emph{是否真的支持物种相似}：地点相同的照片未必同种，同株关系才强地意味着同种。前面反复强调的风险在这里落地——若两张只是背景相似、物种却不同的照片被赋了强连接，错误就会沿这条边向周围扩大；而每当加入一张新照片改变了邻接关系，旧节点的得分也可能随之改变，结果不是一成不变的。

工程上还要计入成本：如何构图、如何稀疏地存储大量边、如何高效求解，都有实际开销；好在前面的线性方程通常可以用迭代或稀疏线性代数求解，而不必真的去计算大规模逆矩阵。最后厘清两件“用了图但含义不同”的事。本节即使让表示和边权也可学习，\emph{类别含义仍然来自已有的物种标签锚点}——是标签把得分解释成了物种。而第~\ref{chap:data-derived-knowledge-acquisition}章的图分组，是在\emph{还没有任何物种标签}时先发现结构上成簇的群组，那些簇尚未被赋予物种语义。同样“在图上计算”，监督信号的性质完全不同，不能混为一谈。

\section{获取过程的迭代改进}
\label{sec:knowledge-acquisition-iterative-improvement}

跑完一轮、得到一大批花卉候选之后，复核往往会暴露出性质各异的问题。有的是\emph{类别边界学错}了（山茶、茶梅系统性混淆）；有的是\emph{观测本身缺失}（关键叶片没拍到，再强的模型也无能为力）；有的是\emph{外部证据关联错}（名录过期或对错了对象）；还有的问题不在单个答案，而在\emph{资源安排}——预算被大量花在重复处理那些模型早就会的简单照片上，难照片却没被照顾到。

这四类问题对应\emph{不同的修改对象}：边界学错要改模型，观测/证据缺失要改预测输入，资源错配则要改“下一轮处理哪些样本”的获取安排。用错了修改对象，比如对着“没拍到叶片”的问题去反复重训模型，只会徒劳。

在动手改之前，必须先\emph{保存并冻结}上一轮所依赖的全部版本：用了哪些样本、哪个模型版本、什么输入、近邻索引或图的哪个版本、产生了哪些预测输出。冻结之后再按成因选择更新对象。改完也不能想当然地认为“重训一次万事大吉”：要沿\emph{版本依赖关系}判断哪些旧标注会因这次改动而失效、需要重算——改了模型，所有直接预测可能都变；只改一条名录，则只有引用它的少数标注受影响。

迭代过程中用来\emph{选择修改方式}的样本（试出合适阈值、挑出输入版本、决定教师如何更新、判断继续还是停止/回退），扮演的是\emph{开发验证和过程审计}的角色。必须守住第2章以来一贯的纪律：一旦某批样本参与了这些决策，它就被“用脏了”，\emph{不能}再充当方案冻结后的最终独立验收——同一批样本不能既当教练又当裁判。图~\ref{fig:knowledge-acquisition-iterative-update}的上半部分专门把“版本决策”与“留出验收”隔成两套数据，下半部分则演示一个教师—学生更新的内部流程。

还要切断一个最诱人的错误因果：\emph{训练目标下降}（损失变小）只说明模型把\emph{当前这轮喂给它的训练信号}拟合得更好了。如果训练信号本身混着错误伪标签，损失下降反而可能意味着在错误的道路上越走越远。因此“损失下降”与“候选质量改善”之间不能直接画等号，后者必须由独立的、面向参照答案的检查来确认。

\subsection{模型的更新}
\label{sec:knowledge-acquisition-model-update}

\subsubsection{基于伪标签与教师目标的更新}
\label{sec:ka-pseudo-label-update}

现在设想一种典型处境：只有少量昂贵的人工物种标签，但还有海量未标注照片，而模型已经能对这些照片给出候选。能不能让模型\emph{用自己的预测}来扩充训练数据、从而自我改进？这就是\term{自训练}（self-training）的基本想法：先用当前模型对未标注照片作预测，再\emph{筛选}出其中相对可信的一部分，把它们的预测\emph{暂时当作目标}，加入下一轮训练去更新模型。

这里引入一个必须正名的概念：这些“模型给自己、又拿来训练自己”的临时目标称为\term{伪标签}（pseudo-label）。“伪”字是关键提醒——它们\emph{源自模型自己的猜测}，即使被选进训练集，也\emph{不会因此变成真实标签}，错的依然是错的。自训练的全部风险都来自这一点。

把选择过程形式化。设人工已标注样本的索引集合为$L$，第$t$轮未标注候选的索引为$U^{(t)}$，两者始终分开保存；模型对候选$j$给出预测类别$\widehat y_j^{(t)}$和一个接收分数$q_j^{(t)}$，约定分数越高越符合“可信、可接收”的偏好。要明确本轮\emph{哪些}候选有资格进入训练，可以用一个在\emph{开发阶段}选定的阈值$\tau^{(t)}$来圈定集合：
\begin{equation}
  P^{(t)}
  =
  \left\{
    \left(\bm{x}_j,\widetilde y_j^{(t)},q_j^{(t)}\right)
    :
    j\in U^{(t)},\
    q_j^{(t)}\geq\tau^{(t)},\
    \widetilde y_j^{(t)}=\widehat y_j^{(t)}
  \right\}.
  \label{eq:knowledge-acquisition-pseudo-label-set}
\end{equation}
集合里的$\widetilde y_j^{(t)}$就取模型当前的预测类别，并连同原样本和接收分数一起保存，供日后追溯。读这条筛选规则要抓住它的能力边界：阈值$\tau^{(t)}$只回答“\emph{哪些}候选被\emph{采用}进训练”，它\emph{不回答、也无法回答}被采用的每一条是否正确——选入是一个置信度决策，不是一次正确性核验。

具体地说，如果分数$q_j^{(t)}$没有经过相应的概率校准，它只能在“当前模型自己的尺度”上比较高低，不能解释成真实正确率；即便整体上分数已校准，也还要检查被接收候选在\emph{各类别、各拍摄条件}上的覆盖是否均衡（自训练很容易只接收模型本就擅长的常见类）。总之，\textbf{高分筛选绝不等于独立核验}：伪标签在进入训练时仍是待验证的候选，这一身份不能因分数高而改变。

确定了“哪些伪标签进训练”，还需要一个目标函数，说明它们如何与\emph{真正可靠}的人工标签\emph{共同}影响参数。一个合理的做法是把两部分损失分开计算：人工标签一部分、伪标签一部分（后者还可逐条加权），再用一个系数控制伪标签相对人工标签的影响力——毕竟人工标签通常更可信，不应让海量伪标签完全盖过它。

式中两处都用了\emph{按各自样本数平均}，而不是把所有损失简单求和。这是有意为之：如果用总和，那么某一轮接收的伪标签\emph{数量}一多，伪标签项的总量就自动变大、凭空压过人工项；取平均后，每类信号的“\emph{单条平均影响}”才可比，数量不再直接决定话语权（实际量级仍受逐条权重和损失大小影响）。一种示意性的训练目标写成
\begin{equation}
  \mathcal{L}^{(t)}(\bm{\theta})
  =
  \frac{1}{|L|}\sum_{i\in L}\ell\!\left(f_{\bm{\theta}}(\bm{x}_i),y_i\right)
  +
  \lambda^{(t)}
  \frac{1}{|P^{(t)}|}
  \sum_{(\bm{x}_j,\widetilde y_j^{(t)},q_j^{(t)})\in P^{(t)}}
  a_j^{(t)}
  \ell\!\left(f_{\bm{\theta}}(\bm{x}_j),\widetilde y_j^{(t)}\right),
  \label{eq:knowledge-acquisition-self-training-objective}
\end{equation}
其中$a_j^{(t)}\geq0$逐条控制单个伪标签的影响（例如对分数越高的给越大权重），$\lambda^{(t)}\geq0$整体控制伪标签项相对人工项的系数。在固定的可更新参数范围内优化这个目标，就得到下一轮模型。

但要看清这个目标\emph{实际在鼓励什么}：它惩罚的是“模型预测偏离\emph{当前训练答案}”，而训练答案里混着伪标签——这意味着\emph{一条错误的伪标签同样会被模型认真拟合}，自训练不仅可能延续错误、还可能通过反复强化让错误越陷越深。两类损失分别按各自样本数归一化，是为了避免“伪标签一旦变多，就在统一平均里悄悄改变两类信号的数量占比”；当然实际的相对量级最终仍由逐条权重$a_j^{(t)}$和损失数值决定。最后处理一个边界情形：当本轮接收集合$P^{(t)}$为空（没有任何候选达到阈值）时，应\emph{直接跳过}伪标签项，并依据开发证据决定是停止自训练、放宽/收紧接收规则、还是补充其他来源；绝不能用某种“分母为零时记为0”的约定，去掩盖“这一轮其实没有任何可接收样本”的事实。

% 本段末尾两条边注不能跨页续排，需连同段落检查可用高度。
\Needspace{110mm}
把“接收”和“训练”首尾相连，就得到自训练的完整循环：\emph{预测候选、选择或复核、用筛出的伪标签更新参数、再用新模型重新预测}，然后进入下一轮。为保证可追溯、可回退，每一轮都要保存当轮使用的人工标签、伪标签、各条权重和模型版本，并明确一条策略：\emph{历史}伪标签在下一轮是固定沿用、用新模型重新评估、还是允许被撤销。

加入节奏也有取舍。一次加入大量伪标签能迅速扩大训练覆盖，但一旦其中有\emph{成类的系统错误}，也会被同时放大、迅速固化；分批小量加入则便于逐轮对比“这一轮到底变好了还是变差了”，代价是要做更多轮预测和训练。继续训练（沿用当前参数和优化器状态接着训）与重新训练（从某个指定起点完整重走训练路径）也会得到不同结果，需要按当前任务实际比较，不能默认其一更好。自训练并非新思想，早在自适应模式识别的早期工作中就已出现\scite{scudder1965selftraining}。

而它最主要的失效模式有一个明确名称——\emph{确认偏差}（confirmation bias）：模型最初猜错的样本，被它自己贴上错误伪标签，又作为“正确答案”回灌训练，下一轮模型对这个错误\emph{更自信}，再用更自信的错误标签强化自身，形成自我实现的闭环。现代深度模型的研究清楚地观察到了这种“错误被自己的预测不断确认、放大”的现象\scite{arazo2020confirmation}。这也解释了为什么伪标签需要筛选、加权、并保留独立核验，而不能无条件全量采信。

自训练循环里还有一个必须明确的设计选择：\emph{每一轮由哪个模型来提供训练目标}？这就是任务教师的\emph{更新安排}，常见有三种，稳定性与适应性各不相同。

\emph{固定教师}：从头到尾都用同一个冻结版本的模型产生伪标签。好处是目标稳定、极易追溯（每轮标签来自同一来源），缺点是它\emph{永远无法吸收学生在训练中变得更好的判断}。\emph{定期替换}：训练若干轮、并经过检查后，把某个表现合格的学生版本\emph{晋升}为新教师。它允许学生采用更大或不同的结构，带来质的提升空间；但“角色互换”本身并不\emph{证明}预测更准——新教师可能只是把旧教师的错误\emph{原样继承}下来，甚至更自信地重复。\emph{随学生更新的教师}：教师频繁地跟随学生当前状态调整，适应性最强，但目标一直在动，训练也最不稳定。三种安排要在“目标稳定可追溯”与“吸收模型进步”之间权衡，没有免费的最优解。

“随学生更新教师”有一个极端做法：每一步都直接把学生\emph{复制}成教师。但这样教师目标会随着学生参数的\emph{每一次抖动}而变化——学生朝错误方向晃一下，目标也跟着晃，训练信号极不稳定。

\term{Mean Teacher}针对这一点做了改进：教师不再是某一步学生的瞬时副本，而是\emph{历史各步学生参数的一个加权汇总（滑动平均）}，变化更平滑；训练目标则是让\emph{同一样本}在学生模型和教师模型下的预测\emph{保持一致}（一致性正则）\scite{knowledge-acquisition-mean-teacher2017}。具体的汇总方式是\emph{指数移动平均}（exponential moving average, EMA）：每更新一步，新教师参数里保留较大比例$\alpha$的旧教师参数、只掺入$(1-\alpha)$比例的新学生参数，于是教师像一个“滞后而稳定”的学生影子。要准确理解EMA的身份：它调节的只是\emph{目标更新的快慢与平滑度}，让目标不被单步噪声带着走；它\emph{完全不提供}“目标正确”的保证——平均一个一直犯错的模型，得到的仍是一个平滑地犯错的教师。

为了算出这一更新，需要同时定义预测差异与模型参数变化。第$t$步用同一样本的两种允许扰动$\widetilde{\bm{x}}_s,\widetilde{\bm{x}}_t$分别送入学生模型和教师模型；扰动须按任务设计，不应改变期望保持的标签含义。差异函数$d$取值越小表示两项预测越一致，加上真实标签的监督损失，得到学生要减小的目标。学生模型参数为$\bm{\theta}^{(t)}$，同结构教师模型参数为$\overline{\bm{\theta}}^{(t)}$，初始化可从同一模型复制；当前梯度只更新学生，然后才按EMA更新下一步教师：
\begin{equation}
  \begin{aligned}
    \mathcal{L}_{\mathrm{con}}^{(t)}
      \bigl(\bm{\theta},\overline{\bm{\theta}}^{(t)}\bigr)
    &=
    d\!\left(
      f_{\bm{\theta}}(\widetilde{\bm{x}}_s),
      f_{\overline{\bm{\theta}}^{(t)}}(\widetilde{\bm{x}}_t)
    \right), \\
    \mathcal{J}^{(t)}
      \bigl(\bm{\theta};\overline{\bm{\theta}}^{(t)}\bigr)
    &=
    \mathcal{L}_{\mathrm{sup}}^{(t)}(\bm{\theta})
    +\lambda_{\mathrm{con}}^{(t)}
      \mathcal{L}_{\mathrm{con}}^{(t)}
      \bigl(\bm{\theta},\overline{\bm{\theta}}^{(t)}\bigr), \\
    \bm{\theta}^{(t+1)}
    &=
    \bm{\theta}^{(t)}
    -\eta^{(t)}
      \left.
      \nabla_{\bm{\theta}}\mathcal{J}^{(t)}
      \bigl(\bm{\theta};\overline{\bm{\theta}}^{(t)}\bigr)
      \right|_{\bm{\theta}=\bm{\theta}^{(t)}}, \\
    \overline{\bm{\theta}}^{(t+1)}
    &=
    \alpha\overline{\bm{\theta}}^{(t)}
    +(1-\alpha)\bm{\theta}^{(t+1)}.
  \end{aligned}
  \label{eq:knowledge-acquisition-mean-teacher}
\end{equation}
逐个解释这些量。$\mathcal{L}_{\mathrm{sup}}^{(t)}$是\emph{真实标签}带来的普通监督损失，它提供一个外部的、正确的任务参照，防止师生一起偏离任务；$\lambda_{\mathrm{con}}^{(t)}\geq0$控制一致性项相对监督项的权重；$\eta^{(t)}>0$是学生梯度下降的步长；$\alpha\in[0,1)$决定EMA中\emph{保留旧教师}的比例，$\alpha$越大教师越稳定、更新越慢。

为什么要加一致性损失？它约束的是“\emph{同一个样本}在两种允许的轻微扰动下，两次任务预测不应差太多”——这是一种合理的平滑性先验；但真正锚定“预测要对”的，仍然是真实标签项。计算次序对理解算法很关键：求梯度时教师参数$\overline{\bm{\theta}}^{(t)}$是\emph{固定}的，梯度只更新学生；直到最后一行才用EMA更新教师，而下一轮的一致性目标才由新教师$f_{\overline{\bm{\theta}}^{(t+1)}}$重新计算。

两个常见误读要避免。其一，EMA平均的对象是\emph{模型参数}，不是“每个样本的历史预测分数”，二者在非线性模型下不等价（下一段的数值例子会说明）。其二，一致性损失\emph{下降}并不必然是好事：如果教师本身是错的，学生向教师靠拢只会让两者\emph{一致地错}，一致性变好了、正确率却没有；因此一致性仍需在独立参照标签上检查，不能自我证明。

用一个只有单个标量参数的玩具例子，把EMA那一行亲手算一遍。设$\alpha=0.9$，旧教师参数为$0.6$，本轮梯度更新后学生参数为$1.0$，则更新后的教师参数为
\[
  \overline{\bm{\theta}}^{(t+1)}=0.9\times0.6+0.1\times1.0=0.64.
\]
注意：算出$0.64$只是得到了一组新参数，\emph{下一轮的预测目标}还要把输入重新喂给$f_{0.64}$才能得到，参数平均和产生目标是两步。更重要的是，由于模型$f$对参数通常是\emph{非线性}的，“先平均参数、再过模型”一般\emph{不等于}“先分别预测、再平均预测”，即
\[
  f_{0.64}\ \neq\ 0.9f_{0.6}+0.1f_{1.0}.
\]
这提醒我们不要把Mean Teacher直观地误想成“在平均模型给出的答案”；它平均的始终是参数空间里的位置。

要厘清一个层次：上面的例子只更新了教师\emph{参数}，并\emph{没有}决定要不要把某个预测当成\emph{硬伪标签}接收。Mean Teacher\emph{原生}的机制只是软一致性（让两边预测接近），它\emph{并不要求}再做一次“分数过阈值就当确定标签”的硬筛选。软一致性和硬伪标签是两种可独立使用、也可组合的东西。

为了演示“若组合会发生什么”，另\emph{人为构造}一组数值：设同一个候选在更新前由$f_{0.6}$给出接收分数$0.72$，更新后由$f_{0.64}$给出$0.78$，而阈值固定为$0.75$。那么更新前$0.72<0.75$不予接收，更新后$0.78\geq0.75$，这一轮才把教师给出的类别作为一个\emph{硬}伪标签放入$P^{(t+1)}$。必须强调$0.72$、$0.78$是\emph{额外假设}的分数，它们\emph{无法}由参数$0.6$、$0.64$推算出来（参数不直接决定接收分数）；这个例子也\emph{不}说明分数会随参数更新单调上升——换个候选完全可能下降。

\begin{figure*}[!t]
  \centering
  % 图4-4：上下两面板共享169 mm版心；模型参数读取、梯度与EMA使用独立路由。
\begin{tikzpicture}[
  x=1mm,
  y=1mm,
  font=\figurefont\footnotesize,
  text=BookInk,
  loop node/.style={book node,draw=FigureModel,fill=FigureModel!7,line width=0.8pt,text width=24mm,minimum height=12mm,inner sep=1.2mm},
  data node/.style={book node,draw=FigureInput,fill=FigureInput!7,line width=0.8pt,text width=24mm,minimum height=12mm,inner sep=1.2mm},
  audit node/.style={book node,draw=BookAmber,fill=BookAmber!7,line width=0.8pt,text width=24mm,minimum height=12mm,inner sep=1.2mm},
  state node/.style={book node,draw=BookMuted,fill=BookPaper,line width=0.7pt,text width=24mm,minimum height=12mm,inner sep=1.2mm},
  loss node/.style={book node,draw=FigureLoss,fill=FigureLoss!7,line width=0.8pt,text width=24mm,minimum height=12mm,inner sep=1.2mm},
  main flow/.style={knowledge arrow,draw=FigureModel,line width=0.9pt,rounded corners=1mm},
  evidence flow/.style={knowledge arrow,draw=BookAmber,line width=0.8pt,rounded corners=1mm},
  external flow/.style={knowledge arrow,draw=BookMuted,line width=0.8pt,rounded corners=1mm},
  gradient flow/.style={knowledge arrow,draw=FigureLoss,line width=1.0pt,dashed,rounded corners=1mm},
  ema flow/.style={knowledge arrow,draw=BookMuted,line width=0.9pt,densely dotted,rounded corners=1mm},
  edge text/.style={inner sep=0.6mm,font=\figurefont\footnotesize,text=BookMuted,align=center}
]
  \draw[BookRule,line width=0.7pt] (1,-74) -- (167,-74);

  \node[anchor=north west,font=\figurefont\bfseries] at (1,1) {A\quad 获取与版本决策};
  \node[data node] (data) at (18,-14) {人工标签与\\待处理样本};
  \node[loop node] (model) at (56,-14) {当前模型\\版本$M^{(t)}$};
  \node[loop node] (candidate) at (94,-14) {候选预测\\及原始观测};
  \node[audit node] (selection) at (139,-14) {选择或复核\\接收训练目标};
  \node[state node] (input-change) at (18,-38) {预测请求修改\\调用安排};
  \node[loop node] (update) at (56,-38) {训练更新\\形成$M^{(t+1)}$};
  \node[audit node] (decision) at (94,-38) {继续、停止\\还是回退？};
  \node[audit node] (audit) at (139,-38) {开发验证\\过程审计};
  \node[state node] (previous) at (56,-61) {上一可用版本\\及旧结果};
  \node[state node,draw=FigureOutput] (freeze) at (94,-61) {冻结选定方案\\模型与预测请求};
  \node[state node,draw=FigureOutput,fill=FigureOutput!6] (holdout) at (139,-61) {留出验收\\独立质量证据};

  \draw[main flow] (data.east) -- (model.west);
  \draw[main flow] (model.east) -- (candidate.west);
  \draw[main flow] (candidate.east) -- (selection.west);
  \draw[evidence flow] (selection.south) -- (audit.north);
  \draw[evidence flow] (audit.west) -- (decision.east);
  \draw[evidence flow] (decision.west) -- node[edge text,above] {继续} (update.east);
  \draw[main flow] (update.north) -- (model.south);
  \draw[external flow] (input-change.east) -- (38,-38) -- (38,-14) -- (model.west);
  \draw[evidence flow] ([xshift=-6mm]decision.south) -- (88,-49) -| (previous.north);
  \node[edge text] at (72,-52) {回退};
  \draw[evidence flow] (decision.south) -- node[edge text,right] {停止} (freeze.north);
  \draw[main flow] (previous.east) -- (freeze.west);
  \draw[main flow] (freeze.east) -- (holdout.west);

  \node[anchor=north west,font=\figurefont\bfseries] at (1,-79) {B\quad Mean Teacher};
  \node[loop node,text width=22mm] (student) at (27,-95) {学生模型参数\\$\bm{\theta}^{(t)}$};
  \node[loop node,text width=22mm] (teacher) at (139,-95) {教师模型参数\\$\overline{\bm{\theta}}^{(t)}$};
  \node[state node,text width=22mm] (student-output) at (27,-114) {学生预测\\第$t$步};
  \node[loss node] (consistency) at (83,-114) {一致性损失\\$\mathcal{L}_{\mathrm{con}}^{(t)}$};
  \node[state node,text width=22mm] (teacher-output) at (139,-114) {教师目标\\不回传梯度};
  \node[loss node,text width=22mm] (supervised) at (27,-133) {监督损失\\$\mathcal{L}_{\mathrm{sup}}^{(t)}$};
  \node[data node,text width=12mm] (labels) at (61,-133) {真实标签};
  \node[loop node,text width=22mm] (optimizer) at (27,-153) {优化器\\学生梯度步};
  \node[loop node,text width=22mm] (student-next) at (70,-153) {学生模型参数\\$\bm{\theta}^{(t+1)}$};
  \node[loop node,text width=12mm] (ema) at (111,-153) {EMA};
  \node[loop node,text width=22mm] (teacher-next) at (149,-153) {教师模型参数\\$\overline{\bm{\theta}}^{(t+1)}$};
  \node[state node,text width=22mm,draw=FigureOutput,fill=FigureOutput!6] (next-target) at (149,-172) {重新预测\\第$t+1$步目标};

  \draw[main flow] (student.south) -- (student-output.north);
  \draw[main flow] (teacher.south) -- (teacher-output.north);
  \draw[main flow] (student-output.south) -- (supervised.north);
  \draw[main flow] (student-output.east) -- (consistency.west);
  \draw[main flow] (teacher-output.west) -- (consistency.east);
  \draw[main flow] (labels.west) -- (supervised.east);
  \draw[main flow] (student.west) -- (6,-95) |- (optimizer.west);
  \node[edge text,anchor=west] at (6.5,-124) {旧值};
  \draw[gradient flow] (supervised.south) -- (optimizer.north);
  \draw[gradient flow] (consistency.south) -- (83,-143) -| ([xshift=6mm]optimizer.north);
  \draw[main flow] (optimizer.east) -- (student-next.west);
  \draw[ema flow] (student-next.east) -- node[edge text,above] {$1-\alpha$} (ema.west);
  \draw[ema flow] (teacher.east) -- (164,-95) -- (164,-136) -| (ema.north);
  \node[edge text] at (137,-132) {旧教师模型参数，权重$\alpha$};
  \draw[ema flow] (ema.east) -- (teacher-next.west);
  \draw[main flow] (teacher-next.south) -- (next-target.north);
\end{tikzpicture}

  \caption{候选监督信号进入下一轮模型更新。面板A的开发验证与过程审计决定继续、停止或回退，选定方案冻结后另做留出验收。面板B实线是计算与模型参数读取，橙色虚线是只更新学生的梯度，灰色点线是EMA。优化器使用旧学生模型参数与两项损失的梯度；EMA结合新学生与旧教师模型参数，再由新教师重新预测下一步目标。}
  \label{fig:knowledge-acquisition-iterative-update}
\end{figure*}

同样用“教师—学生、EMA更新”这套形式，\emph{目标内容}不同，方法的性质就不同，不能只凭公式形式相同就划等号。Mean Teacher里被要求一致的是\emph{具有明确任务含义的预测}（物种类别的得分），所以这里的教师确实在提供“任务目标”。而第~\ref{sec:data-derived-knowledge-cross-view-prediction}节介绍的跨视图方法，其目标分支产生的是一个\emph{表示向量}，让另一分支去预测这个向量，全程\emph{不涉及}“这是什么物种”的名称。两种方法都可能用EMA维护目标网络，但一个在对齐任务答案、一个在对齐表示，差别的根源是\emph{目标内容}，而不是更新公式或“教师”这个称呼。

无论硬伪标签还是软一致性，都共享自训练家族的固有风险：它们都可能\emph{继承模型的初始错误}（没有新的外部信息进来，难以凭空纠错），更新过程也不会\emph{自动补足}稀有类别（模型本就罕见的类别很难产生高质量伪标签）。因此运行时要同时监测接收规则是否合理、各类别覆盖是否均衡、在\emph{人工参照标签}上的真实任务表现、以及所耗资源，不能只看一致性或损失在下降。

与Mean Teacher“每步平滑更新参数”不同，另一种做法是\emph{整轮训练完、检查后再替换}教师，Noisy Student是其代表：教师先为未标注样本生成伪标签，学生同时用真实标签和这些伪标签训练，并在训练学生时刻意\emph{注入噪声}（数据扰动、模型噪声等），训练完成后把这个（带噪声训练出的、通常更强的）学生\emph{晋升}为下一轮教师，如此迭代\scite{knowledge-acquisition-noisy-student2020}。

要把它和Mean Teacher分清楚：Noisy Student是\emph{阶段性的版本替换}（一轮结束、整个学生换上去当教师），而\emph{不是}每一步对参数做EMA移动平均。注入噪声的目的，是迫使学生不能只死记伪标签、而要学到更鲁棒的规律。但同样地，“学生训完、被升为教师”只是完成了一次版本更替，新教师到底\emph{可不可用}仍须经过开发集检查——训练这个动作本身不产生质量保证。

\FloatBarrier
\subsubsection{基于正例与未标注样本的训练}
\label{sec:ka-pu-learning}

自训练虽然不要求全部标签，但它好歹对各类\emph{都}有一些（哪怕是伪的）参照。还有一种信息更少的情形：只确认了\emph{一批正例}，其余样本的类别完全未知。比如花卉清单里只可靠地确认了一批“山茶”照片，其他大量照片\emph{既有山茶也有别的花}、却没有逐一辨认。

这时两种朴素做法都有毛病。其一，直接把所有未标注照片\emph{当成“非山茶”}（负例）来训，但其中其实\emph{混着尚未标出的山茶}，于是把真正例错标成负例，严重扭曲决策边界。其二，只对已确认的山茶照片做单类拟合，却没有任何信号约束模型\emph{该如何对待}其他照片，学不出一个正常的二分类器。

\term{正例与未标注样本学习}（positive-unlabeled learning, PU learning）提供了第三条路：它\emph{不}给任何一张未知照片\emph{指派}确定的正负标签，而是直接用“一批已确认正例”加“一批正负混杂的未标注样本”，在\emph{风险（期望损失）层面}把一个正常二分类器的训练目标\emph{估计}出来。关键要理解它改变的是\emph{损失的估计方式}，自始至终没有逐条宣布任何未知照片的真值。

先把任务和假设严格列清，因为PU结论对假设很敏感。固定一个二分类任务：$Y\in\{+1,-1\}$，正类表示山茶；$X$是照片输入；$g$是模型给出的实值得分函数（据此分正负）；$\ell(g(X),Y)\geq0$是非负、可积的损失。三个关键假设如下。第一，正类\emph{先验比例}$\pi=\mathbb P(Y=+1)\in(0,1)$是\emph{已知}的（或可另行估计）。第二，那批已确认正例是\emph{独立地}从正类条件分布$P_+=P(X\mid Y=+1)$中抽取的。第三，未标注样本独立地取自\emph{总体边缘分布}$P_X=P(X)$，且两批抽样相互独立。再记$P_-=P(X\mid Y=-1)$为负类分布——注意它\emph{只在推导公式时出现}，我们手上\emph{并没有}直接抽得的负类样本，这正是PU要绕过的困难。

为了后面推导，定义四个“带上下标”的风险记号，读法是“\emph{上标}表示这个损失\emph{暂且按哪个标签}算，\emph{下标}表示样本\emph{来自哪个总体}”：
\[
  R_+^+(g)=\mathbb E_{P_+}\ell(g(X),+1),\quad
  R_+^-(g)=\mathbb E_{P_+}\ell(g(X),-1),
\]
\[
  R_-^-(g)=\mathbb E_{P_-}\ell(g(X),-1),\quad
  R_U^-(g)=\mathbb E_{P_X}\ell(g(X),-1).
\]
例如$R_+^-$是“拿\emph{正类总体}$P_+$的样本、却\emph{暂且当作负类}（标签$-1$）去算损失”的期望——它不是真实风险，只是推导中一块可由正例算出的量。

推导的核心是一个简单的“混合—扣除”思想。未标注总体$P_X$按全概率法则，由比例$\pi$的正类和比例$1-\pi$的负类混合而成。把“未标注样本暂且按负类算”的风险$R_U^-$按这个混合比例拆开，就有
\[
  R_U^-=\pi R_+^-+(1-\pi)R_-^-.
\]
这个等式里，$R_U^-$可用未标注样本算、$R_+^-$可用已确认正例算、$\pi$已知，唯一\emph{没有直接样本}的是负类项$(1-\pi)R_-^-$。把它\emph{解出来}（移项扣除正类的贡献），就能用“可观测的量”拼凑出本不可观测的负类风险：
\begin{equation}
  R(g)=\pi R_+^+(g)+(1-\pi)R_-^-(g)
      =\pi R_+^+(g)+R_U^-(g)-\pi R_+^-(g).
  \label{eq:ka-pu-risk}
\end{equation}
这个式子的意义在于：右边的三项期望\emph{都}能从我们\emph{实际拥有}的两批数据（正例、未标注样本）上、用对应的非空样本集求\emph{平均损失}来估计，记为$\widehat R_{\mathrm{PU}}$——于是在\emph{没有一个真负例}的情况下，我们仍然得到了完整分类风险$R(g)$的一个估计量。

但要严格把握它成立的条件与一个统计细节。在$g$固定、$\pi$准确、且前述抽样条件（正例代表$P_+$、未标注代表$P_X$、两批独立）都满足时，这个估计是\emph{无偏}的。然而，一旦我们用\emph{同一批样本}去\emph{优化、挑选}$g$，选出的$g$就已经\emph{依赖}这批数据了；“对某个固定函数无偏”不能自动升级成“学出来的分类器有泛化保证”——无偏估计与模型选择后的保证是两回事。

无偏PU估计在理论上漂亮，落到有限样本训练却有一个工程问题。看$R(g)=\pi R_+^+(g)+R_U^-(g)-\pi R_+^-(g)$的估计，后两项是\emph{两个经验平均相减}。被减出来的那一项$(1-\pi)R_-^-$在\emph{真值}上必为非负（它是“负类上的非负损失”乘以正系数），但在\emph{有限样本估计}中，两个带噪声的平均值相减完全可能得到一个\emph{负数}。

这会给“灵活、表达力强的模型”一个可钻的空子：它可以通过把训练数据拟合得更激进，\emph{不断压低这个本应为负的经验项}，让总目标一路降到没有意义的负值，表面上“损失在稳定下降”，实际是在过拟合估计噪声。为堵住这个漏洞，Kiryo等人提出\emph{非负PU}（nnPU）准则：直接把这部分经验估计\emph{截断在零}处，不允许它为负——
\begin{equation}
  \widehat R_{\mathrm{nnPU}}(g)
  =\pi\widehat R_+^+(g)
   +\max\{0,\widehat R_U^-(g)-\pi\widehat R_+^-(g)\}.
  \label{eq:ka-nnpu-risk}
\end{equation}
截断带来一个权衡，要如实说明：它\emph{引入了轻微的估计偏差}（不再严格无偏），但换来了训练的稳定——模型无法再靠“把一个本该非负的量压成负的”来虚假地降低目标。还要澄清截断\emph{没有}做到的事：它\emph{不是}把未标注样本里的负例\emph{找出来}了，自始至终我们仍不知道任何一张未标注照片的正负身份；nnPU只是修正了\emph{风险估计}的形状，并未产生逐样本标签\scite{ka-kiryo2017-nnpu}。

用花卉任务做一个纯教学的数值演算（这里特意取零一损失，因为它便于按“对/错比例”计数理解）。假设正类先验$\pi=0.2$，三个经验平均损失分别为$\widehat R_+^+=0.2$、$\widehat R_+^-=0.8$、$\widehat R_U^-=0.3$。则被减出的负类贡献估计为
\[
  \widehat R_U^--\pi\widehat R_+^-=0.3-0.2\times0.8=0.14,
\]
而正类项为$\pi\widehat R_+^+=0.2\times0.2=0.04$，合起来总风险估计为$0.04+0.14=0.18$。

再看截断在什么情况下起作用。设想另一个候选模型把$\widehat R_U^-$压到了$0.1$，另外两项不变。此时无偏经验准则会算出
\[
  0.04+\underbrace{(0.1-0.16)}_{-0.06}=-0.02,
\]
一个\emph{负的“风险”}——这在真实风险里不可能发生，正是有限样本估计失控的信号；而非负准则把括号截到0，给出$0.04+0=0.04$。这组数字只是用来演示“截断处理的是哪一部分、怎样阻止目标变负”，它们\emph{不是}在独立评价集上测出的实际错误率。最后提醒，零一损失不可微、无法直接梯度优化，真正训练时通常会另选一个可微的替代损失（如sigmoid或对数损失）来实现同一准则。

PU学习把对数据的要求转移到了\emph{抽样代表性}和\emph{先验$\pi$}上，这两点失准都会动摇式~\eqref{eq:ka-pu-risk}的全部推导。其一，如果确认正例只来自“最容易辨认的近景清晰照片”，它们就\emph{不能代表}全部山茶照片（大量山茶是远拍、遮挡、非常规形态），那么“正例取自$P_+$”的代表性假设失效，风险估计随之有偏。其二，如果正类先验$\pi$本来未知、需要估计，那么$\pi$的估计还要满足额外的可识别条件，估错$\pi$会让扣除的比例系统性偏差。

成本上也要有数：每一步训练都要同时维护两批样本、分别算平均损失和梯度，批次如何构造、截断如何施加、$\pi$如何估计，都要计入训练与验证的投入。还要避免与另一种“缺失标签”方法混淆：第~\ref{sec:ket-sharing-labels}节的逆概率校正，针对的是“在给定输入后标签缺失与真实值独立”的情形；那种缺失机制\emph{不能直接搬用}到PU——PU的未知恰恰是\emph{按真实类别}区分的（正例被挑出、其余混合），缺失依赖于类别，两者的统计假设不同。最后，无论训练时风险估计多么精巧，选定模型后仍应在一批\emph{能独立判定正负}的评价样本上检查真实任务风险；若做不到，就必须明确声明自己采用的是哪一种PU风险估计、依赖哪些条件，而不能含糊地宣称“效果等同全监督”。

\subsection{输入内容的更新}
\label{sec:knowledge-acquisition-input-update}

有些错误无论怎么改模型都无法修复，一个最直白的例子是：若山茶/茶梅的混淆是因为照片\emph{根本没拍到叶片}，那么更新模型参数\emph{不可能把一片不存在的叶片补进当前照片}——模型学不到输入里没有的信息。遇到这类问题，正确的修改对象是\emph{预测输入}，而不是模型。

在动手改之前，应先回头检查模型这一次\emph{实际看到了什么}，按4.2节的四类输入逐项排查：任务要求里类别边界写清了吗？样本上下文是否缺少关键部位（照片缺叶片）？示范是否只给了正例、缺了反例？辅助证据的检索是否只匹配到了背景而非鉴别特征？这四类缺陷分别对应任务要求、样本上下文、示范和辅助证据，而且每一类修改\emph{影响哪些旧调用}的范围也不同（改一条全局规范可能影响所有调用，补一张照片只影响那一个对象），修改时要据此限定重算范围。

确定了错误成因，修改时要把输入的变化保存成\emph{可前后比较的版本}，而不是直接覆盖旧输入（覆盖后就无法判断这次修改到底改了什么、影响了谁）。用四个构造修改对应四类缺陷：规范从$r_1$更新到$r_2$，增加“叶缘锯齿、花期只能在对应部位可见或拍摄时间已知时才允许使用”；示范从$d_1$到$d_2$，补入山茶与茶梅的反例；样本从$x_1$到$x_2$，补入同一植株的叶片；证据从$e_1$到$e_2$，用一份现行名录换掉过时名录。

每一项修改都应\emph{带着出处记录}：是哪一个错误样本触发了它、它适用于哪些对象、哪些旧调用因此需要重算。这样“补充一份观测或证据”就落实成了\emph{具体的内容变更和明确的依赖关系}，而不是用一份新输入笼统地盖掉旧输入——后者会让版本之间无法对账。

要判断某次输入修改\emph{有没有用}，应采用控制变量的比法：固定模型、开发样本和评价口径，\emph{只替换四类输入中的某一类}，若涉及随机生成还采用完全相同的采样配置，然后比较前后差异。如果相对旧版本物种错误减少、且弃权行为仍然符合规范，这才构成“该输入修改\emph{在这批样本上}有效”的\emph{一项}证据（注意只是一项、且限于这批样本）。

比较时不能只盯错误率，还要同时看\emph{作答覆盖率}：模型可能因为新提示让它更频繁弃权，从而“答得少、错得少”，错误数下降只是弃权变多的副产品，并不代表判断能力提升。还要防止错误\emph{归因}过头：补说明、换反例、加叶片分别检验的是规则表达、示范和观测各自的作用，但某一项改动后效果变好，不能\emph{单独}证明它就是唯一成因（可能有交互作用）。最后警惕过拟合到开发集：反复针对同一小批错误样本调整输入，等于把输入设计\emph{拟合}到这些样本；因此仍要在新批次上验证，并检查新内容是否挤占了输入空间、不同说明之间是否冲突、以及额外检索带来的成本。

最后强调\emph{模型变化}与\emph{输入变化}必须分开记录、分别归因。如果一轮里既换了模型版本、又改了提示和证据，那么看到结果变化时，根本无法判断差异应归功（或归咎）于新参数还是新材料——只能得到“这个组合版本如何”。若确实需要区分二者的作用，就应分阶段试验：一次只变一样。只有当完整方案——模型、规范、上下文整理规则、示范、证据的\emph{所有}版本——都冻结之后，才进入第~\ref{chap:knowledge-quality-assessment-correction}章的独立质量评价。本节的职责，是把“诊断出的错误”落实成“可追溯的输入修改”，让每一次输入变更都有因可查、有范围可重算。

\subsection{获取安排的更新}
\label{sec:knowledge-acquisition-schedule-update}

模型和输入的版本更新完，还有一个调度问题：\emph{下一轮有限的资源先花在哪些对象上}？不同对象此刻的处境不同，需要的动作也不同。例如S042仍然缺叶片（问题在观测，需要补拍或转人工），而另一些照片已经有了可用答案、但那个答案是用\emph{旧版本模型}产生的（问题在时效，可能只需重算）；对这两类对象做同一件事显然不合理。\term{获取安排}（acquisition schedule）就是用来规定：处理哪些样本、各调用几次、哪些要补证据、哪些转人工，以及这轮是全量重算还是只局部重算。

查询一个对象的状态时，要把\emph{内容、验证、执行、版本依赖}这几个\emph{不同层面}分开记录，不能把它们压成一个互斥的状态码。“已确认可用”（验证层）、“需要重算”（版本依赖层）、“证据不足”（内容层）、“等待人工”（执行安排层）描述的是不同维度，\emph{可以同时成立}。比如S042完全可以\emph{同时}处于“物种内容未知”和“正在等待人工补观测”两个状态——它们不矛盾，硬要归成一个“当前状态”反而会丢信息。

执行状态仍严格沿用全卷统一的五态协议：尚未发起请求是“未请求”，已排队但未开始是“未处理”，已开始但未提交是“处理中”，操作正常完成是“成功”，异常终止、没跑完是“失败”。两个反复强调的区分在此依旧成立：模型\emph{成功地返回了弃权}（正常运行后判断证据不足）仍是“成功”执行；而程序\emph{失败}中断时根本没有产生“主动弃权”这个动作，不能混为一谈。

特别要澄清“需要重算”的身份：它表示某条旧标注所依赖的版本已经变化、结果可能过期，是一种\emph{版本依赖}关系，\emph{不是}第六种执行状态。处理重算时，应当为新请求\emph{另记}一条进度和新版本，并明确指定\emph{当前采用}哪一个结果，避免新旧两个答案都挂在那里、被下游误认为都是现行有效的标注。

版本变更后，\emph{重算范围}不应靠拍脑袋，而应沿“变更—依赖”关系\emph{计算}出来，不同类型的变更波及面差别很大。整体模型参数更新影响最广：所有由旧模型直接预测的结果原则上都可能变化，通常需要大范围重算。只改某个类别映射，则可把范围缩小到涉及该映射的那批标注。更新一份外部名录，要定位所有\emph{引用该来源}、或其适用范围变化会影响到的候选。撤回一个近邻参考标签，要反向找出\emph{读取过它}的那些查询（近邻结果依赖它）。改动图的一条边，影响可能沿连通区域扩散到\emph{一整片}节点。

可见，要做\emph{局部}刷新，前提是保存了足够完整的“调用依赖清单”（谁的结果用了哪个模型、哪份证据、哪条参考、哪条边）；没有这份清单就算不出准确的受影响范围，此时为了省调用而只重算一部分，会让库里同时残留新旧两种版本的结果，造成不一致——宁可全量重算，也不能留这种隐患。

下一轮预算怎么花，应当\emph{与要补的具体缺口一一对应}，而不是笼统地“多调几次模型”：调用新样本是为了扩大\emph{覆盖}；对同一输入重复生成、取多个样本，是为了观察答案的\emph{随机波动}；对困难样本追加复核，是为了\emph{查错}；刷新旧答案是为了消化\emph{版本变化}；转人工接管则是为了补入模型\emph{根本不具备}的判断依据（比如看不见的部位、需要外部事实的判断）。

这里要破除一个数量幻觉：如果用\emph{完全相同}的模型、输入和证据反复调用，得到的答案高度相关，那么\emph{调用次数增加}并不等于\emph{独立的标注依据增加}——同一个不确定来源问十遍，不会变成十条独立证据，反而只是在重复。表~\ref{tab:knowledge-acquisition-iteration-actions}据此把“反馈情形—可定位的范围—可选动作—额外投入—过程检查”列在一起，提供的是\emph{可选的应对菜单}，而不是对每种现象规定唯一正确的动作；具体选什么仍需结合证据与成本判断。

\begin{table*}[htbp]
  \centering
  \small
  \setlength{\RaggedRightParindent}{0pt}
  \renewcommand{\arraystretch}{1.15}
  \begin{tabularx}{\linewidth}{@{}>{\RaggedRight\arraybackslash}p{22mm}*{4}{>{\RaggedRight\arraybackslash}X}@{}}
    \toprule
    反馈情形 & 定位范围 & 可选动作 & 额外投入 & 过程检查 \\
    \midrule
    模型更新 & 旧模型参数或旧表示的依赖项 & 全量或局部刷新 & 训练、重建与推断 & 损失、覆盖与错误类型 \\
    规范变化 & 语义或字段变化的样本标注 & 重映射、重算或复核 & 规范评审与重算 & 边界案例与新增歧义 \\
    证据缺失 & 缺少必要观测或证据的对象 & 补观测、检索或转人工 & 采集与来源核验 & 对象对应、弃权与错误 \\
    局部错误集中 & 特定类别、来源或图区域 & 定向复核与局部刷新 & 诊断、重建与复核 & 局部收益与其他区域退化 \\
    收益停滞 & 多轮重复的结果 & 停止、回退或换来源 & 比较与回退恢复 & 边际收益、覆盖与成本 \\
    \bottomrule
  \end{tabularx}
  \caption{反馈类型与可调整的获取安排。按反馈定位依赖范围，再比较可选动作、额外投入与过程结果；各行不是自动触发规则。开发验证和过程审计可支持继续、停止或回退，参与选择的样本不再作为冻结方案的最终独立验收依据。}
  \label{tab:knowledge-acquisition-iteration-actions}
\end{table*}

定位重算范围时不能只数\emph{直接}依赖，还要算上\emph{间接}的连锁影响。表示模型一旦更新，依赖该表示构建的近邻索引和图可能都要\emph{重建}；规范变化可能连带引起类别映射和输出解析器的改动；一条证据更新，要追溯回它的原始来源、以及所有用了它的样本标注。即使是一次看起来很小的局部修正，改完也应回头检查\emph{其他区域}：近邻的参考关系、图上的边都可能把这次变化\emph{传播}到一些原本没有错误的对象上——修好了A，却通过依赖关系意外影响了B，这种“修复的副作用”只有沿依赖链检查才能发现。

把这些原则落到花卉任务，一批照片可以\emph{分类安排}：用旧模型产生、依赖已变化的对象安排重算；缺关键观测的照片安排补拍叶片；仍有歧义、模型无法定夺的对象转人工。各走各的通道，而不是统一重跑。

关于\emph{何时停止迭代}，要有据可依：如果新版本只是在那批被反复调试的样本上降低了错误、换到\emph{新的开发批次}却没有带来新增收益，就应结合成本决定\emph{停止}或\emph{回退}到上一版；还要警惕“预测输出逐轮变化越来越小”被误读成收敛变好——它也可能只是模型在\emph{重复同一个错误}、输出稳定地错着，输出变化小与质量提高毫无必然关系。最后再守一次数据纪律：一旦把最终测试样本用于“是否继续/回退”的选择，它就\emph{转为开发依据}，必须\emph{另留}一批方案冻结后的独立验收数据。本节处理的始终是模型获取流程\emph{内部}的执行顺序与版本一致性；需要人工判断的具体方法见第~\ref{chap:knowledge-acquisition-human}章，跨多个来源的资源分配见第~\ref{chap:knowledge-fusion}章。

\section{本章小结}
\label{sec:knowledge-acquisition-model-based-summary}

回顾全章：我们从一批没有物种答案的花卉照片出发，把“用模型获取知识”拆成了一条首尾相连、且每一步都可追溯的链路——先准备一个能用的模型，再为它组织这一次的输入，然后由明确的机制算出候选，最后根据复核反馈决定如何改进。

这条链路的核心精神是\emph{依赖可追溯}：每一条样本标注都必须能回答四个问题——它由一个\emph{经过何种准备}的模型给出、模型当时\emph{实际看到}了哪些任务要求与观测、用了哪些\emph{示范或辅助证据}、以及答案是通过\emph{哪条计算路径}（参数映射/近邻/图）产生的。正因为这些依赖都在，复核发现问题时才能对症决定：是改模型、补输入、还是重新安排下一轮对象；也正因为版本和适用范围被保存下来，每一个物种候选和每一次弃权结果才\emph{有依据被追查、被刷新或被撤回}，而不是一经产生就成了无法核对的黑箱结论。

三条候选计算机制的依据和风险各不相同，值得在收尾时再对比一次。\emph{直接预测}读取的是固化在模型参数里的任务关系，优点是推断时不依赖外部样本，风险在于训练覆盖不到、接口或数据分布变化时会失效。\emph{近邻参考}在调用现场读取带标签的参考样本，知识更新不必重训、即时生效，但它的成立依赖“距离确实反映任务相关特征”，并受参考标签质量、局部覆盖和检索遗漏的制约。\emph{图传播}以任务标签为锚点，能借助未知节点把类别得分传递到更远，天然适合利用样本间的连接，但也因此可能让一个错误锚点或一条跨类别错边的影响\emph{扩散得更广}。

三者可以自由组合在一个系统里，但有两条核查底线不变：不能因为两条路径最终\emph{类别相同}就省略对中间依赖的检查（过程不同、可靠性可能天差地别）；也不能因为图求解或模型训练在数值上\emph{收敛}了，就认定候选\emph{正确}（收敛只意味着达到了所设定目标的解，与是否符合事实是两回事）。

当模型候选被\emph{回流到训练}中时，要始终记得它们没有自带“正确”的保证。任务教师只是\emph{提供目标}的角色，伪标签只是\emph{暂时被采用}的任务答案，二者都可能错，都需要独立核验；弱到强泛化、自训练里的确认偏差都提醒我们，学生不必然纠正教师、自我强化也可能放大错误。在更新教师的具体技术上，Mean Teacher用模型参数的指数移动平均维护一个稳定教师、再让师生的\emph{任务预测}保持一致；Noisy Student则是整轮训练后做版本替换。要把它和\emph{表示学习}里的目标分支区分开：后者让一支网络去预测另一支产生的\emph{表示向量}、全程不给出物种名称，因此不能只因为“都用了EMA、都叫教师”就把二者混同。

材料来源还可以进一步自动化：Self-Instruct让模型\emph{连同任务和输入一起}生成训练材料，于是检查范围也从“固定样本的标签对不对”扩大到“任务是否成立、输入是否合规、材料是否重复、分布是否合理”。在标签极少的极端情形，PU学习提供了另一条合法路径：在正例有代表性、未标注样本代表总体、且正类先验可知的条件下，可以直接\emph{重建}分类风险，而不必先把未知样本\emph{瞎猜成负例}。这一长串方法各管一段——\emph{生成}候选、\emph{筛选}接收、\emph{更新}参数、\emph{核验}质量回答的是四个不同问题，不能用其中任何一步的成功代替其余步骤。

最后回答“要不要继续迭代、何时收手”。决策依据应当是\emph{同一任务范围内}几项实质指标的权衡：错误是否下降、弃权是否恰当、覆盖是否扩大，以及为此新增的投入（算力、调用、人工）是否值得。\emph{不应}用那些与质量没有直接因果的表面量来判断——模型规模更大、输入提示更长、检索回的参考更多、迭代轮次更多，全都不必然意味着结果更好；输出趋于稳定也可能只是在重复错误。

还要守住贯穿全章的数据纪律：那些参与了版本选择（调阈值、选模型、挑输入、定教师更新方式）的样本，提供的只是\emph{开发依据}，不能再反过来证明冻结方案的最终质量，独立验收必须另留数据。归根到底，本章产出的是一批\emph{可追溯的任务候选}，而不是已盖棺定论的标签：多份候选结果如何\emph{综合}成结论，由第~\ref{chap:knowledge-fusion}章展开；它们的\emph{可靠性如何评估、错误如何修正、如何独立验收}，由第~\ref{chap:knowledge-quality-assessment-correction}章负责；而当我们更进一步，想\emph{不借助外部标签、只从观测内部}自动构造表示或分组目标时，就进入了第~\ref{chap:data-derived-knowledge-acquisition}章的内容。

\chapter{基于数据本身的知识获取}
\label{chap:data-derived-knowledge-acquisition}

前几章获取的依据——人工标签、规则结论、模型预测——都在直接回答“这张照片是什么物种”这类既定任务问题，但它们都依赖某种\emph{外部给出的答案}：人要花时间、规则要人写、模型也要先有标签训练。现在面对一个更“拮据”也更根本的情形：我们手上只有\emph{一大批没有逐张确认物种的花卉照片}，没有任何外部答案可用。

是不是没有标签，就完全无法从这批数据里学到东西？并非如此。照片本身仍然携带大量\emph{结构和规律}：花瓣、叶片、背景这些像素\emph{客观地存在于文件里}，可以被读取；同一张照片还能裁出两个视角、序列数据里有时间先后、相似照片之间天然成簇。换句话说，\textbf{缺少“类别答案”，不等于数据里没有任何可以拿来当训练目标的信息}。本章研究的就是：在不借助外部任务标签的前提下，怎样\emph{从观测自身}中“制造”出可训练、同时又可检查的依据。这类思路通常被称为自监督学习。

为了让读者一开始就抓住本章的全貌，先用同一张照片对比三种“无中生有”地制造训练目标的做法，注意它们的\emph{答案从哪里来}完全不同。

第一种：在照片上遮住一块花瓣，让模型根据可见部分\emph{预测被遮住的像素}。这里的答案在“遮住之前”就\emph{原样保留}在我们手里——原始图像本身既是出题者、又是标准答案。第二种：把照片裁成两块，让一个裁剪去\emph{预测另一个裁剪的特征向量}。这时答案不再是现成像素，而要先等编码器把图像\emph{转成特征向量}才会出现，且这个向量会随模型训练不断变化。第三种：先把一批照片按相似度\emph{聚成若干群组}，再让模型预测每张照片“属于第几组”。这里的答案（组号）干脆是由一个\emph{分组算法临时算出来}的，换一轮聚类组号可能全变。

三种做法表面上都在“做预测”，但它们产生的监督目标性质迥异（一个是固定的原始内容、一个是动态算出的向量、一个是算法赋予的临时编号）。因此沿用第~\ref{chap:knowledge-acquisition-overview}章的区分必须强调：这些被构造出来的目标或关系，\emph{首先只是候选监督信号}；只有当它们真正进入损失计算、驱动参数更新时，才承担起监督信号的作用。而“能把它们构造出来”这件事，\emph{远不能}证明它们对“识别物种”这个最终任务有用——这是全章要反复检验的核心问题。

按照“目标的依据从哪来”，本章把这些自监督方法清晰地分成三条路线，也对应三节内容。

第一条是\term{观测内容}（observed content）：素材就是已经采集下来的词元、像素或传感器读数；只要\emph{藏起一部分、保留另一部分}，被保留的部分天然就是答案（遮住像素让模型补、遮住后文让模型猜）。第二条是\term{样本关系}（sample relations）：不直接规定单个样本的答案，而是规定样本\emph{彼此之间}应有的关系——哪些样本本质相同、应保留共同特征（如同一照片的两个裁剪），哪些样本只是对照、允许被区分；再用这些关系去约束学到的表示。第三条是\term{样本分组}（sample grouping）：把多个样本组织成群组，用“组号”或“对某个代表向量的归属权重”作为临时答案。

需要说明，上下文、时间顺序、多模态这些概念\emph{并不是}独立的第四种目标来源：上下文和时间顺序只是\emph{限制观测内容如何切分}，多模态配对只是\emph{规定不同观测如何对应}，它们分别服务于第一、第二条路线。沿这三条路线学习时，我们始终要同时追踪两件事：其一，观测内容或样本关系具体\emph{怎样进入}目标和参数更新（梯度从哪来到哪去）；其二，训练结束后模型\emph{究竟保留了什么表示}、这些表示对下游物种识别是否真的有用。

在进入细节前，图~\ref{fig:data-derived-knowledge-target-sources}先给出一个贯穿全章、值得记住的关键区分：\emph{答案是“现成的”还是“算出来的”}。图左侧那条路径（对应5.1），目标直接从原始观测里\emph{分}出来——遮住的像素、后文的词元，答案在训练前就固定不变，模型只需要做“预测”这一个动作。图右侧两条路径（对应5.2、5.3）则不同：它们都要先经过一个\emph{编码分支}把输入变成向量，再拿这些\emph{训练当步才算出的向量}相互比较或匹配；模型参数一旦更新，这些向量目标\emph{下一步就会跟着变}。

这张图目前只回答了“目标在哪里产生”，\emph{还没有规定}右侧那两个会变的向量在训练时如何更新——是两路一起求导、相互靠拢，还是把其中一路暂时冻充当一个“固定靶子”？这正是\emph{共同求导}与\emph{停止梯度/独立目标分支}两种安排的分野，第~\ref{sec:data-derived-knowledge-training-objectives}节会专门区分；理解了“现成答案 vs 动态向量”这一条，后面 MAE、SimCLR、BYOL、SwAV 等众多方法的差异就有了统一的抓手。

\begin{figure*}[htbp]
  \centering
  \begin{tikzpicture}[
  x=1mm,y=1mm,
  font=\figurefont\footnotesize,
  panel/.style={draw=BookRule,fill=NNPanel,line width=.85pt,rounded corners=1.5mm},
  data/.style={draw=FigureInput,fill=FigureInput!12,line width=1.1pt,minimum height=10mm,align=center,inner sep=2mm},
  model/.style={draw=FigureModel,fill=FigureModel!12,line width=1.1pt,minimum height=10mm,align=center,inner sep=2mm},
  target/.style={draw=FigureOutput,fill=FigureOutput!10,line width=1.1pt,minimum height=10mm,align=center,inner sep=2mm},
  loss/.style={draw=FigureLoss,fill=FigureLoss!12,line width=1.1pt,minimum height=10mm,align=center,inner sep=2mm},
  flow/.style={knowledge arrow,draw=NNWire,line width=1.2pt,rounded corners=1mm},
  keep/.style={knowledge arrow,draw=FigureOutput,line width=1.1pt,rounded corners=1mm},
  annotation/.style={font=\figurefont\scriptsize,text=BookMuted,align=center}
]
  \draw[panel] (1,-76) rectangle (81,12);
  \draw[panel] (87,-76) rectangle (168,12);
  \node[font=\figurefont\bfseries\footnotesize,text=BookInk] at (41,7) {(a) 观测目标：答案随观测保留};
  \node[data,minimum width=24mm] (obs) at (41,-9) {完整观测\\$\bm{x}$};
  \node[data,minimum width=24mm] (visible) at (20,-28) {可见预测输入\\$\bm{x}_{V}$};
  \node[target,minimum width=24mm] (held) at (62,-28) {保留目标\\$\bm{x}_{M}$};
  \node[model,minimum width=24mm] (predict) at (20,-48) {模型预测\\$\widehat{\bm{x}}_{M}$};
  \node[loss,minimum width=24mm] (recloss) at (62,-48) {重建损失};
  \coordinate (obs-split) at (41,-18);
  \draw[NNWire,line width=1.2pt] (obs.south) -- (obs-split);
  \draw[flow] (obs-split) -| (visible.north);
  \draw[keep] (obs-split) -| (held.north);
  \draw[flow] (visible.south) -- (predict.north);
  \draw[flow] (predict.east) -- (recloss.west);
  \draw[keep] (held.south) -- (recloss.north);
  \node[annotation,text width=68mm] at (41,-66) {目标从原始观测中预先保留，\\不需要编码分支生成答案。};

  \node[font=\figurefont\bfseries\footnotesize,text=BookInk] at (127.5,7) {(b) 表示目标：答案在训练中计算};
  \node[data,minimum width=24mm] (raw) at (127.5,-9) {原样本\\$\bm{x}$};
  \node[data,minimum width=24mm] (viewa) at (106,-28) {视图A\\$\bm{v}$};
  \node[data,minimum width=24mm] (viewb) at (149,-28) {视图B\\$\bm{v}'$};
  \node[model,minimum width=24mm] (enca) at (106,-48) {编码分支\\$\bm{z}$};
  \node[model,minimum width=24mm] (encb) at (149,-48) {编码分支\\$\bm{z}'$};
  \node[loss,minimum width=22mm] (match) at (127.5,-66) {匹配目标};
  \coordinate (view-split) at (127.5,-18);
  \draw[NNWire,line width=1.2pt] (raw.south) -- (view-split);
  \draw[flow] (view-split) -| (viewa.north);
  \draw[flow] (view-split) -| (viewb.north);
  \draw[flow] (viewa.south) -- (enca.north);
  \draw[flow] (viewb.south) -- (encb.north);
  \draw[flow] (enca.south) |- (match.west);
  \draw[flow] (encb.south) |- (match.east);

  \node[annotation,anchor=north west,text width=166mm,align=left] at (1,-80) {右侧表示在编码之后产生，随模型参数和训练步变化；本图只比较目标来源，不预设教师或停止梯度。};
\end{tikzpicture}

  \caption{两种目标在何处产生。左侧从完整观测分出答案，只让可见部分进入预测模型；右侧从对应视图分别计算表示，再交给匹配目标。右侧画出的是前向计算，不表示已经指定独立教师或停止梯度。}
  \label{fig:data-derived-knowledge-target-sources}
\end{figure*}
\FloatBarrier

\section{基于观测内容的知识获取}
\label{sec:data-derived-knowledge-observation}

第一节沿三条路线中最直观的一条展开：直接拿\emph{观测自身}的一部分当答案。做法听起来近乎“作弊”，但关键在\emph{时空隔离}：对一张还没有物种标签的照片，我们先\emph{保留}花瓣区域的原始像素，再把这片区域从喂给模型的输入中\emph{遮住}。于是形成一种微妙的局面——模型看不到被遮住的内容（当然无从知道答案），但\emph{学习算法手里握着遮之前的完整原图}，可以拿它核对模型猜得对不对。出题和评分用的是同一份数据的不同部分，不需要任何外部标签。

凡是这样\emph{从已采集的观测内容中取出来}当作预测对象的东西，统称为\term{观测目标}（observation target）：被遮住的像素是最典型的，文本里被盖住的词元、要推测的属性、时序数据里尚未发生的未来片段，都能扮演同一角色。要注意它与5.2节“向量目标”的本质区别：观测目标是\emph{固定的原始内容}，不需要那个正在训练的编码器去生成。当然，自由也是有边界的——“让哪部分当答案、哪部分允许被读、答案放在什么方向”并不是任意的，它由遮蔽位置、预测方向和预处理规则\emph{明确决定}，这些规则设计得好不好，直接决定模型能否真正学到有用的东西。

这类方法最大的风险是\emph{答案泄漏}：本应被藏起来的目标，由于切分或预处理的疏忽，其实以某种形式混进了模型可见的输入，模型不需要真正理解就能“抄到”答案。为了能系统地排查泄漏，必须把一条观测在一次训练中扮演的\emph{几个角色显式分开}。

形式化地，记完整观测为$\bm{x}=(x_j)_{j\in I}$，其中集合$I$枚举了所有位置（对文本是词元序号、对图像是像素块、对结构化数据是字段），$x_j$是该位置的内容。构造规则从$I$中挑出两个不重叠的子集：\emph{可见集合}$V\subset I$（允许模型读）和\emph{目标集合}$M\subseteq I$（要模型预测）；通常严格限制处理后的输入不能以任何形式直接泄露$\bm{x}_M$的原值。模型只读取经过处理的可见部分$\widetilde{\bm{x}}_V$，去预测目标$\widehat{\bm{x}}_M$（或其条件分布）。

工程上务必把\emph{四样东西分别保存}：完整原始观测、处理后的实际输入、被保留下来的目标、以及在哪些位置计损失。只有这样，才能一步步核查“某个预处理算子有没有偷偷把答案带回输入”。反之，如果图省事把它们混成一份不区分角色的“训练样本”，就彻底失去了这条防泄漏的检查路径——模型损失降得再漂亮，你也无法判断它是真学会了恢复，还是只是在抄答案。

根据“从观测里取什么当目标”，常见有四种安排，表~\ref{tab:data-derived-knowledge-observation-tasks}用三个统一的问题来比较它们，读者读表时可带着这三个问题：\emph{模型的输入被改成了什么样}（藏起？加噪？只给前文？只给历史？）、\emph{答案从原始观测的哪里保留}、以及\emph{误差在哪些位置计算}。

这四个名字——掩码恢复、去噪恢复、上下文预测、时序预测——描述的是\emph{构造安排}，并不是四个互不相容的任务：例如同一个“未来片段”，既可以直接当下文预测的目标，也可以先对它加噪、再当去噪恢复的目标，两种构造可以叠加。因此真正决定一种安排属于哪一类的，是\emph{可见观测与目标范围的具体划定}，而不是最后给损失起了什么名字——两种方法即便都叫“MSE”或“交叉熵”，只要可见/目标范围不同，学的东西就可能完全不同。

\begin{table*}[htbp]
  \centering
  \renewcommand{\arraystretch}{1.15}
  \begin{tabularx}{\linewidth}{@{}>{\RaggedRight\setlength{\parindent}{0pt}\arraybackslash}p{22mm}*{3}{>{\RaggedRight\setlength{\parindent}{0pt}\arraybackslash}X}@{}}
    \toprule
    构造方式 & 模型可见部分 & 保留目标 & 损失位置 \\
    \midrule
    掩码恢复 & 隐藏选中图像块后的图像 & 被遮蔽块的原像素 & 通常只在掩码位置 \\
    去噪恢复 & 加入扰动后的图像 & 扰动前的原始观测 & 受扰位置或全部位置 \\
    上下文预测 & 前缀“花朵逐渐” & 后续词元“开放” & 约定的预测位置 \\
    时序预测 & 判断时点以前的历史窗口 & 后续帧或状态 & 一个未来点或一段轨迹 \\
    \bottomrule
  \end{tabularx}
  \caption{按可见观测、保留答案和计分位置比较观测目标。未来片段虽已进入训练样本，在真正预测的时点却不可读取；四种安排可组合，不能据表格分行认定它们互斥。}
  \label{tab:data-derived-knowledge-observation-tasks}
\end{table*}
\FloatBarrier

\subsection{基于遮蔽恢复的目标构造}
\label{sec:data-derived-knowledge-masking-denoising}

在进入两种具体构造前，先用一个操作层面的区分把它们分开。固定一张完整照片和一块需要恢复的区域，我们对这块区域可以做两种性质不同的破坏。一种是\emph{直接拿走}：该位置变成空白或一个“此处缺失”的占位符，模型面对的是\emph{内容缺失}。另一种是\emph{留下位置、却改掉数值}：格子还在、里面的颜色被换成错误值，模型面对的是\emph{内容受损}（信息在、但被污染）。

以缺失为主的构造称为\emph{遮蔽恢复}（5.1.1），以受损为主的构造用于\emph{去噪恢复}（5.1.2）。但要立刻提醒二者边界并不截然：把选中内容替换成一个特殊的缺失符号，本身也可看作一种扰动；所以这里区分的是两种\emph{常用的输入安排和损失位置}，方便分别讲解，而不是声称它们是两个互斥、有严格理论分界的类别——实际方法完全可能混合使用。

图~\ref{fig:data-derived-knowledge-masking-denoising}用同一个$3\times3$的教学网格把两种破坏并排画出来，便于直接对照。上行是遮蔽：两个格子的原值被叉号\emph{整体替换}（信息缺失）；下行是加噪：同样这两个位置\emph{仍有数值}、但已被改动成受损值（信息被污染）。

看图时抓住一个共同点和一个变量。共同点是：\emph{两行的保留目标完全相同}——无论输入被缺失还是被污染，用来核对的都是那两个\emph{未经修改的原始值}，目标不随输入一起改。变量是\emph{计分范围}：图中赭色边框特意指定\emph{本例只核对这两处}的预测；但一个真实的去噪任务也可能选择在\emph{全部}位置计分，此时那些没受损的内容同样参与梯度。计分范围一旦改变，目标的含义就变了，必须在目标定义中如实反映，不能两种做法共用一个损失名就当作一回事。

\begin{figure*}[htbp]
  \centering
  \begin{tikzpicture}[
  x=1mm,y=1mm,
  font=\figurefont\scriptsize,
  flow/.style={knowledge arrow,draw=NNWire,line width=1.1pt},
  grid/.style={draw=BookBlue,line width=.85pt},
  cell/.style={draw=BookRule,line width=.5pt,minimum size=5mm,inner sep=0pt},
  heading/.style={font=\figurefont\bfseries\footnotesize,text=BookInk},
  stage/.style={font=\figurefont\footnotesize,text=BookMuted}
]
  \node[heading,anchor=east] at (13,1) {掩码};
  \node[heading,anchor=east] at (13,-34) {去噪};
  \foreach \x/\name in {23/原始观测,64/处理后预测输入,105/模型预测,146/保留目标}{
    \node[stage] at (\x,13) {\name};
  }

  \begin{scope}[shift={(15,8)}]
    \foreach \r/\c/\v in {0/0/.2,0/1/.4,0/2/.6,1/0/.3,1/1/.8,1/2/.7,2/0/.1,2/1/.5,2/2/.9}{
      \node[cell] at ({\c*5},{-\r*5}) {$\v$};
    }
  \end{scope}
  \begin{scope}[shift={(56,8)}]
    \foreach \r/\c/\v in {0/0/.2,0/1/.4,0/2/{\times},1/0/.3,1/1/{\times},1/2/.7,2/0/.1,2/1/.5,2/2/.9}{
      \node[cell] at ({\c*5},{-\r*5}) {$\v$};
    }
    \fill[FigureOutput!18] (7.5,-2.5) rectangle (12.5,2.5);
    \fill[FigureOutput!18] (2.5,-7.5) rectangle (7.5,-2.5);
    \node[cell] at (10,0) {$\times$};
    \node[cell] at (5,-5) {$\times$};
  \end{scope}
  \begin{scope}[shift={(97,8)}]
    \foreach \r/\c/\v in {0/0/.2,0/1/.4,0/2/.55,1/0/.3,1/1/.70,1/2/.7,2/0/.1,2/1/.5,2/2/.9}{
      \node[cell] at ({\c*5},{-\r*5}) {$\v$};
    }
    \draw[FigureLoss,line width=1.1pt] (7.5,-2.5) rectangle (12.5,2.5);
    \draw[FigureLoss,line width=1.1pt] (2.5,-7.5) rectangle (7.5,-2.5);
  \end{scope}
  \begin{scope}[shift={(138,8)}]
    \foreach \r/\c/\v in {0/0/.2,0/1/.4,0/2/.6,1/0/.3,1/1/.8,1/2/.7,2/0/.1,2/1/.5,2/2/.9}{
      \node[cell] at ({\c*5},{-\r*5}) {$\v$};
    }
    \draw[FigureLoss,line width=1.1pt] (7.5,-2.5) rectangle (12.5,2.5);
    \draw[FigureLoss,line width=1.1pt] (2.5,-7.5) rectangle (7.5,-2.5);
  \end{scope}
  \draw[flow] (31,3) -- (53,3);
  \draw[flow] (72,3) -- (94,3);
  \draw[flow] (113,3) -- node[above,text=FigureLoss] {仅比较框内位置} (135,3);

  \begin{scope}[shift={(15,-27)}]
    \foreach \r/\c/\v in {0/0/.2,0/1/.4,0/2/.6,1/0/.3,1/1/.8,1/2/.7,2/0/.1,2/1/.5,2/2/.9}{
      \node[cell] at ({\c*5},{-\r*5}) {$\v$};
    }
  \end{scope}
  \begin{scope}[shift={(56,-27)}]
    \foreach \r/\c/\v in {0/0/.2,0/1/.4,0/2/.8,1/0/.3,1/1/.6,1/2/.7,2/0/.1,2/1/.5,2/2/.9}{
      \node[cell] at ({\c*5},{-\r*5}) {$\v$};
    }
    \fill[BookAmber!18] (7.5,-2.5) rectangle (12.5,2.5);
    \fill[BookAmber!18] (2.5,-7.5) rectangle (7.5,-2.5);
    \node[cell] at (10,0) {.8};
    \node[cell] at (5,-5) {.6};
  \end{scope}
  \begin{scope}[shift={(97,-27)}]
    \foreach \r/\c/\v in {0/0/.2,0/1/.4,0/2/.65,1/0/.3,1/1/.72,1/2/.7,2/0/.1,2/1/.5,2/2/.9}{
      \node[cell] at ({\c*5},{-\r*5}) {$\v$};
    }
    \draw[FigureLoss,line width=1.1pt] (7.5,-2.5) rectangle (12.5,2.5);
    \draw[FigureLoss,line width=1.1pt] (2.5,-7.5) rectangle (7.5,-2.5);
  \end{scope}
  \begin{scope}[shift={(138,-27)}]
    \foreach \r/\c/\v in {0/0/.2,0/1/.4,0/2/.6,1/0/.3,1/1/.8,1/2/.7,2/0/.1,2/1/.5,2/2/.9}{
      \node[cell] at ({\c*5},{-\r*5}) {$\v$};
    }
    \draw[FigureLoss,line width=1.1pt] (7.5,-2.5) rectangle (12.5,2.5);
    \draw[FigureLoss,line width=1.1pt] (2.5,-7.5) rectangle (7.5,-2.5);
  \end{scope}
  \draw[flow] (31,-32) -- (53,-32);
  \draw[flow] (72,-32) -- (94,-32);
  \draw[flow] (113,-32) -- node[above,text=FigureLoss] {此例比较受扰位置} (135,-32);

  \node[font=\figurefont\scriptsize,text=BookMuted,anchor=west] at (15,-48) {数值与扰动均为教学构造；保留目标是原始观测，不代表无采集噪声的真值。};
\end{tikzpicture}

  \caption{在同一网格上对照隐藏内容与修改内容。叉号标出掩码，浅赭色标出受扰位置，赭色边框指定本例计分范围；两行都用未修改的原始观测核对预测。数值与扰动仅为教学构造，原始观测还可能含有采集噪声，不代表无噪声真值。}
  \label{fig:data-derived-knowledge-masking-denoising}
\end{figure*}
\FloatBarrier

\subsubsection{基于掩码的内容恢复}
\label{sec:data-derived-knowledge-masked-recovery}

先看图的上行。两个格子的值被整体拿走，任务是让模型\emph{补出}它们，而不是从输入里复制（输入里根本没有）。用这种“目标位置被\emph{完全隐藏}”的纯粹情形来定义\term{掩码恢复}（masked reconstruction）：选定待恢复的位置集合$M$，模型只能依据其余可见位置$V=I\setminus M$来预测目标内容$\bm{x}_M$。

“隐藏的单位”是什么，会显著改变任务难度和所需线索，常见的隐藏单位包括：文本中的词元、图像中互不重叠的图像块（patch）、表格里的字段或单元格。单位大小直接决定\emph{恢复它需要看多远}：如果只遮住\emph{单个像素}，那它几乎总能靠相邻像素的平滑插值被轻易补出，模型不必理解任何语义就能拿低分；但如果遮住的是一大片\emph{连续区域}或一整个\emph{完整属性}（比如整朵花、整张叶片），近邻插值就不够了，模型必须利用更远处、甚至跨区域的上下文才能推断。这个“遮蔽单位/比例”的选择因此不是无关紧要的工程细节，而是决定模型被迫学到多深规律的关键旋钮。

要根据这些保留值更新模型，还需把每处恢复偏差汇成一个可计算的分数。对非空掩码集合$M$中的连续像素或像素块，用平方欧氏距离度量预测与原值的偏差，再对选中单位取平均；训练在所用网络结构允许的模型参数范围内减小这一分数。记$\bm{x}_j$为目标单位的像素向量、$\widehat{\bm{x}}_j$为对应预测，常见的掩码均方误差为
\begin{equation}
  \mathcal{L}_{\mathrm{mask}}
  =\frac{1}{|M|}\sum_{j\in M}\lVert \widehat{\bm{x}}_j-\bm{x}_j\rVert_2^2.
  \label{eq:data-derived-knowledge-masked-reconstruction}
\end{equation}
读懂这个损失要注意三点。其一，分母用$|M|$取平均，是为了让\emph{分数高低不直接随隐藏单位的数量变化}——隐藏10块和隐藏100块得到的分数才具有可比性，否则模型可以靠“少藏一点”单纯地把总误差做小。其二，这个平方误差是针对\emph{连续像素}的；文本词元是离散符号，没有“两个词元相差多少”的欧氏距离，因此通常改用\emph{真实词元的交叉熵}（在词表上的分类损失）来计分。

其三是一个容易被忽视的陷阱：如果把\emph{未遮蔽的可见位置}也放进平均误差里，模型就可以靠\emph{大量复制那些它本来就看得见的内容}来拉低总分——可见部分复制得再准也没有难度，却稀释了真正反映“恢复未知”能力的那部分误差。这样的目标当然\emph{可以}使用（某些场景确实希望输出整幅完整图像），但必须意识到它\emph{同时奖励了“复现已知内容”}；此时总分下降就\emph{不能再全部}归因于“恢复未知内容的能力提升”，解读指标时要把两部分区分开。

现在用图~\ref{fig:data-derived-knowledge-masking-denoising}上行复算这一计分。两个保留值为$(0.6,0.8)$，人为设定的预测为$(0.55,0.70)$；把它们代入式~\eqref{eq:data-derived-knowledge-masked-reconstruction}，得到
\[
  \mathcal{L}_{\mathrm{mask}}
  =\frac{(0.55-0.6)^2+(0.70-0.8)^2}{2}
  =0.00625.
\]
算出来的$\mathcal{L}_{\mathrm{mask}}=0.00625$需要正确定位它\emph{衡量了什么、没衡量什么}。它\emph{只}核对那两处被隐藏的值（图中$3\times3$网格共9格，其余7处可见值\emph{没有}参与平均），所以这个数字只描述“在这两个位置、对这批训练观测的恢复偏差”，不能推广成对整图恢复质量的评价。

更重要的是答案的\emph{身份}：模型被要求恢复的$0.6$、$0.8$，只是\emph{训练观测样本中原本记录的数值}，它们\emph{不是}经过额外、更精确测量所确认的“无噪声真值”。原始照片本身可能就带采集噪声、压缩伪影。这一点划出了观测目标的\emph{能力边界}：模型学会的是“复现这份观测本来的样子”，而不是“还原客观世界的真实值”；把恢复观测误当成还原真相，会高估这类目标的可靠性。

这套“遮住—恢复”思想在大规模图像上的一个代表性实现是\emph{掩码自编码器}（masked autoencoder, MAE）。它在工程上有一个醒目的非对称设计：\emph{编码器只处理可见块}（被遮住的大部分图像块根本不进编码器，从而大幅节省算力），再配一个相对轻量的\emph{解码器}，由它结合“掩码标记”（占位符）去补出缺失块的像素。原论文只在\emph{被遮蔽的块}上计算像素均方误差\scite{ka-he2022mae}。

这带来一个干净的互补关系：\emph{编码器读取的块}与\emph{损失核对的块}恰好不重叠——读的是可见部分、考的是遮住部分，从结构上迫使编码器不能靠复制输入得分。至于目标像素的具体形式，除了直接用原始像素值，也可以先对每个块做\emph{归一化}（去掉块间亮度差异）再当目标，MAE对此做过研究。归一化只是改变了数值的\emph{尺度}（让模型更关注图案结构而非平均亮度），但它\emph{没有改变“答案来自原始图像”这一事实}——目标仍然是观测内容，而非外部标签或动态向量。

遮蔽的“比例”和“形状”是一对必须联调的超参数，因为改变它们\emph{改变的是任务本身}，而不只是误差项的多少，这里有两个相反的极端。

遮蔽\emph{比例过低}时，剩下的可见信息太多、太局部，任务可能退化成“插值”——缺一块颜色，看一眼周围就补得出，模型无需理解花的整体结构就能拿低分，学不到深层特征。遮蔽\emph{比例过高}时则走向另一极端：可见信息少到\emph{不足以唯一确定}被遮内容，同一份可见输入可能\emph{同时支持多个合理答案}（这块被遮的花瓣可能是红的、也可能是粉的）；此时用平方误差训练一个\emph{单点预测}，模型的最优倾向是输出这些可能答案的\emph{平均值}，结果往往是一片模糊、什么都不像的“安全色块”。

除了比例，遮蔽的\emph{位置规则}也留下不同的上下文：随机遮蔽（各处零散挖洞）、连续片段遮蔽（挖掉一大块）、结构化遮蔽（按物体或语义区域挖）迫使模型利用的上下文范围各不相同。因此评估时要\emph{同时}考察比例与位置规则，并且牢记一个反直觉的结论：让恢复任务\emph{更容易}（少遮一点、遮得更局部）通常能\emph{降低训练损失}，却\emph{未必}让编码器保留更多对识别物种有用的鉴别特征——好恢复的任务可能恰恰太浅。

这里要补一个非常关键、却常被新手忽略的泄漏途径：\emph{只在“最终喂给模型的那张图”上画掩码，并不足以真正隔离答案}。因为在画掩码\emph{之前}，数据往往已经流过若干预处理算子，而它们可能\emph{早就读过}目标位置。

举例来说，如果在遮蔽\emph{之前}先做了平滑或插值，那么遮住位置的值已经通过邻域运算“扩散”到了周围像素里，模型即使只看掩码后的图，也能从邻居数值把原值反推出来；更隐蔽的是\emph{全图归一化}：若归一化所用的均值、方差是用\emph{包含目标位置的完整图}算出来的，那么这一个统计量里就\emph{编码了被隐藏内容}的信息（知道总和与其余值，就能算出缺失值）。

因此防泄漏不能只看最终输入，而要\emph{逐个核对每个预处理算子读取了哪些位置}：在“输入与目标已经拆分”之后，该算子是否仍能访问到$M$中的内容？需要强调，画上一个掩码占位符\emph{本身}并不能撤销更早步骤已经造成的泄漏——信息一旦在遮蔽前的平滑、插值或统计中泄露出去，事后再怎么遮盖都无济于事，必须从算子执行的先后顺序上保证“目标内容只在允许的环节可见”。

最后指出一个容易混淆的“升级版”做法，它将把我们带向5.2节。掩码操作\emph{可以照用}，但要模型预测的\emph{答案}不再是被遮住的原始像素，而是“另一个编码器对那块内容算出的\emph{特征向量}”。一旦做这个替换，\emph{目标的来源就变了}：答案不再是训练前固定的观测原值，而是某个编码器的\emph{输出}，于是必须追问三个新问题——\emph{那个编码器是固定的还是随训练更新的？它接不接收当前这条损失的梯度？它和正在训练的网络是不是同一个？}

在这些问题没有澄清之前，就\emph{不能}再从“我做了遮蔽”这一形式，想当然地推出“模型在预测观测原值”。这正是5.1（观测目标）与5.2（关系/表示目标）的分水岭：\emph{随训练而变化的表示目标}，由第~\ref{sec:data-derived-knowledge-training-objectives}节接着分析。

\subsubsection{基于扰动的去噪恢复}
\label{sec:data-derived-knowledge-denoising-recovery}

再看图的下行：位置一个都没少，只是那两处的原值被\emph{改成了错误的受损数值}。我们仍然手握\emph{未改动的原始观测}，用它核对模型能否把受损内容\emph{还原回去}，这就是\term{去噪恢复}（denoising reconstruction）。

把它形式化：给定干净观测$\bm{x}$，先按照一个\emph{人为指定的扰动分布}$q(\widetilde{\bm{x}}\mid\bm{x})$随机生成一个受损版本$\widetilde{\bm{x}}$作为输入，再让模型预测扰动\emph{之前}的干净内容。常见的人为损坏手段包括：叠加加性噪声（图像加高斯噪声、音频加背景声）、随机替换词元（把句子里某些词换成随机词）、局部打乱顺序、或随机去掉颜色/遮挡局部。

必须澄清一个概念边界，以免夸大方法的能力：这里要撤销的，是\emph{我们在构造训练数据这一步主动施加的损坏}（我们清楚知道噪声长什么样、怎么加的），并\emph{不是}宣称模型已经学会消除“真实采集过程中产生的物理噪声”——后者的形成机制复杂且未知，用一种教学噪声训出的去噪器未必能迁移过去。

与“加一次噪、恢复一次”相比，真实训练要考虑一个事实：\emph{同一个原始观测可以随机生成许多不同的损坏版本}。只在某一个损坏副本上算一次误差，并不能说明模型对\emph{整套扰动规则}（噪声可强可弱、可落在不同位置）都适应良好；这次碰巧容易，不代表换个噪声模式也行。

因此损失要做\emph{两层平均}。先在\emph{一次}损坏内部，对所有已知受扰位置$R$的恢复偏差取平均；再在\emph{所有}$q$可能生成的损坏版本上取期望，让训练去减小“在这些可能损坏上的\emph{平均}误差”，而不是过拟合到某一个具体副本。设非空的受扰位置集合为$R\subseteq I$，$f_{\bm{\theta}}$是恢复模型，$\ell$按目标是连续像素还是离散词元选择相应的计分方式；固定干净观测$\bm{x}$时，去噪损失的一般形式为
\begin{equation}
  \mathcal{L}_{\mathrm{denoise}}
  =\mathbb{E}_{\widetilde{\bm{x}}\sim q(\cdot\mid\bm{x})}
   \left[\frac{1}{|R|}\sum_{j\in R}\ell\!\left(f_{\bm{\theta}}(\widetilde{\bm{x}})_j,x_j\right)\right].
  \label{eq:data-derived-knowledge-denoising}
\end{equation}
这个损失默认在受扰位置集合$R$上计分，回答的是“\emph{被弄脏的地方}恢复得如何”。计分范围是一个需要明确的设计选择：如果把平均范围从$R$扩大到整个$I$（所有位置），就额外要求模型对\emph{未受扰}内容也保持忠实输出，于是梯度的构成发生了实质变化——一部分梯度来自“修复”，另一部分来自“别动好的部分”，两者权衡不同。

由此得到一条报告规范：汇报去噪结果时，\emph{必须同时说明扰动分布（噪声类型、强度）和损失计分范围}，缺一不可。只写“用了去噪损失/重建损失”这样一个名称，别人无法确定模型实际被训练去做什么——同样叫“去噪”，只在受扰位置计分和在全图计分，学到的行为可能差别很大。

去噪能恢复到什么程度，\emph{根本上受限于输入里还剩多少信息}，优化再充分也无法无中生有，这里有两种对照情形。

当扰动很\emph{轻微}时，输入几乎保留了原值，模型用一个接近“恒等复制”（输入是什么就输出什么）的策略就能把误差压得很低——此时低误差主要说明噪声弱，不代表模型学会了复杂的还原。而当扰动\emph{极强}时，比如一整块花瓣被\emph{完全替换成随机纹理}，残留输入里可能\emph{不再包含}任何能确定那块颜色的信息；这时无论怎么训练，模型都\emph{无法从输入逻辑地推出}真实颜色，因为答案所依赖的信息已被破坏殆尽。

在后一种“信息不可得”的情形下，模型的输出由\emph{所用损失}决定它会给出什么形态的“最优猜测”：用均方误差时，最优解是各种合理颜色的\emph{条件均值}（往往是一块发灰、发糊的平均色）；用适合离散目标的损失时，可能给出\emph{众数}；更讲究的做法是直接\emph{预测一个分布}，把多种可能颜色都保留下来。关键认识是：这些都只是对不确定性的\emph{不同表达方式}，\emph{更充分的优化并不能找回输入中已不存在的答案}；把均值色当成“真实颜色”，是对去噪能力的过度解读。

损坏版本要反复抽取，工程上有两种数据组织方式，各有成本。一种是\emph{在线生成}：不提前存副本，每轮训练时拿原样本加上随机配置临时加噪；它只需保存一份原样本和随机配置、不占额外存储，代价是每轮都要重新执行加噪运算。另一种是\emph{预先离线生成}多个损坏副本存下来；它省掉了训练时重复生成的开销，却要占用存储，而且一旦固定，能用到的扰动就\emph{被限制在已保存的那几个副本之内}，多样性可能不足。

无论哪种方式，\emph{随机种子或扰动分布的配置都应作为构造配置的一部分记录在案}，否则实验无法复现；评价时也要说清：测试用的噪声与训练\emph{是同一种规则}（测的是拟合该噪声的能力），还是\emph{另一种未见的扰动}（更能反映泛化与鲁棒性）。最后守住前面的边界：在\emph{某一种特定教学噪声}上恢复得很准，只\emph{支持“在该噪声规则下有效”}这一条结论，不能外推成“模型能修复真实世界的采集误差”。

本节最后要切断一个最常见的错误因果链：\emph{去噪/重建损失下降，并不自动意味着下游的花卉识别变得更好}。这两件事衡量的对象不同，甚至可能背道而驰。一个模型可能非常擅长撤销\emph{某种固定的压缩伪影}（在这类噪声上恢复误差很低），但它在这个过程中学到的只是“如何反卷积”，并没有保留\emph{花瓣形态这类对鉴别物种有用的特征}。反过来，像MAE这类编码器，它对缺失像素的\emph{恢复结果在像素意义上可能相当模糊、并不唯一}（因为缺失块本来就有多种合理可能），但它\emph{中间学到的表示}却可能\emph{非常适合}下游识别——像素恢复得“不像原图”，特征反而好用。

因此必须\emph{分开取得两类证据、用两套评价回答两个问题}：\emph{恢复质量评价}把输出与原始观测核对，回答“你多大程度复原了输入”；\emph{下游任务评价}冻结编码器、在真正的类别或属性任务上考核，回答“你学到的表示对识别有没有用”。一个低重建误差不能代替下游指标，反之亦然；只有两者都看过，才能判断这个观测目标到底带来了什么。

\subsection{基于上下文预测的目标构造}
\label{sec:data-derived-knowledge-context-prediction}

掩码恢复在选位置时有一定任意性，而对\emph{有自然顺序的序列数据}（最典型是文本），其实存在一种更省事、且位置选择毫不随意的目标构造：直接利用序列\emph{本来的先后顺序}。比如句子“花朵逐渐开放”，我们把前缀“花朵逐渐”给模型看，把紧接其后的“开放”\emph{留作答案}。切分依据不是随机挑选，而是\emph{序列天然的时间/排列顺序}——后文天然就是前文要预测的对象。

凡是“根据约定好的前缀、邻域或其他可见内容去预测某段目标”的构造，统称为\term{上下文预测}（context prediction）。它和掩码恢复可能\emph{都在预测词元}，形式上很像，但两者是\emph{两个不同的设计选择}，关键差别在“允许读哪一侧、目标怎么切”：上下文预测（尤其自回归式的）通常只让模型读目标\emph{之前}的内容，而遮住句中某个词的掩码模型则往往能\emph{同时读左右两侧}。这个可见范围的差别，会直接改变任务难度和模型形态。

要判断“给定前缀，模型该多确信下一个词”，需要模型输出的不只是一个词、而是\emph{候选词表上的条件概率分布}。对序列$x_1,\ldots,x_T$，在第$t$个位置模型只允许读取历史$x_{<t}$，然后在所用词表的所有词元上分配概率：哪个词接在后面的概率最大。计分规则很自然——\emph{真实出现的那个词元}被模型赋予的概率越高，损失越低；用负对数（negative log-likelihood）来衡量，概率越接近1，$-\log$值越接近0。

把序列中所有约定位置的这种损失\emph{加起来}，就把“读懂这一整条序列”的要求汇总成了逐位置的负对数似然；序列开头$t=1$没有历史，此时前缀视为空、或用一个约定好的起始标记来表示。
\begin{equation}
  \mathcal{L}_{\mathrm{context}}
  =-\sum_{t=1}^{T}\log p_{\bm{\theta}}(x_t\mid x_{<t}).
  \label{eq:data-derived-knowledge-context-prediction}
\end{equation}
理解这个式子的关键，是看懂“为什么一个乘积变成了求和”。根据链式法则，整条序列的概率等于各位置条件概率\emph{相乘}：
\[
  p(x_1,\ldots,x_T)=\prod_{t=1}^{T}p_{\bm{\theta}}(x_t\mid x_{<t});
\]
对乘积取\emph{负对数}，$\log$把连乘变成连加、再取负，就得到式中“逐位置负对数\emph{求和}”的形式。这一自回归分解的概率论依据可参考\BookChapterRef{chap:pgm-static-directed}{卷六《模型：概率图模型》的“静态有向模型”章}。

从训练依据的角度看，短短一句“花朵逐渐开放”就能\emph{在多个位置}免费提供监督目标：在“花朵”之后预测“逐渐”、在“花朵逐渐”之后预测“开放”……每一个位置都拿\emph{原始观测里本来就有的词}当答案，无需任何外部标注。但这种“滑动生成位置”的做法是一把双刃剑：增加参与计分的位置确实\emph{扩大了监督覆盖}（一条句子贡献多个训练信号），却也\emph{同步增加了相应的预测和计分计算量}；而把越来越长的全部历史都编码进上下文的总成本，还取决于具体模型（注意力模型是长度的高阶增长），不能简单地只按“位置多了几个”来估算代价。

前缀式预测只能用\emph{左}侧历史，而另一种掩码语言模型允许被遮词同时看到\emph{左右}两侧上下文，二者的可见条件不同。BERT的掩码词元目标是后者的代表，但它有一个新手容易忽略的精细设计：\emph{并不是把所有目标词都换成掩码符}。原方法只随机选中约$15\%$的词元位置参与计分；而在这$15\%$当中，又分成三种处理：$80\%$替换成专用的掩码标记、$10\%$替换成一个\emph{随机词元}、另外$10\%$\emph{保持原词不变}\scite{devlin2019}。

为什么要掺入“随机词”和“原词”而不全用掩码符？这是为了防止模型走捷径：若被选中的位置永远是同一个特殊掩码符，模型只需学会“看见掩码符就使劲猜”，而不会真正利用上下文；掺入随机词迫使模型判断“这个词在这儿合不合理”，保留原词则让模型即便在没被掩码的位置也保持上下文表示。这个设计带来一个记录上的要求：\emph{“被选中计分的位置”与“真正被改动的位置”必须分开保存}，不能简单套用前面“目标位置全部隐藏”的$V=I\setminus M$安排——因为有的计分位置输入仍是原词。

还要分清两种性质不同的“答案泄漏”：一种是\emph{有意}在$10\%$的位置保留原词元（这是方法设计的一部分，模型在这些位置本就应当看到原词）；另一种是\emph{预处理违背了已声明的可见范围}、\emph{意外}把答案带回输入（这是工程错误）。两者表面相似、性质相反。最后强调：双向掩码预测与单向前缀预测即便\emph{答案恰好相同}（都是同一个词），它们的\emph{条件概率}也\emph{不同}——一个是$p(x_t\mid x_{<t})$、一个是$p(x_t\mid x_{\neq t})$，可见条件不一样，使用同一个交叉熵公式并不会让这两个概率变得可以互换。

式~\eqref{eq:data-derived-knowledge-context-prediction}在训练时有一个隐含的安排值得点明：喂给模型的前缀通常是\emph{真实的}$x_{<t}$，即每一步都用\emph{观测样本里实际的前文}，而不是用模型\emph{自己上一步生成的词}。这种训练方式称为\emph{教师强制}（teacher forcing）——相当于不管学生前面答得对不对，下一步都让他看着\emph{标准答案的前文}继续，训练因此稳定、易于并行。

但教师强制会造成训练与使用之间的落差。在\emph{真正连续生成}（推理）时，模型后面位置接到的是\emph{它自己之前预测出的词}，而不再是真实前缀；于是一旦\emph{前一步生成错误}，这个错误词就成了下一步的输入上下文，可能把后续预测也带偏，错误\emph{沿序列累积}（暴露偏差）。因此训练时那种“每一步都站在正确前文上”的\emph{单步误差}，不能直接等同于真实自由生成时的\emph{长序列误差}——前者每步前提都对，后者要为自己早先的错误买单。此外，用截断窗口训练会移除更早的观测，而位置、片段边界的不同切法也可能让训练时看到的上下文偏离真实部署时的上下文，这些都需要在设计时对齐。

上下文预测并不只属于文本，它的内核——“限制可见内容、保留一段观测当答案”——可以推广到其他模态，只要这些模态存在某种天然的相邻或对应结构。图像上可以用\emph{左侧区域预测右侧区域}（空间上的“下文”）；音频（语谱图）上可以用\emph{邻近频段预测被保留的片段}；在多模态情形，还可以用一个模态去预测\emph{同一事件中已经采集到的}另一个模态（用图像预测同期音频的大致内容，或反之）。

跨模态推广时要盯住\emph{什么变了、什么没变}。变的是输入与目标所属的\emph{观测字段类型}（像素、频段、来自另一传感器）；\emph{没变}的判断依据仍然是那句老话——\emph{答案能不能直接从已采集的观测样本里取出来}。只要目标是另一个模态里\emph{现成采集到的原始片段}，它就仍属于5.1的观测目标。但一旦目标变成“\emph{另一个模态经过编码器之后的向量}”，目标来源就转移到了5.2：此时必须追踪那个编码器是否固定、是否更新，不能再仅凭“两个模态原始配对”就断言目标是现成观测。

本节最后提出一个对自监督学习尤其要害的质疑：\emph{上下文预测得好，模型利用的就一定是我们想要的对象特征吗}？未必。可预测性可能来自大量\emph{与任务无关的稳定结构}：网页里固定的模板和边框、图像里周期性重复的背景、视频角落常驻的台标——仅凭这些“脚手架”，模型就可能相当准确地补全内容，却完全没有学到花本身的特征。

更隐蔽的泄漏来自\emph{预处理产物}：分词器留下的痕迹、缓存的统计特征、甚至裁剪坐标本身，都可能\emph{编码了答案}。比如如果不同类别的图像在采集时被系统自动存进了不同尺寸、或文件名/坐标带有类别规律，模型可能靠这些“元信息捷径”把预测分数做得很低，而与花卉形态毫无关系。

要排查模型\emph{到底在依赖什么}，可以做两类干预性检查：一是\emph{改变上下文范围}、或\emph{删除可疑的附属字段}（采集时间、设备编号、坐标等），观察预测能力是否崩塌，以此反推出哪些线索在起作用；二是用\emph{独立的下游任务}去检验这些被利用的规律\emph{是否真的有用}。把两方面结合，才能区分两种截然不同的成就：模型只是\emph{复现了观测序列中的表层规律}（可能是捷径），还是\emph{真的取得了预期任务所需的特征}（对识别有帮助）。

\subsection{基于时序预测的目标构造}
\label{sec:data-derived-knowledge-temporal-prediction}

把上下文预测从“语言顺序”推进到“\emph{物理时间}”，就得到时序预测。设想我们有一段记录植物生长的影像，希望用\emph{今天以及此前}的画面去预测\emph{几天后}植株的状态。

这里与普通语言模型有一个关键的额外约束，必须反复强调：在\emph{整理训练数据}时，我们当然能随意翻看整段影像、拿到后续帧；但这个任务要模拟的是\emph{真实站在今天做判断}的情境，而在今天，未来帧\emph{客观上尚不可得}。因此\term{时序预测}（temporal prediction）以一个\emph{真实判断时点}$t$为\emph{观测可用的硬边界}：只允许使用截至$t$的历史去预测$t$之后的观测。与文本中仅由排列决定的“先后”相比，时序预测多了“\emph{该观测在判断时点必须真的已经发生、可获得}”这一物理约束。未来帧虽然仍像5.1其他目标一样\emph{是现成的观测内容}（训练时用来核对），但在构造每一个训练样本时，它们必须被严格挡在输入之外——这是防止“用未来预测未来”的生命线。

光说“用历史预测未来”还\emph{不足以切出任何一个具体的训练样本}——“历史”是多长？从什么时候开始预测？预测多远、预测多久？这些不定清楚，就既无法构造样本，也无法比较误差。

因此要用三个\emph{窗口参数}把“哪段历史配哪段未来”明确下来，三者都以“采样了多少步”为单位：$L\geq1$是\emph{历史窗口长度}（往回看多少帧）；$h\geq1$是\emph{起始预测跨度}（跳过多少步之后才开始预测，$h=1$表示紧接当前）；$H\geq1$是\emph{目标轨迹长度}（一次预测未来多少步）。记号中的冒号表示\emph{包含两端}的连续索引。于是一个训练样本的输入与目标写为
\begin{equation}
  \bm{x}_{t-L+1:t}
  \quad\longrightarrow\quad
  \bm{y}_{t,h,H}
  =\bm{x}_{t+h:t+h+H-1}.
  \label{eq:data-derived-knowledge-temporal-windows}
\end{equation}
这个式子固定了两件事：模型\emph{允许读取的最后时点}（$t$）和\emph{目标从哪里开始}（$t+h$）。但“步”本身只是一个计数，\emph{必须结合采样间隔才有物理含义}——同样是“预测未来10步”，每分钟采样一帧对应未来10分钟，每天采样一帧却对应未来10天，这显然是\emph{两个完全不同难度的问题}。报告窗口参数时若不说明采样频率，“10步”这个数字几乎没有意义。

$H$的大小则改变预测的\emph{对象}：$H=1$时只核对未来\emph{一个时点}（单步预测）；$H>1$时要一次性核对\emph{一整段生长轨迹}，不仅输出维数增大、计算更重，而且轨迹内相邻时点的误差\emph{彼此相关}（同一株植物连续几帧的偏差不会独立），评估时不能把它们当成互不相关的多个独立样本简单平均。

窗口划定后，还要决定\emph{用什么形式表达未来的答案}，这取决于未来本身的确定性程度。

如果在给定历史下未来\emph{近似唯一}（短期内植物的下一帧状态基本确定），且实际使用中“预测偏差”的代价确实适合用平方误差衡量，那么输出一个\emph{点预测}（单一最佳估计值）是简单可用的选择。但现实中，\emph{同一段历史往往可能通向多条不同的未来轨迹}（之后天气、养护不同，植物会长成不同样子）；这时硬给一个点，模型多半又会输出多条可能轨迹的“平均”，得到一张谁也没发生过的中间状态。更合适的是预测一个\emph{条件分布}、若干个\emph{分位数}、或一个\emph{离散区间}，把“未来有多种可能”这层不确定性保留下来。

需要清醒的是：预测一个分布\emph{并没有让本来不可预测的未来变得可预测}，它只是诚实地\emph{拒绝把多种可能性压缩成一个均值}；而且无论采用哪种表达，\emph{评价规则都必须与最终的实际决策对应}——用于灌溉决策和用于估产的预测，对误差的容忍方式本就不同，不能用同一个分数一概而论。

预测跨度$h$（以及长度$H$）调大，意味着模型必须跨越\emph{更长一段自己看不到的过程}去推测未来，难度并非线性增加。\emph{短跨度}预测常常可以依赖很强的\emph{局部连续性}——下一帧与当前帧差别很小，外推一点点通常很稳。但\emph{长跨度}预测会同时面对三重困难：期间可能发生\emph{外部干预}（浇水、修剪、天气突变），使趋势断裂；逐步外推的\emph{误差会累积}，小偏差滚成大偏差；\emph{环境本身也可能缓慢变化}，使历史规律不再适用。

因此评估时应\emph{固定同一批样本范围}，对\emph{不同跨度分别报告误差}，而不能只给一个“所有跨度的平均值”——平均值可能被短跨度的高精度拉高，从而\emph{掩盖模型在远期其实已经失效}。当然也可以让模型\emph{同时}学习多个跨度以扩展目标覆盖（短、中、长期一起训），但代价是要承担更多的输出计算和目标存储，这是一种需要权衡的投入，不是免费的增益。

“判断时点边界”不能只停留在口头，而必须\emph{贯穿到数据切分和每一个预处理步骤}，否则未来信息会以各种隐蔽方式泄漏，时序任务就名存实亡。这里有两类最典型的泄漏。

一类是\emph{切分泄漏}：如果先在长序列上随机切出大量窗口、再随机划分训练/测试，由于相邻窗口高度重叠，\emph{同一段轨迹的帧很可能同时出现在训练集和测试集}（测试时模型早已见过几乎相同的未来帧）。正确的隔离单位应当是\emph{预期使用中的整体对象}——按时间段、按植株、按设备、或按完整轨迹成组划分，保证测试轨迹不被训练窗口覆盖。

另一类是\emph{预处理泄漏}，更隐蔽：若用\emph{整条序列}的均值做归一化，这个均值里就\emph{包含了未来}；若当前时刻的缺失值用\emph{未来}的观测去插值，更是直接把答案送进了输入。正确做法是：先按使用单位隔离数据，再\emph{只用训练阶段（或在该判断时点之前）允许获得的统计量}去拟合预处理（归一化参数、填充规则），并且\emph{逐个输入检查}：在它声称的判断时点上，其中每一项内容是否\emph{确实已经可以计算}。任何需要“未来数据”才能算出的特征，都属于越界。

即使严格守住了时间边界、也取得了某个预测误差，仍不能据这一个数字直接判定自监督目标的好坏，因为误差高低可能由多种与“表示质量”无关的原因造成，至少要把下面三个问题分开回答。

第一，\emph{未来在客观上能被现有输入确定到什么程度}？这是任务本身的“可预测性天花板”，与模型无关：一段长期生长过程如果随机性很强、受太多未知干预支配，那么再好的模型也会有较高误差，这个高误差\emph{不能}证明编码器没学到东西。第二，\emph{较低的误差是不是来自无关的重复结构}？比如影像里背景（花盆、架子）几乎不动，模型靠\emph{复制背景}就能拿到很低的整体预测误差，却没有捕捉植物状态本身的变化——这与5.1.3的“捷径”警告一脉相承。第三，\emph{学到的表示特征到底支不支持下游用途}？这只能靠独立下游任务检验。

因此合理的评估顺序是：先估计未来的可预测上限（问题一），再分析误差中有多少来自背景等非目标成分（问题二），最后在下游任务上确认表示是否有用（问题三）。\emph{单凭一个时序预测分数高或低，都不足以认定这个监督目标的质量}。

\section{基于样本关系的知识获取}
\label{sec:data-derived-knowledge-relations}

第一节让模型恢复观测\emph{内容}，这一节换一个角度：不再要求补出花瓣像素，而是利用样本\emph{彼此之间的关系}来学习。直觉来自一个朴素观察：同一张花卉照片的两个不同裁剪，拍的\emph{客观上是同一朵花}，那么无论取景、尺度如何变化，模型对它们的“本质描述”理应保持一致；反过来，若要让模型学会\emph{分辨}当前照片与众多照片中的哪一个配对，就得另外给一些\emph{不匹配}的候选作为参照、衬托。

把这种“哪些样本本质相同、哪些只是拿来对照”的可追溯说明，称为\term{样本配对关系}（sample pairing relation）。它由两种方向相反的信息组成：\emph{正对应}明确说出“我们\emph{希望}这两个样本在表示中共享什么”（同一对象的不变特征），\emph{负对照}则给出“当前这次判别\emph{允许}把谁与查询区分开”。

要特别注意关系与表示的\emph{产生时序}：配对关系本身（这两个裁剪来自同一照片）可以由变换规则、采集日志或采样算法在\emph{编码之前}就确定并保存下来，它不依赖模型；但真正参加比较的\emph{表示数值}（两张图各自的特征向量）要等\emph{当步模型前向计算之后}才出现。前者是事先可知的“关系事实”，后者是训练中动态算出的量，5.2节的方法正是把\emph{事先的关系}接到\emph{动态的表示}上去构造损失。

但仅仅\emph{知道}“这两个裁剪是一对”，还远远\emph{没有唯一规定}该怎么训练——同一个正对应关系，至少可以接出三种意图不同的目标，它们比较的量不同、产生的梯度也不同。其一，在一批候选中\emph{提高这一对的相对相似度}（让它比其他不配对的候选更像，即对比判别）；其二，让一路\emph{直接预测}另一路的特征向量（跨视图预测）；其三，\emph{同时调整两路}使它们相互接近，但额外约束这一整批表示不要塌成相同的常量（对称匹配）。

这是一个重要的方法论认识：\emph{关系相同，目标未必相同}。三种方法用的可能是同一批“同照片裁剪对”，但因为比较的对象和施加的约束不同，优化出的表示性质也不同。反过来也要警惕一个反向误区：\emph{不能因为某个损失公式里出现了余弦相似度}，就想当然地认为它用的裁剪关系\emph{一定可靠}——损失形式只说明“怎么比较”，裁剪是否真的保留了共同对象、负对照是否真的语义不同，这些\emph{关系依据}是否成立，必须在构造阶段独立核查，公式本身不提供这种保证。

基于以上区分，本节的讲解遵循一条清晰的三段式主线，建议读者据此组织阅读：\emph{先建立关系}（5.2.1：凭什么说两个样本是正对应或负对照），\emph{再把关系接到模型输出上形成训练目标}（5.2.2：用什么损失、梯度怎么走），\emph{最后检查跨训练步的更新与退化}（5.2.3：动态目标如何更新、表示会不会塌缩）。

这三步分别对应三个\emph{不同的对象}，也需要三种\emph{不同的证据}来检验，不能混为一谈：一是\emph{可事先保存的样本配对说明}（关系本身），要用构造证据去检查“关系假设成不成立”；二是\emph{当步才算出的目标值}（向量、相似度），要沿梯度路径检查“目标有没有按预期方式使用这层关系”；三是\emph{跨步更新的模型参数}，要在下游任务上检查“这样学出的表示差异到底有没有用”。请记住：\textbf{一次损失下降}只是一个优化现象，它\emph{无法同时}回答“关系对不对、梯度对不对、表示有没有用”这三个问题，三者必须分头验证。

\subsection{样本关系的构造}
\label{sec:data-derived-knowledge-relation-construction}

“这两个裁剪属于一对”听起来简单，但要把它\emph{可靠地}交给训练，记录里至少要包含五项信息：它们来自哪个\emph{原对象}、两侧各是什么\emph{观测}、凭什么判定配对（或排除）的\emph{依据}、这层关系的\emph{适用范围}、以及生成它们的\emph{构造配置}。

不同来源的关系还各有需要单独核查的“成立条件”。对来自同一张照片的两个裁剪，要检查它们是否\emph{都覆盖了主体}——一个裁剪若只截到背景，就谈不上与另一个共享花的特征；对同一时刻采集的音频和画面，要规定\emph{允许多大的同步误差}，超出就不再算同一事件。务必区分两件事：“它们出自同一来源文件”的\emph{来源标识}只能帮我们\emph{追溯}，并\emph{不能替代}上面这些成立条件的检查。另外，对那些一时\emph{无法确定}是否成立的关系，可以如实保留一个\emph{置信范围}，而不必强行记成非是即否的二元标签——硬把不确定关系定成“是”，会把噪声直接注入训练。

图~\ref{fig:data-derived-knowledge-relation-construction}特意把“\emph{生成视图}这个动作”与“\emph{判断关系}这个结论”用不同线型分开画，正是为了防止把二者混为一谈。图中\emph{蓝色箭头}只表示一个机械动作——从原图\emph{生成}出一个视图；\emph{青色实线}才表示我们\emph{预期两侧共享对象特征}的正对应关系。左侧\emph{灰色虚线}标出一个\emph{待检查}的负对照（同类另一株，到底该不该被推开还需判断），右侧\emph{赭色虚线}标出一个已知的教学错配（描述与照片对不上）。

沿着这些线型区别，读图时可以分别提出三个独立的问题：这个裁剪\emph{保留了主体}吗？这个候选\emph{允许被当作负对照排斥}吗？这段描述\emph{真的覆盖}了这张照片吗？核心要点在于：\emph{不能}因为一对样本是由程序自动生成的、或是在采集时天然配好的，就直接推出它们的\emph{语义关系}一定正确——生成和采集只是“产生了一对”，是否“真的相关”需要另行判断。

\begin{figure*}[htbp]
  \centering
  \begin{tikzpicture}[
  x=1mm,y=1mm,
  font=\figurefont\footnotesize,
  panel/.style={draw=BookRule,fill=NNPanel,line width=.85pt,rounded corners=1.5mm},
  item/.style={draw=FigureInput,fill=FigureInput!10,line width=1.1pt,text width=25mm,minimum height=11mm,align=center,inner sep=1.5mm},
  textitem/.style={draw=FigureOutput,fill=FigureOutput!8,line width=1.1pt,text width=28mm,minimum height=11mm,align=center,inner sep=1.5mm},
  generate/.style={knowledge arrow,draw=FigureInput,line width=1.1pt,rounded corners=1mm},
  positive/.style={draw=FigureModel,line width=1.4pt},
  negative/.style={draw=BookMuted,line width=.9pt,dashed,rounded corners=1mm},
  warning/.style={font=\figurefont\scriptsize,text=FigureLoss,align=center}
]
  \draw[panel] (1,-77) rectangle (82,12);
  \draw[panel] (87,-77) rectangle (168,12);
  \node[font=\figurefont\bfseries\footnotesize,text=BookInk] at (41.5,5) {(a) 观测变换建立正对应};
  \node[item] (original) at (20,-13) {原照片\\花朵与叶片};
  \node[item] (crop-a) at (61,-13) {裁剪A\\花朵};
  \node[item] (crop-b) at (61,-37) {裁剪B\\花朵与叶片};
  \node[item,draw=BookAmber,fill=BookAmber!10] (same-class) at (20,-63) {同类另一株\\未必是负语义};
  \coordinate (crop-split) at (40,-13);
  \draw[FigureInput,line width=1.1pt] (original.east) -- (crop-split);
  \draw[generate] (crop-split) -- (crop-a.west);
  \draw[generate] (crop-split) |- (crop-b.west);
  \draw[positive] (crop-a.south) -- node[right=1mm,text=FigureModel] {正对应} (crop-b.north);
  \draw[negative] (crop-b.south) |- (same-class.east);
  \node[font=\figurefont\scriptsize,text=BookMuted,anchor=south] at (49,-61) {候选负对照};
  \node[warning,text width=28mm] at (20,-37) {同源裁剪也可能\\丢失共同主体};

  \node[font=\figurefont\bfseries\footnotesize,text=BookInk] at (127.25,5) {(b) 采集日志建立自然配对};
  \node[item] (photo) at (104,-13) {花卉照片\\白色花瓣};
  \node[textitem] (caption) at (148,-13) {原有描述\\“白色花瓣，长叶”};
  \node[textitem,draw=BookAmber,fill=BookAmber!10] (wrong-caption) at (148,-41) {错配描述\\“黄色花瓣”};
  \draw[positive] (photo.east) -- (caption.west);
  \node[font=\figurefont\scriptsize,text=FigureModel,anchor=south] at (125,-11) {观测配对};
  \draw[negative,draw=FigureLoss] (photo.south) |- (wrong-caption.west);
  \node[warning,anchor=south] at (118,-39) {错误配对};
  \node[warning,text width=70mm] at (127.5,-64) {原有描述可能只覆盖部分内容；\\错配样本不能作为正对应。};

  \draw[generate] (3,-84) -- (16,-84);
  \node[anchor=west,text=BookMuted] at (18,-84) {生成视图};
  \draw[positive] (53,-84) -- (66,-84);
  \node[anchor=west,text=BookMuted] at (68,-84) {正对应};
  \draw[negative] (103,-84) -- (116,-84);
  \node[anchor=west,text=BookMuted] at (118,-84) {负对照或错配关系};
\end{tikzpicture}

  \caption{生成视图与确认关系是不同动作。左侧同源裁剪形成候选正对应，同类另一株还需判断是否允许作负对照；右侧已有图文配对与错配描述形成对照。连为负对照没有确认语义不同，已知错配则不能拿来提供正对应。}
  \label{fig:data-derived-knowledge-relation-construction}
\end{figure*}
\FloatBarrier

\subsubsection{基于数据变换的正对应关系构造}
\label{sec:data-derived-knowledge-transformation-positive}

获得正对应最省事的办法，是从同一个样本\emph{人工变换}出两个视图：同一张花卉照片随机裁剪两次，就无须任何逐对人工标注，凭空得到一对“理应相同”的视图。这一做法背后的核心依据可以浓缩成一句话：\emph{变换应当保留我们希望表示共享的内容，同时改变我们不希望表示依赖的因素}。比如随机裁剪、缩放保留了“这是同一朵花”，却改变了取景位置和尺度——我们正希望模型对后者不敏感。

但不同变换改变的内容截然不同，各自需要单独确定“保留什么”和“强度多大”。整体与局部视图改变覆盖范围，尺度变化改变大小，颜色扰动改变色调；换到其他模态，语音加噪改变背景声、文本改写改变措辞。每一种变换都在悄悄声明“这个因素不该影响表示”，因此每一种都要单独审视：这个声明对当前任务成立吗？

用几个具体反例说明“变换不能乱用”。裁剪方面：若两次裁剪\emph{都含有花朵}，要求它们表示一致，就能鼓励模型学会忽略取景差异，这是有益的；但若其中\emph{一路只剩背景}，两侧早已没有足够的共同内容，此时强行要求表示一致，等于教模型“背景就等于这朵花”，反而有害。

颜色方面更有讽刺意味：如果某些品种恰恰要\emph{靠花色}来区分（比如红色与白色品种），而我们又施加了强烈的颜色扰动、并要求扰动前后表示保持不变，就等于\emph{逼着模型忽略花色}——这直接与“按花色识别品种”的下游任务相冲突。同理，当方向或位置本身就是下游答案的一部分时（比如要判断花是朝上还是下垂、或某器官在左还是在右），旋转、翻转、大范围平移就\emph{不能}被无条件地当作“不改变表示”的变换使用，否则模型被训练去丢弃恰恰需要的信息。

这些反例共同揭示了一个关键认识：\emph{变换配置不是中性的工程细节，它实质上是一份声明}——声明了“哪些变化不应该影响表示”。选错了变换，就等于给模型灌输了错误的“不变性”。因此必须把变换的种类、组合顺序、强度、随机参数，以及每对视图的主体覆盖情况都保存下来；只有这样，才能用抽样可视化（肉眼看几十对视图）或自动规则，检查共同内容是否真的还在。

成本上也要算清楚：多生成一些视图固然能\emph{增加正对应的数量}，但每多一个视图就要多做一次编码，计算开销随之增长。所以评估这种方法的投入时，不能只算“省掉了多少人工关系标注”，还要把额外的编码计算一并计入。

还有一种常见扩展：不只是两个对称视图，而是\emph{一个整体视图配多个局部小视图}（多裁剪，multi-crop）。这时关系的\emph{范围}发生了微妙变化：每个局部只覆盖整体的\emph{一部分}，它们“来自同一张图像”这一事实可以继承，但\emph{不能}据此要求某个局部小块\emph{完整复制}整体的全部对象特征——一个只截到花瓣边缘的小块，信息量本就不足以还原整朵花。

局部到整体的匹配到底能不能带来好的表示，最终还依赖具体的多裁剪设计（局部取多小、取几个）和训练经验。仅仅“局部与整体来源一致”这一点，本身并\emph{不}提供方法一定成功的保证。

至此，本小节只完成了一件事：得到了\emph{正对应关系及其适用范围}。它还\emph{没有}规定接下来用“在候选中判别”还是“直接匹配”的方式去用这层关系，也\emph{没有}指定是否设立独立的目标分支——这些属于“怎么用关系”的问题，由第~\ref{sec:data-derived-knowledge-training-objectives}节把这里确定的关系接到具体的损失函数与参数更新路径上。

\subsubsection{基于自然配对的正对应关系构造}
\label{sec:data-derived-knowledge-natural-positive}

人工变换是“自己造”正对应，但很多时候数据\emph{本身就自带}现成的对应关系，无需再造。比如一张花卉照片如果已经附有一段说明文字，那么“这张图”和“这段描述”天然就是一对，可以直接复用这条关联，而不必从同一张图的像素里人为生成两个视图。

这种\emph{采集过程中本来就存在的对应}称为\term{自然配对}（natural pairing），典型例子包括：同一对象在不同时间的多次拍摄、同步录制的音频与视频、配套出现的图片与文字、相邻时刻的传感器读数。与人工变换相比，自然配对有两个显著不同：其一，两侧可以属于\emph{完全不同的观测模态}（图与文、音与画），不再局限于同一张图的两个裁剪；其二，判断“它们是一对”的依据变了——不再是“出自同一张图的变换”，而是采集时记录的\emph{对象标识、时间戳、事件编号或文档结构}。因此可靠性的检查重点也随之转移：要核查这些\emph{采集字段本身是否可靠}（时间戳准不准、图文对应有没有错位）。

自然配对的“天然”并不意味着“等价”，即使是正确配好的一对，两侧覆盖的内容也可能相差甚远。同一张花卉照片配的文字，可能\emph{描述整株形态}（覆盖广），也可能\emph{只提了一句花色}（覆盖窄），还可能是\emph{从另一张照片误配过来的}。前两种都是合法配对，但它们支持的“共同内容”宽窄不同——对整株描述可以要求较多特征一致，对单句花色只能要求颜色相关的特征一致；第三种则是\emph{关系错误}，根本不该作为正对应。

因此训练之前要明确每一对的\emph{对应范围}：是整例对应、局部对应，还是一对多（一张图配多段描述）；也要如实保留“某一侧缺失”和“多重匹配”的情况，而不是把它们丢掉或强行一一对应。最关键的一条：\emph{不能}把“手里有一份配对清单”统一地解读成“每一对的两侧在内容上逐项相同”——配对清单只说明它们相关，不说明它们在哪些方面、多大程度上相同。

% 完整作者和会议条目不能跨页续排。
\Needspace{65mm}
利用自然图文配对最有影响力的例子是CLIP：它把网络上海量已存在的图文配对当作正对应，再拿\emph{同一批次里的其他图文组合}作为参照（负对照），训练模型去分辨“哪一个图文组合才与真实采集的配对相符”\scite{radford2021}。凭借大量这样的配对，模型学会了把图像和文字映射到同一个可比较的空间。

但要准确理解CLIP“省掉了什么、没省掉什么”。它确实\emph{不需要}为目标分类任务\emph{重新逐图标注}类别，这是巨大的节省；但它所复用的那些描述文字，本身就是\emph{人类书写}的——人类监督并没有消失，只是换成了一种“顺带产生、无需专门组织”的形式。因此三种成本不能混为一谈：描述文本从哪来（创作成本）、图文关联是否核对过（质控成本）、以及为某个具体任务专门标注（任务标注成本）；说“无需人工标注”只回答了最后一项。

自然配对中的错误不会凭空消失，它们会\emph{一路带进后续的判别目标}，甚至制造矛盾。举两个典型情形：若多张照片\emph{复用了同一段模板化描述}（例如电商图都配“精美花卉”），那么在批内把其他图文组合当负对照时，这些本应匹配的图文反而被当成“不配对”推开，产生训练上的冲突；若音频轨\emph{落后于画面}，那么“同一时刻”采集的音和画，实际上可能来自两个不同事件，正对应本身就是错的。

时间偏移、描述不完整、重复样本、共享背景——这些问题都需要回到\emph{采集协议和配对清单}去排查根源。要特别破除一个直觉：\emph{用同样有缺陷的方式多采集一些样本，并不会自动纠正这种系统性错配}；数据越多，同类错误反而被复制得越多、越牢固。

最后算一笔实际投入账。自然配对的两侧（如图与文）往往需要\emph{各自独立的编码器}分别处理和存储，跨模态的开销不小。如实报告时应包括：每一侧的缺失率、配对如何核验、以及两种模态各自的处理开销，才能真实反映“复用现成关联”到底花了多少；单单一句“无需新标签”，只回答了“专门标注”这一种成本是否增加，远不是全部成本。

\subsubsection{基于样本采样的负对照关系构造}
\label{sec:data-derived-knowledge-negative-sampling}

正对应告诉模型“这两个应该像”，但对于5.2.2将要讲的对比式目标，光有正对应还不够。设想任务是“在一堆裁剪里，找出与当前照片配对的那一个”——要让这个“找出”有意义，必须先确定\emph{还允许拿哪些其他候选来比}。没有任何参照，“更像”就无从定义。

那些从当前小批量、历史缓存或检索集合中挑出、用来衬托正对应的候选，称为\term{负对照}（negative contrast）。它的角色是“参照物”：训练目标允许（并鼓励）查询样本相对于这些负对照，取得\emph{更高}的匹配分数。但要非常小心这个名字带来的误解：一个样本\emph{被选进负对照集合}，\emph{不等于}已经有人判断过“它与查询在语义上不同”——它只是被算法选来作对照，很可能其实与查询同类（下文会看到这正是主要风险之一）。

负对照怎么选，会显著影响学习效果，常见有两种策略，各有利弊。

\emph{随机抽取}：从数据中随便取一些样本作负对照，好处是覆盖面广、不带偏见；坏处是随机样本往往与查询\emph{差异太大}（查询是山茶，随机抽到的是一辆汽车的照片），模型一眼就能区分，这种“太容易”的对照产生的区分梯度很弱，学不到细致的差别。

\emph{检索难负例}：按模型\emph{当前的表示}去检索与查询\emph{最相似、最难区分}的样本作负对照，这类候选称为\term{难负例}（hard negative）——即模型此刻很难把它和查询区分开的样本。用难负例训练，能迫使模型关注细微的差别（比如区分山茶与茶梅），通常更高效。但必须警惕两个风险：其一，所谓“难”是相对\emph{当前表示}而言的，它\emph{并不证明}这个样本与查询真的语义不同——最难区分的样本，很可能本来就是同类，把它推开是错的；其二，如果反复检索、总是盯着同几个少数样本，还可能\emph{强化}当前表示已有的偏差，形成自我循环。

用花卉照片说明一个更根本的问题：\emph{一个样本该当正对应还是负对照，并没有固定答案，它取决于你要服务的任务}。考虑三类候选：同一植株的另一张照片、同品种的另一株、以及完全不同的植物。

如果任务是\emph{识别具体是哪一株}，那么同一植株的另一张照片是正对应，同品种另一株则是应当区分开的不同个体（可作负对照）。但如果任务换成\emph{识别是什么品种}，情况就反过来了：同品种的另一株\emph{本应与查询聚在一起}，此时若仍把它当负对照推开，就恰恰\emph{违背}了我们想要的聚合方向，模型被训练去区分本该视为同类的样本。

由此得到一条原则：必须\emph{先确定关系所服务的任务和范围}，再决定谁是正、谁是负。尤其要杜绝一种偷懒推理——“这个样本还\emph{没被标成正例}”，不能直接当作“它就\emph{确定是负例}”的依据；未标注不等于不相关。

即使没有专门写一个“负对照采样文件”，负对照的选择也\emph{已经隐含}在每个小批量是由哪些样本组成里了——同一批次的其他样本自动充当负对照。所以批次的组成方式会直接改变参照分布：批里若有\emph{重复样本}，相当于把同一个样本既当正又当负；\emph{类别失衡}会让多数类充斥负对照；\emph{同源样本}（来自同一原图的多个裁剪混在不同位置）也会制造冲突。若负对照来自\emph{缓存}，缓存里存的是\emph{旧模型参数}算出的表示，可能已经与当前模型偏离，参照不再准确。

还要纠正“负对照越多越好”的直觉。扩大负对照集合确实\emph{提高了区分要求}（要在更多候选中胜出），但同时也\emph{增加了混入错误负例的机会}（集合越大，越可能把同类误当负例），还带来更多的通信和存储开销。因此负对照集合\emph{并非无条件越大越有用}，需要在区分强度与错误负例风险之间权衡。

最后处理一个常被忽视的情形：一个查询可能同时有\emph{多个}合理的正对应（比如同一株植物有好几张不同角度的照片）。这时应当把它们\emph{显式地都记为正对应}，或者换用一种\emph{支持多个正例}的目标函数。否则，那些被遗漏、没标成正对应的同类样本，会作为普通候选进入损失的\emph{分母}，在训练中被当成负对照\emph{排斥}——模型被迫推开本该拉近的样本。

这一节可以用一句话总结关系与损失的分工：\emph{关系集合决定训练允许区分谁}（谁是正、谁是负），\emph{损失函数决定怎样施加这种区分}（用什么公式、多大力度）。前者的错误（把同类当负例）是“方向性”的错误，\emph{无法}靠后者的损失下降来排除——损失降得越低，只说明模型越坚决地执行了一套可能本身就错的关系。

\subsection{训练目标的构造}
\label{sec:data-derived-knowledge-training-objectives}

上一小节确定了“哪两路是正对应”，现在要回答“\emph{用什么计分规则}把这层关系转化为参数更新”。前面已经预告过三种思路，这里逐一展开：比较正对应相对其他候选的\emph{相对相似度}（对比判别，5.2.2.1）、让一路\emph{预测另一路的表示}（跨视图预测，5.2.2.2）、以及\emph{同时匹配两路}并约束整批表示的统计量（对称匹配，5.2.2.3）。

为了弄清每种方法训练时\emph{实际改变了什么}，讲解每个方法时都会逐条追踪四件事：喂进去的输入是什么、中间经过了哪些附加映射（投影头、预测头等）、最终比较的是什么量、梯度沿哪条路径回传。同时会反复强调一个实践要点：\emph{训练时用来计分的“训练头”，和最终交给下游的“表示特征”往往不是同一个东西}，要分开看待。

另外请读者先记住一个后文多次用到的细微区分：“\emph{交换两个视图再另算一项损失}”和“\emph{同一项损失同时向两路回传梯度}”是两回事。前者只是把输入角色对调后多算一次，每一项仍可能只更新一路；后者才是真正让两路在同一项里一起被推动。前者并不自动包含后者。

% 标题、表示定义及其边注在同一页内保留足够高度。
\Needspace{80mm}
\subsubsection{基于对比判别的目标构造}
\label{sec:data-derived-knowledge-contrastive}

先看最经典的一种。设想当前的查询是照片$A$的一个裁剪，候选池里既有它的\emph{另一个裁剪}（正对应），也有\emph{其他照片}的视图（负对照）。我们希望模型学到：在这一池候选里，已知的那个配对应当\emph{脱颖而出}、取得相对更高的相似度。\term{对比判别}（contrastive discrimination）就是把这个要求写成一个目标——让正对应相对于所有候选具有更高的相似度。注意它关心的是“\emph{相对排名}”，而非绝对的相似度数值。

SimCLR是这类方法的代表实现，其结构可以分三层理解。数据层：同一批取$n$张原始照片，每张随机生成两个视图，共得到$2n$个视图。编码层：编码器$f_{\bm{\theta}}$把每个视图编码成\emph{特征}$\bm{h}_i$。投影层：再由一个小的投影器$g_{\bm{\theta}}$把特征映射到一个\emph{专门用于训练计分}的空间，得到$\bm{z}_i=g_{\bm{\theta}}(\bm{h}_i)$，对比损失就在这个$\bm{z}$空间里计算。

这里有一个关键且反直觉的实践事实：训练完成后，交给下游任务使用的通常是\emph{投影之前}的特征$\bm{h}_i$，而不是用来算损失的$\bm{z}_i$\scite{chen2020}。原因在于投影层被训练得过度专注于对比任务、可能丢掉了一些对下游有用的信息，而投影之前的$\bm{h}$保留得更完整。因此，“投影空间里配对分数变好了”与“编码特征对下游有用”是\emph{两个不同的环节}，前者不能直接代表后者。

要比较“谁更像”，首先需要一个能给任意两个投影向量打分的统一尺子。这里选用\emph{余弦相似度}，它只比较两个向量的\emph{方向}是否一致，而\emph{不}受向量长度影响——这样模型就不能靠单纯“把向量变长”来作弊式地抬高相似度。以下假定批量$n\geq2$（至少要有负对照），且参与比较的向量都非零（否则分母为零、余弦无定义）：
\[
  \operatorname{sim}(\bm{z}_i,\bm{z}_k)
  =\frac{\langle\bm{z}_i,\bm{z}_k\rangle}
  {\lVert\bm{z}_i\rVert_2\lVert\bm{z}_k\rVert_2}.
\]
有了打分尺子，再把“正对应应脱颖而出”变成一个可优化的损失。思路是把它看作一个\emph{多选一的分类问题}：固定一个\emph{锚点}，即当前被查询的视图$i$；它的候选是本批中\emph{除自己以外}的其余$2n-1$个视图，其中恰好有一个是它的正对应$j$，其余都是负对照。我们要让模型从这$2n-1$个候选里“选中”正对应$j$。

具体做法是：把每个候选的相似度先指数化，再在所有候选上归一化，得到“选中每个候选的概率”（这正是softmax的形式）；然后用负对数惩罚\emph{正对应被选中的概率太低}——正确配对的概率越接近1，损失越小。式中的\emph{温度}$\tau>0$是一个关键旋钮，控制相似度之间的差异在归一化时被\emph{放大}多少：温度越低，微小的相似度差别就会被放得越大。SimCLR使用的这种单向“归一化温度交叉熵”损失为
\begin{equation}
  \ell_{i,j}
  =-\log
  \frac{\exp\!\left(\operatorname{sim}(\bm{z}_i,\bm{z}_j)/\tau\right)}
  {\displaystyle\sum_{k=1}^{2n}\mathbb{I}[k\neq i]\exp\!\left(\operatorname{sim}(\bm{z}_i,\bm{z}_k)/\tau\right)},
  \qquad \tau>0.
  \label{eq:data-derived-knowledge-simclr}
\end{equation}
读懂分母是关键：指示函数$\mathbb{I}[k\neq i]$\emph{只排除锚点自身}，因此分母同时包含正对应$j$和其余$2n-2$个负对照——正对应要在“与负对照的竞争”中胜出，才能让分子占比大。完整的训练目标让\emph{每一个}视图都轮流当一次锚点，把这$2n$个损失取平均；梯度最终回传到\emph{共享的}编码器与投影器。

温度的作用最容易被误解，这里专门澄清。在其余相似度固定的情况下，\emph{降低温度}会\emph{提高那些相似度较高的负例相对于相似度较低负例的权重}——也就是让模型更关注“难”的负例。但这只是\emph{相对}权重的变化。某一个负例实际受到的\emph{绝对梯度}大小，还取决于正对应的得分以及整组的归一化概率。一个反直觉的情形是：当正对应\emph{已经明显胜出}时，降低温度会让它的概率迅速逼近1，整个判别趋于“饱和”，此时负例获得的梯度\emph{反而变小}。所以“低温=更强调难负例”只是相对重加权的意思，\emph{不能}理解为“低温下每个难负例都被推得更狠”。下面的数值例子会具体演示这一点。

用两个原样本可以直接检查候选范围。教学批次取$n=2$，按$A_1,A_2,B_1,B_2$排列四个视图，配对为$(A_1,A_2)$与$(B_1,B_2)$；对角线是自身比较，配对位置是正对应，其余位置在本例中作负对照。人为设定余弦相似度矩阵为
\[
  \bm{S}=
  \begin{bmatrix}
  1&0.8&0.1&0\\
  0.8&1&0&-0.1\\
  0.1&0&1&0.7\\
  0&-0.1&0.7&1
  \end{bmatrix}.
\]
用这个矩阵亲手算一遍锚点$A_1$的损失。取温度$\tau=0.5$，锚点$A_1$只与其余三个视图$A_2,B_1,B_2$比较（排除自身那一项“1”）。从矩阵第一行读出相似度：与$A_2$为$0.8$（正对应），与$B_1$为$0.1$，与$B_2$为$0$。除以温度后分别是$1.6$、$0.2$、$0$，于是正对应被“选中”的概率为
\[
  \frac{\exp(1.6)}{\exp(1.6)+\exp(0.2)+\exp(0)}\approx0.690,
\]
对应的负对数损失为$\ell_{A_1,A_2}=-\log0.690\approx0.371$。

要正确解读这个$0.690$：它衡量的是“\emph{在当前这三个候选之中}，模型多大程度上支持$A_2$是配对”，是一个相对于候选池的配对支持度。它\emph{完全没有}涉及花卉类别的定义，所以\emph{绝不是}“$A_1$与$A_2$属于同一类别的概率”——换一个候选池，这个数就会变，与它们是不是同种花无关。

沿用这三个分数，可以用数字验证前面关于温度的“反直觉”结论。一个有用的事实是：固定其余相似度时，单向损失对\emph{某个负例}相似度的导数，等于该候选的归一化概率再除以$\tau$——负例被“选中”的概率越大，它受到的推力越大。

现在把温度从$\tau=0.5$降到$\tau=0.1$。一方面，正对应的概率从约$0.690$\emph{升高}到约$0.999$，几乎被确定选中；另一方面，在两个负例之间，相似度为$0.1$的那个相对于相似度为$0$的那个，占比确实\emph{更高了}（低温确实放大了“难”负例的相对权重）。但关键是看\emph{绝对}梯度：相似度$0.1$那个负例的导数，从约$0.3405$\emph{骤降}到约$0.00911$。原因在于正对应本就已经明显高于两个负例，低温把整次判别推到了接近饱和的状态，所有负例的概率都被压得极小。这个例子清楚说明：\emph{负例之间的概率比变大了，但绝对梯度并没有随之变大}，“低温推难负例更狠”这一直觉在饱和时并不成立。

这个算例还揭示了对比损失的一个重要特性：它\emph{依赖于小批量的组成}。即使某对正对应自身的相似度丝毫未变，只要批里的\emph{其他}样本（负对照）换了，分母就变，损失也跟着变——损失值并不只由那一对决定。

批量大小因此是一个需要权衡的设计选择。更大的批量提供更多负对照、区分要求更高，但需要计算的两两相似度也更多，而在分布式训练中还要在设备间\emph{交换表示}，通信开销显著；反之，若批量太小、候选缺乏足够区分力，对比目标可能提供的学习信号偏弱。用缓存或检索的方式可以\emph{扩大}参照范围而不必增大批量，但要额外检查两个问题：缓存里是\emph{旧参数算出的表示}（可能已过时）、以及检索带来的\emph{采样偏差}。

对比判别有一个与生俱来、需要正视的局限，它和5.2.1中“负对照未必语义不同”的警告直接相关。SimCLR本质上在做“\emph{实例区分}”：把每张照片当成一个独立的类，要求它只和自己的变换相似、与其他所有照片都区分开。可如果批里恰好有一张\emph{同品种}的花卉照片，它也会作为负对照进入式~\eqref{eq:data-derived-knowledge-simclr}的分母，被模型\emph{推开}。

这样一来，实例区分就可能与“物种识别”的目标\emph{背道而驰}——后者恰恰希望同品种的照片\emph{聚在一起}。缓解手段包括：使用支持多正例的目标、对数据去重、以及用额外信息过滤掉明显同类的负对照。但必须看到，这些手段\emph{都依赖额外的配对依据}（得先知道谁和谁同类）；损失函数\emph{本身}并不知道哪一种语义（实例级还是品种级）应当被保留，它只会执行你给它的那套正负关系。

最后用一个极端情形，检验对比目标是否真的“奖励样本之间的差异”。这个情形就是后文5.2.3要专门讨论的\emph{表示坍塌}：假设模型偷懒，把所有输入的投影都映射成\emph{同一个}非零向量$\bm{u}$。此时任意两个投影的余弦相似度都等于$1$，所有候选看起来一模一样、完全无法区分。把这种情形代入损失，正对应和每个负对照的指数项都相同，于是得到
\[
  \ell_{i,j}=\log(2n-1).
\]
这个结果说明坍塌状态\emph{在对比目标下是不划算的}：常数表示没有任何配对优势，损失停在$\log(2n-1)$。而只要模型和表示空间允许正对应比负对照更相似，就能取得更小的损失——前面那个例子正是如此，单项损失$0.371$远小于常数表示对应的$\log3\approx1.099$。从这个角度看，对比目标确实“偏好”样本之间有差异。

但必须立刻补上一个精细却重要的限定：\emph{“能取得更低的损失”}不等于\emph{“优化过程一定会离开坍塌状态”}。数学上，所有投影都相同时，余弦相似度对任一表示向量的梯度恰好为零——也就是说，\emph{精确的}常数表示是一个\emph{驻点}：在那里，一阶梯度完全推不动表示。如果优化恰好落在（或非常接近）这个点，它可能就停在那儿。另一个边界情形是：若批量$n=1$，候选里只剩唯一的正对应、没有任何负对照，损失恒为零，模型从对比中学不到任何东西。

综上，对比目标\emph{只有在存在有效负对照、且区分确实可实现的条件下}才偏好差异。它\emph{既不保证}从任意初始化出发都能避开坍塌，也\emph{不保证}学到的差异来自花朵本身——差异完全可能来自背景、拍摄设备甚至图片水印，只要这些线索能帮助区分实例，对比目标就会照单全收。

\subsubsection{基于跨视图预测的目标构造}
\label{sec:data-derived-knowledge-cross-view-prediction}

对比判别依赖负对照，而负对照又带来了“同类被误推开”和“需要大批量”等麻烦。一个自然的疑问是：\emph{能不能干脆不要负对照，只用正对应来学}？跨视图预测正是这样的尝试。仍给定同一照片的两个裁剪$A$、$B$，现在\emph{不再}要求模型从一堆候选中挑出配对，而是直接让裁剪$A$\emph{去预测}裁剪$B$的特征向量。这就是\term{跨视图预测}（cross-view prediction）：它的目标由“对应那一侧视图”的编码分支计算出来，全程\emph{只用正对应}，不需要任何负对照。

BYOL是这一思路的代表，它的结构是\emph{两个不对称的分支}。\emph{在线分支}依次执行三步：编码$f_{\bm{\theta}}$、投影$g_{\bm{\theta}}$、以及一个额外的\emph{预测}$q_{\bm{\theta}}$。\emph{目标分支}则只有编码器和投影器，没有预测器，其参数记为$\bm{\xi}$（与在线分支参数$\bm{\theta}$不同）。在线分支多出的那个预测器，作用是把在线投影再映射一次，变成“要去和目标匹配的”向量。对两个视图$\bm{v}_A,\bm{v}_B$，分别记在线预测和目标投影为
\[
  \bm{p}_A=q_{\bm{\theta}}\!\left(g_{\bm{\theta}}(f_{\bm{\theta}}(\bm{v}_A))\right),
  \qquad
  \bm{z}^{\mathrm{tar}}_B=g_{\bm{\xi}}(f_{\bm{\xi}}(\bm{v}_B)).
\]
有了预测$\bm{p}_A$和目标$\bm{z}^{\mathrm{tar}}_B$，还需一个衡量二者偏差的分数，在线网络才能据此改进预测。假定两向量维数相同且都非零。先对它们做$\ell_2$归一化（变成单位长度），只比较方向，再计算平方距离——距离越小，说明当前预测与目标越一致。

这里引入一个对BYOL至关重要的操作：\emph{停止梯度}（stop-gradient），记作$\operatorname{sg}$。它的含义是：\emph{前向}计算时照常使用目标的数值，但\emph{反向}求梯度时把它\emph{当作一个常数}，不让梯度流进目标分支。打个比方，目标像一个“暂时冻结的靶子”，在线分支朝它射箭、调整自己，但这一次射箭本身不会去挪动靶子。单向损失写为
\begin{equation}
  \ell_{A\to B}
  =\left\lVert
  \frac{\bm{p}_A}{\lVert\bm{p}_A\rVert_2}
  -\operatorname{sg}\!\left(
  \frac{\bm{z}^{\mathrm{tar}}_B}{\lVert\bm{z}^{\mathrm{tar}}_B\rVert_2}
  \right)\right\rVert_2^2
  =2-2\frac{\langle\bm{p}_A,\bm{z}^{\mathrm{tar}}_B\rangle}
  {\lVert\bm{p}_A\rVert_2\lVert\bm{z}^{\mathrm{tar}}_B\rVert_2}.
  \label{eq:data-derived-knowledge-byol}
\end{equation}
式子右侧的化简很值得理解：两个\emph{单位向量}之差的平方展开后，由于各自的平方范数都等于$1$，就剩下$2-2\times(\text{余弦相似度})$。这说明\emph{减小归一化平方距离，等价于提高两向量的余弦相似度}——两种说法是同一件事。

BYOL还会把两个视图的角色对调，再算一项$\ell_{B\to A}$（让$B$预测$A$），把两项相加作为总损失。这正好对应5.2.2引言里强调的那个区分：交换视图\emph{另算一项}，使两个视图\emph{都}有机会充当输入；但\emph{每一项}仍然只向它对应的\emph{在线}分支求导，目标端始终是被停止梯度的那一侧。两向求和并\emph{没有}把目标端变成接收梯度的一端\scite{ka-grill2020byol}。

既然目标端在当前这一步\emph{不接受}梯度，就必须另外回答一个问题：\emph{下一步的目标由哪套参数产生}？如果目标参数永远不变，它就无法随学习进步。BYOL的做法是：\emph{不用}损失梯度去更新目标，而是在在线分支完成本步更新\emph{之后}，让目标编码器与投影器的参数，向在线分支\emph{对应部分}的参数\emph{缓慢靠拢}。具体用的是\emph{指数移动平均}（exponential moving average, EMA）：每一步保留大部分旧的目标参数、只混入一小部分新的在线参数，使目标像在线分支的一个“滞后而平滑的影子”。迭代步$t$的更新为
\begin{equation}
  \bm{\xi}^{(t+1)}
  =m\bm{\xi}^{(t)}+(1-m)\bm{\theta}_{\mathrm{enc+proj}}^{(t+1)},
  \qquad 0\leq m\leq1.
  \label{eq:data-derived-knowledge-byol-ema}
\end{equation}
读这个式子要注意两点。其一，$\bm{\theta}_{\mathrm{enc+proj}}^{(t+1)}$\emph{只包含}在线分支\emph{更新后}的编码器和投影器参数；在线分支独有的\emph{预测器在目标侧没有对应部件}，所以它不参加EMA，只按普通梯度下降更新。其二，系数$m$决定新旧参数的混合比例（$m$越大，目标越稳定、变化越慢）。

最重要的是分清BYOL中\emph{两个看似相似、作用完全不同}的机制，它们各管一件事，绝不能混为一谈。\emph{停止梯度}约束的是“\emph{当前这一步}的求导”：它保证这一步的目标答案不会被这项误差反过来推动。\emph{EMA}约束的是“\emph{跨越多步}的参数变化”：它让未来用来生成答案的目标参数，随时间逐步跟随在线网络。一个管“这一步靶子动不动”，一个管“靶子以后怎么慢慢挪”——前者让靶子此刻冻结，后者让靶子长期缓慢跟进。

回头对照本章开头的“现成答案 vs 动态向量”区分，可以准确定位BYOL的目标属于哪一类。BYOL要预测的目标\emph{不是}裁剪$B$的\emph{像素}，而是\emph{当前目标网络}对$B$算出的向量$\bm{z}^{\mathrm{tar}}_B$。像素可以在编码前原样保留，那属于第~\ref{sec:data-derived-knowledge-observation}节的观测目标；而这个表示向量会随网络参数的变化而变化，是典型的“动态向量目标”。

它也必须与第~\ref{chap:knowledge-acquisition-model-based}章的“任务教师”区分开，尽管二者都可能被叫作“教师”或“目标分支”。分类任务里的教师，其输出的每个坐标都\emph{对应一个预设的类别}（第1维=山茶，第2维=茶梅）；而BYOL目标向量的坐标\emph{并没有}事先被命名为任何花卉物种，只是一串未命名的特征数值。所以不能因为名字都叫“教师/目标”，就假设它们的输出含义相同——必须检查输出到底代表什么。

与SimCLR一样，训练用的部件和最终交付的部件并不相同。预训练结束后，BYOL通常\emph{只保留在线编码器}$f_{\bm{\theta}}$输出的表示供下游使用，而把投影器、预测器和整个目标分支都\emph{丢弃}。这意味着训练时被拉得很近的$\bm{p}_A$和$\bm{z}^{\mathrm{tar}}_B$，\emph{并不是}最终交给分类器的那个向量；“它俩在训练中匹配得好”不能直接说明编码特征有用，必须在编码特征上\emph{另行}检验下游任务效用。

BYOL匹配的是两个连续向量的\emph{方向}。另一种思路是：不匹配向量本身，而是把每一路的输出\emph{转换成一个概率分布}，再让两个分布彼此一致。DINO采用的就是这种跨视图目标。它也有一对分支，称为\emph{学生}和\emph{教师}：二者都由骨干编码器加一个投影头组成，\emph{结构相同、参数不同}；与BYOL的一个区别是，DINO\emph{没有}BYOL那种额外的预测头。

具体地，投影头为视图$\bm{v}$输出一个$K$维的分数向量，学生记为$\bm{u}_{\mathrm{s}}(\bm{v})$、教师记为$\bm{u}_{\mathrm{t}}(\bm{v})$；下标$k,r$枚举的是这$K$个\emph{投影坐标}——再次强调，它们\emph{不是}物种类别。为了把分数变成分布，学生和教师都用softmax归一化，但各自用不同的正温度$\tau_{\mathrm{s}}$和$\tau_{\mathrm{t}}$；此外教师在归一化之前还要先\emph{减去一个按维度维护的中心}$\bm{c}$，这是为了调节教师目标集中在哪些维度上（作用见下文）：
\begin{align}
  p_{\mathrm{s},k}(\bm{v})
  &=\frac{\exp(u_{\mathrm{s},k}(\bm{v})/\tau_{\mathrm{s}})}
  {\sum_{r=1}^{K}\exp(u_{\mathrm{s},r}(\bm{v})/\tau_{\mathrm{s}})},\\
  p_{\mathrm{t},k}(\bm{v})
  &=\frac{\exp((u_{\mathrm{t},k}(\bm{v})-c_k)/\tau_{\mathrm{t}})}
  {\sum_{r=1}^{K}\exp((u_{\mathrm{t},r}(\bm{v})-c_r)/\tau_{\mathrm{t}})}.
  \label{eq:data-derived-knowledge-dino-distributions}
\end{align}
这样学生和教师都变成了同一组$K$个维度上的概率分布，就可以用交叉熵来衡量两者是否一致。DINO在视图的使用上做了一个精心的不对称安排：它为每张图生成\emph{两个全局视图}（覆盖大部分画面）和\emph{若干局部视图}（小块）；\emph{教师只看两个全局视图}，\emph{学生则既看全局视图也看局部视图}。训练时，每个教师投影分布去监督\emph{除同一视图以外}的所有学生投影分布。

这个安排的用意可以这样理解：沿用的正对应依据仍是“来自同一张图”，但它新增了一个更强的要求——\emph{学生只看到一小块局部，也要能预测出教师看全局时的分布}，从而迫使学生学会“从局部推断整体”。在更新方面，教师端同样停止梯度，并用学生参数的EMA来更新；而中心$\bm{c}$则\emph{另外}按教师在批量上投影分数的移动平均来更新\scite{ka-caron2021dino}。

前面提到的两个特殊操作，各自针对一种不同的坍塌风险，值得说清。\emph{减去中心}称为\emph{居中}（centering）：如果某个维度长期偏高，它会主导所有样本的分布、导致大家都集中到这一维（一种坍塌）；减去按维度累积的中心，就能抑制这种“一维独大”。\emph{降低教师温度}称为\emph{锐化}（sharpening）：它让教师的分布更尖锐、更集中在少数维度，用以对抗另一种坍塌——所有样本的分布都趋于\emph{均匀}、毫无区分。居中防“全挤到一维”，锐化防“全摊平”，二者方向相反、相互制衡。

但要守住解读的边界：原论文是把居中、锐化与按EMA更新的动量教师、多裁剪\emph{一起}使用，并通过消融实验观察各自作用的；\emph{不能}从其中任何单一组件，推出“任意模型只要加了它就不会坍塌”。此外，这$K$个维度始终\emph{没有}人工赋予的类别名称；交叉熵低，只说明学生与教师的投影分布匹配上了，\emph{不能}据此认定模型已经发现了$K$个真实的花卉物种。

与前面的方法一样，DINO的$K$维分布也只是\emph{预训练阶段}使用的对象，下游真正用的是骨干编码器输出的特征。由于DINO有学生和教师两套编码器，交付时必须明确选哪一套：本章约定取\emph{训练结束时的教师编码器}检查点，并舍弃投影头；使用学生检查点也是一种合理的接口选择，但必须\emph{明确注明}——绝不能把两套不同编码器的结果混在同一个评价里比较，否则数字之间没有可比性。

总结跨视图预测的得失。它\emph{省去了显式的负对照}（不用担心同类被误推开、也不必为负对照维持大批量），代价则是要额外维护一个独立的目标分支，或处理更多视图带来的计算。它还有两个需要警惕的内在问题。其一，当学生只看到局部输入、缺少全局内容时，它能\emph{真正从输入推出}的只是目标的一部分，其余部分只能依赖训练分布去“估计”，并非严格可推断。其二，也是最根本的一点：\emph{如果两路的表示都塌缩成同一个常量}，匹配误差同样可以非常小——没有负对照把样本推开，“大家都一样”反而是一个让损失很低的捷径。

这正是为什么只看更新规则是不够的，必须\emph{额外}检查表示是否发生了退化。第~\ref{sec:data-derived-knowledge-representation-degeneration}节将把这些退化状态放回完整的目标中系统分析。

\subsubsection{基于对称匹配的目标构造}
\label{sec:data-derived-knowledge-symmetric-matching}

BYOL和DINO都依靠\emph{不对称}的设计（一路当靶子停止梯度、另一路追赶）来避免坍塌。第三条思路更“民主”：对同一朵花的两个裁剪，\emph{不指定}哪一路提供固定答案，而是让\emph{两路同时调整}、相互靠近。\term{对称匹配}（symmetric matching）指的正是这样一种安排：同一项配对损失的梯度\emph{同时}回传到两路的表示$\bm{z}_i$与$\bm{z}'_i$；两路可以共享参数，也可以用不同结构，但都\emph{不需要}预先设立一个独立教师、也\emph{不需要}停止某一路的梯度。

这里的“对称”有精确含义，需与BYOL区分：它指的是\emph{在同一项损失中}两路扮演对称的求导角色、都被推动；而BYOL那种“交换两个输入后把两个\emph{不对称}的损失相加”，每一项里仍只有一路接收梯度。二者都有“两个方向”，但对称性所在的层面不同。

VICReg是对称匹配的代表，它把要求施加在一个专门的训练嵌入上。结构上，编码器$f_{\bm{\theta}}$先提取特征，再由一个称为\term{扩展头}（expander）的映射$h_{\bm{\phi}}$把特征\emph{映射到更高维}的计分嵌入空间，得到$\bm{z}_i=h_{\bm{\phi}}(f_{\bm{\theta}}(\bm{v}_i))$（与SimCLR的投影器作用类似，但VICReg常把维度扩大，故名“扩展”）。本章采用两路\emph{共享}这两个映射的安排，因此参数$\bm{\theta}$和$\bm{\phi}$都同时接收来自两路的梯度。

先实现最直接的“让两路接近”的要求。把本批$n$对样本的$d$维嵌入按\emph{样本作行}排成两个矩阵$\bm{Z},\bm{Z}'\in\mathbb{R}^{n\times d}$（需$n\geq2$，下文统计量要用到），然后对\emph{同一行}（同一对）的两个向量求平方距离、再在批内平均：
\begin{equation}
  s(\bm{Z},\bm{Z}')
  =\frac{1}{n}\sum_{i=1}^{n}\lVert\bm{z}_i-\bm{z}'_i\rVert_2^2.
  \label{eq:data-derived-knowledge-vicreg-invariance}
\end{equation}
减小这一项确实会让每对正对应彼此接近，但它暴露出一个致命漏洞——这正是对称匹配必须直面的坍塌问题：如果模型把\emph{所有}样本都映射成\emph{同一个常量向量}，那么每一对两侧都相同，平方距离全为零，这一项损失\emph{恰好达到最小值零}。模型什么都没学，却拿到了完美分数。

这个反例点明了问题的本质：这一项只要求“\emph{每一对内部}一致”，却\emph{完全没有要求}“\emph{不同样本之间}保留差异”。既然对称匹配没有负对照、也没有停止梯度来防坍塌，VICReg就必须另想办法：它在\emph{每一路}分别额外加入两项约束——\emph{方差}约束和\emph{协方差}约束\scite{ka-bardes2022vicreg}——直接从统计上禁止表示塌缩。下面分别介绍。

第一项约束针对“所有样本取同一个值”这种坍塌。要检查某一个坐标是否对所有样本都给出了相同的值，视角就要从“比较一对之内”转到“\emph{看一整批样本在这个坐标上变化有多大}”。对嵌入的第$j$维，先算出批内均值$\bar z_j=n^{-1}\sum_i z_{i,j}$，再用分母为$n-1$的样本方差衡量各样本偏离均值的程度——方差为零就意味着这一维对所有样本都一样（坍塌）。为了数值稳定，加上一个小的正数$\epsilon$后再开平方，得到用于后续阈值比较的标准差：
\[
  \begin{aligned}
    \operatorname{Var}(\bm{Z}_{:,j})
    &=\frac{1}{n-1}\sum_{i=1}^{n}(z_{i,j}-\bar z_j)^2,\\
    \sigma_j(\bm{Z})
    &=\sqrt{\operatorname{Var}(\bm{Z}_{:,j})+\epsilon}.
  \end{aligned}
\]
标准差本身只是一个\emph{诊断}量，告诉我们某一维变化有多大；要让训练\emph{主动惩罚}“变化不足”，还需设定一个正的目标阈值$\gamma$：凡是标准差低于$\gamma$的维度都要受罚，达到或超过$\gamma$的则不罚。这里取$\epsilon>0$且$\gamma>\sqrt{\epsilon}$，这个条件很关键——它保证即便是一个完全坍塌的常数批次（方差为0，标准差为$\sqrt{\epsilon}$），其标准差也\emph{低于}阈值、从而会被惩罚。具体地，对每一维只保留差距$\gamma-\sigma_j$的\emph{正部}（只罚不足、不奖励过量），再在$d$维上平均：
\begin{equation}
  v(\bm{Z})
  =\frac{1}{d}\sum_{j=1}^{d}\max(0,\gamma-\sigma_j(\bm{Z})).
  \label{eq:data-derived-knowledge-vicreg-variance}
\end{equation}
这一项回答的是“\emph{哪些维度在批内仍然变化不足}”。其中$\max(0,\cdot)$的作用是只提取“距阈值还差多少”这个缺口，而\emph{不是}去找某个最坏的样本；变化已经足够的维度，缺口为零，不贡献惩罚。还要明白它的性质：这一项是与其他项\emph{加权}一起优化的\emph{软惩罚}，它鼓励但\emph{并不保证}每一步训练都能让所有维度达到$\sigma_j\geq\gamma$。

方差约束解决了“完全不变”，却挡不住另一种更隐蔽的退化：\emph{每一维都在变化，但它们其实都在重复同一个变化}。举个例子，若两个坐标对所有样本\emph{始终都相等}（都等于某个$a$），那么每一维单看都有充足方差，方差项完全满意；可这两维提供的信息一模一样，名义上是二维、实际只有一维有用。这称为\emph{维度冗余}（后文“维度退化”的根源）。

要识别这种冗余，需要衡量\emph{不同坐标是否一起偏离均值}——如果两维总是同涨同跌，它们就是线性相关的、冗余的。为此令批内均值向量$\overline{\bm{z}}=n^{-1}\sum_i\bm{z}_i$，把每个样本的中心化向量与自身做外积、再平均，得到一个$d\times d$的样本协方差矩阵：
\[
  \bm{C}(\bm{Z})
  =\frac{1}{n-1}\sum_{i=1}^{n}
  (\bm{z}_i-\overline{\bm{z}})
  (\bm{z}_i-\overline{\bm{z}})^{\top}.
\]
协方差矩阵的结构恰好把两类信息分开：它的\emph{对角元素}就是前面已经检查过的各维方差；而\emph{非对角元素}$C_{j,k}$记录的是第$j$维与第$k$维之间的\emph{线性共同变化}——它越大，说明这两维越冗余。为了只惩罚冗余、而\emph{不去压低}单维自身的方差（那会与方差项冲突），VICReg只把\emph{非对角}元素的平方累加起来，并按维度$d$归一化：
\begin{equation}
  c(\bm{Z})
  =\frac{1}{d}\sum_{\substack{j,k=1\\j\neq k}}^{d}C_{j,k}(\bm{Z})^2.
  \label{eq:data-derived-knowledge-vicreg-covariance}
\end{equation}
注意协方差项\emph{单独}使用同样防不住坍塌：一个所有样本都相同的常数表示，其协方差矩阵全为零，协方差项也是零——它“满意”。所以三项必须\emph{联合}使用、互相补位。把它们以正权重组合起来，训练就同时朝三个方向努力：拉近正对应（匹配项）、补足各维变化（方差项）、消除维度间冗余（协方差项）。完整的VICReg目标为
\begin{equation}
  \mathcal{L}_{\mathrm{VICReg}}
  =\lambda s(\bm{Z},\bm{Z}')
  +\mu\!\left[v(\bm{Z})+v(\bm{Z}')\right]
  +\nu\!\left[c(\bm{Z})+c(\bm{Z}')\right],
  \label{eq:data-derived-knowledge-vicreg}
\end{equation}
其中正权重$\lambda,\mu,\nu$分别控制三种要求的相对重要性。读这个式子要注意三项的\emph{统计范围}不同：匹配项是在每一\emph{对}样本上计分的；方差和协方差项则是在\emph{每一路的整个批次}（$n$个样本）上计算的统计量。VICReg的精妙之处在于：它通过把统计范围从“一对”扩大到“一整批”，就实现了防坍塌的效果，\emph{而无需引入任何逐对的负例标签}——这正是它区别于SimCLR的地方。

用两个反例检验这套组合到底挡住了什么。\emph{反例一：完全坍塌。}设$\bm{Z}=\bm{Z}'$且所有行都相同（所有样本同一常量）。此时匹配项为零、协方差项为零，唯独方差项会出手：每一路的方差惩罚为$\max(0,\gamma-\sqrt{\epsilon})$。正是因为前面要求了$\gamma>\sqrt{\epsilon}$，这个惩罚\emph{为正}——在教学设定$\gamma=1,\epsilon=10^{-4}$下，$v(\bm{Z})=1-\sqrt{10^{-4}}=0.99$，坍塌被惩罚了。这也反过来说明了那个条件的必要性：若把参数改成$\gamma=\epsilon=1$，则$\gamma=\sqrt{\epsilon}$，惩罚变成零，三项同时为零，坍塌就逃过了所有检查。

但即便惩罚为正，也有一个与SimCLR相同的精细限定：\emph{正惩罚不等于优化一定能逃离坍塌}。因为在\emph{精确}的常数批次处，方差项对表示的梯度仍然为零，优化可能被困在那里。\emph{反例二：维度冗余。}若各行都沿着$(a,a)$方向变化，那么两个坐标\emph{都有}方差（方差项满意），却在重复同一个方向；这时是协方差项中\emph{非零}的非对角元素揭示了这种冗余。两个反例分别对应方差项和协方差项各自负责的那种退化。

这些统计约束在工程上还依赖批量和嵌入维数的选择。批量太小时，用少数样本估计的方差和协方差会\emph{波动很大}、不稳定；嵌入维数$d$很大时，显式计算和存储一个$d\times d$的协方差矩阵开销显著（随$d^2$增长）。

更根本的是要看清这些统计惩罚\emph{能保证什么、不能保证什么}。即使表示的变化已经很充分、维度间线性相关也很低，这些\emph{变化本身}仍可能来自背景、光照或拍摄设备，而非花的特征——统计惩罚只管“有没有变化、是否冗余”，\emph{并不为表示赋予物种语义}，也\emph{不保证}优化一定收敛。最后与前几种方法一样：预训练后VICReg会\emph{舍弃}扩展头$h_{\bm{\phi}}$，只用编码器$f_{\bm{\theta}}$的特征，所以仍需在\emph{这个}接口上检验下游用途。

学完三种目标，现在可以把它们\emph{放在同一张图上对照}，形成一个整体认识。图~\ref{fig:data-derived-knowledge-relation-objectives}用统一的线型画出同一对视图在三种方法中的三种计分路径，读图时抓住线型含义：\emph{实线}追踪参与损失的前向值，\emph{橙色虚线}追踪当前这一步梯度的去向，\emph{点划线}则另外表示\emph{跨越多步}的参数平均（EMA）。

这样对照后，一个最关键的结构差异立刻凸显：\emph{BYOL的目标端有前向值、却没有损失梯度}（虚线不指向它），而\emph{VICReg的同一项损失却同时向两路回传}。这正是“停止梯度的独立目标”与“两路共同求导”的分水岭。还要注意，是否再交换输入求第二项，\emph{不会改变}这个根本区别——BYOL即使两个方向都算，每项的目标端依然不接收梯度。

\begin{figure*}[htbp]
  \centering
  \begin{tikzpicture}[
  x=1mm,y=1mm,
  font=\figurefont\footnotesize,
  panel/.style={draw=BookRule,fill=NNPanel,line width=.85pt,rounded corners=1.5mm},
  input/.style={draw=NNInput,fill=NNInput!16,line width=1pt,minimum width=7mm,minimum height=7mm,inner sep=1mm,align=center},
  compute/.style={draw=NNHidden,fill=NNHidden!16,line width=1pt,minimum width=10mm,minimum height=8mm,inner sep=1mm,align=center,rounded corners=1mm},
  output/.style={draw=NNOutput,fill=NNOutput!16,line width=1pt,minimum width=8mm,minimum height=8mm,inner sep=1mm,align=center},
  loss/.style={draw=NNGradient,fill=NNGradient!12,line width=1pt,minimum width=10mm,minimum height=8mm,inner sep=1mm,align=center,rounded corners=1mm},
  forward/.style={knowledge arrow,draw=NNWire,line width=1pt},
  gradient/.style={knowledge arrow,draw=NNGradient,line width=1.1pt,dashed,rounded corners=1mm},
  gradient stop/.style={-{Bar[width=2.4mm]},draw=NNGradient,line width=1.1pt,dashed,rounded corners=1mm},
  ema/.style={knowledge arrow,draw=BookTeal,line width=1pt,dash dot},
  note/.style={font=\figurefont\scriptsize,text=BookMuted,align=center}
]
  \draw[panel] (.5,103) rectangle (168.5,153);
  \draw[panel] (.5,48) rectangle (168.5,98);
  \draw[panel] (.5,-10) rectangle (168.5,42);
  \node[font=\figurefont\bfseries\footnotesize,text=BookInk,anchor=west] at (5,148) {(a) 对比判别：SimCLR};
  \node[note,anchor=east] at (163,148) {两路共享编码器与投影器};
  \node[input] (s-v1) at (10,132) {$\bm{v}_1$};
  \node[input] (s-v2) at (10,117) {$\bm{v}_2$};
  \node[compute,minimum width=20mm] (s-f1) at (42,132) {$f_{\bm{\theta}}$};
  \node[compute,minimum width=20mm] (s-f2) at (42,117) {$f_{\bm{\theta}}$};
  \node[compute,minimum width=20mm] (s-g1) at (76,132) {$g_{\bm{\theta}}$};
  \node[compute,minimum width=20mm] (s-g2) at (76,117) {$g_{\bm{\theta}}$};
  \node[output,minimum width=14mm] (s-z1) at (108,132) {$\bm{z}_1$};
  \node[output,minimum width=14mm] (s-z2) at (108,117) {$\bm{z}_2$};
  \node[loss,minimum width=26mm] (s-loss) at (148,124.5) {相似度$\bm{S}$\\与配对损失};
  \draw[forward] (s-v1.east)--(s-f1.west);
  \draw[forward] (s-f1.east)--(s-g1.west);
  \draw[forward] (s-g1.east)--(s-z1.west);
  \draw[forward] (s-v2.east)--(s-f2.west);
  \draw[forward] (s-f2.east)--(s-g2.west);
  \draw[forward] (s-g2.east)--(s-z2.west);
  \draw[NNWire,line width=1pt,rounded corners=1mm] (s-z1.east) -- (126,132) -- (126,124.5);
  \draw[NNWire,line width=1pt,rounded corners=1mm] (s-z2.east) -- (126,117) -- (126,124.5);
  \draw[forward] (126,124.5)--(s-loss.west);
  \draw[gradient] (s-loss.north) -- (148,140) -- (42,140) -- (s-f1.north);
  \draw[gradient] (76,140)--(s-g1.north);
  \draw[gradient] (108,140)--(s-z1.north);
  \draw[gradient] (s-loss.south) -- (148,109) -- (42,109) -- (s-f2.south);
  \draw[gradient] (76,109)--(s-g2.south);
  \draw[gradient] (108,109)--(s-z2.south);

  \node[font=\figurefont\bfseries\footnotesize,text=BookInk,anchor=west] at (5,93) {(b) 跨视图预测：BYOL};
  \node[note,anchor=east] at (163,93) {只画$A\to B$；交换视图另算一项};
  \node[input] (b-va) at (10,78) {$\bm{v}_A$};
  \node[compute,minimum width=20mm] (b-fo) at (42,78) {$f_{\bm{\theta}}$};
  \node[compute,minimum width=20mm] (b-go) at (76,78) {$g_{\bm{\theta}}$};
  \node[compute,minimum width=20mm] (b-qo) at (108,78) {$q_{\bm{\theta}}$};
  \node[input] (b-vb) at (10,59) {$\bm{v}_B$};
  \node[compute,minimum width=20mm] (b-ft) at (42,59) {$f_{\bm{\xi}}$};
  \node[compute,minimum width=20mm] (b-gt) at (76,59) {$g_{\bm{\xi}}$};
  \node[output,minimum width=20mm] (b-sg) at (108,59) {$\operatorname{sg}$};
  \node[loss,minimum width=20mm] (b-loss) at (148,68.5) {$\ell_{A\to B}$};
  \draw[forward] (b-va.east)--(b-fo.west);
  \draw[forward] (b-fo.east)--(b-go.west);
  \draw[forward] (b-go.east)--(b-qo.west);
  \draw[forward,rounded corners=1mm] (b-qo.east) -| (b-loss.north);
  \draw[forward] (b-vb.east)--(b-ft.west);
  \draw[forward] (b-ft.east)--(b-gt.west);
  \draw[forward] (b-gt.east)--(b-sg.west);
  \draw[forward,rounded corners=1mm] (b-sg.east) -| (b-loss.south);
  \draw[gradient] (b-loss.east) -- (164,68.5) -- (164,86) -- (42,86) -- (b-fo.north);
  \draw[gradient] (76,86)--(b-go.north);
  \draw[gradient] (108,86)--(b-qo.north);
  \draw[gradient stop] (164,68.5) -- (164,52) -- (108,52) -- (b-sg.south);
  \draw[ema] (b-fo.south)--(b-ft.north);
  \draw[ema] (b-go.south)--(b-gt.north);
  \node[note,text=BookTeal] at (59,68.5) {EMA};

  \node[font=\figurefont\bfseries\footnotesize,text=BookInk,anchor=west] at (5,37) {(c) 对称匹配：VICReg};
  \node[note,anchor=east] at (163,37) {两路共享编码器与扩展头；无停止梯度};
  \node[input] (v-va) at (10,22) {$\bm{v}_A$};
  \node[compute,minimum width=28mm] (v-fa) at (60,22) {$h_{\bm{\phi}}\circ f_{\bm{\theta}}$};
  \node[output,minimum width=14mm] (v-za) at (100,22) {$\bm{Z}$};
  \node[input] (v-vb) at (10,4) {$\bm{v}_B$};
  \node[compute,minimum width=28mm] (v-fb) at (60,4) {$h_{\bm{\phi}}\circ f_{\bm{\theta}}$};
  \node[output,minimum width=14mm] (v-zb) at (100,4) {$\bm{Z}'$};
  \node[loss,minimum width=34mm] (v-loss) at (148,13) {匹配$s(\bm{Z},\bm{Z}')$\\方差$v(\bm{Z}),v(\bm{Z}')$\\协方差$c(\bm{Z}),c(\bm{Z}')$};
  \draw[forward] (v-va.east)--(v-fa.west);
  \draw[forward] (v-fa.east)--(v-za.west);
  \draw[forward] (v-vb.east)--(v-fb.west);
  \draw[forward] (v-fb.east)--(v-zb.west);
  \draw[NNWire,line width=1pt,rounded corners=1mm] (v-za.east) -- (123,22) -- (123,13);
  \draw[NNWire,line width=1pt,rounded corners=1mm] (v-zb.east) -- (123,4) -- (123,13);
  \draw[forward] (123,13)--(v-loss.west);
  \draw[gradient] (v-loss.north) -- (148,29) -- (60,29) -- (v-fa.north);
  \draw[gradient] (100,29)--(v-za.north);
  \draw[gradient] (v-loss.south) -- (148,-3) -- (60,-3) -- (v-fb.south);
  \draw[gradient] (100,-3)--(v-zb.south);

  \draw[forward] (3,-17) -- (16,-17);
  \node[note,anchor=west] at (18,-17) {前向值};
  \draw[gradient] (69,-17) -- (56,-17);
  \node[note,anchor=west] at (72,-17) {当前步梯度};
  \draw[ema] (120,-17) -- (133,-17);
  \node[note,anchor=west] at (136,-17) {跨步模型参数平均};
\end{tikzpicture}

  \caption{沿三种线型核对关系目标的更新角色。实线传前向值，橙色虚线传当前梯度，青色点划线表示跨步EMA。SimCLR共享编码器与投影器，并使用批内对照；BYOL仅画$A\to B$，sg阻断目标端梯度，EMA只连接编码器和投影器，交换视图需另算一项；VICReg共享$f_{\bm{\theta}}$与$h_{\bm{\phi}}$，匹配项及各路统计项都向相应计算分支回传。}
  \label{fig:data-derived-knowledge-relation-objectives}
\end{figure*}
\FloatBarrier

图给直觉，表给核对。表~\ref{tab:data-derived-knowledge-relation-objectives}从三个维度对四种方法做严格对照：\emph{比较对象是什么}（投影相似度、在线预测对目标投影、投影分布、嵌入匹配加批内统计）、\emph{当前梯度去向}（回传几路）、\emph{跨步如何变化}（共享参数重算、EMA等）。

读这张表要破除一种简单归类的冲动。表面上看，“用没用显式负例”似乎能把方法分成两类：SimCLR用，BYOL/DINO/VICReg不用。但同样\emph{不用显式负例}的三种方法，却选择了截然不同的防坍塌机制——BYOL和DINO靠\emph{独立的目标分支（教师）}，VICReg则\emph{不要教师、靠两路共同求导再加批内统计}。因此“无显式负例”“有教师”“停止梯度”“两向求和”分别描述的是\emph{四个不同属性}，\emph{不能}因为一个方法满足其中一项，就推断它也满足其他项。

\begin{table*}[htbp]
  \centering
  \renewcommand{\arraystretch}{1.15}
  \begin{tabularx}{\linewidth}{@{}>{\RaggedRight\setlength{\parindent}{0pt}\arraybackslash}p{22mm}*{3}{>{\RaggedRight\setlength{\parindent}{0pt}\arraybackslash}X}@{}}
    \toprule
    方法 & 比较对象 & 当前梯度去向 & 跨步变化 \\
    \midrule
    SimCLR & 投影表示的相对相似度 & 两路编码器与投影器 & 共享模型参数更新后重算表示 \\
    BYOL & 在线预测与目标投影 & 每项仅回传在线分支 & 目标编码器与投影器采用EMA \\
    DINO & 学生与教师的投影分布 & 每项仅回传学生分支 & 教师模型参数与投影分数中心分别移动平均 \\
    VICReg & 嵌入匹配与批内统计 & 两路编码器与扩展头 & 共享模型参数更新后重算表示及统计量 \\
    \bottomrule
  \end{tabularx}
  \caption{固定正对应后，比较量与更新方式仍可不同。BYOL在线端有额外预测器，DINO没有该部件，VICReg则没有独立教师。表中组件承担不同职责，完整退化条件和计算成本需结合正文判断。}
  \label{tab:data-derived-knowledge-relation-objectives}
\end{table*}
\FloatBarrier

前面所有损失都在回答“自监督目标优化得好不好”，但读者最关心的问题是：\emph{这些训练依据到底支不支持物种识别？}要公平回答它，需要一套统一、冻结、与训练隔离的评价协议，避免“谁的训练头巧妙谁就显得好”。

标准做法如下：固定同一批任务样本、同一套标签划分和同一个\emph{线性分类协议}——\emph{冻结}每种方法训练后选定的编码器，只在冻结特征$\bm{h}=f(\bm{x})$上\emph{额外训练一个线性分类头}；然后在一批\emph{既没参与这个分类头拟合、也没参与任何超参调整}的任务样本上，统计“分错的样本数除以总样本数”得到误分类率。这个数值越低，说明在\emph{该协议}下，冻结的编码特征对物种识别任务的支持越强。

要理解这种比较比的是什么：比的是\emph{编码特征本身对任务的支撑能力}，\emph{不是}各方法的投影头/预测头在其自身自监督目标上的分数（那些分数高，不代表特征好）。如果评价时\emph{连编码器也一起更新}，那比较对象就变成了“整体微调”而非“特征质量”，应当\emph{另行}报告一种不同的协议。此外，不同方法产出的向量接口不同，也\emph{不能混为一谈}：SimCLR的投影向量、BYOL的预测向量、DINO的$K$维分布都只是训练用的，真正可比的是它们各自编码器在同一接口下的特征。

\subsection{目标更新与表示退化的控制}
\label{sec:data-derived-knowledge-target-update}

本节专门处理那些“目标会随训练变化”的方法（如BYOL、DINO）里一个容易让人糊涂的问题。在BYOL的同一对裁剪上，虽然当前这一步的目标\emph{不接受}梯度，但到了\emph{下一步}、用更新后的网络重新编码时，目标向量很可能\emph{已经变了}。这就引出一连串需要厘清时点的问题：这一步的目标是什么、梯度能沿哪条路回传、哪些参数被更新了、下一步的目标又如何准备？

如果只是笼统地说一句“教师随学生更新”，就把这些关键时点全部跳过了——你既说不清目标在\emph{什么时候保持不变}、也说不清它\emph{什么时候、以多大幅度发生变化}。因此下面按\emph{时间顺序}逐步追踪：前向产生目标、反向允许的梯度路径、接收梯度的参数更新、以及下一步新目标的准备。

除了“目标怎么变”，还有一个与之并列的问题：\emph{一个会平滑变化的目标，是否仍然保留了样本之间的差异}？如前所述，只要两路都给常量，匹配误差照样可以很小——目标平滑更新本身\emph{并不}保证表示不退化。因此还要系统检查：完整目标施加了哪些区分性、方差或去相关的要求？哪些退化被惩罚了？又有哪些退化状态由于梯度为零，仍然可能\emph{把优化停住}？

这两小节回答的是两个互补、不可互相替代的问题：\emph{目标更新规则}回答“目标随时间怎样变化”（多快、是否平滑）；\emph{退化检查}回答“在这个变化过程中可能丢掉哪些信息、会不会塌缩”。更新得平滑不代表信息保留完整，二者必须分别考察。

\subsubsection{目标表示的更新}
\label{sec:data-derived-knowledge-target-representation-update}

先比较不同方法在“下一步如何取得目标”上的安排，大致两类。\emph{共享参数}型（SimCLR和常见VICReg实现）：两路用的是同一套参数，一次优化后参数变成新的，两路的表示都用这同一套\emph{新参数}重新计算，没有独立维护的目标网络。\emph{独立目标分支}型：另存一套目标参数，而它的更新又有几种不同强度——可以\emph{全程固定}（教师从头不变）、\emph{周期性地}把在线参数复制过去、或\emph{逐步地}做移动平均（BYOL/DINO）。

这里必须把三个概念严格分开，它们常被混用：\emph{固定教师}意味着目标参数\emph{整个训练过程都不变}；\emph{停止梯度}只表示“\emph{当前这一步}求导时梯度不沿目标路径回传”，与目标\emph{以后}更不更新\emph{无关}；而\emph{省略目标计算}则更进一步，连前向值都不存在。这三者是不同层面的安排：固定教师不代表用了停止梯度（虽然常一起用），停止梯度也不代表教师以后不更新；它们不能互相代替。

表~\ref{tab:data-derived-knowledge-gradient-and-update}用BYOL式更新检查停止梯度与目标变化如何并存。取教学标量：参与EMA的在线模型参数$\theta^{(t)}=2.00$、梯度$0.50$、学习率$\eta=0.20$，对应目标模型参数$\xi^{(t)}=1.60$、系数$m=0.90$。在线模型参数先变为$2.00-0.20\times0.50=1.90$，目标再按$0.90\times1.60+0.10\times1.90=1.63$改变。目标没有接受这项梯度，仍在跨步平均后改变；在线预测器因没有对应目标模型参数，不参加这次平均。

\begin{table*}[htbp]
  \centering
  \renewcommand{\arraystretch}{1.15}
  \begin{tabularx}{\linewidth}{@{}>{\RaggedRight\setlength{\parindent}{0pt}\arraybackslash}p{25mm}*{3}{>{\RaggedRight\setlength{\parindent}{0pt}\arraybackslash}X}@{}}
    \toprule
    时点 & 在线分支 & 目标分支 & 本步操作 \\
    \midrule
    1. 前向 & $\bm{\theta}^{(t)}$产生预测 & $\bm{\xi}^{(t)}$产生目标 & 比较预测与停止梯度的目标值 \\
    2. 反向 & 接收梯度；示例为$0.50$ & 不接收损失梯度 & 只沿在线路径累计梯度 \\
    3. 在线更新 & 优化器更新；示例为$1.90$ & 仍为$\bm{\xi}^{(t)}$ & 先完成在线模型参数更新 \\
    4. 目标更新 & 保持$\bm{\theta}^{(t+1)}$ & EMA更新；示例为$1.63$ & 用式~\eqref{eq:data-derived-knowledge-byol-ema}平均对应模型参数 \\
    5. 下一步 & 新模型参数产生新预测 & 新模型参数重新计算目标 & 进入$t+1$步前向计算 \\
    \bottomrule
  \end{tabularx}
  \caption{按时点区分损失梯度与目标模型参数变化。固定教师省去第4行，目标模型参数全程不变；共享模型参数方案没有独立$\bm{\xi}$，两路都读取同次更新后的模型参数。VICReg的共同反向路径应按图~\ref{fig:data-derived-knowledge-relation-objectives}理解，不套用独立目标列。}
  \label{tab:data-derived-knowledge-gradient-and-update}
\end{table*}
\FloatBarrier

EMA更新的“跟随速度”由系数$m$控制，形成一对权衡。$m$较\emph{小}时，新的在线参数权重$(1-m)$较大，目标\emph{跟随得快}，几乎立刻反映学生的最新状态（但也更不稳定）；$m$越\emph{接近}1，保留的旧目标比例越高，目标\emph{变化越缓慢、越平滑}（稳定的好处是目标不会被单步噪声带偏，代价是\emph{滞后}——学生已经进步了，目标还停在旧状态）；极端$m=1$时$(1-m)=0$，目标\emph{不再混入任何新值}，等同于一个完全冻结的教师。

需要补充的是，目标的实际变化速度\emph{不只由$m$}决定，还受优化器步长、系数的训练日程（$m$可能随时间变大）、训练总长度影响，因此脱离这些条件\emph{单比一个$m$}无法确定最终结果。最后算一笔成本：维护独立目标端意味着要\emph{额外保存}一整套参数，并在每步多做一次\emph{前向计算}，这是换取目标稳定所付出的代价。

复现这种EMA训练时，\emph{更新顺序}必须严格保持，否则结果无法对齐。注意式~\eqref{eq:data-derived-knowledge-byol-ema}读取的是\emph{已经更新后的}新在线参数$\bm{\theta}^{(t+1)}$（先更新学生，再平均教师）；若错误地改读更新前的$\bm{\theta}^{(t)}$，目标就会\emph{多滞后一步}，得到的是另一种算法。在多卡分布式训练中，若各设备在执行EMA\emph{之前}没有先把在线参数\emph{同步}一致，每张卡可能用不同的新参数维护出\emph{各不相同}的目标网络，训练就发散了。因此必须保存并核对更新顺序与同步时点，才能确认“下一步计算用的究竟是哪一个目标”。

DINO里其实有\emph{两处}移动平均，容易被混为一谈，要分开维护。第一处是\emph{教师网络参数}的EMA，它更新的是“生成特征和投影分数的那个网络”。第二处是\emph{教师投影分数中心}$\bm{c}$的移动平均，它更新的是“在算教师分布时、从分数里\emph{减去}的那个按维度基准”。虽然两处都用了“移动平均”，但它们\emph{保存的状态}（一个是网络权重、一个是一个$K$维统计向量）和\emph{作用的位置}（一个改变整个网络、一个只改变分布的居中量）完全不同，不能因为“都是EMA”就合并成同一个更新。

\subsubsection{表示退化的控制}
\label{sec:data-derived-knowledge-representation-degeneration}

前面多次提到“坍塌”，这里给出它的严格定义和第二种更隐蔽的形式。设想两张\emph{明显不同}的花卉照片，连同它们各自的裁剪，经过编码器后竟然\emph{都得到同一个向量}。从配对匹配的角度看，这简直“完美成功”——每一对正对应两侧都完全一致、距离为零；可这个表示已经\emph{彻底丧失了区分两张不同照片的能力}，对下游识别毫无用处。这种\emph{所有样本的表示都退化成同一个常量}的情形，就是\term{表示坍塌}（representation collapse）最直接的形式。

但退化不止这一种。还有一种更隐蔽的\term{维度退化}（dimensional degeneration）：不同样本的表示\emph{确实在变化}（不是常量），但这些变化\emph{只发生在少数几个方向上}，多个坐标重复着\emph{同一个变化方向}。比如一个二维表示对所有样本都形如$(a,a)$——$a$可以取不同值、看起来在变，但真正能独立变化的量只有\emph{一个}，它名义上是二维、实际只用到了一维。有效维数低于名义维数，等于白白损失了表达能力。这两种退化的共同点是：\emph{只盯着正对应距离（那一对比得像不像）都会漏掉它们}——因为不管哪种，正对应距离都可以很小。必须额外考察“跨样本、跨维度”的统计量才能发现。

\Needspace{6\baselineskip}
图~\ref{fig:data-derived-knowledge-representation-degeneration}用四个教学样本、在固定的二维坐标范围内并排画了三种状态，直观区分“配对成功”与“表示健康”。每个样本用同一种标记表示，实心点与空心圈代表它的两个视图：在三个面板里，每对的两个点都\emph{很接近}，所以\emph{配对误差在三种情况下都可以很小}——这正是只看配对会被骗的原因。

只有当我们转而比较\emph{跨样本}的统计量时，差别才显现：面板(a)所有点挤在一处（跨样本方差为零，常数坍塌）；面板(b)点虽拉开了、却都排在一条斜线上（两维协方差非零，维度退化）；面板(c)点在两个方向上都充分铺开（二维分散，健康）。看图时还应记住：右侧面板只是在\emph{这两个示意维度}上排除了简单退化，它\emph{并没有}给这四个点赋予任何真实的花卉类别——“分散”不等于“分得有语义意义”。

\begin{figure*}[!t]
  \centering
  \begin{tikzpicture}[
  x=1mm,y=1mm,
  font=\figurefont\scriptsize,
  panel/.style={draw=BookRule,fill=NNPanel,line width=.85pt,rounded corners=1.5mm},
  axis/.style={knowledge arrow,draw=BookMuted,line width=.8pt},
  viewone/.style args={#1}{circle,fill=#1,draw=#1,minimum size=3.2mm,inner sep=0pt},
  viewtwo/.style args={#1}{circle,fill=white,draw=#1,line width=1pt,minimum size=3.8mm,inner sep=0pt},
  samplelabel/.style={font=\figurefont\scriptsize,text=BookInk},
  note/.style={font=\figurefont\scriptsize,text=BookMuted,align=center}
]
  % 所有面板用同一映射：(z_1,z_2) -> (中心横坐标+8z_1,37+8z_2) mm。
  % 坐标范围均为[-2,2]^2；点位直接使用矩阵中的教学数值。
  \draw[panel] (.5,.5) rectangle (51.5,66.5);
  \draw[panel] (58.5,.5) rectangle (110.5,66.5);
  \draw[panel] (117.5,.5) rectangle (168.5,66.5);
  \foreach \cx in {26,84.5,143}{
    \begin{scope}[shift={(\cx,37)},x=8mm,y=8mm]
      \draw[axis] (-2,0) -- (2,0) node[right] {$z_1$};
      \draw[axis] (0,-2) -- (0,2) node[above] {$z_2$};
      \draw[BookMuted,line width=.6pt] (1,-.04) -- (1,.04)
        node[below=.4mm] {$1$};
      \draw[BookMuted,line width=.6pt] (-.04,1) -- (.04,1)
        node[left=.4mm] {$1$};
    \end{scope}
  }
  \node[font=\figurefont\bfseries\footnotesize,text=BookInk] at (26,62) {(a) 常数表示};
  \node[viewtwo=BookAmber,minimum size=7mm] at (26,37) {};
  \node[viewtwo=NNHidden,minimum size=5.8mm] at (26,37) {};
  \node[viewtwo=BookTeal,minimum size=4.6mm] at (26,37) {};
  \node[viewone=BookBlue] at (26,37) {};
  \node[samplelabel,anchor=west] at (30,40) {$A=B=C=D$};
  \node[note] at (26,12) {$\bm{Z}=\bm{Z}'=\begin{bmatrix}0&0\\0&0\\0&0\\0&0\end{bmatrix}$\\配对误差为$0$，各维方差为$0$};

  \node[font=\figurefont\bfseries\footnotesize,text=BookInk] at (84.5,62) {(b) 单方向变化};
  \begin{scope}[shift={(84.5,37)},x=8mm,y=8mm]
    \draw[BookRule,line width=.6pt,dashed] (-1.7,-1.7) -- (1.7,1.7);
    \foreach \x/\y/\sampleid/\samplecolor in {-1.5/-1.5/A/BookBlue,-.5/-.5/B/BookTeal,.5/.5/C/NNHidden,1.5/1.5/D/BookAmber}{
      \node[viewtwo=\samplecolor] at ({\x+.08},{\y-.05}) {};
      \node[viewone=\samplecolor] at (\x,\y) {};
      \node[samplelabel,anchor=south east,xshift=-1mm,yshift=1mm] at (\x,\y) {\sampleid};
    }
  \end{scope}
  \node[note] at (84.5,10.5) {$\bm{Z}\approx\bm{Z}'=\begin{bmatrix}-1.5&-1.5\\-.5&-.5\\.5&.5\\1.5&1.5\end{bmatrix}$\\样本有方差，但两个坐标重复同一变化};

  \node[font=\figurefont\bfseries\footnotesize,text=BookInk] at (143,62) {(c) 示意维度中分散};
  \begin{scope}[shift={(143,37)},x=8mm,y=8mm]
    \foreach \x/\y/\sampleid/\samplecolor in {-1/-1/A/BookBlue,-1/1/B/BookTeal,1/-1/C/NNHidden,1/1/D/BookAmber}{
      \node[viewtwo=\samplecolor] at ({\x+.08},{\y-.05}) {};
      \node[viewone=\samplecolor] at (\x,\y) {};
      \node[samplelabel,anchor=south east,xshift=-1mm,yshift=1mm] at (\x,\y) {\sampleid};
    }
  \end{scope}
  \node[note] at (143,10.5) {$\bm{Z}\approx\bm{Z}'=\begin{bmatrix}-1&-1\\-1&1\\1&-1\\1&1\end{bmatrix}$\\两维有方差且示例协方差为$0$};

  \node[viewone=BookInk] at (5,-5) {};
  \node[note,anchor=west] at (8,-5) {视图A};
  \node[viewtwo=BookInk] at (37,-5) {};
  \node[note,anchor=west] at (40,-5) {对应视图B};
  \node[note,anchor=west,text=FigureLoss] at (76,-5) {教学构造；分散不等于形成有用类别};
\end{tikzpicture}

  \caption{低配对误差没有排除样本或维度退化。相同标记代表同一样本，实心点与空心圈代表其两视图，三个面板使用相同的坐标范围与单位刻度。(b)的数值跨度更大，却仍只沿一条方向变化。手工矩阵依次给出常数、重复方向与二维分散；分散只描述示意坐标，不代表真实语义类别。}
  \label{fig:data-derived-knowledge-representation-degeneration}
\end{figure*}

逐一看每种方法在防坍塌上的真实保证。对\emph{SimCLR}，前面已算过：常数表示时的损失是$\log(2n-1)$，而存在有效负对照、且模型\emph{能实现}配对优势时，可以把损失降得更低——所以从\emph{目标取值}上，它\emph{偏好}区分而非常量。但正如5.2.2.1强调的，\emph{精确的}常数状态仍是一个梯度为零的驻点，目标偏好区分\emph{不等于}优化必然能逃离常数。

更要警惕SimCLR区分性的\emph{内容}：它要求的只是“在\emph{当前这组正负关系}中排出相对名次”。如果负对照集合里有错误负例，它就会\emph{排斥语义上其实相近的样本}；而如果样本间的背景差异比花朵差异更显著，模型完全可以\emph{靠背景}就取得很好的排序优势。因此SimCLR的“区分偏好”\emph{既不保证}优化会逃离常数，也\emph{没有}对“区分的必须是物种语义”作任何直接约束。

对\emph{BYOL}，防坍塌不能归因于任何单一零件，而要看\emph{预测头、停止梯度、移动平均目标以及优化动态共同组成的整套非对称训练}。原论文中完整的BYOL\emph{没有使用显式负例}，它通过一系列消融实验来研究这些组件各自的作用\scite{ka-grill2020byol}。

这里要纠正两个常见的错误归因。其一，\emph{不能}因为用了停止梯度，就推断“所以BYOL不会坍塌”——停止梯度只阻断当前这一步的求导，它本身是个中性的操作，\emph{单凭这一事实}推不出任何防坍塌保证。其二，EMA是把\emph{旧目标与新在线参数}做加权平均，它\emph{并不是}沿着“目标损失的梯度”去更新目标（目标根本没有损失梯度），两者的更新逻辑完全不同。结论是：停止梯度和EMA各自有明确作用，但BYOL能否稳定不坍塌，是\emph{整套训练安排（连同具体优化器动态）}的性质，不能把成功简单归功于某一个零件。

对\emph{DINO}，由于它操作的是$K$维\emph{分布}而非连续向量，诊断退化的对象也要相应改变，此时有两种对称的退化形态。一种是\emph{所有样本都给出近乎均匀的分布}（每个维度都约$1/K$，样本之间无法区分）；另一种是\emph{所有样本都集中到同一个维度}（大家都压向一维，同样无法区分）。前面介绍的两个操作恰好各防一种：\emph{教师锐化}（低温）抵抗前者（不允许分布太平），\emph{居中}（减去中心）抵抗后者（不允许一维独大），二者在动量教师的条件下配合使用\scite{ka-caron2021dino}。

要注意这与连续向量的“重合/恒定”\emph{不是同一个诊断口径}：判断DINO有没有退化，必须直接去\emph{查看投影分布}的形状（是否过平或过尖），\emph{不能}只凭一张特征散点图就下结论排除退化。

对\emph{VICReg}，可以直接用图中的三个教学矩阵验证它的两项附加计分是否真能抓住对应退化（均用分母$n-1$计算，条件为$n\geq2$、$\gamma>\sqrt{\epsilon}$）。

面板(a)是常数状态：所有样本相同，方差项$\max(0,\gamma-\sqrt{\epsilon})$为正、施以惩罚，从而挡住完全坍塌。面板(b)沿单一方向变化：计算其协方差矩阵，各元素为$5/3$，非对角元素非零，于是协方差项给出非零惩罚$c(\bm{Z})=25/9$，抓住了“两个坐标重复同一方向”。面板(c)在二维上充分分散，其协方差矩阵恰好是$(4/3)\bm{I}_2$（$\bm{I}_2$为二维单位矩阵），非对角元素全为0，因此$c(\bm{Z})=0$，不罚。

这组计算\emph{验证了}VICReg的计分确实能\emph{区分}“常数、单方向冗余、二维分散”三种状态。但要守住边界：\emph{能在目标函数里给退化打分，并不证明从任意初始化出发的优化动力学都一定能避开退化}——前面已经看到精确常数处方差梯度为零，打分与“优化一定躲开它”仍是两回事。

即使把上述数学退化都排除了，还要回到花卉任务追问一个更本质的问题：\emph{这些“健康的差异”就一定是我们想要的吗}？未必。较高的方差，可能\emph{只来自照明}条件的变化（白天/黄昏），而非花的种类；较低的协方差，也可能只是\emph{背景特征}恰好彼此去相关，与花瓣形态无关。表示在统计上“充分分散”，并\emph{没有}证明其中\emph{保留了花瓣的鉴别特征}——它完全可能在一组与任务无关的维度上分散得很漂亮。

因此严谨的检查至少要分三层、分别观察：\emph{匹配误差}（正对应是否一致）、\emph{每维方差与协方差（或有效秩）}（差异是否存在、维度是否冗余），以及\emph{独立下游任务上的指标}（这些差异到底有没有用）。其中\emph{有效秩}是一个有用的概括量，它近似刻画表示\emph{实际利用了多少个}变化方向（而非名义维度数）。三层的分工要记清：前几项回答“差异\emph{存不存在}”，最后一项回答“这些差异\emph{有没有用}”；\emph{训练总损失一个数字，替代不了这三层检查}。

\section{基于数据分组的知识获取}
\label{sec:data-derived-knowledge-grouping}

前两条路线处理的都是\emph{成对或局部}的关系：5.1逐个位置恢复，5.2大多约束一对样本。但如果想让\emph{一整批}照片共享某种更高层的结构目标，只保存若干孤立的样本对就不够了——成对关系\emph{无法决定整个批次该如何划分成群}。

于是有了第三条路线：先根据当前表示的相似性，把一批样本\emph{组织成若干群组}，再让模型去预测“每张照片属于哪个组”（组号）或“它对各组的归属权重”。这种\emph{由当前表示、比较规则和分组配置临时生成}的答案，称为\term{分组目标}（grouping target）。

与前两条路线相比，它有两个关键变化。其一，\emph{作用范围}从“一对样本”扩大到“一批样本”——目标取决于整批数据的全局组织。其二，\emph{目标的身份}变成了\emph{算法临时产生的组号}，这带来一个必须时刻记住的特性：\emph{组号没有预设的物种含义，而且不稳定}。这一轮叫“组1”的簇，下一轮重新聚类后可能改名叫“组3”；两张照片同组只说明它们在当前表示下相似，\emph{不}说明它们是同一个物种。

从数据流上看清分组目标的产生顺序：原始观测\emph{先}经过一个表示映射变成特征，\emph{再}由分组规则在这些特征上产生组号或归属权重；模型则读取\emph{原始输入}去预测这个被构造出来的目标。这是一个带有“循环”色彩的过程——用来\emph{生成目标}的表示，本身又是模型\emph{正在学习}的东西，因此目标会随表示一起变化。

必须强调其中的假设链条：分组结果\emph{依赖}于所用的特征、相似度定义和结构假设；“样本\emph{能}被分成若干个组”这件事，\emph{完全没有}确定这些组\emph{就是}花卉的物种类别——它们完全可能是按背景颜色、拍摄地点聚成的。聚类相关的假设与求解方法可参考\BookChapterRef{chap:clustering}{卷四《模型：经典模型》的“聚类”章}。本章不再重复聚类算法本身，而是聚焦四个使用问题：分组\emph{用哪种表示}、产生的是\emph{硬编号还是软分配}、目标\emph{何时重算}、以及重算后如何\emph{同步}依赖它的预测接口。

图~\ref{fig:data-derived-knowledge-grouping-targets}固定\emph{同一批表示}，并排对照三种把它变成分组目标的机制：\emph{整体聚类}（直接在全部特征上聚成簇并给簇号）、\emph{代表向量（原型）分配}（维护少量代表点，把样本分配过去）、以及\emph{图连接}（先在样本间建图，再由连通结构形成组）。

三者并非互斥，它们可以\emph{共用同一个编码器}，也可以相互\emph{组合}：比如原型方法既用分组决定目标权重，又往往同时沿用第~\ref{sec:data-derived-knowledge-relations}节的跨视图正对应关系（如SwAV）。本章把它们分成三小节，只是\emph{按主要的目标构造机制}归类，并不意味着用了样本关系的方法就不属于分组——实际方法常常是多种机制的混合体。

\begin{figure*}[!t]
  \centering
  \begin{tikzpicture}[
  x=1mm,y=1mm,
  font=\figurefont\scriptsize,
  panel/.style={draw=BookRule,fill=NNPanel,line width=.85pt,rounded corners=1.5mm},
  source/.style={draw=FigureInput,fill=FigureInput!12,line width=1.1pt,minimum height=9mm,minimum width=28mm,align=center,inner sep=1.5mm},
  target/.style={draw=FigureOutput,fill=FigureOutput!10,line width=1.1pt,minimum height=8mm,minimum width=23mm,align=center,inner sep=1.5mm},
  model/.style={draw=FigureModel,fill=FigureModel!12,line width=1.1pt,minimum height=8mm,minimum width=23mm,align=center,inner sep=1.5mm},
  flow/.style={knowledge arrow,draw=NNWire,line width=1.1pt},
  refresh/.style={knowledge arrow,draw=BookTeal,line width=.9pt,dashed,rounded corners=1.5mm,preaction={draw=white,line width=2.5pt,solid}},
  edge/.style={draw=NNLine,line width=.8pt},
  warningedge/.style={draw=FigureLoss,line width=1.2pt,dashed},
  point/.style={circle,draw=BookBlue,fill=BookBlue!18,minimum size=3.4mm,inner sep=0pt},
  note/.style={font=\figurefont\scriptsize,text=BookMuted,align=center}
]
  \draw[panel] (.5,-34) rectangle (52.5,34.5);
  \draw[panel] (58.5,-34) rectangle (110.5,34.5);
  \draw[panel] (116.5,-34) rectangle (168.5,34.5);
  \node[source] (batch) at (84.5,57) {同一批未标注样本$\bm{X}$};
  \node[model] (repr) at (84.5,44) {共享编码器$f_{\bm{\theta}}$\\$\bm{Z}=f_{\bm{\theta}}(\bm{X})$};
  \draw[flow] (batch)--(repr);
  \draw[NNWire,line width=1.1pt] (repr.south) -- (84.5,38);
  \draw[NNWire,line width=1.1pt] (5,38) -- (164,38);

  \node[font=\figurefont\bfseries\footnotesize,text=BookInk] at (26.5,30) {(a) 聚类编号};
  \draw[flow] (5,38) |- (26.5,24);
  \foreach \x/\y in {10/16,14/20,18/14,34/17,39/21,43/14}{\node[point] at (\x,\y) {};}
  \draw[BookBlue,line width=.8pt,dashed] (14,17) ellipse (9mm and 7mm);
  \draw[BookTeal,line width=.8pt,dashed] (39,17) ellipse (9mm and 7mm);
  \node[text=BookBlue] at (14,8) {组1};
  \node[text=BookTeal] at (39,8) {组2};
  \node[target] (cluster-target) at (26.5,-5) {临时目标$c_i$};
  \node[model] (cluster-predict) at (26.5,-24) {预测组号并更新};
  \draw[flow] (26.5,7)--(cluster-target);
  \draw[flow] (cluster-target)--(cluster-predict);
  \draw[refresh] (cluster-predict.east) -- (54,-24) |- (repr.west);
  \node[note,text=BookTeal,rotate=90] at (47,-9) {可选重编码};

  \node[font=\figurefont\bfseries\footnotesize,text=BookInk] at (84.5,30) {(b) 原型软分配};
  \node[note,text=BookInk] (assignments) at (84.5,15) {$\begin{array}{c|cc}
    &\text{样本1}&\text{样本2}\\
    \bm{c}_1&0.690&0.310\\
    \bm{c}_2&0.310&0.690
  \end{array}$};
  \draw[flow] (61,38) |- (assignments.west);
  \node[target] (proto-target) at (84.5,-5) {分配向量$\bm{q}_{s,b}$};
  \node[model] (proto-predict) at (84.5,-24) {视图$t$预测并更新};
  \draw[flow] (assignments.south)--(proto-target.north);
  \draw[flow] (proto-target)--(proto-predict);
  \draw[refresh] (proto-predict.east) -- (112,-24) |- (repr.east);

  \node[font=\figurefont\bfseries\footnotesize,text=BookInk] at (142.5,30) {(c) 邻域／图分组};
  \draw[flow] (119,38) |- (142.5,26.5);
  \node[point,minimum size=4.3mm] (g1) at (125,17) {1};
  \node[point,minimum size=4.3mm] (g2) at (132,23) {2};
  \node[point,minimum size=4.3mm] (g3) at (136,12) {3};
  \node[point,minimum size=4.3mm] (g4) at (150,12) {4};
  \node[point,minimum size=4.3mm] (g5) at (156,23) {5};
  \node[point,minimum size=4.3mm] (g6) at (163,17) {6};
  \draw[edge] (g1)--(g2);
  \draw[edge] (g2)--(g3);
  \draw[edge] (g3)--(g1);
  \draw[edge] (g4)--(g5);
  \draw[edge] (g5)--(g6);
  \draw[edge] (g6)--(g4);
  \draw[warningedge] (g3.east)-- node[above=1mm,text=FigureLoss] {桥接} (g4.west);
  \node[target,text width=38mm,inner sep=1mm] (graph-target) at (142.5,-5) {前：$(1,1,1,2,2,2)$\\后：$(1,1,1,1,1,1)$};
  \node[model] (graph-predict) at (142.5,-24) {预测目标并更新};
  \draw[flow] (142.5,7)--(graph-target);
  \draw[flow] (graph-target)--(graph-predict);
  \draw[refresh] (graph-predict.east) -- (167,-24) |- (repr.north east);

  \node[note,anchor=north west,text width=166mm,align=left] at (1,-38) {原型面板只展开两个样本，列内权重和为1；图面板展开六个节点。数值均为教学构造，分组和原型索引不具有预设物种含义。};
\end{tikzpicture}

  \caption{由同批表示生成三种分组目标。聚类给临时组号，原型面板的两列为第~\ref{sec:data-derived-knowledge-prototype-targets}节算例的归一化分配，图面板六个节点用于第~\ref{sec:data-derived-knowledge-graph-targets}节的连通计算；两列与六节点不是同一展开规模。赭色虚线表示拟加入的错误边$(3,4)$，青色回线表示更新后重算目标，不表示向分组求解器回传梯度。}
  \label{fig:data-derived-knowledge-grouping-targets}
\end{figure*}

表~\ref{tab:data-derived-knowledge-grouping}提供了读这三种机制的三个角度：各自\emph{构造出什么目标}（硬簇号、软权重、图上组关系）、依赖\emph{多大的统计范围}（全量样本、一个小批量、局部邻域或整张图）、以及\emph{刷新时需要同步什么}。

最后一点尤其关系到工程正确性：三种机制所需的样本范围差别很大——全量聚类往往要一次性看到所有（或大量）样本，原型分配常在一个小批量内完成（可辅以队列），图方法则依赖局部近邻或全局图。而\emph{一旦重新计算分组}，所有\emph{依赖旧分组的下游部件}都要跟着检查：簇编号可能重排、原型位置移动、图的邻接关系改变，连\emph{预测输出头}（输出维度对应组数）都可能需要重新对应。如果只图省事\emph{替换一个目标文件}而不做这些同步，就可能留下一个“仍在输出旧含义组号”的预测头，造成隐蔽的版本错乱。

\begin{table*}[htbp]
  \centering
  \renewcommand{\arraystretch}{1.15}
  \begin{tabularx}{\linewidth}{@{}>{\RaggedRight\setlength{\parindent}{0pt}\arraybackslash}p{25mm}*{3}{>{\RaggedRight\setlength{\parindent}{0pt}\arraybackslash}X}@{}}
    \toprule
    分组方式 & 构造的目标 & 统计范围 & 刷新时同步对象 \\
    \midrule
    聚类分组 & 硬簇编号$c_i$ & 常需全样本或较大样本 & 簇中心、样本目标、预测头与编号对应 \\
    原型分配 & 硬编号或软分配$\bm{q}_i$ & 小批量，可辅以表示队列 & 原型、分配结果与队列 \\
    邻域／图分组 & 组号或组内关系 & 局部近邻或全局图 & 邻接表、分组结果与训练目标 \\
    \bottomrule
  \end{tabularx}
  \caption{比较分组结果及其刷新所依赖的状态。三种机制可以组合，硬编号和软权重都由算法构造，没有预设类别含义；统计范围和同步对象不同，不能用一套刷新规则代替全部机制。}
  \label{tab:data-derived-knowledge-grouping}
\end{table*}

% 首段的DeepCluster边注从段末开始，避免贴近纸张底边。
\Needspace{65mm}
\subsection{基于聚类的分组目标构造}
\label{sec:data-derived-knowledge-cluster-targets}

最直接的分组目标来自\emph{聚类}。对一批没有物种标签的照片，先用模型\emph{当前}提取出的特征把它们分组，于是每张照片都\emph{暂时}获得了一个组号，这个组号就被拿来当训练目标。具体地，先提取特征$\bm{z}_i=f_{\bm{\theta}}(\bm{x}_i)$，按某种聚类规则把它分到$K$个簇之一，得到组号$c_i\in\{1,\ldots,K\}$；然后反过来，从\emph{原始输入}$\bm{x}_i$训练模型去预测这个$c_i$。

DeepCluster是这种“自举式”交替机制的代表：它用$k$均值对网络特征聚类、再用聚类得到的簇分配当作标签去监督网络更新\scite{knowledge-acquisition-deepcluster2018}。注意这里形成了一个\emph{自我强化的循环}——网络产生特征、特征决定簇号、簇号又反过来训练网络，下一轮网络特征因此改变、簇号也可能重排。

实现时，组号可以直接接到一个普通的\emph{分类式损失}（像预测$K$个类别那样去预测$K$个簇）。但这是一个\emph{形式上的借用}，必须反复强调：\textbf{用了分类损失，并没有把组号变成物种类别标签}。二者的本质区别在于编号的来源——真正的类别标签由外部世界固定（山茶就是山茶），而簇编号完全由“\emph{当前这一版表示}”和“\emph{当前的聚类配置}”决定，表示一变、重跑一次聚类，编号就可能整个重排。

下面把这一类交替训练抽象成三个清晰步骤：\emph{提取}特征、\emph{分组}得到编号、再\emph{预测}编号。围绕这个框架还有三个使用时绕不开的工程问题：跨轮次的\emph{编号如何对齐}、\emph{多久刷新}一次分组、以及如何\emph{检查}分组是否失效。这些是使用该机制时需要自行处理的，不应把它们都归功或归咎于某一篇代表性论文。

在开始这种交替训练前，必须先做一个策略选择：\emph{生成出来的组号要不要、以及多久刷新一次}？这直接影响训练的稳定性。

\emph{固定分组}：在开始时只聚类一次得到$c_i$，之后无论网络怎么变，始终预测这一套固定编号。它简单稳定，但有个明显问题——网络越训越好，可它追逐的目标却是用“最初那个很差的表示”聚出来的簇，目标很快就过时了。\emph{周期性分组}：每隔一段就用\emph{当前}较新的表示重新聚类一次、刷新编号；它能让目标跟上表示的进步，但为了能把“生成目标”和“拟合目标”两件事分开追溯，必须\emph{逐轮记录}当时使用的模型参数版本、生成的编号、以及训练样本集。对固定的$n$个样本，这个交替过程形式化地写为
\[
  \begin{aligned}
    \{\bm{z}_i^{(r)}=f_{\bm{\theta}^{(r)}}(\bm{x}_i)\}_{i=1}^{n}
    &\ \longrightarrow\
    \bm{c}^{(r)}
    =\operatorname{Cluster}(\{\bm{z}_i^{(r)}\}_{i=1}^{n}),\\
    \mathcal{D}^{(r)}=\{(\bm{x}_i,c_i^{(r)})\}_{i=1}^{n}
    &\ \longrightarrow\
    \bm{\theta}^{(r+1)}=\operatorname{Train}(\mathcal{D}^{(r)}),
  \end{aligned}
\]
读这个交替式要注意几个要点。上标$r$编号的是\emph{分组轮次}（重新聚一次类算一轮），\emph{不是}单个梯度下降步——一轮分组之后通常要做很多步梯度训练。$\operatorname{Cluster}$负责按\emph{当前}特征生成一轮编号，$\operatorname{Train}$负责用本轮的输入—编号对去拟合网络。

这个“先训练、后重分组”的顺序意味着：\emph{前一轮训练改变了表示，下一轮重新分组就可能改写答案}——同一张照片这一轮在簇2、下一轮可能跑到簇5，这是周期性分组的固有现象。刷新频率构成一对权衡：刷新\emph{越频繁}，目标相对表示的\emph{滞后越小}（目标更新、更贴合新表示），但每轮都要对全部样本\emph{重新编码、重新聚类、并重新存储目标}，计算和存储开销显著增大。

刷新分组时最容易犯的错误，就是\emph{把簇的数字编号当成群组的固定身份}。即使背后照片的\emph{几何分簇完全没有变}，聚类算法这一轮也可能仅仅因为初始化或遍历顺序不同，就把原来的“组1”\emph{改名}成“组3”——内容相同、标签名却换了。

因此跨轮\emph{保留预测头}（那个输出$K$个簇分数的分类层）时，必须先解决编号对齐问题：要么按\emph{样本重叠}或\emph{簇中心匹配}，把新旧两套编号对应起来；要么干脆\emph{重新初始化}这个预测输出层，不做继承。更麻烦的是簇还可能发生\emph{分裂}（一个簇裂成两个）或\emph{合并}（两个簇并成一个），此时新旧簇之间\emph{根本不存在一一对应}，更\emph{不能}仅仅因为“数字恰好相同”（新旧都有个“组2”）就把旧的输出参数错误地继承过去。

聚类是“无监督”的，它只负责让“簇内相似、簇间相异”，但\emph{相似的维度由特征和算法决定，并不指向物种}。同一批照片，完全可能这一次按\emph{背景颜色}聚成组（亮背景一组、暗背景一组），下一次按\emph{花瓣形态}聚成组；这两种划分在\emph{各自采用的聚类计分}下都可能是“好的聚类”，但它们为物种识别提供的鉴别特征\emph{天差地别}——前者几乎没用，后者才可能有用。

有大量因素会改变最终的簇：簇数$K$怎么定、特征用什么尺度、有没有异常点、各组大小是否均衡、以及\emph{初始表示本身的偏向}（早期网络可能主要响应背景）。更要警惕自举循环的副作用：模型被训练去\emph{预测}这些未必合任务需要的簇，下一轮特征就会\emph{更加迎合}这种分组，从而\emph{强化}一个与物种无关的划分，错误被自我固化。

由此得到使用聚类目标时最重要的一条判据：\emph{聚类损失低、或模型能准确预测簇号，只说明它很好地形成或复现了“当前这一版分组”}，仅此而已。这既不证明分组对应物种，也不证明特征有用。

要想\emph{把某个簇解释成一个有意义的任务类别}，必须有\emph{额外}的证据支撑，证据来源可以是：少量\emph{外部标签}的核对、\emph{领域专家}的解释、分组在\emph{独立样本}或多次运行下的\emph{稳定性}（真正的结构应可重复）、或在\emph{下游任务}上的实际效用。在取得这些证据之前，簇编号\emph{始终只是}一个临时的训练目标，不能当作知识结论使用。

\subsection{基于原型的分配目标构造}
\label{sec:data-derived-knowledge-prototype-targets}

聚类要对\emph{所有样本}计算群组边界，规模一大开销很高。原型方法换了一个更轻量的思路：不直接维护“样本之间谁和谁一类”的完整边界，而是只维护\emph{少量有代表性的向量}，让每张照片去和这些“代表”比较、决定归属。

这些可更新的代表向量$\bm{c}_1,\ldots,\bm{c}_K$就称为\emph{原型}（prototype），你可以把它们想象成$K$个“虚拟中心样本”。对一张照片的表示$\bm{z}_i$，有两种分配方式：找出\emph{最近}的那个原型、给一个\emph{硬编号}（非此即彼）；或者不做硬性选择，而是给出它对各个原型的\emph{支持权重}，形成一个\emph{软分配}$\bm{q}_i\in\mathbb{R}^{K}$（各权重非负、总和为1，表示“多大程度属于各组”）。

要立刻澄清一个常见误读：这些软权重描述的只是\emph{算法把样本分配给了哪些代表}，它们\emph{不会自动成为}“该样本属于某物种类别的后验概率”——原型没有物种含义，权重自然也不具备概率语义。

SwAV把“原型分组”和5.2的“正对应关系”结合了起来：它\emph{不再}直接给样本一个固定簇号，而是要求\emph{同一张照片的不同裁剪，能够互相预测对方的原型分配}——沿用“两裁剪是正对应”的关系，但把预测的\emph{答案}从像素或向量，换成了\emph{分组权重}。

计算上，先把表示和原型都固定为\emph{单位长度}（便于用内积衡量方向一致性），用\emph{内积}计量一个表示对每个原型的“支持程度”，再在$K$个原型上做指数化和归一化，得到一个可以放进交叉熵的\emph{预测分布}。对同一张图的两个视图$t,s$，其归一化表示记为$\bm{z}_t,\bm{z}_s$，待构造（目标）的分配记为$\bm{q}_t,\bm{q}_s$；在温度$\tau>0$下，一个视图表示$\bm{z}$“预测原型$k$”的概率为
\[
  p_k(\bm{z})
  =\frac{\exp(\bm{z}^{\top}\bm{c}_k/\tau)}
  {\sum_{r=1}^{K}\exp(\bm{z}^{\top}\bm{c}_r/\tau)}.
\]
这里的$p_k$只是模型\emph{在线给出的预测}，我们\emph{还没有}规定作为\emph{目标}的那个分配$\bm{q}$要怎样产生（这是SwAV最核心、稍后专门讨论的部分）。暂时假设目标分配已知：把\emph{另一视图}的归一化分配$\bm{q}$当作这一视图\emph{当步}的答案，按这$K$个原型上的目标权重计算负对数（交叉熵）计分。

要正确理解这个交叉熵在鼓励什么：在目标分配固定时，减小交叉熵是让\emph{整个预测分布去逼近}给定的目标分配（目标权重大的原型概率升高、权重小的降低），\emph{并不是}要求每个原型的预测概率都同时增大。再把$t\to s$和$s\to t$两个方向的项交换相加，得到SwAV的交换损失：
\begin{equation}
  \mathcal{L}_{\mathrm{swap}}
  =-\sum_{k=1}^{K}q_{s,k}\log p_k(\bm{z}_t)
   -\sum_{k=1}^{K}q_{t,k}\log p_k(\bm{z}_s).
  \label{eq:data-derived-knowledge-swav-swap}
\end{equation}
这个“交换”体现在：第一项中$\bm{q}_s$是从\emph{$s$侧}生成的分配，却拿来核对\emph{$t$侧}的预测$p_k(\bm{z}_t)$；第二项正好反过来。它利用的正是“两视图来自同一张图、应有相同分组”这一正对应关系。注意两项用的是\emph{求和}（而非对两个方向取平均），所以在与其他方法比较损失数值时，求和口径会让数值相差一个倍数，必须先固定\emph{是否对方向再取平均}\scite{ka-caron2020swav}（后面的数值算例会具体展示这个差异）。

现在回到SwAV\emph{最关键}的问题：作为目标的分配$\bm{q}$究竟\emph{怎样构造}？为什么不能图简单、让每个样本\emph{独立}地选最近原型？

如果每个样本都独立地归给最近的原型，会发生一种\emph{塌缩}：由于所有样本都在追求“最匹配”，它们很可能\emph{一窝蜂地集中到同一个}（或少数几个）原型上——那个原型短暂地最受欢迎，于是所有分配都倒向它，其余原型形同虚设。为避免这种“赢家通吃”，SwAV在\emph{同一个分配批次内}加了一个\emph{均衡约束}：要求$K$个原型\emph{分担}总的分配权重（每个原型大约分到$1/K$的样本），同时又\emph{保留}“应当尽量相似才分配过去”的要求。

这就不能逐样本独立决策了，而要\emph{联合}求解一个分配方案。用一个非负矩阵记录“每个样本向每个原型分配多少权重”，在\emph{满足行、列总量约束}的候选矩阵中挑出整体最优的一个，而非每个样本各自取最大分。先约定记号（注意与VICReg的排法\emph{不同}，避免混淆）：对正整数$B,K$，把$B$个样本\emph{按列}排成$\bm{Z}=[\bm{z}_1,\ldots,\bm{z}_B]\in\mathbb{R}^{d\times B}$，把$K$个原型\emph{按列}排成$\bm{C}=[\bm{c}_1,\ldots,\bm{c}_K]\in\mathbb{R}^{d\times K}$；这里的$\bm{C}$是\emph{原型矩阵}，而VICReg里的$\bm{C}$是协方差矩阵，二者只是同符号、含义完全不同。

最优分配要同时兼顾两个目标，方向恰好相反：一是\emph{相似度总分}越大越好（内积加权越大越匹配）；二是\emph{熵正则}（entropy regularization）额外\emph{奖励不过分集中}的软权重（熵越大、分配越均匀越好），让“匹配相似度”和“分配的集中/均匀程度”\emph{共同}参与选择。SwAV在下面的允许集合$\mathcal{Q}$内\emph{最大化}这二者之和：
\begin{align}
  \max_{\bm{Q}\in\mathcal{Q}}\quad&
  \operatorname{Tr}(\bm{Q}^{\top}\bm{C}^{\top}\bm{Z})
  +\varepsilon H(\bm{Q}),\\
  \mathcal{Q}
  ={}&
  \left\{
    \bm{Q}\in\mathbb{R}_{+}^{K\times B}
    \ \middle|\
    \begin{aligned}
      \bm{Q}\bm{1}_B&=\frac{1}{K}\bm{1}_K,\\
      \bm{Q}^{\top}\bm{1}_K&=\frac{1}{B}\bm{1}_B
    \end{aligned}
  \right\}.
  \label{eq:data-derived-knowledge-swav-assignment}
\end{align}
逐项解释这个优化。第一项$\operatorname{Tr}(\bm{Q}^{\top}\bm{C}^{\top}\bm{Z})$展开后就是$\sum_{k,b}Q_{k,b}\bm{c}_k^{\top}\bm{z}_b$：每个样本$b$向原型$k$分配的权重$Q_{k,b}$，乘上二者的相似度（内积）$\bm{c}_k^{\top}\bm{z}_b$，再全部加总——它直接奖励“把权重大的分配放在确实相似的样本—原型对上”。第二项是熵$H(\bm{Q})=-\sum_{k,b}Q_{k,b}\log Q_{k,b}$（约定$0\log0=0$），系数$\varepsilon>0$控制它相对相似度项的权重；熵越大表示分配越均匀、越不集中。

下面集合$\mathcal{Q}$里的等式约束是SwAV均衡思想的核心：要求\emph{每个原型（每行）的权重之和为}$1/K$、\emph{每个样本（每列）的权重之和为}$1/B$，整个矩阵总和为1。理解这两个约束：行和约束强制\emph{各原型被均匀使用}（避免全挤到一个原型），列和约束则保证\emph{每个样本的总分配权重相同}。务必强调：这是\emph{人为施加的均衡先验}，\emph{不是}从数据里观测到的真实类别频率——现实中各类别本就不均衡，是SwAV为了防塌缩而\emph{假设}训练时等频使用原型。

由于这些约束在数学上是等式、在数值上只能近似满足，还有一个归一化细节：矩阵里每列的和是$1/B$（而不是概率分布要求的1），所以\emph{在把列向量送进交换损失之前必须再归一化}：
\begin{equation}
  \bm{q}_b
  =
  \frac{\bm{Q}_{:,b}}
       {\bm{1}_K^{\top}\bm{Q}_{:,b}}
  =
  B\bm{Q}_{:,b},
  \qquad
  \bm{1}_K^{\top}\bm{q}_b=1.
  \label{eq:data-derived-knowledge-swav-code-normalization}
\end{equation}
式~\eqref{eq:data-derived-knowledge-swav-code-normalization}把第$b$列放大$B$倍、变成一个\emph{和为1}的\term{分配向量}（assignment vector），它才是该样本对$K$个原型的\emph{目标权重}$\bm{q}_b$。交换损失里出现的$\bm{q}_t,\bm{q}_s$始终指这种归一化后的向量。术语上，原论文把它称为“code”，要注意这里的code是“编码/分配码”的意思，\emph{并非程序代码}。如果偷懒直接使用矩阵里和为$1/B$的原始列，交叉熵会额外被缩小一个$1/B$的因子，算出来的就不再是前面定义的那个归一化目标，数值和梯度都会失真。

上面那个带边际约束的联合优化，还需要一个\emph{实际可计算的求解方法}。SwAV借用了一个经典的矩阵平衡算法。先把“样本—原型”的匹配分数记为$\bm{F}=\bm{C}^{\top}\bm{Z}$（每个元素是一个原型与一个样本的内积）。熵正则项在这里带来一个实际好处：它使最优分配可以\emph{从正的指数权重出发}（因为$\exp$恒正，所有初始权重都大于0），然后只需通过\emph{反复缩放行列}来满足边际总量，而不必求解复杂的离散分配。第一步，先对所有分数做逐元素指数化、并把\emph{总和归一化为1}：
\[
  Q_{k,b}^{(0)}
  =\frac{\exp(F_{k,b}/\varepsilon)}
  {\sum_{r=1}^{K}\sum_{a=1}^{B}\exp(F_{r,a}/\varepsilon)}.
\]
但这个初始矩阵只保证\emph{总和}为1，\emph{行和、列和}还\emph{未必}正好等于要求的$1/K$和$1/B$。\emph{Sinkhorn--Knopp迭代}用一个非常朴素的思想来修正：\emph{交替地、反复地}缩放——先统一缩放每\emph{行}，使每行总量变成$1/K$；这一步通常会破坏刚得到的列和，于是再统一缩放每\emph{列}，使每列总量变成$1/B$；列缩放又可能破坏行和……如此\emph{交替}进行，在所有权重为正、各分母非零的条件下，它会\emph{收敛}到同时满足两组边际的矩阵。用$t$编号缩放轮次，一轮（先行后列）写为
\begin{equation}
  \begin{aligned}
    R_{k,b}^{(t)}
    &=\frac{Q_{k,b}^{(t-1)}}{K\sum_{a=1}^{B}Q_{k,a}^{(t-1)}},\\
    Q_{k,b}^{(t)}
    &=\frac{R_{k,b}^{(t)}}{B\sum_{r=1}^{K}R_{r,b}^{(t)}}.
  \end{aligned}
  \label{eq:data-derived-knowledge-sinkhorn-scaling}
\end{equation}
由于列缩放会反过来破坏行和，单轮通常不够，需要多轮重复，直到行、列边际都足够接近目标。停止准则要设计得严谨：下面的教学算例约定，当“\emph{所有行和}偏离$1/K$”与“\emph{所有列和}偏离$1/B$”这两组偏差中的\emph{最大绝对值}都小于$10^{-10}$时才停止。为什么取\emph{最大}偏差而不是平均偏差？因为平均误差可能很小、却\emph{掩盖住某一行（列）严重不达标}——用最大值能保证\emph{每一条}边际都合格，不被其他达标项稀释。一旦列和在数值精度内达到$1/B$，就可以用$\bm{q}_b=B\bm{Q}_{:,b}$得到归一化的分配向量。工程上若为省时间\emph{固定只跑很少几轮}，则必须\emph{报告最终残留的行/列边际误差}，让别人知道这个分配离满足均衡约束还差多少。

用两个样本和两个原型，我们可以从分数一直复算到交换损失。以下全是教学构造，不是训练实验：$K=B=2$，单位原型$\bm{c}_1=(1,0)^{\top}$、$\bm{c}_2=(0,1)^{\top}$，视图$s$的单位表示为$\bm{z}_{s,1}=(1,0)^{\top}$、$\bm{z}_{s,2}=(0.8,0.6)^{\top}$，$\varepsilon=\tau=0.5$。据内积和指数初始化，得到
\[
  \bm{F}_s=\begin{bmatrix}1&0.8\\0&0.6\end{bmatrix},
  \qquad
  \bm{Q}^{(0)}\approx
  \begin{bmatrix}0.443462&0.297262\\0.060016&0.199260\end{bmatrix}.
\]
将初值各行先调到$0.5$，再将各列调到$0.5$，首轮结果分别为
\[
  \begin{aligned}
    \bm{R}^{(1)}&\approx
    \begin{bmatrix}0.299344&0.200656\\0.115738&0.384262\end{bmatrix},\\
    \bm{Q}^{(1)}&\approx
    \begin{bmatrix}0.360584&0.171525\\0.139416&0.328475\end{bmatrix}.
  \end{aligned}
\]
行缩放后的$\bm{R}^{(1)}$列和约为$0.415081,0.584919$，尚不均衡；列缩放后的$\bm{Q}^{(1)}$行和又变成约$0.532109,0.467891$。继续用未舍入值交替缩放，12轮后的最大边际误差约为$1.82\times10^{-11}$，达到本例停止条件，归一化列后得到
\[
  \bm{Q}_s\approx
  \begin{bmatrix}0.344987&0.155013\\0.155013&0.344987\end{bmatrix},
  \qquad
  \bm{q}_{s,1}\approx\begin{bmatrix}0.689974\\0.310026\end{bmatrix},
  \quad
  \bm{q}_{s,2}\approx\begin{bmatrix}0.310026\\0.689974\end{bmatrix}.
\]
这个算例最有教益的地方，是它展示了“\emph{联合均衡}”如何改变单个样本的目标，值得在此点破（详细矩阵数值见上文推演）。如果让两个样本\emph{各自独立}地按最大内积分配，它们\emph{都会}选原型1（因为两个样本与原型1的内积都不小）——这正是没有均衡时的“塌向一个原型”。但在SwAV的\emph{联合均衡}分配下，由于强制要求两个原型\emph{平均分担}权重，最终样本1主要归原型1（约0.69/0.31），而\emph{样本2反而更多地支持原型2}（约0.31/0.69）。

这个结果可能让人意外，但\emph{不是相似度算错了}：它是因为允许集合里\emph{加入了“批次内各原型均衡”这一全局要求}，一个样本的最优分配\emph{不再只由它自己决定}，而\emph{依赖于同一批里还有哪些别的样本}。这带来一个重要认识：你\emph{无法}仅凭某个样本自身和它的最近原型，就还原出SwAV给它的目标——同一张照片，换一个批次（同批其他样本不同），分配就可能改变。

保持样本1、2的正对应，再构造$t$侧表示$\bm{z}_{t,1}=(0.8,0.6)^{\top}$、$\bm{z}_{t,2}=(1,0)^{\top}$。数值位置虽然交换，样本身份没有交换，不能据此重配两侧。对$t$侧独立运行同一分配，得到$\bm{q}_{t,1}=\bm{q}_{s,2}$、$\bm{q}_{t,2}=\bm{q}_{s,1}$；模型对样本1的两侧预测概率为
\[
  \bm{p}(\bm{z}_{t,1})\approx(0.598688,0.401312)^{\top},
  \qquad
  \bm{p}(\bm{z}_{s,1})\approx(0.880797,0.119203)^{\top}.
\]
用另一侧分配核对这些概率，并采用自然对数，样本1的两向计分及其和为
\[
  \begin{aligned}
    -\bm{q}_{s,1}^{\top}\log\bm{p}(\bm{z}_{t,1})&\approx0.637025,\\
    -\bm{q}_{t,1}^{\top}\log\bm{p}(\bm{z}_{s,1})&\approx1.506877,\\
    \mathcal{L}_{\mathrm{swap},1}&\approx2.143902.
  \end{aligned}
\]
在这个对称构造里，样本2的两项恰好与样本1交换，因此对两个样本取平均、再对两个方向\emph{求和}，总损失仍是$2.143902$；但若在此基础上再对\emph{两个方向取平均}，数值就变成$1.071951$（恰好一半）。这两个数字\emph{描述的是同一个损失}，仅仅因为\emph{聚合口径}（对方向求和还是再平均）不同就相差一倍。这给所有跨方法比较提了个醒：\emph{在比较损失数值之前，必须先固定“对样本、对方向到底是求和还是取平均”的口径}，否则差一倍纯属口径差异、与方法优劣无关。

SwAV中“求解分配”和“更新网络”是\emph{两条不同的路径}，梯度怎么走需要特别清楚。在\emph{当前这一步}，Sinkhorn算出的目标分配$\bm{q}_t,\bm{q}_s$是\emph{停止求导}的——它们被当作固定的目标，编码器和原型只从交换损失里\emph{预测}$p_k(\bm{z})$的\emph{那一侧}接收梯度，梯度不会反向穿过Sinkhorn求解过程去“推着分配变”。

另一个更新发生在\emph{优化步之后}：原方法会把每个\emph{非零}原型列\emph{重新投影回单位长度}$\bm{c}_k\leftarrow\bm{c}_k/\lVert\bm{c}_k\rVert_2$（因为前面假设了原型是单位向量），然后再进入下一批计算分数。要准确区分这是一个“更新\emph{之后}施加的\emph{归一化约束}操作”，它\emph{并不是}对分配目标做反向传播。还要注意SwAV与BYOL的结构差异：SwAV\emph{没有}BYOL那种独立的、用EMA维护的目标网络；因此如果擅自让梯度\emph{穿过}Sinkhorn分配求解器、或\emph{省略}原型归一化，得到的都是\emph{另一条计算路径}、已不再是原方法。

最后要把两个层次的问题分开：“\emph{可行的分配是否存在}”与“\emph{分配能不能区分样本}”，它们不是一回事。先看可行性：对\emph{任意}正整数$K,B$，取所有$Q_{k,b}=1/(KB)$（每个位置都相等）就能\emph{满足全部边际约束}——也就是说，哪怕所有样本的表示\emph{完全相同}，这个软约束集合也\emph{不会为空}，均衡分配\emph{永远存在}一个（平凡的）解。

但这个平凡解恰恰暴露了真正的困难：当$K>1$时，均衡约束\emph{排除}了“所有样本集中到一个原型”的塌缩，却\emph{接受}了另一种无区分的状态——\emph{每个样本都均匀地分配给所有原型}（每个样本的权重都是$1/K$），这样谁也没被区分开。真正困难、也真正有用的，是既满足均衡、又\emph{能把不同样本区分开}的“近硬”分配。多样性不足时，强行均衡还有副作用：它可能\emph{迫使}本就相似的样本被\emph{分摊}到不同原型，或与“按相似度本应聚到一起”的分组\emph{冲突}。原论文用一个\emph{缓存前几批表示的队列}来扩大统计范围、补足小批量内可能缺少的多样性——这是为了让均衡有足够多样的对象可分，\emph{不是}在修补可行性；队列同时也带来“缓存的是旧表示”的问题，而熵权重$\varepsilon$过大则会\emph{过度偏好均匀}分配、加剧“全摊平”。

SwAV的目标还会随一整套超参数改变：原型个数$K$、相似度温度$\tau$、熵权重$\varepsilon$、均衡的严格程度、以及批量由哪些样本组成。要公平比较“最近原型硬编号”与“软分配”两种做法，应当\emph{固定原型和表示不变}、只换分配规则，才能把观察到的差别干净地归因于\emph{分配方式}本身——硬分配删除了“一个样本支持多个原型”的信息、给出更脆的决策，软分配保留了这种不确定性，却增加了存储和交叉熵计算量。最后再次声明：本书的SwAV采用上面给出的\emph{严格边际要求}；其他实现若放宽了均衡约束，应当\emph{另行说明}，不能默认等价。而无论怎么分配，\emph{任何原型编号都不预设物种含义}。

作为本小节收尾，把SwAV的训练依据\emph{拆成两个可分别追溯的来源}，这也呼应本章开头的三条路线：其一，“同一张图两个裁剪应预测相同分组”的\emph{正对应关系}，来自\emph{样本变换}（5.2的路线）；其二，作为答案的\emph{分组分配}，来自\emph{带均衡约束的原型求解}（5.3的路线）。所以单看SwAV，它同时是“非对比式的匹配”（描述它\emph{怎样使用}正对应关系、不需要负对照）和“原型式的分组”（描述它\emph{怎样产生}目标分配）。只用其中任何一个名称，都无法完整说明这条训练路径——这正体现了不同机制在实际方法中的交叉。

\subsection{基于邻域与图结构的分组目标构造}
\label{sec:data-derived-knowledge-graph-targets}

原型方法用\emph{少量代表点}概括全部样本，而图方法提供了另一种组织方式：不维护少数中心，而是\emph{先保存每张照片与哪些照片相近}（即样本两两之间的连接关系），再\emph{从这些连接结构中自然地形成群组}。它不依赖“存在几个有代表性的中心”，而是让群组从局部连接的\emph{拓扑结构}里涌现。

具体地，按\emph{$k$近邻}（每个样本连最近的$k$个）、\emph{半径阈值}（距离小于阈值就连）或\emph{已有关系}（如同株、采集关联）建立一张图$\mathcal{G}=(V,E)$：\emph{顶点}对应样本，\emph{边}表示某种规则认定的“邻近”。这里符号$V$的含义相比前几章\emph{变了}——它现在指\emph{样本顶点}的集合，而\emph{不再是}5.1里那个“观测位置”的集合，阅读时注意区分。建图之后，最简单可以直接用\emph{邻居集合}形成局部目标；更进一步，则用\emph{连通分量、谱划分或社区发现}等方法，把这些\emph{局部的两两联系汇总}成\emph{整体的分组编号}。

从“局部的边”走向“全组的目标”，中间隐含了一个必须审视的\emph{传递规则}，这是图分组最主要的风险来源。以最直接的\emph{连通分量}为例：它默认“邻近”是\emph{可传递}的——只要$A$连$B$、$B$连$C$，就把$A$、$C$也\emph{视为同组}。但“近邻相似”在现实中\emph{并不保证传递}：$A$与$B$很像、$B$与$C$很像，\emph{推不出}$A$与$C$还共享同一语义（相似性沿链条会逐渐漂移）。

因此，图~\ref{fig:data-derived-knowledge-grouping-targets}(c)中那条\emph{赭色的“桥接边”}$(3,4)$特别危险：仅\emph{一条}这样的边，就可能把\emph{两个本来各自密集、内部一致的群组合并成一个}大组。相比“一对样本配错”，图结构会\emph{放大}错误——\emph{单条边的错误，其影响不再只限于这两个端点}，而会沿连通路径波及整条链、整片区域。下面用一个教学算例定量演示这种放大。

固定规则可以复算这种扩大。用图中六个教学节点$V=\{1,\ldots,6\}$，对不同节点在对称相似度不低于$0.8$时连无向边，再按连通分量分组，以每分量最小节点的顺序命名。设$\{1,2,3\}$与$\{4,5,6\}$各自内部两两分数为$0.9$，跨组为$0.2$，规则产生
\[
  \begin{aligned}
    E_{\mathrm{before}}
    &=\{(1,2),(1,3),(2,3),(4,5),(4,6),(5,6)\},\\
    \bm{c}_{\mathrm{before}}&=(1,1,1,2,2,2).
  \end{aligned}
\]
现在做一个“只改一条边”的对照实验，其余完全不动。仅把节点3与4之间的相似度从$0.2$\emph{错误地}改成$0.85$（越过了0.8阈值），于是新增一条边$(3,4)$；这条边恰好\emph{桥接}了原来的两个分量，使全部6个节点连成一体，分组目标立刻变成
\[
  \bm{c}_{\mathrm{after}}=(1,1,1,1,1,1).
\]
这个变化会\emph{连锁}放大到关系层面。若进一步把“同组内不同节点的无序对”都当作正对应（像5.2那样用于训练），那么原来两组各有3个节点、组内无序对共$3+3=6$对；合并成6个一组后，组内无序对总数变成$\binom{6}{2}=15$对——\emph{凭空新增了9对跨越原两组的“正对应”}。具体看：样本$(\bm{x}_4,c_4)$的组号从2被改成1；原本毫无关系的节点1与6也成了“正对应”，可它们之间\emph{直接}的相似度仍然只有$0.2$（根本不相近）。

这组数字非常直观地说明：\emph{连通规则会沿路径放大关系的范围}，一条错误边最终让一堆并不相近的样本被宣称为“同组、正对应”。它\emph{只演示了这种放大机制有多强}，\emph{完全没有}证明这六个对象在真实语义上属于同一类。

上面的结论依赖“对称阈值＋连通分量”这套\emph{具体}规则；换一套建图/划分规则，就必须\emph{重新明确}它的各个要素：节点代表什么、用什么距离度量、邻居数$k$或半径阈值多大、边有没有方向、边权怎么定、最终用什么算法划分。不能因为换了规则，还\emph{照搬}“全部连通”算例里的结论。

不同规则各有典型问题：\emph{固定$k$近邻}会在样本\emph{稀疏}的区域强行连出很远的邻居（那里本没有近的点，为凑够$k$只好连远样本）；\emph{固定半径}则会在\emph{密集}区域产生\emph{过多}的边（点太密，半径内全是邻居）。采用\emph{互近邻}（必须互相都把对方当邻居才连）、\emph{局部自适应尺度}或\emph{密度}类规则可以缓解这些问题，但无论怎么调，选择都\emph{依赖具体的样本分布}，不存在一套对所有数据都最优的统一配置。

图分组\emph{刷新}一次，要依次完成：重新编码所有样本、检索近邻、建立边、再做划分，成本和依赖都不少。先把计算与存储复杂度理清楚。对$n$个样本，如果直接做\emph{两两全比较}，精确检索需要$O(n^2)$次距离计算。存储则要看存什么：只有\emph{显式保存全部两两距离}才需要$O(n^2)$空间；如果最终只保留每个点的$k$个近邻，邻接表的规模是$O(nk)$（稀疏得多）；另外近邻检索所用的索引要单独算成本。为加速常用\emph{近似最近邻检索}，但它\emph{可能漏掉本应存在的边、或加入错误的边}——对图分组来说，漏边/错边都会直接改变群组，需谨慎。最后，重算后群组编号同样可能\emph{置换、分裂或合并}，依赖它的训练目标要\emph{同步刷新}；那些\emph{无法确定}的关系仍应如实保留为不确定，而不是强行连边或断开。

本节的群组还可以\emph{反过来}转成训练用的正对应：把同组的样本两两视为正对应，就能接入第~\ref{sec:data-derived-knowledge-relations}节的关系目标。但必须清楚这条依赖链的方向：这样产生的新关系\emph{继承的是“图分组假设”}（它们被连在同一分量里），而\emph{不是}经过额外确认的语义关系——分组可能本就被错误桥接边污染，转成正对应只会把这个假设继续传递下去。

尤其要与第~\ref{chap:knowledge-acquisition-model-based}章的图标签传播\emph{划清界限}，二者\emph{都用了图}，依据却根本不同：本节是\emph{先没有任何标签}、纯粹\emph{从样本之间的关联里发现群组}（图上没有类别锚点）；而第4章的图传播\emph{已经有}一批带任务标签的锚点，是\emph{借图把已有答案传播}给未知节点。正因为本节群组\emph{没有}标签锚点，它们天然不具备类别含义；只有拿到\emph{额外证据}（标签核对、领域解释等），才能把这里发现的群组解释成任务类别。

图分组在迭代训练中还会形成一个\emph{自我强化的反馈回路}，值得警惕：\emph{当前表示决定了谁是近邻}（建图依赖表示），而\emph{基于近邻的分组目标又会训练着把这些表示拉得更近}（训练依赖图）。这意味着，一旦早期因表示不成熟而\emph{建错了一条边}，这条错误连接会让两端表示被训练得越来越像，而越像就越会\emph{稳定地保持近邻关系}——一个\emph{稳定的错误}就此被持续锁定、不断强化，难以自行纠正。

为了能事后追查这类问题，应保留每条边的\emph{来源}、它与判定阈值之间的\emph{距离裕量}（贴着阈值的边最不可靠）、以及边在\emph{跨轮次}中的稳定性（时有时无的边可疑）。但这些记录主要用于诊断“边可不可疑”；连接\emph{到底有没有用}，仍要靠\emph{独立的下游用途检验}来回答。再次强调：\emph{图损失下降}这一件事，\emph{无法独自}完成“边是否正确”和“连接是否有用”这两项判断。

\section{本章小结}
\label{sec:data-derived-knowledge-summary}

回到最初那批没有物种标签的花卉照片，本章展示了一件关键的事：\emph{即使没有任何外部类别答案，数据自身也足以构造出三类不同的训练依据}，它们对应全章三条路线，结尾值得再并排回顾一次。

第一条是\emph{观测内容}：只要\emph{保留}被遮住的花瓣像素、或尚未到来的后续影像，模型之外就存在一个固定的内容答案（5.1）。第二条是\emph{样本关系}：保存“同对象两个视图应共享特征”的正对应、采集天然形成的自然配对、以及用来对照的负对照，就得到了“表示应当共享什么、允许区分什么”的关系约束（5.2）。第三条是\emph{样本分组}：用当前表示形成簇、把样本分配给原型、或在图上形成连通群组，就得到算法生成的临时目标（5.3）。

这三类依据都能驱动预测、训练模型，但它们各自\emph{依赖的前提}完全不同：观测目标依赖“被保留的观测”、关系目标依赖“成立的关系假设”、分组目标依赖“分组规则与表示”。因此\textbf{不能只凭损失函数的形式去推断目标的来源和可靠性}——同一个“余弦相似度损失”，背后既可能是可靠的同图裁剪、也可能是错误的负对照；同一个“分类式损失”，其目标既可能是真实类别、也可能是随时重排的临时簇号。

要让这些自造的依据\emph{可检查}，本章为每类构造都立了具体的核查要求。对\emph{观测目标}，核心是\emph{隔离答案、固定损失位置}：必须保证目标没有通过切分或预处理（平滑、全图归一化、未来信息）泄漏进输入，并明确只在哪些位置计分。对\emph{关系目标}，核心是\emph{说清比较对象、候选集合和梯度路径}：正对应是谁、负对照是谁、梯度回到一路还是两路。

三种关系目标的机制也在此对照：\emph{对比判别}在一批候选中奖励正对应的\emph{相对}优势（需要负对照）；\emph{跨视图预测}把另一路算出的\emph{表示向量或投影分布}作为当步答案（BYOL用停止梯度加EMA目标、DINO比较分布，都不需要显式负对照）；\emph{对称匹配}则\emph{同时调整两路}，并额外用批内方差、协方差检查\emph{跨样本}的变化以防坍塌（VICReg）。

这里有两个全书层面的概念不能混用。其一，\emph{停止梯度}只约束“\emph{当前这一步}不求导”，而目标\emph{以后}仍可能因共享参数重算、周期性复制、或EMA而改变——当步冻结与跨步更新是两回事。其二，第4章“\emph{任务教师}”的输出坐标\emph{对应预设类别}，而本章表示目标（BYOL向量、DINO分布）的坐标是\emph{未命名}的特征维度；不能因为都叫“目标/教师”就把它们当成同一种答案。

\emph{分组目标}相比前两类，多出一个特殊的维护要求：必须\emph{同步维护目标的“身份”}，因为算法生成的组号天然不稳定。聚类编号会\emph{跨轮置换}（同样的簇这一轮叫1、下一轮叫3），预测头不能凭数字相同就继承；原型分配\emph{受批次均衡和熵权重影响}，一个样本的目标会\emph{依赖同批其他样本}，且存在“均匀摊平”这种无区分的可行解；图\emph{连通}则会\emph{放大局部错误边}——一条错误桥接边可让大量本不相近的样本被当成同组正对应。在取得额外证据前，\textbf{这些组号都还不等同于物种}。

SwAV是一个很好的“机制交叉”范例，值得在结尾记住：它把\emph{原型分组}（均衡分配产生目标）和\emph{跨视图正对应}（同图两裁剪预测一致）\emph{联合}在一个方法里，说明目标的\emph{来源}和\emph{利用方式}本来就可以交叉组合。但越是这种联合目标，越要\emph{分别}说清每一部分的依据是什么、更新顺序如何（先在线更新再EMA、Sinkhorn停止求导）、以及各组件兼容的条件（均衡假设、队列多样性），不能用一个笼统的方法名掩盖内部的多重依赖。

最后必须把全章的立场讲到底，避免对自监督产生不切实际的期待：本章回答的只是“\emph{怎样在没有标签时，从已有样本中\emph{建立}可训练的依据}”，它\textbf{并没有保证这些依据会自动支持正确的物种判断}。

这些方法里的每一个旋钮，本质上都是在\emph{引入一个假设、并消耗一份计算}：遮蔽比例（假设缺失可由上下文推断）、数据变换（声明某种变化不应影响表示）、负对照采样（假设“未标正即不同”的近似成立）、目标更新方式、方差/协方差等统计要求、原型个数、建图的邻居尺度……没有一个是免费或中立的。而三类\emph{结构性边界}始终需要专门检查：优化可能停在\emph{精确常数驻点}（梯度为零推不动）、均衡分配可能\emph{均匀摊平}（不区分任何样本）、模型可能走上\emph{背景捷径}（靠设备、水印、背景而非花朵拿低分）。

因此评价必须\emph{分三层、分别确认}：\emph{构造条件}是否成立（关系对不对、有没有泄漏、时间边界守没守住）、\emph{表示差异}是否存在且非冗余（方差、协方差、有效秩）、以及这些差异在\emph{下游任务}上是否真的\emph{有效用}——训练损失一个数字替代不了这三层。本章至此解决了“单靠数据自身如何造依据”；当我们要把这些自监督得来的依据，与人工、规则、模型等\emph{多种来源组织在一起共同使用}时，进入第~\ref{chap:knowledge-fusion}章；而要把这些五花八门的证据\emph{接成一套可复核的质量协议}、判断它们到底可不可靠，则由第~\ref{chap:knowledge-quality-assessment-correction}章完成。

\chapter{知识的融合}
\label{chap:knowledge-fusion}

经过前几章，我们手里关于\emph{同一张花卉照片}的信息一下子多了起来，而且彼此\emph{并不一致}：人工判断它是茶梅；两条用同一套模板写出来的规则都报山茶；一个模型这次干脆没有作答。系统里还另外存着两个分别训练好的分类器，以及专门识别颜色、识别形态的小模块。这些来源各自都“懂”一点，但它们给出的东西互相重复、互相打架。

材料多了，能做的事情自然也多了。我们可以\emph{综合已有的样本标注}，把它们变成下一轮训练更可靠的依据；可以让现成的多个模型\emph{在预测时一起表态}；也可以干脆\emph{造一个新模型}，把这些老模型的能力承接过去。

可是“来源更多”这件事本身，并\textbf{没有}自动解决四个棘手的问题：\emph{重复证据}（两条同模板规则其实只提供了一份独立判断，却看起来像两票）、\emph{答案冲突}（山茶、茶梅、弃权同时出现，听谁的）、\emph{接口错位}（颜色模块和形态模块输出的维度、顺序对不对得上），以及\emph{能力损失}（把多个模型压成一个之后，会不会把它们各自的长处弄丢）。本章要回答的共同问题正是：\textbf{哪些样本标注或哪些已有模型可以合在一起用、应当在哪一个环节汇合、又怎样判断汇合之后是否真的满足了任务要求}？

先把“融合的对象”分清楚，这是全章最重要的一条线索。前几章取得的\emph{类别标签、样本配对关系、预测约束、恢复目标}，本质上都是\emph{数据/标注层面}的东西：它们可以直接进入训练样本、参与损失计算。而训练好的\emph{分类器和计算模块}，其行为是由\emph{模型参数}固定下来的，是另一种形态的资产。

正因为对象不同，“融合”这个词在本章实际上指代五件\emph{很不一样}的事，我们要把它们放在\emph{同一个任务要求和同一份预算}下逐一比较：综合多份样本标注、让多个模型共同预测、把计算模块连接起来、对模型参数做融合、以及用教师模型去蒸馏一个学生模型。

还要先纠正一个常见的误解：\textbf{融合并不意味着最后只能剩下一条答案、或一个模型}。完全可以保留多个来源模型、保留完整的类别概率、或保留一组候选标签。“怎么用”融合结果，是最后一节才专门讨论的独立选择。

本章的推进顺序如下。先用花卉系统\emph{比较融合的对象和路线}，让读者对全局有张地图；随后针对\emph{样本标签}，建立“来源到底靠不靠谱”的模型，并核对不同标注指向的\emph{对象、语义、形式}是否一致。

进入具体机制后，我们严格\emph{按照“样本标注或模型在哪一个位置参与计算”}来组织，从外到内、从浅到深依次是：样本标注与模型的\emph{预测输出}（最表层，谁都能拿到）、内部的\emph{计算模块}、再到\emph{模型参数}、最后是作为监督信号的\emph{蒸馏}。讲完机制，再说明融合结果\emph{怎样使用、怎样验证}。

这里要预先提醒两个\textbf{不能互相替代}的方向：其一，对来源作答所建模的可靠性、混淆和依赖，是\emph{服务于相应概率推断}的，\textbf{并非}所有参数融合、模块融合都得先估出一张来源混淆矩阵；其二，反过来，“我能把几个模型跑起来、或者把参数平均一下”，也\textbf{替代不了}标注对象的对齐和证据核验这两件更基础的事。除文献另行注明外，本章所有数值都是为讲解而\emph{教学构造}的，不是真实实验结果。

\section{知识融合的对象与路线}
\label{sec:knowledge-fusion-routes}
\label{sec:ket-integration-composition}

为了让“路线”这个概念落地，我们\emph{先固定住问题本身}：始终是\emph{同一张照片、同一个物种判断}。在这个问题上，人工和规则提供的是\emph{类别标签}，而来源模型除了标签，往往还能返回\emph{整类的概率分布}。

当我们把这些标注或预测输出\emph{汇总}起来时，直接得到的东西无非三种形态：一个\emph{标签}、一个\emph{类别分布}、或者\emph{未决}状态。而这份汇总结果有\emph{两种去处}要分清：它可以\emph{直接用于当前这一次预测}，也可以\emph{当作下一轮训练的监督目标}——后者就从“用模型”滑向了“造训练数据”。

特别要注意\emph{蒸馏}里的一个易混点：教师模型返回的是\emph{预测输出}，是\emph{学生的学习算法}拿这些输出去\emph{构造训练目标}。教师的“预测过程”和学生的“训练过程”是\textbf{两个独立的过程}，不能糊里糊涂当成同一次映射。

然而，光有“类别票”有时\emph{根本不够用}。设想任务不是“这是什么花”，而是“图中\emph{哪一朵}花\emph{同时}呈红色\emph{且}是钟形”。这个问题需要的不只是一个物种名，而是\emph{属性}（颜色、形态）以及\emph{属性属于哪一个对象}之间的关系。当现有类别票里没有这些关系时，我们就必须往模型\emph{内部}走，去寻找可以继续参与计算的\emph{中间表示}。

比如让颜色模块和形态模块\emph{共同}支持属性判断——但前提是必须把“红色”和“钟形”\textbf{绑定到同一朵花}上，不能取这朵的颜色、配那朵的形态。这类做法的特点是\emph{保留各个来源单元}，在预测时按事先约定把它们\emph{连接}或\emph{选择}起来。

\term{模型组合}（model composition）这个术语既涵盖这种\emph{模块连接}，也涵盖\emph{保留多个完整模型、再组合它们预测}的做法；它描述的是\textbf{部署（运行）时如何让来源模型或计算模块共同工作}。

但如果\emph{部署预算只允许保存并运行一份模型}（比如花卉识别终端算力、内存都很紧），前面那种“运行时同时带着好几个模型”的组合就行不通了，这时有两条路线。

一条是\term{模型参数融合}（model merging）：在\emph{准备阶段}直接对已经\emph{对应好}的来源参数（或它们相对公共底座的增量）做运算，\emph{拼出一份新的模型参数}。另一条是\term{知识蒸馏}（knowledge distillation）：拿教师的\emph{类别票或中间表示}当\emph{监督}，去\emph{训练}一个承接能力的新模型（学生）。

两者的前提很不一样，务必分清：\emph{参数融合}要求来源模型的参数\textbf{能够对应}（同一个位置承担同一种计算）；\emph{蒸馏}则要求你能\textbf{查询}到预测输入、拿到\textbf{教师信号}、并\textbf{真正训练}学生。蒸馏本质上也是“把已有成果转成新的学习依据”，而且要记住：\textbf{最终只部署学生，并不意味着准备阶段可以省掉访问教师和训练学生所花的成本}。

图~\ref{fig:ket-integration-routes}把上面三条“利用现成模型”的路线并排画在一起，便于对照：\emph{组合}是\textbf{保留来源、运行时一起跑}；\emph{参数融合}是\textbf{读取来源参数、离线构造一份新模型}；\emph{蒸馏}是\textbf{用教师当监督、训练出一个学生}。注意图中双重轮廓标出的，是\emph{部署时仍需保存并运行的单元}——组合要带着所有来源和聚合器，融合只跑新模型，蒸馏只跑学生。

表~\ref{tab:ket-integration-route-choice}进一步从\emph{来源访问、是否需要额外训练、最终部署形态、首要失败边界}四个角度比较它们。这三条路线\textbf{不是互斥的}，完全可以\emph{串联}：例如先把几个教师的预测\emph{组合}起来，再用组合后的结果去\emph{蒸馏}一个学生（带路由的组合系统本身也能当教师）。这张图和表只比较“来源模型/计算模块怎么被利用”；至于\emph{样本标注层面}的投票和概率推断，是另一套机制，放到后文专门推导。

\begin{table*}[htbp]
  \centering
  \small
  \begin{tabularx}{\linewidth}{@{}p{22mm}XXXX@{}}
    \toprule
    路线 & 来源访问 & 额外材料与训练 & 部署形态 & 首要失败边界 \\
    \midrule
    模型组合 & 部署时能够执行来源；模块连接还需中间接口 & 固定规则可直接使用；学习聚合或路由需要相应材料 & 保留所需来源、聚合器或路由器 & 接口错位、共同错误、错误路由与重复计算 \\
    模型参数融合 & 准备时能够读取并对应来源模型参数及状态 & 可只用模型参数，也可用预测输入统计、激活或标签决定运算 & 通常保存一份模型参数，也可能保留任务分支 & 坐标不对应、平均经过高损失区域与任务冲突 \\
    知识蒸馏 & 训练时能够取得教师响应；表示约束还需内部激活 & 需要查询预测输入、学生训练和独立模型选择 & 通常只保留学生的预测路径 & 预测输入覆盖不足、教师偏差与学生容量不足 \\
    \bottomrule
  \end{tabularx}
  \caption{三条模型融合路线的访问、训练和部署条件。各行比较的是来源如何参与结果，条件可以交叉，例如带路由的组合系统也能充当蒸馏教师。部署只保留一个模型，不表示准备阶段无需样本、教师存储或额外计算。}
  \label{tab:ket-integration-route-choice}
\end{table*}

\begin{figure*}[htbp]
  \centering
  % Three independent integration routes; schematic, no empirical data.
% Final canvas: 164 mm x 131 mm. Double outlines mark deployed units.
% Solid: forward computation; dashed: training; dotted: parameter reading.
\begin{tikzpicture}[
  x=1mm,y=1mm,font=\figurefont\footnotesize,text=BookInk,
  box/.style={rectangle,draw=FigureModel,fill=FigureModel!7,
    line width=0.8pt,minimum width=25mm,minimum height=9mm,
    inner sep=1mm,align=center},
  deployed/.style={box,double,double distance=0.7pt},
  input/.style={box,draw=FigureInput,fill=FigureInput!7,
    minimum width=15mm},
  result/.style={box,draw=FigureOutput,fill=FigureOutput!7,
    minimum width=15mm},
  loss/.style={box,draw=FigureLoss,fill=FigureLoss!7},
  flow/.style={-{Latex[length=2mm,width=1.3mm]},draw=BookInk,
    line width=0.8pt,rounded corners=1mm},
  train/.style={flow,dashed,draw=FigureLoss},
  weights/.style={flow,dotted,draw=BookBlue},
  construct/.style={flow,draw=BookBlue},
  note/.style={font=\figurefont\footnotesize,text=BookMuted,inner sep=1mm}
]
  \path[use as bounding box] (0,-13) rectangle (164,118);
  \node[anchor=west,font=\figurefont\small] at (1,113) {(a) 保留来源计算};
  \node[input] (xa) at (12,94) {$\bm{x}$};
  \node[deployed] (a1) at (47,101) {来源$f_1$};
  \node[deployed] (a2) at (47,87) {来源$f_2$};
  \node[deployed] (agg) at (99,94) {聚合器};
  \node[result] (ya) at (150,94) {预测};
  \draw[flow] (xa.east) -- (25,94) |- (a1.west);
  \draw[flow] (25,94) |- (a2.west);
  \draw[flow] (a1.east) -- (76,101) |- (agg.west);
  \draw[flow] (a2.east) -- (76,87) |- (agg.west);
  \draw[flow] (agg.east) -- (ya.west);
  \draw[BookRule,line width=0.6pt] (0,78)--(164,78);

  \node[anchor=west,font=\figurefont\small] at (1,73) {(b) 读取来源模型参数};
  \node[box] (b1) at (34,60) {来源$f_1$\\$\bm{\theta}_1$};
  \node[box] (b2) at (34,45) {来源$f_2$\\$\bm{\theta}_2$};
  \node[box] (rule) at (76,53) {融合规则};
  \node[deployed] (merged) at (122,53) {融合模型};
  \node[result] (yb) at (151,53) {预测};
  \node[input] (xb) at (92,37) {$\bm{x}$};
  \draw[weights] (b1.east) -- (53,60) |- (rule.west);
  \draw[weights] (b2.east) -- (53,45) |- (rule.west);
  \draw[construct] (rule.east) -- node[note,above] {构造新模型参数} (merged.west);
  \draw[flow] (xb.east) -| (merged.south);
  \draw[flow] (merged.east) -- (yb.west);
  \draw[BookRule,line width=0.6pt] (0,29)--(164,29);

  \node[anchor=west,font=\figurefont\small] at (1,24) {(c) 用来源监督训练};
  \node[input] (xc) at (12,8) {$\bm{x}$};
  \node[box] (c1) at (45,15) {来源$f_1$};
  \node[box] (c2) at (45,1) {来源$f_2$};
  \node[loss] (kd) at (85,8) {蒸馏约束};
  \node[deployed] (student) at (124,8) {学生};
  \node[result] (yc) at (152,8) {预测};
  \draw[flow] (xc.east) -- (24,8) |- (c1.west);
  \draw[flow] (24,8) |- (c2.west);
  \draw[train] (c1.east) -- (65,15) |- (kd.west);
  \draw[train] (c2.east) -- (65,1) |- (kd.west);
  \draw[train] (kd.east) -- node[note,above] {更新} (student.west);
  \draw[flow] (student.east) -- (yc.west);
  \draw[flow] (student.north) -- (124,21) -| (kd.north);
  \node[note] at (105,23) {学生预测输出};
  \draw[flow] (xc.south) -- (12,-9) -| (student.south);
  \node[note,fill=white] at (86,-9) {同一样本};
\end{tikzpicture}

  \caption{同一对来源模型参与融合的三种方式。实线表示前向计算，虚线表示监督约束及学生更新，点线表示模型参数读取，蓝色实线表示融合规则生成新模型参数。双重轮廓标出部署时仍需保存并运行的单元：组合需要来源与聚合器，融合运行新模型，蒸馏运行学生。三个面板互相独立，不表示操作先后。}
  \label{fig:ket-integration-routes}
\end{figure*}

无论选哪条路线，有一套\emph{共同的台账}必须记录：任务是什么、用的是哪个\emph{来源版本}、能访问哪些\emph{样本}、拿到的是\emph{参数还是表示}、以及各自的\emph{成本}。

也要清楚本章\emph{不}负责回答的问题，避免越界：某个来源\emph{值不值得复用}，由第~\ref{chap:ket-reuse-selection}章检验；模型被修改后\emph{怎样保住原有能力}，由第~\ref{chap:ket-update-maintenance}章处理；多个当前任务\emph{共同}塑造表示和参数的\emph{共享}问题，属于第~\ref{chap:ket-sharing}章。本章只聚焦一件事：\textbf{在已经允许使用的这些来源之间，它们怎样共同发挥作用}。

最后破除一个简单的分类错觉：\textbf{不能用“这里面有没有训练”来判断它属于哪一类问题}。参数融合可以学习系数、组合可以训练一个聚合器/路由器、蒸馏当然也要训练学生——“是否训练”区分不了路线，真正的区别在于\emph{融合发生在预测的哪一个位置、动的是什么对象}。

\section{知识来源的建模}
\label{sec:knowledge-fusion-source-modeling}

面对人工、规则和模型给出的一大堆标签，\emph{在急着把它们加起来之前}，我们必须先做一件更基础的事：\textbf{把每一份判断追溯回“产生它的那个过程”}。同样一个“山茶”，出自资深园艺师、出自一条只看叶缘的规则、出自一个模型，其分量完全不同。

\term{监督来源}（supervision source）就是对“产生样本标注的主体或过程”的统称：它可以是\emph{一名标注者、一条规则、一个模型，或模型的一种固定配置}。这里藏着一个贯穿全章的陷阱——\textbf{列名不同，不等于来源独立}：两列类别票如果\emph{共用同一套模板、同一本词典、同一批训练样本或同一个基础模型}，那它们本质上只是一个判断的两份拷贝。

所以在描述来源时，模型要\emph{同时}读取四样东西：来源的\emph{身份}、它的\emph{覆盖状态}（这题它到底答没答）、\emph{已对齐的结果}（下一节处理），以及来源之间的\emph{共享关系}。用这些样本标注去刻画每个来源的\emph{能力}和彼此的\emph{依赖}，才能为后面的证据组合提供可靠的模型参数。

那么，应当从哪几个方面去“解释”一份来源标注？本章固定用\textbf{三个互补的角度}，缺了任何一个都会出问题：

\begin{itemize}\item \term{可靠性}（reliability）：在本章的标签任务里，它回答“这个来源\emph{一旦作答，有多大概率答对}”，是一个总体比例。\item \term{混淆}（confusion）：光知道“八成对”还不够，还要知道“\emph{真实类别给定}时，它\emph{容易报成哪一类}”，即错误的\emph{方向}。\item \term{相关性}（correlation）：考察\emph{多个来源的票或错误是否一起变化}，尤其是它们因为\emph{共享证据}而产生的依赖——这直接决定了几票能不能当成几份独立证据。\end{itemize}

这三个角度之所以要分开，是因为\textbf{它们彼此不能替代}：一个\emph{正确率}说明不了错误\emph{方向}；把每个来源\emph{分别}估得再准，也说明不了它们\emph{合在一起}的行为。两个总体正确率相同的来源，可能一个老把茶梅报成山茶、另一个老把茶梅报成月季；而四个正确率都很高的来源，也可能只是\emph{同一个错误机制的四份副本}——表面四票，实际一份证据。下面三小节就依次把这三个角度建模清楚。

\subsection{来源可靠性的建模}
\label{sec:knowledge-fusion-reliability}

先解决一个容易被掩盖的问题：\textbf{“答没答”和“答对没”是两件事，必须分开记}。如果一个来源\emph{专挑容易的照片}回答，那么把“未作答”和“答错”揉进同一个指标，你会高估它的能力——它显得很准，只是因为难题它根本不碰。

所以我们引入两个量。设在\emph{已完成对象对齐与语义对齐}的$n$个样本上有$m$个来源，共同类别集合为$Y=\{1,\ldots,K\}$。来源$s$在样本$i$上给出的有效类别记为$\lambda_{i,s}\in Y$；再用一个\emph{覆盖掩码}$o_{i,s}$：给出有效类别时取$1$，否则取$0$。

要特别强调这个掩码的\textbf{语义边界}：它\emph{只是为了计算而挑出“可比较的标签”}，并\textbf{不会}改写原始样本标注——原始记录里的\emph{执行状态、来源动作、取得条件}都原样保留。最典型的区分是：一条规则\emph{运行成功但没有触发}，此时动作是空、内容未知，触发与否的结果要另外记录，\textbf{绝不能}把它改记成一次“主动弃权”。

\subsubsection{利用参考标签估计可靠性}
\label{sec:knowledge-fusion-reference-reliability}

最理想的情形，是\emph{有一部分照片已经经过独立审计、拿到了参考标签}。这时我们可以\textbf{直接核对}来源的答案对不对，而\textbf{不必}先靠“来源之间互相赞同”去\emph{猜}真值（那是没有参考时的无奈之举，见后文）。记被审计样本的索引集合为$I_{\mathrm{audit}}$、其参考标签为$y_i$。

在核对时需要\textbf{两个分母不同的量}，初学者最容易把它们混为一谈：\emph{覆盖率}回答“在\textbf{全部}目标样本里，这个来源\textbf{回答了多少}”；\emph{覆盖条件下的正确率}回答“在\textbf{既有参考、又确实作答}的样本里，它\textbf{答对了多少}”。两者各自在\emph{不同的样本集合}上取比例：
\begin{equation}
  \widehat c_s
  =\frac{1}{n}\sum_{i=1}^{n}o_{i,s},
  \qquad
  \widehat a_s
  =\frac{\sum_{i\in I_{\mathrm{audit}}}
    o_{i,s}\mathbb{I}\{\lambda_{i,s}=y_i\}}
  {\sum_{i\in I_{\mathrm{audit}}}o_{i,s}}.
  \label{eq:knowledge-fusion-audited-reliability}
\end{equation}
请盯住这两个\textbf{分母的区别}：覆盖率的分母是\emph{全部}$n$个样本——哪怕没答案，也照样占一个位置，所以弃权、失败都会拉低覆盖；而正确率的分母\emph{只数}审计样本集合里\emph{真正作答}（$o_{i,s}=1$）的样本，分子再数其中与参考一致的。

这里有两个必须守住的规矩。第一，当第二个分母为$0$（这个来源在审计样本上一题都没答）时，正确率是\textbf{尚未估计}，\textbf{不能}随手填成$0$或$1$。第二，\emph{未请求、执行失败、主动弃权}都会\emph{降低覆盖率}，但它们并\textbf{不直接构成一次类别错误}；如果你图省事把“弃权”当成一个负类（比如当成“不是山茶”），那你统计的\textbf{已经不是原来那个分类任务}了。

表~\ref{tab:knowledge-fusion-source-records}把这两个分母放在同一组教学标注中。规则$R_1,R_2$共享模板，人工来源$H$独立作答，模型$M$可以弃权；前四张照片有独立参考，后四张没有参考。我们先用前四行核对正确性，再用全部八行计算覆盖，后文仍用这些类别票比较融合规则。表中标注均为人工构造，不是实验结果。

\begin{table*}[htbp]
  \centering
  \small
  \renewcommand{\arraystretch}{1.15}
  \begin{tabular}{@{}clccccc@{}}
    \toprule
    \textbf{样本} & \textbf{区域} & \textbf{参考} & \textbf{$R_1$} & \textbf{$R_2$} & \textbf{$H$} & \textbf{$M$} \\
    \midrule
    $x_1$ & 简单 & 山茶 & 山茶 & 山茶 & 山茶 & 山茶 \\
    $x_2$ & 困难 & 茶梅 & 山茶 & 山茶 & $\varnothing_{\text{未请求}}$ & 茶梅 \\
    $x_3$ & 简单 & 山茶 & 山茶 & 山茶 & 山茶 & 茶梅 \\
    $x_4$ & 困难 & 茶梅 & $\varnothing_{\text{失败}}$ & $\varnothing_{\text{失败}}$ & $\varnothing_{\text{未请求}}$ & 山茶 \\
    \midrule
    $x_5$ & 简单 & 未审计 & 山茶 & 山茶 & 茶梅 & $\bot$ \\
    $x_6$ & 简单 & 未审计 & 山茶 & 山茶 & 茶梅 & 茶梅 \\
    $x_7$ & 未覆盖 & 未审计 & $\varnothing_{\text{失败}}$ & $\varnothing_{\text{失败}}$ & $\varnothing_{\text{未请求}}$ & $\bot$ \\
    $x_8$ & 简单 & 未审计 & 茶梅 & 茶梅 & 山茶 & 山茶 \\
    \bottomrule
  \end{tabular}
  \caption{多来源标签与作答覆盖清单。$\varnothing_{\text{未请求}}$表示没有发起请求，$\varnothing_{\text{失败}}$表示执行失败，不能用来代指仍在等待或处理的任务；$\bot$表示成功执行后主动弃权，内容仍为未知。这三种情况都不计类别票。$x_1$至$x_4$有独立参考，用于教学与开发诊断，不充当最终留出评价；所有对象和类别语义均假设已经对齐。}
  \label{tab:knowledge-fusion-source-records}
\end{table*}
\FloatBarrier

我们用表~\ref{tab:knowledge-fusion-source-records}把这套算法\emph{逐行走一遍}，体会“分母不同”会带来什么结论。就\emph{覆盖}而言，$R_1,R_2,M$的有效覆盖率都是$6/8$，只在简单区作答的$H$是$5/8$。只看\emph{审计区}（前四行）的正确性：两条规则各答对$2/3$，$H$答对$2/2$，$M$答对$2/4$。

但请注意：\textbf{这些比例的分母根本不是同一批对象}。$H$压根没回答困难区，所以“$H$答对$2/2$”和“$M$答对$2/4$”不能直接比。要做\emph{公平}比较，必须把范围\emph{限制到两者共同覆盖}的样本上：在简单审计样本$x_1,x_3$上，$H$是$2/2$、$M$是$1/2$；而困难区$x_2,x_4$只有$M$答了（$1/2$），对$H$则\textbf{无法估计}。

所以这张表\textbf{真正支持}的结论，仅仅是“\emph{在这两个简单样本上}$H$答得更准”，\textbf{绝不能}夸大成“$H$在全范围都更可靠”。一般地，想报告某个区域$g$的正确率，就要把式~\eqref{eq:knowledge-fusion-audited-reliability}的分子、分母\emph{同时}限制到$I_{\mathrm{audit}}\cap I_g$，并\emph{如实报告有效分母}；分组越细，所需的审计投入也越多。

最后要给这种核对本身\emph{降降温}：它的一切结论都\textbf{以参考标签可靠、抽样范围对题为条件}。参考答案自己也可能错；审计样本是\emph{怎么抽}的、按什么\emph{规范}、哪个\emph{版本}形成的，都必须\emph{对应你最终要使用的范围}，否则估出来的能力会“水土不服”。像$H$这样只答对两次，\textbf{绝不能}当成一个精确的能力值。

实操上：参考样本很少时，应保留\emph{置信区间}或参数的\emph{后验}，而不是报一个孤零零的比值；\emph{没有}参考的区域，老老实实记为\emph{未估计}。对那种\emph{选择性作答}的来源，还必须额外检查$o_{i,s}$如何随\emph{类别、难度、时间}变化——只有这样，才能分清它的正确率升高，到底是因为\emph{判断力真的变强了}，还是仅仅因为它\emph{越来越只挑容易题作答}。

\subsubsection{利用来源关系估计可靠性}
\label{sec:knowledge-fusion-relational-reliability}

现在\emph{抽掉参考标签}这根拐杖，只留下三列\emph{始终作答}的二分类标注，看看会撞上什么本质困难。假设一半样本三列都报山茶、另一半都报茶梅。我们能\textbf{确认}的只有一件事：\emph{三个来源高度一致}。但\textbf{一致不等于正确}，我们\textbf{无法直接核对}它们到底对不对。

原因在于，同样这批观测至少存在\textbf{两种都解释得通}的说法：解释甲认为三列\emph{全都报对}，于是山茶、茶梅各占一半；解释乙认为三列\emph{全都报反}（系统性地把两类颠倒），结果\emph{也是}两类各占一半。\textbf{这两种解释产生的观测分布完全相同}，但它们对\emph{每一个样本真值}的判断却\textbf{恰好相反}。光看票，你根本无法判断自己处在哪一个世界。

这个例子还揭示一个反直觉的事实：\textbf{单纯增加样本数量，并不能解套}。哪怕再来十万个样本、三列依旧齐刷刷一致，甲乙两种解释\emph{照样}都成立——因为障碍\textbf{不在观测数量，而在观测里压根没有“类别方向”这个信息}。这正是为什么要引入\term{可识别性}（identifiability）这个概念：它专门检查\textbf{不同的参数解释，能不能由观测分布区分开}。如果两组参数\emph{无论看多少数据都给出同一个分布}，那就说真实参数\emph{不可识别}，样本再多也无法唯一确定它。

怎么破？必须从数据之外\emph{额外}加进一点“方向信息”：\emph{少量独立参考}可以给潜在类别\emph{命名}（告诉模型“这一类才叫山茶”）；或者干脆\emph{假设}“来源总体上倾向于报对”，用这个先验把“全报反”的解释排除掉。但要清醒：此时排除报反解释，靠的是\textbf{新增的约束/假设}，\textbf{不是}因为答案已经被数据验证了。

把这件事用更形式化的语言说清楚：当我们想从标注样本估计来源能力时，\emph{看得见}的只有标签矩阵$\bm{\Lambda}$和覆盖掩码$\bm{O}$；\emph{看不见}的未知量有三类——来源的模型参数、类别的真实比例、以及每个样本的潜在答案。是\emph{共同样本上的赞同与分歧}把这些未知量\emph{纠缠}在一起，让我们有可能反推；但\textbf{能不能真的反推出参数}，还取决于\emph{重叠样本够不够、票与票之间有没有差异}。

这里有三种\textbf{互不相同、且都不是“多采点样本”能解决}的障碍，要能分清：
\begin{itemize}\item \emph{缺少重叠（桥梁断裂）}：如果两个来源\textbf{从不在同一批样本}上作答，又没有第三个来源同时和它们重叠，那两者之间\textbf{根本没有可比较的桥梁}，谁也无法借对方校准自己。\item \emph{复制票（证据冗余）}：十列如果\textbf{逐项照抄}同一条规则，那它们只提供了\textbf{一份}证据，不是十份。\item \emph{类别置换（方向缺失）}：就是前面“全对/全报反”的歧义，票的模式对、但类别的名字可以整体对调。\end{itemize}

如果非要在没有参考的情况下继续估计，常用手段有三类：\emph{减少待估参数}（让模型更简单）、\emph{合并重复来源}（把抄同一规则的几列并成一列）、或\emph{设置明确先验}（比如“大家总体偏靠谱”）。但每一种处理\textbf{都在悄悄改变“允许哪些解释”}，它们是\emph{建模选择}，\textbf{替代不了对真值的核验}。

第~\ref{sec:knowledge-fusion-latent-label-inference}节会用一个“来源如何生成标签”的生成模型（Dawid--Skene），把“算后验”和“更新参数”正式接起来，并分析在什么\emph{充分条件}下能排除额外歧义。这里先记住一个关键区分：\textbf{迭代算法能不能稳定停下来}是\emph{数值/优化}问题；而\textbf{观测到底能不能区分不同参数}是\emph{可识别性}问题——EM收敛得再漂亮，也不代表参数被唯一确定、更不代表类别方向正确。

\subsection{来源混淆与偏差的建模}
\label{sec:knowledge-fusion-confusion}

现在进入第二个角度。在山茶、茶梅、月季的三分类任务里，设想两个来源\emph{总体都答对八成}。单看正确率，它们一样好；但面对\emph{同一次分歧}，它们能提供的证据可能\textbf{截然不同}：A一犯错就老把\emph{茶梅报成山茶}，B一犯错却老把\emph{茶梅报成月季}。\textbf{总体正确率把这种“错误方向”平均没了}。

可融合时恰恰要用方向：当看到票时，我们真正想算的是“\emph{假如真值是$k$，出现当前这张票$\ell$的可能性有多大}”。这就需要为每个来源$s$配一张\emph{混淆矩阵}。以$Y_i$表示样本$i$的真实类别，矩阵\emph{逐行}记录“在\emph{给定真值$Y_i=k$、且来源确实作答}的条件下，它报成各个类别”的条件概率：
\begin{equation}
  \theta^{(s)}_{k,\ell}
  =\mathbb{P}(\lambda_{i,s}=\ell
    \mid Y_i=k,o_{i,s}=1),
  \qquad
  \sum_{\ell=1}^{K}\theta^{(s)}_{k,\ell}=1.
  \label{eq:knowledge-fusion-confusion-matrix}
\end{equation}
务必先把\textbf{行列方向}记牢，这是后面所有计算不出错的前提：\emph{行}索引$k$固定\textbf{真实类别}（条件），\emph{列}索引$\ell$遍历\textbf{实际报出的票}（结果），所以\emph{每一行}是一个条件分布、概率之和为一。读这张矩阵时：\emph{对角项}是“该类被报对”的概率（即分类别正确率），\emph{非对角项}则清楚显示了“错误\emph{流向}了哪一类”。

后文一旦\emph{看到某张票}$\ell$，就要反过来\emph{去各候选真值所在的行、取同一列}$\ell$来比较似然。因此无论写代码还是画图表，都\textbf{必须沿用“行=真值、列=票”这个方向}，行列一颠倒，整个推断就会错。

有参考标签时，估计这张矩阵就是\emph{数计数}：在审计样本集合内\emph{只数来源实际作答}的样本，把“参考为$k$、却报成$\ell$”的次数记为$N^{(s)}_{k,\ell}$。在“各项非负、每行和为一”的约束下最大化这些标注的似然，得到的就是简单的频率比
$\widehat\theta^{(s)}_{k,\ell}=N^{(s)}_{k,\ell}/\sum_rN^{(s)}_{k,r}$。
分母是“真实为$k$”这一行下的\emph{有效作答总数}；分母为零，说明这一类从没在审计样本里出现过，该行\emph{未估计}，不能硬算。

小样本下还会遇到\emph{零计数}问题：某个稀有类别的样本太少，很多格子是$0$，被当成“绝不可能”往往过强。这时可以用\emph{Dirichlet平滑}（给整行先加入一个明确的先验伪计数再估计），或\emph{合并过细类别}、\emph{报告区间}。但要记住：平滑只是\textbf{缓解零计数带来的过强断言}，它\textbf{并不会凭空增加任何真实观测}。

回到A、B两个来源，设它们在同一目标范围上总体正确率\emph{都是}$0.8$，但A主要把茶梅报成山茶、B主要把茶梅报成月季。现在观察到一次分歧：\textbf{A报山茶、B报月季}。这时真值\emph{仍然可能是茶梅}——因为“茶梅”恰好能同时\emph{解释}这两张票：要验证它，就该去\textbf{读A矩阵的“茶梅行、山茶列”}这一项，再读\textbf{B矩阵的“茶梅行、月季列”}这一项，然后在相应的\emph{依赖假设}下把它们组合起来。

请注意：手里只有“$0.8$、$0.8$”两个数字，\textbf{根本无法完成这个计算}，也无法单独定出最终类别——这正是必须保留\emph{分类别似然}（混淆矩阵）而非仅留正确率的根本理由。（这里的比例与错误方向都是教学设定，只为说明“为什么需要分类别似然”。）

混淆矩阵还能帮我们避免一个常见的误判。假设你发现“票里报\emph{山茶}的比例\emph{升高}了”，能不能就此断定这个来源退化了、变得偏爱山茶？\textbf{不能}，因为至少有\textbf{三个完全不同的机制}都能推高这个比例：
\begin{itemize}\item 来源\emph{更多地挑山茶样本来作答}——这改变的是\textbf{覆盖机制}；\item 在\emph{真值相同}的情况下它\emph{更偏向}报山茶——这才是\textbf{混淆矩阵}变了；\item 根本是\emph{目标样本里山茶本身变多}了——这改变的是\textbf{类别先验}$\mathbb{P}(Y=k)$。\end{itemize}
三种变化都会抬高“山茶票”的\emph{边缘比例}，但含义天差地别。所以\textbf{只看票的边缘分布，不足以认定来源退化}；必须\emph{结合覆盖情况和分类别的参考}才能定位到底是哪一环变了。

一张“全局”混淆矩阵，隐含假设是：\emph{在所建模的整个范围内，来源的行为稳定不变}。但如果它在\emph{困难照片}和\emph{简单照片}上犯错的模式不同，一张全局矩阵就把这种差异平均糊掉了。这时应改用\emph{分组}参数$\theta^{(s,g)}_{k,\ell}$（按区域$g$分别建矩阵），或让概率\emph{依赖样本特征}。注意行列的含义\textbf{不变}，只是\emph{每条估计所依据的范围变小了}。

分组是一把\textbf{双刃剑}：组划得越细，参数越多，\emph{每一行能分到的参考样本和重叠票就越少}，估计越不稳。所以应\emph{按实际使用需要}来划组，并核对每组的\emph{有效样本量}和\emph{更新频率}，而不是越细越好。另外，如果想让\emph{弃权}本身也参与推断，可以把它当成一个\emph{额外的票状态}，或者对覆盖\emph{单独建模}——但无论哪种，都\textbf{不会把弃权当成一个真实的物种类别}。

\subsection{来源相关性的建模}
\label{sec:knowledge-fusion-correlation}

第三个角度是\emph{相关性}，它处理的正是章首那个“两条同模板规则”的隐患。表~\ref{tab:knowledge-fusion-source-records}里$R_1,R_2$经常\emph{同时}报山茶，看起来互相印证；但它们\textbf{共用同一套模板}，所以第二张票到底有没有带来\emph{新的独立判断}，必须\textbf{另查}，不能想当然。

这种“看似多源、实则同源”的情况在实践中极其普遍：\emph{共享训练样本集、共享基础模型、共享检索结果、共享提示模板、甚至共享同一本标注指南}，都会让多个来源\emph{系统性地一起对、一起错}。正确的做法是：\textbf{先把这些已知的共享关系登记进“来源目录”}，再用已对齐的结果去\emph{检验是否还藏着没被描述的依赖}。核心纪律一句话：\textbf{绝不从“有几列票”直接推断“有几份独立证据”}。

为什么“一致率高”不能直接当好事？因为\textbf{高一致率背后可能是两种截然相反的机制}：其一是两个来源都\emph{忠实地随真实类别而变}（真值是山茶就都报山茶），这是我们想要的“真独立、真印证”；其二是它们在\textbf{真值已经固定之后}，仍因为\emph{共享的那个错误源}而\textbf{一起犯错}（都被同一套坏模板带偏）。这两种情形表面一致率都很高，价值却完全不同。

要把二者分开，正确的问法是：\textbf{把真值固定住以后}，两个来源的票还会不会互相影响？只有在“给定真值”的条件下，才能检验“把每个来源\emph{分别}描述”是否已经足够。对两个始终作答的来源$s,t$，如果它们\emph{联合}出票的概率，恰好等于\emph{各自条件概率的乘积}，就称它们在给定$Y_i$后\term{条件独立}：
\begin{equation}
  \mathbb{P}(\lambda_{i,s},\lambda_{i,t}\mid Y_i)
  =
  \mathbb{P}(\lambda_{i,s}\mid Y_i)
  \mathbb{P}(\lambda_{i,t}\mid Y_i).
  \label{eq:knowledge-fusion-conditional-independence}
\end{equation}
这条等式虽然简短，却是\textbf{后面“把多个来源的似然\emph{相乘}”这一操作的唯一合法依据}：正因为给定真值后两者独立，联合概率才能拆成乘积。务必注意它的条件——它\textbf{只要求“给定真值之后”独立}，并\textbf{不}要求两个来源在\emph{还不知道真值}时（边缘上）独立；事实上它们都会随真值变化，边缘上通常是相关的，这完全正常。

反过来，一旦这条等式\textbf{不成立}（来源确实共享证据），那么\textbf{光有各自的混淆矩阵就不够了}：你无法用两张“各自的”表还原出它们\emph{联合}犯错的概率。实践中有时没有参考真值，只能借助潜在标签去检查一种“\emph{残余相关性}”——即\emph{扣掉当前模型所能解释的类别影响之后}，还剩下的那部分共同变化；但这种检查\textbf{仍然受模型假设约束}，单凭一个经验相关系数，既\textbf{不能唯一确定}它们到底共享了什么机制，也\textbf{指不出因果方向}。

图~\ref{fig:knowledge-fusion-source-dependence}把这一区别画成两种假设。左侧每份类别票只按各自条件分布响应真实类别，右侧三条规则还受同一模板或证据影响。我们应观察右侧的虚线和来源组：它们提示三张赞同票可能应作为一组证据处理，而不是自动计作三份独立判断。

\begin{figure*}[htbp]
  \centering
  \begin{tikzpicture}[
  x=1mm,y=1mm,
  every node/.style={font=\figurefont\small,text=BookInk},
  truth/.style={circle,draw=BookTeal,fill=FigureModel!12,line width=1pt,minimum size=11mm,inner sep=1pt},
  source/.style={book node,draw=BookBlue,fill=FigureInput!9,line width=0.8pt,minimum width=18mm,minimum height=9mm,inner sep=2pt},
  shared/.style={book node,draw=BookAmber,fill=BookAmber!10,line width=0.8pt,minimum width=24mm,minimum height=9mm,inner sep=2pt}
]
  \node[font=\figurefont\bfseries] at (35,26) {给定类别后独立};
  \node[truth] (yl) at (35,10) {$Y$};
  \node[source] (l1) at (8,-10) {来源1};
  \node[source] (l2) at (35,-10) {来源2};
  \node[source] (l3) at (62,-10) {来源3};
  \draw[book arrow,knowledge arrow,draw=BookTeal,line width=1pt] (yl) -- (l1);
  \draw[book arrow,knowledge arrow,draw=BookTeal,line width=1pt] (yl) -- (l2);
  \draw[book arrow,knowledge arrow,draw=BookTeal,line width=1pt] (yl) -- (l3);
  \node[font=\figurefont\scriptsize,text=BookMuted,align=center] at (35,-26) {三条标注路径分别提供类别票};

  \draw[BookRule,line width=0.7pt] (73,-40) -- (73,29);

  \node[font=\figurefont\bfseries] at (124,26) {共享机制引入相关错误};
  \node[truth] (yr) at (124,10) {$Y$};
  \node[source] (r1) at (86,-10) {规则1};
  \node[source] (r2) at (109,-10) {规则2};
  \node[source] (r3) at (132,-10) {规则3};
  \node[source] (human) at (155,-10) {标注者$H$};
  \node[shared] (u) at (109,-38) {共享模板或证据};
  \draw[book arrow,knowledge arrow,draw=BookTeal,line width=1pt] (yr) -- (r1);
  \draw[book arrow,knowledge arrow,draw=BookTeal,line width=1pt] (yr) -- (r2);
  \draw[book arrow,knowledge arrow,draw=BookTeal,line width=1pt] (yr) -- (r3);
  \draw[book arrow,knowledge arrow,draw=BookTeal,line width=1pt] (yr) -- (human);
  \draw[book arrow,knowledge arrow,draw=BookAmber,dashed] (u.north west) -- (r1.south);
  \draw[book arrow,knowledge arrow,draw=BookAmber,dashed] (u.north) -- (r2.south);
  \draw[book arrow,knowledge arrow,draw=BookAmber,dashed] (u.north east) -- (r3.south);
  \draw[BookAmber,dashed,line width=0.8pt,rounded corners=1.2mm] (75,-21) rectangle (143,0);
  \node[font=\figurefont\scriptsize,text=BookAmber,anchor=south west] at (78,3) {同一来源组};
  \node[font=\figurefont\scriptsize,text=BookAmber,anchor=west] at (78,-28) {共享影响};
\end{tikzpicture}

  \caption{独立判断依据与共享错误来源。左面板表示给定真实类别后来源独立的模型；右面板包含三条同模板规则与一名独立标注者，虚线和分组框标出三条规则的共享影响。图展示两种待检验的模型假设，不表示可以只从类别票一致率唯一恢复真实依赖结构。}
  \label{fig:knowledge-fusion-source-dependence}
\end{figure*}
\FloatBarrier

怎么\emph{量化}“相关性会改变结论”？我们固定图~\ref{fig:knowledge-fusion-source-dependence}右面板的那组票：\textbf{三条规则都报山茶，独立的$H$报茶梅}。下面用\textbf{三种不同的处理方式}去看同一组票，体会结论如何随“是否承认重复”而翻转。

\emph{第一种，假装四列独立、逐列计票}：山茶得$3$、茶梅得$1$，山茶以$3:1$胜。\emph{第二种，承认三条规则完全重复}：它们应\textbf{合并成一组}、只算\emph{一份}证据，于是变成“山茶组一票、茶梅一票”，$1:1$——结果只能是\textbf{未决}。\emph{第三种，如果还知道规则组和$H$各自的能力}，就可以\emph{比较“真值是山茶”与“真值是茶梅”分别有多大概率产生这组观察}，从而得到一个\textbf{不同于单纯票数}的支持程度（这正是概率推断，下一节细算）。

我们把第三种处理\textbf{算到底}，这是一个完整的迷你贝叶斯推断，请跟着走一遍。设两类\emph{先验各为}$0.5$；规则组与$H$都始终作答、且给定真值后相互独立；再做一个简化的\emph{对称误差}假设：规则组、$H$的二分类正确率分别为$0.7$、$0.8$（即报错率为$0.3$、$0.2$，且两个方向的错误率相同）。

\emph{假设真值是茶梅}：那么规则组\emph{报错}（报成山茶）的概率是$0.3$，$H$\emph{报对}（茶梅）的概率是$0.8$，连先验一起，这个世界的权重为
$0.5\times0.3\times0.8=0.12$。
\emph{假设真值是山茶}：规则组\emph{报对}概率$0.7$，$H$\emph{报错}概率$0.2$，权重为
$0.5\times0.7\times0.2=0.07$。
这两个权重还不是概率（和不为一），把茶梅的权重除以两者之和，才得到后验
$\mathbb{P}(Y=\text{茶梅}\mid\bm{\lambda})=0.12/(0.12+0.07)\approx0.632$。

结论很有冲击力：\textbf{表面$3:1$的多数，在承认“三条规则其实是一组”并纳入各自能力之后，后验反而更偏向茶梅（$0.632$）}。之所以能直接用“总体正确率$0.7/0.8$”代入，全靠\emph{两类错误对称}这个简化；一旦错误不对称（这正是混淆矩阵存在的理由），就必须\textbf{读取混淆矩阵的对应格子}，不能再用单个正确率。最后强调：这三种结果\emph{用的是同一组观察}，只改了“如何处理重复票与能力”，\textbf{还没有用任何独立真值证明哪一种更准}——它演示的是\emph{机制}，不是\emph{正确率排名}。

知道了依赖的存在，工程上有\textbf{由粗到细}的三种应对，复杂度和所需证据递增：
\begin{itemize}\item \emph{完全复制}时，直接\textbf{把这些来源合并成一个}，避免重复计票（最简单）。\item 当组内\emph{既有共享证据、又有各自独立证据}时，可以\textbf{先给整个来源组分配一个总权重，再在组内分配各成员的贡献}。\item 当证据足够充分时，还可以在联合模型里\textbf{加入“成对”甚至“高阶”的依赖因子}——这里的\term{因子}（factor）就是“描述多张票如何\emph{共同}起作用”的函数，后文标签模型会给出它的数学形式。\end{itemize}
这同样是一把双刃剑：\textbf{合并得太狠}，会把来源之间\emph{真实存在的差异}也一起丢掉；\textbf{依赖刻画得太细}，又会让参数和推断成本\emph{急剧增加}。要特别警惕过度承诺：弱监督里的依赖结构学习本身受\emph{数据稀疏、样本量、模型条件}的限制，\textbf{“能学出某类依赖”绝不等于“无标签样本能恢复任意关系”}\scite{varma2019dependencies}。

\section{知识的对齐}
\label{sec:knowledge-fusion-alignment}

在把票“加起来”之前，还有一道\textbf{比可靠性更靠前}的关卡——\emph{这些标注到底说的是不是同一件事}。设想两份标注分别写着“人物”和“组织”，你能立刻断定两个来源\emph{意见冲突}吗？\textbf{不能}：它们可能根本在标\emph{不同实体}（一个标“林舟”、一个标“海岬科技”，于是都对）；也可能用了\emph{不同的任务规范}（同名的标签，定义却不一样）。如果不先排除这些“假冲突”，后面的投票、加权全是在\emph{张冠李戴地做加法}。

如果说\emph{来源建模}回答的是“\textbf{谁、以怎样的能力作答}”，那么\emph{标注对齐}回答的就是另一个问题：“\textbf{哪些标注可以放在一起解释}”。它要动用对象的\emph{定位、内容、来源、所依据的规范、以及时间版本}。对齐的\emph{结果}不应该只有“配上了/没配上”两种，而应包括一张\textbf{对应映射及其状态}：可能是\emph{确定对应}、\emph{部分对应}、\emph{一对多候选}、\emph{待复核}、或\emph{无法转换}。

这里有一条硬纪律：\textbf{投票和概率标签模型，只直接接收那些“已经确定落入共同答案空间”的标注}。如果对应关系本身还没确定，\textbf{绝不能为了凑齐表格的列数而硬造一个对应}——宁可保留候选或未决。下面分三层来对齐：\emph{对象、语义、粒度与形式}。

\subsection{知识对象的对齐}
\label{sec:knowledge-fusion-object-alignment}

\emph{第一层是对“对象本身”的对齐}。先看同一句话里的“林舟”和“海岬科技”：人物标签若落在\emph{林舟}、组织标签若落在\emph{海岬科技}，那么这两个标签\textbf{可以同时正确}，根本不冲突。\term{对象对齐}（object alignment）要确定的就是：不同来源的标注，\textbf{究竟指向哪一个样本、实体、文本跨度、图像区域或时间区间}。怎么对？最可靠的是\emph{唯一标识}：如果对象 ID 在\emph{同一命名范围、同一版本}内唯一且一致，就能精确对应。一旦\emph{没有共享标识}，就得退而比较\emph{边界、位置、时间、上下文}来判断。对\emph{关系标注}还要更小心：必须\textbf{同时对齐两端实体\emph{以及}关系的方向}。“海岬科技任命林舟”和“林舟任命海岬科技”用的是\emph{同样两个实体}，但主宾一颠倒就是\textbf{完全相反}的命题，\textbf{绝不能因为实体相同就合并}。

实操上应当\textbf{按“由准到松”的顺序逐级尝试}：先做\emph{精确标识匹配}；匹配不上或不唯一，再核对\emph{边界/时间是否重叠}；\emph{仍有多个候选}时，最后才比较\emph{内容相似度}。无论走到哪一步，都要\textbf{保留原始坐标、判断依据和确定程度}，而不是只留下一个拍脑袋的结果。

尤其要\textbf{允许“不是一一对应”}：一整段标注可能对应\emph{好几条}局部标注，两个时间区间也可能\emph{只重叠一部分}。在这种情况下，如实地记成\emph{一对多}或\emph{部分对应}，往往比\textbf{强行选一条}更能准确表达现有证据。还要理解搜索范围的\emph{权衡}：把搜索范围扩大，能让更多\emph{真实对应}进入候选，但也会放进更多\emph{同名干扰对象}互相竞争，同时\emph{增加计算量}。

图~\ref{fig:knowledge-fusion-object-alignment}固定同一句短文本，只改变两个标签的对象对应。左侧虚线把不同实体配在一起，制造了虚假冲突；右侧两项标注都指向“林舟”，此时“人物”和“组织”才是同一对象上的类别分歧。关系箭头另存端点与方向，承担的不是实体类别判断，不能混入这两张类别票。

\begin{figure*}[htbp]
  \centering
  \begin{tikzpicture}[
  x=1mm,y=1mm,
  every node/.style={font=\figurefont\small,text=BookInk},
  token/.style={book node,draw=BookRule,fill=white,minimum height=8mm,inner sep=2pt},
  span a/.style={draw=BookBlue,fill=FigureInput!10,line width=0.9pt,minimum height=7mm,inner sep=1.5pt},
  span b/.style={draw=BookAmber,fill=BookAmber!10,line width=0.9pt,minimum height=7mm,inner sep=1.5pt}
]
  \node[font=\figurefont\bfseries] at (36,28) {对象错配造成伪冲突};
  \node[token,minimum width=14mm] (ll1) at (9,13) {林舟};
  \node[token,minimum width=9mm] (join1) at (24,13) {在};
  \node[token,minimum width=13mm] (sh1) at (37,13) {上海};
  \node[token,minimum width=9mm] (verb1) at (49,13) {加入};
  \node[token,minimum width=22mm] (co1) at (66,13) {海岬科技};
  \node[span a,minimum width=20mm] (a1) at (9,-2) {来源A：人物};
  \node[span b,minimum width=26mm] (b1) at (66,-16) {来源B：组织};
  \draw[BookBlue,line width=0.9pt] (ll1.south) -- (a1.north);
  \draw[BookAmber,line width=0.9pt] (co1.south) |- (b1.north);
  \draw[book arrow,knowledge arrow,draw=BookAmber,dashed] (a1.south east) -- node[below,font=\figurefont\scriptsize,text=BookAmber] {错误对应} (b1.north west);
  \node[font=\figurefont\scriptsize,text=BookMuted,align=center] at (36,-29) {标签落在不同实体上\\不能作为同一对象的类别冲突};

  \draw[BookRule,line width=0.7pt] (80,-35) -- (80,32);

  \node[font=\figurefont\bfseries] at (121,28) {对齐后识别真实分歧};
  \node[token,minimum width=14mm] (ll2) at (94,13) {林舟};
  \node[token,minimum width=9mm] (join2) at (109,13) {在};
  \node[token,minimum width=13mm] (sh2) at (122,13) {上海};
  \node[token,minimum width=9mm] (verb2) at (134,13) {加入};
  \node[token,minimum width=22mm] (co2) at (151,13) {海岬科技};
  \node[span a,minimum width=20mm] (a2) at (94,-2) {来源A：人物};
  \node[span b,minimum width=20mm] (b2) at (94,-16) {来源B：组织};
  \draw[BookBlue,line width=0.9pt] (ll2.south) -- (a2.north);
  \draw[BookAmber,line width=0.9pt] (a2.south) -- (b2.north);
  \draw[BookTeal,line width=1pt,knowledge arrow] (ll2.south east) to[out=-25,in=-155] node[below,font=\figurefont\scriptsize,text=BookTeal] {加入：主体$\to$组织} (co2.south west);
  \node[font=\figurefont\scriptsize,text=BookMuted,align=center] at (128,-29) {同一跨度上的“人物/组织”\\才是需要融合的类别分歧};
\end{tikzpicture}

  \caption{对象错配与真实标签冲突。上、下两行分别表示两个来源的标注跨度；左面板的虚线把不同实体错误配成同一对象，右面板先用实线对齐“林舟”，再显示“人物”与“组织”的类别分歧。关系箭头单独标明两端和方向。}
  \label{fig:knowledge-fusion-object-alignment}
\end{figure*}
\FloatBarrier

所有“标识之外”的匹配都要\textbf{保留一份疑义}，别配得太果断：\emph{同名}的对象可能压根是两个东西（叫“林舟”的不止一人）；\emph{同一实体}在\emph{不同时刻}的状态也可能不同（这家公司换过主营）。如果用\emph{相似度阈值}来决定是否接收对应，就必须\textbf{同时统计“漏配”和“错配”两类错误}，并\textbf{专门复核那些紧贴阈值}的样本——阈值附近的判断最不稳定。

最后厘清一个\textbf{不能互相替代}的指标：\emph{映射正确}只说明“我\textbf{找对了同一个对象}”；但两条\emph{已经准确对上}的标注，仍\textbf{可能一起标错}（对象是同一个，可惜两个来源都把类别判错了）。所以\textbf{对象对齐的成功率，替代不了标签正确率}——前者管“有没有指对东西”，后者管“对这个东西说的结论对不对”。

\subsection{知识语义的对齐}
\label{sec:knowledge-fusion-semantic-alignment}

就算对象已经对上了，\textbf{名字一样的标签仍可能在回答不同的问题}。同样是“正类”，A 来源可能只表示“\emph{文本里提到了}某内容”，B 来源却表示“\emph{该内容在现实中确实成立}”——一个是“提及”、一个是“属实”，是\textbf{两个命题}。\term{语义对齐}（semantic alignment）要做的，就是核对来源的\textbf{判定规范、任务条件、有效版本}，确认它们到底是不是\emph{同一个命题}，并明确类别之间该\textbf{怎样对应}。

记住一句话：\textbf{标签名称只是“查找线索”，不是“含义相同”的充分证据}。看到两个标签都叫“相关”，不能默认它们同义。\textbf{给两个本质不同的命题去累加票数，是无法靠“多数”消除定义差异的}——你得到的只是一个把两种口径混在一起的数字。

为了把这种“常常不是一一对应”的语义关系\textbf{如实保存}下来，需要一个比“类别对类别”更宽的记号。设来源$s$自己的类别集合为$Y_s$，目标共同空间为$Y^\star$，用映射
$h_s:Y_s\to 2^{Y^\star}$
把来源的\emph{一个}类别映到\emph{目标空间的一个候选集合}。这里$2^{Y^\star}$表示$Y^\star$的\textbf{全部子集}（即“候选可以是多个目标类别”）。

于是各种口径关系都能表达：\emph{定义相同}的类别映到\emph{单元素}集合；来源给的是\emph{上位（粗）类别}时，只能映成\emph{若干细类组成的候选集合}（“花”不能凭空指定是山茶还是茶梅）；若两个定义\emph{部分交叠、互不包含}，还得\emph{保留交叠的条件}以及\emph{无法映射的那部分}。此外还要记录这个映射\textbf{是否可逆、是否依赖上下文、对应哪个规范版本}——只有这样，将来才知道转换结果\emph{能不能还原}，以及一旦\emph{修订映射，会波及哪些历史标注}。

真正把映射\emph{用到标注上}时，顺序应当是：\textbf{先转换类别表/候选集合，再决定它能不能进入融合}。沿用上面的例子：A 的“相关”只要求“\emph{出现关键词}”，B 的“相关”要求“\emph{文本能证实命题}”，那么 A 的一张票，\textbf{最多只能为 B 提供一点旁证}，\textbf{绝不能直接算作 B 口径下的一票}。

这个障碍在处理\emph{机器自动生成}的标签时尤其隐蔽：\textbf{没有预设语义的簇编号、表示向量}同样存在语义对齐问题。第5章里那些自监督\emph{簇号}是会随轮次置换的、特征维度也是\emph{未命名}的——\textbf{编号相同、或向量维度相同，完全不建立语义对应}。要让它们可比，必须借助\emph{生成机制、一批已知对应的锚点样本、或下游目标}来\emph{标定}含义，而不能看到形状对得上就直接相加。

语义还会\emph{随时间漂移}：标注\emph{规范一旦修订}，可比较的范围就变了。最典型的是“\emph{旧类别被拆成两个新细类}”：旧规范下标的一张“茶梅属”标签，\textbf{通常只能约束“它属于这两个新细类组成的集合”}，你\textbf{无法反推}它到底是哪一个细类。于是原来那张对应关系表，\textbf{不再足以支撑新任务}。

正确的长期做法是：\textbf{保存原始类别、规范版本、以及完整的转换链}，迁移到新任务时\textbf{重新核对映射}。这样才能把两类东西分清楚——旧标注\emph{至今仍能支持的“集合级”判断}，和它\emph{给不了、必须新增证据}的“细类级”判断。反过来，最危险的偷懒是“\emph{只把编码统一成一样}”：\textbf{如果底层命题是错配的，那么概率模型建得再精巧，也只是在更优雅地组合错误命题}。

\subsection{知识粒度与形式的对齐}
\label{sec:knowledge-fusion-granularity-form-alignment}

过了对象、语义两关，还剩“\emph{粗细}”和“\emph{长什么样}”的问题。两项标注即使\emph{含义相容}，也可能一个只给\textbf{父类}（“这是山茶属”），另一个却给\textbf{细类的概率}（“山茶$0.7$/茶梅$0.3$”）。\term{粒度对齐}（granularity alignment）处理的就是这种\textbf{类别层级、对象范围粗细不同}的情形；\term{形式对齐}（form alignment）则处理\textbf{硬标签、概率、任意分数、集合}等\emph{表达方式不同}的情形。对每一种转换，我们都要问两件事：\textbf{允许做哪些转换}、以及\textbf{转换之后我到底还知道什么、丢了什么}。

这里有一条\textbf{方向性}极强的规律。当若干子类\emph{互斥且穷尽}父类时，\textbf{从细到粗}是安全的：把各子类概率\emph{相加}就得到父类概率（山茶属概率 = 山茶概率 + 茶梅概率 + …）。但\textbf{反方向不成立}：\textbf{只有一个父类标签时，没有任何依据能把这份质量“摊派”给某个细类}——知道“是山茶属”，推不出“山茶$0.7$”。信息只能从细向粗无损汇总，不能无中生有地向细拆分。

“对象范围”的转换同样\textbf{不能只看形式，要核对任务语义}。举例：如果“文档为正”的定义是“\emph{至少出现一段正向证据}”，那么任一处局部正例都能\emph{合法地推出}文档为正；可一旦定义改成“\emph{全文都支持该命题}”，\textbf{同样那一条局部标注就不够用了}。局部到整体能不能推，\textbf{取决于聚合规则，而不是算术}。

几类常见的形式转换也各有前提：\emph{概率转硬标签}需要一个明确的\textbf{接收阈值或决策损失}；\emph{任意分数转概率}则需要\textbf{尺度变换和校准依据}——这里的\term{校准}，特指检查“报出的概率\textbf{是否真的对应事件发生的频率}”（报$0.8$的那些样本，是否约有八成真的发生）。最后别忘区分三种“没有类别值”：\emph{没调用、调用失败、运行后主动弃权}，它们\textbf{数值上都表现为“无类别”，生成条件却完全不同}，转换时必须\textbf{分开保存}，不能合并成一个“空”。

表~\ref{tab:knowledge-fusion-form-conversions}把这些转换\emph{逐行}列清楚：每一行都写明\textbf{在什么条件下能转}、转换后\textbf{保留了}哪些类别/概率/对应关系、以及\textbf{哪些细类或不确定性是无法再恢复}的。建议读这张表时\textbf{重点盯最后一列}——它明确告诉你信息\textbf{丢在哪}。

要点是：\textbf{不可逆转换依然可以做}（接口往往只收硬标签），但做的时候\textbf{必须留下原始标注和所用配置}，这样日后才\textbf{追溯得到“在哪一步、把什么细类或不确定性丢掉了”}。而如果某类标注\textbf{根本找不到进入共同答案空间的合理映射}，正确做法不是硬转，而是\textbf{保留它原来的监督对象}，交给第~\ref{sec:knowledge-fusion-heterogeneous-use}节的\emph{异构知识联合训练}去处理——那里允许不同形态的监督各走各的损失，而不必压成同一个标签。

\begin{table*}[htbp]
  \centering
  \small
  \renewcommand{\arraystretch}{1.15}
  \begin{tabularx}{\linewidth}{@{}>{\raggedright\setlength{\parindent}{0pt}\arraybackslash}p{28mm}>{\raggedright\setlength{\parindent}{0pt}\arraybackslash}X>{\raggedright\setlength{\parindent}{0pt}\arraybackslash}X>{\raggedright\setlength{\parindent}{0pt}\arraybackslash}X@{}}
    \toprule
    \textbf{转换} & \textbf{操作与条件} & \textbf{保留的标注属性} & \textbf{丢失或待补的标注属性} \\
    \midrule
    细类到粗类 & 合并互斥且穷尽的子类，概率求和 & 父类归属及总概率 & 原细类及其概率分配 \\
    粗类到候选细类 & 保留父类下全部合法子类 & 真值所属父类 & 具体子类，需新增证据 \\
    概率到硬标签 & 按阈值或决策损失选类 & 行动或训练类别 & 类别支持程度与冲突强度 \\
    局部到整体 & 按任务规定的聚合规则转换 & 规则定义的整体结论 & 未保留的局部位置与数量 \\
    异构表示到共同空间 & 依据配对锚点、生成关系或可验证映射 & 映射覆盖的共同结构 & 尚未确定的语义对应 \\
    \bottomrule
  \end{tabularx}
  \caption{标注形式转换的条件与标注属性损失。“丢失”指转换结果不再保留的细类、概率或局部位置，“待补”指预测输入原本就不足以确定的细类或局部位置；二者都不能从转换后的结果凭空恢复。}
  \label{tab:knowledge-fusion-form-conversions}
\end{table*}
\FloatBarrier

\section{基于记录与预测的融合}
\label{sec:knowledge-fusion-inference}

完成了来源建模和三层对齐，我们终于可以\textbf{正式“做加法”}了。本节先做一个明确的前提约定：把表~\ref{tab:knowledge-fusion-source-records}里的类别票\textbf{视为已经对齐的标签}——同一行就是同一张照片，所有类别都已落在共同空间$Y$里。

在“\emph{固定住这批待融合标签}”的条件下，我们并排比较\textbf{三种由浅入深的组合办法}，请特别注意它们\emph{各自用了多少、用了哪些假设}：
\begin{itemize}\item \term{投票}：给\textbf{每一张有效票相同的贡献}，只数人头；\item \term{加权}：承认来源有强有弱，\textbf{改变各张票的贡献大小}；\item \term{概率推断}（probabilistic inference）：反过来问“\textbf{假如真值是某类，有多大概率产生眼前这组票}”，用生成模型+混淆矩阵来推断。\end{itemize}
此外，当一张样本包含\emph{多个相互关联的局部答案}时（如一句话里的多个实体标签），还会受到\emph{整体结构约束}，这是再往后的话题。

比较这三种办法时请保持一个平常心：\textbf{它们的区别是“计算方式和所依赖的假设不同”，绝不是“算法越复杂质量越高”}。而且它们\textbf{不是互斥的}，实际系统里常常组合使用。

\subsection{基于投票的知识融合}
\label{sec:knowledge-fusion-voting}

最朴素、也最该先当\textbf{基线}的办法，是在“\emph{暂时还没有任何来源能力参数可用}”时，只问一个最简单的问题：\textbf{每个类别各得到了多少来源的支持}？\term{多数投票}（majority voting）把\emph{每一份已对齐的有效票计为一票}，再从共同集合$Y$里\textbf{选出票数最多}的类别。对样本$i$，要先\emph{逐个候选类别}$k$\emph{把票数理出来}，然后才能做“取最大”这个比较：
\begin{equation}
  v_i(k)
  =\sum_{s=1}^{m}o_{i,s}\mathbb{I}\{\lambda_{i,s}=k\},
  \qquad
  \widehat y_i\in\argmax_{k\in Y}v_i(k).
  \label{eq:knowledge-fusion-majority-vote}
\end{equation}
读懂这个式子的几个细节：$o_{i,s}$负责\textbf{把“没有有效答案”的标注排除在外}（弃权、失败都不进票箱）；指示函数$\mathbb{I}\{\cdot\}$\textbf{只给来源实际报的那个类别加一}，其余类别不加；最后的$\argmax$选出最高票类别，但要注意\textbf{最高票可能不止一个}（平票），此时集合里有多个元素，必须另有规则才能定夺。

还要区分两种“多数”：\term{相对多数}（plurality）\textbf{只要求票数最多}，哪怕不到一半也算赢；\term{绝对多数}（majority）则\textbf{要求超过有效票的一半}。换任务时规则也要换：\emph{多标签}任务可以\textbf{逐标签独立计票}，但选出的多个标签\textbf{组合}起来仍要检查\emph{互斥、层级、序列}等约束；至于\emph{关系标签}，必须先\textbf{对齐成同一个有向命题}才能计票，\textbf{不能}直接套用实体类别的计票方式。

投票有两个\textbf{天生给不出答案}的情形，公式自己解决不了，必须\textbf{另外规定“接收规则”}：一是\emph{最高票并列}（平票），二是\emph{一张有效票都没有}（零票）。表~\ref{tab:knowledge-fusion-comparison}统一采用这样一条相对多数规则：\textbf{唯一最高票才输出该类别；一旦平票或零票，就记为“未决”}，并且\textbf{不额外设置最低票数门槛}。

读表中“票数比例”时要清醒：它的分母是\emph{有效票数}（弃权、失败\textbf{不计}），所以它表达的只是“\textbf{来源之间的赞同比例}”，\textbf{不是}一个经过校准的、“真实正确的概率”。“三票里两票山茶”等于$2/3$赞同，\textbf{不等于}“山茶有$2/3$概率为真”——后者需要概率模型和校准，投票本身给不了。

\begin{table*}[htbp]
  \centering
  \small
  \renewcommand{\arraystretch}{1.15}
  \begin{tabularx}{\linewidth}{@{}ccccc>{\raggedright\setlength{\parindent}{0pt}\arraybackslash}X@{}}
    \toprule
    \textbf{样本} & \textbf{有效票} & \textbf{山茶/茶梅票} & \textbf{等权结果} & \textbf{加权结果} & \textbf{关键原因} \\
    \midrule
    $x_1$ & $4$ & $4/0$ & 山茶 & 山茶 & 全部来源一致 \\
    $x_2$ & $3$ & $2/1$ & 山茶 & 茶梅 & 模型权重超过两条规则之和 \\
    $x_3$ & $4$ & $3/1$ & 山茶 & 山茶 & 人工与规则共同支持山茶 \\
    $x_4$ & $1$ & $1/0$ & 山茶 & 山茶 & 只有模型作答 \\
    $x_5$ & $3$ & $2/1$ & 山茶 & 茶梅 & 两张山茶票共享模板 \\
    $x_6$ & $4$ & $2/2$ & 未决 & 茶梅 & 人工与模型共同支持茶梅 \\
    $x_7$ & $0$ & $0/0$ & 未决 & 未决 & 全部来源缺失或弃权 \\
    $x_8$ & $4$ & $2/2$ & 未决 & 山茶 & 人工与模型共同支持山茶 \\
    \bottomrule
  \end{tabularx}
  \caption{同一组候选的融合结果对照。加权列使用教学权重$w_{R_1}=w_{R_2}=0.25$、$w_H=1.2$、$w_M=0.6$，只用于复核运算与展示假设差异，不表示这些权重经过真实审计，也不证明加权方法普遍优于投票。}
  \label{tab:knowledge-fusion-comparison}
\end{table*}
\FloatBarrier

\emph{逐行复核}表~\ref{tab:knowledge-fusion-comparison}，能清楚看到三种结局：\emph{唯一最高票}（直接出类）、\emph{平票}（$x_6,x_8$的$2:2$，记未决）、\emph{零票}（$x_7$全部缺失或弃权，记未决）。

而$x_5$这一行\textbf{专门用来暴露“等权投票”的要害}：两条\emph{共享模板}的规则都报山茶，于是这一张山茶票被\textbf{重复计成了两票}，可它\textbf{并没有新增任何独立判断}——$x_5$里真正独立的人工票（茶梅）就这样被“复制票”淹没了。这正好呼应6.2节的相关性。

那么投票\textbf{为什么仍值得作为第一个基线}？因为它\textbf{不需要拟合任何来源参数、计算透明、每一张票都能人工复核}，是衡量后续复杂方法“到底有没有净收益”的参照。但代价同样明确：它\textbf{同时忽略了三件关键的事——来源能力的强弱、错误的方向、以及来源之间的依赖}。所以结论要讲全：\textbf{新增一个来源是否能改善正确性，取决于它的能力、错误方向、覆盖以及是否独立}，\textbf{绝不能仅凭“票数变多了”就推出“更准了”}。

工程上还应\textbf{把“计票”和“选类”两段成本分开核算}，它们复杂度不同。若直接遍历整张$n\times m$稠密标签矩阵来累计票数，需要$O(nm)$；累计完再逐类比较$K$个候选，总时间为$O(nm+nK)$。如果只\emph{稀疏地}保存有效票（利用$o_{i,s}=1$），则降为$O(N_{\mathrm{obs}}+nK)$，其中$N_{\mathrm{obs}}=\sum_{i,s}o_{i,s}$是实际有效票总数。内存方面：\emph{逐样本}边读边算只需$O(K)$工作内存；若要\emph{保留全部}样本的票数则需$O(nK)$。

请注意这些复杂度界\textbf{只覆盖“累计票数+选类别”这一步}，并\textbf{不包含}前面更费劲、也更贵的\emph{语义对齐}和\emph{相关来源去重}——那两块的成本往往才是大头，不能因为计票是线性的就以为整个融合很便宜。

\subsection{基于加权的知识融合}
\label{sec:knowledge-fusion-weighting}

如果我们\emph{已经通过审计或来源模型知道：几个来源的贡献并不相同}（比如$H$明显比同模板规则可信），那“每列一票”就\textbf{白白浪费了这条能力信息}。\term{加权融合}（weighted fusion）的思路很直接：\textbf{给每张有效票乘上一个明确的贡献系数}，再比较各类别累计的\emph{加权得分}。沿用前面那批待融合标签，把“来源$s$在样本$i$上支持类别$k$”这件事的权重记为$w_{i,s,k}$，硬标签下的得分与选类规则为：
\begin{equation}
  s_i(k)
  =\sum_{s=1}^{m}o_{i,s}w_{i,s,k}
    \mathbb{I}\{\lambda_{i,s}=k\},
  \qquad
  \widehat y_i\in\argmax_{k\in Y}s_i(k).
  \label{eq:knowledge-fusion-weighted-vote}
\end{equation}
和等权投票相比，选类动作没变（仍在$Y$中挑累计支持最高者），但“支持的\textbf{单位}”从“一票”变成了“\textbf{一份权重得分}”。权重$w_{i,s,k}$设计得很灵活：可以\emph{全局固定}一个值，也可以\emph{随类别、区域、时间}变化（比如某来源只在简单区给高权重）。

但灵活意味着必须\textbf{交代来历}：每个权重都要\textbf{说清它来自哪一批审计、哪个来源模型、或哪一条人工设计规则}，以及得分的\textbf{量纲}（不同量纲的权重不能混用）。还有一条防止“自欺欺人”的纪律：\textbf{用来挑选/拟合权重的那批样本，属于开发依据，绝不能再拿去当最终质量评价}——否则就是用考题调参、再用同一套考题报成绩。

硬标签有个浪费：它只在一个类别上给支持，把来源“\emph{其实也有点怀疑别的类}”的信息全丢了。如果来源能给出\emph{类别概率}$p_s(k\mid x_i)$，我们就能\textbf{保留它分配给所有类别的全部质量}。做法是改用\emph{同一照片、同一类别定义}下的\textbf{概率向量}，用非负的来源权重$w_{i,s}$做\emph{线性汇总}（此时$o_{i,s}=1$表示该来源给出了\emph{可用的完整分布}）。最后\textbf{除以有效权重之和}，是为了让混合结果\textbf{仍然是一个各类概率之和为一}的合法分布：
\begin{equation}
  \widetilde p_i(k)
  =\frac{\sum_{s=1}^{m}o_{i,s}w_{i,s}p_s(k\mid x_i)}
  {\sum_{s=1}^{m}o_{i,s}w_{i,s}}.
  \label{eq:knowledge-fusion-weighted-probability}
\end{equation}
使用这个公式有几个\textbf{不可省略的前提}：各来源必须使用\textbf{可比较的类别语义、概率尺度和校准含义}（回到6.3的对齐！），权重非负、分母为正。\textbf{一旦没有任何有效正权重（分母为零），正确动作是保留“未决”，而不是硬算一个$0/0$}。

还要从概念上分清两件常被混淆的事：式~\eqref{eq:knowledge-fusion-weighted-probability}做的是\textbf{对各来源预测分布的“线性混合（求平均）”}，它\textbf{不是}后面那种“把来源\emph{似然相乘}”的\emph{生成式推断}——平均别人的概率，和在生成模型里反推真值，是两套逻辑。最后重申校准：\textbf{模型随便输出的一个分数不能直接当概率}；即便它加起来等于一，也只\textbf{满足了概率的“形式”}，并\textbf{不自动保证}它和真实发生频率\textbf{校准}。

回到表~\ref{tab:knowledge-fusion-comparison}的硬标签那一列，我们人为预设一组教学权重$w_{R_1}=w_{R_2}=0.25$、$w_H=1.2$、$w_M=0.6$，其意图正是\textbf{让那两条同模板规则的“合计”贡献压低}。于是在$x_5$上：山茶只得$0.25+0.25=0.5$，茶梅得$1.2$，加权结果\textbf{翻转为茶梅}——纠正了等权投票被复制票带偏的问题。表中其他行也都可以照此复算。

但请务必注意这组数字的\textbf{身份}：它们\textbf{只是教学配置}，\textbf{并不是从前四行审计正确率拟合出来的}。所以这里“加权改判成茶梅”，\textbf{只能演示“加权能利用能力差异、能抵抗复制票”这个机制}，\textbf{不能据此宣称}“加权法普遍更好”、更\textbf{不能说}这就\textbf{证明了人工答案正确}。

权重该\textbf{从哪来}？这里有一个\textbf{非常诱人、但错误}的取巧办法：直接拿来源当前输出的\textbf{最大概率}当权重（它越自信，权重越高）。这等于\textbf{把“类别票的集中程度”偷换成了“来源的真实能力”}，而\textbf{自信和正确是两回事}。

用一个算例看清这个错误。设一张\emph{参考真值是茶梅}的照片：自信但错向山茶的$C$给$p_C=(0.99,0.01)$，不那么自信但更准的独立来源$D$给$p_D=(0.25,0.75)$（分量依次为山茶、茶梅）。
\begin{itemize}\item 若\textbf{用最大概率}$0.99,0.75$\textbf{当权重}：山茶的未归一化得分$0.99\times0.99+0.75\times0.25=1.1676$，茶梅$0.99\times0.01+0.75\times0.75=0.5724$，于是\textbf{选山茶——错}。\item 若权重来自\textbf{独立审计}（$C,D$分别为$0.2,0.8$）：山茶得分$0.2\times0.99+0.8\times0.25=0.398$，茶梅$0.2\times0.01+0.8\times0.75=0.602$，于是\textbf{选茶梅——对}。\end{itemize}
两次计算\textbf{固定了完全相同的分布和真值，只改变“权重的依据”}，结论就相反。这说明：\textbf{一个异常集中（高置信）的概率，完全可能正信心十足地指向错误类别}。审计权重仍需核对其\textbf{代表性}，概率本身仍需\textbf{校准检验}，\textbf{不能从这个教学例推出某种加权法有普遍收益}。

把权重\emph{细化到“分类别、分区域”}固然更精准，但每一分细化都\textbf{要求有对应的参考样本和作答覆盖清单}来支撑，同时\textbf{增加更新和维护成本}——又回到6.2节“分组越细、每行证据越少”的权衡。至于\textbf{负权重}，它的含义是把某张票当成\textbf{反向证据}（“它越支持，我越该警惕”），这\textbf{需要额外的模型来证明这种反向关系}，而且\textbf{不能}用在前面那个要求非负、和为一的\emph{概率线性混合}里。

怎么检验一组权重靠不靠谱？可以做\textbf{敏感性分析}：\emph{删掉}某一个来源、\emph{合并}相关来源、\emph{更新}审计数据、或小幅\emph{扰动}权重，看最终选类会不会\textbf{轻易翻转}；如果结论被一个权重的微小改动左右，就说明它不稳健。最后点明加权在全章脉络中的位置：\textbf{加权能表达“来源贡献不同”，但它仍然没有回答“潜在真值究竟是如何\emph{产生}出这些观察票的”}——补上这个\emph{生成关系}的，正是下一小节的概率模型。

\subsection{基于概率模型与约束的知识融合}
\label{sec:knowledge-fusion-probabilistic-constrained}

投票和加权都是在“\emph{票}”这个表层做文章，而概率推断要\textbf{换一个更根本的视角：反过来建模“票是怎么被产生出来的”}。这一小节目处理两类\emph{不同}的问题，先把它们分清楚，避免混成一个方法。

\emph{第一类}：当看到\textbf{A报山茶、B报月季}这种\emph{彼此矛盾}的票时，6.2节的混淆矩阵就能派上用场——我们可以\textbf{比较“真值到底是哪一类时，最可能产生眼前这组观察”}。这是\term{概率模型}处理的问题，它把\textbf{真实标签设成一个“未观测变量（潜变量）”}来反推。

\emph{第二类}：当一句话\emph{每个位置}都已经有了各自的最高分标签后，这些标签\textbf{拼在一起可能根本不合法}（比如实体内部标记$I$前面没有实体开始标记$B$）。这种“\emph{整体组合是否合法}”由\term{结构约束}处理，它\textbf{限定或惩罚}不合法的整体组合。

两者\textbf{合用}时的正确分工是：\textbf{先用概率模型算出每个局部位置的支持，再在硬/软约束下挑出一整组彼此相容的标签}。千万\textbf{不能把“局部概率高”和“整体组合合法”当成同一条要求}——一个处处局部最优的序列，完全可能是个非法序列。

\subsubsection{潜在真实标签的概率推断}
\label{sec:knowledge-fusion-latent-label-inference}

现在正面解决“\emph{标签已经对齐、却没有参考真值}”的局面：我们手上只有一张标签矩阵，却想\textbf{同时}搞清两件纠缠在一起的事——\emph{每个样本的真实类别}、以及\emph{每个来源的犯错规律}。这有点像“鸡生蛋、蛋生鸡”：知道真值就能统计来源错误，知道来源错误又能反推真值。\term{Dawid--Skene模型}（Dawid--Skene model）正是破解这个循环的经典工具。

它的核心动作是把真实类别$Y_i$设为一个\term{潜变量}——即\textbf{没有被直接观测到、只能由模型推断}的量；再用6.2节那张\textbf{混淆矩阵}$\theta^{(s)}_{k,\ell}$描述“每个来源\emph{倾向于怎样错报}这个潜变量”。模型的零件有：类别\textbf{先验}$\pi_k=\mathbb{P}(Y_i=k)$、沿用式~\eqref{eq:knowledge-fusion-confusion-matrix}的混淆参数$\theta^{(s)}_{k,\ell}$，以及$S_i$（\textbf{实际}给样本$i$投了有效票的来源集合，因为不同样本被哪些来源作答可能不同）。

但天下没有免费的午餐，这个模型\textbf{带着一串必须摆明的假设}：\textbf{样本之间独立同分布}；\textbf{每个来源的混淆行为在建模范围内稳定}；最关键的是\textbf{给定真值$Y_i$之后各来源条件独立}（就是6.2节那条能把似然相乘的依据）。关于“谁来作答”还有个讲究：如果\textbf{不另外为“缺失/弃权”建模}，就要求$S_i$是\textbf{事前由外生设计固定}的（比如随机分派），或这种缺失\textbf{对当前似然可以忽略}。\textbf{反例}是：如果一条规则\emph{更容易在山茶照片上触发}，那么“它\emph{没}投票”本身就泄漏了物种信息，这时\textbf{事后看到的}$S_i$\textbf{不能被当作可忽略}。

把上面那条“谁作答”的假设说得更具体些：如果\emph{事前随机}决定每张照片交给哪些来源，分派\emph{不依赖未知类别}，那我们就可以\textbf{老老实实在这份分派名单$S_i$之内}计算出票的概率，名单外的来源不参与。反过来，如果触发与否\textbf{依赖类别}（规则更常在山茶上触发），那么“\textbf{没有投票}”这件事\textbf{本身就在改变对物种的判断}，你若\textbf{只读取已投出的票}，就会\textbf{漏掉“覆盖机制”所提供的那条类别线索}。

下面的推导\textbf{建立在“覆盖可忽略”等前述假设都成立}的前提下。逻辑分两步：先在“\emph{假定真值为某类}”下，用\textbf{类别先验}\emph{乘以}\textbf{各个来源在该真值下出票的似然}，得到“真值+观察”的\textbf{联合概率}；但真值我们并不知道，所以再\textbf{对全部$K$个候选真值求和}，才得到一个\textbf{只含看得见的标签}的概率：
\begin{equation}
  p(Y_i=k,\bm{\lambda}_i)
  =\pi_k\prod_{s\in S_i}\theta^{(s)}_{k,\lambda_{i,s}},
  \qquad
  p(\bm{\lambda}_i)
  =\sum_{k=1}^{K}\pi_k
    \prod_{s\in S_i}\theta^{(s)}_{k,\lambda_{i,s}}.
  \label{eq:knowledge-fusion-dawid-skene}
\end{equation}
这个公式的两项要分开理解：\emph{第一项}（等号左）把“这个类别有多常见（先验）”和“在这个类别下\emph{各来源恰好这样投票}有多可能（似然）”\textbf{合并}起来，是针对\emph{某个具体候选真值}$k$的；\emph{第二项}（等号右）因为真值未知，\textbf{把所有$K$种可能的真值加总}，得到“\textbf{不管真值是什么、总之看到当前这组票}”的\textbf{总概率}（边缘似然）。这个总概率正是我们\textbf{比较不同参数谁更能解释数据}的依据。

于是估计目标就清楚了：在“先验非负且和为一、每张混淆矩阵逐行归一化”的\emph{合法参数集合}里，我们要\textbf{让所有样本观测概率的乘积（似然）最大}；由于连乘求极值不便，通常\textbf{等价地最大化它的对数（对数似然）}。

难点在哪？注意\textbf{求和号是套在对数\emph{里面}}的（要先对潜在真值求和、再取对数），这使得我们\textbf{没法像“真值已知”那样直接数计数}——真值是隐藏的。Dawid 与 Skene 的解决方案，就是用下一小节的\textbf{EM 算法}来\emph{交替}地估计“观察者误差率”和“潜在类别”\scite{dawid1979observer}。

这里正式引出\term{期望最大化}（expectation-maximization, EM）算法，它专门对付“\emph{含潜变量的极大似然}”这种\textbf{真值未知}的局面。\textbf{核心思想是“鸡和蛋交替着来”}，两步循环：
\begin{itemize}\item \textbf{E 步（猜真值）}：\emph{先用当前这一版参数}，给每个\emph{未知类别}\textbf{分配一个“它属于各类的概率”}（软的，不是非黑即白）；\item \textbf{M 步（改参数）}：再把这些概率当成\textbf{分数权重（软计数）}，去\textbf{更新先验和混淆矩阵}。\end{itemize}
然后反复迭代，直到稳定。这种“交替更新如何保证不降低观测似然”的一般结论，可参考\BookChapterRef{chap:pgm-parameter-learning}{卷六《模型：概率图模型》的“参数学习”章}。

E 步真正要算的量叫\term{责任度}（responsibility）$q_{i,k}^{(t)}$：在第$t$轮、\textbf{已经看到这组票之后}，模型认为“样本$i$属于类别$k$”的\textbf{后验概率}。它的算法就是\textbf{贝叶斯公式}：拿“候选类别$k$的联合概率”（上一条公式的第一项）\textbf{除以“所有$K$个类别联合概率之和”}做归一化，得到一组\textbf{和为一、可直接当作下一步软计数权重}的数：
\begin{equation}
  q_{i,k}^{(t)}
  =
  \frac{\pi_k^{(t)}
    \prod_{s\in S_i}
    \left(\theta^{(s)}_{k,\lambda_{i,s}}\right)^{(t)}}
  {\sum_{r=1}^{K}\pi_r^{(t)}
    \prod_{s\in S_i}
    \left(\theta^{(s)}_{r,\lambda_{i,s}}\right)^{(t)}}.
  \label{eq:knowledge-fusion-dawid-skene-e-step}
\end{equation}
读这个 E 步公式有几个要点。\textbf{分母必须为正}（至少要有可能的解释）；而且它\textbf{比较的是共同集合$Y$里的\emph{全部}$K$个类别}，\textbf{不是}只在“\emph{被投过票}”的几个类别里比——某个类别这次没人投，不代表它不可能是真值。一个边界情形：若$S_i$为空（这个样本\emph{没有任何来源作答}），公式里的\emph{空乘积约定为一}，于是后验\textbf{直接退回类别先验}——这准确表达了“\textbf{这一行没提供任何类别证据，只能靠各类的常见程度}”。

接着是 M 步：把上一步算出的$q^{(t)}$\textbf{固定住}，去最大化“含潜在标签的联合对数似然”的后验期望。直觉上它就是\textbf{把硬计数换成软计数}：每个样本\textbf{按$q_{i,k}^{(t)}$这么多的“分数份”}贡献给类别$k$（而不是整张算给某一类）；在这份软归属下，\textbf{报了票$\ell$}的样本再相应地贡献到来源混淆矩阵的第$(k,\ell)$格。在“先验、每行混淆都归一化”的约束下，\textbf{加权计数除以相应总权重}，就得到更新公式：
\begin{equation}
  \pi_k^{(t+1)}
  =\frac{1}{n}\sum_{i=1}^{n}q_{i,k}^{(t)},
  \qquad
  \left(\theta^{(s)}_{k,\ell}\right)^{(t+1)}
  =
  \frac{\sum_{i:s\in S_i}
    q_{i,k}^{(t)}\mathbb{I}\{\lambda_{i,s}=\ell\}}
  {\sum_{i:s\in S_i}q_{i,k}^{(t)}}.
  \label{eq:knowledge-fusion-dawid-skene-m-step}
\end{equation}
M 步这两个更新的\textbf{分母又不一样}，正好呼应6.2节对“分母”的强调：\emph{先验}更新是在\textbf{全部$n$个样本}上平均类别责任度（每个样本无论谁作答都参与）；而\emph{混淆矩阵}更新\textbf{只使用来源$s$实际作答}的那些样本，其分母是“这些样本归属到真值$k$的软计数之和”。当这个分母为零（来源$s$从没在“疑似$k$类”的样本上作答），这一行\textbf{没有得到有效权重}，应\textbf{保留未估计或加入明确先验}，\textbf{不能硬算$0/0$}。

再补三个使用要点。其一，\textbf{初始化}可以用多数票、少量审计标签、或近似对角（假设大家基本靠谱）的矩阵，但它们\textbf{各自注入了不同的方向性假设}，可能影响最终落点。其二，\textbf{停止条件}要\textbf{同时}设“\emph{最大迭代次数}”和“\emph{似然/参数变化小到阈值}”，防止无限迭代或过早停在抖动中。其三，也是最容易被算法成功感掩盖的：\textbf{EM 的更新公式不会自动修复模型缺陷}——\emph{共享错误}（来源相关）、\emph{选择性覆盖}（谁作答依赖类别）、\emph{从未观测到的类别-来源组合}，这些\textbf{模型设定与证据上的缺口}都必须\textbf{另外检查}，\textbf{不会因为迭代收敛而消失}。

\textbf{光看公式容易飘，我们用一组完整的小数字，把一次 E 步、一次 M 步亲手算一遍}，顺便看清“概率推断”和“数票”到底差在哪。另取\emph{三个始终作答}的来源，仍是山茶、茶梅二分类；为便于手算，设它们\emph{初始对称正确率}分别为$0.8,0.7,0.6$（报错率即$0.2,0.3,0.4$），两类先验各$0.5$。三个教学样本的票依次是：
\begin{itemize}\item 样本1：$(\text{山茶},\text{山茶},\text{山茶})$（三票全山茶）；\item 样本2：$(\text{山茶},\text{茶梅},\text{茶梅})$（一票山茶、两票茶梅）；\item 样本3：$(\text{茶梅},\text{茶梅},\text{山茶})$（两票茶梅、一票山茶）。\end{itemize}第一行在山茶与茶梅下的权重为$0.168$与$0.012$，分别来自$0.5\times0.8\times0.7\times0.6$和$0.5\times0.2\times0.3\times0.4$；第二行两类权重为$0.048,0.042$，第三行为$0.018,0.112$。逐行归一化，山茶责任度为$0.9333,0.5333,0.1385$，茶梅为各自补数。请\textbf{特别盯住样本2}：它明明有\textbf{两票茶梅、只有一票山茶}，按多数投票必然判茶梅；可在当前这版“来源能力参数”下，它给\textbf{山茶的后验（$0.5333$）反而略高于茶梅（$0.4667$）}。为什么？因为那\emph{唯一}投山茶的是\emph{最可信的来源1（正确率$0.8$）}，而两个投茶梅的来源可信度更低（$0.7,0.6$），它们\textbf{同时看错}的可能性，压过了“最可信来源单独看错”的可能性。这就是\textbf{来源似然与单纯票数最直观的区别}：\emph{概率推断让“谁投的”参与了判断，而投票只看“投了什么”}。（当然，这个反转依赖初始的对称能力假设，只是教学演示，不代表真值就是山茶。）

用未舍入责任度更新，山茶与茶梅的软计数总量约为$1.6051,1.3949$，除以三个样本得到先验$(0.5350,0.4650)$。来源1在前两行报山茶，这两行分给真值山茶的权重合计为$0.9333+0.5333$，除以山茶总软计数得$0.9137$；同样，两行分给真值茶梅的权重合计约$0.5334$，除以茶梅总软计数得$0.3824$。这些分数不是参考计数，而是当前后验下的期望计数。记第1次M步后的矩阵为$\bm{\Theta}_s^{(1)}=\left(\left(\theta^{(s)}_{k,\ell}\right)^{(1)}\right)_{k,\ell}$，按行是真值山茶/茶梅、列是类别票山茶/茶梅排列，得到
\[
  \begin{aligned}
    \bm{\Theta}_1^{(1)}
    &=
    \begin{pmatrix}
      0.9137 & 0.0863 \\
      0.3824 & 0.6176
    \end{pmatrix},\\
    \bm{\Theta}_2^{(1)}
    &=
    \begin{pmatrix}
      0.5815 & 0.4185 \\
      0.0478 & 0.9522
    \end{pmatrix},\\
    \bm{\Theta}_3^{(1)}
    &=
    \begin{pmatrix}
      0.6677 & 0.3323 \\
      0.6654 & 0.3346
    \end{pmatrix}.
  \end{aligned}
\]
核对一下：每张矩阵\textbf{逐行求和都是一}（它们仍是合法的条件分布）。注意这一版更新的是\textbf{先验加上三个来源的“完整混淆矩阵”}——每个来源都有“山茶行、茶梅行”两行，\textbf{而不是只更新一个笼统的总体正确率}。还有一个细节：我们\emph{初值}假设的是\emph{对称误差}，但 M 步公式\textbf{并没有强制更新后仍对称}，所以你能看到来源1的“山茶行正确率$0.9137$”和“茶梅行正确率$0.6176$”\textbf{明显不同}——数据把对称假设打破了。最后务必清醒：这一轮\textbf{没有加平滑、只有三个样本}，数字\textbf{仅仅用来复核“计算流程对不对”}，\textbf{远不足以稳定地恢复}来源的真实能力（真实使用需要多得多的样本）。

图~\ref{fig:knowledge-fusion-latent-label-em}把这轮计算接成一个\textbf{闭环}，建议对照记忆：\textbf{标签矩阵$\bm{\Lambda}$始终固定不动}（原始观测不会变），\textbf{E 步只改责任度}$q$，\textbf{M 步只改参数}$\pi,\theta$，然后\textbf{检查停止条件}，不满足就带着新参数回到 E 步。

关于 EM 有三句必须分清的“它能保证什么、不能保证什么”，请\textbf{刻在脑子里}：
\begin{itemize}\item \textbf{能保证单调不降}：在精确执行 E、M 步时，\emph{当前模型的观测似然不会被一轮更新降低}（这是 EM 的经典性质）。\item \textbf{不保证全局最优}：目标函数一般\emph{非凸}，算法\textbf{可能停在局部最优或鞍点}附近，换个初始化也许得到另一个解。\item \textbf{“收敛”不等于“正确”}：\textbf{更新稳定}只回答了“\emph{数值迭代是否趋稳}”这个\emph{优化}问题；它\textbf{既不回答参数是否\emph{可识别}}（还记得“全对/全报反”吗），\textbf{更不证明推断出的类别符合现实真值}。\end{itemize}

\begin{figure}[!t]
  \centering
  \begin{tikzpicture}[
  x=1mm,y=1mm,
  every node/.style={font=\figurefont\small,text=BookInk},
  data/.style={book node,draw=BookBlue,fill=FigureInput!10,line width=0.9pt,minimum width=42mm,minimum height=11mm,inner sep=2pt},
  state/.style={book node,draw=BookTeal,fill=FigureModel!10,line width=0.9pt,minimum width=45mm,minimum height=11mm,inner sep=2pt},
  check/.style={book node,draw=BookAmber,fill=BookAmber!10,line width=0.9pt,minimum width=34mm,minimum height=10mm,inner sep=2pt},
  output/.style={book node,draw=BookInk,fill=BookPaper,line width=0.9pt,minimum width=39mm,minimum height=10mm,inner sep=2pt}
]
  \node[data] (matrix) at (0,34) {固定观测$\bm{\Lambda},\bm{O}$};
  \node[state] (init) at (0,17) {初始化$\bm{\pi}^{(0)},\bm{\Theta}^{(0)}$};
  \node[state] (estep) at (0,0) {E步：计算$q_{i,k}^{(t)}$};
  \node[state] (mstep) at (0,-17) {M步：更新$\bm{\pi}^{(t+1)},\bm{\Theta}^{(t+1)}$};
  \node[check] (stop) at (0,-34) {似然或模型参数已稳定？};
  \node[output] (result) at (0,-51) {后验分布与来源模型参数};

  \draw[book arrow,knowledge arrow,draw=BookBlue,line width=1pt] (matrix) -- (init);
  \draw[book arrow,knowledge arrow,draw=BookTeal,line width=1pt] (init) -- (estep);
  \draw[book arrow,knowledge arrow,draw=BookTeal,line width=1pt] (estep) -- node[right,font=\figurefont\scriptsize] {软类别计数} (mstep);
  \draw[book arrow,knowledge arrow,draw=BookTeal,line width=1pt] (mstep) -- (stop);
  \draw[book arrow,knowledge arrow,draw=BookInk,line width=1pt] (stop) -- node[right,font=\figurefont\scriptsize] {是} (result);
  \draw[book arrow,knowledge arrow,draw=BookBlue] (matrix.east) -- ++(16,0) |- node[pos=0.22,right,font=\figurefont\scriptsize,text=BookBlue] {每轮读取} (estep.east);
  \draw[book arrow,knowledge arrow,draw=BookAmber,dashed] (stop.west) -- ++(-16,0) |- node[pos=0.18,left,font=\figurefont\scriptsize,text=BookAmber,align=right] {否：携带\\新模型参数} (estep.west);

  \node[anchor=west,font=\figurefont\scriptsize,text=BookMuted,align=left] at (32,-42) {分别检查：\\可识别性\\数值收敛\\答案正确性};
\end{tikzpicture}

  \caption{来源模型参数与潜在标签的交替估计。观测标签矩阵在迭代中保持不变；E步用当前来源模型参数计算类别后验，M步用后验软计数更新类别先验和来源模型参数，再检查停止条件。实线表示样本与计算流，虚线表示更新后的模型参数回流。}
  \label{fig:knowledge-fusion-latent-label-em}
\end{figure}

还记得6.2节那个“\emph{全部报对}还是\emph{全部报反}”的幽灵吗？这里要明确：\textbf{Dawid--Skene 同样甩不掉这种类别置换歧义}。设想我们做一个整体对调：\textbf{同时}交换先验里的元素、交换\emph{所有}混淆矩阵的行、再把潜在类别的名字（山茶$\leftrightarrow$茶梅）一并对调——你会发现\textbf{观测到的标签分布$p(\bm{\Lambda})$纹丝不动}。也就是说，\textbf{无标签数据只能把类别“聚成两组”，却无法给这两组\emph{命名}}。

给类别\textbf{命名}必须依赖外部信息：\textbf{参考锚点}（少量已知标签）或\textbf{“来源总体报对”的方向假设}。而且除了类别置换，\textbf{重叠不足、退化的混淆矩阵}（比如某来源对两类反应完全一样）还会留下\emph{其他}形式的歧义。\textbf{另一类性质不同的问题}是：来源\emph{实际相关}，你却\emph{按独立去把似然相乘}——这\textbf{破坏了生成模型的核心假设}，后果是\textbf{即使拟合得再稳定，也在\emph{重复使用同一份证据}}，得到过度自信的结论。

既然 EM 收敛不保证参数唯一，一个自然的\textbf{更深刻的问题}是：\textbf{在什么条件下，参数原则上就是“可以被唯一恢复”的（除了无法避免的类别置换）}？这就是\emph{可识别性}的\textbf{充分条件}研究。直觉上，要分辨不同类别，它们必须在来源投票中\textbf{留下\emph{彼此可区分}的模式}——如果两个类别在所有来源看来反应都一模一样，数据当然无法把它们分开。下面给出\textbf{针对“一类来源模型”的}充分条件（请记住它是\emph{充分}、且\emph{只适用于这一类模型}，\textbf{不是}所有融合方法的必要条件）。将先验与来源模型参数统记为$\bm{\theta}$，严格可识别要求相同的观测分布只对应同一模型参数：
\[
  p_{\bm{\theta}}(\bm{\Lambda})
  =p_{\bm{\theta}'}(\bm{\Lambda})
  \quad\Longrightarrow\quad
  \bm{\theta}=\bm{\theta}',
\]
若预先把类别名称置换视为等价，结论相应放宽到这类变换。还须每类先验$\pi_k>0$：先验为零的类别不参与观测分布，无论把它的类别票概率改成什么，观察都不变，因此无法恢复该行模型参数。

% 保留矩阵定义的引导段与首个公式，避免前一幅浮动图插入句中。
\Needspace{45mm}
为允许组内共享而保留组间独立，我们把来源分成三个非空组$G_1,G_2,G_3$，要求给定$Y_i$后三组联合类别票独立，组内可以依赖。单列类别票不再是当前比较单位，改为组内所有类别票构成的联合状态；这些状态组成$\Omega_j$，也可包含模型明确保留的弃权。要衡量第$j$组能否区分类别，先定义$\bm{M}_j\in\mathbb{R}^{K\times|\Omega_j|}$：
\[
  [\bm{M}_j]_{k,u}
  =\mathbb{P}(\bm{\lambda}_{i,G_j}=u\mid Y_i=k),
  \qquad
  k\in Y,\ u\in\Omega_j.
\]
矩阵$\bm{M}_j$的\textbf{每一行，是“当真值固定为某一类时，这一组来源联合投票的概率分布”}。于是判读很直观：\textbf{如果有两行（两个类别）几乎相同，就说明这组来源对这两个类别的反应无法区分}，它们在这一组里是“混在一起”的。

为了\textbf{严格地}刻画“任意两类都能被分开”，这里用的不是普通矩阵秩，而是更严格的\term{行Kruskal秩}（row Kruskal rank）$\operatorname{krank}(\bm{M}_j)$：\textbf{任取$r$行都保证线性无关}的最大$r$。为什么不用普通行秩？因为\textbf{普通行秩只要求“\emph{能找到}若干独立行”，而我们需要的是“\emph{无论选哪几行}都独立”}——后者才能保证\emph{每一对}类别都可分，不留下“恰好有两类撞车”的死角。书里这个小例子很说明问题：三行$(1,0)$、$(0,1)$、$(1,0)$的\emph{普通秩}是$2$（确实能找到两行独立），但\emph{第一、第三行完全相同}，所以\emph{Kruskal秩只有$1$}——准确反映出“\textbf{有两类在这组票下根本无法区分}”。

固定三组联合类别票$u_j\in\Omega_j$，$j=1,2,3$后，每个类别贡献其先验乘三个组条件概率，再对类别求和。把观测到这三个状态的联合概率记为$p(u_1,u_2,u_3)$，得到
\[
  p(u_1,u_2,u_3)
  =\sum_{k=1}^{K}\pi_k
    \prod_{j=1}^{3}[\bm{M}_j]_{k,u_j}.
\]
这就是以类别为混合成分的三因子分解。三组状态分别索引联合概率表的三个方向；组内可能出现的类别票状态越多，相应矩阵就有越多列，但是否能区分类别仍由这些行的实际差异决定。下面的秩和门槛为这种分解提供除成分置换外的唯一性条件：
\begin{equation}
  \operatorname{krank}(\bm{M}_1)
  +\operatorname{krank}(\bm{M}_2)
  +\operatorname{krank}(\bm{M}_3)
  \geq 2K+2,
  \label{eq:knowledge-fusion-kruskal-condition}
\end{equation}
这就是本小节的\textbf{核心结论}，先记住它的“保障范围”：在\textbf{每类先验都为正、三组之间给定真值后条件独立、且三组的 Kruskal 秩之和满足式~\eqref{eq:knowledge-fusion-kruskal-condition}}这几个条件下，\textbf{类别先验和三个组条件矩阵在“除类别置换外”的意义上是可识别的}。

这句话有两处需要\textbf{正确理解}，别过度解读：其一，“概率行归一化”这个约束\textbf{消掉了张量分解里常见的“缩放歧义”}（同一成分可以任意乘除常数），但它\textbf{消不掉“类别置换”}，所以\textbf{类别叫什么名字，仍要靠锚点另外定}；其二，这是一个关于\textbf{“真实总体参数能不能被唯一区分”}的结论，\textbf{完全没有承诺“你手上有限的几百个样本能把它估得多准”}——\emph{可识别}和\emph{有限样本估计误差小}是两回事\scite{allman2009identifiability}。

直接核对实际矩阵的秩，得到的是当前模型参数点的结论。若只知道每组联合空间有多大，还可以问这些维度是否容许足够高的秩；每个矩阵的最大行Kruskal秩至多为$\min\{K,|\Omega_j|\}$。当维度满足
\[
  \sum_{j=1}^{3}\min\{K,|\Omega_j|\}\geq2K+2,
\]
在一般位置上，三个条件概率矩阵各自达到允许的最大行Kruskal秩，前面的逐点门槛因而一般成立。“一般位置”排除了代数退化构成的零测集合，不等于每个模型参数点都成立，也不排除有限样本不稳定。实际模型有额外模型参数约束时，必须检查可产生的矩阵；仅数类别票空间大小可能高估区分能力。

一个你可能正想问的问题：\textbf{把同一个来源复制成好几份（列变多），到底会不会损害可识别性}？先给结论：\textbf{复制不会\emph{新增}独立票，但是否会\emph{降低}已有的区分能力，不能凭直觉，要看实际的行分布}。下面用具体矩阵验证这一点——这也回应了“多一列到底有没有用”。设一个二分类来源对称正确率为$0.8$，组内第二个来源逐项复制它；以$0,1$编码两类，联合类别票依次为$00,01,10,11$。不一致的两列概率为零，该组条件矩阵为
\[
  \bm{M}
  =\begin{pmatrix}
    0.8 & 0 & 0 & 0.2\\
    0.2 & 0 & 0 & 0.8
  \end{pmatrix}.
\]
由于两来源逐项复制，\emph{不一致}的两列（$01$和$10$）概率都是$0$，矩阵只在$00$、$11$两列上有值。取第1、第4列算一个二阶子式：$0.8^2-0.2^2=0.6\neq0$，所以\textbf{两行仍然线性独立、Kruskal秩仍为$2$}。这个结论可能出乎意料：若三个这样的组给定真值后独立，秩和为$6=2K+2$，满足二分类门槛；复制没有增秩，也没有降低本例已有的秩。把正确率换成$0.5$，两行相同，秩才降到$1$；与类别无关的恒定类别票同样只有秩$1$。组间依赖则破坏三因子分解前提，不是一个单纯的秩问题。来源数量、组空间大小与EM收敛都不能替代这些检查。

到目前为止，我们的概率乘积\textbf{只纳入了“实际投了票”的来源}。但有两种常见情况它还没覆盖：其一，我们\textbf{想把“成功执行后却\emph{不给类别}（弃权）的\emph{倾向}”本身也建模}（弃权可能携带信息）；其二，我们\textbf{已经知道某些来源的票\emph{共享影响}}，想把这种\emph{关系}直接写进分布。

这就需要一个更一般的框架。\term{标签模型}（label model）是\textbf{“描述弱来源\emph{如何产生票}”的一\emph{家族}生成模型}的统称——请特别注意：\textbf{它\emph{不是} Dawid--Skene 的另一个名字}，Dawid--Skene（每来源一张完整混淆矩阵）只是这个家族里\emph{一种}参数化。为统一表达“弃权”，约定标注函数返回$\lambda_{i,s}\in Y\cup\{\bot\}$，其中$\bot$是\textbf{模型接口里“不投票”的编码}。

这个编码与原始样本标注中的“主动弃权”动作需要分开。对已成功执行且对象落在共同任务适用范围内的来源，规则未触发和主动弃权都可以按明确的映射进入$\bot$，原始样本标注仍分别保留触发结果、动作、适用范围和原因。Snorkel把标注函数未匹配时返回的空值也称为abstain，指的正是没有给出类别票，而非本书样本标注层中单独区分的主动弃权动作\scite{ratner2017}
未请求、执行失败、范围外或适用范围尚未对齐的样本标注不按这一映射并入$\bot$，须另按观测设计或生成机制处理。若需要分别估计未触发与主动弃权的发生规律，还应扩展模型状态，不能在合并编码后假定两种机制相同。

在确定类别票状态以后，用$R$个函数$\phi_r$描述类别与类别票的特定关系，再以模型参数$\theta_r$控制其强度。一个通用对数线性形式为
\begin{equation}
  p_{\bm{\theta}}(Y_i,\bm{\lambda}_i)
  =\frac{1}{Z_{\bm{\theta}}}
  \exp\left(
    \sum_{r=1}^{R}\theta_r
    \phi_r(Y_i,\bm{\lambda}_i)
  \right).
  \label{eq:knowledge-fusion-label-model}
\end{equation}
先看懂这个通用形式的两个部件：指数函数$\exp(\cdot)$给\textbf{每一种“真值+票”的组合}$(Y_i,\bm{\lambda}_i)$\textbf{配一个正权重}——某个因子$\theta_r$为正、且该因子$\phi_r$在这个组合上被激活，就会\textbf{抬高这个组合的相对权重}；而分母$Z_{\bm{\theta}}$是把这些权重在\textbf{所有可能的类别、所有可能的出票组合}上求的总和，\textbf{除以它（归一化）之后，整堆权重才成为一个合法的概率分布}（这就是“配分函数”的作用）。

举个最常用的因子：令$\phi_r=\mathbb{I}\{\lambda_{i,s}=Y_i\}$（这张票\emph{报对了}吗），那么对应的$\theta_r>0$就表示“这是个\emph{倾向报对}的来源”，它会\textbf{系统性地提高“票与真值一致”这类配置的权重}。换不同的$\phi_r$，就能灵活表达各种先验知识。这个因子框架的\textbf{表达力}正在于$\phi_r$可以五花八门：除了“是否报对”，还能放入\textbf{类别先验}（某类更常见）、\textbf{偏好}（来源偏爱某类）、\textbf{覆盖/弃权}（它在什么情况下不给票）、以及最关键的\textbf{“来源对”的相关行为}（两条规则总是一起错）。于是模型可以\textbf{由简到繁}：最简版本只保留“先验 + 各来源条件独立的准确性因子”（这时近似回到 Dawid--Skene 的独立假设）；\textbf{需要建模依赖时，再显式加入描述成对/成组关系的因子}。程序化监督（Data Programming）将规则、实体关系库和模型写成可弃权的标注函数，再利用重叠与冲突拟合生成模型参数\scite{ratner2016data-programming}

用式~\eqref{eq:knowledge-fusion-label-model}无标签拟合时，仍要对真值求和，使用$p_{\bm{\theta}}(\bm{\lambda}_i)=\sum_kp_{\bm{\theta}}(Y_i=k,\bm{\lambda}_i)$，然后计算类别后验。标准用法假设固定模型参数下样本独立、来源行为在建模范围内稳定；它只读取类别票矩阵，不自动从原始$\bm{x}_i$识别困难照片。质量随样本或时间变化时，需要分组或加入条件因子。

% 段末文献需要与本段一同检查可用高度。
\Needspace{70mm}
现在可以\textbf{精确地}说清 Dawid--Skene 和程序化弱监督里常用标签模型的\textbf{区别}了——它们\textbf{不是新旧名称的替换，而是“对同一个生成过程的两种\emph{参数化}”}：
\begin{itemize}\item \textbf{Dawid--Skene}：通常给\emph{每个来源一整张完整混淆矩阵}（$K\times K$个格子），\textbf{刻画最细、参数最多}；\item \textbf{程序化弱监督的标签模型}：通常\emph{只用较少的“准确性因子”}（比如一个“总体靠不靠谱”的标量），并\textbf{显式地描述弃权、来源间依赖}，\textbf{参数更省}。\end{itemize}
两者\textbf{都在推断潜在类别}，但\textbf{改变的是“票的生成模型长什么样”以及“采用了哪些生成假设”}。理解到这一层，你就不会误以为“标签模型”只是换了个时髦说法。Snorkel将这类生成模型得到的概率标签用于训练下游判别模型，后者再学习从样本特征预测类别\scite{ratner2017}

这又回到那条贯穿全章的\textbf{“表达力 vs 可估计性”权衡}：\textbf{保留完整混淆矩阵}，能\textbf{区分每一种错误方向}（茶梅误成山茶 vs 误成月季分得清），但代价是\textbf{需要多得多的“分类别、有重叠”证据}才能把每个格子估准；\textbf{改用较粗的准确性因子}，参数少了、好估了，代价却是\textbf{被迫加入“误差对称”“来源共享某结构”等额外限制}——而这些限制一旦不成立，结论就会偏。所以\textbf{没有绝对更好的那种}：选哪种，要看\textbf{当前的关键分歧是否真的需要区分这些方向}，以及\textbf{现有样本量能不能撑得起}这么细的估计。

增加因子也会增加无法区分模型参数的风险。除类别置换外，对数线性表示可能冗余，使不同模型参数给同一分布；条件独立版本需要足够重叠与有效独立来源，依赖版本还需因子结构正确且证据能区分其模型参数。Data Programming的恢复结论只适用于论文规定的模型和条件，不意味着任意规则都能恢复真值。高度相关、类别方向未知或模型参数超过有效证据时，应补锚点、简化或合并来源，必要时保留未决；后验收敛不是验证证据。

\subsubsection{结构约束下的联合推断}
\label{sec:knowledge-fusion-structured-inference}

前面的融合都是\textbf{对“一个对象、一个标签”独立做判断}。但有一大类任务，\textbf{一张输入里的多个标签之间存在硬性的合法关系}，这时\emph{每个位置各自挑最高分}，拼起来可能是个\textbf{根本不合法}的东西。最经典的是命名实体的\term{BIO标注}（beginning--inside--outside）：$B$表示实体\emph{开始}、$I$表示\emph{同类实体内部}、$O$表示\emph{实体外部}。这里有条不可违反的语法：\textbf{$I$不能凭空出现——它前面必须有一个对应的$B$}。如果逐位置独立取最大，很容易得到“$O$后面直接跟$I$”这种非法序列。除了 BIO，\emph{类别层级}（子标签不能违背父标签）、\emph{互斥规则}（同一位置不能既是A类又是B类）也都在限制多个答案的\textbf{组合}。

\term{联合推断}（joint inference）正是为此而生：\textbf{把这一组\emph{相互牵连}的标签放在一起、作为一个整体来选择}，而\textbf{不是“各自 independently 选完、再一厢情愿地假定拼起来也合法”}。

为避免概念串台，这里要和\textbf{第3章的“规则获取”划清界限}——两者\textbf{都会用到业务规则，但处理的对象和时机完全不同}：
\begin{itemize}\item \emph{规则获取}（第3章）：规则的作用是\textbf{从条件\emph{生成}候选标签或约束}，是在“\emph{无中生有地造证据}”；\item \emph{本节的联合推断}：\textbf{多源融合的局部得分已经摆在眼前}，规则的作用是\textbf{在众多候选\emph{组合}里挑出一个整体合法的赋值}。\end{itemize}
同一条业务规则（比如“$I$必须接在$B$后”）完全可以在\emph{两处}都出现，但它\textbf{前者在评价“生成的候选好不好”，后者在评价“整体组合可不可行、得分高不高”}。\textbf{不能因为“都用了约束”就把这两步混为一谈}。

联合推断的核心是在两股力量间\textbf{取舍}：一股是每个局部位置\textbf{自己的支持分}（都想选自己的最高分），另一股是\textbf{结构要求}（不允许/不鼓励某些组合）。为此首先要\textbf{区分两种强度截然不同的约束}——\textbf{哪些组合是“\emph{根本不允许}”（硬约束，违反即作废），哪些只是“\emph{最好别这样}”（软约束，违反要扣分但尚可）}。设待选标签为$z_i\in Y$、局部得分为$a_i(z_i)$；任务硬约束定义可行向量集合$\bm{Z}_{\mathrm{feas}}$，软约束$r$的违反程度为$\rho_r(\bm{z})\geq0$，罚值权重$\eta_r\geq0$须使代价与得分可比较。在允许集合中最大化支持总分并减去违反代价，得到
\begin{equation}
  \widehat{\bm{z}}
  \in\argmax_{\bm{z}\in\bm{Z}_{\mathrm{feas}}}
  \left\{
    \sum_i a_i(z_i)
    -\sum_r\eta_r\rho_r(\bm{z})
  \right\}.
  \label{eq:knowledge-fusion-structured-objective}
\end{equation}
这个目标函数的读法：\textbf{硬约束}表现为\emph{取值范围}$\bm{z}\in\bm{Z}_{\mathrm{feas}}$——\textbf{集合之外的组合直接不可选}（连被比较的资格都没有）；\textbf{软约束}表现为\emph{目标里减去的惩罚项}$\eta_r\rho_r(\bm{z})$——它\textbf{不删候选，只在候选之间“谁违反得多谁扣分多”}。惩罚项前面是\textbf{减号}，纯粹因为这里在做\emph{最大化}（违反得越狠，总目标越低）。

几个使用细节：如果局部给的是\emph{概率}$q_{i,k}$而非任意分数，可以取$\log q_{i,k}$作为得分（概率越高、对数越大），此时\emph{零概率对应$-\infty$}，天然等价于一条硬约束。但要\textbf{防止过度解读}：\textbf{如果你并没有“联合概率能分解成局部乘积”的依据，那么这个式子就\emph{只是一个评分目标}，优化它\textbf{不等于}自动求得了“联合后验}”；同理，\textbf{未经校准的加权分数，只表达“当前这套评分体系的偏好”}，不能当成真实概率。

我们用三个词元枚举一次联合选择。共同标签为$O,B\text{-ORG},I\text{-ORG}$，后两者表示组织实体的开始与内部；第三个位置只允许$O$。表内$-\infty$排除对应位置标签，其余数值是人为设定的局部支持分：
\[
  \begin{array}{c|ccc}
    & O & B\text{-ORG} & I\text{-ORG} \\
    \hline
    1 & 1.2 & 0.8 & -\infty \\
    2 & 0.5 & 0.6 & 1.4 \\
    3 & 0.7 & -\infty & -\infty
  \end{array}.
\]
\textbf{先看“不懂结构”会怎样}：逐位置各自取最大，会选出$(O,I\text{-ORG},O)$，总分高达$3.3=1.2+1.4+0.7$。但它\textbf{违反了“$I$不能接在$O$后面”这条硬规则}（位置2是$I$，位置1却是$O$，凭空冒出一个实体内部），是个\textbf{非法序列}——这就是“局部全最优、整体不合法”的活例子。有限分数的合法组合只有五个：$(O,O,O)$为$2.4$，$(O,B\text{-ORG},O)$为$2.5$，$(B\text{-ORG},O,O)$为$2.0$，$(B\text{-ORG},B\text{-ORG},O)$为$2.1$，$(B\text{-ORG},I\text{-ORG},O)$为$2.9$。在五个合法组合里，总分最高的是$(B\text{-ORG},I\text{-ORG},O)$，得$2.9$。请\textbf{注意硬约束到底“做了什么动作”}：它\textbf{不是简单地把那个冲突的$I$标签删掉}，而是\textbf{回过头去修改了\emph{第一个位置}}（把位置1从$O$改成$B\text{-ORG}$），\textbf{好让位置2那个很强的$I$证据（$1.4$）能合法地保留下来}。换言之，\textbf{联合推断为了保住一个强局部证据，可能主动调整别的位置去“配合”它}，这是逐点独立决策做不到的。

保持局部分数不变，只把这条转移规则改成一次违反记$1$的软代价，非法序列便以$3.3-\eta$参加比较，其余五个候选罚值为零。$\eta=0.2$时它得$3.1$，仍胜过合法最优$2.9$；$\eta=0.5$时只得$2.8$，转而选择$(B\text{-ORG},I\text{-ORG},O)$。对比两次结果能\textbf{把“硬/软约束的本质差别”看得清清楚楚}：在\emph{完全相同的局部分数}下，\textbf{硬约束改变的是“\emph{可选集合}”}（非法序列根本不许入场，于是选$2.9$的合法最优）；\textbf{软约束改变的是“\emph{入场后的得分}”}（非法序列允许参赛，但带罚分参赛，罚轻了它仍能赢、罚重了合法组合才反超）。\textbf{同一份局部证据，仅因约束软硬不同，最终决策就可能相反}。互斥分类也可以保留得分损失较小的一项；若互斥规则可能不适用，则需明确是否改为允许违反的软约束。

最后给联合推断\textbf{去去魅}，它不是“加了约束就一定对”。首先，\textbf{选出合法组合的\emph{前提}是“约束本身正确且适用于当前样本”}：约束可能\emph{过时}（业务规则早改了）、可能\emph{互相冲突}（两条规则不容）；当\textbf{求解失败}时，原因也不止一种——也许是\emph{确实没有可行赋值}、也许是\emph{你给当前样本错套了不适用的规则}、又或许只是\emph{求解器在资源限额内没来得及找到解}，这三者要能区分。

工程伦理上：\textbf{删除/消解冲突时要保留原始来源信息}（不能把证据抹掉），\textbf{放松硬约束为软约束时要记录违反程度}，给\textbf{近似解则必须说清“尚未保证的是可行性还是最优性”}。求解成本也别想当然：\emph{链式序列}（如BIO）有高效的\textbf{动态规划}（维特比）；但\textbf{更一般的离散耦合}可能要动用\textbf{整数规划、消息传递、或近似搜索}，规模一大就可能算不动——\textbf{不能因为“模型写得出来”就推断“求解成本一定扛得住”}。

\subsection{基于模型预测的融合}
\label{sec:ket-integration-prediction}

前面的投票、加权、概率推断，融合对象主要是\emph{已经产出的标签或票}。现在视角换到\textbf{模型的“预测输出”本身}：我们\emph{保留选定的若干来源模型、让它们在运行时持续地就新照片做预测}，重点考察三件\emph{标签层融合没有充分暴露}的事——\textbf{它们的错误是\emph{互补}还是\emph{重合}、由此带来的收益、以及同时运行多个模型的\emph{成本}}。

开始前先重申两条前置纪律：其一，\textbf{各来源仍必须使用共同的类别语义}（回到6.3的对齐），把“山茶”和一种内部码当同类相加毫无意义；其二，\textbf{只检查概率向量“和为一”，并不能确认它已校准}——形式合法不等于频率可信。至于这些成员模型\emph{最初如何产生、如何训练}（Bagging、Boosting 等），可参考\BookChapterRef{chap:ensemble-learning}{卷4《模型：经典模型》的“集成学习”章}；\textbf{本节只聚焦“\emph{已经在手}的这些来源，怎样在预测时共同工作”}。

这一小节先区分两种“\emph{组合规则随不随输入变}”的大类型，判据要抓准：
\begin{itemize}\item \term{固定组合}：在\textbf{声明好的任务和输入范围内}，\textbf{用哪些来源、各自权重多少，事先固定、对所有样本一视同仁}；\item \term{条件组合}：\textbf{依据“当前是什么任务、或当前这张输入长什么样”，动态地改变来源、权重或连接}（模糊照片多信模型二，清晰照片多信模型一）。\end{itemize}
这里有一个\textbf{容易误判}的灰色地带，判据是“\textbf{运行时的规则\emph{是否依赖}任务/输入条件}”，而\textbf{不是“这套规则有没有经过训练”}：哪怕权重是\emph{学出来}的，只要\textbf{学好后冻结、对所有样本都用同一组}，它仍属\textbf{固定组合}；反过来，哪怕路由规则是\emph{人工拍脑袋指定}的，只要它\textbf{按输入选择来源}，它就是\textbf{条件组合}。

\subsubsection{基于固定规则的预测集成}
\label{sec:ket-integration-fixed-prediction}

最直接的固定组合就是前面讲过的\textbf{投票}（见第~\ref{sec:knowledge-fusion-voting}节）；如果\textbf{保留的是模型的类别概率}，则直接\textbf{复用式~\eqref{eq:knowledge-fusion-weighted-probability}的线性混合}（加权平均概率）。下面把它用“预测集成”的语言再写一遍，并把记号与标签模型的类别数$K$区分开。这里固定非零权重来源均可用，将其个数记为$K_{\mathrm{s}}$，与前面标签模型的类别数$K$分开；以$\alpha_k$表示归一后、在声明范围内不随预测输入改变的权重。以$p_k(y\mid\bm{x})$表示来源$k$给出的类别概率，于是$\bar p(y\mid\bm{x})=\sum_{k=1}^{K_{\mathrm{s}}}\alpha_kp_k(y\mid\bm{x})$，其中$\alpha_k\geq0$且$\sum_k\alpha_k=1$。关于这个加权平均，先记住一个\textbf{最容易被高估}的事实：\textbf{归一化（除以权重和）只保证“输出是个合法分布”，它\emph{本身不保证集成更准}}。如果你打算\textbf{用数据去\emph{选}这些权重}，就必须遵守\textbf{样本隔离}：用一份与“来源训练”\emph{隔离}的验证材料选权重，再用\textbf{完全没参与选择}的样本来报最终成绩，否则就是在验证集上过拟合。

接下来这个算例\textbf{专门隔离出一个变量}：\textbf{来源数量不变、各自的正确率也不变，只改变“它们的错误\emph{落在不在同一批照片上}”}。你会看到，\textbf{仅仅这一个差别，就足以决定集成到底有没有收益}。

\begin{example}[互补错误与共同错误]
  \label{ex:ket-integration-weighted-prediction}
  将花卉任务简化为月季与非月季二分类。正类$y=1$表示月季，正类概率达到$0.5$时判为月季，平局也归正类。四张教学照片的真实标签与来源概率为
  \[
    (y_1,y_2,y_3,y_4)=(1,0,1,0),\qquad
    \begin{aligned}
      (p_1(\bm{x}_i))_{i=1}^{4}&=(0.9,0.4,0.6,0.6),\\
      (p_2(\bm{x}_i))_{i=1}^{4}&=(0.4,0.1,0.9,0.3).
    \end{aligned}
  \]
  这里$p_k(\bm{x}_i)$简写$p_k(y=1\mid\bm{x}_i)$。模型一错第四张，模型二错第一张，两者正确率都是$3/4$。给两个来源权重$0.6$与$0.4$，组合概率为
  \[
    (\bar p(\bm{x}_i))_{i=1}^{4}
    =(0.70,0.28,0.72,0.48).
  \]
  第四张由$0.6\times0.6+0.4\times0.3=0.48$落到阈值以下，四张照片因而全部分类正确。现在保持标签、模型一和权重不变，只将模型二概率换成$(0.8,0.2,0.7,0.7)$。模型二正确率仍是$3/4$，但它与模型一都错第四张；组合概率变为$(0.86,0.32,0.64,0.64)$，第四张仍错。成员数量和各自正确率相同，错误发生在哪些预测输入上却不同，组合收益也随之改变。
\end{example}

上面的分类算例已经点出了\textbf{关键：收益来自“错误\emph{不}重叠”}。但“按$0.5$阈值判对错”是个\emph{非黑即白}的操作，\textbf{不方便从数学上拆解“平均到底消掉了哪一部分误差”}。为了\textbf{看清机制}，我们\textbf{暂时把问题换成实值回归 + 平方损失}——在这个更光滑的设定下，误差可以干净地分解为“偏差”和“方差/协方差”，结论再反过来帮助理解分类。设$y=f_\star(\bm{x})+\varepsilon$，其中$f_\star(\bm{x})=\mathbb{E}[y\mid\bm{x}]$，噪声满足$\mathbb{E}[\varepsilon\mid\bm{x}]=0$。固定预测输入$\bm{x}$，只让训练样本和训练随机性变化，假设所需二阶矩有限，并假设测试噪声与学习算法在该预测输入下条件独立。

将第$k$个预测器相对回归函数的误差记为$e_k=\hat f_k(\bm{x})-f_\star(\bm{x})$，误差均值为$b_k=\mathbb{E}[e_k]$，协方差为$\Sigma_{k,j}=\operatorname{Cov}(e_k,e_j)$。务必注意这些期望\textbf{是对“什么”求的}：它描述的是\textbf{“重复地用不同训练集/随机性去训练这些模型”所产生的波动}（把训练过程当成随机的），\textbf{不是}对你眼前\emph{这几个已经训好的固定预测值}求平均。明确这一点，后面的“方差、协方差”才有意义。在权重和为一时，集成在点$\bm{x}$的误差就是各来源误差的加权和$\bar e=\sum_k\alpha_ke_k$。把每个误差拆成“均值（偏差）+ 中心化波动”，可以把集成的\emph{期望平方误差}展开为：
\begin{equation}
  \begin{aligned}
    \mathbb{E}[\bar e^2]
    &=\biggl(\sum_k\alpha_k b_k\biggr)^2+
      \mathbb{E}\biggl[\biggl(\sum_k\alpha_k(e_k-b_k)\biggr)^2\biggr]\\
    &=\biggl(\sum_k\alpha_k b_k\biggr)^2+
      \sum_k\alpha_k^2\Sigma_{k,k}
      +2\sum_{k<j}\alpha_k\alpha_j\Sigma_{k,j}.
  \end{aligned}
  \label{eq:ket-integration-bias-covariance}
\end{equation}
这个展开是理解集成的\textbf{钥匙}，逐项读：\emph{第一行}是偏差的平方——把各来源误差的均值加权求和后整体平方；展开平方时本会出现大量“偏差$\times$中心化波动”的\emph{交叉项}，但因为\textbf{中心化误差的均值恰好为零}，这些交叉项\textbf{期望都消失}了。\emph{第二行}是纯粹的\textbf{方差/协方差}：每个来源自己一项$\alpha_k^2\Sigma_{k,k}$，\textbf{每\emph{一对}不同来源再贡献一个协方差项}$2\alpha_k\alpha_j\Sigma_{k,j}$。

这个公式最有价值的地方，正是\textbf{它同时保留了两样东西}：一是组合的\textbf{偏差}（大家共同的、系统性的偏，平均掉不掉得看下面），二是\textbf{来源错误“一起变化”的程度}（协方差）——后者就是“错误重不重叠”的数学化身。两个使用上的补充：第一，如果你\textbf{计分的对象是“带噪声的观测标签”$y$}（而非干净的$f_\star$），期望误差里还要\textbf{额外加上一项不可消除的噪声方差}$\operatorname{Var}(\varepsilon\mid\bm{x})$；好在测试噪声\emph{零均值}且与学习算法\emph{条件独立}，所以噪声与模型误差的\emph{交叉项}期望为零，只是凭空多出这一项“地板”。第二，以上都在\textbf{固定一个点}$\bm{x}$上分析；若要评价\textbf{整个输入范围}的总体误差，还需\textbf{再对输入分布$\bm{x}$求一次期望}。

为了\textbf{最直观地看出“加模型”到底有没有用}，我们代入一个干净的\textbf{对称设定}：设每个来源都\textbf{无偏}（$b_k=0$，没有共同的系统性偏差）、单个方差都为$\sigma^2>0$、且\textbf{任意两个来源误差的相关系数都是同一个}$\rho$（错误“同涨同跌”的程度），然后做\textbf{等权平均}。将这些值代入协方差展开，得到
\begin{equation}
  \operatorname{Var}(\bar e)
  =\frac{K_{\mathrm{s}}\sigma^2+K_{\mathrm{s}}(K_{\mathrm{s}}-1)\rho\sigma^2}{K_{\mathrm{s}}^2}
  =\sigma^2\left(\rho+\frac{1-\rho}{K_{\mathrm{s}}}\right).
  \label{eq:ket-integration-equal-correlation}
\end{equation}
代入后的式子能读出\textbf{集成学习最核心的两条洞见}，请记住这两个极端：
\begin{itemize}\item \textbf{错误完全相关（$\rho=1$）}：集成方差$=\sigma^2\bigl(1+(1-1)/K\bigr)=\sigma^2$，\textbf{平均\emph{一点都没}降低方差}。这对应“\emph{所有模型犯一模一样的错}”——典型如它们都用同一套坏特征、或根本是同一模型的复制品，加再多也没用。\item \textbf{错误互不相关（$\rho=0$）}：集成方差$=\sigma^2/K_{\mathrm{s}}$，\textbf{方差随来源数线性下降}。这正是集成\emph{理想}的工作状态：各模型的\emph{随机}错误方向不同，一平均就互相抵消。\end{itemize}
中间值$\rho$给出部分收益（下降到$\sigma^2\bigl(\rho+(1-\rho)/K\bigr)$）。还有两点不能漏：其一，这个“等相关”模型在数学上要\textbf{协方差矩阵合法}，要求$-1/(K_{\mathrm{s}}-1)\leq\rho\leq1$（不可能让很多误差两两强负相关）；其二，也是最关键的提醒——\textbf{平均只能消“方差”，消不掉“共同的偏差”}：\textbf{如果所有来源都带着\emph{同一个}非零偏差（都被同一个背景捷径骗），等权平均后那个偏差\emph{原封不动}}。最后，\textbf{回到分类任务要打个折扣}：分类对错还取决于“概率是否跨过决策阈值”，所以\textbf{不能把回归里这条干净的方差下降，直接当成“准确率必然提高”的证明}——它解释了\emph{为什么}常常有益，但阈值处的非线性需另算。

在分类里做平均，还有一个\textbf{极易踩的坑：在“哪一层”汇总，结果是不一样的}。来源给的若是\emph{还没归一化的分数（logits）}$\bm{a}_k$，你既可以“先把分数平均、再统一 softmax”，也可以“各自 softmax 成概率、再平均概率”。由于 softmax 是\emph{非线性}的，这两者\textbf{一般不相等}：
$$\operatorname{softmax}\Bigl(\sum_k\alpha_k\bm{a}_k\Bigr)\ \ne\ \sum_k\alpha_k\operatorname{softmax}(\bm{a}_k).$$这两种做法各有脾气：在当前非负归一权重下，\textbf{“先平均分数再 softmax”等价于“对各来源的\emph{严格正}类别概率做\emph{加权几何平均}再归一化”}（几何平均对“某来源给某类极低概率”更敏感，一个近零就能压垮该类）；而\textbf{直接平均概率做的是\emph{算术平均}}。前者\textbf{更容易受分数\emph{尺度}影响}（logits 整体放大/缩小会改变 softmax 的尖锐度）；再加上前面的\textbf{投票}则更极端——它\textbf{连概率的幅度信息整个丢弃}，只留类别名。也因此再次强调：\textbf{未经校准的“高置信度”，不能直接当成“给这个来源提高权重”的理由}（呼应6.4.2那个自信却错的算例）。

最后算一笔\textbf{容易被忽略的“成本账”}：固定集成\textbf{通常意味着要\emph{执行所有}非零权重的来源}，\textbf{省不掉}。若来源是彼此独立的模型，它们的\textbf{权重存储和运算量大致是“相加”}的；想减少\emph{等待时间}可以\textbf{并行}执行，但并行\textbf{需要多份设备}，还受\textbf{通信开销和“最慢那个来源”（木桶效应）}限制，而\textbf{串行}则让延迟\textbf{逐段累加}。若来源是\textbf{共享同一底座的适配器}，可能\textbf{复用一部分前向计算}，但具体省多少\textbf{必须按真实计算路径核算}，不能想当然。

收尾也是两个\textbf{必须分别验证、不能互相顶替}的结论：\textbf{“这样做是否\emph{节约了资源}”}与\textbf{“这样做是否\emph{降低了任务错误}”}是\emph{两个独立的问题}——完全可能更便宜但没更准、或更准但贵得多，要\textbf{分开测量、分开报告}。

\subsubsection{基于条件规则的预测集成}
\label{sec:ket-integration-conditional-prediction}

固定权重有个\textbf{盲区}：它对\emph{所有}照片都用同一套配比。可现实中常常是“\textbf{模型一在\emph{清晰}照片上可靠、模型二在\emph{模糊}照片上可靠}”——这种\emph{各有所长、且长处随输入而变}的情形，一组\emph{全局}权重根本表达不了。\term{条件规则}（conditional rule）的思路是：\textbf{引入一个“选择器”，让它读取\emph{预测时拿得到}的信息}（比如\textbf{图像质量、任务说明、或其他无需知道答案就能观测的证据}$\bm{q}$），\textbf{据此\emph{当场}分配各来源的权重}。注意这里把\textbf{任务索引}$u$和\textbf{来源索引}$k$\textbf{分开}，是因为“在做哪个任务”和“听哪个模型”是两个维度。用$u$索引目标任务$\mathcal T_u$，与来源索引$k$分开；$\bm{q}(\bm{x},u)$表示选择器可观察的图像质量、任务标识与证据。同一任务内，只有预测输出含义相容且可执行的来源进入候选，组合分布写为
\begin{equation}
  \bar p(y\mid\bm{x},u)
  =\sum_k\alpha_k(\bm{q}(\bm{x},u))p_k(y\mid\bm{x},u).
  \label{eq:ket-integration-conditional-average}
\end{equation}
这个式子要当\textbf{概率混合}来解释，仍需\textbf{权重非负、对每个输入和为一}。条件组合还能分成两种极端形态：\textbf{只让权重集中到一个来源（硬选择/路由）}，或\textbf{让多个来源按权重共同贡献（软组合）}。

选择器（路由器）本身也得\textbf{学}，学的时候有两条路和一条铁律：可以在\textbf{与来源训练隔离}的有标签材料上，\textbf{分别统计各来源在不同质量区域的损失}来定权重；也可以\textbf{直接训练一个“门控函数”}去最小化最终组合损失——所谓\term{门控函数}（gating function），就是\textbf{“从看得见的条件$\bm{q}$到来源权重”的那个映射}。\textbf{铁律是样本覆盖与隔离}：训练门控的材料\textbf{必须覆盖你将来打算服务的各种条件}，而且\textbf{绝不能拿“来源在\emph{它自己训练样本}上的过拟合高分”去证明“换个条件它仍可靠”}。

条件路由\textbf{不是稳赚}，它把“融合错误”换成了“路由错误”的可能，用月季二分类（阈值仍$0.5$）看这个\emph{双刃}效果。设一张\emph{清晰的月季}、一张\emph{模糊的非月季}：擅长清晰的模型一给$0.9,0.95$，擅长模糊的模型二给$0.2,0.2$。
\begin{itemize}\item \textbf{等权平均}：得$0.55,0.575$，第二张（模糊非月季）因$0.575>0.5$\textbf{误判成月季}。\item \textbf{路由正确}（清晰度指标分得清两张）：清晰的走模型一、模糊的走模型二，得$0.9,0.2$，\textbf{两张都对}。\item \textbf{路由颠倒}（质量分组恰好搞反）：清晰的错送给模型二、模糊的错送给模型一，得$0.2,0.95$，\textbf{两张\emph{全错}}——比等权平均还糟。\end{itemize}
最后一句是路由的\textbf{根本约束}：选择器\textbf{能依据的只有“预测时\emph{看得见}的质量”}，\textbf{不可能依据“预测时\emph{尚不知道}的真实答案}”。所以路由质量的上限，取决于\textbf{这些可见特征与“哪个来源更靠谱”之间到底有多强、多稳定的相关}。

条件“\textbf{在哪一层}发生”会\textbf{显著改变调用成本}，要区分两种粒度：
\begin{itemize}\item \textbf{任务层条件}：根据任务说明或少量示例，\textbf{先定好一组权重$\bm{\alpha}_u$，然后整个任务期间冻结}。选择器\textbf{每个任务只跑一次}，摊薄后很便宜。\item \textbf{输入层条件（逐样本路由）}：\textbf{每一张照片都要跑一次选择器}，路由开销按样本数累积。\end{itemize}
还有一个\textbf{决定“路由能不能真正省钱”}的关键细节：\textbf{你能否\emph{在运行来源\emph{之前}}就做出选择}？如果清晰度可以\textbf{直接从原图用很便宜的方式估出}，那么\textbf{只运行被选中的那个来源即可}，真省了算力；但如果你的路由依据是“\textbf{先看看各个模型的置信度谁高}”，那就\textbf{必须先把\emph{所有}模型都跑一遍}才能比——此时\textbf{来源的计算已经全部发生}，\textbf{事后只保留一个答案，并不能撤销这些已花掉的费用}。

一个\textbf{负责任的条件系统必须安排“出路”}：对那些\textbf{超出已验证条件、或选择器证据不足}的输入，应\textbf{预先设计弃权、转人工复核、或走备用流程}，而不是硬着头皮乱路由。评价这类带弃权的系统时，\textbf{必须同时报告两个数}：\textbf{被接纳的输入占全部输入的\emph{比例}（覆盖率）}，以及\textbf{被接纳部分的\emph{错误率}}。只报错误率会骗人——一个\textbf{只挑极少数最容易样本作答}的系统，错误率可以极低，却\textbf{完全不代表系统整体有用}（大量样本被它回避了）。

还要\textbf{警惕路由关系随环境漂移}：\textbf{地区、季节、拍摄设备}一变，“图像质量”和“哪个来源更准”之间的关系\textbf{可能就变了}，需要\textbf{重新检验}路由规则，不能一次调好永久使用。最后划清能力边界：\textbf{这里组合的是来源\emph{已经给出的答案}}——\textbf{对答案加权，并不等于你能去组合各来源\emph{内部}的推理步骤}（那是下一节“模块融合”才能做的事）。

\section{基于计算模块的融合}
\label{sec:ket-integration-modules}

到目前为止，无论是投票、加权还是概率集成，融合的都是\textbf{“同一问题的\emph{同义答案}”}（都在回答“这是什么花”，只是口径不同）。但有一类目标\textbf{光靠汇总最终答案根本做不到}：比如不问“这是什么花”，而问“\textbf{图中\emph{哪一朵}花\emph{同时}是红色\emph{且}钟形}”。这个问题需要的是\textbf{中间的属性、以及属性与对象的绑定关系}，最终的“物种概率”里\textbf{已经把这些信息压没了}。

于是融合要\textbf{再往模型内部走一层}：从“汇总答案”转向“\textbf{连接计算模块}”——让一个模块的\textbf{中间表示}成为下一个模块的输入、继续参与计算。这里立刻冒出一个\textbf{多对象难题}：对\emph{单朵花}的图，传一个属性向量就行；但对\emph{多朵花}的图，必须\textbf{把“颜色、形态、位置”绑定到\emph{同一朵}花上}，绝不能\textbf{取A花的颜色、配B花的形态}（张冠李戴）。图~\ref{fig:ket-integration-interfaces}正好对照了这两种接口：左边是汇总分布，右边是中间表示继续流动。图~\ref{fig:ket-integration-interfaces}对照了这两种接口：左侧汇总分布，右侧让中间表示继续参与计算。

\begin{figure*}[htbp]
  \centering
  % Interface comparison, two fixed tasks; schematic, no empirical data.
% Canvas: 164 mm x 69 mm. Dashed arrows carry control, not predictions.
\begin{tikzpicture}[
  x=1mm,y=1mm,font=\figurefont\footnotesize,text=BookInk,
  box/.style={rectangle,draw=FigureModel,fill=FigureModel!7,
    line width=0.8pt,minimum width=22mm,minimum height=9mm,
    align=center,inner sep=1mm},
  input/.style={box,draw=FigureInput,fill=FigureInput!7,
    minimum width=15mm},
  result/.style={box,draw=FigureOutput,fill=FigureOutput!7,
    minimum width=20mm},
  control/.style={box,draw=BookAmber,fill=BookAmber!7},
  flow/.style={-{Latex[length=2mm,width=1.3mm]},draw=BookInk,
    line width=0.8pt,rounded corners=1mm},
  choose/.style={flow,dashed,draw=BookAmber},
  edge label/.style={fill=white,inner sep=0.6mm,
    font=\figurefont\footnotesize}
]
  \path[use as bounding box] (0,0) rectangle (164,69);
  \node[anchor=west,font=\figurefont\small] at (0,65) {(a) 同义预测的集成};
  \node[anchor=west,text=BookMuted] at (0,58) {任务：识别物种};
  \node[input] (image-a) at (9,33) {图像};
  \node[box] (model-a) at (39,45) {来源$f_1$};
  \node[box] (model-b) at (39,21) {来源$f_2$};
  \node[result] (aggregate) at (68,33) {聚合};
  \node[result] (species) at (68,6) {物种};
  \draw[flow] (image-a.east) -- (22,33) |- (model-a.west);
  \draw[flow] (22,33) |- (model-b.west);
  \draw[flow] (model-a.east) -- (68,45) -- (aggregate.north);
  \draw[flow] (model-b.east) -- (68,21) -- (aggregate.south);
  \node[edge label] at (65,51) {$p_1(y\mid\bm{x})$};
  \node[edge label] at (64,15) {$p_2(y\mid\bm{x})$};
  \draw[flow] (aggregate.east) -- (80,33) |- (species.east);
  \draw[BookRule,line width=0.6pt] (85,1)--(85,68);

  \node[anchor=west,font=\figurefont\small] at (90,65) {(b) 中间计算的组合};
  \node[anchor=west,text=BookMuted] at (90,58) {任务：属性条件判断};
  \node[input] (image-b) at (101,43) {图像};
  \node[box] (attributes) at (144,43) {属性模块$g$};
  \node[control] (selector) at (104,21) {条件选择器};
  \node[box] (reasoner) at (144,21) {判断模块$h$};
  \node[result] (decision) at (144,5) {判断};
  \draw[flow] (image-b.east) -- (attributes.west);
  \draw[flow] (attributes.south) -- (reasoner.north);
  \node[edge label,anchor=west] at (147,32) {$\bm{z}$};
  \node[edge label,anchor=east] at (142,32) {中间表示};
  \draw[flow] (image-b.south) -- (101,32) -| (selector.north);
  \draw[choose] (selector.east) -- node[edge label,above] {选择信号} (reasoner.west);
  \draw[choose] (selector.west) -- (89,21) -- (89,52)
    -| (attributes.north);
  \draw[flow] (reasoner.south) -- (decision.north);
\end{tikzpicture}

  \caption{答案汇总与中间计算连接的差别。左侧固定物种分类，两条来源接口携带类别顺序相同的分布；右侧固定属性条件判断，属性模块向判断模块传递中间表示，虚线仅传递模块选择信号。单花示例中，若上游按“红色、钟形”预测输出，下游却按反序读取，两维接口虽然能连接，实际判断已改变。多花场景还需额外保存属性所属对象。}
  \label{fig:ket-integration-interfaces}
\end{figure*}

\subsection{基于固定连接的模块组合}
\label{sec:ket-integration-fixed-modules}

先从最简单的\emph{单花}情形入手（先不碰多对象绑定难题）：输入只有一朵花，目标是\textbf{判断它是否\emph{同时}满足“红色”和“钟形”两个属性}。我们让一个属性模块先吐出两个置信度，再交给一个判断器去做“与”组合。属性模块$g$预测输出$\bm{z}=g(\bm{x})\in[0,1]^2$，依次表示红色与钟形的置信度；固定判断器采用教学规则$h(\bm{z})=\mathbf{1}\{z_1>0.8,\ z_2>0.6\}$，指示函数在两条件都满足时返回真。预测输出$(0.9,0.7)^\top$经过$h$会返回真。若下游反序读取为$(0.7,0.9)^\top$，红色阈值不再满足；若顺序正确却把钟形错估成$0.4$，也会返回假。结果相同，前一种是接口错误，后一种是属性估计错误。

这种做法就是\term{固定连接}（fixed connection）：在声明好的范围内，\textbf{对每个输入都雷打不动地计算}$h(g(\bm{x}))$，连接关系不随输入改变。也可以把\textbf{颜色、形态拆给两个不同来源}，先\emph{拼接}它们的输出再判断，即$h([g_{\mathrm{color}}(\bm{x});g_{\mathrm{shape}}(\bm{x})])$（分号表示向量拼接）。

这里有两个\textbf{新手必踩}的认识点：其一，\textbf{拼接}\emph{本身}\textbf{不会自动完成“两个条件同时成立”的判断}——把两个向量接成长向量只是凑齐了输入，\textbf{“与（AND）”这层逻辑仍要由下游}$h$\textbf{显式实现}。其二，\textbf{真实接口要约定的远不止维度}：\textbf{单位、归一化尺度、空间位置、缺失值、状态标记}都得对齐；尤其在视频场景，若某模块\textbf{需要连续帧和历史缓存}，那你\textbf{不能因为“单帧张量形状碰巧相同”就用单帧去替换它的输入}——形状对、语义不对。

当\textbf{上游输出的表示和下游“看得惯”的表示\emph{对不上}}（维度或坐标系不同），可以在两者之间\textbf{插入一个\emph{学出来}的适配映射}$r_{\bm{\phi}}$，整条链变成$h(r_{\bm{\phi}}(g(\bm{x})))$。怎么训练这个映射，取决于你\textbf{手里有什么监督}，两种情形差别很大：
\begin{itemize}\item \textbf{有“参考表示”可对}：如果每张训练照片都有一个\emph{目标表示}$\bm{z}^{\mathrm{ref}}_i$（比如已知的属性向量），就能\textbf{直接用平方残差}$\sum_i\lVert r_{\bm{\phi}}(g(\bm{x}_i))-\bm{z}^{\mathrm{ref}}_i\rVert_2^2$\textbf{逐维约束}“映射后的表示要逼近参考”，监督最直接。\item \textbf{只有最终标签}：那就\textbf{固定住来源模块}，用\textbf{整条链的任务损失}反向去学$\bm{\phi}$；这要求链\textbf{可微}（或另有可行的求解手段）。\end{itemize}
但第二种情形有个\textbf{天花板}：\textbf{最终标签只约束“最终判断对不对”，并\emph{不能保证}映射的\emph{每一维}都恢复了你\emph{预想}的那个属性含义}（多解！只要最终判断对，中间维怎么编码都行）。此外还\textbf{必须另行声明}：训练映射时，\textbf{上游、下游的参数到底更不更新}（只训适配层，还是端到端一起调）。

还有一个\textbf{“训练分布”与“使用分布”错位}的坑：就算接口语义完全正确，\textbf{下游模型也可能“\emph{不习惯}”真实的上游输出}。设想下游判断器当初\textbf{只用“真实属性的零一标签”}（非0即1、永远完整、互相独立）训练；可一旦接上真实的属性模块，它收到的却是\textbf{软概率（0.7这种）、有缺测、且各属性错误相关}的输入——\textbf{分布一变，原本很好的下游可能失灵}（这是一种典型的暴露偏差）。

怎么\textbf{定位问题出在哪}？应当做\textbf{受控对照}：\textbf{分别}用“真实中间值”和“来源实际预测的中间值”去测下游，两者差距就\textbf{隔离出了“输入变脏”这项因素}，而不会和“下游判断本身的错误”混淆。也可以\textbf{干脆拿来源的（含噪）预测去训练下游}（让它提前适应软输入），再用\textbf{独立样本}检查整链。最后一个重要认知：\textbf{冻结来源参数，只保证它在“\emph{原来的输入、状态、调用方式}”下行为不变}；一旦你\textbf{改了输入分布、加了路由、或改了最终决策方式}，\textbf{冻结并不能保护系统在新条件下仍然正确}。

\subsection{基于条件选择的模块组合}
\label{sec:ket-integration-conditional-modules}

固定连接假设“模块该怎么连”是\emph{事先定死}的。但判断“红色钟形花”和判断“病斑叶片”，可能\textbf{需要不同的提取器、不同的连接方式}——我们希望系统能\textbf{根据新任务\emph{自己挑}该用哪些模块、怎么连}。在让它自由搜索之前，必须先\textbf{立规矩}：\textbf{限定“允许使用哪些模块”以及“什么算合法接口”}，否则搜索会把\textbf{任意两个模块乱连}（比如把需要图像的模块接到文本模块上）都当成候选，既无意义又会让候选数爆炸。\term{组合语法}（composition grammar）规定模块可出现的位置及允许的模块自变量与计算结果的连接。\term{模块化元学习}（modular meta-learning）先在相关任务中学习可重组模块，再用新任务的少量样本选择结构。本节\textbf{假定模块库已经训练好}，只考察\emph{后一步“怎么选结构”}；任务分布的准备与留出条件分别由第~\ref{sec:ket-transfer-preparation}~节和第~\ref{sec:ket-transfer-task-splits}~节展开\scite{ket-alet2018-modular}

用一个极简的教学场景看清“选结构”在做什么。设模块库里有\textbf{两个属性提取器}$g_1,g_2$和\textbf{两个判断器}$h_1,h_2$（$h$能接受$g$的输出格式）。组合语法规定：\textbf{必须“一个提取器 + 一个判断器”}，不许别的连法。于是所有合法结构只有$2\times2=4$种，组成候选集合$\mathcal{G}=\{h_b\circ g_a:a,b\in\{1,2\}\}$，其中$\circ$是函数复合（先$g$后$h$）。

新任务会给一个\textbf{非空的“支持集”}$S_u$——即\textbf{专门供“在任务内挑选结构”用的少量带标签样本}。做法朴素直接：\textbf{让这四个候选结构都在支持集上跑，按单样本损失}$\ell$\textbf{比较，选平均损失最小的那个}：
\begin{equation}
  G_u^\star\in\argmin_{G\in\mathcal{G}}
  \frac{1}{|S_u|}\sum_{(\bm{x},y)\in S_u}\ell(G(\bm{x}),y).
  \label{eq:ket-integration-structure-search}
\end{equation}
读懂这个选择问题的性质：搜索\textbf{全程\emph{不改变}模块内部的参数}，$g,h$都冻住，\textbf{改变的只是“谁接谁”这个\emph{计算结构}}——这是一种\textbf{离散的结构选择}，不是梯度更新参数。由于候选集\textbf{有限且非空}，最小值\textbf{一定取得到}；若有多个结构并列最优，则\textbf{事先约定}并列时怎么挑（比如固定编号优先）。书里的数字例很直白：四个候选在四个支持样本上分别错$2,1,3,0$个，零一损失均值即$0.5,0.25,0.75,0$，于是选零错误的第四个$h_2\circ g_2$。

但这里\textbf{藏着一个和“训练/测试分离”同源的关键纪律}：\textbf{支持集$S_u$是用来“挑结构”的，不能既当裁判又当运动员}。挑完之后，必须\textbf{另用一个\emph{独立的查询集}$Q_u$}去检验它在\textbf{没见过的输入}上表现如何。\textbf{在支持集上零错误，绝不等于这个结构在新任务上风险为零}——四个样本上全对太容易靠运气或记忆得到。

四个模块时穷举很轻松，但\textbf{模块一多，候选数组合爆炸，把每个都跑一遍不现实}。这时要靠\textbf{启发式搜索}：模块化元学习常\textbf{从一个合法结构出发}，每一步\textbf{替换一个模块、或改动一条合法连接}得到邻居，并用\term{模拟退火}（simulated annealing）这类策略决定是否接受。

理解模拟退火的\textbf{哲学}很重要：它\textbf{不只在“损失下降”时接受}，当新结构\textbf{损失反而上升}时，也\textbf{以一定概率接受}——正是这种“偶尔允许走下坡路”的宽容，让搜索\textbf{有机会跳出局部最优、不至于每一步都被眼前的小改善锁死}。接受概率有两条规律：\textbf{同一温度下，损失涨得越多，接受概率越低}（小扰动更敢试）；\textbf{对同样大小的恶化，温度越低，接受概率越低}（\textbf{降温}意味着搜索从“广撒网探索”逐渐转向“精细开发”）。即便如此也要现实：四模块例子里依次替换判断器、提取器能找到不错的结构，但\textbf{在有限预算下，这类局部搜索\emph{不保证}找到全局最优}。另外要分清两种部署代价：\textbf{任务内固定结构}只需执行选中的那几个模块；\textbf{若逐输入动态改结构}，则\textbf{每次都要承担在线路由和模块切换}的开销。

“离散地二选一/四选一”有时太硬——为什么不能\textbf{在每一层\emph{软地、按比例}地同时利用所有来源模块}？AdapterFusion就是这种思路：它\textbf{冻结一个共享底座和若干\emph{已训练好的适配器}}（每个适配器对应一种能力/任务），\textbf{不改变它们}，转而\textbf{学习一个“融合层”}，让\textbf{当前样本的表示\emph{动态决定}“此刻该多用哪个适配器的表示、少用哪个”}。和前面选唯一结构不同，它\textbf{无需先挑出唯一模块}，而是\textbf{把所有候选适配器的表示向量\emph{投影后加权求和}}，再喂给后续网络。这其实就是一种\emph{注意力}机制。设底座表示$\bm{h}\in\mathbb{R}^{d_h}$、适配器表示$\bm{z}_k\in\mathbb{R}^{d}$，查询与键投到$d_a$维，值投到$d_v$维，其核心注意力可概括为
\begin{equation}
  \begin{aligned}
    \bm{q}&=\bm{W}_{Q}\bm{h},&
    \bm{k}_k&=\bm{W}_{K}\bm{z}_k,&
    \bm{v}_k&=\bm{W}_{V}\bm{z}_k,\\
    a_k&=\frac{\exp(\bm{q}^{\top}\bm{k}_k)}
                  {\sum_j\exp(\bm{q}^{\top}\bm{k}_j)},&
    \bm{z}_{\mathrm{fuse}}&=\sum_k a_k\bm{v}_k.
  \end{aligned}
  \label{eq:ket-integration-adapter-attention}
\end{equation}
这个公式就是\textbf{注意力的标准三件套}，对应关系要点破：
\begin{itemize}\item \textbf{查询$\bm{q}$来自“当前底座表示”}$\bm{h}$（它代表“我现在在处理什么、需要什么信息”）；\textbf{键$\bm{k}_k$、值$\bm{v}_k$来自“第$k$个适配器的表示”}$\bm{z}_k$（键用于“被匹配”，值是“真正取用的内容”）。\item \textbf{查询与每个键做内积}$\bm{q}^{\top}\bm{k}_k$，就是在给“此刻第$k$个适配器有多相关”\textbf{打分}；再经\textbf{指数归一化（softmax）}得到一组\textbf{非负、和为一}的权重$a_k$（这就是注意力系数）。\item 最后\textbf{用$a_k$对各适配器的“值向量”加权求和}，得到融合表示$\bm{z}_{\mathrm{fuse}}$。\end{itemize}
几个关键认知：融合表示是$d_v$维，\textbf{必须和后续网络的接口对得上}；这里省略了多头、位置编码、残差等工程细节。\textbf{训练时}，目标标签\textbf{只更新这些投影矩阵和任务头}，\textbf{底座和适配器都冻住}——所以这是一种\textbf{“激活（输出）层面的运行时组合”，而\emph{不是}第6.6节那种“离线把参数平均”}，两者作用的对象完全不同。另外别误以为它省钱：\textbf{软注意力通常仍需把\emph{所有}候选适配器的表示都算出来}才能加权，\textbf{这部分计算一般省不掉}。最后注意证据边界：原论文的主要实验里\textbf{适配器是在\emph{包含目标任务}的数据上训练过的}，因此\textbf{不能据此夸大声称“对完全没见过的任务也同样有效”}\scite{ket-pfeiffer2021adapterfusion}。

为避免和一个“长得很像”的机制混淆，要把 AdapterFusion 和第~\ref{sec:ket-sharing-conditional}节的\textbf{多门控混合专家（MoE）}区分开：\textbf{AdapterFusion 的来源适配器是\emph{预先固定}的}（它们在融合阶段不再共同更新）；而\textbf{混合专家里的“专家”是\emph{由多个任务\emph{共同训练}出来}的}。两者\textbf{都会学一组随输入变化的权重}，但\textbf{这层形式上的相似，抹不掉“专家到底有没有被联合训练”这个实质差异}。

更重要的是，无论结构选择还是表示融合，都\textbf{不能默认带来“组合泛化”}，这项能力必须\textbf{专门设计测试}。所谓\term{组合泛化}，指的是\textbf{用“见过的\emph{成分}”去完成“没见过的\emph{组合}”}——它\textbf{不等于}“在来源各自训练过的任务上答得好”。怎么测才公平？关键是\textbf{在训练里故意藏起某些组合、留到测试}：比如属性任务，训练时\textbf{让“红色”和“钟形”都\emph{单独}出现过，却把“红色\emph{且}钟形”整体留到测试}；指令任务里，\textbf{“筛选”和“排序”分别都教过，但“先筛选再排序”这条组合留出}。评测时还应\textbf{把“单一来源、固定组合、动态重组”三者并排比较}，并\textbf{检查下游是否真的用到了中间模块的结果}（而不是绕开它）。Conklin 等人提出用\emph{相似性驱动的任务采样}来构造“训练任务—推广任务”配对，为培养这种能力提供了一种训练安排；但要注意边界：\textbf{它并不保证任意合法的模块组合都能成功推广}\scite{ket-conklin2021compositional}。

前面的模块都在处理\emph{静态输入}（一张图）。一旦对象变成机器人那样的\textbf{一连串“技能/动作”}，接口的含义就\textbf{多出一维：执行动作之后\emph{环境变成了什么状态}}。比如花卉采样机器人要“\emph{先}接近植株、\emph{再}对准相机”，那么\textbf{前一个技能结束时留下的环境状态，必须恰好落在后一个技能\emph{能够接手处理}的状态范围内}，否则两个单独都很好的技能也接不上——这是“顺序组合”特有的接口约束。

强化学习里还有一种\textbf{不是“顺序串接”、而是“综合各旧策略对动作的\emph{价值判断}”}的融合，它依赖\term{后继特征}（successor features）这个概念，这里只做直觉介绍。直觉上，后继特征在\textbf{固定一个策略}的前提下，记录“\textbf{从现在起、沿这个策略走下去，未来能收获的各种特征的\emph{折扣累计}期望}”；所谓\emph{折扣}，就是\textbf{越遥远的未来，权重越小}（眼前的特征更重要）。它的巧妙之处在于一个\textbf{可迁移的分解}：如果\textbf{环境动力学、特征定义、状态-动作接口是共享的}，并且\textbf{新任务的奖励恰好能写成“这些特征的\emph{线性函数}”}（新任务 $=$ 旧特征的一组新权重），那么\textbf{拿“旧策略的后继特征”和“新任务的奖励权重”做内积}，就能\textbf{不重新训练、直接估算“旧策略在新任务上值多少”}。在此基础上，\term{广义策略改进}再\textbf{横向比较各个旧策略对当前动作给出的价值、择优选动作}。

但边界要说清：\textbf{这种“价值层面的融合”并不保证任意两个技能的“终止状态”都能物理衔接}（那是前面顺序组合的约束，两件事）；它还\textbf{强依赖“奖励是特征线性、动力学共享”等假设}。完整推导见\BookChapterRef{chap:rl-transfer}{卷7《学习范式：强化学习》第6章“任务泛化与适应”}\scite{ket-barreto2016successor}。

模块融合要\textbf{真正落地、长期运行}，还绕不开“\textbf{模块库管理}”这件看似琐碎、实则致命的事。当模块库\textbf{越攒越多}，结构候选\textbf{高度依赖“版本管理”和“检索”}：\textbf{冗余模块}（功能重叠）会\textbf{无谓地增大搜索空间}；可\textbf{删除}一个模块又可能\textbf{悄悄破坏所有引用它的组合}；而\textbf{检索漏检}（明明库里有合适模块却没搜出来）会让\textbf{本来可行的组合压根没进入候选}。

因此每一个\textbf{实际采用的组合}都应\textbf{存档}：用了哪些\textbf{模块版本}、各自的\textbf{接口}、\textbf{依赖}关系，以及出问题时的\textbf{回退清单}。成本核算也要\textbf{分段}，不能只算最后的执行：\textbf{检索、结构搜索、模块加载、实际执行}四段的开销都要分别计入。本章开头要求记录的“版本、访问权限、成本”台账在这里\textbf{依然适用}；本节额外增加的\textbf{一项检查}是：\textbf{这些\emph{已知单元}，在\emph{指定的连接路径}上，到底能不能\emph{协同}完成任务}（接口兼容不等于功能可组合）。至于多轮维护中如何保证版本一致性，由第~\ref{sec:ket-update-versions}节继续展开。

\section{基于参数的融合}
\label{sec:ket-integration-merging}

前面的“模型组合”无论软组合还是模块连接，\textbf{部署时都得\emph{带着}两个或更多模型/模块一起跑}。可如果花卉识别\textbf{终端的内存、算力、功耗只允许\emph{保存并运行一份}模型}，这条路就走不通了。于是出现第三条路线：\term{模型参数融合}（model merging）——\textbf{不在运行时组合，而在\emph{准备阶段}直接\emph{读取各个来源模型的参数}、做运算，\emph{离线造出一份\emph{新的}模型参数}，部署时只跑这一份}。

这个想法看似简单（“把参数平均一下不就行了”），但全节会反复证明它\textbf{暗藏一个致命前提}：\textbf{数组\emph{形状相同}，绝不代表“同一个坐标位置承担同一种计算”}。神经网络的隐藏单元可以\emph{任意换序}而不改变功能，于是两份“功能相同”的参数，其坐标可能\emph{对不上}，直接平均会产生垃圾。所以参数融合必须\textbf{先解决“\emph{哪些坐标可以放在一起运算}”（参数对应）}，然后才能：其一，选\textbf{融合规则}（平均完整参数？只平均相对底座的增量？）；其二，用\textbf{允许访问的样本或统计量确定融合系数}。最后永远记住风险：\textbf{“部署模型从$K$个减到1个”只是\emph{资源}上的成功，它\emph{不保证}新模型保住了所有来源的能力}。

\subsection{参数对应关系的建立}
\label{sec:ket-integration-alignment}

要真正理解“\textbf{为什么坐标对不上、平均就会坏}”，最好的办法是亲手做一个思想实验：\textbf{把一个网络的隐藏单元\emph{重新编号（交换顺序）}，但\emph{完全不改变它所计算的函数}}。同一个函数因此可以有两套“坐标完全不同”的参数——这直接暴露了“参数坐标”和“功能”之间不是一一对应的。设输入$\bm{x}\in\mathbb{R}^d$、隐藏层$m$个单元、输出$c$维、激活$\sigma$逐元素作用：设预测输入$\bm{x}\in\mathbb{R}^d$，隐藏层有$m$个单元，预测输出为$c$维，激活$\sigma$逐元素作用：
\[
  f(\bm{x})=\bm{W}_2\sigma(\bm{W}_1\bm{x}+\bm{b}_1)+\bm{b}_2,
  \quad \bm{W}_1\in\mathbb{R}^{m\times d},\
  \bm{W}_2\in\mathbb{R}^{c\times m}.
\]
偏置为$\bm{b}_1\in\mathbb{R}^m$、$\bm{b}_2\in\mathbb{R}^c$。引入一个\term{置换矩阵}$\bm{P}\in\mathbb{R}^{m\times m}$：它是一个\textbf{每行每列恰有一个1、其余为0}的正交矩阵，作用是把隐藏单元\textbf{整体重新排列}，所以满足$\bm{P}^{\top}\bm{P}=\bm{I}$（换过去还能换回来）。由于激活$\sigma$是\textbf{逐元素独立}作用的，它和这种重排\textbf{可交换}：先换序再激活，等于先激活再换序，即$\sigma(\bm{P}\bm{v})=\bm{P}\sigma(\bm{v})$。

关键在于：\textbf{只重排第一层会把连接关系弄乱}。为了让“换序后的每个隐藏单元\textbf{仍然连回它原来对应的输出权重}”，必须\textbf{同时、配套地变换两层}——第一层把单元换过去，第二层就要相应地把列换回来：
\begin{equation}
  \begin{aligned}
    \widetilde{\bm{W}}_1&=\bm{P}\bm{W}_1,&
    \widetilde{\bm{b}}_1&=\bm{P}\bm{b}_1,\\
    \widetilde{\bm{W}}_2&=\bm{W}_2\bm{P}^{\top},&
    \widetilde{\bm{b}}_2&=\bm{b}_2
  \end{aligned}
  \label{eq:ket-integration-permutation}
\end{equation}
把新参数代回网络，输出变成
$\bm{W}_2\bm{P}^{\top}\bm{P}\sigma(\bm{W}_1\bm{x}+\bm{b}_1)+\bm{b}_2=f(\bm{x})$，
因为中间的$\bm{P}^{\top}\bm{P}=\bm{I}$正好抵消。结论是：\textbf{第一层的行换了序、第二层同步换了列，\emph{中间的坐标命名}全变了，但\emph{对外的函数}一点没变}。反过来，如果你\textbf{只改第一层、不配套改第二层}，置换\textbf{通常无法抵消}，函数就真的被破坏了。

这个看似简单的等价关系，正是\textbf{整个参数融合领域的“理论警钟”}：\textbf{同一个函数，能对应无穷多套“坐标排列不同”的参数}；因此\textbf{仅凭“函数相同”，无法唯一确定一份“可以逐坐标拿来平均”的参数数组}。两个独立训练的模型，哪怕功能一模一样，也很可能各自选用了\emph{不同的隐藏单元排列}——直接对坐标做算术，等于是在\textbf{把“苹果的第3维”和“橘子的第3维”相加}。

\begin{example}[完全相同的函数也不能直接平均]
  \label{ex:ket-integration-permutation}
  取一个标量自变量、标量预测输出的教学网络，偏置为零，$\operatorname{ReLU}(v)=\max(v,0)$。两层权重和隐藏单元交换为
  \[
    \bm{W}_1=\begin{bmatrix}1\\-1\end{bmatrix},\quad
    \bm{W}_2=\begin{bmatrix}1&1\end{bmatrix},\quad
    \bm{P}=\begin{bmatrix}0&1\\1&0\end{bmatrix}.
  \]
  \textbf{这是全节最值得亲手算一遍的例子}，它用极简的网络证明“功能相同也能平均出零”。两个来源\textbf{计算的是\emph{完全相同}的函数}$\operatorname{ReLU}(x)+\operatorname{ReLU}(-x)=|x|$（绝对值）。但它们的第一层参数分别是$(1,-1)^\top$和$(-1,1)^\top$——后者其实就是前者把两个隐藏单元\emph{对调}（置换$\bm{P}$）的结果。

现在\textbf{天真地直接平均}第一层：$[(1,-1)+(-1,1)]^\top/2=(0,0)^\top$，得到\textbf{零向量}；后一层保持$[1\ 1]$不变，则融合后的网络\textbf{无论输入什么都恒输出0}。具体在$x=1$处：\textbf{两个来源各自都输出1（正确），融合模型却输出0}——把两个“绝对值函数”平均，竟得到一个“恒零函数”！

\textbf{正确做法}是：\textbf{先对第二个来源施加\emph{逆置换}}（用$\bm{P}^{\top}$把它的单元顺序拧回来），让两份数组在“同一坐标语义”下恢复一致，\textbf{然后再平均}，结果才仍是$|x|$。所以请牢记这个例子的判词：\textbf{融合失败\emph{不是因为来源函数不同}（它们明明相同），而是因为“\emph{平均之前没有先恢复参数对应}”}。
\end{example}

既然问题是“隐藏单元对不上”，自然要想办法\textbf{自动把两个网络的单元\emph{配对}}——这正是 \textbf{Git Re-Basin} 这类工作的主题。配对主要靠两类信息：
\begin{itemize}\item \term{激活匹配}（activation matching）：喂入\textbf{双方共同的输入}，比较\textbf{每个单元被激活的统计/向量}，\emph{激活行为越像}的单元越可能是同一个；在相似度矩阵上求一个\textbf{一一指派}（每个单元只许配给对方一个单元）。它\textbf{需要有代表性的输入数据}。\item \term{权重匹配}（weight matching）：\textbf{不看数据、只读参数}，\textbf{逐层调整置换}，目标是让“对齐后双方对应权重尽可能一致”。它\textbf{无需训练样本}。\end{itemize}
无论哪种，\textbf{一一指派}都是硬约束，而且\textbf{相邻层、残差连接的重排必须\emph{配套}}（呼应前面“两层要一起换”）。

这里要澄清一个常被夸大的结论——“找到好置换”和“\textbf{线性连通}”不是一回事。所谓\term{低损失线性连通}（linear mode connectivity），是指\textbf{对齐之后，在两份参数之间\emph{连一条直线}，沿线上\emph{任意}插值出的模型都\emph{没有明显的损失障碍}（损失都很低）}。\textbf{仅仅找到一个让端点对齐的置换，\emph{不能推出}整条线段都低损失}；而且原论文还\textbf{构造了一些“无论怎么置换都无法线性连通”的解}（它们不是随机梯度下降通常找到的那类解）。所以准确的说法是：该结论\textbf{针对“SGD 找到的解之间往往可连通”}，\textbf{并不能推广成“任意两个解都能靠置换连通”}\scite{ket-ainsworth2023rebasin}。

“强行一一配对”还有个\textbf{语义错配}风险：如果两个来源\textbf{学的是不同任务}（一个识别\emph{物种}、一个识别\emph{病害}），它们的特征\textbf{用途本来就不同}，硬把单元配成一一对应，可能是在\textbf{强行撮合“本不是一回事”的特征}。ZipIt! 对此放宽了限制：它利用\textbf{无标签输入上的激活\emph{相关性}}，不只\textbf{跨来源}匹配特征，还能\textbf{合并“同一来源\emph{内部}”冗余的特征}（即允许“多对一/部分合并”，不再死守一一对应）；并且可以\textbf{只融合前面若干共享层、保留后面各自的任务分支}。

这在性质上\textbf{不只是“重排第二个模型的单元”}，而是\textbf{改变了匹配的约束本身}（允许合并、允许部分融合）。也要算清部署账：\textbf{一旦保留了多个任务分支}，部署时\textbf{可能仍需知道“当前是哪个任务”、并携带对应的预测头}，所谓“融成一份”就要打折扣——实际的\textbf{计算和存储都必须按这张真实的、带分支的结构去核算}\scite{ket-stoica2024zipit}。

“参数对不上”在\textbf{低秩适配器}里还有\textbf{更隐蔽的一层}。先回顾\term{低秩适应}（low-rank adaptation, LoRA）：为了便宜地微调，它\textbf{不直接改整个权重矩阵}，而是把“相对底座的权重变化”\textbf{表示成两个小矩阵的乘积}$\Delta\bm{W}_k=\bm{B}_k\bm{A}_k$，用一个\textbf{较小的内维数}$r$来\textbf{限制这个增量的秩}（参数量远小于全矩阵）。某层写为$\Delta\bm{W}_k=\bm{B}_k\bm{A}_k$，其中$\bm{B}_k\in\mathbb{R}^{d_{\mathrm{out}}\times r}$、$\bm{A}_k\in\mathbb{R}^{r\times d_{\mathrm{in}}}$，$d_{\mathrm{in}},d_{\mathrm{out}}$分别为层表示与层响应维数，$r$为内维数。（下面固定相同的缩放约定，并把缩放常数吸收进因子；LoRA 的完整训练机制与开销见第~\ref{sec:ket-transfer-added-parameters}节。）问题在于这个乘积分解\textbf{天然带着“内部坐标自由度”}：对\textbf{任意一个可逆矩阵}$\bm{R}\in\mathbb{R}^{r\times r}$，都有
$$(\bm{B}_k\bm{R})(\bm{R}^{-1}\bm{A}_k)=\bm{B}_k\bm{R}\bm{R}^{-1}\bm{A}_k=\bm{B}_k\bm{A}_k=\Delta\bm{W}_k.$$
也就是说，\textbf{同一个增量$\Delta\bm{W}_k$，可以用无穷多组\emph{不同的因子}$(\bm{B},\bm{A})$表示}（中间随便插一个$\bm{R}$再抵消）。更何况\textbf{不同任务的适配器，可能根本处在\emph{不同的更新子空间}里}。结论因此很明确：\textbf{LoRA 因子即使\emph{形状完全相同}，其\emph{坐标也\emph{不}自动可比}}——不能看到两个$\bm{B}$都是$d_{\mathrm{out}}\times r$就直接相加。

既然 LoRA 因子$\bm{B}\bm{A}$的内部坐标可以任意旋转、直接对因子做运算不可靠，一个更稳的思路是\textbf{绕开“各自的因子”，回到“\emph{完整增量}”这个唯一、坐标无关的对象上再建立共同坐标}。KnOTS 的做法是：\textbf{先把每个任务实际的增量$\Delta\bm{W}_k$还原出来}，把同一层所有任务的增量\textbf{横向拼接}成一个大矩阵，再对这个大矩阵做\textbf{联合奇异值分解（SVD）}。这样一来，\textbf{所有任务的增量都被投影到\emph{同一组基}上来表达}，而\textbf{不是}分别采信每个来源各自那套可以自由旋转的因子。对某一层可写为：对某一层可写为
\begin{equation}
  \begin{aligned}
    \relax[\Delta\bm{W}_1\ \cdots\ \Delta\bm{W}_{K_{\mathrm{s}}}]
      &=\bm{U}\bm{\Sigma}[\bm{V}_1^\top\ \cdots\ \bm{V}_{K_{\mathrm{s}}}^\top],\\
    \Delta\bm{W}_k&=\bm{U}\bm{\Sigma}\bm{V}_k^\top.
  \end{aligned}
  \label{eq:ket-integration-knots}
\end{equation}
读懂这个分解里各块的角色：$q$是拼接矩阵的秩；$\bm{U}$的列是\textbf{所有任务共享的“层响应侧”正交基}，$\bm{\Sigma}$集中保存奇异值，而每个任务$k$的\textbf{个性}体现在各自的块$\bm{V}_k^\top$上。在\textbf{不截断}奇异值时这是一个\textbf{精确的恒等分解}；于是融合规则可以\textbf{先在这些“共同坐标块”$\bm{V}_k^\top$上做运算，最后再统一乘回$\bm{U}\bm{\Sigma}$}还原成真实增量。

要注意它的\textbf{能力边界}，别神化：这是“\textbf{为了在一致坐标上做融合}”而做的坐标变换，\textbf{不是“融合完为了压缩再去分解”}，顺序和目的都不同。而且\textbf{如果后续操作只是\emph{纯线性相加}}，那么在精确分解下\textbf{结果与直接相加并无区别}（换个基不改变线性加和）；共同基真正发挥作用的地方，是\textbf{裁剪（按奇异值截断）、符号选择等\emph{非线性}操作}——这些操作\textbf{依赖坐标怎么选}，统一基才有意义。一句话：\textbf{共同基只提供“一套一致的运算坐标”，它\emph{不会仅因为做了分解}就给模型增加任何新的任务能力}\scite{ket-stoica2025knots}。

实践中\textbf{最可靠、最常见的“天然对应”来自\emph{共同底座}}：当所有微调都\textbf{从\emph{同一个底座}出发}时，它们的增量\textbf{引用的是同一版本、同样的层、以及层前后一致的表示坐标}，这就\textbf{免费提供了“参数可以逐坐标对应”的起点}（这也是任务向量类方法好用的根本原因）。但反过来有一条硬约束：\textbf{如果你手里只有适配器（增量），就\emph{必须同时保留它所对应的那个\emph{准确底座版本}}}——适配器$\Delta\bm{W}$离开底座$\bm{W}_0$就没有意义。

很多“看起来能融”的情形其实不满足前提，要警惕：\textbf{两个\emph{各自独立从头训练}的底座}之间\textbf{没有这种共同坐标}，\textbf{不能把A底座上训出的适配器硬加到B底座上}。结构差太远时，可能得\textbf{读入完整模型去重新找层/激活对应}，或者干脆\textbf{放弃参数融合、改用蒸馏}（用教师行为训练承接模型）。如果\textbf{深度、宽度、算子都不同}，那么“逐参数运算”在数学上\textbf{根本没有定义}；这时\textbf{靠补零、或按“同名层”匹配}都\textbf{只是表面凑形状}，\textbf{证明不了语义兼容}。最后哪怕对应能建立，\textbf{仍要检查\emph{任务冲突}}（两个能力可能在参数上打架）——\textbf{共同起点只是“好对应”的\emph{必要}便利，\emph{不是}融合成功的\emph{充分}条件}。

\subsection{融合规则的构造}
\label{sec:ket-integration-merging-rule}

对应问题解决后，下一个问题是“\textbf{具体对\emph{什么}做平均}”：是平均\textbf{完整参数}，还是只平均\textbf{相对底座的变化（增量）}？把参与融合的某个参数块\textbf{展平}成向量$\bm{\theta}_k\in\mathbb{R}^{d_\theta}$，共同底座记为$\bm{\theta}_0$，$d_\theta$是这个块的坐标总数。

先做一个\textbf{记号上的提醒}：这里的系数$\alpha_k$\textbf{改指“参数组合系数”}，它\textbf{不再承担概率混合里“非负、和为一”的约束}——参数不是概率，后面会看到它可以为负、和也不必为一。那么“\textbf{直接加权\emph{完整}参数}”和“\textbf{加权\emph{相对底座的增量}}”到底什么关系？用一个\textbf{在每项里同时加、减底座}$\bm{\theta}_0$的恒等变形就能看清：
\begin{equation}
  \sum_k\alpha_k\bm{\theta}_k
  =\left(\sum_k\alpha_k\right)\bm{\theta}_0+
    \sum_k\alpha_k(\bm{\theta}_k-\bm{\theta}_0).
  \label{eq:ket-integration-rule-equivalence}
\end{equation}
这个恒等式把两种做法的关系讲透了，分两种情况读：
\begin{itemize}\item \textbf{当系数和为一}$\sum_k\alpha_k=1$时：底座那一项的系数恰好是1，于是\textbf{“直接平均完整参数”}和\textbf{“底座不动、只把各增量按同样系数加进去”}（$\bm{\theta}_0+\sum_k\alpha_k(\bm{\theta}_k-\bm{\theta}_0)$）\textbf{完全等价}。\item \textbf{当系数和\emph{不}为一}时：两者就\textbf{不等价}了——直接加权完整参数还会\textbf{额外地把底座\emph{缩放}}$(\sum_k\alpha_k)$倍，可能偏离原始底座的尺度。\end{itemize}
还要强调适用边界：这只是一个\textbf{线性运算}的恒等关系；一旦融合规则涉及\textbf{裁剪、或对低秩因子做相乘}等非线性操作，这个等式\textbf{不能无条件照搬}。下面分两小节分别讨论“\emph{完整参数}”和“\emph{参数增量}”的融合，至于“\textbf{系数到底取多少、由什么证据决定}”，留给再下一小节。

\subsubsection{完整参数的融合}
\label{sec:ket-integration-full-weights}

先看最自然、也最好用的场景：\textbf{同一个花卉分类目标}，且所有来源都是\textbf{从\emph{同一个预训练模型}出发、经过不同方式微调}得到的（它们坐标天然对得很齐）。Model Soups 的做法朴素得出奇：\textbf{用不同的学习率、数据增强、或训练随机种子}去微调出很多份权重，\textbf{把这些权重当“配料”都留着，最后对权重做平均}（像把多种汤料炖成一锅“汤”）。

它给了几种\textbf{怎么配料}的策略：\term{均匀汤}（uniform soup）最简单，对所有来源\textbf{等权}$1/K_{\mathrm{s}}$；\term{贪心汤}则\textbf{先按验证准确率给配料排序，从最好的那个开始，\emph{逐个}试着加进来——只有当“加它之后、平均模型的验证准确率\emph{没有下降}”才真正保留}，否则丢弃；也可以\textbf{去学习}一组插值系数而非取等权。

为什么这么简单的平均常常有效？因为\textbf{共同的预训练起点}让这些微调解\textbf{大概率落在同一个低损失“盆地”里、坐标也对应}，平均相当于在盆地里找一个更居中、更平滑的点。但\textbf{边界要讲清}：\textbf{贪心汤的“接纳/拒绝”是\emph{依据验证集结果}做的选择}，它\textbf{只保证验证表现不降}，\textbf{并不保证换到\emph{独立测试集}也不降}（验证集选过了就带选择偏差）\scite{ket-wortsman2022-soups}。

等权平均有个“\textbf{一视同仁}”的盲点：\textbf{不同参数坐标对任务的\emph{重要性}差别很大}，有的坐标动一点损失就剧变、有的坐标怎么动都无所谓，等权平均\textbf{没有利用这种差异}。Fisher merging 引入\term{Fisher信息}来\textbf{按坐标重要性加权}。

先把这个略显抽象的量讲明白：Fisher 矩阵的\textbf{第$i$个对角元}，是在\textbf{来源的输入分布下、按来源模型自己抽取输出}时，“\textbf{对数预测概率关于第$i$个参数的\emph{梯度}的平方}”的\textbf{期望}。它的\textbf{直观含义}是“\textbf{这个概率模型对该参数\emph{局部扰动}有多敏感}”：梯度平方大，说明\textbf{轻微挪动这个参数就会让预测概率大幅变化}，即这是个\textbf{关键、确定}的坐标；在适当正则条件下，它还\textbf{与“期望负对数似然”在该点的\emph{曲率}对应}（曲率大＝山谷陡、不该偏离）。Fisher merging 于是用这些\textbf{非负的对角估计}$F_{k,i}$为每个来源$k$构造一个\textbf{局部二次近似}：直觉上，它给融合参数$\theta_i$和来源原值$\theta_{k,i}$之间的偏差加一个惩罚$F_{k,i}(\theta_i-\theta_{k,i})^2$，\textbf{某个坐标在某来源里 Fisher 值越大（越关键），就\emph{越不允许}融合结果偏离该来源在这个坐标上的值}。最小化所有来源的这种加权偏离之和，驻点（令导数为零）给出一个加权平均：最小化$\sum_k F_{k,i}(\theta_i-\theta_{k,i})^2$，驻点给出
\[
  \theta_i^{\mathrm{F}}
  =\frac{\sum_k F_{k,i}\theta_{k,i}}
         {\sum_k F_{k,i}},
\]
当分母为正时，这个比值就是该二次目标的解析最小点（\textbf{本质是“按 Fisher 重要性做加权平均”}）；若\textbf{分母为零}（所有来源在这个坐标上都不敏感），二次目标对该坐标\textbf{没有任何约束}，必须\textbf{预先规定一个回退值}（比如退回等权或底座值），不能放任。

使用时有几个\textbf{不能混淆的细节}：如果为了省事，用\textbf{观测到的真实标签}去替代“按模型分布抽取输出”，得到的叫\term{经验Fisher}，它\textbf{并不无条件等于}严格的模型 Fisher、也不等于样本 Hessian（外积形式少了损失二阶项）。这些 Fisher 统计量\textbf{必须用来源自己的样本估计、或让来源方随模型一并保存}（部署方未必有它的数据）。而且\textbf{“只取对角”这个近似，\emph{忽略了参数之间\emph{相关}的方向}（重要的曲率可能在斜向上）}。所以结论要克制：\textbf{在“局部二次近似”这个目标下最优，\emph{不等于}在真实的多任务风险上最优}——它是个有明确前提的、便宜的加权方案。同类“参数敏感性估计用于约束更新”的思想，第~\ref{sec:ket-update-objective}节还会从“持续学习防遗忘”的角度再用\scite{ket-matena2022-fisher}。

Fisher 是从“损失曲率/参数重要性”出发；RegMean 换了一个\textbf{更贴近“这一层实际在干什么”}的角度：\textbf{线性层的“响应（activation）偏差”}。它的目标非常具体——\textbf{让融合后的那一层，在\emph{每个来源\emph{自己}的数据}上，都尽量\emph{复现}该来源原来那一层\emph{产生的中间响应}}（即不希望融合改变下一层收到的输入）。为保持矩阵维度明确，此段采用行样本约定：$\bm{X}_k\in\mathbb{R}^{n_k\times d_{\mathrm{in}}}$，$\bm{W}_k\in\mathbb{R}^{d_{\mathrm{in}}\times d_{\mathrm{out}}}$，批预测输出为$\bm{X}_k\bm{W}_k$。取$n_k>0$，层自变量二阶矩阵为$\bm{G}_k=\bm{X}_k^\top\bm{X}_k/n_k$。最小化$\sum_k n_k^{-1}\lVert\bm{X}_k(\bm{W}-\bm{W}_k)\rVert_{\mathrm{F}}^2$，正规方程给出
\[
  \bm W^{\mathrm{R}}
  =\left(\sum_k\bm G_k\right)^{-1}
    \left(\sum_k\bm G_k\bm W_k\right),
\]
闭式解成立有前提：这里\textbf{假设}$\sum_k\bm G_k$\textbf{可逆}；一旦矩阵\textbf{病态或奇异}（某些输入方向数据里几乎不出现），就必须\textbf{另定正则化（比如加阻尼项）和求解规则}，不能直接求逆。实现时还有个\textbf{容易转置错}的坑：\textbf{这里为了把样本放进行里，权重写成了$d_{\mathrm{in}}\times d_{\mathrm{out}}$（行样本约定）}，和前文“列向量层权重”的方向\textbf{是转置关系}，写代码时\textbf{必须全程保持同一约定}。

式中除以$n_k$的归一化也有用意：它让\textbf{每个来源贡献的是“\emph{单位样本的平均}二阶偏差”}，而\textbf{不会因为某来源样本多就\emph{天然}霸占更大权重}（若想按样本数加权应显式表达，而非偷偷由归一化决定）。数据隐私/可得性上有个好消息：\textbf{真正的原始样本可以用“事先保存好的二阶统计量}$\bm G_k$\textbf{”替代}，不必传数据本身——但这些统计量\textbf{终究还是要用来源的输入数据才能产生}，并非凭空可得。最后是适用范围：RegMean 是\textbf{对每一层\emph{独立}求最优}，\textbf{逐层都最优\emph{不保证}整个非线性网络端到端的预测偏差最小}（各层最优的组合不必然是整网最优）\scite{ket-jin2023-regmean}。

Fisher merging 和 RegMean 都是\textbf{“在做任何预测\emph{之前}，就把参数改掉”}的方法，所以它们\textbf{属于参数融合，\emph{不能}当成“对来源\emph{预测}的加权平均”}。这俩概念必须彻底分清，记号上对照：
\begin{itemize}\item \textbf{参数平均（model merging）}：先把参数融成$\sum_k\alpha_k\bm{\theta}_k$，运行的是\textbf{一个模型}$f(\bm{x};\sum_k\alpha_k\bm{\theta}_k)$；\item \textbf{预测平均（ensemble）}：\textbf{先分别运行}$K$个模型$f(\bm{x};\bm{\theta}_k)$，再把它们的输出平均。\end{itemize}
这两者\textbf{什么时候才相等}？只有当\textbf{函数$f$对被平均的参数\emph{恰好是线性}的}，“先平均参数”和“先算再平均输出”才能交换次序。而真实神经网络里\textbf{多层矩阵相乘、非线性激活（ReLU/softmax）\textbf{通常会破坏这种线性关系}}，所以\textbf{一般不相等}。下面的例子\textbf{固定来源、输入、权重都不变，只挪动“平均发生的位置”}，让你亲眼看到差异。

\begin{example}[非线性使两种平均不相等]
  \label{ex:ket-integration-nonlinear-average}
  取教学分数模型$f(x;w)=\operatorname{ReLU}(wx)$，来源模型参数$w_1=2,w_2=-1$，预测输入$x=1$，两权重都为$1/2$。先只看激活前的分数（暂不做 softmax）：两个来源在$x=1$处的分数分别是$f(1;2)=\operatorname{ReLU}(2)=2$和$f(1;-1)=\operatorname{ReLU}(-1)=0$。
\begin{itemize}\item \textbf{预测平均}（先算后平均）：$(2+0)/2=1$。\item \textbf{参数平均}（先平均参数再算）：$\bar w=(2+(-1))/2=0.5$，$f(1;0.5)=\operatorname{ReLU}(0.5)=0.5$。\end{itemize}
同样两个来源、同样等权，\textbf{只因“平均”发生在 ReLU 的前后两侧，答案就一个是1、一个是$0.5$}。这就是\textbf{非线性激活让“平均次序不可交换”}的最干净证据。

  将分数向量$(f(x;w),0)$归一化为二分类概率，第一坐标表示正类。两个来源正类概率为$e^2/(1+e^2)\approx0.8808$和$0.5$，平均约为$0.6904$；融合模型参数后的正类概率约为$0.6225$。先平均分数再做softmax得到约$0.7311$，又不同于平均概率。再把分数归一化成“二分类正类概率”（第一维为正类），差异更丰富：两个来源的正类概率分别约为$0.8808$和$0.5$，\textbf{先算概率再平均}得约$0.6904$；而\textbf{先平均参数（$\bar w=0.5$）再算概率}约为$0.6225$；至于\textbf{先平均\emph{分数}（logits 平均为1）再统一 softmax}，又得到约$0.7311$。于是出现\textbf{三个不同的数}。请记住这个现象的\textbf{唯一原因}：\textbf{来源、输入、权重全都一样，区别仅在于“平均”和“非线性（ReLU、softmax）”的\emph{先后顺序}}。工程上“到底在哪一层融合”因此不是无关紧要的实现细节，而是会改变结果的实质选择。
\end{example}

前面 Model Soups 的成功有个\textbf{隐含前提：所有来源在解\emph{同一个}分类目标}。把\textbf{物种模型和病害模型}往一起融，难度上了一个台阶，因为\textbf{多了“任务差异”}：它们的\textbf{输出头类别含义完全不同}（一个头在山茶/茶梅/月季上做 softmax，另一个头在各种病害上做 softmax），\textbf{这些输出坐标\emph{没有对应关系}，绝不能按同一坐标平均}。

实际可行的做法是\textbf{“融主干、留分头”}：\textbf{只融合那些\emph{可对应}的\emph{共享特征层}（主干）}，而\textbf{保留每个任务\emph{各自独立}的预测头}。但要意识到这\textbf{并没有完全实现“一份模型走天下”}：部署时\textbf{可能仍需先知道“当前是什么任务”}，才能\textbf{选用对应的那个头}（\textbf{任务身份}问题还在）；形成一个共享主干，\textbf{并不自动解决“面对未知任务该选哪个头”}。所以\textbf{Model Soups 在“同一目标多个解”上的好结果，\emph{不能直接外推}成“任意不同技能都能随便融合”}。

融合不同任务后，\textbf{怎么算“好不好”}也要改：不能只看一个任务的分数，而应\textbf{在每个来源任务上\emph{分别}测量风险}，再\textbf{按事先声明的偏好}汇总（见6.6.4的帕累托视角），\textbf{用一个平均分掩盖某任务严重退化是自欺欺人}。成本也要\textbf{按真实结构}算：如果最终\textbf{只跑一个共享主干 + 当前任务头}，计算确实可能少于两个完整模型；但\textbf{若保留了多个任务分支、或需要重估批归一化统计、重新校准}，又会\textbf{增加额外开销}。

这里专门点一个\textbf{极易出错的“伪参数”}：\textbf{批归一化（BatchNorm）里保存的“运行均值/方差”}。它们\textbf{不是学出来的普通权重，而是\emph{强依赖输入数据分布}的\emph{统计量}}，\textbf{绝不能未经核查就像普通可训练参数那样拿来平均}——融合后数据分布一变，这些统计量往往要\textbf{在新数据上重新估计}，否则归一化就是错的。总之，\textbf{部署结构、状态（归一化统计等）处理、以及各任务能力是否保留}，都需要\textbf{一份版本严格对应的模型配置清单}来管理，不能想当然。

Self-Soupervision进一步改变来源的训练目标，让共同底座在不同自监督目标或样本上继续训练。其方案既有之后再用标签微调的来源，也有直接融合自监督来源、用少量标签选系数的安排，以及无监督选择变体。它扩展了可融合来源的条件，但“继续训练不用标签”不等于“选择不用标签”，更不等于“评价不用标签”。训练或选择使用目标分布的无标签预测输入，需要声明这种目标任务样本访问，不能作为完全未见目标域的标准领域泛化。若先用独立适应集准备模型，再冻结模型预测与之分离的新预测输入，仍可评价适应后的归纳预测；若这批待预测输入本身也参与模型准备，则属于直推式（亦称转导式）预测，或按具体更新方式声明测试时适应。原论文分别报告适应样本与评价样本分离和重合的设置，不能把“来自同一分布”直接等同于“使用同一批样本”。这里仍按预印本的限定证据使用这项工作\scite{ket-fuller2026selfsoup}

\subsubsection{参数增量的融合}
\label{sec:ket-integration-deltas}

当我们的\textbf{目标不是“在同一任务上找更稳的点”，而是“\emph{保住一个通用底座}，同时把\emph{多个不同任务}的微调成果组合起来”}时，直接平均完整参数就不够贴切了——更自然的做法是\textbf{只对“相对底座的\emph{变化量}”做运算}。

\term{任务向量}（task vector）因此被定义为\textbf{同一坐标下，微调后参数与共同底座参数之差}$\bm{\tau}_k=\bm{\theta}_k-\bm{\theta}_0$。它的几何含义是“\textbf{从底座指向该任务解的一个\emph{参数位移向量}}”。\textbf{先给这个概念去魅}：任务向量\textbf{只是一个\emph{参数空间里的位移}}，它\textbf{并不等于“一项已经被干净剥离出来的、纯语义的\emph{技能}”}（$\bm{\tau}$里混着该任务的各种变化，不能保证“减去它就精确遗忘某技能”，这点后面还会强调）。\term{任务算术}（task arithmetic）的操作就是三步：\textbf{先作差}得到各任务向量、给每个增量\textbf{乘一个缩放系数}、再\textbf{加回底座}：
\begin{equation}
  \bm{\tau}_k=\bm{\theta}_k-\bm{\theta}_0,\qquad
  \bm{\theta}_{\mathrm{merge}}(\bm{\lambda})
    =\bm{\theta}_0+\sum_{k=1}^{K_{\mathrm{s}}}\lambda_k\bm{\tau}_k.
  \label{eq:ket-integration-task-arithmetic}
\end{equation}
式中$\bm{\tau}_k$是第$k$个任务的位移，$\bm{\lambda}$把所有缩放系数收集在一起。\textbf{原始的任务算术为简单起见常让所有$\lambda_k$取同一个值}（统一缩放），但完全可以\textbf{扩展成每个来源各自一个系数}（接6.6.4的系数学习）。

这里务必\textbf{再次厘清“参数系数”和“概率权重”的区别}：\textbf{模型参数\emph{不是}概率}，所以这里的$\lambda_k$\textbf{既不需要非负、也不需要和为一}。你可以用\textbf{负系数}、可以用大于1的系数。但要\textbf{警惕对负系数的过度解读}：出现负值和大尺度，只\textbf{在几何上意味着“可以沿更多方向移动参数、步幅可大可小”}，\textbf{绝不能由此想当然地推出“能力可以像数字一样精确地\emph{相减}（负系数＝遗忘某技能）、或可以\emph{任意放大}不失真”}——参数空间的距离和“语义上的技能加减”之间没有那种简单对应\scite{ket-ilharco2022taskarithmetic}。

朴素地把任务向量相加会遇到两个麻烦：其一，\textbf{同一坐标上不同任务的增量\emph{正负相反}，会\emph{互相抵消}}（你拉我推，结果内耗）；其二，\textbf{大量坐标只有\emph{微不足道}的小幅变化}，把它们都加进来既稀释信号、又可能是噪声。TIES-Merging 用\textbf{三个有顺序的步骤}来解决（方法名 TIES 正对应 Trim、Elect Sign、Disjoint Merge 等意）：
\begin{itemize}\item \textbf{裁剪（Trim）}：每个任务向量\textbf{只按绝对幅度保留\emph{给定比例}最重要的坐标}，其余置零，得到$\widetilde{\bm{\tau}}_k$（小变化视为不重要）。\item \textbf{选方向（Elect Sign）}：对每个坐标，把各任务保留下来的（\emph{带正负号、带幅度}的）增量求代数和，\textbf{用和的符号作为该坐标的“\emph{选举方向}”}$s_j$。\item \textbf{同号合并}：\textbf{只把\emph{与当选方向同号}的非零增量拿来平均}（异号的视为冲突票，\textbf{不参与}这一步合并）。\end{itemize}
注意这里的“\textbf{方向}”\textbf{不是按“有几个任务投正、几个投负”的\emph{人头}投出来的}，而是\textbf{按“带幅度的代数和”决定}（一个大幅增量能压过几个小幅增量）；合并的\textbf{平均分母也只数真正参与的同号来源}。完整写法为：
\begin{equation}
  \begin{aligned}
    s_j&=\operatorname{sign}\!\left(\sum_k\widetilde{\tau}_{k,j}\right),\\
    I_j&=\{k:\widetilde{\tau}_{k,j}\neq0,\
                    \operatorname{sign}(\widetilde{\tau}_{k,j})=s_j\},\\
    \tau_{\mathrm{TIES},j}
      &=\frac{1}{|I_j|}\sum_{k\in I_j}\widetilde{\tau}_{k,j},\\
    \bm{\theta}_{\mathrm{merge}}
      &=\bm{\theta}_0+\lambda\bm{\tau}_{\mathrm{TIES}}.
  \end{aligned}
  \label{eq:ket-integration-ties}
\end{equation}
读懂这几行：$\operatorname{sign}(\cdot)$返回正、负或零；$s_j$是坐标$j$选出的方向；$I_j$是“在该坐标上\textbf{非零且与$s_j$同号}”的来源集合；最后在$I_j$内做平均得到$\tau_{\mathrm{TIES},j}$，再加权$\lambda$加回底座。为了让结果\textbf{完全可复现}，必须\textbf{事先规定边界情形}：本章约定\textbf{当}$s_j=0$（正负恰好抵消）\textbf{或}$I_j$为空（没有同号者）时，该坐标取零；当\textbf{按幅度裁剪出现“并列”}（第$r$名和第$r+1$名幅度相等）时，\textbf{按固定的来源索引顺序}决定取舍，而不是随机。

再次强调两个\textbf{不能误读}之处：\textbf{选方向凭的是带幅度的和，不是来源个数}；\textbf{平均分母只计实际同号参与的来源}（异号票被排除，不能让它们把分母撑大）。而且\textbf{保留比例设多少、并列怎么裁、最终缩放$\lambda$取几，每一个都会改变答案}，这些\textbf{超参数和约定都应写进版本清单}，保证可追溯\scite{ket-yadav2023ties}。

用三维教学模型参数复算这一操作，底座与两来源为
\[
  \begin{aligned}
    \bm{\theta}_0&=(1,1,1)^\top,\\
    \bm{\theta}_1&=(1.4,0.8,1.01)^\top,\\
    \bm{\theta}_2&=(1.2,1.3,0.98)^\top.
  \end{aligned}
\]
按任务向量定义相减，得到
\[
  \begin{aligned}
    \bm{\tau}_1&=(0.4,-0.2,0.01)^\top,\\
    \bm{\tau}_2&=(0.2,0.3,-0.02)^\top.
  \end{aligned}
\]
每个来源保留两个最大幅度坐标后，增量为$(0.4,-0.2,0)^\top$和$(0.2,0.3,0)^\top$。前两维的和为$0.6,0.1$，都选正方向：第一维平均得$0.3$，第二维排除$-0.2$后只剩$0.3$。合并增量为$(0.3,0.3,0)^\top$，简单任务算术的和则为$(0.6,0.1,-0.01)^\top$。两者缩放都取一，加回底座分别得到$(1.3,1.3,1)^\top$与$(1.6,1.1,0.99)^\top$。跟着算完，重点\textbf{不在记住这些数字，而在看出两种方法“\emph{处理冲突}”的哲学差异}：\textbf{简单任务算术}让第二维的$+0.3$和$-0.2$\emph{部分抵消}成$0.1$，且忠实地保留了第三维那个微小的$-0.01$；\textbf{TIES}则\emph{先把每个任务最不重要的第三维整个剪掉}（两个增量第三维都置零），再在第二维上\textbf{排除掉与当选方向（正）冲突的}$-0.2$\textbf{、只保留同号的}$+0.3$（于是该维得$0.3$而非$0.1$）。两者最终得到\textbf{明显不同的参数}。但务必清醒：\textbf{这些算术只确定了“参数变成了什么”，它\emph{完全没有}提供“哪种结果在真实任务上风险更小”的证据}——要回答谁更好，仍须在各自任务的验证数据上实测。

TIES 是“\textbf{按重要性（幅度）\emph{确定性}地}剪枝”，DARE 则提供了另一种很不一样的稀疏化思路：\textbf{\emph{随机}地丢弃增量坐标}。它以概率$p$\textbf{随机把某些增量坐标“丢掉”（置零）}，再\textbf{对幸存下来的值做\emph{放大补偿}}；DARE 的结果\textbf{之后可以接任务算术、也可以接 TIES}（它是个可插拔的前置稀疏化步骤）。

为什么丢掉之后还要“放大”？因为随机置零会让\textbf{期望幅度系统性地变小}（只剩$1-p$的存活率）。为了让\textbf{“随机丢弃之后、该坐标的\emph{期望}值仍等于原增量”}（保持\emph{无偏}），就必须\textbf{用保留概率$1-p$去\emph{补偿}}。具体地，令$M_{k,j}\sim\operatorname{Bernoulli}(1-p)$（取1表示保留），且$0\leq p<1$：固定原增量，令$M_{k,j}\sim\operatorname{Bernoulli}(1-p)$，掩码为一时保留，且$0\leq p<1$，得到
\begin{equation}
  \begin{aligned}
    \widehat{\tau}_{k,j}&=\frac{M_{k,j}}{1-p}\tau_{k,j},\\
    \mathbb{E}_M[\widehat{\tau}_{k,j}]&=\tau_{k,j},&
    \operatorname{Var}_M(\widehat{\tau}_{k,j})
      &=\frac{p}{1-p}\tau_{k,j}^{2}.
  \end{aligned}
  \label{eq:ket-integration-dare}
\end{equation}
这两个统计量要\textbf{分开理解}：\textbf{“期望不变”}来自$\mathbb{E}[M_{k,j}]=1-p$，所以$\mathbb{E}[\widehat{\tau}_{k,j}]=\tau_{k,j}$（\emph{无偏}）；而\textbf{方差}由伯努利方差$p(1-p)$除以$(1-p)^2$推出为$\frac{p}{1-p}\tau_{k,j}^2$，它\textbf{衡量“某\emph{一次}随机抽样，结果可能偏离期望多远”}。

这里藏着 DARE 最需要\textbf{正确理解}的两点。其一，\textbf{“无偏”是“\emph{对所有随机掩码求平均}才成立”的性质}：对\textbf{固定输入}，\textbf{单层线性更新}确实能在期望上保持响应；但\textbf{一旦进入\emph{非线性网络}、或用\emph{任务损失}评价}，\textbf{“求期望”和“实际计算”不能随便交换}（非线性不穿过期望），所以\textbf{不能声称“任意一次丢弃后模型行为都不变”}。其二，注意方差项$\frac{p}{1-p}$：\textbf{丢弃率$p$\emph{越接近1}，这个因子\emph{反而越大}}——意味着\textbf{剪得越狠，\emph{单次}结果的波动（方差）越大}（幸存的少数坐标被放得极大）。所以 DARE 是否在实践中可接受，\textbf{取决于微调增量本身的\emph{尺度}、网络里的\emph{冗余程度}、以及丢弃率}三者的平衡；它有明确的经验适用范围，\textbf{绝不能当成“对\emph{任意}独立训练模型都成立”的保证}\scite{ket-yu2024dare}。

用前面那个$\bm{\tau}_1=(0.4,-0.2,0.01)^\top$、取$p=0.5$来\emph{具体感受“无偏并不等于“这次没事””}：假设这一次随机抽出的保留掩码是$(1,0,1)^\top$（第二维被丢、第一三维幸存），由于$1-p=0.5$、补偿因子为$2$，结果变成$(0.8,0,0.02)^\top$——\textbf{它和原增量$(0.4,-0.2,0.01)^\top$\emph{差别很大}}（第二维整个没了、第一维翻倍）。

这清楚说明：\textbf{“无偏”只描述“在\emph{无穷多种随机掩码}上取平均”的性质，\emph{绝不保证}你手上\emph{这一次}抽样后模型行为还保持不变}。实践中要靠\textbf{多次尝试、验证集筛选}来挑一个好掩码。最后把 TIES 和 DARE 的\textbf{角色差异}钉死：\textbf{TIES 是“按幅度和符号\emph{确定性}筛选+解决冲突”，DARE 是“\emph{随机}筛选 + 补偿期望”}；两者\textbf{可以串联使用}（先随机稀疏、再按符号合并），但它们\textbf{改变的是不同的运算、依赖不同的假设}，别当成同一种方法。

在结束“增量融合”前，要用前面“置换对称性”的洞见\textbf{纠正两个非常普遍的想当然}，它们都源于“把参数坐标当成了语义”。

\textbf{误区一：把“坐标上的符号冲突”直接当成“语义冲突”}。回想隐藏单元可以任意重排——\textbf{同一种功能，换一套单元排列，某个坐标上增量的\emph{正负号}就可能翻转}。所以两个任务在\emph{同一坐标}上一个$+$、一个$-$，\textbf{未必}表示它们在“语义上打架”，可能纯粹是\textbf{内部坐标对不齐}造成的假象（该先做的是对应，而不是判定冲突）。

\textbf{误区二：以为“加上负任务向量”就等于“精确忘掉某项技能”}。$\bm{\theta}_0-\lambda\bm{\tau}_k$\textbf{只是在参数空间里沿$\bm{\tau}_k$的\emph{反方向}走了一步}，它\textbf{几何上是个反向位移，\emph{语义上}并不等于“把某个指定任务的能力干净地删除”}——因为任务向量本就不是一项隔离的纯技能（前面已强调），而且能力在网络中是\emph{分布式、纠缠}的。真正意义上“\textbf{消除指定训练样本影响}”的\term{机器遗忘}，需要第~\ref{sec:ket-update-unlearning-validation}节的参照评估、以及\BookChapterRef{chap:trust-privacy}{卷3《基础、理论与可信性》第16章“隐私保护”}中讨论的\textbf{正式保证}，\textbf{绝不能由“向量取个负号”这种简单操作推出}。

前面融合增量时，我们大多是先得到完整的$\Delta\bm{W}_k$再相加，那会\textbf{破坏“低秩”结构}（多个秩$r$增量之和的秩可能涨到$Kr$）。如果我们\textbf{希望融合\emph{之后}的产物\emph{仍然是一个低秩模块}}（保持省参数的好处），就得\textbf{换个运算位置：直接在\emph{因子}$\bm{B},\bm{A}$上做加权，\emph{最后才相乘}，而\emph{不是}先把完整增量加起来}。

LoRAHub 正是这种做法，但它有一\textbf{串严格的适用前提}：各来源必须\textbf{基于同一底座、插在同一层、秩$r$和缩放约定都相同}。在此前提下，它\textbf{分别}$\bm{B}_k$加权求和、$\bm{A}_k$加权求和，再把两个加权结果相乘，得到一个\textbf{仍是秩$r$的新模块}：
\begin{equation}
  \begin{aligned}
    \bm{B}(\bm{w})&=\sum_k w_k\bm{B}_k,&
    \bm{A}(\bm{w})&=\sum_k w_k\bm{A}_k,\\
    \Delta\bm{W}(\bm{w})&=\bm{B}(\bm{w})\bm{A}(\bm{w}).
  \end{aligned}
  \label{eq:ket-integration-lorahub}
\end{equation}
式中$\bm{w}=(w_1,\ldots,w_{K_{\mathrm{s}}})^\top$是来源系数，\textbf{同样可以为负、无需和为一}（参数系数的老规矩）。\textbf{这里有个\emph{极其}关键、且名字最具迷惑性的问题}：虽然方法叫“加权”，但\textbf{“先把因子各自加权、再相乘”得到的结果，\emph{一般不等于}“把完整增量加权相加”}。要确认这件事，不能望文生义，必须\textbf{老老实实展开乘积}。两个来源（$K_{\mathrm{s}}=2$）时展开为：
\begin{equation}
  \begin{aligned}
    \Delta\bm{W}(\bm{w})
      ={}&w_1^2\bm{B}_1\bm{A}_1+w_2^2\bm{B}_2\bm{A}_2\\
       &+w_1w_2(\bm{B}_1\bm{A}_2+\bm{B}_2\bm{A}_1).
  \end{aligned}
  \label{eq:ket-integration-lorahub-cross-terms}
\end{equation}
展开后能看到和“完整增量加权”$\sum_k w_k\Delta\bm{W}_k=w_1\bm{B}_1\bm{A}_1+w_2\bm{B}_2\bm{A}_2$\textbf{的两点本质差异}：
\begin{itemize}\item \textbf{每个来源\emph{自己}的增量，乘的是系数的\emph{平方}}$w_k^2$，而不是$w_k$；\item 更要紧的是\textbf{凭空多出了\emph{交叉项}}$w_1w_2(\bm{B}_1\bm{A}_2+\bm{B}_2\bm{A}_1)$——它们\textbf{混用了“来源1的一个因子”和“来源2的另一个因子”}。\end{itemize}
所以\textbf{除非这些交叉项恰好为零，否则两种融合结果\emph{一般不相等}}。而交叉项\textbf{强依赖各来源因子的\emph{具体内部坐标}}——还记得前面$\bm{B}\bm{A}=(\bm{B}\bm{R})(\bm{R}^{-1}\bm{A})$的坐标自由度吗？\textbf{即便一个可逆变换能保持“单个增量$\Delta\bm{W}_k$不变”，它也\emph{可能改变}LoRAHub 融合的结果}（因为交叉项变了）。这再次说明：\textbf{因子表示本身也是融合算法的“输入条件”}，不是无关紧要的内部细节。（说明：原论文用$\bm{A}$命名层响应侧因子，本章为统一维度记号交换了字母，但乘法次序保持不变。）\scite{ket-huang2024lorahub}

\textbf{用最简单的标量（$1\times1$）因子，就能\emph{无可辩驳}地验证“结果依赖因子写法、而非只依赖完整增量”}。设两个来源最初都取$B_k=A_k=1$，于是\emph{各自的完整增量}都是$B_kA_k=1$，系数都取$1/2$。此时因子加权后$B=(1/2)(1)+(1/2)(1)=1$、$A=1$，组合结果$B\cdot A=1\times1=1$。

现在\textbf{只对第二个来源做一次“乘积不变的改写”}：把$B_2,A_2$从$(1,1)$改成$(2,1/2)$。注意它的\textbf{完整增量仍然是}$B_2A_2=2\times(1/2)=1$\textbf{，丝毫未变}。可重新做因子组合：加权后$B=(1/2)(1)+(1/2)(2)=1.5$、$A=(1/2)(1)+(1/2)(1/2)=0.75$，结果变成$1.5\times0.75=1.125\neq1$。

\textbf{同样两个完整增量（都是1），仅因为第二个来源的因子换了一种“乘积等价”的写法，LoRAHub 融合结果就从1变成了$1.125$}。这正是交叉项在作祟，也\textbf{实证}了上一条结论：\textbf{不能只凭“完整增量相同”就推断“因子融合结果相同”}——因子怎么分解，是会影响最终模型的真实条件。

再用一个\textbf{秩一（$2\times2$）}的例子，揭示“\textbf{在因子上融合}”带来的另一个结构性后果——\textbf{秩（表达能力）可能被压低}。取$\bm{B}_1=(1,0)^\top,\bm{A}_1=(1,0)$（则$\Delta\bm{W}_1=\operatorname{diag}(1,0)$）和$\bm{B}_2=(0,1)^\top,\bm{A}_2=(0,1)$（则$\Delta\bm{W}_2=\operatorname{diag}(0,1)$），权重都取$1/2$。
\begin{itemize}\item \textbf{完整增量平均}：$0.5\,\operatorname{diag}(1,0)+0.5\,\operatorname{diag}(0,1)=\operatorname{diag}(1/2,1/2)$，这是个\textbf{秩为2}的对角矩阵。\item \textbf{先平均因子再相乘}：平均$\bm{B}=(1/2,1/2)^\top$、$\bm{A}=(1/2,1/2)$，相乘得到一个\textbf{四个元素全是$1/4$的矩阵}，它\textbf{秩为1}。\end{itemize}
差异有两层：其一，\textbf{非对角线上冒出的$1/4$正是交叉项}（$\bm{B}_1\bm{A}_2$、$\bm{B}_2\bm{A}_1$的贡献）；其二，也是更深刻的，\textbf{结果秩从2被压到了1}——这说明\textbf{“坚持在秩$r$因子内做融合”会\emph{限制融合后模型可表达的秩}}，可能\textbf{不足以表达“完整增量之和”本可表达的内容}，这是保持低秩形态所付出的\textbf{表达力代价}。最后对照一下部署路径：LoRAHub\textbf{选定系数后只执行这一个新模块（或写回底座）}，是\textbf{一份模型}；而前面的\textbf{AdapterFusion 是\emph{保留多个适配器}、运行时组合各自激活}，要跑多条路径——\textbf{两者省/费的资源和运行方式都不同}，别因都“加权”就混淆。

\subsection{融合系数的确定}
\label{sec:ket-integration-coefficients}

无论用哪条融合规则，只要系数$\alpha,\lambda,\bm{w}$可变，就会\textbf{对应一\emph{整族}候选融合模型}，于是绕不开一个元问题：\textbf{系数取多少？必须\emph{另有依据}来选，不能拍脑袋}。不同依据的“含金量”差别很大，先排个序：
\begin{itemize}\item \textbf{预设等权（如全$1/K_{\mathrm{s}}$）}：\textbf{不消耗任何目标标签}，最省事，是合理的默认基线，\textbf{但不能因为效果不错就声称它“最优”}。\item \textbf{有标签材料}：可以\textbf{直接在上面比较真实任务损失}，依据最硬。\item \textbf{只有无标签材料}：没法算真实损失，只能用\textbf{熵（预测集不集中）、一致性}等\textbf{代理量}间接挑，依据较弱（下面会看到熵会骗人）。\end{itemize}
开始选之前，必须\textbf{先声明三件事}：\textbf{允许的系数\emph{范围}}（能否为负、多大）、\textbf{数据\emph{访问权限}}（能不能用目标标签、用多少无标签数据）、以及\textbf{构造候选的\emph{计算预算}}。最后是数据纪律：\textbf{测试标签一旦被你用来\emph{反复}调系数，它就\emph{已经参与了开发}，失去了“独立裁判”资格}，\textbf{不能再拿它当泛化证据}报成绩。

LoRAHub 给出了一种“\textbf{有少量示范时如何选系数}”的具体流程。在\textbf{新任务}上准备\textbf{少量“输入—示范答案”对}，然后\textbf{冻结底座和所有来源因子，\emph{只}把系数$\bm{w}$当变量去搜索}。每评估一组候选系数：先\textbf{用它构造出候选模块}，再让它在示范上做预测、计算\textbf{目标输出的负对数似然}（希望正确答案概率大），并\textbf{额外加一项对系数幅度的惩罚}（防止权重乱飙）。

由于目标对$\bm{w}$可能不好直接求导（涉及离散选择等），它用\textbf{无梯度优化器}：\textbf{根据一批候选各自的评分，决定下一批该试哪些系数}。原论文用名为 Shiwa 的策略选择器挑选搜索方法，多数设置采用\term{协方差矩阵自适应进化策略}（covariance matrix adaptation evolution strategy, CMA-ES）——直觉上就是\textbf{维护一个“系数采样分布”，哪片区域的候选评分好，就把采样分布往哪片挪、并收窄}，逐步逼近好区域。\textbf{到达预算后系数就\emph{固定}下来，预测时不再为每个样本重新搜索}。

别被“\textbf{无梯度}”这个词误导而以为很便宜：它\textbf{省掉的只是“搜索过程中的\emph{反向传播}”}，但\textbf{准备示范、成百上千次\emph{前向}评分、以及整个搜索循环的成本\emph{一样不少}}。

为了明确评分对象，令非空$S_u$为任务$u$的适应示例集，$\Omega$为预定系数范围，$p_{\bm{w}}$为底座加候选模块后的预测输出概率。取$\Omega$非空、闭且有界，并假设有限样本上的概率为正且随系数连续，目标就能取到最小值。用负对数概率使正确预测输出概率越大、损失越小，再以$\eta\geq0$限制系数幅度，理想搜索目标为
\begin{equation}
  \bm{w}_u^\star\in\argmin_{\bm{w}\in\Omega}
   \left\{-\frac{1}{|S_u|}
   \sum_{(\bm{x},y)\in S_u}\log p_{\bm{w}}(y\mid\bm{x})
   +\eta\lVert\bm{w}\rVert_1\right\}.
  \label{eq:ket-integration-lorahub-search}
\end{equation}
几个实现与边界细节：对\textbf{序列输出}，一句话的对数概率是\textbf{各词元条件对数概率之和}；而\textbf{具体实现里是否按长度归一化、正则项占多大比例}，都会改变数值，应\textbf{记录在案}。也要注意本式只是对“负对数似然 + 幅度惩罚”这个\textbf{目标的概括}，不声称逐行复刻某个库的实现。

更重要的是\textbf{管理预期}：\textbf{有限预算下的无梯度搜索，只能给你一个“试过的候选里较好的\emph{近似}解”}，\textbf{不保证}真的找到全局$\argmin$；而且这个搜索有两个现实困难——\textbf{来源越多，系数维数越高}（搜索空间指数级变大）；\textbf{示范样本太少时，极易“\emph{选择过拟合}”}（挑出一组只在这几例上好、换数据就垮的系数）。所以\textbf{最终任务评价必须与选择系数用的$S_u$严格隔离}。

如果\textbf{连少量示范标签都没有}，还能不能选系数？AdaMerging 给出一条“\textbf{无标签}”路线：用\textbf{允许访问的无标签输入}，去\textbf{最小化模型预测的\emph{熵}}——直觉是“\textbf{好模型应当\emph{笃定}，所以找一组让预测分布\emph{更集中（更不犹豫）}的系数}”。它还提供两种\textbf{系数粒度}：
\begin{itemize}\item \textbf{任务级}：给每个来源的\emph{整个}任务向量配\textbf{一个}系数$\lambda_k$（参数少、好优化）；\item \textbf{层级}：更细，给“来源$k$在第$\ell$层”的向量\textbf{单独}配系数$\lambda_{k,\ell}$（自由度大、潜力高但更难调、更易过拟合）。\end{itemize}
整个过程\textbf{底座和任务向量都固定，唯一在优化的变量就是这些系数}。用$\Omega_{\mathrm{A}}$记录实验允许的系数范围，一种带显式范围的目标写为
\begin{equation}
  \begin{aligned}
    \bm{\theta}_{\ell}(\bm{\lambda})
      &=\bm{\theta}_{0,\ell}+\sum_k\lambda_{k,\ell}\bm{\tau}_{k,\ell},\\
    \bm{\lambda}^{\star}
      &\in\argmin_{\bm{\lambda}\in\Omega_{\mathrm{A}}}\frac{1}{|U|}
        \sum_{\bm{x}\in U}H(p_{\bm{\lambda}}(\cdot\mid\bm{x})).
  \end{aligned}
  \label{eq:ket-integration-adamerging}
\end{equation}
$U$为非空无标签适应集，$H(p)=-\sum_y p(y)\log p(y)$为预测熵，$p_{\bm{\lambda}}$由候选模型参数和该任务接口给出。原研究从预设系数出发，用无标签测试样本的前向计算与熵梯度更新系数，按预算或停止规则结束。式中的$\Omega_{\mathrm{A}}$是为说明实验范围补入的记号，不把某一固定范围强加给所有实现。非空紧范围与连续概率足以保证最小点存在，但迭代仍不保证取得全局最小值；不具备存在条件时，只报告有限预算中得到的候选。任务级指每个来源向量共用一个系数，不是逐预测输入生成权重\scite{ket-yang2024adamerging}

\textbf{“用熵选系数”有个致命盲区，必须当成重点}：熵只衡量“\textbf{分布集\emph{不集中}}”，它\textbf{完全不评价“\emph{对不对}”}。一个\textbf{自信地\emph{全错}}的模型，熵可以极低！举个极端例子：如果模型\textbf{无论输入什么都笃定预测“月季”}，它的预测分布是$(1,0,\ldots)$，\textbf{熵接近零（看起来无比自信、完美）}，可\textbf{所有非月季照片全错}。

因此：\textbf{熵下降\textbf{绝不等于}真实风险下降}；用熵选出的系数，\textbf{最终仍须用真实标签评价}，并且要\textbf{专门检查模型是否“\emph{塌缩}成只会预测单一类别”}（这是最小化熵的一个退化解）。还有一个\textbf{方法论上必须如实声明}的点：如果你\textbf{直接拿“\emph{当前这批要预测的测试数据}”当作适应集$U$去调系数}，这属于\textbf{直推式（转导）预测 / 测试时适应}，\textbf{不能}包装成“对未来未知输入也有效”；\textbf{要声称泛化能力，就必须先在适应数据上\emph{冻结}系数，再在\emph{另外}的测试集上评测}。适应时\textbf{能否看到任务身份、迭代几轮、是否持续更新、何时重置}，这些条件都应\textbf{预先记录}。

当一次要融合\textbf{多个\emph{目标不同}的任务}时，系数搜索会撞上一个\textbf{没有“唯一标准答案”}的处境：\textbf{这些任务的要求可能\emph{无法同时}满足}（提升物种准确率的系数，可能拉低病害准确率）。这时用一个标量总分硬压平差异，往往掩盖了真实的取舍。

更诚实的语言是\textbf{多目标优化}。设有$m$个目标任务，$R_j(\bm{\theta})$是任务$j$在\emph{它自己}的数据分布和损失下的\textbf{平均风险}（越小越好），于是一个融合模型对应一个\textbf{风险向量}$(R_1(\bm{\theta}),\ldots,R_m(\bm{\theta}))$。在此基础上定义\term{帕累托最优}：\textbf{一个可行参数$\bm{\theta}$，如果\emph{不存在}另一个可行参数能“在\emph{所有}任务上都\emph{不差}、且\emph{至少一个}任务\emph{严格更好}”，就称它是帕累托最优的}。

这个概念有\textbf{三层必须说准}的含义：其一，它表示“\textbf{已经到了边界——想再改进任何一个任务，就\emph{不得不牺牲}别的任务}”；其二，它\textbf{不表示每个任务都达到了它\emph{各自单独训练}的最优}（帕累托最优点上各任务通常都做了妥协）；其三，它\textbf{只是一个“理想边界”的定义，\emph{更不}表示你的搜索算法\emph{真的找到了}这样的点}（算法可能停在远离边界的地方）。

Pareto Merging 把这个思想落成算法：它构造一族候选，其中\textbf{“共享参数”部分与偏好无关}（大家共用的底座式能力），再叠加\textbf{“与偏好\emph{相关}的低秩修正”}——于是\textbf{给定不同的偏好权重，就沿着帕累托前沿得到\emph{不同的折中}模型}（更看重物种、或更看重病害，各取一个）。在\textbf{没有额外样本}时，它用\textbf{参数空间距离}当代理目标；在\textbf{有无标签输入}时，用\textbf{分任务的熵}目标。但同样要守住那条边界：\textbf{这些\emph{代理}目标上的“帕累托性质”，\emph{不能直接等同}于“真实风险上的全局帕累托最优”}（距离近/熵低不等于真实错误最优）。

对花卉系统，\textbf{落地建议}很具体：与其报一个糊在一起的综合平均分，不如\textbf{把“物种错误率”和“病害漏检率”\emph{分别画出来}（甚至画成一条帕累托曲线）}，\textbf{按\emph{事先声明}的用途}（比如“漏检病害代价极高”）\textbf{在曲线上选一个点}。这样\textbf{折中是明明白白做出的}，而不是\textbf{被一个加权平均分悄悄掩盖了“病害漏检严重退化”}\scite{ket-chen2025pareto}。

最后把本章的“\emph{事后融合}”和另一处的“\emph{共同训练}”区分开，避免术语混用：\textbf{本节（Pareto Merging）的候选空间是\emph{被已有来源和既定融合规则框死}的}——你只能在这些现成模型的组合里挑折中；而第~\ref{sec:ket-sharing-directions}节讨论的是\textbf{从头\emph{共同训练}多个任务时，如何\emph{协调各任务梯度的\emph{下降方向}}（梯度冲突怎么处理）}，那里\textbf{参数可以自由地往任意方向学}，自由度大得多。两者\textbf{都在处理多目标}，但\textbf{能改动的参数范围、能用的训练证据、计算费用都不同}，\textbf{不能因为术语相近就互相借用对方的最优性保证}。

表~\ref{tab:ket-integration-merging-operations}汇总运算对象与样本访问条件。完整模型参数、任务增量与低秩因子是不同运算对象，基于样本统计的加权可作用于其中一些对象，对应调整又可先于多种规则。表的行因此允许交叉，不是一组互斥方法。第~\ref{sec:ket-integration-alignment}~节决定哪些坐标相容，不能由表中一种算术直接跳过对应检查。

来源存储、预测输入统计、校准材料以及反复构造候选和前向评分都属于融合成本。多轮吸收来源时，裁剪、匹配或重新选系数可能使顺序影响结果；应保存准确底座、来源、规则与系数，逐任务回归检查，并按有效维护要求选择回退版本。未来可塑性需要同预算的新任务学习实验，未见组合能力需要成分与关系分开的测试，都不是单份模型参数自动具有的性质。

Tam等把来源任务与未见组合任务分开，在图像分类、图像生成和跨语言任务的共同协议中比较融合方法。其预印本报告，来源任务成绩与未见任务成绩的关系随应用而变，增加来源时还可能出现来源任务保持下降、未见任务泛化上升的相反趋势。这项结果支持分开评价两者，不支持跨架构的固定排名\scite{ket-tam2024-realistic-merging}实际对照应包括底座、相关单来源、可执行的预测集成，以及样本和预算允许时的联合训练；来源能力保持也不应代替组合泛化测试。

\begin{table*}[htbp]
  \centering
  \small
  \setlength{\tabcolsep}{4pt}
  \renewcommand{\arraystretch}{1.2}
  \begin{tabularx}{\linewidth}{@{}>{\raggedright\arraybackslash}p{24mm}
      >{\raggedright\arraybackslash}X
      >{\raggedright\arraybackslash}X
      >{\raggedright\arraybackslash}X@{}}
    \toprule
    操作 & 运算对象与可执行条件 & 系数所用样本或统计量 & 部署与主要风险 \\
    \midrule
    完整模型参数平均 & 全部权重或共享层；要求层、坐标和预测输出含义对应 & 预设系数，或在独立开发材料上选择、学习 & 通常单模型，也可多头；平均模型可能落入高损失区域 \\
    \addlinespace
    基于样本统计的加权 & 模型参数坐标或线性层；Fisher merging用局部曲率，RegMean用层自变量二阶矩 & 来源统计或生成统计的预测输入；仍可另选训练配置 & 单模型；统计失配、对角近似和病态矩阵限制结论 \\
    \addlinespace
    任务增量融合 & 相对同一准确底座的位移；可接裁剪、符号选择或随机稀疏 & 固定缩放、标签选择或无标签适应 & 底座加增量；任务冲突、过大位移或随机波动可能损害行为 \\
    \addlinespace
    低秩因子组合 & 同底座、插入层、秩与缩放约定；分别加权因子再相乘 & LoRAHub用少量任务示例作无梯度搜索 & 一个模块或写回底座；交叉项、因子坐标与样本过拟合影响结果 \\
    \midrule
    对应调整（前置步骤） & 置换、特征匹配或共同子空间；结构异构时仍可能无法运算 & 权重匹配只读模型参数，激活匹配另需预测输入 & 可后接融合或保留分支；对应成立不保证任务能力兼容 \\
    \bottomrule
  \end{tabularx}
  \caption{模型参数融合的运算对象、所需样本或统计量和部署结果。上部各行允许交叉，下部的对应调整可以先于多种规则执行。无标签只限制标签访问，不表示没有预测输入、任务身份或适应计算。比较来源保留、未见组合和资源节省时，均须以最终实际模型而非操作名称为准。}
  \label{tab:ket-integration-merging-operations}
\end{table*}

\section{基于知识蒸馏的融合}
\label{sec:ket-integration-distillation}

参数融合有个\textbf{硬门槛：参数必须能对应}。可如果\textbf{两个花卉分类器的\emph{结构根本不同}}（一个是卷积网络、一个是 Transformer，层数、单元数都不一样），那它们的参数\textbf{在数学上就\emph{无法}逐坐标对应}，前面的 Model Soups、任务向量全都用不了。但请注意：\textbf{参数对不上，\emph{不代表能力无法转移}}——因为它们\textbf{对同一张照片，仍然可以用\emph{同一套类别顺序}给出预测}。

于是有了\textbf{最后一条、也是适用面最广的}融合路线：\term{知识蒸馏}（knowledge distillation）。它\textbf{完全绕开“参数对应”}，做法是——\textbf{去\emph{查询}来源模型、把它的\emph{回答（输出/内部表示）}当成\emph{训练目标}}，然后\textbf{从头训练一个\emph{新}模型去\emph{模仿}这些行为}。提供监督的来源叫\term{教师}（teacher），接受监督、被训练的新模型叫\term{学生}（student）。

厘清几个常见误解：\textbf{教师不必是单个模型}，它\textbf{可以是前面讲过的组合系统、甚至融合后的模型}（先集成、再蒸馏）；\textbf{学生也\emph{不一定}更小}——它可以更小（为了压缩部署），也可以\textbf{换成另一种结构}（为了换平台）。决定蒸馏\textbf{可不可行}的，是\textbf{“监督接口能拿到什么、查询输入覆盖全不全、学生容量\emph{能不能表达}教师的行为”}这三件事，而\textbf{不是“教师和学生谁大谁小”这种表面关系}。

\subsection{教师与蒸馏输入的准备}
\label{sec:ket-integration-teachers-inputs}

蒸馏的\textbf{第一步不是急着训学生，而是“盘点教师到底\emph{吐得出什么}”}——因为\textbf{监督信号的\emph{类型}，完全取决于教师接口的开放程度}，直接决定学生能学走什么：
\begin{itemize}\item 接口\textbf{只返回一个类别}（硬标签）：学生\textbf{只能拿到“最终答案”}，学不到教师的犹豫。\item 接口\textbf{返回概率或分数}：学生\textbf{才拿得到“软类别分布”}（教师觉得次像什么），信息量大得多。\item 接口\textbf{连内部激活都能读}：才\textbf{有可能}设置“模仿中间表示”的约束（见后文FitNets）。\end{itemize}
多教师时还要先\textbf{统一口径}：若两位教师都答\emph{物种}，应按第~\ref{sec:ket-integration-prediction}节\textbf{统一类别顺序、概率尺度、聚合规则}后再当监督；若两位\textbf{分别答物种和病害}，则必须\textbf{保留任务条件（告诉学生当前问哪个）、或用\emph{各自独立的监督头}}，\textbf{绝不能把两个不同问题的答案\emph{硬拼成一个向量}让学生去拟合}（语义对不齐，回到6.3）。

别以为“多找几个教师”就一定更可靠：教师之间\textbf{同样可能\emph{一起犯错}（共享偏差）或\emph{答案冲突}}（呼应6.2的相关性）。所以\textbf{蒸馏前必须先\emph{明确“学生到底在模仿谁”}}：是模仿\textbf{等权集成}的输出、\textbf{条件集成}（按输入选教师）、\textbf{任务专用的单个教师}、还是\textbf{融合后的那一个模型}？这个“\emph{模仿对象（聚合教师）}”要先\textbf{固定}下来。

评测时要\textbf{把收益来源\emph{拆开}归因}，尤其在“先集成、再蒸馏”时，应\textbf{同时报告三档成绩}：\textbf{单个教师}、\textbf{教师集成}、\textbf{蒸馏出的学生}。这样才能区分“\textbf{提升到底来自‘多教师\emph{聚合}’这一步，还是来自‘训练学生’这一步}”——学生超过单教师，可能主要是集成的功劳而非蒸馏魔法。第~\ref{chap:ket-reuse-selection}章那套“可靠性、收益”判断在这里\textbf{照样适用}：\textbf{教师数量多，\emph{不等于}独立判断依据等比例增多}（它们可能是同一错误的副本）。

教师不像静态教材——\textbf{它\emph{只在你拿输入去\emph{查询}它的那一刻}才提供监督}。这带来一个\textbf{蒸馏特有的、且经常被忽视的要素：\emph{查询覆盖}}，即“你到底在\emph{哪些输入}上问过教师”。\textbf{学生能学到的行为，\emph{不可能超出}你查询教师的输入范围}，而不同来源的查询数据各有陷阱：
\begin{itemize}\item 用\textbf{教师原来的训练数据}查询：覆盖的是\textbf{教师最熟悉的区域}，但教师在这些点上可能只是\textbf{“\emph{背答案}”（记忆性响应）}，学生学到的是记忆而非泛化。\item \textbf{另采无标签新数据}查询：能\textbf{补充新的条件/区域}，但可能\textbf{冲出教师\emph{可靠}的范围}（教师在陌生分布上乱答，学生还一本正经地学）。\item 用\textbf{合成数据}查询：能\textbf{补足现实中难采样的场景}，但\textbf{生成器的偏差、重复模板}会\textbf{限制覆盖的多样性}。\end{itemize}
一个极端但极有说服力的例子：\textbf{如果你只在月季、百合的照片上查询教师，\emph{从来没有}兰花的输入}，那么\textbf{即使蒸馏损失\emph{降到零}，也\emph{丝毫没有}约束学生在兰花上会怎么答}——那里是监督的真空地带。所以\textbf{“在查询区域拟合得完美”\textbf{不能证明}“教师的\emph{全部}能力都被保持”}；要看能力保持，查询分布\textbf{必须覆盖你关心的全部区域}。

蒸馏还允许一种\textbf{师生“信息\emph{不对等}”}的巧妙设定：\textbf{训练教师时，可以让它看到\emph{学生在部署时\emph{看不到}}的额外观测}。这叫\term{广义蒸馏}（generalized distillation）。典型例子：\textbf{训练时}教师\textbf{同时读到“照片 + 实验室检测报告”}（信息很全），而\textbf{部署时学生只能拿到照片}（报告在真实使用中来不及或成本太高）；我们让\textbf{“见过报告的教师”为这些“(照片, 报告)配对”生成监督，再教给“只看照片的学生”}，希望学生学会\textbf{仅从照片里就推断出那些“通常需要报告才能知道”的线索}\scite{ket-lopezpaz2016generalized}。

使用它有两个\textbf{严格的对应/边界}要求。其一，\textbf{配对必须\emph{严格对应同一植株、同一观测时刻}}——照片和报告一旦\textbf{错配}（张三只的报告配给李四的照片），教师给的就是\textbf{另一个对象/另一个状态}的判断，监督直接变成噪声。其二，也是能力上限：\textbf{如果某些差异\emph{在照片里\emph{根本没有}任何痕迹}、\emph{只}存在于报告中}（比如两种病原在照片上完全无法区分，只有化验能分辨），那么\textbf{学生\emph{无论如何}都不可能“对每一个例子”恢复这种差异}——\textbf{信息不在它的输入里，再强的教师也变不出来}。下一条公式会精确说明，这种情况下学生\textbf{最好能做到的，只是学到“对不可见报告的\emph{平均}判断”}。

为明确这种额外观测的作用边界，固定学生照片$\bm{x}$，将教师额外观测记为$\bm{x}^{+}$，教师给类别$y$的软概率为$q_y(\bm{x},\bm{x}^{+})$。学生只能给同一$\bm{x}$分配一个概率向量。若允许选择任意概率分布，最小化额外观测条件下的期望交叉熵，最优预测是
\begin{equation}
  p^\star(y\mid\bm{x})
   =\mathbb{E}\!\left[q_y(\bm{x},\bm{x}^{+})\mid\bm{x}\right].
  \label{eq:ket-integration-privileged-limit}
\end{equation}
这个理想解值得\textbf{逐字读懂}，它精确刻画了“广义蒸馏\emph{最多}能给学生什么”：
$$p^\star(y\mid\bm{x})=\mathbb{E}\bigl[q_y(\bm{x},\bm{x}^+)\mid\bm{x}\bigr].$$
意思是，面对一张只看得见$\bm{x}$的学生，它\textbf{最优的预测，是“把所有可能的隐藏报告$\bm{x}^+$上、教师的判断\emph{平均}起来”}。为什么是平均？推导上，最优目标可写成$-\sum_y\mathbb{E}[q_y\mid\bm{x}]\log p(y\mid\bm{x})$，它与“条件平均教师分布的熵”两者相减后，\textbf{剩下的正是一个\emph{非负}的 KL 散度}；KL 散度\textbf{当且仅当两个分布\emph{相等}时取到最小值零}，所以最优解就是让学生分布\textbf{等于这个“平均后的教师分布”}（零概率项按通常的极限约定$0\log0=0$处理）。

这个结论\textbf{既令人鼓舞、又划定了天花板}：鼓舞在于，学生\textbf{确实能吸收“隐藏报告\emph{统计上}带来的判断倾向”}（比如某类照片纹理总体上更对应某病原）；天花板在于，\textbf{它学到的是“对看不见的报告的\emph{平均}反应”，\textbf{不可能对每一个具体例子}\emph{逐例恢复}报告里的私有信息}——同一张照片对应两份不同报告时，学生只能给同一个平均答案，无法分身。更何况，\textbf{学生容量有限、训练样本有限}，实践中连这个理想平均分布\textbf{都未必能达到}。

\subsection{蒸馏约束的构造}
\label{sec:ket-integration-distillation-losses}

明确了“能从教师拿到什么、在哪些输入上拿”，下一个选择是\textbf{“\emph{让学生匹配教师的哪一部分}行为”}。大方向有两层，由表及里：
\begin{itemize}\item \term{预测输出约束}（output constraint）：只比较学生和教师的\textbf{最终答案/输出分布}（“你答得像不像我”）。这是最外层、接口要求最低的约束。\item \term{表示约束}（representation constraint）：进一步比较两者的\textbf{内部计算（中间层激活）}（“你\emph{思考过程中的内部特征}像不像我”）。\end{itemize}
这两类\textbf{可以\emph{叠加}使用}（既学输出、又学中间层）。除此之外，还能从这些基本量\textbf{派生}出更丰富的监督，比如“\textbf{样本之间的关系矩阵}”（教师认为哪些照片彼此相似，学生也应保持这种相似结构）、“位置注意力图”、甚至“输出对输入的\emph{导数}”；但每一种派生信号都必须\textbf{额外定义清楚“\emph{比较的对象是什么}”以及“\emph{需要什么访问权限}”}。本节\textbf{集中把“输出”和“表示”这两种\emph{最基本}的约束讲透}，而\textbf{不}把各种派生信号当成与它们并列、互斥的“第三类载体”——它们本质上仍是输出或表示的变换。

\subsubsection{基于输出行为的蒸馏}
\label{sec:ket-integration-output-distillation}

先理解“\textbf{为什么\emph{软}标签（概率）通常比\emph{硬}标签更值得学}”，这是 Hinton 蒸馏的核心直觉。当教师把一张照片\textbf{判成月季}时，它\textbf{内心里可能还有细腻的差别\emph{倾向}}：“它\textbf{挺像百合、但\emph{很不像}菊花}”。这条“月季$>$百合$>$菊花”的\textbf{次级相似性结构}，蕴含着教师对“类别之间如何关联”的丰富认知。

可是，一旦把答案压成\textbf{硬类别}$(1,0,0)$，这些次要倾向\textbf{被\emph{全部抹掉}}——百合和菊花都变成了一模一样的$0$，学生无从知道“错也有错得远近之分”。\textbf{概率（软）目标则\emph{保留}了这些暗知识}。问题来了：教师通常很自信，正确类概率接近1、其余类概率小到梯度几乎可忽略，\textbf{软信息虽在、却太微弱，学生学不到}。经典蒸馏\scite{hinton2015distilling}的解决办法极其简洁：引入一个\term{温度}（temperature）$T$去\textbf{“软化”softmax，\emph{压平}概率分布}，让那些非最大类别\textbf{获得更明显、可被学习的监督贡献}。

固定同一照片及$c$个类别顺序，教师和学生分数为$\bm{a}^{\mathrm{T}},\bm{a}^{\mathrm{S}}\in\mathbb{R}^{c}$，上标$\mathrm{T},\mathrm{S}$标记角色。将分数除以共同温度$T>0$再归一化，得到
\begin{equation}
  q_j^{(T)}=\frac{\exp(a_j^{\mathrm{T}}/T)}
                    {\sum_{v=1}^{c}\exp(a_v^{\mathrm{T}}/T)},\qquad
  p_j^{(T)}=\frac{\exp(a_j^{\mathrm{S}}/T)}
                    {\sum_{v=1}^{c}\exp(a_v^{\mathrm{S}}/T)}.
  \label{eq:ket-integration-temperature}
\end{equation}
先排雷：式子里的上标$(T)$是\textbf{“在温度$T$下算的”}这个标记，\textbf{千万\emph{不要}和训练里的“第$T$次迭代”搞混}，两者毫无关系。温度的行为很直观：\textbf{$T=1$就是普通softmax}；当分数不全相同时，\textbf{$T$越大，分布越\emph{平坦、均匀}}（各类概率被拉得越接近），\textbf{$T\to\infty$时趋近均匀分布}；反之$T<1$会让分布更尖锐。

训练目标上，为了让学生分布$\bm{p}^{(T)}$\textbf{逼近冻结的教师分布}$\bm{q}^{(T)}$，最直接的做法是\textbf{最小化“以教师概率为权重的、学生负对数概率”}，这正是\textbf{交叉熵}：
\begin{equation}
  L_{\mathrm{out}}
  =-\sum_j q_j^{(T)}\log p_j^{(T)},\qquad
  D_{\mathrm{KL}}(\bm{q}^{(T)}\Vert\bm{p}^{(T)})
  =L_{\mathrm{out}}-H(\bm{q}^{(T)}).
  \label{eq:ket-integration-output-loss}
\end{equation}
这里要把\textbf{交叉熵和 KL 散度的关系、以及 KL 的“方向”}讲清楚。KL 散度定义为$D_{\mathrm{KL}}(\bm{q}\Vert\bm{p})=\sum_jq_j\log(q_j/p_j)$，它\textbf{衡量“从教师分布$\bm{q}$到学生分布$\bm{p}$”的偏离}。\textbf{方向在蒸馏里被\emph{固定}为“教师在前（$\bm{q}\Vert\bm{p}$）”}——这不是随意的：因为\textbf{教师$\bm{q}$是\emph{冻结}的目标、只有学生$\bm{p}$在更新}，我们要优化的是学生。KL\textbf{非负、且两分布相等时为零}，天然适合当“差距”。

为什么常用交叉熵而非直接写 KL？因为\textbf{交叉熵 = KL散度 + 教师分布的熵}，而\textbf{教师是冻结的，它的熵对学生参数而言是个\emph{常数}}，\textbf{求梯度时常数项消失}——所以\textbf{最小化交叉熵和最小化 KL，对学生产生\emph{完全相同}的梯度}，用交叉熵算起来更方便。两个工程细节：其一，\textbf{多教师时“先平均\emph{分数}再软化”和“各自软化成概率后再平均\emph{概率}”，结果\emph{通常不同}}（softmax 非线性，回到6.4.4），要明确选哪种；其二，\textbf{如果接口只给了你\emph{概率}、没给原始分数}，你就\textbf{只能拿这个分布当目标}，\textbf{不能假装自己拿到了 logits}去做“先平均分数”。

一个实际常遇到的问题：如果教师接口\textbf{只吐给你“某个温度$T_0$下的\emph{完整、且严格为正}的概率”}（没给原始分数），我\textbf{还能不能换一个温度}$T$去软化它）?答案是\textbf{可以\emph{恢复}（在严格正的前提下）}。由$q_j^{(T_0)}\propto\exp(a_j/T_0)$可知分数和概率对数挂钩，把每一项取$T_0/T$次幂再归一化即可：由$q_j^{(T_0)}\propto\exp(a_j/T_0)$可知，将每项取$T_0/T$次幂后再归一化，就得到$q_j^{(T)}=(q_j^{(T_0)})^{T_0/T}/\sum_v(q_v^{(T_0)})^{T_0/T}$，共同分数常数在归一化中消去。截断概率、未知温度或只返回最大类别的接口不满足这项恢复条件。

\Needspace{10\baselineskip}
\begin{example}[三个类别携带的信息]
  \label{ex:ket-integration-soft-labels}
  一张教学照片的教师分数按月季、百合、菊花排列为$(2,1,0)$。用上述归一化分别计算温度一与温度二，得到
  \[
    \begin{aligned}
      T=1:\quad &(0.6652,0.2447,0.0900),\\
      T=2:\quad &(0.5065,0.3072,0.1863).
    \end{aligned}
  \]
  这个例子有\textbf{三个值得驻足的点}。其一，\textbf{硬目标}$(1,0,0)$\textbf{只说“是月季”}；而\textbf{软目标额外\emph{保留了“百合$>$菊花”这条排序}}——学生被要求不仅选对月季，还要\textbf{觉得它更像百合而非菊花}，这是硬标签给不了的“暗知识”。其二，软目标\textbf{能区分两个“\emph{硬标签完全相同}”的教师}：另一个教师若给分数$(2,0,1)$，它\textbf{最大类\emph{仍是}月季}（硬标签都是“月季”），但\textbf{次要类别\emph{反序}}（这次菊花$>$百合），于是\textbf{两者提供的软监督\emph{截然不同}}——只看硬标签会把它们当成同一个老师。其三，要守住批判态度：\textbf{升温只是\emph{放大}了次要类别的贡献、让这些相对关系更可学}，它\textbf{\emph{并不}保证这些教师的次要倾向\emph{本身符合真值}}——教师也可能以一种\emph{自洽但错误}的方式混淆百合与菊花，\textbf{软标签忠实地传递了教师，也可能\emph{连同偏见一起}传递}。
\end{example}

温度除了压平分布，还\textbf{悄悄改变了训练步的“\emph{尺度}（步长大小）”}，如果不处理，软目标项和真实标签项在总损失里的\textbf{相对权重会随温度失真}。为看清这一点，我们对\textbf{学生的分数}求交叉熵的梯度（这是少数值得亲手推一遍的求导）：
\[
  \frac{\partial\log p_j^{(T)}}{\partial a_v^{\mathrm{S}}}
    =\frac{\mathbf{1}\{j=v\}-p_v^{(T)}}{T}.
\]
将该导数代入交叉熵，并利用教师概率之和为一，得到
\begin{equation}
  \frac{\partial L_{\mathrm{out}}}{\partial a_v^{\mathrm{S}}}
   =-\frac{1}{T}\sum_jq_j^{(T)}
      \bigl(\mathbf{1}\{j=v\}-p_v^{(T)}\bigr)
   =\frac{p_v^{(T)}-q_v^{(T)}}{T}.
  \label{eq:ket-integration-temperature-gradient}
\end{equation}
梯度已经显式含一个$1/T$，但事情没完——\textbf{概率差$(p_v^{(T)}-q_v^{(T)})$本身也会随温度升高而\emph{缩小}}（高温下两个分布都趋近均匀、彼此接近）。两个效应要\textbf{合起来}看。为了看清高温下的净尺度，记$\bar a^{\mathrm{S}},\bar a^{\mathrm{T}}$分别为学生、教师各类别分数的均值；当\textbf{温度远大于“中心化后分数的幅度”}（即分数差异相对于$T$很小）时，对指数做一阶泰勒展开再归一化，可得到渐近式：
\begin{equation}
  \begin{aligned}
    p_v^{(T)}
      &=\frac{1}{c}+\frac{a_v^{\mathrm{S}}-\bar a^{\mathrm{S}}}{cT}
          +O(T^{-2}),\\
    \frac{\partial L_{\mathrm{out}}}{\partial a_v^{\mathrm{S}}}
      &=\frac{(a_v^{\mathrm{S}}-\bar a^{\mathrm{S}})
              -(a_v^{\mathrm{T}}-\bar a^{\mathrm{T}})}
              {cT^2}+O(T^{-3}).
  \end{aligned}
  \label{eq:ket-integration-temperature-square}
\end{equation}
读这个渐近结果：学生、教师概率\textbf{各自偏离均匀分布的程度约为$1/T$}，再\textbf{乘上导数里那个$1/T$}，于是\textbf{高温时整条梯度\textbf{大致按$1/T^2$的速度\emph{缩小}}}。这意味着：你\textbf{温度开得越高，软目标项对参数的实际推动力\textbf{反而以平方级\emph{变弱}}}——若不加补偿，软目标会在总损失里\textbf{被真实标签项“淹没”}，形同虚设。

这正是经典实现里\textbf{为什么要给软目标项乘一个$T^2$}（即用$T^2L_{\mathrm{out}}$或$T^2D_{\mathrm{KL}}$）的原因：\textbf{用$T^2$恰好\emph{大致抵消}那个$1/T^2$的尺度衰减}，让\textbf{软目标项的梯度量级，与“温度为1的\emph{真实标签}项”保持\emph{大致可比}}，这样你设的损失权重$\lambda_{\mathrm{D}},\lambda_{\mathrm{L}}$才有稳定的意义。措辞要准确：它是一个\textbf{高温下的\emph{渐近}尺度补偿}，\textbf{并不是“在任意温度下两种写法梯度\emph{精确}等价”}。例~\ref{ex:ket-integration-soft-labels}~中学生初始分数全零、$T=2$时，月季梯度约为$(1/3-0.5065)/2=-0.0866$，乘温度平方后约为$-0.3463$。梯度下降都会提高月季分数，但步幅不同。

蒸馏的“输出”\textbf{不限于分类概率}，在强化学习里可以是\textbf{动作}。Policy Distillation 比较了三种把教师策略传给学生的形式：\textbf{模仿硬动作}（教师选哪个动作就学哪个）、\textbf{回归动作价值}、以及\textbf{匹配动作\emph{分布}}。这里的\term{动作价值}（action value）描述“在当前状态选这个动作后，预期能拿到的\emph{折扣累计回报}”，把各动作价值过一次 softmax 就能形成一个“动作分布”当教师目标。

迁移到动作域有\textbf{特殊的对齐要求}，不能照搬分类：\textbf{动作编号、当前\emph{可用}的动作集合、以及\emph{状态观测}都必须严格对应}——分类里“编号3”在两个模型间通常稳定指代同一类，但动作的编号\textbf{极易因实现不同而错位}，\textbf{绝不能把“编号相同但含义不同的动作”当成同一个监督信号}。还有一个\textbf{方向\emph{相反}的温度用法}值得注意：该研究在\textbf{策略蒸馏里反而倾向用\emph{低}温}，目的是\textbf{让分布更尖锐、突出那几个\emph{高价值}动作}（动作选择往往要果断）；所以\textbf{分类蒸馏里“升温软化”的经验\textbf{不能机械地套到}策略蒸馏上}——温度高低\textbf{服务于目标}，没有统一方向\scite{ket-rusu2015policy}。

策略蒸馏还带来一个\textbf{分类任务里没有的、动态的覆盖难题}：\textbf{动作会\emph{改变环境}，从而改变接下来看到的状态}。学生\textbf{一旦选了和教师\emph{不同}的动作，就可能\emph{走入教师轨迹\emph{从未到过}的状态}}；而教师的\textbf{演示轨迹\emph{不一定}覆盖学生实际执行时会抵达的所有状态}。这是一个\textbf{闭环、自举}的覆盖问题，\textbf{不能用“分类输入覆盖”那种静态眼光想当然}，必须\textbf{在“学生\emph{真正交互执行}”的设定下，专门检查它访问到的状态分布有没有被教师监督覆盖}。

回到分类也有对应的“条件”纪律：\textbf{同一张照片，问“物种”和问“病害”应有\emph{不同}答案}，如果\textbf{不给学生任务标识}，就等于用两个互相冲突的答案去监督同一个输入，冲突依然无解（教师准备时讲过的任务条件要一路带到学生输入）。最后用一句话给整个蒸馏去魅，它适用于所有变体：\textbf{蒸馏损失下降，\emph{只}证明“学生\emph{接近了}你\emph{所选}的那个教师信号”}；至于\textbf{真实任务上的准确率，\emph{仍然}必须用独立标签去检验}——模仿得惟妙惟肖，不等于模仿的对象本身正确、或模仿对了任务。

\subsubsection{基于中间表示的蒸馏}
\label{sec:ket-integration-feature-distillation}

只在\emph{最终输出}上约束学生，有个空白：它\textbf{没有告诉学生“\emph{中间}各层应当形成\emph{怎样的内部表示}”}。结果可能是学生\textbf{最终答案蒙对了，但内部绕了一条和教师很不一样、且不利于继续学习（比如不利于加深网络）的弯路}。FitNets 的想法是给学生一个\term{提示}（hint）：\textbf{挑教师\emph{某个中间隐藏层}的表示，让学生\emph{某个中间层}直接去\emph{预测/逼近}这个教师表示}，相当于在网络\textbf{中段也设一个“\emph{对齐点}”}，而不只是两头对齐。

立刻有个工程障碍：\textbf{教师那层通常更\emph{宽}（维度大）、学生那层更\emph{窄}（维度小），两者维度\emph{对不上}，没法直接相减}。FitNets 的处理是\textbf{先在学生一侧\emph{加一个可学习的映射}（通常是一层线性变换），把学生的窄表示\textbf{升维/投影到与教师\emph{相容}的空间}，再在那个空间里算距离}\scite{ket-romero2015fitnets}。注意这个映射\textbf{是学生侧的“适配垫片”}，只为了让比较成为可能。

教师第$\ell_{\mathrm{T}}$层表示为$\bm{z}^{\mathrm{T}}\in\mathbb{R}^{d_{\mathrm{T}}}$，学生第$\ell_{\mathrm{S}}$层为$\bm{z}^{\mathrm{S}}\in\mathbb{R}^{d_{\mathrm{S}}}$，映射$r_{\bm{\phi}}:\mathbb{R}^{d_{\mathrm{S}}}\to\mathbb{R}^{d_{\mathrm{T}}}$。在固定预测输入上，平方残差惩罚各对应坐标的偏离，提示损失为
\begin{equation}
  L_{\mathrm{rep}}(\bm{x})
   =\frac{1}{2}\left\|
      r_{\bm{\phi}}(\bm{z}^{\mathrm{S}}(\bm{x}))
        -\bm{z}^{\mathrm{T}}(\bm{x})
     \right\|_2^2.
  \label{eq:ket-integration-feature-loss}
\end{equation}
读这个损失的几个要点：\textbf{参与学习（被优化）的是“映射参数$\bm{\phi}$”和“生成学生表示的那些模型参数”}，而\textbf{教师表示$\bm{z}^{\mathrm{T}}$是\emph{固定}的目标}（教师冻住，不被学生牵动）。前面的系数$1/2$纯属\textbf{数学便利}——平方项求导会产生因子2，乘个$1/2$正好抵消，\textbf{它\emph{不改变}最小点的位置}（有没有$1/2$，最优解都一样），只影响步长尺度。

推广到卷积网络时，“对应”要求更细致：不能只比向量整体，还要\textbf{对齐\emph{空间位置}（哪个像素对哪个像素）、\emph{通道}、以及考虑\emph{感受野}是否匹配}。虽然可学习映射能\textbf{调整形状}，但如果映射中用到了\textbf{插值或聚合}（比如把大特征图缩小），\textbf{由此\emph{丢掉}的空间细节必须\emph{明确说明}}，不能假装信息无损。最后提醒：\textbf{距离的数值\emph{强依赖}各维的\emph{尺度}}（某一层数值范围是$[0,1]$、另一层是$[0,100]$，距离会被后者主导），往往需要\textbf{归一化}；\textbf{绝不能只因为“两个向量\emph{等长}”就直接比较}——等长只说明维度凑巧相同，不说明坐标语义对应。

用三维教师提示与二维学生表示计算一次教学更新：$\bm{z}^{\mathrm{T}}=(1,0,2)^\top$、$\bm{z}^{\mathrm{S}}=(0.5,1)^\top$，取线性映射
\[
  r_{\bm{\phi}}(\bm{z}^{\mathrm{S}})=\bm{R}\bm{z}^{\mathrm{S}},
  \qquad
  \bm{R}=\begin{bmatrix}2&0\\0&1\\0&2\end{bmatrix}.
\]
同一预测输入分别经过教师和学生，再经映射得到$(1,1,2)^\top$，残差为$(0,1,0)^\top$，故损失为$1/2$。对$\bm{R}$求导得到残差与学生表示的外积，对学生表示求导得到$\bm{R}^{\top}$乘残差；提示梯度因此只进入映射和学生此前各层。暂时固定学生、以学习率$0.1$更新映射，第二行由$(0,1)$变成$(-0.05,0.9)$，映射后的第二坐标为$0.875$，损失为$0.3828125$。这展示了怎样学习维度对应，不是把短向量补零。

FitNets 的\textbf{原始\emph{训练流程}分两个\emph{先后}阶段}，不要和后来的常见变体搞混：\textbf{第一阶段（提示阶段）先\emph{只}训练“学生到提示层为止的部分 + 那个适配映射”}，让学生中段先学会模仿教师中段；\textbf{第二阶段再\emph{丢掉这层中段约束}、用常规的\emph{预测输出蒸馏}（软标签）去训练\emph{完整}学生}。\textbf{“把表示损失和输出损失\emph{从头就\emph{同时}优化}”是\emph{另一种}可选安排}，效果可能不错，但\textbf{那\emph{不是}FitNets 原论文的两阶段做法}，引用时要分清。

这里还有一个\textbf{影响部署成本}的关键认知：那个适配映射$r_{\bm{\phi}}$，\textbf{如果它\emph{纯粹}是“为了训练时能比较维度而搭的脚手架”}（只出现在损失计算里），那么\textbf{训练一结束、部署时\emph{完全可以把它删掉}}，最终预测\textbf{仍然只走学生\emph{自己}的预测输出路径}（学生窄表示直接接它自己的后续层），\textbf{不增加}部署开销。但要\textbf{警惕例外}：如果某个方案\textbf{把这个映射\emph{接进了真实的预测路径}}（推理时也要靠它变换表示），那它\textbf{就成了必需组件、\emph{不能}同样移除}——判断标准是“它在\emph{推理}时是否被用到”。

表示蒸馏有一对\textbf{内在张力}，理解它才算真懂 FitNets 的局限。
\begin{itemize}\item \textbf{匹配\emph{太死}（逐坐标硬配）}：要求学生第$j$维等于教师第$j$维，这\textbf{隐含假设了“内部坐标\emph{对应}”}——可第~\ref{sec:ket-integration-alignment}节的置换例子早已说明\textbf{内部坐标\emph{不是}语义本身}（同一功能有无数种坐标排列）。\textbf{强迫学生\emph{逐维}复制教师，可能\emph{束缚}学生、逼它放弃对自己更优、却和教师坐标不同的表示}（学生本可能用更少单元达成同样功能）。\item \textbf{匹配\emph{太松}（映射容量太大）}：学一个高容量映射虽\textbf{允许学生用“另一种表示”}（只要经映射能预测教师提示即可，更尊重学生自由），但这个\textbf{大映射可能\emph{背下}训练集上“学生表示到教师表示”的\emph{有限对应}}（死记硬背），于是\textbf{提示损失很低、\emph{却不代表}在\emph{新预测输入}上还能从学生表示读出同样的教师特征}——泛化是另一回事。\end{itemize}
所以评估要\textbf{两端都查}：在\textbf{独立样本}上既看\textbf{提示损失}（中段是否真的可迁移地对齐），也看\textbf{最终任务效果}；并且\textbf{务必按“\emph{实际部署路径}”（删掉脚手架映射后）去测学生}，而不是测带着训练脚手架的版本。当师生\textbf{跨结构、尤其\emph{跨模态}}（如图像教师、文本学生）时，连\textbf{空间位置、哪些信息可观测（图像/位置/时间）}都要\textbf{重新核对}，不能照搬同构网络里的对齐假设。

\subsection{承接模型的训练}
\label{sec:ket-integration-student-training}

两类约束都确定后，最后一步是\textbf{把它们和“真实标签”一起，组装成一个\emph{真正能跑的训练目标}}。先固定记号：令$S_{\mathrm{L}}$为\textbf{带真实标签}的训练样本集，$U$为\textbf{用于蒸馏的（常可无标签）输入集}；两者\textbf{允许重叠}，但\textbf{都\emph{绝不}包含最终测试材料}（否则评测失效）。要训练的是\textbf{学生参数$\bm{\theta}$}以及（如果用表示蒸馏）\textbf{提示映射参数$\bm{\phi}$}；而\textbf{教师和聚合规则在整个过程中\emph{冻结固定}}。

联合目标里\textbf{三类监督各管一段、互补而不重复}：\textbf{真实标签项}\emph{纠正“最终任务答案”}（锚定正确方向，防止被错误教师带偏）；\textbf{软目标项}\emph{限制输出分布\emph{偏离}教师}（学习暗知识）；\textbf{提示项}\emph{约束中间层残差}（对齐内部表示）。为了让三类\textbf{样本数不同的损失可比}，每项都\textbf{先在各自样本集上\emph{平均}}，再\textbf{乘一个非负权重}控制相对重要性。一种联合目标写为：
\begin{equation}
  \begin{aligned}
    L(\bm{\theta},\bm{\phi})={}&
      \frac{\lambda_{\mathrm{L}}}{|S_{\mathrm{L}}|}
      \sum_{(\bm{x},y)\in S_{\mathrm{L}}}
        -\log p_{\bm{\theta}}(y\mid\bm{x};1)\\
      &+\frac{\lambda_{\mathrm{D}}T^2}{|U|}
      \sum_{\bm{x}\in U}
        D_{\mathrm{KL}}\!\left(
          \bm{q}^{(T)}(\bm{x})\Vert
          \bm{p}_{\bm{\theta}}^{(T)}(\bm{x})\right)\\
      &+\frac{\lambda_{\mathrm{R}}}{|U|}
      \sum_{\bm{x}\in U}L_{\mathrm{rep}}(\bm{x};\bm{\theta},\bm{\phi}).
  \end{aligned}
  \label{eq:ket-integration-joint-distillation}
\end{equation}
读懂这个三项式的\textbf{使用规则}：$\lambda_{\mathrm{L}},\lambda_{\mathrm{D}},\lambda_{\mathrm{R}}$分别是\textbf{真实标签、蒸馏、表示}三项的权重，\textbf{设为0即关闭该项}；$p_{\bm{\theta}}(y\mid\bm{x};1)$是\textbf{温度取1的学生概率}（真实标签项用普通交叉熵，不升温）。\textbf{软目标项\emph{师生必须用\emph{同一个}温度}$T$}，并\textbf{带上前面推导的$T^2$尺度补偿}，否则高温软项会被淹没。

两个\textbf{健壮性}细节：其一，\textbf{启用某项的\emph{前提}是对应集合\emph{非空}}——\textbf{没有标签就把第一项整项去掉、没有内部提示就去掉第三项}，\textbf{绝不能让程序对一个空集求平均（$0$作分母）}；纯无标签蒸馏就是$\lambda_{\mathrm{L}}=0$的合法情形。其二，\textbf{温度$T$、三个权重、以及学生容量}这些\textbf{超参数，都要在“允许使用的\emph{验证}材料”上挑选}，不能碰测试集。最后要有合理预期：这个联合目标\textbf{本质是在“\emph{听真实标签}”和“\emph{听教师}”之间做\emph{折中}}，当\textbf{学生容量不足以同时满足所有要求、或几类信号\emph{互相矛盾}}（教师和真实标签系统性不一致）时，\textbf{不存在让三项\emph{各自}都达最小的解}，优化只会找到一个权衡点。

\begin{algorithm}[从冻结教师到可部署学生]
  \label{alg:ket-integration-distillation-loop}
  \begin{algorithmic}[1]
    \Require 固定版本的教师与相容接口，可用的$S_{\mathrm{L}}$、$U$，独立验证样本集，损失设置、优化器与训练预算
    \Ensure 所选学生模型参数、部署配置和验证报告
    \State 初始化学生与所需提示映射，冻结教师及其依赖路径
    \For{学习步骤$t=0,1,\ldots$，直到预算耗尽或满足停止条件}
      \State 按预定任务与类别比例，从各启用目标的训练样本集合抽取小批次
      \State 查询冻结教师或读取兼容缓存，取得软目标与可用提示
      \State 沿学生路径计算所需概率、提示映射与各项损失
      \State 按式~\eqref{eq:ket-integration-joint-distillation}汇总启用项
      \State 用预定优化器只更新学生及提示映射
      \State 按验证任务损失和保留要求检查停止条件，保存候选检查点
    \EndFor
    \State 按预定验证规则选择并冻结学生，移除仅训练需要的教师与映射
    \State 用未参与选择的测试材料检查任务、来源保持、组合能力与实际部署成本
  \end{algorithmic}
\end{algorithm}

算法流程看似机械，但有几个\textbf{决定成败的实操点}。\textbf{采样}环节：多任务时\textbf{每个小批次从各类数据里抽多少，直接决定各任务\emph{实际}拿到的梯度份额}——如果\textbf{常见任务样本比例天然就大}，又不做平衡，\textbf{少见任务的能力会在训练中被\emph{悄悄削弱}}（这和6.7.4、以及集成里的覆盖问题同源），所以常要按任务做重采样或加权。

\textbf{教师缓存}：为省算力常把教师回答预先算好存起来，但\textbf{缓存必须\emph{绑定}教师版本、预处理方式、温度}等配置；一旦\textbf{训练时对照片做了\emph{数据增强}}（裁剪、变色），\textbf{存的旧教师软标签\emph{未必}还对得上“增强\emph{后}的当前输入”}（增强前算的答案，配增强后的图，可能错配），要么对每种增强实时查教师、要么谨慎设计。\textbf{停止条件}应\textbf{事先定好}（预算耗尽、验证损失不再改善、或某项“能力保留”指标开始恶化即停），不能事后挑一个对自己最有利的点。\textbf{部署}时\textbf{一般把温度\emph{调回1}}输出正常概率；如果训练后还\textbf{额外做了概率校准}，那是\textbf{独立的一步}，其\textbf{参数和所用选择材料要\emph{单独记录}}。

怎么\textbf{公平地评判“蒸馏到底有没有用”}？要\textbf{固定任务、输入范围、测试样本}，并\textbf{有层次地报告}：先报\textbf{教师}、再报\textbf{蒸馏学生}，看学生逼近了多少；\textbf{如果有真实标签}，\textbf{关键是再加一个“\emph{同样结构、同样参数预算、但\emph{只}用真实标签训练}的学生”作对照}——蒸馏学生相对这个对照的\emph{增量}，才真正归因于“教师提供的额外信息”，而不是学生结构本身的功劳。如果\textbf{没有训练标签}，这个对照做不了，就应\textbf{在“相同可用输入、相同预算”下，改去和其他\emph{可行}基线（如自监督、其他弱监督）比较}。

还要\textbf{如实承认蒸馏\emph{不能克服}的三类客观限制}：\textbf{学生容量更小}会\textbf{增加逼近困难}（教师的某些行为学生根本装不下）；\textbf{异构结构}会\textbf{增加表示匹配和优化的难度}；若教师是个\textbf{只能远程调用的黑箱}，还会受\textbf{查询次数配额、输出精度（可能只给粗概率）、内部激活\emph{不可见}}等约束。请记住：\textbf{蒸馏/匹配损失只是“\emph{希望}学生靠近教师”的一个目标，它\emph{不会改变}上述容量、求解、访问权限这些\emph{客观}条件}——条件不具备时，损失设计得再巧也补不上。

底座升级将模型参数不对应的问题具体化。Trans-LoRA把旧底座加旧适配器作为教师，在新底座上训练新适配器，而不是把旧因子直接装上去。论文的迁移阶段保留五个任务种子示例，并使用原任务阶段训练的判别器：新底座生成候选预测输入，判别器筛选，旧模型为保留预测输入提供预测输出，新适配器学习这些预测输出。原完整训练样本迁移时不再读取，但种子、生成、筛选、查询与训练仍有样本和计算成本。论文的近乎不读取原训练样本设定不保证任意任务无损迁移\scite{ket-wang2024translora}

当蒸馏被用于\textbf{持续学习}（一个接一个地学很多任务）时，会出现\textbf{“压缩误差会不会\emph{逐轮累积}”}的新问题——每蒸馏一次就丢一点，很多轮后旧能力可能面目全非。Progress \& Compress（进展与压缩）用一个\textbf{双组件交替}架构来管理这件事：
\begin{itemize}\item 一个\term{活动分支}（active column）：\textbf{学\emph{当前}新任务}，它在学习时还能\textbf{读取冻结的“汇总网络”的表示}（站在历史肩膀上）。\item 一个\term{汇总网络}（knowledge base / 知识库）：\textbf{保存到目前为止的\emph{压缩}知识}。\end{itemize}
两者的\textbf{角色周期性\emph{互换}地协作}：学新任务时活动分支是学生、汇总网络是\textbf{被读的}参考；而到了\textbf{“压缩”阶段，方向\emph{反过来}}——\textbf{以刚学好的活动分支为\emph{教师}，把新行为\emph{蒸馏回}汇总网络}，让汇总网络在纳入新能力的同时，靠\textbf{在线弹性权重巩固（EWC）\emph{限制对重要旧参数的改动}}（防遗忘）。于是“当前学习”和“历史压缩”由两部分\textbf{交替承担}，\textbf{网络结构始终\emph{保持固定规模}}（不像给每个任务都并列挂一列、越积越臃肿）\scite{ket-schwarz2018progress}。

这套机制有几个\textbf{必须警惕的“\emph{自我强化}”风险}。其一，\textbf{监督目标会\emph{悄悄漂移}}：压缩阶段的教师（活动分支）若\textbf{依赖某个旧标注库/旧缓存}，就必须\textbf{把那个\emph{旧版本固定}住}；否则一旦你在压缩的同时又更新了汇总网络或底层数据，\textbf{教师所模仿的目标\emph{本身}在移动}，学生会去追一个不断变化的靶。其二，\textbf{蒸馏\emph{补不回}从未查询的区域}，也\textbf{不能突破容量上限}：如果\textbf{每一轮都只在\emph{常见花卉}上查询}，那么\textbf{少见物种的能力会\emph{一轮轮}地被稀释、偏离}；更糟的是\textbf{学生学成之后\emph{又去当\emph{下一任}教师}}——于是\textbf{早期的错误会被\emph{代际复制、越放大}}（信息逐代损失的“复制品”效应）。

因此长期运行要\textbf{逐任务复测三件事}：\textbf{旧任务能力\emph{保留}、\emph{通用}能力有没有退化、以及在\emph{相同预算}下学习\emph{新}任务的能力（可塑性）}；并\textbf{把\emph{教师存储、数据生成、反复查询、多轮重复训练}的费用\emph{累计}入账}。持续验证的完整方法由第~\ref{sec:ket-update-retention-validation}~节补充，而\textbf{压缩算子和部署实现能否落地，还必须结合\emph{具体模型结构和运行系统}逐项检查}。

\subsection{训练样本的凝练}
\label{sec:ka-sample-condensation}

前面的蒸馏，改变的是\textbf{模型}（把教师行为压进学生参数）。6.7.4 要讲的则是另一个维度的“压缩”：\textbf{改变所\emph{保存的训练数据}本身}。动机来自现实痛点：如果\textbf{同一批花卉照片要被\emph{反复}用来重训、做大量结构对比、或回放旧任务}，那么\textbf{每次都在成千上万张原始照片上训练，代价极高}。能不能\textbf{只保留/制造\emph{少得多}的样本，却让“在它们上训练”达到接近原数据集的\emph{某种}效果}？

这里有\textbf{两种程度不同}的做法，先分清：\textbf{其一，选子集（coreset/选择）}：从原始照片里\textbf{挑出一\emph{小部分真实}照片}保留——它\textbf{仍然是真实图像}，但\textbf{信息量受限于原样本}（挑不出原集里没有的东西）。\textbf{其二，合成凝练}：\textbf{允许\emph{自由调整像素}、\emph{凭空优化}出少量（甚至看起来不像花的）图像}，目标是让这一小撮合成数据\textbf{对\emph{指定的学习算法}特别“高效”}。这就是\term{训练样本凝练}（dataset distillation / dataset condensation）。

理解这个概念有\textbf{两个要害，缺一不可}：其一，它\textbf{追求的是“\emph{保留\emph{特定}的训练\emph{效果}}”（在合成集上训练后，模型在某些指标上接近用全集训练），\textbf{而\emph{不是}}“让合成照片\emph{看起来逼真}”}——凝练出的图像\textbf{常常是人类看不懂的抽象图案}，这完全正常，\textbf{视觉真实\emph{不是}它的目标}。其二，也是常被误解的：\textbf{合成数据\emph{没有新增任何关于物种的\emph{事实证据}}}——它的全部信息都\textbf{源自}原始数据集，\textbf{不可能凝练出原集里不存在的新知识}，只是\textbf{把原有训练信号“榨浓、重组”}。

设原训练样本集为$\mathcal T$，待优化的合成训练样本集为$\mathcal S$，$|\mathcal S|\ll|\mathcal T|$；$L_{\mathcal T}$和$L_{\mathcal S}$分别为两集上的平均任务损失。固定模型结构、合成样本数与类别配额，并规定允许的像素范围和正步长范围。令模型参数初值$\bm\theta_0$来自指定分布$P_0$，$A_J(\mathcal S;\bm\theta_0,\eta)$表示以步长$\eta$在合成样本上做$J>0$步梯度下降。Wang等把更新后的模型参数视为合成像素的函数，再通过更新步骤反向求导，优化合成像素与步长；以共同步长写出的多步目标为
\begin{equation}
  \min_{\mathcal S,\eta}\quad
  \mathbb E_{\bm\theta_0\sim P_0}
  L_{\mathcal T}\!\left(A_J(\mathcal S;\bm\theta_0,\eta)\right).
  \label{eq:ka-condensation-objective}
\end{equation}
这个目标里有一个\textbf{“双层优化（bi-level）”}的嵌套结构，是理解全节所有方法的钥匙，务必看清\textbf{内外两层各自在干什么}：
\begin{itemize}\item \textbf{内层（更新模型）}：$A_J(\mathcal{S};\bm{\theta}_0,\eta)$表示“\textbf{从初值$\bm{\theta}_0$出发，在\emph{合成集}$\mathcal{S}$上\emph{真的做$J$步梯度下降}}”——合成样本\textbf{在内层\emph{推动模型参数更新}}。\item \textbf{外层（更新合成样本）}：内层训出一个模型后，\textbf{在外层把它放到\emph{原始集}$\mathcal{T}$上计算损失}$L_{\mathcal{T}}$，用“学完之后在\emph{真数据}上表现好不好”来\textbf{反过来调整合成图像的像素和步长}。\end{itemize}
也就是说，\textbf{合成像素的好坏，是由“它训练出来的模型在\emph{原集}上的表现”\emph{间接}定义的}——这是个“\textbf{为了‘教得好’而优化‘教材’}”的元优化。还要注意\textbf{数据纪律}：$\mathcal{T}$和$\mathcal{S}$\textbf{都属于\emph{开发过程}}，\textbf{最终测试集必须\emph{另外}留出}，不能让凝练过程窥见测试数据。当然，式~\eqref{eq:ka-condensation-objective}只是\textbf{规定了理想目标}，它\textbf{不保证最小值一定能取到、或被找到}（这是个困难的非凸双层问题）。关于初值和适用性还要补两句：模型初值$\bm{\theta}_0$\textbf{可以\emph{固定}成一个}（为这一个起点量身定制教材），也可以\textbf{按约定分布$P_0$\emph{随机重新抽取}}（希望教材对不同初始化都管用，更鲁棒但更难）；原论文\textbf{固定住分类标签、只优化像素}，并\textbf{允许\emph{每一步}用不同学习率}（分步学习率也当成可优化变量）。\textbf{适用边界}要清醒：\textbf{在\emph{某一个}模型结构、某一个初值、某一套训练步数下凝练成功，\emph{不保证}换成\emph{别的结构/初值/步数}后这些合成样本仍然有效}——凝练出的教材往往\textbf{和“\emph{特定学习算法}”绑定}，通用性是它的主要软肋之一\scite{ka-wang2018-dataset-distillation}。

上面的方法有个\textbf{致命的工程瓶颈}：目标要求\textbf{“\emph{穿透}内层那$J$步梯度下降去\emph{求导}”}（为了用外层损失调整像素，必须对整条“$J$步更新链”做反向传播）。当$J$稍大，这种\textbf{“对训练过程求导”}\emph{计算图极深、显存和算力开销爆炸}，几乎不可行。

Zhao 等人（Gradient Matching）的巧妙之处是\textbf{“\emph{斩断}长链、只看\emph{一步}”}：既然“穿透很多步”太贵，那就\textbf{不追求最终模型一致，转而要求“在训练路径的\emph{每一个当前参数点}上，合成集和原集\emph{产生的\emph{梯度方向}相近}”}——\textbf{让两者\emph{先推动一个\emph{相近的局部更新}}}。如果\textbf{每一小步}的方向都接近，那么\textbf{累积起来}走出的训练轨迹自然也近似。用$D$表示预先选定的梯度差异度量，沿合成样本更新路径的一种写法是
\begin{equation}
  \begin{aligned}
    \min_{\mathcal S}\quad&
    \mathbb E_{\bm\theta_0\sim P_0}
    \sum_{t=0}^{J-1}
    D\!\left(\nabla L_{\mathcal S}(\bm\theta_t),
             \nabla L_{\mathcal T}(\bm\theta_t)\right),\\
    &\bm\theta_{t+1}
    =\bm\theta_t-\eta\nabla L_{\mathcal S}(\bm\theta_t).
  \end{aligned}
  \label{eq:ka-condensation-gradient-matching}
\end{equation}
读这个目标时一个\textbf{最容易绕晕}的地方是“\textbf{到底在对谁求导、又在更新谁}”，一定分清两层导数：
\begin{itemize}\item 式中$\nabla L_{\mathcal{S}}$和$\nabla L_{\mathcal{T}}$\textbf{都是对\emph{模型参数}$\bm{\theta}_t$求的梯度}（这是“训练\emph{会把模型往哪推}”的方向）；\item 但\textbf{外层优化器\emph{更新}的对象却是\emph{合成样本}（像素）}——它通过“改变像素会如何改变合成梯度”这条链，去最小化两种梯度的差异。\end{itemize}
还有一个\textbf{关键的比较方式}：两种梯度是在\textbf{同一个、\emph{共享}的参数点}$\bm{\theta}_t$\textbf{上\emph{当场}分别算出来再比}的，\textbf{而\emph{不是}“让模型在合成集和原集上\emph{各走各的}、走出两条\emph{已经分道扬镳}的路径后，再去\emph{逐坐标硬配}两个终点参数”}——后者正是第5章SwAV等处反复警告过的“不可逐维硬配”的错误。原方案还\textbf{按\emph{类别}分别匹配梯度}、对\textbf{逐层梯度做归一化}后再算距离（防止某层范数主导），并\textbf{交替地更新合成样本和模型参数}。最后是边界：\textbf{“\emph{局部}梯度每步都像”\emph{并不自动保证}“\emph{长路径}累积后仍像、或\emph{最终风险}相同”}——小步方向一致只是长程一致的\emph{必要而非充分}条件\scite{ka-zhao2021-gradient-matching}。

用可复算的一维回归例说明这种匹配在做什么。固定预测输入为1、模型预测为$\theta$，两条原训练样本的目标值为1和3，单条损失为$(\theta-y)^2/2$；平均梯度是$\theta-2$。此教学例允许设置合成目标值，区别于上述固定分类标签、只优化像素的实现。把两条样本凝练成一条目标值为2的合成样本，梯度仍为$\theta-2$。从$\theta_0=0$以步长$1/2$更新，两套训练样本集都依次得到$1,1.5$。这个均值充分的二次损失例只证明当前模型与更新规则等价，不证明一张合成花卉照片能保留所有物种差异。

梯度匹配只约束“\emph{一步}之近”，那能不能\textbf{直接约束“\emph{更长}一段训练的\emph{净效果}}”？轨迹匹配（Trajectory Matching）换了个更有全局观的参照系。它\textbf{先在\emph{原始数据集}上\emph{完整训练}出若干个模型，并把训练过程中\emph{一连串\emph{参数检查点}}保存下来——这些点连成“专家\emph{训练轨迹}”}。然后在每一\emph{轮}：\textbf{取某条真实轨迹上的一个起点}$\bm{\theta}_t^\star$，\textbf{让一个\emph{同起点}的学生模型在\emph{合成集}上走$N$步}得到$\widehat{\bm{\theta}}_{t+N}$，\textbf{要求这个终点\emph{逼近}“真实模型\emph{再走$M$步}后到达的”}$\bm{\theta}_{t+M}^\star$。其\term{轨迹匹配}（trajectory matching）用归一化距离
\begin{equation}
  L_{\mathrm{traj}}
  =\frac{\lVert\widehat{\bm\theta}_{t+N}
                   -\bm\theta_{t+M}^\star\rVert_2^2}
         {\lVert\bm\theta_t^\star
                   -\bm\theta_{t+M}^\star\rVert_2^2},
  \qquad
  \widehat{\bm\theta}_t=\bm\theta_t^\star,
  \label{eq:ka-condensation-trajectory-matching}
\end{equation}
读懂这个归一化距离：\textbf{分子}是“学生走$N$步的终点”与“参照走$M$步的终点”之差的平方；\textbf{分母}是“参照模型\emph{自己}从$\bm{\theta}_t^\star$走到$\bm{\theta}_{t+M}^\star$移动了多远”的平方。\textbf{除以分母的作用是“\emph{按轨迹\emph{自身步长}归一化}”}——在专家本身移动很大的区间，容忍的绝对偏差也相应放大，使不同训练阶段的匹配误差\textbf{尺度可比}；它要求分母为正（即该区间专家确实在移动）。优化时同样要\textbf{穿透学生那$N$步更新去反向调整合成像素和学习率}；注意\textbf{$M$和$N$\emph{不必相等}}（合成集可以用更少的步数、更大的跳跃去追平真实训练更多步的净效果，这正是“加速训练”的来源）。原方案还使用\textbf{可微的参数变换、并在\emph{多条}专家轨迹上匹配}以增强鲁棒性。边界依旧：\textbf{“某一段端点对上了”是个\emph{局部}目标，\emph{不能}夸大成“合成集与任意训练过程\emph{全局等价}”}\scite{ka-cazenavette2022-trajectory-matching}。

最后\textbf{算清这笔账、也认清这些坑}，别把数据集凝练当成“免费午餐”。
\textbf{成本是“\emph{提前}、且常常很贵”}：合成阶段要付出\textbf{在原集上训练专家、存储大量轨迹检查点、反复做梯度匹配/穿透展开求导、以及大量验证}的开销。所以\textbf{它\emph{省不省钱}，完全取决于“这批合成数据\emph{之后被\emph{复用}多少次}”}——只用一次，前期投入往往收不回来；要\textbf{反复重训/结构搜索/旧任务回放}，摊薄后才可能划算。
\textbf{它还会\emph{继承原数据集的缺陷}}：原集里\textbf{少数类样本缺失}，合成集\textbf{大概率照样缺}（无中生有不出来）；原集\textbf{标签有错}，凝练过程可能\textbf{把错误也\emph{高效}地浓缩进去}。务必破除三个想当然：\textbf{“样本更少”\emph{并不自动}带来}“\textbf{质量更高}”（可能丢信息）、“\textbf{隐私自动受保护}”（合成图可能仍\emph{记忆/泄漏}原图，隐私需专门保证）、或“\textbf{跨结构通用}”（换架构常常失效）。
\textbf{公平的选择方法}是：在\textbf{相同的可用训练数据、相同的合成预算、相同的后续训练协议}下，\textbf{把“真实子集、合成凝练集、完整原集”三者\emph{并排比较}}，并\textbf{额外在\emph{新结构、新初值}上检验合成集的泛化}。这些合成样本如何用于\textbf{旧任务回放}由第~\ref{sec:ket-update-memory}节接续，而\textbf{跨任务比较}仍需第~\ref{chap:ket-evaluation}章那套\emph{独立评价}来背书。

\section{融合结果的表达、使用与验证}
\label{sec:knowledge-fusion-result-use}

走到这里，融合的\emph{计算}基本完成了——我们手上有了票数、加权得分，或一个像“山茶$0.62$、茶梅$0.28$、月季$0.10$”这样的融合分布。但\textbf{“算出结果”}和\textbf{“\emph{怎么用}这个结果”}是\textbf{两件\emph{必须分开}的事}，本节专门处理后者。

先用一个反差引出问题：\textbf{同一个$0.62$}，在不同用途下\textbf{命运可以\emph{完全不同}}——如果它是\textbf{一个训练任务}的软标签，系统完全可以\textbf{把“$0.62/0.28/0.10$”整份分布\emph{保留}下来}去训练；可如果它要用于\textbf{高风险的真实决策}（比如误判成某有毒近缘种会造成严重后果），使用方可能\textbf{认为$0.62$\emph{根本不够格}，于是\emph{任何类别都不接收}、直接弃权转人工}。

\textbf{请注意：这两个决定\emph{互不矛盾}，差异来自“\emph{用途和代价}”，而\emph{不是}模型对同一张照片突然有了两种确信程度}——模型给的支持始终是那一个$0.62$，变的只是使用门槛。因此本节的核心纪律是\textbf{严格区分两个层}：\textbf{推断结果层}（票数、得分、后验分布、整体赋值——“证据\emph{说了什么}”）与\textbf{使用动作层}（选一个类、拿去训练、保留候选集、还是弃权——“你\emph{决定做什么}”）。同时，为了能审计和复现，\textbf{来源摘要、对齐映射、融合配置、结果版本}都要\textbf{随结果一起保存}。

\subsection{单一标签与概率分布的使用}
\label{sec:knowledge-fusion-label-probability-use}

最常见的使用动作是“\textbf{接收端\emph{必须}吐出\emph{一个}类别}”。直觉上选后验最大的那个就行了，但这里有个\textbf{常被忽略的前提}：\textbf{“选概率最大的”只有在“\emph{所有错误\emph{代价都一样}}”时才天然合理}。现实中错误往往\textbf{不对称}：\textbf{把有毒的茶梅误报成可观赏的山茶}（可能被误食），和\textbf{漏掉一个本应送去复核的稀有类别}，代价可能\textbf{差很多倍}。\textbf{这种“犯不同错误损失不同”的\emph{取舍}，单看一个光秃秃的最大概率\emph{根本无从表达}}——概率只说“多可能”，没说“错了多疼”。设$q_i(k)$为融合类别分布，$A$为用途允许的行动集合，$L(a,k)$为真值$k$下行动$a$的损失；我们对未知真值按$q_i$加权平均损失，再在允许行动中选最小者，得到贝叶斯决策：
\begin{equation}
  \widehat a_i
  \in\argmin_{a\in A}
  \sum_{k=1}^{K}L(a,k)q_i(k).
  \label{eq:knowledge-fusion-bayes-decision}
\end{equation}
这个公式就是\term{贝叶斯决策}（Bayesian decision），读法很直白：$\sum_k L(a,k)q_i(k)$这一\textbf{求和，是在给“采取行动$a$”算\emph{预期会\emph{平均}付出多少损失}}——它\textbf{走遍所有可能的真值$k$}，每一种的损失$L(a,k)$\textbf{按该真值出现的后验概率}$q_i(k)$加权；外面的$\argmin$则\textbf{在所有允许行动里挑“预期损失\emph{最小}”的那个}（注意这里是\textbf{最小化}损失，和前面最大化支持相反）。

现在可以\textbf{看清“最大后验（MAP）”到底是\emph{什么特例}}了：当\textbf{行动集合恰好就是“选一个类别”}、并且采用\term{零一损失}（\textbf{选对损失0、选错损失1，所有错误等代价}）时，行动$a$的期望损失等于
$$\sum_k L(a,k)q_i(k)=\sum_{k\neq a}q_i(k)=1-q_i(a).$$
\textbf{最小化}$1-q_i(a)$\textbf{恰好等价于\emph{最大化}$q_i(a)$}——于是“选最大后验”不过是\textbf{贝叶斯决策在“错误等代价”这个\emph{特殊}假设下的退化形式}。一旦错误不等价，就必须老老实实用完整的损失矩阵，MAP 不再最优。用一个小算例\textbf{坐实“同一份概率、不同代价、相反决策”}。设二分类后验$q_i=(0.62,0.38)$（分量依次为山茶、茶梅），乍看山茶概率高、该选山茶。现在\textbf{规定非对称代价}：\textbf{把茶梅误报成山茶}（行动选山茶、真值却是茶梅）的损失是\textbf{2}（比如茶梅需特殊处理、漏掉代价高）；\textbf{把山茶误报成茶梅}的损失是\textbf{1}；报对损失0。
\begin{itemize}\item \textbf{行动“选山茶”}：只在真值为茶梅（概率$0.38$）时错，期望损失$=2\times0.38=0.76$。\item \textbf{行动“选茶梅”}：只在真值为山茶（概率$0.62$）时错，期望损失$=1\times0.62=0.62$。\end{itemize}
比较期望损失：$0.62<0.76$，于是\textbf{贝叶斯决策选择\emph{茶梅}}——\textbf{尽管它的概率（$0.38$）\emph{更低}}！原因正是“\textbf{漏判茶梅代价更高}”，值得为规避这个高风险而押在低概率类上。\textbf{同一份后验：在零一损失下选山茶、在这组非对称损失下选茶梅}。请反复体会这句结论——\textbf{改变的是\emph{用途的代价结构}，\emph{不是}模型的概率或确信程度}；把决策的差异错误归因于“模型更相信谁”，是混淆了推断层和使用层。

当接收端\textbf{不是要“立刻采取行动”，而是“\emph{训练另一个分类模型}”}时，使用目标又变了：此时\textbf{我们\emph{不}替它做决策，而是希望它学出来的预测分布$p_{\bm{w}}$去\emph{逼近}融合分布$q_i$}。最自然的损失就是\textbf{软标签交叉熵}
$$-\sum_k q_i(k)\log p_{\bm{w}}(k\mid\bm{x}_i),$$
其中$p_{\bm{w}}$是待训练模型的类别概率。它的好处是\textbf{$q_i$分配给\emph{每一类}的支持\emph{都}进入了损失}——$0.28$的茶梅、$0.10$的月季都会给学生相应的梯度，学生学到的是“一整张分布的形状”，而不只是一个名字。

作为对比，如果你\textbf{只保留$\argmax_k q_i(k)$}（硬化成一个标签“山茶”），接口和存储\textbf{确实更简单}，但\textbf{把其余$0.38$的质量\emph{整份丢弃}}，\textbf{也\emph{藏起了冲突强度}}（下游再也看不出“这张其实挺纠结，只有$0.62$”）。最后是质量底线：\textbf{软训练只是“\emph{忠实地传递}$q_i$”，它\emph{不会自动把一个坏$q_i$变好}}。如果$q_i$本身来自\textbf{未校准的分数、或设错了的来源模型}，那么它表达的只是“\emph{当前这套模型的（可能有偏的）支持}”，\textbf{再怎么用它做软标签，也\emph{不会}把它 magically 变成真实频率}——垃圾进、垃圾出。

为了\textbf{事后能审计、能追溯“结论为什么变了”}，分布和使用规则\textbf{都必须\emph{分别}版本化}，因为它们\textbf{是两类\emph{独立}的变化来源}：
\begin{itemize}\item \textbf{改变“推断”}：\textbf{新增来源、更新混淆矩阵、修订语义映射}——这些会\textbf{真正改变融合分布$q_i$本身}（证据变了）。\item \textbf{只改变“使用”}：\textbf{调整决策阈值、换一个损失矩阵}——这些\textbf{分布一个数都没动}，只是\textbf{最终选的动作可能变了}（代价/门槛变了）。\end{itemize}
\textbf{审计时，只有\emph{同时}拿到“原始分布”和“这两类配置”，才能定位一次结论翻转\emph{究竟}是“证据更新导致的”还是“使用规则调整导致的”}；若只存了最终标签，这个归因就永远丢失了。还有一个常见的工程妥协：如果\textbf{某个训练接口\emph{只接受硬标签}}（喂不进概率），可以\textbf{在接口边界处\emph{临时}选类}来适配，但\textbf{上游那份完整概率向量仍应\emph{保留}下来供复核}——适配接口不该成为丢弃不确定性信息的理由。

\subsubsection{含噪标签下的损失校正}
\label{sec:ka-noise-correction}

到目前为止，用多源噪标签训练的主流思路是“\textbf{先融合}”：先推断每张照片的候选真值（软/硬标签），\textbf{再拿这个“\emph{净化过}”的标签去训练}最终分类器。但这有个隐患——\textbf{融合出的标签可能\emph{本身}就估错了}，而你\textbf{已经把原始票里的不确定性丢掉}了。

还有一条\textbf{思路相反}的训练路径，叫“\textbf{不替换、\emph{显式建模报错过程}}”：\textbf{保留\emph{原始观测到的那张噪标签票}（不把它改写成猜测真值）}，转而\textbf{在模型里\emph{明确}写出“\emph{真实类别}\emph{如何经过一道\emph{标注错误}通道}\emph{变成}我们看到的票}”。训练时让分类器\textbf{去拟合“\emph{在我的真类预测下、\emph{可能观察到}什么样的噪标签}”}，而\textbf{不是}生硬地把票当真值。这样做的好处是：\textbf{模型\emph{不必先}给出一个可能错误的“确定答案”，而是把标注噪声\emph{显式吸收}进损失}，避免\textbf{分类器\emph{将错就错}地把“报错规律”学成“真实类别规律”}（比如系统性地学会把茶梅标成山茶）。

你应该已经见过描述这个“真类$\to$观测票”方向的工具了——正是6.2的\textbf{来源混淆矩阵}。下面就把它\textbf{从“融合推断”里搬过来、\emph{接进训练损失}}。

设真类$Y$与观测标签$\widetilde Y$均取$1,\ldots,K$，$X$为照片的预测输入。暂采用类别条件噪声假设$\mathbb P(\widetilde Y=b\mid Y=a,X=x)=T_{a,b}$，即给定真类后，报错概率不再随照片特征变化。\term{噪声转移矩阵}（noise transition matrix）$\bm T$的第$a$行给出真类$a$会被报成各类的概率，元素非负、每行和为1；它沿用式~\eqref{eq:knowledge-fusion-confusion-matrix}的行列方向，但此处固定一个训练标签通道。以$\bm\theta$表示分类器模型参数，令$\bm p_{\bm\theta}(x)$为其预测的真类概率列向量。先经通道得到观测标签概率，再与实际类别票$b$比较，交叉熵的\term{前向损失校正}（forward loss correction）为
\begin{equation}
  \bm q_{\bm\theta}(x)=\bm T^{\top}\bm p_{\bm\theta}(x),
  \qquad
  \ell_{\mathrm F}(\bm\theta;x,b)
  =-\log q_{\bm\theta,b}(x).
  \label{eq:ka-forward-correction}
\end{equation}
理解前向校正的\textbf{关键}是“\textbf{让模型\emph{对观测}负责，而不是\emph{假装}观测等于真值}”：训练时\textbf{优化改变的是分类器参数$\bm{\theta}$（即真类预测$\bm p_{\bm{\theta}}$）}，而\textbf{观测标签$b$\emph{原样保留}}——\textbf{既没有被\emph{声明}为正确，也没有被\emph{改写}}。模型为了让“\emph{经过$\bm T^{\top}$后、在$b$这一类上的概率}”变大，必须\textbf{反过来把内部的真类预测$\bm p_{\bm{\theta}}$摆到\emph{正确}位置}（因为它知道有一部分真山茶会被观测成茶梅）。

那么它\textbf{为什么能\emph{纠正}}？在一组\textbf{理想条件}下——$\bm T$\textbf{准确且\emph{可逆}}、模型族\textbf{有能力\emph{表示}真实的类别后验}、且优化\textbf{能到达\emph{总体}极小点}——\textbf{最小化这个前向校正损失的解，恰好\emph{恢复}真实类别后验}（直观上，既然模型必须同时解释所有类别的观测频率、而$\bm T$可逆，唯一对得上的真类分布就是真实分布）。但这三个条件\textbf{缺一就可能失效}，要\textbf{逐条打折}：\textbf{模型容量\emph{受限}}可能根本表达不出真实后验；\textbf{$\bm T$是\emph{估计}出来的、本身带误差}；\textbf{有限步的梯度下降\emph{未必}到达总体极小点}。所以这是个“\emph{理想情形可纠正}”的结论，\textbf{不是任意实践中自动成立的保证}。

也可以直接修正损失而非概率。以下取严格正的类别预测概率，并假设所需对数损失可积，避免无穷大损失之间出现未定义的线性组合。把按各个真类计算的损失排成列向量$\bm\ell(\bm\theta;x)$，第$a$项为$-\log p_{\bm\theta,a}(x)$；$\bm T$可逆时，\term{后向损失校正}（backward loss correction）取
\begin{equation}
  \widetilde{\bm\ell}=\bm T^{-1}\bm\ell,
  \qquad
  \mathbb E[\widetilde\ell_{\widetilde Y}\mid X=x,Y=a]
  =(\bm T\widetilde{\bm\ell})_a=\ell_a.
  \label{eq:ka-backward-correction}
\end{equation}
这个式子为什么叫“\emph{无偏}校正”？看条件期望那一行：当真值为$a$时，观测标签$\widetilde Y$会\textbf{按通道的第$a$行}$T_{a,b}$\textbf{随机抽到各类}$b$，于是校正后损失$\widetilde\ell_b$的期望就是$\sum_b T_{a,b}(\bm{T}^{-1}\bm\ell)_b=(\bm{T}\bm{T}^{-1}\bm\ell)_a=\ell_a$——\textbf{矩阵和它的逆\emph{正好相消}}，期望\textbf{恰好还原干净损失}$\ell_a$。再对照片和真类求平均，就能\textbf{在期望意义上恢复“用干净标签训练”的风险}，所以它对噪声抽样是\emph{无偏}的。

但这里有个\textbf{反直觉、必须接受}的现象，第一次见常会以为算错了：\textbf{校正后的\emph{单条}损失$\widetilde{\bm{\ell}}=\bm{T}^{-1}\bm\ell$\textbf{可能出现\emph{负数}}（后面数值例会看到$-0.026$、$1.54$这种）}。这\textbf{在数学上完全正常}——因为$\bm{T}^{-1}$里有负元素，而无偏性是对“\emph{噪声标签的\emph{随机抽样}求期望}”才成立的性质，\textbf{并不要求每一条具体损失都非负}。这些负值在\textbf{平均}后仍给出正确的干净风险，但\textbf{单条出现负值可能\emph{增大优化波动}}（损失不再被0下界约束）。最后澄清：\textbf{后向校正是“改变\emph{训练目标}使其\emph{期望无偏}}”，它\emph{自始至终\emph{没有改动、也没有“修好”任何一张标签}}——原始噪标签还是那张噪标签。Patrini等区分了前向校正的理想极小点性质与后向校正的风险无偏性质\scite{ka-patrini2017-noise-correction}

仍以山茶、茶梅作二分类教学例，设
\[
  \bm T=\begin{bmatrix}0.8&0.2\\0.1&0.9\end{bmatrix},
  \qquad
  \bm p=(0.75,0.25)^{\top}.
\]
观测到山茶票的概率为$0.8(0.75)+0.1(0.25)=0.625$，茶梅票为$0.375$。山茶票的前向损失约为$-\log0.625=0.4700$，而直接把它当真类的损失为$-\log0.75=0.2877$。后向校正的两项损失约为$(-0.0262,1.5432)$；若真类为山茶，按该行通道平均得$0.8(-0.0262)+0.2(1.5432)\approx0.2877$。负的单条值并未破坏期望等式，却可能增加优化波动。

前面的前向、后向校正都\textbf{假设噪声通道$\bm{T}$\emph{已知}}，这是个很强的前提。现实中$\bm{T}$往往\textbf{也要从数据里\emph{估计}}，而这里\textbf{藏着一个和6.2“可识别性”\emph{同源}的根本困难}：\textbf{只看观测到的类别票\emph{边缘比例}，\emph{无法自动反推出}通道}——同一组观测比例，可能由“高准确率通道”产生，也可能由“乱翻通道”配不同真类分布产生，数据本身分不出来。假设某一预测输入上的山茶票概率为$0.6$，它既可由真类概率$0.6$和无报错通道产生，也可由真类概率$0.7$与两类均以$0.25$概率翻转的通道产生，因为$0.75(0.7)+0.25(0.3)=0.6$。具体看这个“\textbf{同一观测、多种解释}”：山茶票概率$0.6$，既可能来自“真类概率就是$0.6$、通道\emph{完全不报错}（单位阵）”，也可能来自“真类山茶概率为$0.7$、通道以$0.25$概率在两类间乱翻”——验算$0.75\times0.7+0.25\times0.3=0.6$，\textbf{两种解释给出\emph{完全相同}的观测}，却要求\textbf{不同的真类概率}。这正是6.2讲过的\textbf{不可识别}在噪声学习里的翻版：\textbf{必须靠\emph{外部信息}来打破歧义}——\textbf{干净锚点}（一批确知正确的标签）、\textbf{独立参考}、或\textbf{额外的观测}。\textbf{绝不能因为“训练损失\emph{收敛}了”就反推“通道$\bm{T}$已被正确识别”}——错误的通道下损失照样可以收敛到一个稳定但有偏的解。

Liu等进一步分析实例相关通道的可识别性：对同一预测输入$x$，若三份噪声标签在给定$x$与真类后条件独立、共享同一满秩通道$\bm T(x)$，且各类别均有$\mathbb P(Y=a\mid X=x)>0$，则其联合分布可将该通道与真类分布识别到类别置换。这里讨论总体分布的可识别性，不是只收集三张已观察类别票就能估准通道或恢复真值。不同来源有各自混淆、共同复制或几乎不重叠时，仍须回到本章的来源假设逐项检查\scite{ka-liu2023-noise-identifiability} 识别条件成立后，估计精度还是有限样本问题。每轮概率变换、通道估计与验证增加费用；后向矩阵求逆还可能因病态性放大估计误差。实例依赖噪声若不满足此处固定通道假设，应另建模型或保留未决，不得直接套用同一矩阵的保证。

还有不显式估计$\bm T$的训练路线。Co-teaching让两个网络各自在同批训练样本中按小损失选择预定比例，再把所选子集交给另一网络更新；保留率日程依赖噪声率等设定。DivideMix逐轮用二成分高斯混合拟合逐样本损失，以另一网络形成的划分区分较可信与可疑样本，再把可疑部分当无标签训练样本参与半监督学习，而非全部删除。两者都把小损失当筛选线索，相关错误、稀有难例与错误的混合假设仍会造成误选；它们不证明每条保留标签正确，也不保证任意实例相关噪声均可消除\scite{ka-han2018-coteaching,ka-li2020-dividemix}

\subsection{候选集合与弃权的使用}
\label{sec:knowledge-fusion-set-abstention-use}

贝叶斯决策逼你“\emph{即使没把握，也得\emph{选一个}行动}”，但很多真实场景里，\textbf{更诚实的做法是“\emph{承认我现在分不出来}”}。当融合分布\textbf{还\emph{不足以区分}两个类别}（比如山茶$0.51$、茶梅$0.49$）时，\textbf{硬选一个会\emph{白白丢掉}这份“势均力敌”的歧义信息}，还给人一种“已经确定”的假象。更好的选择是\textbf{先输出一个\emph{候选集合}}$C_i\subseteq Y$——不假装知道唯一答案，而是声明“真值\textbf{很可能落在这\emph{几个}类别之中}”。一种构造按$q_i(k)$降序加入类别，直到累计质量达到预设$\tau\in(0,1]$，相同概率的排序也要有固定规则；另一种保留所有$q_i(k)\geq\gamma$的类别，$\gamma\in[0,1]$是逐类接收阈值。前者在有效分布下会得到非空集合，后者可能为空；单元素集合相当于硬标签，过大集合则少有区分作用。两种构造法各有性格，要根据用途选：\textbf{累计质量法}因为从最高分往下加，在概率有效时\textbf{一定能得到非空集合}（最差也含最高类），且能\textbf{控制“集合\emph{合计}多可信”}；\textbf{逐类阈值法}更\textbf{直接}（保留所有“单独看就够格”的类），但\textbf{可能\emph{没有}任何一类越过阈值（集合为空）}，这本身是个有用信号（“没有一类足够可信”）。两种方法的退化情形也直观：\textbf{单元素集合就\emph{等价于硬标签}}；\textbf{集合\emph{过大}}（把几乎所有类都纳入）则\textbf{几乎没有区分作用}，等于没说。

\textbf{最关键、也最容易被一厢情愿误解的一点}：\textbf{“累计质量达到$\tau$”或“逐类概率超过$\gamma$”，\emph{只描述}了\emph{模型自己分布}的形状，它\emph{绝不}等于“\emph{保证}真值以$\tau$（或相应高概率）\emph{落在集合里}}”。前者是“\emph{模型\emph{自称}多有把握}”，后者是“\emph{事实上\emph{覆盖}得准不准}”的频率性质——\textbf{两者之间隔着“\emph{校准}}”这座桥。要让“$90\%$候选集真的有约$90\%$的时候盖住真值”，\textbf{必须\emph{独立}校准、并依赖相应的分布假设}，不能凭阈值名字想当然。

保留了候选集，并不意味着事情结束——\textbf{使用端可能\emph{仍然}需要一个\emph{确定}判断}（下游系统没法处理“可能是A或B”）。这时需要一条\textbf{升级/兜底}规则：当\textbf{综合全部证据后，仍\emph{无法形成当前用途\emph{可接收}的结论}}时（最高支持太低、候选太多、或关键来源弃权），就把结果\textbf{明确记为“\emph{未决}”}，而不是勉强给个答案；\textbf{使用端\emph{主动}决定\emph{不提交}这个决策}的\textbf{动作}，沿用\term{弃权}（abstention）的约定，后续可\textbf{转人工复核、追加观测、或追加标注}。

务必\textbf{区分\emph{两个不同环节}的“弃权”}，它们在本章台账里是\emph{不同字段}，不能混记：一个是\textbf{上游“\emph{提供标签的来源}”在\emph{运行后弃权}}（6.2讲的来源层$o_{i,s}=0$，描述的是“某个来源没给票”）；另一个是\textbf{下游“\emph{融合结果的使用端}弃权”}（本节描述的是“系统\emph{综合后}决定不下最终结论”）。\textbf{一个发生在证据产生阶段、一个发生在证据使用阶段，必须\emph{分别记录}}。系统可以因$\max_kq_i(k)<\eta$、有效独立来源不足或对齐未确认而记未决，再请求复核、追加观测或标注或由使用端弃权；这些接收条件需预先明确。系统可以依据\textbf{预先明确的接收条件}记未决，比如：最高支持$\max_k q_i(k)<\eta$（没有一类够笃定）、\textbf{有效\emph{独立}来源数不足}（怕被复制票撑场面）、或\textbf{对齐尚未确认}（对象/语义还没对上，回到6.3）。这些\textbf{触发条件要\emph{事先}定好}，不能事后挑对自己有利的。还有一条统计纪律：\textbf{弃权\emph{不是}一个负类}——弃权的照片\textbf{仍然属于目标总体}（它还是一张要处理的花卉照片，只是这次没下结论而已），所以在做\textbf{覆盖率、总体质量}统计时，\textbf{不能把这些弃权样本\emph{永久}地\emph{剔除、略去}不计}，否则总体指标会被“只报敢答的简单样本”系统性美化。

图~\ref{fig:knowledge-fusion-output-forms}固定教学分布$q=(0.62,0.28,0.10)$，分量依次为山茶、茶梅、月季，再只改变使用规则。最大后验选山茶，软训练保留三项；累计阈值$\tau=0.90$时，$0.62+0.28=0.90$，候选为$\{\text{山茶},\text{茶梅}\}$。若高风险用途要求最大支持至少$0.70$，当前$0.62$不足，记为未决并由使用端弃权。四个分支是不同用途的对照，不是一次使用中同时执行互斥行动。

\begin{figure*}[htbp]
  \centering
  \begin{tikzpicture}[
  x=1mm,y=1mm,
  every node/.style={font=\figurefont\small,text=BookInk},
  evidence/.style={book node,draw=BookInk,fill=BookPaper,line width=1pt,minimum width=61mm,minimum height=12mm,inner sep=2pt},
  result/.style={book node,minimum width=33mm,minimum height=18mm,inner sep=2pt,text width=30mm,align=center}
]
  \node[evidence] (q) at (75,19) {教学支持分布（山茶、茶梅、月季）\\$q=(0.62,\ 0.28,\ 0.10)$};

  \node[result,draw=BookBlue,fill=FigureInput!9,line width=0.9pt] (hard) at (12,-14) {硬标签\\山茶};
  \node[result,draw=BookTeal,fill=FigureModel!9,line width=0.9pt] (soft) at (54,-14) {概率目标\\$(0.62,0.28,0.10)$};
  \node[result,draw=BookAmber,fill=BookAmber!9,line width=0.9pt,dashed] (set) at (96,-14) {候选集合\\$\{\text{山茶},\text{茶梅}\}$};
  \node[result,draw=BookMuted,fill=BookPaper,line width=0.9pt,dotted] (abstain) at (138,-14) {暂缓使用\\弃权};

  \draw[book arrow,knowledge arrow,draw=BookBlue,line width=0.9pt] (q.south west) -- ++(0,-4) -| node[pos=0.74,left,font=\figurefont\scriptsize,text=BookBlue] {$\argmax$} (hard.north);
  \draw[book arrow,knowledge arrow,draw=BookTeal,line width=0.9pt] ([xshift=15mm]q.south west) -- ++(0,-6) -| node[pos=0.76,left,font=\figurefont\scriptsize,text=BookTeal] {完整保留} (soft.north);
  \draw[book arrow,knowledge arrow,draw=BookAmber,dashed,line width=0.9pt] ([xshift=-15mm]q.south east) -- ++(0,-6) -| node[pos=0.76,right,font=\figurefont\scriptsize,text=BookAmber] {累计权重$\geq0.90$} (set.north);
  \draw[book arrow,knowledge arrow,draw=BookMuted,dotted,line width=0.9pt] (q.south east) -- ++(0,-4) -| node[pos=0.74,right,font=\figurefont\scriptsize,text=BookMuted] {$\max q<0.70$} (abstain.north);

  \node[font=\figurefont\scriptsize,text=BookMuted] at (12,-31) {离散决策};
  \node[font=\figurefont\scriptsize,text=BookMuted] at (54,-31) {软标签训练};
  \node[font=\figurefont\scriptsize,text=BookMuted] at (96,-31) {部分排除};
  \node[font=\figurefont\scriptsize,text=BookMuted] at (138,-31) {触发复核};
\end{tikzpicture}

  \caption{同一融合证据的不同预测输出方式。三类别支持分布为人工教学设定；四个分支分别采用最大后验、完整概率、累计概率权重阈值和最低接收阈值，展示硬化、保留不确定性、保留候选与暂缓使用的差异，不赋予该分布额外校准或覆盖保证。}
  \label{fig:knowledge-fusion-output-forms}
\end{figure*}
\FloatBarrier

候选集合除了“供人决策”，还能\textbf{当作一种\emph{训练监督}}用，但这里\textbf{必须\emph{严格区分两种\emph{出身不同}的集合}}，它们的训练含义天差地别：
\begin{itemize}\item 一种是前面按分数\textbf{“\emph{模型自己\emph{筛}出来}的候选}”$C_i$——它\textbf{可能\emph{漏掉}真值}（模型不确定时尤其不可靠），只是个\emph{猜测}。\item 另一种是\textbf{“\emph{已经\emph{独立确认}真值就在其中}”的集合}$S(\bm{x}_i)$——比如经过可靠核查，确知这张照片\textbf{非山茶即茶梅}（但分不清具体哪个），月季被\emph{确定}排除。这是一条\textbf{有\emph{确定性}的事实}。\end{itemize}假设照片已确认属于山茶或茶梅，但无法再细分，完整类别表还包括月季；将这个已确认集合记为$S(\bm{x}_i)$，不同于前面按分数构造的$C_i$。当前监督要求提高集合内总概率，而没有指定山茶和茶梅各占多少。对模型参数$\bm{w}$的分类模型预测输出$p_{\bm{w}}(k\mid\bm{x}_i)$，先在确认集合内求和，再以负对数惩罚总质量过小，得到集合约束损失：
\begin{equation}
  \ell_{\mathrm{set}}(\bm{w};i)
  =-\log\sum_{k\in S(\bm{x}_i)}
    p_{\bm{w}}(k\mid\bm{x}_i).
  \label{eq:knowledge-fusion-set-loss}
\end{equation}
集合内概率越大，负对数损失越小；总质量为零则损失发散。若三类概率为$(0.6,0.3,0.1)$，本例为$-\log(0.6+0.3)=-\log0.9\approx0.10536$。降低这项损失要求减少集合外月季的质量，却不区分集合内两类，也不要求各占一半。把两类分别当确定正类，或指定各一半软标签，都改变了原来的集合事件监督。

理解集合损失的\textbf{语义}：降低它，只要求\textbf{“\emph{别把概率放在集合\emph{外面}那些\emph{已被排除}的类}上}”，是个很\emph{弱、很诚实}的监督——\textbf{它\emph{不指定}山茶、茶梅各占多少，更\emph{不要求}各一半}。若你\textbf{把这两类\emph{分别}当成确定正类}（一会儿说它是山茶、一会儿说它是茶梅），或\textbf{擅自指定一个$0.5/0.5$的软标签}，那都是\textbf{添加了原监督里\emph{根本不存在}的信息}，\textbf{改变了集合事件的本来含义}（“在集合内”并不等于“均匀分布在集合内”）。

还要认清式~\eqref{eq:knowledge-fusion-set-loss}的\textbf{适用边界}：它\textbf{只\emph{使用}了“\emph{已确认}的集合事件”}这一条监督，\textbf{并没有\emph{描述}$C_i$是\emph{怎么生成}的}，所以\textbf{不能把它当成“\emph{任意}候选生成机制”的完整似然}。尤其\textbf{不能拿模型自己筛的$C_i$直接当“\emph{确定}约束”}：因为$C_i$\textbf{可能\emph{漏掉}真值}，一旦真值恰好在集合外，强行最大化集合内概率就会\textbf{逼迫模型\emph{压低}真类概率}，越训越错。对这类\textbf{出身不确定}的候选，应\textbf{保留其“候选”状态}，再\textbf{按它的生成条件}另行规定\textbf{接收、降权、或强制复核}的处理办法，而不是当作铁证。

调整接收阈值还有一个\textbf{极易制造“虚假进步”}的陷阱，核心是\textbf{“阈值改变的是\emph{评价对象（在哪些样本上评分）}”}。设想\textbf{固定模型和分数不动，只\emph{提高}最低接收阈值}：那么\textbf{被接收的样本集合\emph{只会缩小或不变}}（更挑剔了）。直觉上“留下的都是高分、应该更准”，但\textbf{数学上\emph{并不保证}“每一个更高阈值对应的\emph{接收组}里，错误比例都\emph{单调}下降}”——因为\textbf{被阈值卡掉的不一定是错的样本}（可能只是模型对其不够自信的\emph{正确}样本），错误率随阈值的变化取决于校准的具体形状。

更要紧的是\textbf{“\emph{选择性评价}”陷阱}：如果\textbf{困难区域的样本因为达不到高阈值而被\emph{大量排除}}，那么你会看到\textbf{“总体准确率\emph{可喜地}上升}”，但这\textbf{可能完全是靠\emph{收缩服务范围}换来的}——\textbf{系统对最难、往往也最重要}的那部分样本\textbf{直接\emph{拒答}}了，数字好看了、用处却小了。\textbf{用“敢答的简单题”上的高分，掩盖“难题大面积弃权”，是这类系统最常见的自我美化}。因此要同时报告总体与分组覆盖、已接收错误率和集合大小，并把阈值选择样本与最终独立评价分开。

用一组人工计数固定这个比较：简单区与困难区各$100$个样本，阈值$0.60$时分别接收$90,60$个，错$9,18$个，总体已接收错误率为$27/150=18\%$。阈值改为$0.80$后接收$70,20$个，错$4,5$个，总体错误率为$9/90=10\%$。错误率下降了$8$个百分点，但总体覆盖也从$150/200=75\%$降到$90/200=45\%$，困难区覆盖从$60\%$降到$20\%$，已接收错误率为$5/20=25\%$。本例两区错误率都下降，不是反转例子；它说明只报总体错误改善会遗漏困难区的大幅覆盖损失。数字均为教学设定。

\subsection{异构知识的联合使用}
\label{sec:knowledge-fusion-heterogeneous-use}

一张花卉照片的类别、两张视图来自同一对象的关系，以及遮蔽花瓣的恢复内容，不能在同一类别表内计票。它们要求模型预测或保留不同对象。\term{异构监督信号的联合使用}就是保留各类监督对象，让它们通过共享模型参数、多个损失或结构约束共同影响训练；要组织的是联合训练安排，而非强制产生一个融合标签。

要把这些要求交给同一模型，先为每类监督信号定义与其内容匹配的损失，再确定哪些训练单元具有该类监督信号。设本次训练有$n$个单元、$R$种监督信号，共享模型参数为$\bm{w}$，$\ell_r(\bm{w};i)$是监督信号类型$r$的单项损失，$m_{i,r}\in\{0,1\}$表示单元$i$是否适用，$\alpha_r\geq0$控制该项贡献。这里一个单元可以包含单张图像，也可以包含配对视图及其索引，不把关系误作单图类别。我们在各类有效单元上平均损失，再加权形成待最小化的经验目标：
\begin{equation}
  \mathcal{L}(\bm{w})
  =
  \sum_{r=1}^{R}
  \alpha_r
  \frac{\sum_{i=1}^{n}m_{i,r}\ell_r(\bm{w};i)}
  {\max\left(1,\sum_{i=1}^{n}m_{i,r}\right)}.
  \label{eq:knowledge-fusion-heterogeneous-objective}
\end{equation}
没有相应监督信号时，该单元不参与该项求和，不能把缺失内容填成零目标。某类有效数为零时，$\max(1,\cdot)$防止零分母，整项贡献为零，但该类监督信号仍属缺失；这不是用零标签监督模型。逐项平均消除了有效数的直接倍数影响，却未统一不同损失尺度，$\alpha_r$还需结合目标重要性、梯度尺度、采样频率与成本设定。

图~\ref{fig:knowledge-fusion-heterogeneous-training}沿这三种预测输入追踪计算：类别标签与分类预测输出比较，已知视图对应与两份表示比较，被遮蔽内容与恢复预测输出比较。图内$m_r\alpha_r\mathcal{L}_r$简写了各分支的样本选择与贡献，正文目标还包含逐项有效数归一化。三条路径共享表示模型参数，不共享监督语义；不同编码器的向量没有锚点映射时，也不能因维度相同就逐维相加。

\begin{figure*}[htbp]
  \centering
  \begin{tikzpicture}[
  x=0.8mm,y=1mm,
  every node/.style={font=\figurefont\small,text=BookInk},
  input/.style={book node,draw=BookBlue,fill=FigureInput!10,line width=0.9pt,minimum width=22mm,minimum height=10mm,inner sep=2pt},
  model/.style={book node,draw=BookTeal,fill=FigureModel!10,line width=1pt,minimum width=28mm,minimum height=15mm,inner sep=2pt},
  branch/.style={book node,draw=BookTeal,fill=FigureModel!7,line width=0.8pt,minimum width=27mm,minimum height=9mm,inner sep=2pt},
  knowledge/.style={book node,draw=BookAmber,fill=BookAmber!9,line width=0.8pt,minimum width=45mm,minimum height=10mm,inner sep=0pt},
  knowledge text/.style={font=\figurefont\small,anchor=west,align=left,inner sep=0pt},
  edge label/.style={font=\figurefont\scriptsize,text=BookAmber,inner sep=0pt},
  loss/.style={book node,draw=BookInk,fill=BookPaper,line width=0.8pt,minimum width=31mm,minimum height=10mm,inner sep=2pt}
]
  \node[input] (image) at (0,0) {图像与增强视图};
  \node[model] (encoder) at (37,0) {共享表示模型\\$f_{\bm{w}}$};
  \draw[book arrow,knowledge arrow,draw=BookBlue,line width=1pt] (image) -- (encoder);

  \node[branch] (pred) at (78,32) {分类预测};
  \node[branch] (compare) at (78,0) {视图表示比较};
  \node[branch] (recover) at (78,-32) {局部内容重建};
  \node[loss] (l1) at (130,32) {$m_1\alpha_1\mathcal{L}_{\mathrm{cls}}$};
  \node[loss] (l2) at (130,0) {$m_2\alpha_2\mathcal{L}_{\mathrm{rel}}$};
  \node[loss] (l3) at (130,-32) {$m_3\alpha_3\mathcal{L}_{\mathrm{rec}}$};
  \node[knowledge] (label) at (105,52) {};
  \node[knowledge] (pair) at (105,20) {};
  \node[knowledge] (patch) at (105,-52) {};
  % Reserve the left 13 mm for icons; text uses the remaining column.
  \node[knowledge text] at ([xshift=13mm]label.west) {外部任务类别标签};
  \node[knowledge text] at ([xshift=13mm]pair.west) {外部视图对应关系};
  \node[knowledge text] at ([xshift=13mm]patch.west) {外部局部恢复目标};
  \node[rectangle,draw=BookAmber,fill=white,line width=.8pt,minimum size=5mm,inner sep=0pt,font=\figurefont\bfseries\scriptsize] at ([xshift=6.5mm]label.west) {A};
  \node[circle,draw=BookAmber,fill=white,line width=.8pt,minimum size=3.5mm,inner sep=0pt] (pair-left) at ([xshift=4mm]pair.west) {};
  \node[circle,draw=BookAmber,fill=white,line width=.8pt,minimum size=3.5mm,inner sep=0pt] (pair-right) at ([xshift=9mm]pair.west) {};
  \draw[BookAmber,line width=.8pt] (pair-left) -- (pair-right);
  \node[rectangle,draw=BookAmber,fill=white,line width=.8pt,minimum size=6mm,inner sep=0pt] (patch-icon) at ([xshift=6.5mm]patch.west) {};
  \draw[BookAmber,line width=.5pt] ([xshift=-3mm]patch-icon.center) -- ([xshift=3mm]patch-icon.center);
  \draw[BookAmber,line width=.5pt] ([yshift=-3mm]patch-icon.center) -- ([yshift=3mm]patch-icon.center);
  \fill[BookAmber!45] (patch-icon.center) rectangle ([xshift=3mm,yshift=3mm]patch-icon.center);
  \node[loss,draw=BookTeal,fill=FigureModel!12,line width=1pt,minimum width=20mm] (sum) at (180,0) {总损失};

  \draw[book arrow,knowledge arrow,draw=BookTeal] (encoder.east) -- ++(3,0) |- (pred.west);
  \draw[book arrow,knowledge arrow,draw=BookTeal] (encoder.east) -- (compare.west);
  \draw[book arrow,knowledge arrow,draw=BookTeal] (encoder.east) -- ++(3,0) |- (recover.west);
  \draw[book arrow,knowledge arrow,draw=BookTeal] (pred) -- node[below,font=\figurefont\scriptsize,text=BookTeal] {预测输出} (l1);
  \draw[book arrow,knowledge arrow,draw=BookTeal] (compare) -- node[below,font=\figurefont\scriptsize,text=BookTeal] {预测输出} (l2);
  \draw[book arrow,knowledge arrow,draw=BookTeal] (recover) -- node[above,font=\figurefont\scriptsize,text=BookTeal] {预测输出} (l3);
  \draw[book arrow,knowledge arrow,draw=BookAmber] (label.south) -- ++(0,-5) -| (l1.north);
  \draw[book arrow,knowledge arrow,draw=BookAmber] (pair.south) -- ++(0,-5) -| (l2.north);
  \draw[book arrow,knowledge arrow,draw=BookAmber] (patch.north) -- ++(0,5) -| (l3.south);
  \node[edge label,anchor=south,yshift=1mm] at ($(label.south)!0.5!(l1.north)$) {监督目标};
  \node[edge label,anchor=south,yshift=1mm] at ($(pair.south)!0.5!(l2.north)$) {监督关系};
  \node[edge label,anchor=north,yshift=-1mm] at ($(patch.north)!0.5!(l3.south)$) {监督目标};
  \draw[book arrow,knowledge arrow,draw=BookTeal,line width=1pt] (l1.east) -| (sum.north);
  \draw[book arrow,knowledge arrow,draw=BookTeal,line width=1pt] (l2.east) -- (sum.west);
  \draw[book arrow,knowledge arrow,draw=BookTeal,line width=1pt] (l3.east) -| (sum.south);

  \node[anchor=north,font=\figurefont\scriptsize,text=BookMuted,align=center] at (130,-61) {掩码$m_r$在损失处选择具有相应标注的样本};
\end{tikzpicture}

  \caption{异构监督信号形成联合训练约束。图像经共享表示模型进入三个损失分支；任务标签、视图关系和局部恢复目标分别提供不同监督对象，经适用样本掩码与权重后组成总训练目标。共享训练不要求把三类监督信号压缩成同一类别标签。}
  \label{fig:knowledge-fusion-heterogeneous-training}
\end{figure*}
\FloatBarrier

这些损失通过同一模型参数更新相互影响，改权重会改变合成梯度，若同时改采样或训练安排，还会改变资源分配。需要分别检查适用样本、尺度、采样比例和单项贡献；梯度方向长期冲突时，可以调整权重、交替批次、分阶段训练或减少共享层，但没有一种处理自动保证全部任务改善。各项仍须有独立验证指标，总损失下降不能定位哪种监督信号有用。

\subsection{追加来源与跨来源预算的分配}
\label{sec:knowledge-fusion-source-budget}

融合之后常会遇到“\textbf{证据还不够、要不要\emph{再}花钱追加判断}”的决策，比如$x_5$上规则和人工吵得不可开交。直觉上“\textbf{哪个\emph{便宜}先问哪个}”，但这\textbf{忽略了关键一点：\textbf{最便宜的那一次调用，可能\emph{根本不增加}独立判断依据}}。

用6.2的眼光看几个选项就很清楚：\textbf{把同一个\emph{确定性}模型、在\emph{完全相同}的输入上\emph{再跑一遍}}——它\textbf{给出的必然是\emph{同一个答案}（纯复制，零新信息）}，这种钱\textbf{纯属浪费}；而\textbf{追加一位\emph{独立}专家}，他\textbf{可能\emph{改变}候选排序}（带来真正新证据，但单价高）；\textbf{补充一条新规则、或新采一批样本}，又有\textbf{各自\emph{不同}的\emph{覆盖范围}}（可能对当前这张帮助不大、但对一批同类样本有长期价值）。所以\textbf{跨来源预算分配的正确依据}，不是单次价格，而是：\textbf{在“\emph{当前这个\emph{证据缺口}}”下，这一行动预期能\emph{减少多少类别不确定性}、或带来多少\emph{任务收益}}（信息增益的思路），\textbf{再与其成本权衡}——\textbf{贵但能一锤定音的独立判断，可能远\emph{值}过十次免费的复制}。

比较之前先规定允许的追加行动集合$A$，包括可执行的来源调用或开发方案。一个行动可以是单次调用，也可以是已经明确调用步骤和停止条件的组合方案；比较范围取决于$A$实际纳入了哪些选择。固定融合配置，将当前样本标注集合记为$\mathcal{D}$，行动$a$可能返回$Z_a$，$\mathcal{D}\cup Z_a$表示取得新结果后的样本标注集合；组合方案的$Z_a$包含它取得的全部新增样本标注。任务效用$U$越高越好，例如可以取既定决策损失的负期望，并固定评价对象与损失规则。结果尚未知，所以按$Z_a\mid\mathcal{D}$的预测分布平均追加后的效用，再减去当前效用，衡量期望增量。人工时间、费用、计算量与反馈时延单位不同，分别记成本$C_j(a)$和预算$B_j$，$j=1,\ldots,J$；成本须覆盖整个行动或方案。未给换算规则时，只在每项资源允许的行动中最大化增量：
\begin{equation}
  a^\star
  \in\argmax_{a\in A}
  \left\{
    \mathbb{E}_{Z_a\mid\mathcal{D}}
    [U(\mathcal{D}\cup Z_a)]
    -U(\mathcal{D})
  \right\}
  \quad
  \text{s.t.}\quad
  C_j(a)\leq B_j,\ j=1,\ldots,J.
  \label{eq:knowledge-fusion-source-budget}
\end{equation}
这个框架有\textbf{一个\emph{量纲}设计上的讲究}，是它区别于“拍脑袋性价比”的关键：\textbf{目标函数里只有“\emph{任务效用}”这一种单位}（最大化期望效用增量），而\textbf{每一种\emph{资源}约束\emph{各自}保留\emph{自己}的单位、\emph{分开}写在约束里}（时间用小时、金钱用元、计算用量用算力单位，各比各的预算$B_j$）。之所以\textbf{不把“$-0.3$元”这种项直接塞进效用目标里}，是因为\textbf{小时、元、效用}\textbf{三者\emph{量纲不同}，\emph{没有}天然的换算率}，\textbf{不能未经声明就直接相减}（“减0.3元”等于减多少效用？没人知道）。当然，业务上\textbf{常常确实想把成本折算成钱、和收益合并成一个“\emph{净效用}”来直接选最大}（比如商业系统要算账）。这时\textbf{不能偷偷折算}，必须\textbf{先\emph{显式}规定一组换算系数}$\beta_j$——它的含义是“\textbf{每一单位资源$j$，\emph{在决策者眼里值多少个效用单位}}”（比如人工一小时值多少元、延迟一秒钟值多少效用）。用$\beta_jC_j(a)$把成本\textbf{换算到与$U$\emph{同量纲}}之后，才能合法地写出“期望效用增量$-$$\sum_j\beta_jC_j(a)$”这种净效用目标。这组$\beta_j$\textbf{本质是一个\emph{业务价值判断/偏好}}，应由决策者明确给出、并承担其合理性，而不是分析师替他假设。

最后给整个公式\textbf{去魅}：要真正\textbf{算出}“期望效用增量”，你得\textbf{先知道“行动$a$\emph{可能返回什么}$Z_a$”的\emph{预测分布}}（不查怎么知道独立专家会说什么？），这又\textbf{需要历史\emph{审计数据}、\emph{覆盖}模型、\emph{来源依赖}模型、或\emph{明确的业务取舍}}来支撑——这些输入\textbf{往往本身就难估}。所以\textbf{式~\eqref{eq:knowledge-fusion-source-budget}的价值在于提供了一个“\emph{该\emph{比较什么}}”的\emph{决策框架}}（逼你列出行动、收益、各维成本、预算），它\textbf{\emph{不会}凭空替你算出“真实最优动作”}——喂给它的分布和换算系数不可靠，最优也就是纸上谈兵。

回到表~\ref{tab:knowledge-fusion-source-records}的$x_5$具体决策，三种行动的账一目了然：\textbf{在\emph{相同照片、相同规则版本}上\emph{重跑}一遍确定性规则$R_1$}——结果必然还是山茶，\textbf{期望信息增量为零，几乎不该花这个钱}；\textbf{追加一位\emph{独立}人工判断}——有机会直接核对“规则组 vs 人工”的异议，\textbf{单点信息价值高}（虽然单价贵）；\textbf{开发一条\emph{新规则}}——\textbf{当前这一张}要等规则开发完才有答案（慢），但一旦建成，\textbf{之后可能\emph{批量}覆盖大量同类照片}（长期摊薄成本）。三者性质不同，正对应“复制 / 独立即时证据 / 有复用价值的建设投入”。

还要\textbf{把\emph{真实}成本算全}，别只盯着调用单价：\textbf{新来源/规则\emph{长期维护}的费用、调用的\emph{时延}、为保证质量做的\emph{审计}、以及\emph{失败重试}}都要计入预算。两个常见错觉要点破：\textbf{“单次调用价格低”\emph{不等于}“\emph{有效标注}的\emph{单位成本}低”}（便宜来源若需要大量冗余调用、清洗、复核才得到一条可用证据，摊下来未必便宜）；\textbf{“能\emph{批量复用}}”\textbf{\emph{也不能代替}\emph{独立的质量验证}}”（用得再广，也得先证明它这次提供的是独立、可靠的证据，否则只是规模化地复制错误）。

另一个关键决策是“\textbf{什么时候\emph{停止}追加来源}”，但停止结论\textbf{严格\emph{相对于你所考虑的行动集合}}$A$而言，不能夸大。\textbf{如果你的$A$\emph{只}列了“\emph{单独}再调用某一个来源”这种\emph{单步}行动}，然后你\textbf{发现每一个\emph{单独}调用的期望增量都$\leq0$}，于是停手——这\textbf{只\emph{支持}“按这种\emph{单步贪心}策略可以停”}，它\textbf{\emph{没有}排除“\emph{两个或多个}来源\emph{组合}起来”可能有巨大\emph{互补}收益}。用一个极简例子\textbf{无可辩驳地}证明“\emph{单步无收益}并不等于\emph{组合无收益}”。设要判断“\textbf{两个独立的$0/1$读数\emph{是否\emph{不同}}}”（异或任务），两个数各自等概率取$0$或$1$。
\begin{itemize}\item \textbf{什么都不读}：两数相同/不同各半，最优瞎猜错误率$1/2$。\item \textbf{只读\emph{第一个}数}：知道$a$后，$b$仍是$0/1$各半，于是“$a,b$是否不同”\textbf{依然只能瞎猜}，最小零一损失\textbf{还是$1/2$}——\textbf{读第一个数的（单步）信息增量为\emph{零}}！\item 同理\textbf{单独读第二个数，增量也是零}（对称）。\item 但\textbf{\emph{读齐}两个数}：一旦$a,b$都知道，是否不同\textbf{完全确定}，错误率\textbf{降到$0$}。\end{itemize}
于是在允许读两次的预算下：\textbf{两个行动\emph{各自单步}的边际增量都是零，但\emph{合在一起}却带来了\emph{从$1/2$到$0$}的全部收益}。这就是\textbf{来源间\emph{互补}（组合）价值}：\textbf{每个单独的碎片}都\textbf{不充分}，只有\textbf{凑齐}才揭晓答案。所以用“逐个问、每个都没用”来论证“别再问了”，是典型的短视。

那怎么\textbf{正确地}决定停止？如果你\textbf{预期到存在这种互补、或存在“\emph{根据上一步返回结果}\emph{再决定}问什么”的\emph{自适应}查询}（先问一个便宜的、看结果再决定要不要问贵的），就必须\textbf{把这些\emph{多步/组合/自适应}方案\emph{整体}作为一个行动\emph{纳入}$A$}，然后\textbf{按“\emph{整个方案}”的返回分布和\emph{整个方案}的总成本}去评估，而不是逐段短视地算。以分项预算判断，在已经纳入的方案中都没有正预期增量，或没有方案满足预算，才支持相对于这一范围停止；若已给成本换算，则比较相应的净增量。证据达到接收标准后进入使用，关键对象无法自动解决时转人工，预算不足而证据仍不足时保留未决并由使用端弃权。小批试用可先估计新来源的互补性，再决定扩大投入；人工内部如何安排新增与重复标注，仍由相应获取策略处理。

最后把表~\ref{tab:knowledge-fusion-source-records}的$x_2$写成可交给质量闭环的标注版本清单K17。该行的$R_1,R_2$同模板且报山茶，$M$报茶梅，$H$未请求，对象和物种语义已对齐；K17固定采用表~\ref{tab:knowledge-fusion-comparison}的等权相对多数，唯一最高票选类，平票或零票未决。因此山茶以$2:1$获选，却与独立参考茶梅不符。为定位这个错误，样本标注要保留原始类别票或类别概率、共享关系、缺失原因、票数及决策配置；$2/3$仅是有效票中的山茶比例，不是后验。

K17虽然保存了共享关系，计票仍把两条规则分别当一票，标注字段完整并不等于规则已经利用依赖。这一行可用于诊断与选择修正，但已看过的$x_2$参考不能再次充当最终留出证据。第~\ref{chap:knowledge-quality-assessment-correction}章将接续诊断形成K18配置，并在第~\ref{sec:knowledge-quality-assessment-correction-repair-validation}节分别检查已改结果与冻结策略的独立留出质量。本章的一次错误和教学权重只揭示待核查机制，不足以决定新版本是否发布。

\subsection{融合结果的验证}
\label{sec:knowledge-fusion-validation}

\textbf{验证的第一原则是“\emph{融合发生在\emph{哪一层}，就先查那一层特有的风险}}”，因为不同层的典型故障完全不同，不能用同一张检查表：
\begin{itemize}\item \textbf{样本标注层融合}（投票/概率）：回到源头，核对\textbf{来源是否可靠、对象/语义是否对齐、“不投票”状态有没有被误当类别}，并用\textbf{独立参考}去检查综合出的标签和概率对不对。\item \textbf{模块层融合}：重点核对\textbf{接口（维度/顺序/尺度）和\emph{对象绑定}}（别张冠李戴），且必须\textbf{测试\emph{整条}连接路径}，不能只单测每个模块。\item \textbf{模型参数融合}：核对\textbf{底座版本是否准确、坐标\emph{对应}}有没有建立、用的是哪条\emph{运算规则}、系数怎么来的。\item \textbf{蒸馏}：核对\textbf{查询输入的\emph{覆盖范围}}、教师是否\emph{固定}、学生实际学的是\emph{哪个信号}。\end{itemize}
还要分清一个界限：\textbf{“计算过程\emph{符合约定}}（流程跑通、公式没写错）”只说明你\emph{正确实施}了所声明的方法，它是\textbf{\emph{必要但不充分}}的——\textbf{最终还必须在\emph{任务标签和明确的成功标准}下检验结果对任务\emph{有没有用}}，“跑通了”绝不等于“有效”。

做方法对比时，要\textbf{先把\emph{目标、可用材料、预算}三者\emph{固定}住}（否则不公平），再\textbf{设置两个\emph{必须可执行}的基线}：一个是\textbf{最好的\emph{单}来源}（不融合能到哪），一个是\textbf{干脆\emph{不融合}}（原始状态），没有基线就说不清融合有没有净收益。当多个模型\textbf{承担\emph{不同}任务}时，不能只报一个混合分，而应\textbf{分别报告“\emph{原任务能力\emph{保持}}了多少”和“\emph{联合请求}表现如何”}；声称“\emph{未见组合}能力”时，还要像6.5那样\textbf{专门留出“成分见过、组合没见过”的成分组合测试}，\textbf{不能把各来源单任务的正确率\emph{相加}就当作组合能力}。成本归因也要诚实：\textbf{运行时组合（ensemble）通常\emph{继续承担}所有来源的调用开销}，而\textbf{参数融合和蒸馏是把\emph{一部分}成本\emph{转移到准备阶段}}（选来源、查教师、尝试失败候选、多轮重复训练的费用都要计入总账），部署省了不等于全周期省了。

本章只讲了“\emph{机制}”和“\emph{直接}验证条件”，一套完整体系还需要其他章节接力：对\textbf{已有学习依据的\emph{审计与修正}}（怎么系统地发现并纠正像$x_2$那样的错、怎么防止改了A处坏B处）由第~\ref{chap:knowledge-quality-assessment-correction}章展开；\textbf{来源\emph{选择}}（值不值得复用）、\textbf{长期\emph{维护}}（改一处如何不毁别处）、\textbf{跨任务\emph{评价}}分别由第~\ref{chap:ket-reuse-selection}、\ref{chap:ket-update-maintenance}和\ref{chap:ket-evaluation}章补充。请以一句\textbf{不过度承诺}的话收尾本章：\textbf{一次融合操作的完成，\emph{绝不等于}已经自动获得了“可靠的融合标签”或“可靠的预测模型”}——可靠性是\emph{验证}出来的，不是融合这个动作本身赋予的。

\section{本章小结}
\label{sec:knowledge-fusion-summary}
\label{sec:ket-integration-summary}

回到章首那张花卉照片上互相矛盾的一堆答案，现在我们有一张\textbf{按“证据\emph{从哪一层进入计算}”组织的完整地图}，这是理解全章最值得带走的骨架。人工标签、规则的类别票、已有模型，虽然都提供信息，但它们\textbf{进入后续计算的\emph{位置}完全不同}，由此对应\textbf{四条不可混为一谈的路线}：
\begin{itemize}\item \textbf{样本标注/预测输出层}：\emph{综合的是“回答\emph{同一个问题}”的证据}——投票、加权、概率推断，都在“票/最终答案”这一最表层做文章。\item \textbf{计算模块层}：让\textbf{中间表示、属性、子步骤\emph{继续参与}后续计算}（解决“最终答案里已被压没”的关系型问题）。\item \textbf{模型参数层}：\textbf{在\emph{做任何预测之前}，离线构造出一份新参数}（只许部署一份模型时的选择）。\item \textbf{蒸馏层}：\textbf{把已有模型的成果，\emph{转化成}训练一个新模型（学生）所需的\emph{学习依据}}。\end{itemize}
尤其再强调一次蒸馏的易混点：教师的\textbf{预测输出可以\emph{成为}学生训练样本里的\emph{监督目标}}，但\textbf{教师的“预测过程”和学生的“训练过程”是\emph{两个独立的学习算法}，各自训练自己的模型}，不是同一次映射。

在“标签层”这条路上，\textbf{来源建模、三层对齐、概率推断}三件事环环相扣，目的都是让\textbf{每一张票、每一个权重、每一个似然都有\emph{可追溯}的含义}，而不是一堆无出处的数字。全章反复出现的几条纪律，结尾值得再钉一遍：
\begin{itemize}\item \textbf{“列数”不等于“证据数”}：重复模板不会因为多记一列就变成独立证据（$R_1,R_2$合起来是一份）；\textbf{“没有票”也不等于“负类”}（弃权/失败/未请求要分开记，不能当成“不是山茶”）。\item \textbf{三件事是\emph{三个}问题，互不蕴含}：概率模型的\textbf{可识别}（数据原则上能否唯一区分参数，还得除类别置换）、EM 的\textbf{数值收敛}（迭代停没停稳）、以及\textbf{答案\emph{正确}}（推断的类别符不符合现实）——\textbf{各需各的条件，收敛/稳定都不保证另外两个}。\item \textbf{两种集合不能混用}：“模型自己\emph{筛}出的候选$C_i$（可能漏真值）”与“已经\emph{确认}真值在其中的集合$S$（确定性事实）”是两种监督，后者才能当硬集合约束。\item \textbf{噪声校正的边界}：保留观测标签时，一张\emph{准确已知}的报错通道$\bm T$可用于\emph{前向}或\emph{后向}损失校正，理想情形下能恢复干净风险/真类后验；但它\textbf{既\emph{不会自动}把通道$\bm T$识别出来}（要锚点），\textbf{也\emph{没有逐条}把噪标签修好}（单条校正损失甚至可能为负）。\end{itemize}

在“模型层”这几条路线上，反复出现的主题是\textbf{“接口/坐标/监督\emph{对得上}”是一切融合的前提，且\textbf{表面的相似常常骗人}}：
\begin{itemize}\item \textbf{集成（输出层）}：收益来自\textbf{错误\emph{互补}}（回归方差按$1/K$下降），但\textbf{错误\emph{共同}}（$\rho=1$时毫无收益）、共同偏差、以及\textbf{错误路由}平均不掉，还照样留下。\item \textbf{参数融合}：\textbf{同一个函数可以有\emph{无数种参数化}}（隐藏单元可置换），所以“功能相同”的模型\textbf{直接逐坐标平均可能\emph{完全失效}}（那个平均出恒零的$|x|$例子）；而且\textbf{低秩\emph{因子}组合\emph{不等于}完整增量相加}（有系数平方、交叉项，还可能压低秩）。\item \textbf{蒸馏}：它的最大优点是\textbf{\emph{不要求}师生参数坐标相同}（跨结构也能传能力），但仍受\textbf{查询输入\emph{覆盖}}（没查过的区域学不到）、\textbf{教师偏差}（会连同错误一起传）、\textbf{学生容量}（装不下的学不会）三重约束。\item \textbf{样本凝练}：通过匹配\textbf{训练后损失、局部梯度、或较长一段参数移动}，把训练信号浓缩进很小的合成集；但它\textbf{费用提前且昂贵、会继承原集缺陷、且常与特定结构绑定}，是否划算取决于复用次数和后续协议。\end{itemize}
一句话戳破三类常见的“数字迷信”：\textbf{“融成了\emph{单份}参数”、“无标签选系数时\emph{熵很低}”、“蒸馏/凝练的\emph{匹配误差很小}”，这三者\emph{都不能代替}“在真实任务上\emph{能力是否保持/提升}”的\emph{独立验证}}——单份不代表无损、低熵可能是自信地全错、匹配好不等于泛化好。

把全章收成一个可操作的\textbf{决策清单}：无论选哪条融合路线，都必须\textbf{同时说清五件事}——
\begin{itemize}\item \textbf{保留\emph{什么}}：希望最终保住哪些监督信号、哪些中间表示特征、或哪些预测行为（不能只说“融一下”）；\item \textbf{能访问\emph{什么}}：允许使用哪些样本标注、哪些来源模型、哪些计算模块（数据权限、能否读参数/激活、能否查教师）；\item \textbf{在\emph{哪里}汇合}：输出层、模块层、参数层、还是蒸馏监督（对象不同，机制不同）；\item \emph{花多少成本}：来源调用、准备阶段训练、存储、查询、维护，分别计入；\item \emph{结果以什么形态}交付、如何使用。\end{itemize}
特别要记住\textbf{融合结果\emph{不必}只剩一条答案或一个模型}：它可以是\textbf{标签、完整概率分布、候选集合、仍在运行的组合系统、或一份新模型/学生}；并且一个成熟的系统\textbf{必须\emph{保留}“未决、弃权、拒绝候选”这些诚实的出口}，而不是逼系统对没把握的事硬下结论。最后，判断“多来源到底\emph{有没有}带来预期收益”，\textbf{不能靠来源数量、票数或算法名字}，只能靠\textbf{在\emph{当前目标}下，对“来源能力、组合行为、资源代价”做\emph{独立}检查}。带着这份纪律，下一章我们将进入“如何系统地审计、评估并纠正这些学习依据”的质量闭环。

\chapter{知识的质量评估与修正}
\label{chap:knowledge-quality-assessment-correction}

前面六章都在讲“\textbf{怎么\emph{取得和融合}知识}”，本章要回答一个更尖锐、却最容易被乐观跳过的问题：\textbf{这些辛苦得到的标签、关系、训练目标，\emph{到底可不可靠}？发现了错又该\emph{怎么修、修完怎么证明确实变好了}}？先用一个贯穿全章的思想实验，把三个\textbf{最常被混淆}的证据概念分开。

设想你有\textbf{一万条}候选标签。一个“风险检测程序”按某条规则（比如模型和标签分歧大）\textbf{挑出其中\emph{一百条}最可疑的}，人工复核后\textbf{确认其中\emph{四十条}确实错了}。面对这个结果，很多人会脱口而出“\textbf{错误率是$40\%$}”，甚至进一步说“整批数据有$40\%$的错”。这两句话\textbf{都不严谨}，请逐句辨析：
\begin{itemize}\item 你\textbf{真正能说的}只有：\textbf{“在\emph{被选中}的这$100$条里，确认了$40$条错”}，即\textbf{入选样本中的命中率（命中$40\%$）}。它衡量的是“这个检测器\emph{挑得准不准}”。\item 你\textbf{\emph{不能}说的}是：\textbf{其余$9900$条的错误比例}——它们\textbf{根本没按同样条件被检查}（它们是“没被风险规则挑中”的，不代表它们就对），\textbf{整批一万条的总体错误率仍然未知}。\end{itemize}
第二个区别：\textbf{“检测分数高”不是“这四十条有错”的\emph{依据}}。检测器只是\textbf{喊了一声“这一百条可疑”}，\textbf{真正判定其中四十条错的，是\emph{复核时拿到的独立证据}}（原图、规范、采集日志）。\textbf{线索（signal）}和\textbf{裁决（verdict）}是两件事。

第三个区别出现在“修”之后：就算你把这四十条\textbf{全部改对了}，也\textbf{还不能收工}——必须\textbf{另问一句}：“\textbf{我用来修的这条规则，在\emph{没参与规则设计}的新对象上还成立吗}？”在这一百条上调出来的修正规则，可能只是“\emph{专为这一百条量身定做}”，换一批就失效。

这三个区别\textbf{正是全章的三条主线}，先记住这三句“\textbf{什么\emph{不等于}什么}”：\textbf{重点排查（挑可疑的）\emph{不等于}总体审计（科学抽样推全体）；检测线索\emph{不等于}错误证明（要独立证据裁决）；“改好的结果被复核确认”\emph{不等于}“修正策略在\emph{留出}新数据上有效}”。把这三件事分清，质量改进里的证据才不会被夸大。

再明确\textbf{本章要质检的“对象”有多杂}——它们来自前几章，形态完全不同，因此\textbf{不能用一把尺子（比如一个分类正确率）量到底}：
\begin{itemize}\item \textbf{任务答案}（第2–4章的类别、数值标签）：这类\textbf{有明确正确答案}，原则上\textbf{可以和独立参考直接比较}对错。\item \textbf{局部结构/关系}（实体框、样本配对等，第5章和第6章）：要\textbf{分别核验两端对象是谁、关系方向对不对、成立条件满不满足}，不是一个“对/错”字能概括。\item \textbf{训练中\emph{生成}的表示目标}（自监督目标、教师软标签等）：它们\textbf{通常\emph{没有}“逐向量标准答案”}——你无法指着神经网络中间某一维说它该等于几；对它们只能\textbf{拆成两件事查}：\textbf{“\emph{构造过程}守没守约定”}（输入对不对、教师版本、变换规则）和\textbf{“\emph{下游用途}有没有用”}（拿它训练后任务有没有变好）。\end{itemize}
由此自然得到本章的\textbf{两步走逻辑}：\textbf{先确定“要核验的、\emph{可核验}的属性是什么”}——属性没定义清楚，从有限样本算出的任何百分比都没有意义（分母、对错口径都是糊的）；\textbf{再区分“我这是在\emph{发现问题}（挑出具体错例）还是在\emph{估计总体}（推断没查的那一大批）”}——目的不同，抽样方法、修正范围、验证方式全都不同。

本章\textbf{沿用第~\ref{chap:knowledge-fusion}章那张有争议的照片$x_2$作为\emph{连贯的教学案例}}，先把案情再固定一遍：两条\emph{共享同一本旧词典}的规则都报“山茶”，模型报“茶梅”，人工当时尚未请求。在标注版本\textbf{K17}里，规则是“每个有效来源一票、唯一最高票获胜、平票或无票记未决”，于是$x_2$以\textbf{$2:1$被判为山茶}。

但第6章我们已经埋下伏笔：这两票其实\emph{同源}，可能只是一份证据被算了两次。本章的关键动作是——\textbf{我们\emph{不是}因为“来源意见不一致”就凭感觉改标}，而是\textbf{回到\emph{照片本身}和\emph{采集日志}去查}，发现独立证据其实支持“茶梅”，并顺藤摸瓜追到“\textbf{共享词典失效 + 重复计票}”这两个\textbf{系统性根因}。这正示范了“\emph{用证据定位问题，而非用票数压人}”。

还要特别强调这张照片在\textbf{开发流程}里的\textbf{身份}：$x_2$和它的近重复样本\textbf{只用于“讲解 + 开发诊断”}（用它们发现词典问题、设计修正策略），\textbf{\emph{绝不}进入最终的“留出验收总体”}。为什么这么严格？因为\textbf{一旦你\emph{看过}某条数据、又围绕它\emph{设计和调整}了修正规则，它就\emph{失去了当裁判}的资格}——用它再去验收，等于用“做过的考题”报成绩（这条纪律在7.4会反复用到）。

错误的“长相”远不止“标错类别”一种。同一批数据里还可能：\textbf{漏标}（照片里明明有两朵花只标了一朵）、\textbf{对象误配}（把两株不同植物错绑成同一对象）、\textbf{视图错误}（裁剪后只剩背景，却被当成“保留花朵语义”的正样本）。要让这些五花八门的问题\textbf{事后查得清、责任理得顺}，每一份审计报告都必须\textbf{能往回\emph{连接}到五样东西}：\textbf{标注版本、来源路径、抽样设计、参考证据、裁决理由}——这就是“质量\emph{可追溯}”的最低要求。

这里要\textbf{引入第一组需要分清的状态词}，它们\textbf{属于\emph{两个不同的层}，千万不要混用}。
\begin{itemize}\item \textbf{处理阶段}（流程走到哪了）：一条线索要么是\term{待复核线索}（review lead）——\textbf{只表示“某项检查\emph{提出了怀疑}”，\emph{还没}下质量结论}（它是个待办、一个问号，不是定罪）；要么是“\textbf{已完成}”（复核流程走完了）。\item \textbf{质量结论}（查完到底怎样）：\textbf{正确、错误、部分正确、缺失、不适用、无法裁决}。\end{itemize}
两个\textbf{反直觉但重要}的点：其一，\textbf{“处理已完成”\emph{完全可以}得出“无法裁决”的结论}——流程走完了不等于一定能判对错；其二，\textbf{“无法裁决”必须\emph{写明原因}并保留}，常见原因是“\textbf{观测/证据不足}”（照片太糊、无参考）或“\textbf{规范本身有歧义}”（两种标注按现行规范都成立）。\textbf{把无法裁决偷偷改成“正确”或“错误”，都会让质量指标系统性失真}。

\section{质量评估依据的建立}
\label{sec:knowledge-quality-assessment-correction-basis}

\textbf{“算错误率”看似只是数数，其实在“数”之前有三件事\emph{必须先定下来}}，否则同一个数据能算出互相矛盾的“质量”：
\begin{enumerate}\item \textbf{检查\emph{什么属性}}：这次到底查\textbf{物种对不对、框的位置准不准、还是配对关系成不成立}？这几者是\emph{不同的问题}，混在一起数，分母分子都对不上。\item \textbf{用什么材料\emph{裁决}}：凭什么说一条标签错？必须\textbf{提前准备好能下判断的材料}（原图、规范、独立复标等），不能边查边改口裁判规则。\item \textbf{比例\emph{代表谁}}（抽样范围）：你数出的这个百分比，\textbf{说的是“这一百条”、还是“常见区”、还是“全部一万条”}？范围不划定，一个百分数就没有归属。\end{enumerate}
这三者合在一起，构成一份\term{审计方案}（audit plan）——也就是对“\textbf{这次审计\emph{查什么}、\emph{怎么判}、\emph{怎么选样}}”的\textbf{事先书面约定}。它的意义在于\textbf{可复现、可追溯}：只有当“检查标准、裁决依据、抽样设计”三者\textbf{彼此对应、且事先冻结}，事后把\textbf{一条条裁决\emph{汇总}起来}，才能得到一个\textbf{“对明确用途有效”的质量结论}；否则只是一堆各说各话的对错。

\subsection{评估标准的确定}
\label{sec:knowledge-quality-assessment-correction-criteria}

先破除一个想当然：\textbf{“标签栏里填了一个合法物种名”，\emph{远不等于}它满足了用途要求}。一个格式完全正确、拼写无误的“山茶”，可能\textbf{判错了物、可能漏标了同框的另一株花、可能只细到属而任务需要种、还可能用了已废止的旧分类}。所以需要\term{评估标准}（assessment criteria）来\textbf{把模糊的“质量好不好”翻译成“复核者\emph{能逐条照着执行}的具体判断”}。它至少要回答四类问题：
\begin{itemize}\item \textbf{类别}是否\textbf{符合\emph{任务定义}}（任务问物种就不能用属名充数）；\item \textbf{需要的对象是否\emph{齐全}}（该标的花有没有漏）；\item \textbf{答案是否\emph{细到规定层次}}（要亚种就不能停在种）；\item \textbf{事实和来源版本是否\emph{适用于目标时间}}（旧规范、过期名录不能用于新任务）。\end{itemize}
这里有\textbf{四个层次\emph{必须分开记录、不能互相顶替}}，因为它们可能一好一坏：\textbf{格式合规}（能不能被系统解析）、\textbf{证据充分}（有没有足够材料下判断）、\textbf{语义正确}（内容到底对不对）、\textbf{下游效果}（拿去训练有没有用）。请\textbf{记住两组反例}：\textbf{能被解析的标签仍可能\emph{错类}}（格式对、内容错）；\textbf{证据很充分的标签也可能据此\emph{判错}}（材料齐全但结论就是错）。至于“\textbf{与参考一致的标签\emph{能否改善训练}}”，那更是\textbf{另一个问题}，必须\textbf{在指定的训练流程和评价样本上\emph{另测}}，标签正确不自动等于训练有用。

\textbf{属性定了，还要定“\emph{一次裁决的\emph{单位}}是什么”}——单位选错，对错就没法一致地记。不同监督信号的裁决单位差别很大：
\begin{itemize}\item \textbf{单条类别标签}：通常可以\textbf{整体记一次}对/错。\item \textbf{对象检测框}：\textbf{不能整框一刀切}，要\textbf{拆成三件分别查}——\textbf{对象有没有\emph{遗漏}}（该框的东西在不在）、\textbf{边界合不合格}（框得准不准）、\textbf{类别正不正确}。可能“对象对、框偏、类错”三件事同时发生，必须各记各的。\item \textbf{关系}：单位是“\textbf{两个对象 + 方向 + 时间 + 上下文}”这一整套，单看一端无意义。\item \textbf{内容预测目标}（如图像修复）：检查的是“\textbf{重建目标是否仍\emph{对应}被遮蔽的原观测}”。\item \textbf{表示目标}：检查的是“\textbf{训练视图}”和“\textbf{目标\emph{生成分支}}”的条件有没有被违反（见第6章蒸馏）。\end{itemize}
还有一个常被忽略的情形：\textbf{当任务本身\emph{允许多个正确答案}}（开放集、排序、偏好类任务）时，评估单位\textbf{可以是“答案\emph{集合}、一个\emph{排序}、或一组\emph{偏好关系}”}。这时\textbf{绝不能\emph{未经约定}就强行按“单一类别”去判它错}——人家本来就该有多个解。

用同一个“山茶”标签把四类毛病\textbf{一次看清}：一张写着“山茶”的照片，可能分别\textbf{其一，错了物种（内容错）、其二，漏标了另一株花（不完整）、其三，只到种没到任务要求的亚种（粒度不够）、其四，用了已废止的分类版本（时间不适用）}。这\textbf{四种问题需要\emph{四种不同的依据}}：内容对不对找独立物种判断、完整性查原图、粒度查任务规范、时间查分类版本——\textbf{拿一种证据盖四种问题是盖不住的}。

对象类任务还有个直观例子：\textbf{两张裁剪图即使\emph{来自同一张原图}}，\textbf{只剩背景的那一张}也\textbf{不能算满足“保留花朵语义”}——来源对不等于内容对。表~\ref{tab:knowledge-quality-assessment-correction-criteria}把五类监督信号“\textbf{该找什么参考、该测什么属性}”列成对照表，它最大的用处是\textbf{防止一种偷懒的错误}：\textbf{因为“来源对应正确/格式能解析”，就误以为“语义已经保留、质量过关”}。

\begin{table*}[htbp]
  \centering
  \small
  \begingroup
  \setlength{\tabcolsep}{2.5pt}
  \begin{tabularx}{\linewidth}{@{}>{\RaggedRight\setlength{\parindent}{0pt}\arraybackslash}p{26mm}>{\RaggedRight\setlength{\parindent}{0pt}\arraybackslash}p{42mm}>{\RaggedRight\setlength{\parindent}{0pt}\arraybackslash}p{47mm}>{\RaggedRight\setlength{\parindent}{0pt}\arraybackslash}X@{}}
    \toprule
    \textbf{监督对象} & \textbf{可能问题} & \textbf{独立参考} & \textbf{可测属性} \\
    \midrule
    任务类别 & 错类、强制单解、时效失配 & 原始对象、规范与独立复标 & 参考一致性、错误类型、无法裁决比例 \\
    局部结构 & 漏标、边界偏移、类别错配 & 原图或原文、对象级规范与独立定位 & 对象召回、边界重合与局部完整性 \\
    样本关系 & 错配、漏配、方向或条件错误 & 采集标识、对象身份与采集时间 & 错配率、关系覆盖与条件违例率 \\
    内容预测目标 & 内容泄漏、对象错误、可见范围越界 & 原始观测、遮蔽位置说明与对象对应 & 目标一致性、泄漏率与对象错配率 \\
    表示目标 & 教师退化、语义被破坏、分组不稳定 & 构造日志、变换条件与冻结的用途评价 & 构造违例、稳定性、覆盖与下游效用 \\
    \bottomrule
  \end{tabularx}
  \endgroup
  \caption{不同监督信号的质量问题与判断依据。独立参考应与待核验属性相匹配；构造条件检查与下游用途评价分别提供证据，不能用一种结果替代另一种。}
  \label{tab:knowledge-quality-assessment-correction-criteria}
\end{table*}

同样一个“不符”，\textbf{严重程度\emph{完全取决于用途}}，要会分级而不是一律当错误：\textbf{错写物种}\textbf{直接改变分类监督}（高严重）；\textbf{别名/格式不统一}可能\textbf{只需做一次接口转换}就能修（低严重，不影响语义）；而\textbf{证据不足}的正确处理是\textbf{记“无法裁决”}（悬置，不是凑一个对错）。还要\textbf{把“范围外对象”\emph{另记为}“不适用”}——它\textbf{既不是缺失、也不是错误}，\textbf{混进错误率或缺失率都会污染分母}（比如一张根本不含花的风景照，谈不上“漏标花朵”）。

最后是“\textbf{规范变更}”带来的判断陷阱：如果\textbf{新规范重新划定了物种边界}，那么新旧标签不一致，\textbf{首先意味着“\emph{评价的尺子（基准）变了}”}，\textbf{不能\emph{据此}就说“旧流程\emph{退化}了”}——它们是在两把尺子下的结果，不构成时间上的优劣比较。为了\textbf{事后还能\emph{恢复}这些判断条件}，数据应\textbf{保存样本组成、采集方式、处理过程、用途和维护背景}（这也是数据文档/datasheet 的作用）\scite{gebru2021datasheets}。

\subsection{参考依据的建立}
\label{sec:knowledge-quality-assessment-correction-reference}

“查物种”这件事定了，下一个问题更关键：\textbf{\emph{凭什么}材料，才足以\emph{推翻}或\emph{保住}原标签}？注意，这里要的是\textbf{能\emph{裁决}该属性的“材料 + 程序”}，而\textbf{\emph{不是}“再多投一票”}。\term{参考依据}（reference evidence）因此被定义为：\textbf{用于下结论的\emph{证据}（原始采集日志、原始照片、独立复标结果、专家裁决、或一个可自动执行的检查）以及配套的\emph{判定程序}}。它和“又找个人标一遍”的根本区别是：\textbf{它要\emph{独立于}原标签}。

实践中最能保证独立性的程序叫\term{盲化复标}（blind re-labeling）：\textbf{复核者先\emph{只看}原始对象（照片）和当前规范}，\textbf{\emph{刻意不看}原标签写了什么、模型建议什么、置信分数多少}，先\textbf{独立写下自己的判断并\emph{保存}}；\textbf{然后}才把原标签、模型建议等冲突信息连同证据\textbf{展示给他对照}。为什么要这样排序？因为\textbf{人一旦先看到一个“貌似合理”的答案，判断会被它\emph{锚定、牵引}}——先入为主地为原答案找理由；先独立、后对照，才能\textbf{最大限度减少原答案对新判断的污染}。

但要清醒：\textbf{盲化只\emph{隔离了“已有答案}”，并\emph{不保证}复核者用的是\emph{独立的信息来源}”}。这是两种不同的独立：他可能没看原标签，却\textbf{仍然和原标注者\emph{查了同一本错误名录}}；模型和人工也可能\textbf{共用同一本\emph{过时}名录}（回到第6章“来源相关”）。因此参考裁决报告必须\textbf{把这几样都\emph{存下来}}：\textbf{证据\emph{来源}}（查了哪本名录/哪个库）、\textbf{适用\emph{时间}}（名录版本、新旧）、\textbf{人员\emph{角色}}（初标/复核/专家）、\textbf{冲突\emph{意见}}（各方当时怎么看）、\textbf{裁决\emph{理由}}（最后凭什么定），这样裁决才可追溯、可二审。

两条“不要硬来”：当\textbf{规范本身\emph{允许多解}}时（一张照片确实既像A也像B），应\textbf{保留答案集合或“成立条件”}，而不是逼一个唯一标签；当\textbf{材料不足}时，保留\textbf{无法裁决}。\textbf{强行让每条都只剩一个标签，等于把“\emph{没解决的歧义}”\emph{偷偷藏进}参考集}——表面干净，实则埋下隐患。

一个例子能说明“\textbf{证据}推翻答案”和“\textbf{票数}推翻答案”的差别。原标签、模型建议\textbf{居然\emph{都}说是山茶}（多数票一边倒），但\textbf{采集编号（照片和实物植株的绑定记录）却\emph{指向另一株植物}}。复核者按盲化流程：\textbf{先看照片和同株植物的其他视图\emph{独立}判断}，再\textbf{调采集日志}，最终发现是\textbf{照片在“写回数据库”时\emph{对象错配}}（照片被挂到了错误的植株条目下）。

请注意\textbf{推翻原结果的\emph{不是}“第三票”}——票数明明是二比零支持山茶；真正推翻它的是\textbf{一条\emph{独立的对象证据}}（采集绑定）。这再次呼应章首“检测线索并不等于错误证明、票数并不等于真相”。同时也要谦虚：\textbf{参考标签/复核结论\emph{自己也可能出错}}（复核者也会看走眼、名录也会错），所以对\textbf{高风险结论可以安排\emph{二次复核}}，并且\textbf{始终保留“能\emph{回到原始材料}重新查”的路径}，而不是只存一个最终结论。

类别标签可以请专家“再看一眼判对错”，但\textbf{第5、6章里那些\emph{表示向量/模型生成目标}，通常\textbf{\emph{根本没法}交给专家“逐维判对错”}}——没人能看着一串浮点数说“第7维应该是$0.31$”。对这类目标，参考依据要\textbf{换成“查\emph{生成过程守不守约定}}”，可查的是：
\begin{itemize}\item \textbf{视图/输入}：两张视图是否\textbf{确实来自同一对象}、有没有泄漏；\item \textbf{变换}：裁剪/变色等变换\textbf{是否保留了任务所需语义}（如裁剪后花还在不在）；\item \textbf{时间顺序}：用到序列/视频时，时间关系合不合规；\item \textbf{教师分支生成}的目标：查\textbf{教师模型版本、学生可见的输入、更新规则}（教师有没有看到学生看不到的信息、是不是冻结的）。\end{itemize}
但要\textbf{严格区分这两类证据\emph{各自能回答什么}}：上面这些\textbf{只回答了“\emph{构造过程}有没有遵守约定”}，它们\textbf{\emph{没有}回答“\emph{拿去学习}有没有\emph{受益}}”。\textbf{构造合规 并不等于 训练有用}。后一个问题（学了到底好不好）必须\textbf{固定下游任务、使用方式、评价样本集，\emph{另行}做对比实验}，不能用“生成日志齐全”冒充“效果良好”。

\subsection{审计样本的选择}
\label{sec:knowledge-quality-assessment-correction-sampling}

回到章首那一百条：\textbf{那种“挑最可疑的”选法，只适合回答“\emph{预算有限，先查哪里最可能有问题}}”（\textbf{提高找到错误的\emph{效率}}），但它\textbf{在统计上\emph{没有资格}回答“一万条\emph{整体}有多少错”}——因为它\textbf{故意}把样本往可疑处偏，入选概率既不相等也未知，无法反推全体。于是需要把两种目的\textbf{明确分成两类审计}：
\begin{itemize}\item \term{总体审计}（population audit）：\textbf{用\emph{已知的概率机制}选样}（每条被抽中的概率都明确、为正），目的是\textbf{用样本\emph{无偏地推断}指定总体的错误率}——回答“整体质量怎样”。\item \term{重点排查}（targeted inspection）：\textbf{把有限预算\emph{集中}投给冲突大、损失高、影响面广的对象}，目的是\textbf{尽快\emph{捞出}具体问题}——回答“先修哪条”。\end{itemize}
两者的\textbf{分界线（也是能否推总体的硬条件）}是：\textbf{要估计全体，抽样框（你打算代表的那批对象）里\emph{每一条相关样本}都必须有\emph{正的入样概率}}（这样每一条都“有机会被代表”，才能加权还原总体）。\textbf{重点排查\emph{可以}违反这一条}（它可以让大多数对象入样概率为0、专挑可疑的），\textbf{但代价是：它算出的错误率\emph{只能描述“入选范围}”，\emph{不能}外推到总体}。

要把“入样机会”量化，设目标总体$U=\{1,\ldots,N\}$，实际抽出的样本集为$S\subseteq U$，每条对象$i$的\term{入样概率}（inclusion probability）记为$\pi_i=\mathbb{P}(i\in S)$——即“在这套抽样方案下，对象$i$有多大概率被选中审计”。最朴素的\term{简单随机无放回抽样}（simple random sampling without replacement）（抽$n$条、抽过不放回）下，\textbf{每个对象的入样概率都相等}$\pi_i=n/N$，因此样本里的错误率可以相当直接地代表总体。

但等概率抽样在数据\textbf{不均匀}时很浪费或不公平（稀有类可能一条都没抽到），于是有两种常见的“成结构”抽样：
\begin{itemize}\item \term{分层抽样}（stratified sampling）：\textbf{先按某个\emph{已知}特征（类别、来源、时间段、风险等级）把总体切成若干\emph{互不重叠、合起来覆盖全体}的“层”}，\textbf{再在\emph{每一层内部}随机抽}。好处是\textbf{保证每一层（尤其少数类）都有样本、估计更准}。\item \term{整群抽样}（cluster sampling）：当“同组对象绑在一起”更自然时（同一文档的句子、同一用户、同一批次），\textbf{先\emph{随机抽簇}（抽文档/用户/批次），再\emph{检查被抽中簇内}的所有（或再抽样）样本}。\end{itemize}
一个\textbf{关键纪律}：为了让少数类估计更准，你\textbf{可以\emph{故意多抽}稀有层}（不等概率），但\textbf{必须\emph{如实保存每个对象实际的入样概率}}——这样汇总时\textbf{才能给“被过量抽取的组”\emph{压低权重}}，否则\textbf{多抽的那组会在总体估计里\emph{占据虚高的份额}}（稀有组看着错误很多，其实只是被你抽得多）。

图~\ref{fig:knowledge-quality-assessment-correction-audit-sampling}把\textbf{同一个教学总体}\textbf{分别}送进这两种审计，\textbf{读图时\emph{不要}数谁挑得多，要比“被挑中对象的\emph{组成}”}：
\begin{itemize}\item \textbf{左侧（分层随机审计）}：每个形状（层）都\textbf{按\emph{层内概率}被均匀代表}（保留了各层的真实比例$\pi_h=n_h/N_h$），所以它的结果\textbf{能\emph{按设计}推回总体}。\item \textbf{右侧（高风险重点排查）}：\textbf{选择条件\emph{主动}把可疑对象（叉号）的占比\emph{抬高}}——它\textbf{捞出叉号（错误）的效率更高}，但\textbf{样本组成已被人为改变}，不能再说它代表总体。\end{itemize}
同一张图、同一堆对象，仅因“怎么选”不同，一个能推总体、一个只能用于排查——这就是两种审计的直观对照。

\begin{figure*}[htbp]
  \centering
  % 图7-1：同一教学总体上的概率审计与重点排查，按169 mm全版心宽度设计。
\begin{tikzpicture}[
  x=1mm,
  y=1mm,
  font=\figurefont\footnotesize,
  text=BookInk,
  panel/.style={draw=BookRule,fill=white,line width=0.8pt},
  common item/.style={circle,draw=FigureInput,fill=FigureInput!9,line width=0.7pt,minimum size=4.2mm,inner sep=0pt},
  key item/.style={rectangle,draw=FigureModel,fill=FigureModel!9,line width=0.7pt,minimum size=4.2mm,inner sep=0pt},
  audit ring/.style={circle,draw=BookBlue,line width=1.2pt,minimum size=6.8mm,inner sep=0pt},
  focus ring/.style={circle,draw=BookAmber,line width=1.2pt,dashed,minimum size=6.8mm,inner sep=0pt},
  problem/.style={font=\figurefont\bfseries\scriptsize,text=BookAmber},
  note/.style={font=\figurefont\scriptsize,text=BookMuted,align=center}
]
  \draw[panel] (0,0) rectangle (80,58);
  \draw[panel] (86,0) rectangle (166,58);
  \draw[BookRule,line width=0.6pt,dashed] (2,20) -- (78,20);
  \draw[BookRule,line width=0.6pt,dashed] (88,20) -- (164,20);

  \node[anchor=north west,font=\figurefont\bfseries] at (2,56) {(a) 总体审计};
  \node[anchor=north west,font=\figurefont\bfseries] at (88,56) {(b) 重点排查};
  \node[note,anchor=west] at (3,17) {重点类别};
  \node[note,anchor=west] at (89,17) {重点类别};
  \node[note,anchor=west] at (3,49) {常见类别};
  \node[note,anchor=west] at (89,49) {常见类别};

  \foreach \shift in {0,86}{
    \foreach \x/\y in {14/42,25/42,36/42,47/42,58/42,69/42,14/32,25/32,36/32,47/32,58/32,69/32,14/24,25/24,36/24,47/24,58/24,69/24}{
      \node[common item] at ({\shift+\x},\y) {};
    }
    \foreach \x/\y in {17/12,31/12,45/12,59/12,73/12,24/5,38/5,52/5,66/5}{
      \node[key item] at ({\shift+\x},\y) {};
    }
    \foreach \x/\y in {25/42,58/32,36/24,31/12,66/5}{
      \node[problem] at ({\shift+\x},\y) {$\times$};
    }
  }

  \foreach \x/\y in {14/42,47/42,25/32,58/32,36/24,69/24,17/12,45/12,66/5}{
    \node[audit ring] at (\x,\y) {};
  }
  \node[note,text=BookBlue,anchor=east] at (77,17) {层内随机：$\pi_h=n_h/N_h$};

  \foreach \x/\y in {25/42,47/42,58/32,69/32,36/24,58/24,31/12,45/12,66/5}{
    \node[focus ring] at ({86+\x},\y) {};
  }
  \node[note,text=BookAmber,anchor=east] at (163,17) {优先选择：风险分数$s_i\geq\tau$};

  \node[common item] at (7,-6) {};
  \node[note,anchor=west] at (11,-6) {常见类别};
  \node[key item] at (47,-6) {};
  \node[note,anchor=west] at (51,-6) {重点类别};
  \node[audit ring] at (94,-6) {};
  \node[note,anchor=west] at (99,-6) {审计入样};
  \node[problem] at (136,-6) {$\times$};
  \node[note,anchor=west] at (140,-6) {复核后确认错误};
\end{tikzpicture}

  \caption{总体审计与重点排查的抽样差异。两个面板复用同一教学总体，形状表示分层，叉号表示实际存在但抽样前未知的错误。左侧在每层随机入样并保存$\pi_h=n_h/N_h$，可按设计估计总体；右侧按冲突分数超过阈值优先选择，能提高排查命中，却改变了样本组成。图中点数仅用于说明选择机制，不按真实总体规模绘制。}
  \label{fig:knowledge-quality-assessment-correction-audit-sampling}
\end{figure*}

下面用一个\textbf{可亲手复算}的教学例子，看清“\textbf{为什么\emph{多抽}了重点层、汇总时反而要给它\emph{更小}权重}”。设独立验收总体含常见类别$N_1=800$条、茶梅易混淆区域$N_2=200$条，不含开发对象$x_2$及其近重复样本标注。用于比较的K17和K18配置均在本次抽样前冻结，这里先看K17。两层独立进行简单随机无放回抽样，分别审计$n_1=80$条和$n_2=40$条，全部可裁决，确认K17的错误数为$4$和$8$。两层错误率为$\widehat p_1=0.05$和$\widehat p_2=0.20$，按总体占比$W_1=0.8,W_2=0.2$汇总得到$\widehat p_{\mathrm{st}}=0.08$。现在\textbf{对比三种算法，它们在同一批数据上给出三个不同的数}，务必分清各自在回答什么：
\begin{itemize}\item \textbf{直接把样本错误数除以样本数}$12/120=0.10$：\textbf{错}。它\textbf{默认每个抽中的样本权重相同}，但易混淆区被\emph{过量抽取}了（实际只占总体两成，却占了样本的$1/3$），于是\textbf{把问题严重的重点区\emph{看得过重}}，高估了总体。\item \textbf{把两层错误率\emph{简单平均}}$(0.05+0.20)/2=0.125$：\textbf{也不是总体错误率}。它回答的是“\textbf{一个‘两层\emph{等权}’的平均质量}”——即“随便选一个层、再在层里随便看”，\textbf{让只占两成的重点层和占八成的常见层\emph{平起平坐}}，\emph{高估得更厉害}。\item \textbf{按总体占比加权}$\widehat p_{\mathrm{st}}=0.8\times0.05+0.2\times0.20=0.08$：\textbf{这才是“\emph{总体}里有多少比例错”的无偏估计}。\end{itemize}
直觉口诀：\textbf{权重只由“\emph{这一层在总体里占多少}”决定，\emph{不}由“你在这一层抽了多少}”决定。这些数字都是教学构造（且该独立验收总体\textbf{不含}$x_2$及其近重复样本），后面7.4还会用\textbf{同一批独立样本}比较K17、K18两个版本。

抽样里还有一个\textbf{容易被统计术语绕晕、但必须先分清的前提}：你\textbf{要估计的，到底是“\emph{一份\emph{已经存在}的固定名单}的质量”，还是“\emph{未来\emph{还没出现}的新对象}的质量”}？这两者的“随机性来自哪”完全不同，结论的适用范围也不同。\textbf{第一种：固定名单}。这份名单里\textbf{每条标注的对错在抽样前就已经\emph{客观确定}}（要么对要么错，只是你不知道），所以\textbf{随机性\emph{唯一}来自“这一轮\emph{抽中了哪些编号}}”——这是一种“对固定有限总体的抽样推断”，结论只对这份名单负责。这里要注意\textbf{抽样单位（簇）}：
\begin{itemize}\item 如果你确实是\textbf{按“全体\emph{样本}编号”简单随机抽}，那么哪怕同一文档内的句子共享来源和上下文，只要它们各自被独立编号、独立入样，仍可对这份有限总体做标准推断。\item 但如果流程是\textbf{先抽“\emph{文档}}”、再审被抽中文档里的句子，那么\textbf{文档才是第一阶段的\emph{抽样单位}}，句子只是簇内单元；这时\textbf{方差计算和重抽样都\emph{必须保留这个“先文档、后句子”的簇结构}}，把句子当独立样本会\textbf{严重低估}不确定性（同文档句子高度相关，没那么多独立信息）。\end{itemize}
两个适用边界要记住：其一，\textbf{这种“固定名单”推断\emph{不会自动}告诉你“\emph{未来新文档/新照片}”的质量}——名单外是另一个总体；其二，\textbf{近重复照片、同一用户标的大量数据、同一批次生成物，\emph{不能}因为“条数多”就当成“同样多份\emph{独立}证据}”（它们高度相关，呼应第6章$\rho$和来源依赖）。

最后是“\textbf{到底审多少条}”的样本量问题——它\textbf{不是拍脑袋，也不是越多越好}，而要同时联系四件事：\textbf{其一，允许的估计误差}（你能接受多宽的区间，要求越精样本越多）、\textbf{其二，需要\emph{单独}报告的分组}（每个组都要有足够样本，少数类往往要加抽）、\textbf{其三，预期错误率}（错误越稀少，要抓到足够错误所需样本越多）、\textbf{其四，复核成本}（每条复核都花钱）。这里有个反直觉点：\textbf{只按总体固定比例抽取，\emph{不保证}少数类别能攒够证据}——一个占$0.5\%$的稀有类，哪怕抽$1000$条也可能一条没有。

还有一个\textbf{“无响应/没查成”陷阱}：\textbf{抽中了但\emph{没能完成复核}的对象，也必须\emph{保留在记录里}}，不能默默删掉。因为\textbf{往往\emph{难例}更容易“无响应”}（太糊、太难、专家没空判），如果你\textbf{把没查成的删掉}，留下的就\textbf{系统性偏向“\emph{容易判}”的样本}，错误率被人为压低。最后澄清两类审计能否合并：\textbf{概率审计和重点排查可以\emph{同时并行}做，但要\emph{分别标明用途}}；只有当你\textbf{给“合并样本”也\emph{建立了覆盖目标总体的入样机制和配套估计量}}时，才能把它们\emph{合并进同一个推断}——\textbf{绝不能图省事直接把“重点排查的错误数”加到“随机审计的错误数”上、分母也直接相加}（两者抽样机制不同，相加就是把选择性偏差掺进总体估计）。

\section{知识质量的测量与估计}
\label{sec:knowledge-quality-assessment-correction-measurement}

审计样本拿到手、复核也做完，接下来是\textbf{“把\emph{证据}变成\emph{数字}”}的两步：\textbf{先把每条证据落成一个\emph{单条裁决}}（这条到底是对、错、还是无法判），再把单条裁决\textbf{汇总到“分组”和“总体”}。这里有两个\textbf{性质完全不同}的结论，必须在脑子里分开：
\begin{itemize}\item \textbf{\emph{直接核验}}：在当前参考依据下，\textbf{确认\emph{某一条具体标注}有错（或对）}。这是\textbf{针对个体的“事实认定”}，审过才有，确定。\item \textbf{\emph{有限样本估计}}：\textbf{只查了\emph{一部分}样本，却要\emph{推断没查的那一大批}的错误比例}。这是\textbf{从样本到总体的“统计推断”}，\textbf{依赖你\emph{怎么选的样}}，天然带不确定性、可能偏。\end{itemize}
正因为后一种结论\textbf{依赖选样方式}，所以\textbf{报告任何比例时，都不能只甩一个百分数}，必须\textbf{同时给出}：\textbf{分母是谁}（哪批对象、哪个范围）、\textbf{标注版本}、\textbf{有效裁决数}（多少条真的判出了对错）、以及\textbf{无法裁决的比例}。离开这四样，一个“错误率$8\%$”无法被解释、也无法复算。

\subsection{样本级质量的测量}
\label{sec:knowledge-quality-assessment-correction-item-measurement}

一次单条审计要\textbf{备齐四类材料}：\textbf{待审标注本身、原始对象、任务规范、参考证据}；产出的审计结论则要\textbf{分字段记录}：\textbf{处理阶段}（待复核/已完成）、\textbf{质量结论}、\textbf{问题类型}、\textbf{严重程度}、\textbf{证据}。这种结构化记录是为了后面能\emph{按问题类型汇总、按版本追溯}，而不是一句“这条有问题”。

怎么判，要\textbf{先固定“\emph{评价单位}”和“\emph{计分规则}”}，否则对错数会随心情变：
\begin{itemize}\item \textbf{单一物种答案}：通常直接判\textbf{正确/错误}。\item \textbf{含多个待测属性的结果}：需要“\textbf{部分正确}”这个档——\textbf{结构化标注要把“对象、边界、类别、关系”\emph{拆开}各判各的}。典型例子：\textbf{一个花卉框\emph{位置（边界）对}、但\emph{物种}错}，这两个结论\textbf{必须\emph{分别保存}}，不能合成一个“对”或“错”。\end{itemize}
算二值错误率\textbf{之前}，要先白纸黑字写死计分规则，比如规定“\textbf{对象存在、边界、类别\emph{全部}合格才记这一条正确}（任一出错即错）”。\textbf{切忌事后把“部分正确”\emph{临时}折算成“半个错误}”——折算规则必须\emph{事先}定，不同任务不可通用。\textbf{缺失、不适用、无法裁决}这三类\textbf{要按\emph{预定的分母口径}单列}，\textbf{既不强填成“正确”、也不强填成“错误}”（填正确会掩盖漏洞，填错误会冤枉“证据不足”的样本）。

如果一条标注给的不是硬类别、而是\textbf{一整个类别概率分布$\bm{q}_i$}（比如融合/教师软标签），那么\textbf{只挑“最高概率那个类”判对错，会\emph{丢掉}“概率\emph{怎么分配}}”的信息——两个都把最高分给山茶的标注，可能一个$(0.9,0.1)$很笃定、另一个$(0.51,0.49)$很犹豫，硬判都“对”，却完全没有区别。所以要评价概率本身的质量。设$Y=\{1,\ldots,K\}$，$\bm{q}_i=(q_{i,c})_{c\in Y}$且$q_{i,c}\geq0,\ \sum_cq_{i,c}=1$，独立参考给出的真实类别为$y_i$。设$Y=\{1,\ldots,K\}$，样本标注$i$给出类别概率$\bm{q}_i=(q_{i,c})_{c\in Y}$，其中$q_{i,c}\geq0$且$\sum_cq_{i,c}=1$，独立参考类别为$y_i$。对数损失（log loss）$-\log q_{i,y_i}$惩罚给实际类别过低的概率，概率为零时损失为无穷；\term{Brier评分}（Brier score）$\sum_{c\in Y}(q_{i,c}-\mathbf{1}\{y_i=c\})^2$则计算概率向量与实际类别指示向量的平方距离，其中指示函数在$y_i=c$时为1，其余为0。读懂这两个分数的\textbf{不同脾气}：
\begin{itemize}\item \textbf{对数损失}：\textbf{对“\emph{把真实类概率压到极低}”惩罚极重}——若$q_{i,y_i}=0$，损失是$+\infty$（你断言真类不可能，错得无可挽回）；它\textbf{无界}，越自信地错、罚得越狠，因此鼓励“\emph{没把握别押死}”。\item \textbf{Brier 评分}：本质是\textbf{概率向量与“真值独热向量”的\emph{平方欧氏距离}}，它\textbf{有界}（二分类时在$0$到$2$之间），更温和、对异常值不那么敏感。\end{itemize}
两者都是\textbf{越小越好}。为什么它们能“逼”模型说真话？因为它们属于\term{严格适当评分规则}（strictly proper scoring rule）：有一个很漂亮的性质——\textbf{当你\emph{确实知道}真实条件分布时，\emph{唯一}能让期望分数最小的做法，就是\emph{如实}报出这个分布}；任何系统性地夸大、押注都会在期望上吃亏\scite{gneiting2007scoring}。所以用它们评软标签，不是只看“中没中”，而是看“概率分配诚不诚实”。

但单条分数再低，也\textbf{回答不了另一个更宏观的问题——\term{概率校准}（probability calibration）}：\textbf{你\emph{一批}预测里报出的概率，\emph{整体上}和这些事件\emph{实际发生的频率}对得上吗}？“这次$0.8$对了”没有意义，校准永远是\textbf{对“\emph{一组}报了相同概率的预测”做统计}。\textbf{校准检查的第一件事是“\emph{先指定你说的是哪个事件}}”，因为“$0.8$”可以指两种完全不同的事件，混淆它们会得到错误结论：
\begin{itemize}\item \textbf{查“最高分决策”对不对（\emph{置信度/精度}校准）}：把所有“\textbf{预测类的最大概率约为$0.8$}”的样本聚成一组，核对其中“\textbf{预测\emph{正确}}”的比例是否约$0.8$。这里事件是“\emph{我的最高票答案中没中}”。\item \textbf{查“某个\emph{固定类别}}”：把所有“\textbf{分给\emph{茶梅}}的概率约$0.8$”的样本聚成一组，核对其中“\textbf{真实类别\emph{就是茶梅}}”的比例是否约$0.8$。\end{itemize}
\textbf{这两种分组、两个事件\emph{不一样}}：\textbf{前者成立（置信度校准好），\emph{不代表}每个类别的概率都已校准}（可能某类系统性高估、另一类低估，平均后最高分仍看似准）。要全面校准，得\emph{逐类}做第二种检查。概率是连续值，实践中把它\textbf{按\emph{事先约定}的概率区间}分桶（如$[0,0.1),[0.1,0.2),\ldots$），每桶报告“桶内平均预测概率 vs 实际频率”，并\textbf{同时报告每桶\emph{样本数}和\emph{不确定性}}（只含3条的桶，频率再准也不可信）。一个\textbf{不可妥协的数据纪律}：这些校准检查\textbf{必须用\emph{没参与}概率拟合/温度缩放/调参的\emph{独立评价}样本}——在调参用的数据上检查校准，结论天然偏乐观\scite{guo2017calibration}。

这里要给一组\textbf{“\emph{不是}质量裁决”的常见信号}正名，它们最容易被误用：\textbf{模型置信度低、损失高、来源之间分歧大}——这些\textbf{都\emph{不能直接}判定这条标签错}，它们\textbf{只能\emph{触发一次复核}}（把它排进待查队列），但\textbf{无法\emph{独立}回答“原标签/约束到底对不对}”（模型可能自己错、高损失可能只是样本罕见，见第6章和7.3）。两个同理的区分：\textbf{文件字段齐全 到 只能记“格式合规”}（不代表内容对）；\textbf{多个来源一致 到 只能记“意见\emph{一致}}”（不代表一致地对，它们可能共享同一错误）。

总原则一句话：\textbf{每一个状态都必须\emph{对应明确的证据}}；而\textbf{“无法裁决”的样本，\emph{不能}为了让数字好看就\emph{悄悄从分母里删掉}}。

对那些\textbf{“构造类”目标（样本关系、恢复/表示目标）}，样本级测量要\textbf{诚实——只\emph{报告可观察、可核对的条件}，\emph{不编造}无法核对的对错}。各类能直接查的条件是：
\begin{itemize}\item \textbf{正对应关系（配对/正样本）}：核对\textbf{对象标识是否一致}（两张视图是否真属同一对象）、\textbf{变换有没有越界}（裁剪有没有把配对破坏）。\item \textbf{内容恢复目标}（如图像补全）：核对\textbf{重建目标是否确实来自“被遮蔽的那部分”}（不能用别处可见内容冒充）。\item \textbf{表示目标}（蒸馏提示）：核对\textbf{教师版本是否存在、表示有没有\emph{退化成常数}}（所有样本提示都一样＝没信息，见第6章FitNets）。\end{itemize}
\textbf{如果根本没有“逐向量标准答案”可比，就\emph{不要}硬贴一个“表示正确/错误”的标签}——那是无法裁决的构造问题；它的\textbf{“下游到底有没有用”应\emph{留给独立的用途评价}}（在下游任务上实测），两层不互相冒充。

\subsection{数据集与分组质量的估计}
\label{sec:knowledge-quality-assessment-correction-dataset-estimation}

要把一条条状态\textbf{汇总成数据集级指标}，第一步是\textbf{用四个$0/1$开关，把“一条标注所处的\emph{状态}”\emph{拆开}记录}，避免“没采到、判不了、缺失、错误”被混成一锅。对样本$i$定义：
\begin{itemize}\item $a_i=1$：\textbf{取得}了候选标注（这条进了我们实际拿到的集合）；\item $j_i=1$：这条监督信号\textbf{\emph{可裁决}}（证据够，能判对错）；\item $e_i=1$：裁决\textbf{确认它与参考\emph{不一致}}（即确认有错）；\item $m_i=1$：\textbf{任务\emph{要求}有这条监督信号、它却\emph{缺失}}。\end{itemize}

这四个开关怎么组合，\textbf{决定了不同指标的\emph{分子分母}，最容易出错的就是“分母放谁}”：
\begin{itemize}\item \textbf{错误标志$e_i$只在“\emph{已取得}$a_i=1$\textbf{且}\emph{可裁决}$j_i=1$”的标注上计分}。\textbf{没取得、或无法裁决的，\emph{不能}按“错误=$0$”处理}——它们不是“判对了”，只是“这次没结论”，塞进去会人为压低错误率。\item \textbf{条件错误率}：分母是“\textbf{取得\emph{且}可裁决}”的标注（在“能下结论的”里有多少错）。\item \textbf{缺失率}：分母是“\textbf{任务要求\emph{本应}有该信号的\emph{总体}}”（该有的里面缺了多少），跟错误率分母不同。\item \textbf{获取覆盖率}：还得\textbf{先选口径}——“\emph{有任意一类票}”算覆盖、“\emph{可裁决}”算覆盖、还是“\emph{达到接收阈值}”才算覆盖？\end{itemize}
一句话：\textbf{分母一改，你回答的就是另一个问题}。所以\textbf{绝不能只甩一个光秃秃的“$8\%$”而不写口径}——同样$8$条错，除以“可裁决的$100$条”和除以“应有的$125$条”是两个完全不同的结论。

表~\ref{tab:knowledge-quality-assessment-correction-metrics}把\textbf{五个常用指标的“分子 / 分母 / 汇总时还要别报什么”}一次列清（覆盖条件下错误率、总体缺失率、获取覆盖率、关系错配率、构造违例率）。读这张表要抓一条：\textbf{这些指标\emph{描述的是不同对象}}——类别错误说的是答案、关系错配说的是两端对象和方向、构造违例说的是生成条件，它们\textbf{应当\emph{并列}报告，让读者看到各自的短板}，而\textbf{\emph{不应}加权相加成一个“含义不明的\emph{总分}}”（把“类别错$3\%$”和“关系错$5\%$”加成“$8\%$”没有物理意义，还会互相掩盖）。

\begin{table*}[htbp]
  \centering
  \small
  \begingroup
  \setlength{\tabcolsep}{2.5pt}
  \begin{tabularx}{\linewidth}{@{}>{\RaggedRight\setlength{\parindent}{0pt}\arraybackslash}p{29mm}>{\RaggedRight\setlength{\parindent}{0pt}\arraybackslash}p{43mm}>{\RaggedRight\setlength{\parindent}{0pt}\arraybackslash}p{44mm}>{\RaggedRight\setlength{\parindent}{0pt}\arraybackslash}X@{}}
    \toprule
    \textbf{指标} & \textbf{分子} & \textbf{分母} & \textbf{汇总时另报} \\
    \midrule
    覆盖条件下错误率 & 已获得且经参考确认错误的监督信号 & 已获得且可裁决的监督信号 & 无法裁决比例、参考标准与版本 \\
    总体缺失率 & 任务要求但未取得监督信号的对象 & 总体中应获得监督信号的对象 & 完全遗漏区域与不适用对象 \\
    获取覆盖率 & 达到所选接收条件的对象 & 目标总体中的适用对象 & 任意类别票、可裁决或可用的接收口径 \\
    关系错配率 & 对象、方向、时间或条件错误的关系 & 已声明且可核验的关系 & 关系单位与共享对象的依赖 \\
    构造违例率 & 违反已声明生成条件的目标或关系 & 具备检查条件的已生成对象 & 变换、教师、批次与构造版本 \\
    \bottomrule
  \end{tabularx}
  \endgroup
  \caption{质量指标的分子、分母与汇总范围。未知、不适用和无法裁决状态应单列；表中指标可以并列报告，但类别错误、关系错配和构造违例描述不同对象，不能相加成含义不明的总分。}
  \label{tab:knowledge-quality-assessment-correction-metrics}
\end{table*}

为了先把“\textbf{怎么从样本错误率\emph{无偏}地估总体}”讲清楚，本小节\textbf{先做一个简化限定}：只考虑“\textbf{每条\emph{都有}候选标注、且\emph{都能}按同一套二值规则裁决}”的适用样本（即暂时令$a_i=j_i=1$、没有无法裁决的情形）。\textbf{“无法裁决怎么办”这个更麻烦的情形，\emph{留到下一小节}单独处理}——先在干净情形下建立估计量，再处理复杂口径。\textbf{简单随机抽样}时，直觉也成立：\textbf{用样本里错误指示的均值，就能估总体错误率}（每条权重相同）。但\textbf{分层（且各层\emph{不等比例}抽取）时，不能直接平均}——因为样本的层比例被抽样改变了，必须\textbf{把各层“\emph{恢复}回它在总体里的真实份额}”再加权。设互不重叠的$H$层\textbf{覆盖整个总体}，第$h$层总体有$N_h$条、占比$W_h=N_h/N$，在该层随机审计$n_h>0$条、样本索引集$S_h$，层内错误率$\widehat p_h$为层内$e_i$的均值。设互不重叠的$H$层覆盖总体，第$h$层有$N_h$条，随机审计$n_h>0$条，样本索引为$S_h$，总体权重$W_h=N_h/N$。层内先求均值，再按总体权重相加，得到
\begin{equation}
  \widehat p_{\mathrm{st}}
  =
  \sum_{h=1}^{H}W_h\widehat p_h,
  \qquad
  \widehat p_h
  =
  \frac{1}{n_h}\sum_{i\in S_h}e_i.
  \label{eq:knowledge-quality-assessment-correction-stratified-rate}
\end{equation}
再强调一次权重的来源：\textbf{$W_h$由“第$h$层在\emph{总体}里占多少条”决定，\emph{不}由“你在该层\emph{抽了}多少条”决定}——这就是前面教学例得到$8\%$（而非样本直算的$10\%$）的原因。

算出分层错误率后，\textbf{怎么\emph{用}它做决策也有讲究}：如果验收要求是“\textbf{每个类别（层）都要达标}”，就必须\textbf{把\emph{各层}$\widehat p_h$\emph{分别}与该组标准比}，\textbf{一层不达标就是不达标}；\textbf{不能只看“\emph{按类等权}平均”的总错误率}——平均可能\textbf{被表现好的大类拉下来，\emph{掩盖}某一小类严重不合格}（这和6.6帕累托、6.8分别报告各任务是同一种纪律）。

分层估计量解决了“层间比例失衡”，但还有一种更一般的失衡：\textbf{如果\emph{每条对象}的入样概率$\pi_i$\emph{彼此不同}}（重点排查、不等概率抽样时必然如此），那\textbf{直接把抽中样本的错误率平均，会\emph{偏向}那些“\emph{更容易被选中}}”的对象——常被抽中的类，无论对错都被反复代表，估计就歪了。针对上面那个“都可二值裁决”的总体，在\textbf{总体规模$N$已知、且\emph{每条}对象入样概率$\pi_i>0$}的前提下，\term{Horvitz--Thompson 估计}（Horvitz--Thompson estimator，霍维茨–汤普森）用一个极简洁的想法纠偏：\textbf{每一条“被抽中”的错误，都\emph{除以它自己的入样概率}$1/\pi_i$}——\textbf{越难被抽中的对象（$\pi_i$小），一旦被抽中、就“\emph{以一当多}”，\emph{补偿}它代表的那一大群没被抽中的同类}；最后再除以总体规模$N$归一化：
\begin{equation}
  \widehat p_{\mathrm{HT}}
  =
  \frac{1}{N}
  \sum_{i\in S}\frac{e_i}{\pi_i}.
  \label{eq:knowledge-quality-assessment-correction-ht-rate}
\end{equation}
为什么除以$\pi_i$恰好纠偏？关键恒等式是
$$\mathbb{E}\!\left[\mathbf{1}\{i\in S\}\frac{e_i}{\pi_i}\right]=\frac{e_i}{\pi_i}\mathbb{P}(i\in S)=\frac{e_i}{\pi_i}\pi_i=e_i.$$
也就是说：\textbf{固定总体里每条对错（$e_i$已定），只对“抽不抽中”取期望，被抽中时乘的$1/\pi_i$正好和入样概率$\pi_i$抵消，\emph{还原出这条错误在总体里本应有的贡献}$e_i$}。这就是“无偏”的来源。

但它有\textbf{三条必须知道的边界}：其一，\textbf{依赖\emph{准确已知}的设计概率$\pi_i$}——$\pi_i$是多少必须来自抽样设计本身，估错了纠偏也错；其二，\textbf{抽样框\emph{漏掉}的区域$\pi_i=0$}（那些对象从一开始就不可能被抽中），\textbf{再怎么给其他样本加大权重也\emph{补不回来}}——权重只能补偿“有机会被抽中但没中”的，补不了“根本不在框里”的；其三，\textbf{某\emph{一次}抽样算出的估计\emph{可能落在$[0,1]$之外}}（比如超过100\%），\textbf{不必惊慌也别截断}——\textbf{“无偏”是“\emph{反复按同一设计}抽样、估计值平均等于真值”的性质}，\textbf{不是“每一次估计都落在合理范围”的保证}。另外，如果\textbf{连总体规模$N$（分母）也未知}，可改用\term{比率估计}（ratio estimator，加权错误数/加权样本数），但那\textbf{是（近似）无偏的比值，\emph{不能}沿用这里“精确无偏”的解释}。该构造来自 Horvitz 与 Thompson\scite{knowledge-quality-assessment-correction-horvitz1952}。

最后一个\textbf{名字像质量指标、其实很弱}的量：\term{来源一致率}（source agreement rate）只描述“\textbf{两个来源给出\emph{相同}结果的\emph{频率}}”。它\textbf{\emph{不是}正确率}，两个反例要记住：\textbf{没有独立参考时，两个来源可能\emph{一致地一起错}}（共享同一错误名录/模板，呼应第6章$\rho=1$）；\textbf{任务\emph{允许多解}时，它们给出\emph{不同但都合理}的答案}也完全正常（一致反而怪）。

所以做来源比较时，要先\textbf{对齐“对象、任务定义、适用范围”}（否则比的不是一回事），并把结果\textbf{分成三类\emph{分开}记：一致、冲突、无法比较}——一致不等于对、冲突不等于某方错。最后回到报告纪律：\textbf{任何质量估计都要附上“审计样本数、有效裁决数、无法裁决数、标注版本”}，让读者能判断“这个百分数\emph{实际建立在多少证据上}”——一个基于$3$条有效裁决的$0\%$和基于$300$条的$0\%$，分量天差地别。

\subsection{质量估计不确定性的分析}
\label{sec:knowledge-quality-assessment-correction-uncertainty}

前例教学算出的$8\%$，\textbf{只是“\emph{这一次}随机抽样”得到的\emph{一个}数值}。请意识到：\textbf{如果你\emph{重新}按同一方案随机抽一批}（抽到的编号不同），错误数会变，算出的百分比也会变——可能这次是$8\%$，下次是$6\%$、$11\%$。所以单点估计旁边必须带一个“\textbf{它\emph{会波动多大}}”的量，这就是\term{设计方差}（design variance）：\textbf{固定总体、固定标注版本，\emph{重复执行同一套抽样方案}时，估计值在不同次抽样之间的波动}。注意它\textbf{刻画的是“抽样随机性”，不是标注本身对错}——总体没变，只是看的样本不同。

先推导最基础的情形——\textbf{简单随机、\emph{无放回}抽样}（抽过的不再放回去）。设$2\leq n\leq N$，$e_i$是前面定义的二值错误指示，$\widehat p$是样本错误率均值，并记\term{抽样比例}（sampling fraction）$f=n/N$（你查了总体的多大部分），样本内部错误指示的（无偏）样本方差为
$$s_e^2=\frac{1}{n-1}\sum_{i\in S}(e_i-\widehat p)^2.$$用这个样本方差去估计总体内部错误指示的离散程度，再\textbf{按“样本量$n$越大越准”和“抽样比例$f$越大越准”两件事\emph{调整}}，就得到估计量$\widehat p$的方差：
\begin{equation}
  \widehat{\operatorname{Var}}(\widehat p)
  =
  (1-f)\frac{s_e^2}{n}.
  \label{eq:knowledge-quality-assessment-correction-srs-variance}
\end{equation}
这个公式要读出两个要点。其一，\textbf{分母上的$1/n$}：\textbf{查得越多（$n$越大），估计越稳、方差越小}——这是“多样本平均抵消波动”。其二，前面的因子$(1-f)$叫\term{有限总体修正}（finite population correction）：\textbf{因为是“\emph{无放回}”抽样，\textbf{同一条不会被重复抽中}，当你\textbf{已经查了总体的很大一部分}时，剩余不确定性更小——\textbf{相比“有放回（可能重复抽到同一条）”的抽样，它\emph{额外}压低了方差}}。

两个边界：\textbf{当你\emph{审遍全部}标注（普查）时$f=1$，这个抽样方差\emph{正好为零}}——因为没有“没看到”的抽样随机性了；\textbf{但“抽样波动为零”\emph{不等于}“没有错误”}——参考依据本身可能错、无法裁决仍存在，只是不再有“用样本推总体”的那层不确定。最后，方差开平方得到\term{标准误}（standard error）$\widehat{\operatorname{se}}$，它\textbf{把波动换回“错误率”同一个尺度}（百分点），报告区间时直接用它。有限总体、分层与整群抽样的设计公式可参见Cochran的整理\scite{knowledge-quality-assessment-correction-cochran1977}

分层抽样下方差怎么合成？关键直觉：\textbf{每一层是\emph{独立}选样的（层间抽样互不影响），而各层在总体里占的权重$W_h$是\emph{固定}常数}，所以\textbf{“加权和”的方差＝“方差的加权和”，而\textbf{常数权重\emph{要平方}}}——即各层方差按$W_h^2$相加，\textbf{\emph{不是}按$W_h$相加}（独立随机变量和的方差里权重是平方进去的）。设层内也是简单随机无放回抽样（$n_h\geq2$、$f_h=n_h/N_h$），层内错误指示样本方差为$s_h^2$，则：设层内简单随机无放回抽样，$n_h\geq2$，$f_h=n_h/N_h$，错误指示样本方差为$s_h^2$，则
\begin{equation}
  \widehat{\operatorname{Var}}(\widehat p_{\mathrm{st}})
  =
  \sum_{h=1}^{H}
  W_h^2(1-f_h)\frac{s_h^2}{n_h}.
  \label{eq:knowledge-quality-assessment-correction-stratified-variance}
\end{equation}
要让上面的分层均值公式和分层方差公式\textbf{自洽}，必须\textbf{两式用\emph{同一套}分层定义、同一个抽样单位}（不能均值按“类别”分层、方差按“来源”分层；也不能一个按文档、一个按句子）。前面警告过“抽样单位可能是\emph{簇}”，这里看它对方差的影响：
\begin{itemize}\item \textbf{如果是“先抽\emph{文档}、再审被抽中文档里的标注”}：那么\textbf{第一阶段（抽没抽中某文档）带来的波动，必须\emph{按文档为单位}估计}，不能把同文档的句子当独立。\item \textbf{如果簇内还\emph{再抽样}}（抽中文档后不查全部句子、再抽句子）：这是\textbf{两阶段整群抽样}，\textbf{两个阶段的波动\emph{都}要计}，方差通常比把句子当独立更大。\end{itemize}
当抽样设计很复杂、解析公式难推时，有一条\textbf{更通用的路：\term{复制权重}（bootstrap replicate weights / 复制权）}——\textbf{\emph{严格照原设计}保留“层”和“第一阶段抽样单位（簇）”，\emph{整体地}复制整个簇/层，生成很多组“重加权样本”，每组算一个估计值，再用这些复制估计值之间的差异\emph{经验地}估方差}。这里\textbf{最忌讳}的是\textbf{把所有标注\emph{逐条独立}重抽的普通自助法}——它\textbf{直接丢掉“同文档/同簇相关”这层结构}，会\textbf{严重\emph{低估}真实方差}（假装独立的句子比实际多），不能替代“尊重设计”的重抽样\scite{knowledge-quality-assessment-correction-rao1988}。

把方差公式\textbf{套回前面$N_1=800,N_2=200$、抽$n_1=80,n_2=40$的教学例}，亲手算一遍。二值指示有个方便性质：层内（无偏）样本方差可写成$s_h^2=\dfrac{n_h\widehat p_h(1-\widehat p_h)}{n_h-1}$（只由该层错误率和样本量决定）。代入两层：
\[
  \begin{aligned}
    \widehat{\operatorname{Var}}(\widehat p_{\mathrm{st}})
    &=
    0.8^2(1-0.1)\frac{0.05(1-0.05)}{79}\\
    &\quad+0.2^2(1-0.2)\frac{0.20(1-0.20)}{39},\\
    \widehat{\operatorname{se}}(\widehat p_{\mathrm{st}})
    &\approx 0.02185.
  \end{aligned}
\]
若再\textbf{用正态近似}给一个$95\%$区间：$8\%\pm1.96\times2.185\%$约为$[3.72\%,12.28\%]$。怎么\textbf{正确解读}这个区间？它只表示“\textbf{在正态近似成立时，点估计周围的\emph{抽样波动}范围}”。本例\textbf{每层错误条数很少（4条、8条），正态近似本身就勉强}，所以这个区间\textbf{只能当“近似演算”，\emph{不能}当成“精确的$95\%$覆盖保证}”。更稳的做法：\textbf{靠近$0$或$1$边界、或要做\emph{严格验收}时，应改用与抽样设计相符的\emph{离散}精确区间（见下页超几何法），或\emph{事先}规划一个样本量更充分的审计}，而不是用一个勉强的正态区间下“达标”结论。

分层、简单随机里，各对象入样相对独立；但\textbf{不等概率、尤其整群抽样时，波动还取决于“\emph{两条}是否\emph{容易一起}被抽中”}——同一文档的句子\textbf{要么一起进、要么一起不进}，这种“\textbf{成对的同进同出}”会\textbf{放大波动}。此时\textbf{单条入样概率$\pi_i$\emph{不足以}描述这种关系}（它只说每条自己多容易进，没说两条一起进的概率），还要二阶量。引入\term{二阶入样概率}（second-order inclusion probability）$\pi_{i,j}=\mathbb{P}(i,j\in S)$（$i,j$两条\textbf{同时}被抽中的概率），约定$\pi_{i,i}=\pi_i$。那么量$(\pi_{i,j}-\pi_i\pi_j)$正好是“两条入样指示的\textbf{协方差}”：它\textbf{为正＝两条倾向\emph{同进同出}}（同簇，正相关、增大方差），为负＝此进彼出。为估计式~\eqref{eq:knowledge-quality-assessment-correction-ht-rate}的方差，把这些协方差贡献加总，并用联合入样概率补偿只观察到入样样本标注对的事实：
\begin{equation}
  \widehat{\operatorname{Var}}(\widehat p_{\mathrm{HT}})
  =
  \frac{1}{N^2}
  \sum_{i\in S}\sum_{j\in S}
  \frac{\pi_{i,j}-\pi_i\pi_j}{\pi_{i,j}}
  \frac{e_i}{\pi_i}\frac{e_j}{\pi_j},
  \label{eq:knowledge-quality-assessment-correction-ht-variance}
\end{equation}
这个 HT 方差公式\textbf{有个硬前提：所需的二阶概率$\pi_{i,j}$}\emph{都为正、且真的保存了、可取得}。实践里若\textbf{没存$\pi_{i,j}$}（只记了单条$\pi_i$），就\textbf{不能硬套这个公式}，只能\textbf{改用能\emph{反映真实设计}的方差方法}：比如\textbf{一阶（线性化）近似}（用一阶泰勒展开近似估计量的抽样变化），或前面讲的\textbf{复制权重}法。\textbf{绝不能图省事\emph{退回}“把样本当独立”的公式}——那等于假设$\pi_{i,j}=\pi_i\pi_j$（无簇相关），会系统性低估方差。

选哪种方差方法，还\textbf{必须有配套的“抽样设计字段”支撑}（分层定义、抽样单位、入样概率等都要在日志里）；\textbf{设计日志缺失时，\emph{不能}只凭“我用的是整群抽样”这个\emph{方法名}就\emph{补}出一个方差数}——没有$\pi_{i,j}$、没有簇结构，方差在数学上无法恢复。这套“尊重抽样设计”的估计框架，系统的整理可参见 Särndal 等人的抽样调查教材\scite{knowledge-quality-assessment-correction-sarndal1992}。

当\textbf{样本量足够、抽样分布可用正态近似、且估计\emph{不靠近}$0/1$边界}时，可以用一个简洁的对称区间报告：
$$\widehat p\ \pm\ z_{1-\alpha/2}\,\widehat{\operatorname{se}}(\widehat p),$$
其中$0<\alpha<1$是你容许的误差概率（如$\alpha=0.05$对应$95\%$），$z_{1-\alpha/2}$是标准正态分位数（$95\%$时约$1.96$）。它的含义是“在正态模型成立时，估计值大致落在点估计上下这么多个标准误内”。但对称正态区间有两个\textbf{失效区}：其一，\textbf{错误很少}时，对称区间可能\textbf{越出$[0,1]$}（算出负的错误率），不合常理；其二，更隐蔽的是\textbf{“样本里\emph{零个错误}}”——此时经验方差甚至也算成$0$，正态区间退化成一个点$[0,0]$，\textbf{但这\emph{绝不}意味着“总体里没有错”}。没抽到错误，可能只是错误稀少、样本不够大。

要正确处理“零错误”，得用\textbf{离散的超几何模型}精确算：固定总体规模$N$、并\textbf{假设其中真有$M$条错}（$M$未知），令$X$为“简单随机无放回抽$n$条后，\textbf{恰好发现的错误数}”。事件“\textbf{一条错都没发现}（$X=0$）”意味着抽出的$n$条\textbf{全部来自那$N-M$条正确标注}，其概率（超几何）为：
\begin{equation}
  \mathbb{P}(X=0\mid M)
  =
  \frac{\binom{N-M}{n}}{\binom{N}{n}}.
  \label{eq:knowledge-quality-assessment-correction-zero-hypergeometric}
\end{equation}
这个公式怎么用来做“零错误时的保守上界”？逻辑是反问：\textbf{“假设总体里\emph{真有}$M$条错，那我\emph{一条都没抽到}的概率有多大？”}$M$越大，这个概率越小（错误越多，越不可能全躲掉）。于是在整数候选$0\leq M\leq N$里，\textbf{找出“使‘零发现’概率仍\emph{不小于}$\alpha$”的\emph{最大}那个$M$}——它就是置信水平$1-\alpha$下的\textbf{单侧“错误总数上界”}，再除以$N$得到\textbf{错误率上界}。

为什么取“\textbf{最大}的$M$”？因为我们要\textbf{保留所有“\emph{尚未被‘零错误’这个观察\emph{排除}}’的候选真值”}，给出一个\emph{保守、不冒进}的上界，而\textbf{不是挑那个最好看（最小）的$M$}自欺欺人。一个常用对照：\textbf{有放回、各条独立（Bernoulli/二项）抽样}时，由$(1-p)^n\geq\alpha$反解，错误率上界恰为$1-\alpha^{1/n}$——它是\textbf{精确二项（有放回）区间的边界情形}\scite{knowledge-quality-assessment-correction-clopper1934}。注意\textbf{这两个区间对应\emph{不同抽样机制}}（无放回超几何 vs 有放回二项），而且\textbf{都\emph{不能}直接搬到分层或整群设计上}（簇相关改变分布，需用对应设计的区间）。

在结束抽样区间前，要\textbf{清醒地列出它\emph{管不到}的东西}：无论区间多漂亮，它\textbf{都\emph{不}自动覆盖}“\textbf{参考依据本身错了}（复核也会错）、\textbf{无响应偏差}（难判的系统性不答）、\textbf{抽样框遗漏}（某些对象$\pi_i=0$根本进不来）、\textbf{规范歧义}（现行规范允许多解）”这四类问题——它们不是“多抽点样、区间算准点”能解决的。

\textbf{“无法裁决”还带来另一层、与抽样无关的不确定}：注意，\textbf{即使你把每条标注\emph{全部}审一遍（没有抽样了），只要还有“无法裁决”的条目，真实错误率\emph{仍然}无法唯一确定}——因为那些条目可能对、也可能错。一个极简例子把“抽样不确定”和“裁决不确定”彻底分开：$10$条标注\textbf{全部审过}，确认\textbf{$1$条错、$7$条对、$2$条无法裁决}。这里\textbf{没有“没抽到”的对象}（普查了），所以$10\%$～$30\%$这个范围\textbf{\emph{完全}来自那$2$条结论未定}，\textbf{不是抽样波动}：
\begin{itemize}\item \textbf{最乐观}：那$2$条其实\textbf{都对}，则错误率$=1/10=\mathbf{10\%}$；\item \textbf{最悲观}：那$2$条其实\textbf{都错}，则错误率$=3/10=\mathbf{30\%}$。\end{itemize}
在“\textbf{已裁决结论可信}”这个假设下，真实错误率只能被夹在$[10\%,30\%]$之间，无法再缩窄。

把上面的思想推广到一般总体：令$U$为“\textbf{已有候选标注的\emph{适用}标注总体}”，规模$N$；仍\textbf{假定已裁决的结论是可信的}，用$j_i=1$表示该条可裁决、$j_i=0$表示无法裁决。做法是\textbf{把无法裁决条目\emph{分别}按“\emph{全对}”和“\emph{全错}”两个极端处理}，得到错误率的\term{识别范围}（identification bounds）——注意这里的“识别/identification”是计量经济/统计里“\textbf{不加额外假设时，参数\emph{能被确定}到什么范围}”的意思，与第6章可识别性同源：
\[
  L=\frac{1}{N}\sum_{i:j_i=1}e_i,
  \qquad
  U_{\mathrm{err}}=L+\frac{1}{N}\sum_i(1-j_i).
\]
读懂两个端点：\textbf{下端点$L$只数“\emph{已确认}错误”}（无法裁决的先当全对），\textbf{上端点$U_{\mathrm{err}}$再把“\emph{全部}无法裁决条目”\emph{都算成错}加进去}。这个$U_{\mathrm{err}}$是“错误率”的上界，\textbf{记号上别和总体集合$U$混淆}（一个是样本总体、一个是错误率上端）。

两个适用要点：其一， 上面$L,U_{\mathrm{err}}$是在“普查、所有可裁决都判完”时的确定范围；若\textbf{只查了\emph{概率样本}}，则要把$e_i,j_i$按设计权重（$\pi_i$、$W_h$）加权得到$\widehat L,\widehat U_{\mathrm{err}}$，此时\textbf{两端点\emph{还各带抽样误差}}，它们\textbf{只是“端点的\emph{估计}}”，不再是确定范围。其二， 要把它变成“覆盖保证”的置信区间时，\textbf{应给“下端点、上端点”\emph{分别}分配误差概率}（如各分$\alpha/2$），构造一个能\textbf{同时覆盖真实$L$和真实$U_{\mathrm{err}}$}的区间，不能只对中点做一个对称正态区间。把“零错误”和“识别范围”合起来看一个极易误判的例子：从$1000$条里\textbf{简单随机抽$100$条}，这$100$条\textbf{全部可裁决、且\emph{一条错都没有}}。此时点估计和识别端点似乎都是$0$，\textbf{但绝不能报告“总体错误率的$95\%$置信区间是$[0,0]$”}——用式~\eqref{eq:knowledge-quality-assessment-correction-zero-hypergeometric}（无放回超几何）反算，\textbf{在“抽$100$条零错”的观察下，$95\%$单侧“总体错误\emph{总数}上界”仍有约$28$条，即错误率上界约$2.8\%$}。意思是：\textbf{总体里\emph{完全可能}还藏着多达约$2.8\%$的错，只是这次没抽到}；“抽$100$条全对”远不足以证明“整体干净”。

除了错误率，还可以\textbf{对一部分参考结果\emph{再}做一次独立复核}，用来检查\textbf{裁决本身的稳定性}（不同复核者/不同证据，结论一不一致、有没有主观性）。当发现\textbf{某个分组只审了寥寥几条、结论却很极端}时，有\textbf{三种合理但不同的选择，应\emph{按用途}决定}，没有默认答案：\textbf{合并分组}（但\emph{会改变正在估计的问题}，两组一合问题口径就变了）、\textbf{追加审计}（提高把握、但增加成本）、\textbf{暂缓下结论}（诚实保留不确定性）。

最后两条\textbf{“数据窥探”陷阱}，在质量审计里尤其要守纪律：
\begin{itemize}\item \textbf{不要“反复看区间、\emph{直到}它达到门槛”才停}：区间的$95\%$覆盖概率是在“\emph{事先固定样本量、只看一次}”下成立的；\textbf{反复根据当前结果决定抽不抽、停不停}，会让这个覆盖解释\textbf{失效}（你总能凑到一个看似达标的瞬间）。\item \textbf{同时在\emph{很多}类别/分组里搜索}时，\textbf{纯偶然冒出“极端好/极端坏”结果的机会\emph{自然增多}}（多重比较问题）——某层碰巧$0$错误，可能只是运气，不代表它特殊。\end{itemize}
所以\textbf{审计方案应\emph{事先}约定}：主要指标是哪些、要单独报告哪些分组、什么条件下停止；\textbf{探索性}的额外发现要\textbf{明确标注“这是探索结果”，并用\emph{新证据}复核}，不能当成事先假设的确认。

反过来\textbf{规划样本量}则可以\textbf{从“我需要多窄的\emph{区间半宽}}”\emph{反推}该抽多少，但代入的\textbf{必须是“实际抽样设计”的方差公式和“预期的各组错误率”}（不能用简单随机公式套整群设计）。这里有个常用量叫\term{设计效应}（design effect）：\textbf{“本设计（分层/整群/不等概率）的方差”相对“同样本量下简单随机抽样方差”的\emph{比值}}。\textbf{整群相关强时设计效应$>1$（要抽更多才能达到同等精度），高效分层可能$<1$}，它\textbf{定量描述“你的抽样安排\emph{改变了多少抽样波动}}”，是把样本量和实际设计挂钩的关键参数。

\section{知识问题的检测与定位}
\label{sec:knowledge-quality-assessment-correction-detection}

再次把两件\textbf{目的不同}的事并排：\textbf{总体审计}回答“\emph{目标总体质量\emph{怎么样}}（有多少错）”——求“面”，要概率抽样、能推总体；而本章\term{问题检测}（issue detection）回答“\emph{复核预算有限，\emph{先看哪些对象}最划算}”——求“点”，要\textbf{尽快捞出具体问题}，不追求代表性。

一个\textbf{检测系统}的产物是一份\textbf{“待复核清单”}：对每个可疑对象列出\textbf{其一，可疑对象是谁、其二，\emph{触发依据}（哪个信号报的警）、其三，可能的\emph{影响范围}、其四，建议走哪个核验入口}。\textbf{请特别注意它\emph{不做}什么：它\emph{不直接裁决}“这条确实错”}——它只“报警”，裁决永远是7.1那套独立证据的事（呼应章首“检测线索并不等于错误证明”）。

\textbf{怎么\emph{评价}一个检测方法好不好}？不能只看谁的风险分数高，而要在\textbf{固定复核预算}下比较三件可量化的事：\textbf{其一，确认了多少\emph{真}问题（命中/精确率）、其二，在\emph{不同分组}（各类别、区域、来源）覆盖得均不均、其三，对“\emph{没入选}区域”随机抽查后，\emph{估计还漏掉多少}（漏报）}。一个只把最显眼问题挑出来、却整片忽略冷门区的检测器，分数再高也不算好。

\subsection{基于模型行为的问题检测}
\label{sec:knowledge-quality-assessment-correction-model-detection}

最自然的检测信号是“\textbf{模型和标签\emph{意见不一致}}”：一张标着“山茶”的照片，模型却预测“茶梅”。但\textbf{分歧本身是\emph{多义}的——它\emph{不能}告诉你“到底是标签错了，还是模型错了}”，两种解释都成立，所以它\textbf{只能让这条\emph{进入复核候选}}。同类的“模型行为异常”信号还有：
\begin{itemize}\item \textbf{高损失}：在现有标签下，模型怎么学都在这条上犯大错（可能标签误导，也可能样本本身就是难例）。\item \term{遗忘}（forgetting）信号：训练过程中模型\textbf{早先\emph{一直}和观测标签一致、后来却变得不一致}（注意：这里“遗忘”是\textbf{以“现有标签”为参照}的行为描述，\textbf{\emph{不}代表已经知道真实类别}——也可能模型只是变得更对了）。\item \textbf{表示离群}：该样本的表示向量\textbf{明显远离同类标注的表示簇}（可能错标，也可能只是罕见样本）。\end{itemize}这些模型信号有个\textbf{内生偏差}：模型\textbf{是拿这些标签\emph{训练}出来的}，它\textbf{倾向于“\emph{记住}训练标签”}——错标签它也会慢慢拟合，于是\textbf{在训练数据上评估}“模型同不同意标签”会\textbf{系统性地\emph{低估}分歧}（自己考自己）。为了\textbf{减少“模型见过待查标签”的影响}，有两种标准做法：
\begin{itemize}\item \textbf{留出预测}：\textbf{用\emph{没拿该样本}训练的模型}来给它打分（它没记住这条，分歧更真实）。\item \term{交叉验证评分}（cross-validated scores）：\textbf{把样本分成若干\emph{折}，每一条都由“\emph{只在其他折}训练”的模型评分}，于是每条都有一个“对它而言是外部”的模型判断。\end{itemize}
做交叉折有一条\textbf{不能违反的“\emph{同组同折}}”纪律：\textbf{同一对象、以及它的\emph{近重复}样本（近重复照片、同图多框）必须\emph{放进同一折}}——否则一张图在训练折、它的近重复在测试折，等于\textbf{测试时泄题}，分数又不可信。最后是数据纪律：\textbf{如果之后又\emph{用最终审计结果反复}挑模型、挑特征、调阈值}，那么那份审计就\textbf{被“用进了方案选择”，\emph{不能再}当独立评价}——最终评测必须另留（贯穿全书的留出纪律）。

% 两条边注连续排放，按整组高度为本段预留空间。
\Needspace{95mm}
但即使把“模型见过标签”这个干扰\textbf{用交叉折\emph{隔离}掉}，\textbf{信号的\emph{多义性}仍然\emph{没有}消除}——这是7.3最核心的认知，请务必记住：\textbf{同一个“高损失/分歧/离群”信号，至少有四种互斥解释}：
\begin{itemize}\item \textbf{标签真的错了}（这是我们想找的）；\item \textbf{标签对、但样本\emph{罕见/难}}（罕见正样本本来就学不像同类，损失高是正常的）；\item \textbf{模型\emph{能力不足}}（它就是学不会这个难但正确的模式）；\item \textbf{分布变了}（样本来自模型没覆盖好的新区域，不是标签问题）。\end{itemize}
所以\textbf{光凭模型信号\emph{无法}区分这四种}，必须回到\textbf{独立证据（原图、规范、采集记录）}去判。两个有代表性的方法，注意它们\textbf{都只“排序/圈区域”、不直接定罪}：
\begin{itemize}\item \term{置信学习}（confident learning）：\textbf{用模型的\emph{预测概率}和“观测到的标签”去\emph{估计}潜在的错标、再给样本\emph{排序}}。它依赖第6章讲过的“\textbf{类别条件噪声}”设定——\textbf{给定真实类别后，错标机制\emph{不再依赖}输入特征}（错标只看真类、不看长相）；而且它的\textbf{推断还\emph{依赖}“预测概率估得准（校准等）”这些条件}，条件不满足结论要打折\scite{knowledge-quality-assessment-correction-northcutt2021}。\item \textbf{训练动态地图（data maps）}：\textbf{利用\emph{训练过程中每一轮}的预测输出}，把样本按“学得\emph{多快、多稳}”分成“易学/模糊/困难”等区域。其中\textbf{“困难/模糊”区域\emph{可能富集}标签错误}，但\textbf{\emph{不能}把整个困难区域\emph{整体}当成错误集合}——里面混着大量“\emph{正确但天生难学}”的样本\scite{knowledge-quality-assessment-correction-swayamdipta2020}。\end{itemize}

图~\ref{fig:knowledge-quality-assessment-correction-detection-signal}把\textbf{同一个高损失信号\emph{画成三条岔路}}，正是上面“多义性”的可视化：\textbf{每条虚线都\emph{只}表示“这条\emph{值得核查}}”（一个候选），\textbf{不能}在虚线上直接写“错”；要走到右边，用\textbf{不同的证据检查}，才能\textbf{分别}决定：该监督信号是“\textbf{应修正}（标签错）”、“\textbf{应保留}（罕见/正确难例）”、还是“\textbf{问题在\emph{检测模型}本身}（模型能力/分布），从而\emph{限制}检测器解释范围”。

% 独立浮动页避免宽图回浮至长引文段上方，挤掉边栏的实际可用高度。
\begin{figure*}[p]
  \centering
  % 图7-2：检测信号产生候选解释，独立证据才形成裁决。
\begin{tikzpicture}[
  x=1mm,
  y=1mm,
  font=\figurefont\footnotesize,
  text=BookInk,
  signal/.style={draw=BookAmber,fill=BookAmber!9,line width=1.1pt,rounded corners=1.2mm,text width=24mm,minimum height=19mm,align=center,inner sep=1.8mm},
  explanation/.style={draw=BookMuted,fill=BookPaper,line width=0.8pt,rounded corners=1.2mm,text width=30mm,minimum height=12mm,align=center,inner sep=1.5mm},
  check/.style={draw=FigureInput,fill=FigureInput!7,line width=0.8pt,text width=30mm,minimum height=12mm,align=center,inner sep=1.5mm},
  verdict/.style={draw=FigureOutput,fill=FigureOutput!6,line width=0.8pt,text width=23mm,minimum height=12mm,align=center,inner sep=1.4mm},
  candidate/.style={knowledge arrow,draw=BookAmber,line width=0.8pt,dashed,rounded corners=1mm},
  evidence/.style={knowledge arrow,draw=BookBlue,line width=0.9pt},
  decision/.style={knowledge arrow,draw=BookTeal,line width=1.0pt},
  edge label/.style={fill=white,inner sep=0.8pt,font=\figurefont\scriptsize,text=BookMuted,align=center}
]
  \node[signal] (signal) at (15,20) {高损失或\\预测分歧\\{\scriptsize 检测线索}};

  \node[explanation] (wrong) at (58,40) {候选解释一\\样本标注确有错误};
  \node[explanation] (rare) at (58,20) {候选解释二\\样本正确但罕见};
  \node[explanation] (model) at (58,0) {候选解释三\\模型能力或分布不匹配};

  \node[check] (wrong-check) at (108,40) {核对原始对象、规范\\与独立参考证据};
  \node[check] (rare-check) at (108,20) {检查类别频率、覆盖\\与专业复核结果};
  \node[check] (model-check) at (108,0) {更换模型或预测输入条件\\比较留出行为};

  \node[verdict,draw=BookAmber,fill=BookAmber!9] (repair) at (151,40) {确认错误\\进入修正};
  \node[verdict] (keep) at (151,20) {样本标注可保留\\维护稀有覆盖};
  \node[verdict,draw=BookMuted,fill=white] (limit) at (151,0) {不能据此改标\\限制检测解释};

  \draw[candidate] (signal.east) -- (wrong.west);
  \draw[candidate] (signal.east) -- (rare.west);
  \draw[candidate] (signal.east) -- (model.west);
  \node[edge label,anchor=south] at (34,22) {待核查解释};

  \draw[evidence] (wrong) -- (wrong-check);
  \draw[evidence] (rare) -- (rare-check);
  \draw[evidence] (model) -- (model-check);
  \draw[decision] (wrong-check) -- (repair);
  \draw[decision] (rare-check) -- (keep);
  \draw[decision] (model-check) -- (limit);

  \node[anchor=north,font=\figurefont\scriptsize,text=BookMuted,align=center] at (84.5,-11) {同一个排序分数不能区分三种原因；裁决来自右侧检查，而不是来自分数本身。};
\end{tikzpicture}

  \caption{同一检测信号可能对应不同原因。高损失或预测分歧只把样本送入候选解释；错误标签或约束、正确但罕见的样本和模型能力不足分别需要原始证据核验、分组与覆盖检查、模型对照。只有独立证据确认监督信号错误后才进入修正，箭头不表示三个解释已经由分数证实。}
  \label{fig:knowledge-quality-assessment-correction-detection-signal}
\end{figure*}
\FloatBarrier

所以\textbf{模型在这一步的角色被严格限定为“\emph{排序器}}”：它\textbf{读取原始监督信号、留出预测、损失或训练日志}，吐出的是\textbf{“待复核队列 + 触发理由”}，\textbf{不产出对错结论}。这给“开发检测器”立了条规矩：评估集\textbf{必须\emph{同时}包含三类区域——真错标、\emph{罕见但正确}的样本、\emph{模型失效}区域}，不能只在“人为翻转标签”上测。因为\textbf{只测人工翻转的标签，\emph{只能}说明检测器会不会抓“人造噪声”}，\textbf{说明不了}它\textbf{怎么处理\emph{真实的歧义}（本来就模棱两可）和\emph{分布变化}（新区域）}——而后者才是实际部署里最要命的误报来源。

这个排序思路对其他监督信号也成立：\textbf{关系/配对目标}可用“\emph{配对相似度}异常”排序，\textbf{表示目标}可用“\emph{表示退化}、对扰动的\emph{响应}异常”排序；但无论信号是什么，\textbf{最后的裁决仍须回到“对象标注对不对、构造条件守没守”的独立核对}，模型分数不替代证据。

检测阈值怎么定？\textbf{顺序不能反}：应\textbf{先规定“复核预算（能查多少条）”和“\emph{绝不能漏}的风险类型”}（比如某种高危病害漏检代价极大，必须全进），\textbf{再去比较“不同阈值\emph{会送进}哪些对象}”，在预算内尽量覆盖高风险，而不是先看分数随意切。阈值定下来后，用\textbf{两个互补指标}评价这次排查：
\begin{itemize}\item \textbf{命中率/精确率}：\textbf{入选对象里，经\emph{证据确认}真错的比例}（章首那$40/100$）——衡量“选得准不准”。\item \textbf{未入选区域的遗漏}：\textbf{对“\emph{没被选中}”的那大批做\emph{概率抽查}}，估计还漏了多少错——衡量“有没有大面积放过”。\end{itemize}
两个成本提醒：其一，\textbf{留出/交叉验证评分不是免费的}，要算\textbf{训练折、推断}的开销，\textbf{不能只计“人工复核时间”}；其二，\textbf{模型\emph{无法解析/无法响应}的预测请求（模型拒答、超范围），\emph{不能}因此就从审计总体里“消失”}——它们应\textbf{另安排一条检查路径}（转人工或专门规则），否则模型不答的对象永远不会被质检。

\subsection{训练样本贡献的评价与归因}
\label{sec:ka-sample-valuation}

7.3.1 的高损失回答了“\textbf{\emph{哪些标签值得复核}}”（怀疑它错），但它\textbf{\emph{没有}回答另一个对“要不要动手改/删”同样重要的问题}：\textbf{“改动\emph{哪一条}训练样本，会真正\emph{改变}任务结果？”}这是两个\textbf{独立维度}，而且常常不一致：
\begin{itemize}\item \textbf{两张\emph{都正确}的近重复照片}：它们\textbf{谁也没错}，但对训练高度冗余，\textbf{可以\emph{相互替代}}（删一张对结果几乎无影响）。\item \textbf{一张\emph{正确}的稀有物种照片}：标签完全对，但它可能\textbf{只对极个别类别有用}、对总体指标贡献很小。\item \textbf{一条\emph{错}标}：它当然该改，但\textbf{也可能早已被模型\emph{忽略}}（模型根本没拟合这个噪声点），改了结果也不变。\end{itemize}
所以必须把\textbf{“标签\emph{对不对}}”（质量问题）和\textbf{“这条样本对\emph{指定任务结果}有多大\emph{作用}}”（影响/贡献问题）\textbf{分开}——错的未必重要，重要的未必错。由此定义\term{训练样本贡献}（training-sample contribution）：\textbf{固定\emph{学习算法}和\emph{评价规则}之后，“加入、删除或改动这条训练样本，会\emph{怎样改变}训练后的表现”}。请注意\textbf{它\emph{没有脱离这些条件}的“永久价值}”——换个算法、换个任务，同一张照片的贡献可以完全不同。

要把“贡献”定义得既公平又能处理“冗余、互补”，借用博弈论里的 Shapley 值思想。先把记号立起来：训练样本索引集$N=\{1,\ldots,n\}$（$n>0$），$A(S,U)$表示“\textbf{用\emph{子集}$S\subseteq N$、加上算法随机种子$U$训练出的模型}”。两个约定保证它处处有定义：\textbf{空集$S=\varnothing$约定为一个\emph{固定基线预测器}}（不训练时的默认模型，而不是随机乱猜）；\textbf{即使子集里缺某些类别，仍保留\emph{完整类别头}、按\emph{同一规则}训练}（避免“缺类就没法定义效用”）。

然后\textbf{先固定一批\emph{独立}验证样本和一条\emph{效用规则}}（比如$V=$“验证集上平均损失的相反数”，效用越大越好）；训练随机性$U$要么\textbf{固定成同一组预定随机数}（得到“给定种子下”的贡献），要么\textbf{对约定的$U$分布求平均}（得到“平均意义上”的贡献），\textbf{所有子集都用同一约定}。注意\textbf{“重复训练很多次”只能\emph{近似}后一种平均，不能把\emph{没控制住}的随机波动当成“样本的作用}”。

为什么不能简单用“\textbf{从\emph{完整}集里删掉$i$}”的效用差$V(N)-V(N\setminus\{i\})$？因为它\textbf{只看了$i$在“\emph{所有其他样本都在}”这\emph{一种}上下文里的作用}，\textbf{看不到$i$在\emph{别的}训练子集里}（比如没有那些冗余近重复时）的作用，而贡献恰恰随上下文变。\term{Data Shapley}的修正是\textbf{“\emph{公平地平均所有\emph{加入次序}}”}：设想样本按随机顺序一个个加入训练集，看$i$\textbf{在它\emph{被加入的那一刻}带来的\emph{边际增量}}，再对所有次序平均。令$\sigma$等概率遍历全部$n!$种排列，$S_{\sigma,i}$为排列中排在$i$前面的样本集合，则
\begin{equation}
  \phi_i
  =\mathbb E_{\sigma}
    \bigl[V(S_{\sigma,i}\cup\{i\})-V(S_{\sigma,i})\bigr].
  \label{eq:ka-data-shapley}
\end{equation}
理解这个定义的两个要点：其一，\textbf{每次比较都\emph{重用同一}学习算法和同一验证效用}，\textbf{只改变“用哪些训练样本（子集）”}——这样效用差才干净归因于样本；其二，\textbf{Data Shapley 是把经典 Shapley 的“公平分配”思想\emph{用到训练数据}上}，\textbf{分配的是“\emph{训练样本}对任务效用的贡献”，\emph{不是}“特征对某条当前预测的作用}”（那是特征归因/SHAP），两者别混。计算上$n!$种排列不可能枚举，实践用\textbf{随机采样一部分排列来\emph{近似}}；\textbf{一旦用“\emph{截断}训练（每个子集只训几步）”或“\emph{只近似}训练”去省算力，近似本身\textbf{还会额外改变误差}}，要单独评估\scite{ka-ghorbani2019-data-shapley}。

用\textbf{两张“\emph{可相互替代}的正确照片}”把Shapley的“公平”算明白。设$V(\varnothing)=0$（没有照片时效用为0，即基线）、$V(\{1\})=V(\{2\})=V(\{1,2\})=1$（每张单独都能把效用带到满值1，两张一起也还是1——它们完全冗余）。
\begin{itemize}\item \textbf{“留一/删一”值}：从\emph{完整}集删任一张，效用都还是1、不变，所以\textbf{两条的“留一”贡献\emph{都是}0}。\item \textbf{Shapley（跨所有次序平均）}：看两种加入顺序——排列$(1,2)$中，1先加入带来$1-0=1$、2后加入带来$1-1=0$，增量对为$(1,0)$；排列$(2,1)$对称为$(0,1)$。\textbf{两种次序一平均，每张各得$1/2$}。\end{itemize}
这里有个\textbf{看似矛盾、实则关键}的点：既然“\textbf{在\emph{已有照片1}时再加照片2，增量是0}”，为什么照片2的 Shapley 还是$1/2$而不是0？因为\textbf{Shapley 还要算“\emph{照片1\emph{不在}}（排列2先到）”的那种次序}——那时照片2是\emph{单独}把效用从0带到1的功臣，贡献为1；\textbf{跨“它在/不在”两种上下文平均，才得到$1/2$}。这正好说明：\textbf{“删一试验（leave-one-out）”回答“在\emph{当前完整}模型里它多不多余”，“Shapley 跨子集分配”回答“\emph{平均而言}它对训练有多大功劳”，是两个不同问题}，数值不同不矛盾。

Data Shapley 要为大量子集\textbf{反复\emph{完整重训}}，大模型上算不起。如果\textbf{问题只针对“\emph{某一次具体预测}会怎么变”}（而非全体任务效用），可以\textbf{不重训、用\emph{训练后模型的局部近似}}：\term{影响函数}（influence function）通过\textbf{对“训练样本权重”做微小扰动}，\textbf{局部线性地近似“模型参数、进而预测”会怎么变}（用梯度/Hessian 之类）。但它\textbf{带着一串数学条件}：损失要\textbf{光滑}、要停在\textbf{局部最优解}、\textbf{曲率（Hessian）要可逆}等；\textbf{“\emph{整条}样本删除”或“训练路径切换”这种\emph{大改动}，往往\emph{破坏}局部近似}（小扰动线性近似不了大跳变）。推导见\BookChapterRef{chap:trust-interpretability}{卷三《基础、理论与可信性》的“可解释性”章}。

两条不能越界的结论：其一，\textbf{这种局部敏感性\emph{不等于}Shapley 的“跨所有联盟（子集）平均”}——一个是局部微分近似、一个是全局公平分配；其二，\textbf{无论影响多大，都\emph{不能}由“这条样本\emph{有影响}”推出“它的\emph{标签有错}}”——高影响只说明“改了会变”，变的方向可好可坏。

Data Shapley 和影响函数在\textbf{超大模型}上还有个现实障碍：\textbf{模型参数梯度\emph{维数极高}}（几十上百亿维），存不下、算不动。\term{TRAK}（Tracing with the Randomly-projected After Kernel）的核心思路是\textbf{先\emph{降维}再归因}：\textbf{把“指定预测对模型参数的梯度”\emph{投影到一个低维空间}}（用随机投影把维度压到可算），再配合\textbf{局部线性化 + 一步 Newton（牛顿步）近似}，得到“训练样本对这条预测的归因分数”。以分类为例，给定样本$z=(x,y)$，取模型对类别$y$的对数几率$f(z;\bm\theta)=\log[p_{\bm\theta}(y\mid x)/(1-p_{\bm\theta}(y\mid x))]$，要求概率严格位于0与1之间。对一个固定的已训练模型，令随机投影$\bm P\in\mathbb R^{D\times r}$，$D$为模型参数维数，$r$为投影维数；$\bm\varphi(z)=\bm P^\top\nabla_{\bm\theta}f(z;\bm\theta)\in\mathbb R^r$。将训练样本的投影梯度按行排成$\bm\Phi$，并令$\bm Q=\operatorname{diag}(1-p_{\bm\theta}(y_i\mid x_i))$，其单模型近似的核心为
\begin{equation}
  \bm\tau(z)
  =\bm\varphi(z)^\top
   (\bm\Phi^\top\bm\Phi)^{-1}\bm\Phi^\top\bm Q.
  \label{eq:ka-trak-core}
\end{equation}
读这个核心式（式~\eqref{eq:ka-trak-core}）三块的直觉（这里假设$\bm{\Phi}^{\top}\bm{\Phi}$可逆，实际\textbf{要靠加正则项保证}，换求解规则需另说近似）：
\begin{itemize}\item \textbf{投影梯度$\bm{\varphi}(z)$}：刻画“\textbf{这条指定预测\emph{对模型参数}的局部敏感方向}（降到低维后）”——预测最吃哪些参数变化。\item \textbf{中间的$(\bm{\Phi}^{\top}\bm{\Phi})^{-1}$}：这是一个\textbf{“重复度/协方差”校正}——\textbf{调整“这些敏感方向在\emph{训练样本}里\emph{重复出现}的程度}”：如果一批训练样本在某方向上高度雷同（冗余），就不能各算一份，要用逆矩阵压权重。\item 最后的$\bm Q$：按\textbf{模型给出的类别概率加权}（对角阵，对角线是$1-p$）。\end{itemize}
还要知道这个式子\textbf{只是 TRAAK 的“\emph{单模型核心}”，\emph{不是}全部实现}：原方法还\textbf{组合“\emph{多个}子集训练（类似折/集成）得到的梯度关系”和类别权重}，并可做\textbf{软阈值}处理以去噪/降本\scite{ka-park2023-trak}；省略这些步骤时要说明自己用的是简化版。

用一个一维教学例子看清 TRAK 分数的量纲：投影只有一维、两条训练样本的梯度都为1、目标梯度也为1，且$\bm Q=\operatorname{diag}(1/2,1/2)$。此时$\bm{\Phi}^{\top}\bm{\Phi}=(1,1)\cdot(1,1)^\top=2$，代入得每条分数$\tau(z)=\tfrac12\cdot\tfrac12\cdot(1\cdot\tfrac12+1\cdot\tfrac12)=\tfrac14$，即两条各$1/4$。

\textbf{请务必不要把这个$1/4$和前面 Shapley 的$1/2$当成“同一东西的两种算法”}——它们数字不同是因为\textbf{计量的对象根本不同}：\textbf{TRAK 计量的是“\emph{指定（已训练）模型}对某条预测的\emph{局部线性}响应”}（站在训好的模型附近做微分）；\textbf{Shapley 分配的是“\emph{训练后任务效用}在数据上的\emph{全局}公平份额”}（要跨子集重训）。两者不可互证。那怎么\textbf{验证 TRAK 这类近似归因可不可信}？做法是：\textbf{按“子集归因分数\emph{之和}”去\emph{预测}“用这个子集训练后模型响应怎么排”}，然后\textbf{真的\emph{按预定比例抽子集、\emph{真正}重训}}，比较两者响应\emph{排序}的\emph{秩相关}；\textbf{不能}只凭“排名靠前的照片\emph{看起来}像相关样本”这种主观印象就采信。

贡献分数用起来的\textbf{共同纪律}（和检测信号一样）：
\begin{itemize}\item \textbf{留出纪律}：一旦\textbf{验证样本被用来“\emph{筛选}训练样本”或“\emph{设计}修正方案”}，它就参与了方案选择，\textbf{最终测试必须\emph{另留}}。\item \textbf{成本要算全}：Shapley 的\textbf{反复重训排列}、TRAK 的\textbf{梯度计算、随机投影、矩阵求逆、还要为\emph{多个模型（模型集）}重复}，都不便宜，别只看人工。\item \textbf{低贡献 并不等于 错标}：一条样本贡献低，\textbf{可能只是\emph{冗余}（有近重复）、可能是\emph{评价范围没覆盖它的稀有类}、也可能是\emph{模型能力不足}没用上它}——这些都\textbf{不直接裁定“它标错了”}。\item \textbf{归因 并不等于 事实支持}：在大模型里，能\textbf{解释“模型为什么生成某句话”（高归因）}，\textbf{\emph{不}等于“训练文本里\emph{确实支持}这句话这个事实”}（影响高不等于出处正确，更不等于内容属实）。\end{itemize}
一句话收口：\textbf{贡献分析只负责“为\emph{复核和干预}提出\emph{候选}（先看谁、可先删谁）”；至于“原始标签到底对不对”，\textbf{仍要回到\emph{独立证据}裁决}}，两个流程各司其职。

\subsection{基于多源关系的问题检测}
\label{sec:knowledge-quality-assessment-correction-multisource-detection}

看到“\textbf{两条标注给同一植株\emph{不同}答案}”，新手最容易立刻记为“冲突/有错”，但\textbf{第一反应应是“它们\emph{到底是不是在回答同一个问题}}”——很多“分歧”其实是问题不同，根本不矛盾：
\begin{itemize}\item \textbf{时间不同}：同一植株\textbf{春季的开花状态和秋季可以\emph{不同}}（春天开花、秋天结果），两条都对。\item \textbf{粒度不同}：一条给\textbf{物种}、另一条给\textbf{属}，\textbf{可能\emph{同时成立}}（属是上位，物种落在该属里）。\end{itemize}
所以\textbf{形成“冲突线索”\emph{之前}必须先回到第6章的三层对齐：对象、语义、粒度、上下文全部对上}。只有当\textbf{“\emph{同一对象、同一时点、同一问题}”}（比如同一时刻问“这朵花开没开”，一个说是、一个说否），\textbf{才是真正\emph{不能同时成立}的冲突}。反过来，\textbf{同一条件下本就允许多个合理答案}，两答不同\textbf{不构成冲突}。\textbf{重复标签之间的差异、以及结构约束违例，也都按这条原则先核对，而\emph{不直接}记成“标注错误”}。

这条\textbf{在第6章已反复强调、但质量检测时最容易忘}：\textbf{“来源\emph{数量多}”\emph{绝不是}证据独立性的保证}。十个网页可能\textbf{都复制了同一份错误名录}，于是“十比零多数票”只是\textbf{把同一个错误\emph{数了十遍}}，一个独立的原始采集日志（哪怕只有一份）反而可能是更强证据。因此\textbf{冲突表不能只存“谁赞同、几票”}，必须同时保存\textbf{来源依赖关系、来源版本、时间、证据路径}（否则看不出十票其实是一票）。

还有一个\textbf{“一致”的盲区}：\textbf{当所有来源\emph{碰巧}一致时，系统最容易\emph{取消}审计（觉得没冲突=没问题）——这是危险的}。因为\textbf{共享同一个错误来源时，大家会\emph{一起错、还不产生任何冲突信号}}（异常检测正是靠分歧报警，一致的错它根本看不见）。所以\textbf{无论来源一不一致，\emph{随机抽查}都不能省}——专门用来抓“\textbf{沉默的共同错误}”。

关系类的“冲突”还要\textbf{额外检查\emph{方向}和\emph{条件}}，不能只看两端对象对不对：
\begin{itemize}\item \textbf{方向不能反}：若$A$是照片、$B$是植株，$(A,\text{采集自},B)$和\textbf{把两端\emph{交换}}得到的$(B,\text{采集自},A)$\textbf{\emph{不是}同一个命题}（植株不会“采集自”照片），判错时要指出是方向错。\item \textbf{条件/时间不能无条件复制}：关系上带的时间、上下文属性，\textbf{换个版本、换个对象不能照搬}（春天的“位于”不等于秋天的）。\item \textbf{结构违例\emph{不等于}标签错}：一个\term{违例}（violation）只说明“\textbf{当前这组答案\emph{和声明的约束}\emph{不相容}}”（如BIO非法、互斥类同时被选）。到底是“某条\emph{标注}错”还是“\emph{约束本身}不适用这里”，\textbf{必须\emph{同时}复核“标注”和“约束的适用性”两边才能定}——可能是约束写错了。\end{itemize}
工程上来源很多时，可\textbf{先按“任务规范 + 对象标识 + 版本”把同源/同对象的标注\emph{聚合}}（先消除重复和明显对齐的），\textbf{减少大量无意义的逐条两两比较}，再\textbf{集中查“影响大的冲突”和“\emph{无法靠对齐解释}的冲突”}。

\subsection{基于来源记录的问题定位}
\label{sec:knowledge-quality-assessment-correction-provenance-localization}

当\textbf{多条错误被\emph{分别}确认}后，下一步不该一条条孤立修，而要\textbf{顺着它们的\emph{生成过程}往回找“\emph{共同的根因}}”——如果一批错都指向同一个词典版本、同一条规则、同一个模型，那是系统性问题，改一处胜过改一片。要做到这一点，前提是有\term{溯源清单}（provenance manifest）：它\textbf{记录每条结果是“\emph{谁、用什么版本、在什么配置下、什么时候}”产出的}，至少包括——
\begin{itemize}\item \textbf{样本/对象}、\textbf{人员或模型}（谁标的、哪个模型）；\item \textbf{规则与词典的版本}（用的哪一版）；\item \textbf{预测请求配置、执行状态、批次、时间戳}。\end{itemize}
有了它，结果才能\textbf{回溯到\emph{具体的处理活动}}（复现“当时到底跑了什么”）。

\textbf{定位根因的方法}：先\textbf{列出错误样本\emph{共同}用了什么}（共同的词典？同一批？同一模型版本？），再\textbf{和\emph{没受影响}的样本\emph{对照}}（它们缺了哪个共同因素），然后\textbf{\emph{复现}那个候选变更（切回旧版/新版）看错误是否随之出现/消失}，并\textbf{圈定影响\emph{范围}}（波及哪些下游结果）。\textbf{两个“别冤枉好人”的提醒}：其一，\textbf{流水线上\emph{最先}报警/最先暴露异常的环节，\emph{不一定}就是产生错误的环节}——\textbf{解析器}最先崩，可能只是上游喂了脏数据；其二，\textbf{对象错配可能在\emph{更早}的“写回数据库}”步骤就发生了（如前面S042挂错植株），报警的是下游类别，根因却在写回。

一个\textbf{“怎么\emph{排除干扰}、确证根因”的范例}。现象：\textbf{某次词典更新后，“茶梅”被\emph{大量}错映成“山茶}”。直觉会说“就是词典的锅”，但\textbf{“时间上\emph{同时发生}”只能提供一个\emph{候选}原因，不能定罪}——因为\textbf{同一时期可能\emph{还}发生了别的变化}：样本换了地区（分布变了）、评价规范改了（尺子变了），这些也能造成错误率同时上升。

\textbf{要\emph{隔离}出词典这个变量}，必须做\textbf{受控对照}：\textbf{挑\emph{同一批}、\emph{冻结}的样本}，\textbf{只切换词典版本（旧 vs 新）、\emph{保持模型、规范等其他一切不变}}，分别跑。\textbf{只有当差异“\emph{只随词典版本}出现”}，才能把根因\textbf{锁定到词典}；然后再\textbf{用原始名录\emph{核验}那条具体映射}（坐实“茶梅到山茶”这条映射是错的）。最后是\textbf{影响\emph{范围}}：\textbf{共享词典被\emph{多个规则、多个模型的预测输入}共用}，所以一旦确认词典错，\textbf{受影响的是\emph{全部}依赖它的下游结果}，不能只修“最先报警的那一个来源”——同源污染要成片排查。

溯源\textbf{强依赖日志是否完整}，所以\textbf{日志缺失会\emph{直接}限制你能下的结论}：如果\textbf{无法恢复“当时的预测请求、模型版本、随机配置”}，你就\textbf{无法受控复现}，这时只能诚实写“\textbf{这是个\emph{与该批次相关}的\emph{候选}原因}”，\textbf{\emph{不能}声称“已经\emph{找到根因}}”。

两个严谨性边界：其一，\textbf{“调用可\emph{重现}}（每次跑结果一样）”\textbf{\emph{不}自动证明\emph{因果}}——稳定复现只说明确定性，要证“是A导致B”仍需上面那种“\emph{对照变更}（切版本）”、或\emph{独立证据}、或\emph{可执行检查}；其二，定位结论本身\textbf{要\emph{分三档}写清影响面}：\textbf{\emph{已确认}受影响的、\emph{可能}受影响的、\emph{尚未覆盖}还没查的}。这个影响范围清单\textbf{正好为下一节“\emph{修正该动多大粒度}}（单条/整批/整条流水线）”提供依据——根因波及多广，修正范围就该多大。

\section{知识的修正与验证}
\label{sec:knowledge-quality-assessment-correction-repair}

进入修正阶段，先立一条\textbf{反直觉的总原则}：\textbf{修正\emph{不是}“把检测分数最高的标签\emph{换掉}}”。检测器只负责\textbf{把可疑对象和\emph{证据线索}排到面前}，\textbf{“改成什么、改不改、改多大范围”必须由\emph{已确认的问题}（7.3的根因）、\emph{可用证据}、和\emph{影响范围}共同决定}——分数从来不是修改的依据。

针对一个确认的问题，\textbf{可采取的行动\emph{不止}“改一个字”}，按问题类型可以是：\textbf{改正内容}（错标改对）、\textbf{补齐遗漏}（漏标补上）、\textbf{重建关系}（对象错配重新绑定）、\textbf{修改\emph{构造条件}}（规则/教师/变换错了就改生成程序）；而当\textbf{证据不足}时，\textbf{正确的“行动”恰恰是\emph{保留不确定性}}（标记未决、降权、隔离），而不是硬改一个答案。

无论采取哪种行动，\textbf{每一次变更都要可追溯、可回退}：必须\textbf{保存“原始值、候选（改成什么）、\emph{采用依据}（凭什么这么改）、\emph{未决问题}、\emph{回退入口}（怎么撤回）”}。最后\textbf{给“改了就会变好”泼冷水}：\textbf{“\emph{改变了}监督信号或生成程序”只说明“\emph{系统发生了变化}”，它\emph{绝不}自动等于“\emph{质量已经改善}}”——改完可能更糟（改错、引入新错），\textbf{是否改善要靠7.4.2的独立验证说了算}。

\subsection{修正方案的选择与实施}
\label{sec:knowledge-quality-assessment-correction-repair-selection}

用$x_2$看清“\textbf{修正\emph{粒度}要和问题\emph{范围}匹配}”：
\begin{itemize}\item \textbf{如果只是$x_2$这\emph{一条}经独立材料确认为茶梅}（个别错标），那就\textbf{只对这一条做\emph{人工改正}}，范围精确到单条。\item \textbf{如果根因是“\emph{同一本词典的某个映射}错了，而它影响了\emph{整批}预测输出”}（系统性问题），那\textbf{只改$x_2$这一条}就是\textbf{“按下葫芦”——\emph{其他受同一错误影响的样本全被留下}}，治标不治本；正确动作是\textbf{修词典/修规则、再\emph{批量}重算受影响结果}。\end{itemize}
所以\textbf{选方案时要同时判断三件事}：其一，\textbf{“改成新值”的\emph{证据}充不充分}（没有独立证据就别自动替换）；其二，\textbf{需不需要\emph{新增专业判断}}（是不是超出当前规范、要请专家/查外部资料）；其三，\textbf{这是“\emph{任务定义/生成条件}变了”，还是“旧定义下\emph{做错了}}”——前者要换标准、后者是纠错。

把四类“处理方式”的\textbf{分工}理清，它们\textbf{可以组合、\emph{不是}互斥的四个抽屉}：\textbf{人工复核}说明“\emph{由谁执行判断}”（人而非机器）；\textbf{自动推断}说明“\emph{谁执行}”（规则/模型自动）；\textbf{补充与调整}说明“\emph{改变了什么}”（加证据/改配置）；\textbf{“无法可靠修正时怎么办}”规定“\emph{怎么保留未决}”。最后是版本纪律：\textbf{如果是\emph{任务定义本身}改变（新规范重新划界），必须\emph{另建新版本}}，\textbf{\emph{不能}把“按新标准重标的结果”\emph{混进}“旧标准下的错误率”}——新旧标准是两把尺子（呼应7.1）。

\subsubsection{基于人工复核的修正}
\label{sec:knowledge-quality-assessment-correction-human-repair}

\textbf{人工修正不能“看着模型建议打勾”}：它所需的材料\textbf{必须包括\emph{原始对象}、\emph{适用规范}、\emph{独立证据}}——\textbf{而\emph{不只是}一条模型建议}（否则只是给模型背书）。流程上沿用7.1的盲化顺序：\textbf{复核者先\emph{独立}给出新判断（不看原标签/模型），\emph{保存}后再查看原标签和冲突判断及其证据对照}，最后提交\textbf{四类结论之一并附理由}：\textbf{保留}（原判对）、\textbf{修改}（改成什么+依据）、\textbf{维持多解}（本来允许多答案）、\textbf{无法裁决}（证据/规范不足）。

还要\textbf{分清“操作错”和“\emph{领域事实}错”该谁处理}：\textbf{操作层面的错误}（框拖偏了、填错字段、可按规范机械纠正）\textbf{可以由\emph{受训复核者}处理}；但一旦涉及\textbf{需要\emph{领域事实}判断（这个物种到底是不是茶梅）或\emph{规范解释}（这条规则在这种边缘情形算不算适用）}，就\textbf{应交给\emph{相应领域专家}}，不能让受训人员猜。并且\textbf{专家\emph{可推广}的裁决（“这类情形以后统一这么判”）应\emph{写回规范}}——让同类问题以后不必反复请教，把一次性判断沉淀成可复用的规则。

对象错配这类问题\textbf{有“\emph{牵一发动全身}}”的特点——\textbf{一条绑定错误，往往\emph{连带}多条标注/关系}，不能只改眼前字段。举例：照片S042被错挂到植株P017名下，正确的修正顺序是：
\begin{enumerate}\item \textbf{先\emph{恢复照片的真实对象}}（把S042从P017解绑、挂回正确植株）——这是根；\item 再\textbf{检查S042自己的\emph{类别}}（对象对了类还对不对）；\item 检查\textbf{它和“\emph{同株}其他照片”的\emph{正对应关系}}（同一植株的多视图要一致）；\item 还要\textbf{检查所有“\emph{以P017为锚}生成的\emph{邻居/下游结果}}”（邻居关系、训练关系可能都建立在错绑上）。\end{enumerate}
\textbf{只把类别字段改对、却留着错误绑定，会留下一整片错误关系}。所以\textbf{修正范围应\emph{沿来源依赖关系}展开}（顺着7.3溯源清单波及到哪就查到哪），但\textbf{每一条受影响标注\emph{仍要}保留它\emph{自己的变更依据}}——批量改不等于“一刀切共享同一条理由”，每条仍可追溯。

\textbf{为什么坚持“\emph{先}独立复标、\emph{后}知情裁决”这个顺序不能颠倒}？因为两步在做\emph{不同}的事：\textbf{独立复标}产生的是“\emph{不受原答案影响的\emph{新增判断}}”（一份新鲜证据）；\textbf{知情裁决}则是在\textbf{已有这份新判断\emph{之后}，再\emph{比较各方意见和证据}做仲裁}。如果\textbf{先看了各方争议再“独立”判断，那份判断已被污染，\emph{不能}倒过来算成一份“盲化}判断”——顺序反了，证据等级就变了。

\textbf{多人复核仍分歧}时怎么办？\textbf{不必靠“\emph{无限}加人投票”硬凑一致}：若规范允许多解，\textbf{保留“答案\emph{集合}”和“各自\emph{成立条件}}”即可。\textbf{要不要继续加人复核}，判据不是“已经复核了几次”，而是三件实事：\textbf{其一，\emph{新证据}还能不能\emph{排除}某些候选（能排除才值得查）、其二，\emph{当前用途的风险}（高风险值得多查）、其三，\emph{追加复核的成本}}——证据不再改变结论时，继续复核只是烧钱。

\subsubsection{基于自动推断的修正}
\label{sec:knowledge-quality-assessment-correction-automatic-repair}

自动修正听起来省力，但它\textbf{有个\emph{硬前提}}：\textbf{只有当“\emph{转换规则}、\emph{外部证据}、或\emph{经过验证的预测}”能够\emph{支持}“新结果确实成立”时，才允许自动替换}。没有这个支撑的“自动改”只是把模型的猜测当成事实。先分清两种自动程度完全不同的情形：
\begin{itemize}\item \textbf{\emph{确定性}的格式转换}：比如\textbf{类别编码从旧编号\emph{无歧义}地一一映射到新编号}（旧码X恒等于新码Y，纯重命名）——这是“\textbf{在\emph{限定条件}下\emph{保持语义}的确定性转换}”，\textbf{不涉及判断对错}，可以高把握自动执行。\item \textbf{\emph{模型对疑似错标重新预测}}：模型给个新答案——这\textbf{\emph{永远只是}“生成一个\emph{候选}}”，它可能重复原标签的错、可能自己没把握，\textbf{必须再有证据/验证环节}，不能直接落库。\end{itemize}
因此这两类\textbf{\emph{不能}共用同一个“\textbf{置信度高于阈值＝修正成功}”的标准}：\textbf{确定性映射的把握来自“\emph{规则无歧义}}”，\textbf{模型重测的把握需要\emph{外部验证}}——前者不看置信度（看规则），后者不能只看置信度（要看证据）。

一个\textbf{可审计的自动修正流程}应当\textbf{顺序清晰、每步留痕}：\textbf{其一，生成候选}（规则/模型给新值）到\textbf{其二，检查\emph{适用条件}}（这个对象/范围里这条规则成立吗）到\textbf{其三，决定“\emph{自动接收}”还是“\emph{转人工}}”（高把握、确定性的自动；有歧义的转人）到\textbf{其四，\emph{保存修改前后差异}和决策理由}。不同步骤要\textbf{存不同的东西}：
\begin{itemize}\item \textbf{确定性规则}：要\textbf{写明“\emph{为什么}它\emph{保持对象和语义}”}（说明这是一一、无损、无歧义的转换，而不是判断）。\item \textbf{预测替换}：要\textbf{分别保存“模型分数、\emph{替换风险}（错改概率/影响面）、\emph{复核要求}}（哪些情况必须人工回看）”。\end{itemize}
两个“一致不等于对”的陷阱（第6、7.3反复出现，这里同样适用）：其一，\textbf{模型和原标签若来自\emph{相同训练样本}或\emph{共享来源}}，那么“模型重测结果和原标签\emph{多次一致}”\textbf{仍可能只是\emph{重复同一个错误}}（同源）；其二，\textbf{用\emph{同一个教师}重新生成表示目标}，教师和“学生”共享同样的错，\textbf{\emph{不构成}对“教师模型预测输出”的\emph{独立}证明}——自证不算证。

自动修正的\textbf{规则/阈值必须在\emph{开发样本}上确定、然后\emph{冻结}，再进入\emph{留出}审计}——绝不允许边看留出结果边调（数据纪律同前）。在冻结前的\textbf{小范围试用}阶段，要\textbf{\emph{同时}检查两组对象，它们各暴露一类不同风险}：
\begin{itemize}\item \textbf{被修改的对象}：检查\textbf{“\emph{改错}”风险}——规则有没有把本来对的改成错的（精确率）。\item \textbf{\emph{未}被修改的对象}：检查\textbf{“规则\emph{漏修}”和“\emph{作用边界}}”——该改的有没有没动、规则在什么样本上失效（召回率/适用范围）。\end{itemize}
只看被改的，会盲目乐观。最后一条\textbf{反“过拟合调参”}纪律：如果自动修正\textbf{出现\emph{系统性}副作用}（比如某一类被成片改错、某来源整体异常），\textbf{正确动作是\emph{回退到旧版本、查根因}，而\emph{不是}“继续\emph{提高阈值}，直到\emph{开发样本看起来}正确}”——后者只是让规则在开发集上自洽，副作用在新数据上依旧，甚至更隐蔽。

\subsubsection{基于问题类型的补充与调整}
\label{sec:knowledge-quality-assessment-correction-type-repair}

不同类型的问题，要求\textbf{不同性质的修正动作}，先把“\emph{不改变任务定义}就能解决”的和“\emph{连定义都要动}”的分开：
\begin{itemize}\item \textbf{照片\emph{漏标}}（该有的花没框）到\textbf{\emph{补齐}监督信号}（补框/补关系），任务定义不变。\item \textbf{\emph{对象绑定}错}（S042挂错植株）到\textbf{\emph{重建对应关系}}（解绑重挂），定义不变。\item \textbf{答案\emph{过粗}}（只到属、任务要种）到\textbf{\emph{增加细节}}（细分/补更细标签），也不改变任务。\end{itemize}
这三类\textbf{都能在\emph{既有任务定义}内修}。但如果\textbf{歧义的根源是\emph{规范本身}}，就要先分清两种性质完全不同的动作：
\begin{itemize}\item \textbf{\emph{澄清}原有规则}（原规则没说清、大家理解不一致，补一个解释）——“正确”的标准其实没变，只是明确了；\item \textbf{\emph{重新规定}类别边界}（把原来一类拆成两类/合并两类）——这\textbf{\emph{改变了“什么算正确”的含义}}，旧标注在新标准下无所谓对错，\textbf{必须\emph{分版本}评价}（7.1讲过）。\end{itemize}
还有一类问题在“\emph{构造条件}”而不在标签本身：比如\textbf{一个视图变换\emph{删掉了主体}，却还声明“语义不变”}（裁剪后只剩背景还当正样本）——这时\textbf{要修的是“\emph{变换规则}和\emph{目标构造}}”（重新生成对应、改裁剪约束），\textbf{\emph{不是}更换\emph{某一个标签}}就能解决。\textbf{一句话：处理范围由“\emph{问题发生在\emph{哪一层}}”决定}——标签错改标签、绑定错重建绑定、规则错改规则、定义变建版本。

表~\ref{tab:knowledge-quality-assessment-correction-repair-scope}是\textbf{“核查结果 到 该改\emph{什么}、改\emph{多大范围}”的\emph{对照表}}：前五行（单点错标、对象错配、标注缺失、定义变化、构造条件失效）都是\textbf{\emph{已有证据支持}、可以\emph{动手}}的问题，每一行对应不同的“改变对象与范围”；\textbf{最后一行“证据不足}”则\textbf{\emph{不是}一种修正，而是\emph{保留“无法裁决”状态}}（暂缓/降权/隔离/补证据）。

使用这张表有个\textbf{不可跳过的闸门}：\textbf{只有当“核查依据”已经是\emph{独立证据}}时，才能按表批量处理。\textbf{如果你手上\emph{仍然只有}“检测分数高”或“来源之间有冲突”}，那\textbf{还停在“线索”阶段}，\textbf{必须先去\emph{找参考证据}（原图/规范/日志/复标）}，\textbf{\emph{不能}据此就认定“某类全错”并\emph{批量替换}}——用线索当证据批量改，是修正阶段最大的误伤来源。

\begin{table*}[htbp]
  \centering
  \small
  \begingroup
  \setlength{\tabcolsep}{2.5pt}
  \begin{tabularx}{\linewidth}{@{}>{\RaggedRight\setlength{\parindent}{0pt}\arraybackslash}p{26mm}>{\RaggedRight\setlength{\parindent}{0pt}\arraybackslash}p{47mm}>{\RaggedRight\setlength{\parindent}{0pt}\arraybackslash}p{43mm}>{\RaggedRight\setlength{\parindent}{0pt}\arraybackslash}X@{}}
    \toprule
    \textbf{核查结果} & \textbf{改变对象与范围} & \textbf{可选处理} & \textbf{后续核验} \\
    \midrule
    单点错标 & 当前答案及引用它的派生样本标注 & 独立复标后替换或保留多解 & 复核改后结果，抽查关联样本标注 \\
    对象错配 & 对象绑定、同对象标签与关系 & 恢复对象后重建关联结果 & 对象、方向与影响范围 \\
    标注缺失 & 缺失监督信号及同类、同批次对象 & 补标、追加获取或保留弃权 & 覆盖变化与新增部分质量 \\
    定义变化 & 规范版本及旧定义下的相关样本标注 & 分版本保留、有限映射或重标 & 共同定义域内的前后差异 \\
    构造条件失效 & 来源规则及同配置生成的目标 & 修改规则，重生成受影响批次 & 构造条件与独立下游用途 \\
    证据不足 & 状态、候选集合与使用方式 & 暂缓、降权、隔离或追加证据 & 覆盖损失与再次触发条件 \\
    \bottomrule
  \end{tabularx}
  \endgroup
  \caption{核查结果与处理范围的对应。前五行处理已有证据支持的问题，末行处理仍无法裁决的情况；高损失或来源冲突本身不能确立问题类型。}
  \label{tab:knowledge-quality-assessment-correction-repair-scope}
\end{table*}

一个“\textbf{构造条件错、要动整条流水线}”的具体例子：某条\textbf{裁剪规则把花朵主体\emph{裁掉了}}（只剩背景），但流程却\textbf{仍把裁剪图和原图记为“\emph{正对应}}”（声称语义保留）。修正\textbf{不是改某张图的标签}，而是一套动作：\textbf{其一，\emph{限制裁剪区域}}（规则上保证主体在框内）到\textbf{其二，\emph{重新生成}“原图—裁剪图”的对应关系}（旧的正对应作废）到\textbf{其三，\emph{检查所有用旧关系}训练/评价过的\emph{派生版本}}（哪些下游模型/数据集建立在错关系上，要重算）。

这类\textbf{补标/重生成之后}还要追三个连带变化：\textbf{其一，对象有没有\emph{新增}（新框出新花）或\emph{退出}（原来误标的对象被移除）}；\textbf{其二，\emph{适用条件}有没有变}（新裁剪适用的范围和旧的不同）。一旦这些发生，\textbf{“\emph{抽样框}和\emph{审计范围}”都要\emph{更新}}（新对象进框、退出对象出框），并\textbf{\emph{单列}报告“新增覆盖部分”的类别\emph{组成}}（防止新旧混在一起后指标含义漂移）。\textbf{如果只是“\emph{对同一总体}\emph{重新计算}目标”}（对象、范围都没变，只是重新生成像素/关系），则\textbf{总体可以保留}，但\textbf{必须保存\emph{新的构造版本}}以便区分“同一总体的不同构造产物”。

\subsubsection{无法可靠修正时的处理}
\label{sec:knowledge-quality-assessment-correction-unresolved}

有些情形\textbf{不是“再努力一点就能改对”}，而是\textbf{本质上\emph{证据不足}}：比如一张\textbf{严重遮挡的照片}，可能\textbf{无论如何都\emph{无法区分}山茶与茶梅}。这时\textbf{把它\emph{强行}改成任一单标签，都\emph{缺少依据}}——硬改只会制造一个看似确定、实则猜测的新错误。正确做法不是“选一个”，而是一组\textbf{“诚实处理不确定性”的动作}：
\begin{itemize}\item \textbf{保留候选集合}（山茶/茶梅都留着，不假装唯一）；\item \textbf{标记“待复核”}（等更好照片/专家）；\item \textbf{\emph{降低}训练权重}（让这条不太可靠的样本少影响训练）；\item \textbf{\emph{暂时隔离}}（高风险时干脆先不用于训练）；\item \textbf{\emph{追加}证据}（补拍、补检测、调其他视图）。\end{itemize}
要明白这些动作\textbf{改变的是“\emph{表达形式}和\emph{使用方式}”，\emph{未必}纠正了内容}——它没让你知道真值，只是不让缺证据的样本危害系统。两个具体区分：其一，若\textbf{独立证据已\emph{确认}“真值在某\emph{集合}内”}（确定非此即彼），可按第~\ref{sec:knowledge-fusion-set-abstention-use}节\textbf{用集合约束训练}；但\textbf{仅由\emph{模型自己筛}出的候选}\textbf{仍可能漏掉真值}，\textbf{不能当“确定集合约束}”（第6章反复强调）。其二，\textbf{降权和隔离不一样}：\textbf{降权\emph{只减小}该样本对训练的贡献（保留覆盖、降低风险），隔离则\emph{同时移除}风险\emph{和}覆盖}（更安全但这条等于不用了）。

当\textbf{不是个别、而是\emph{一整批}样本都疑似有标注噪声}、又\textbf{无法逐条核出真值}时，还有第6章那条\textbf{不逐条改标签、而在\emph{训练损失}里校正}的路：若能为这批样本\textbf{建立可信的“报错通道”（噪声转移矩阵$\bm T$）}，可按第~\ref{sec:ka-noise-correction}节\textbf{保留原始观测标签、用前向/后向校正损失训练}。

但要\textbf{严格记得它\emph{做了什么、没做什么}}：它\textbf{改变的是“训练时\emph{对噪声的\emph{统计处理}}”（让模型不被噪声带偏）}，\textbf{\emph{没有}逐条裁定真值、也\emph{没有}把任何一条标签改对}——所以它是“\textbf{带着已知不确定性\emph{稳健训练}}”，不是“消除了未决”。它的两个失效条件也照搬到这：\textbf{通道$\bm T$\emph{不可识别}时}（没有锚点估不出，见第6章）、或\textbf{通道\emph{随照片特征}变化}（不是类别条件噪声）时，\textbf{矩阵校正都\emph{不能}消除这些未决}，还得回到补证据/隔离。

一个\textbf{用“删除低分样本”会\emph{系统性作恶}}的场景：\textbf{当某\emph{少数类别}\emph{持续}拿到低置信度时，\textbf{批量删除这些样本\emph{尤其危险}}。为什么？回忆第6章——\textbf{低分\emph{不一定}是标签错}，它\textbf{很可能来自“\emph{模型本身\emph{缺少}识别这个少数类的能力}}”}（模型不会，不是标签错）。这时\textbf{把低分样本\emph{删掉}，等于\emph{把模型\emph{本来就学不好}的那类样本}进一步抽走}——\textbf{训练数据在这个任务区域\emph{进一步}流失}，模型更学不会，下次给该类的分更低，形成“越删越不会、越不会越删”的恶性循环。

正确顺序：\textbf{先\emph{抽查独立证据}}（这些低分到底是错标，还是模型无能/样本罕见？），\textbf{再\emph{比较}几种处理的代价}：保留候选集合、允许弃权转专家、追加专业复核，哪个代价/风险更合适。\textbf{如果当前任务实在\emph{不能容忍}“未决”状态}（下游必须确定答案），可以\textbf{\emph{暂时隔离}}这些样本，但\textbf{必须\emph{如实报告覆盖率因此下降}}（你放弃了多少、哪一类），\textbf{\emph{绝不能}把“\emph{剩余}样本上的\emph{较高正确率}”\emph{解释成}“\emph{原总体}质量\emph{提高}了”}——那只是“把难题都不算之后”的数字美化（第6.8阈值陷阱同款）。

“待复核/未决”\textbf{不能变成一个“\emph{只进不出}的黑洞”}——很多系统的待复核队列越积越多、再没人看，等于把不确定性藏起来了。所以\textbf{每条未决标注都要有\emph{明确的“去向”和“重新触发”规则}}：
\begin{itemize}\item \textbf{保存原始内容和候选证据}（不能在搁置中丢失）；\item \textbf{事先规定\emph{什么条件下重新审计}}：相关\textbf{来源更新}了（新词典/新模型可能解决它）、\textbf{风险累积超阈值}（同类未决多到影响使用）、或\textbf{积累到一定\emph{数量}}（批量送专家更划算）；\item \textbf{\emph{定期}检查这些未决是否\emph{集中在某些特定群体}}（某地区/某来源/某难类长期悬而未决，说明那一块有系统性盲区）。\end{itemize}
连“\textbf{什么时候\emph{停止}追加证据、就此搁置}”也要\textbf{有\emph{事先明确}的理由}，比如“\textbf{追加证据对\emph{降低使用风险}的作用已\emph{低于门槛}}（再查也改变不了用途决策）”或“\emph{边际复核成本过高}、风险可接受”——\textbf{而不是让“待复核”状态\emph{永久}挂着、在报表里\emph{变成不可见}（既不算错也不算对、无人负责）}。

\subsection{修正效果的验证}
\label{sec:knowledge-quality-assessment-correction-repair-validation}

“修正完了，怎么\emph{证明}它确实更好？”——这里\textbf{藏着\emph{两个}必须分开回答的问题}，混为一谈就会出现“改的那批看着挺好、上线却没用甚至更糟”：
\begin{itemize}\item \textbf{修正\emph{结果}的独立复核}：问的是“\textbf{\emph{这一条}被改动的对象，\emph{现在}对不对}”。它针对的是\emph{具体改动点}（$x_2$改后真的是茶梅了吗），\textbf{可以复核\emph{同一对象}的前后}，但\textbf{下判断用的\emph{证据}或\emph{复核者}必须和“修正过程”\emph{适当隔离}}（不能让设计修正规则的人又来当裁判）。\item \textbf{修正\emph{策略}的\emph{留出}验证}：问的是“\textbf{我这套\emph{冻结}后的处理规则（新词典/新阈值/新流程），\emph{遇到\emph{没参与设计}的\emph{新对象}}时，\emph{还}有效吗}”。它针对的是\emph{规则的泛化}，\textbf{必须把“策略\emph{开发}用的样本”和“\emph{评价}用的样本”\emph{严格隔离}}。\end{itemize}
两者最容易踩的坑：\textbf{不能\emph{先看了留出结果}、发现不达标、再回去\emph{继续调}规则，却还\emph{沿用}原来那份“这是独立留出”的声明}——一旦你为了留出结果改了规则，留出集就变成开发集了（数据窥探，独立性作废），必须换新的留出集。

图~\ref{fig:knowledge-quality-assessment-correction-evidence-isolation}把这\textbf{两种“隔离”画成上下两条路径}，对照着看就不会混：
\begin{itemize}\item \textbf{上板（结果复核）}：\textbf{允许\emph{同一对象}形成“修正前 vs 修正后”的\emph{配对}}（当然要看同一张照片改对没），\textbf{这里被隔离的是“\emph{判断证据/裁判}}”——要由没参与修正、或独立的新证据来判改后对不对。\item \textbf{下板（策略验证）}：\textbf{隔离的是“\emph{策略开发}”和“\emph{推广评价}}”两个\emph{阶段}——用一批\emph{开发审计}定策略、定阈值，\textbf{在\emph{看任何留出结果之前}就把策略\emph{冻结}}，再用\emph{另外}的留出审计评价；\textbf{冻结线之后，留出结果\emph{不能再送回去}改策略}（否则越界）。\end{itemize}

\subsubsection{修正结果的独立验证}
\label{sec:knowledge-quality-assessment-correction-independent-validation}

先看“结果复核”具体怎么做。\textbf{要检查改动\emph{是不是真的纠了错}，必须让“修正前、修正后”\emph{可比}}——所以要\textbf{在\emph{同一任务定义}、\emph{同一参考依据}下，把\emph{同一对象}的前后结果\emph{配对}}（同一张照片、同一套规范，只看改前改后），不能拿“改后这张”去和“另一张照片”比。对\textbf{前后两版\emph{都}有监督信号、且\emph{都}能二值裁决}的同一对象$i$，记改前错误指示$e_i^{(0)}$、改后$e_i^{(1)}$（都是1=错、0=对），定义\term{改进量}（improvement）
$$d_i=e_i^{(0)}-e_i^{(1)}\in\{-1,0,1\}.$$
它有三种取值，但对应\textbf{四种必须\emph{分开}记录的“对错转移”}（这是关键，零值里藏着两种完全不同的情况）：
\begin{itemize}\item $d_i=+1$：\textbf{错到对}（真正改进）；\item $d_i=-1$：\textbf{对到错}（\emph{改坏了}，新引入的错误）；\item $d_i=0$包含两种：\textbf{对到对}（本来就对、没动坏）和\textbf{错到错}（\emph{改了但还是错}，漏修）。\end{itemize}
所以\textbf{四类转移（错到对、对到错、对到对、错到错）都要\emph{分别保存}计数}，不能只存一个$d_i$——否则“对到对”和“错到错”被混成同一个0，你就看不出“有多少改了白改”。而\textbf{那些“改前/改后有一版\emph{退出使用}、\emph{缺失}、或\emph{无法裁决}”的对象}，它们\textbf{\emph{不能}硬塞进这个$\{-1,0,1\}$二值配对}（你没有两个可判对错的版本可比）。对它们要\textbf{\emph{另列}报告}（多少条没法纳入配对、原因是什么），\textbf{不能为了凑数把“无法裁决”当“对”或“错”填进$d_i$}——否则配对率被系统性污染（前面样本级讲过的纪律）。

单条配对$d_i$只说明“\textbf{被抽中、被复核的这些对象\emph{实际}变好了多少}”；要\textbf{由它\emph{估计总体}下降多少}，还得加一个抽样前提：\textbf{“\emph{各对象的前后结果}\emph{不由‘评价入样}’决定}”。换句话说，\textbf{用哪些对象来做这个评估（抽$n_h$、分层），必须\emph{在看到改前改后答案之前}就定好}，不能“先看哪些改对了、再专挑那些来评估”。

具体地：\textbf{两个版本（或将要推广到总体的策略）应在\emph{评价抽样之前}就\emph{确定（冻结）}}；\textbf{策略允许“\emph{按对象特征}决定\emph{改不改}某类”}（比如只修高风险层，这是事先规则），\textbf{但\emph{不能}“看了\emph{本次参考答案}”才\emph{挑对象}或\emph{临时改规则}}”——前者是事先可复现的策略，后者是用评价数据做开发，估出的“总体改善”必然虚高。在满足上述前提（前后都可二值裁决、抽样事先固定）的同一总体上，若\textbf{各\emph{层}\emph{独立}简单随机无放回抽样}，令$\overline d_h$为第$h$层差值$d_i$的层内均值、$s_{d,h}^2$为层内差值的（无偏）样本方差（$n_h\geq2$），沿用7.2的总体层权重$W_h=N_h/N$和层抽样比例$f_h=n_h/N_h$。关键技巧是：\textbf{直接对“\emph{差值}$d_i$”做分层估计，而\emph{不是}“先算改前错误率、再算改后错误率、两个方差各自算一遍相减}”——因为\textbf{前后是\emph{同一对象}，两者\emph{强相关}}，对差值估计能自动利用这种相关、方差更小。公式为：
\begin{equation}
  \widehat\Delta_{\mathrm{st}}
  =
  \sum_{h=1}^{H}W_h\overline d_h,
  \qquad
  \widehat{\operatorname{Var}}(\widehat\Delta_{\mathrm{st}})
  =
  \sum_{h=1}^{H}W_h^2(1-f_h)\frac{s_{d,h}^2}{n_h},
  \label{eq:knowledge-quality-assessment-correction-paired-difference}
\end{equation}
读这个配对差值公式的两个要点：其一，\textbf{对“差值$d_i$”而不是“两次独立错误率”求方差}，正好\textbf{\emph{保留}了“\emph{同一对象}前后”的\emph{相关性}}（减掉了“对象本身难易不同”这部分无关波动，所以配对估计通常比两次独立估计更稳）；其二，方差里的\textbf{平方层权重$W_h^2$和有限总体修正$(1-f_h)$}都\textbf{来自\emph{同一个}抽样设计}，不能和别的设计混用。

一个最常用的特例：\textbf{简单随机、二值配对}时，设$b$条“错到对”、$c$条“对到错”、其余差值为0，则平均改进量直接是
$$\widehat\Delta=(b-c)/n.$$
即\textbf{改进量＝“改对的”减去“\emph{改坏的}”，再除以样本数}——\textbf{只报$b$（改对多少）会掩盖$c$（新引入多少错）}，必须净值。针对这种“只有两类\emph{不一致对}（$b,c$）提供信息”的配对设计，有专门的\term{McNemar 检验}（McNemar test）：它\textbf{用$b$和$c$检验“修正\emph{前后}两个边际错误率是否\emph{真的}不同”}（区分“改善”和“随机抖动”）。但要注意它的边界：\textbf{McNemar 只回答“\emph{有没有}差异”的\emph{显著性}，\emph{不替代}“效应\emph{多大}（$\widehat\Delta$）”和“\emph{不确定区间}”}，而且它针对的是简单配对设计，\textbf{不能直接套到分层/整群等复杂抽样上}\scite{knowledge-quality-assessment-correction-mcnemar1947}。

光“随机入样”\textbf{还\emph{不足以}保证配对估计无偏}——还有一个更隐蔽的陷阱：\textbf{“\emph{改哪些}”如果\emph{由入样决定}}，配对估计就被污染。一个极端但极有说服力的例子：假设1000条原标注\textbf{全部错误}，你随机抽100条、\textbf{\emph{只把这被抽中的100条}改正}，再请独立专家确认这100条\emph{全部}改对（$b=100,c=0$）。
\begin{itemize}\item 若你\textbf{错误地把“样本里$d_i=1$”直接套进配对公式}，会报“总体下降$\widehat\Delta=100\%$、方差为0”——仿佛整体全好了。\item \textbf{实际只改了100条，总体错误率只从$100\%$降到$90\%$，\emph{真实}下降只有$10$个百分点}。\end{itemize}
差在哪？因为\textbf{“\emph{修改对象}”是由“\emph{入样}”决定的——\textbf{评价样本（那100条）\emph{已经参与了修正方案的执行}}，它不再是“对一个\emph{预先冻结、推广到全体}的策略”的独立评价}（你只修了抽到的，没抽到的900条策略根本没作用过）。正确报告法：\textbf{老实报“这\emph{被修的100条}里，确认纠正了100条（确认纠正数）、它们占总体的\emph{已知份额}$100/1000=10\%$”}；\textbf{若要估计“\emph{修正后总体}会怎样”或“策略\emph{推广}效果”，必须\emph{另取}\emph{没参与修正}的独立评价样本（或用预先冻结策略+HT加权）}。

配对差值的估计\textbf{仍要\emph{匹配抽样设计}}，不同设计有不同注意点：
\begin{itemize}\item \textbf{按簇抽样}（先抽文档/批次）：\textbf{配对差值要在\emph{簇层}估计、或用尊重簇的重抽样}，不能把同文档样本当独立（同7.2）。\item \textbf{只评价“高风险\emph{范围}}”（开发/排查只在$x_2$这类冲突区做修正）：结论\textbf{\emph{严格限于该范围}}——你只能说“高风险区改善了多少”，不能外推到全体。\item \textbf{策略\emph{事先}规定“只修高风险对象}”（冻结的规则），但\textbf{评价样本\emph{覆盖整个总体}}（在全体上随机抽，含大量策略根本不会动的低风险对象）：这时反而\textbf{可以\emph{估计全体}变化}——因为未修对象的$d_i=0$，它们“零变化”被\textbf{如实计入}，于是得到的是“这个只修高风险的策略\emph{推广到全体}的净效果”。\end{itemize}
最后是“规范改变”的配对口径：\textbf{只有\emph{新旧两个定义\emph{共同覆盖}}（含义没变、只是对错可能变）的对象，才能做“\emph{同一口径}”的前后配对}；那些“旧类被拆/并、在新规范里没有对应物”的对象，\textbf{无法同口径配对，应\emph{分版本}报告}（7.1的尺子问题）。

\begin{figure*}[htbp]
  \centering
  % 图7-3：结果独立复核与策略留出验证使用不同的隔离要求。
\begin{tikzpicture}[
  x=1mm,
  y=1mm,
  font=\figurefont\footnotesize,
  text=BookInk,
  record/.style={draw=FigureInput,fill=FigureInput!8,line width=0.8pt,text width=21mm,minimum height=13mm,align=center,inner sep=1.4mm},
  action/.style={draw=FigureModel,fill=FigureModel!8,line width=0.9pt,rounded corners=1.2mm,text width=23mm,minimum height=13mm,align=center,inner sep=1.4mm},
  evidence/.style={draw=BookAmber,fill=BookAmber!8,line width=0.8pt,text width=25mm,minimum height=12mm,align=center,inner sep=1.4mm},
  result/.style={draw=FigureOutput,fill=FigureOutput!6,line width=0.9pt,text width=22mm,minimum height=13mm,align=center,inner sep=1.4mm},
  flow/.style={knowledge arrow,draw=BookTeal,line width=1.0pt,rounded corners=1mm},
  evidence flow/.style={knowledge arrow,draw=BookAmber,line width=0.9pt,rounded corners=1mm},
  boundary/.style={draw=BookMuted,line width=0.8pt,dashed},
  note/.style={font=\figurefont\scriptsize,text=BookMuted,align=center}
]
  \draw[BookRule,line width=0.7pt] (0,37) -- (166,37);
  \node[anchor=north west,font=\figurefont\bfseries] at (1,93) {(a) 已改结果的独立复核：对象可以相同，判断证据需要独立};
  \node[record] (before) at (14,58) {开发对象$x_2$\\修正前标注};
  \node[action] (repair) at (48,58) {依据开发证据\\实施修正};
  \node[record,draw=FigureModel] (after) at (82,58) {开发对象$x_2$\\修正后标注};
  \node[action,draw=BookAmber,fill=BookAmber!8] (review) at (122,58) {未参与修改者\\独立复核};
  \node[result] (verdict) at (153,58) {确认改对／改错\\或无法裁决};
  \node[evidence] (new-evidence) at (122,79) {新证据或盲化材料};

  \draw[flow] (before) -- (repair);
  \draw[flow] (repair) -- (after);
  \draw[flow] (after) -- (review);
  \draw[flow] (review) -- (verdict);
  \draw[evidence flow] (new-evidence) -- (review);

  \node[anchor=north west,font=\figurefont\bfseries] at (1,32) {(b) 总体或策略效果评价：在评价抽样前冻结策略};
  \node[record] (development) at (15,9) {开发审计样本\\已查看问题};
  \node[action] (design) at (50,9) {设计规则、模型\\或接收阈值};
  \node[action,draw=BookInk,fill=BookPaper] (freeze) at (85,9) {冻结修正策略\\和评价口径};
  \node[evidence] (holdout) at (123,9) {独立抽样审计\\比较前后版本};
  \node[result] (effect) at (153,9) {策略效果\\及适用范围};

  \draw[flow] (development) -- (design);
  \draw[flow] (design) -- (freeze);
  \draw[flow] (freeze) -- (holdout);
  \draw[flow] (holdout) -- (effect);
  \draw[boundary] (104,-4) -- (104,22);
  \node[note,anchor=south] at (104,22) {评价边界};
  \node[note,anchor=north,text width=162mm] at (83,-6) {查看评价结果后若继续调整策略，需另取独立证据；开发对象不进入本例的验收总体。};
\end{tikzpicture}

  \caption{修正与验证的证据隔离。上面板复核同一被修改对象，独立性来自未参与修改的复核者或新证据；下面板用开发审计确定策略，在评价抽样前冻结，再用留出审计评价。上面板只确认已改结果；用于总体推断的配对样本还须满足下面板的隔离条件。}
  \label{fig:knowledge-quality-assessment-correction-evidence-isolation}
\end{figure*}

独立验证时\textbf{看什么、不看什么}有几条纪律：
\begin{itemize}\item \textbf{先对可二值裁决对象\emph{统计齐四类对错转移}}（错到对/对到错/对到对/错到错），\textbf{未裁决、退出的\emph{另列}}。\item \textbf{还要\emph{抽查“未修改区域}}”——这是\textbf{查“规则\emph{漏修}”}的\emph{唯一}办法。只看被改对象，你能算清“改对/改错比例”，却\textbf{\emph{看不见}“该改却没被规则触及的那批”}（规则召回率问题）。\item 反过来说，\textbf{只看“总体平均改善”也不够}——一个$4\%$的平均改善，可能\textbf{\emph{掩盖}某\emph{少数类别}\emph{恶化}}（常见类大改善、稀有类被改坏，平均仍好看），所以要\emph{分组}看（呼应6.6/7.2）。\item \textbf{不同监督信号用不同的复核法}：\textbf{关系重建}须\textbf{核查“对象绑定+成立条件”}（改完的关系两端、方向对不对）；\textbf{表示目标更新}须\textbf{\emph{分开}检查“构造条件”和“下游用途}”（构造合规不代表下游变好，见7.1、第6章）。\end{itemize}
最后，验证不只是“打勾通过”，还要\textbf{回到决策}：\textbf{是否回退版本、追加审计、或限制策略范围}，应\textbf{依据\emph{事先预定的分组要求}和\emph{证据充分程度}决定}（比如某层改善证据不足就限制该层或回退），而不是看总体数字碰巧好不好看。

\subsubsection{修正收益与成本的比较}
\label{sec:knowledge-quality-assessment-correction-benefit-cost}

就算\textbf{确实“把一批标注\emph{改对了}”}（结果复核通过），事情还\textbf{没结束}——还要问一句\textbf{“这次改进\emph{值不值得}投入}”，并\textbf{检查“\emph{改错}（对到错的新错）”和“\emph{覆盖}损失（隔离/收紧范围丢了多少）”}。所以比较的是一个\textbf{多方面的账单}，至少含四项：\textbf{其一，净纠错（改对$-$改坏，不是只看改对）、其二，关键分组（各类/各区域分别有没有变好变坏）、其三，下游用途（拿去训练到底有没有用）、其四，本轮\emph{完整}成本}。

其中“下游贡献”最容易被冒领。\textbf{若要测“\emph{标签修正}对下游的贡献}”，必须做\textbf{控制变量}：\textbf{固定\emph{学习算法}、固定\emph{样本范围}、固定\emph{评价样本集}}，让\textbf{两次训练之间\emph{唯一}的主要变化就是“标签换了”}，并\textbf{另计训练随机性}（用同种子或多种子平均）。如果你\textbf{一边换标签、一边\emph{还换了模型、加了算力预算、加了训练数据}}，那\textbf{下游变好\emph{无法归因}给“标签修正”}（差异可能来自其他任何一项）。最后守住一句克制：\textbf{本节先只比较“\emph{标注/训练目标本身}的质量”变化，\emph{不预设}它已经带来模型收益}——标签更对不等于模型更好，下游要另测。

下面把前面$x_2$/K17 的开发诊断，\textbf{接成一次\emph{完整、可复算}的具体验收流程}（这是本章教学案例的收官）。先回顾背景：\textbf{开发复核}在$x_2$等对象上确认“\emph{共享旧词典失效}”，于是\textbf{新版本K18采取一种\emph{隔离}措施：\emph{暂时}把那两条同模板规则的权重置零}（不让它们投票），\textbf{其他来源仍按等权合并}；\textbf{无可用答案或平票时仍保留未决}。要准确理解这个动作的性质：\textbf{它只是“\emph{先把\emph{失效来源}从融合中\emph{摘除/隔离}}”的处置，\emph{还\emph{没有}“修好词典本身}}”（词典里那条错映射还原样在，只是暂时不采信依赖它的票）。具体到$x_2$：两条规则票被摘除后，\textbf{模型那张“茶梅”成为\emph{剩下唯一}的有效票}，$x_2$从而改判茶梅。\textbf{K17、K18\emph{两版都保留原始来源类别票}}（不删历史），\textbf{以便以后\emph{恢复}和\emph{比较}}（能看到规则被摘除前后各投了什么）。

为了让验收\textbf{不被数据窥探污染}，\textbf{所有“策略、分层方式、阈值”都在“\emph{打开那1000条独立验收总体的参考答案}”\emph{之前}就\emph{先固定}好}，并\textbf{在“\emph{评价抽样}”\emph{之前}冻结}（先定规则、再看结果）。还有一个\textbf{范围纪律}：\textbf{K18要\emph{按同一规则\emph{重算\emph{全部}1000条}}，而\emph{不是}只重算“抽中的那120条}”——因为策略变了（摘除两规则），\textbf{没被抽中的880条的结果\emph{也}可能变}，只重算抽中的会让“版本对比”漏掉未抽查部分的变化（呼应iv7）。这批教学样本的\textbf{可比性前提}要交代清楚：\textbf{两版\emph{都}给出1000条“可接收”的类别答案，\emph{没有}新增、退出或未决对象}——即两版\textbf{评价的是\emph{同一批}1000个对象、同样都有答案可比}（没有“这版多了/少了对象”造成的口径漂移）。但措辞要\textbf{精确}：\textbf{“可接收”只表示“\emph{满足作答协议与阈值}（有有效答案、够接收线）”，\emph{\emph{不}表示“已\emph{确认正确}}”}——“能给出答案”和“答案对”是两回事（后者正是接下来独立复核要查的）。然后请\textbf{\emph{独立}复核者（与K18的开发隔离）复核\emph{同一组}$80$常见+$40$易混淆的样本}，记录每一条的“K17对错 vs K18对错”，得到表~\ref{tab:knowledge-quality-assessment-correction-benefit-cost}里的\textbf{四类配对转移计数}（注意是同一批样本的前后对照，不是两批不同样本）。

常见层有3条改对、1条改错，易混淆层有6条改对、2条改错，因此
\[
  \overline d_1=\frac{3-1}{80}=0.025,\qquad
  \overline d_2=\frac{6-2}{40}=0.10,
\]
\[
  \widehat\Delta_{\mathrm{st}}=0.8(0.025)+0.2(0.10)=0.04.
\]
\textbf{先把表读成“\emph{纠正}了多少”和“\emph{改坏}了多少”两笔账}，且都按总体权重汇总（不能按抽样数$80/40$加权，理由同7.2）：
\begin{itemize}\item \textbf{纠正率（错到对的净\emph{正向}贡献）}：$0.8(3/80)+0.2(6/40)=6\%$；\item \textbf{新增错误率（对到错的\emph{负向}贡献）}：$0.8(1/80)+0.2(2/40)=2\%$。\end{itemize}
\textbf{两者\emph{相减}$6\%-2\%=4$个百分点，才是“\emph{净}改善”}——只看“改对$6\%$”会漏掉“同时还新引入了$2\%$的错”。再直接核对K18改完后的错误率：常见层剩$2$条错（原4条里改对3、又把1条对的改错）、易混淆层剩$4$条（原8条改对6、改坏2），故K18总体错误率点估计$0.8(2/80)+0.2(4/40)=4\%$，\textbf{相对K17的$8\%$，\emph{净}下降$4$个百分点}。措辞要准：是“\textbf{从$8\%$\emph{降到}$4\%$，\emph{相差}4个百分点}”，\textbf{\emph{不是}“\emph{下降了}$4\%$”}（后者会被理解成相对降幅，量纲不同）。

更重要的是\textbf{不要只看净值，要看四类原始转移}：样本里还有\textbf{$3$条“错到错”（\emph{始终错误}}）——它们\textbf{暴露“\emph{漏修}}”（K18规则没碰到它们）；另有\textbf{$3$条“对到错”（\emph{被改错}}）——它们\textbf{暴露“\emph{新引入风险}}”（新规则误伤）。\textbf{$4\%$这个净值是$+6$和$-2$抵消后的结果，\emph{不能替代}这两类原始计数}；只报净值，你既不知道规则漏了多少、也不知道误伤了多少。

要给“净改善$4\%$”配不确定性，注意\textbf{不能把“K17错误率方差”和“K18错误率方差”\emph{简单相加}}——因为两次测量来自\emph{同一批对象}，\textbf{前后结果\emph{强相关}}，必须用\textbf{对“\emph{差值}$d_i$}”直接估方差（就是式~\eqref{eq:knowledge-quality-assessment-correction-paired-difference}的作用）。由$d_i^2$只在改变对错的对象上等于1，得到
\[
  s_{d,1}^2=\frac{3+1-80(0.025)^2}{79}=0.05,\qquad
  s_{d,2}^2=\frac{6+2-40(0.10)^2}{39}\approx0.19487.
\]
代入式~\eqref{eq:knowledge-quality-assessment-correction-paired-difference}得到标准误约0.02271。正态近似95\%区间为
\[
  0.04\pm1.96(0.02271)
  \approx[-0.00452,0.08452].
\]
用百分点表示，这次净改善量约落在$\boldsymbol{[-0.45,\ 8.45]}$个百分点之间。怎么\textbf{诚实地}读这个区间？
\begin{itemize}\item \textbf{区间\emph{跨过了$0$}（下界$-0.45<0$）}：意味着“\textbf{这次修正\emph{可能}\emph{完全没改善、甚至轻微变差}}”在统计上\emph{不能排除}——所以\textbf{\emph{没有}证据支持“改善\emph{确定}为正}”。\item 区间这么宽，\textbf{根因是“\emph{真正改变对错的配对}只有$b+c=3+1+6+2=12$条}”（其余都没变、不含信息），\textbf{有效样本小，\emph{近似}误差不可忽略}。\end{itemize}
所以这次演算\textbf{\emph{不足以}提供“严格的$95\%$发布保障}”——点估计好看（$+4\%$），但不确定性太大。

\begin{table*}[htbp]
  \centering
  \small
  \begingroup
  \setlength{\tabcolsep}{4pt}
  \begin{tabular*}{\linewidth}{@{\extracolsep{\fill}}lrrrrrr@{}}
    \toprule
    \textbf{分层} & \textbf{抽查数} & \textbf{错$\to$对} & \textbf{对$\to$错} & \textbf{对$\to$对} & \textbf{错$\to$错} & \textbf{退出} \\
    \midrule
    常见类别 & 80 & 3 & 1 & 75 & 1 & 0 \\
    易混淆区域 & 40 & 6 & 2 & 30 & 2 & 0 \\
    \bottomrule
  \end{tabular*}

  \medskip

  \begin{tabularx}{\linewidth}{@{}>{\RaggedRight\setlength{\parindent}{0pt}\arraybackslash}p{29mm}>{\RaggedRight\setlength{\parindent}{0pt}\arraybackslash}X>{\RaggedRight\setlength{\parindent}{0pt}\arraybackslash}p{56mm}@{}}
    \toprule
    \textbf{验收项} & \textbf{预先给定的要求} & \textbf{本次结果} \\
    \midrule
    错误改善 & 点估计至少2个百分点，且区间下界大于0 & 点估计4个百分点；近似下界为负，证据不足 \\
    分组表现 & 两层错误率点估计均不升高 & 5\%降至2.5\%；20\%降至10\% \\
    接收覆盖 & 总体及每层均不低于99\% & 两版均为100\%；退出为0 \\
    本轮投入 & 人工不超过32工时，计算不超过1机时 & 开发4、复核24、维护2工时；重算0.5机时 \\
    \bottomrule
  \end{tabularx}
  \endgroup
  \caption{K17到K18的教学验收。上表是独立样本中的原始配对数，汇总需使用总体层权重；下表连接预定要求与结果。全部审计对象均可裁决，且没有退出，故四类对错转移覆盖全部样本。点估计改善而不确定性仍大，本次暂缓发布。}
  \label{tab:knowledge-quality-assessment-correction-benefit-cost}
\end{table*}

\textbf{标签/目标质量改善，\emph{不等于}下游模型/任务就改善——必须等“监督信号\emph{真正进入任务}（拿去训练/使用）\emph{以后}，才有\emph{相应的下游收益}可言}。而且\textbf{不同监督信号要\emph{用不同的下游评测}，不能用一个分类准确率包打}：
\begin{itemize}\item \textbf{类别修正}：\textbf{固定学习算法}，比较“用K17标签训 vs 用K18标签训”的评价结果（才是标签修正的贡献）。\item \textbf{关系重建}：比较\textbf{关系任务本身、或\emph{使用这些关系}的下游模型}（如配对/检索任务）。\item \textbf{表示目标}：在\textbf{约定好的使用方式下}，比较\textbf{独立的下游任务}（不能用构造损失代替）。\end{itemize}
这正呼应7.1那个根本分工：\textbf{“构造检查}”回答“\emph{条件/约定有没有\emph{落实}}（流程对不对）”，\textbf{“用途评价}”回答“\emph{任务有没有\emph{受益}}（实际好不好）”——两者\textbf{不能互换}（流程合规并不等于任务变好）。成本也要\textbf{算\emph{全周期}}：\textbf{人工复核、自动调用、重生成、再训练、再次审计、后期维护}都要计入。在“预算口径、效果口径”都明确后，才可以\textbf{在“继续修正、追加监督信号、\emph{保留现状}}”之间理性比较——\textbf{\emph{不必}把“继续修改”当作唯一或默认选项}，证据不足时暂停/维持现状完全可以是最优决策。

把表~\ref{tab:knowledge-quality-assessment-correction-benefit-cost}当“\textbf{\emph{预先}定好的发布门槛 vs 实际结果}”的验收单读，结论很清楚：
\begin{itemize}\item \textbf{覆盖、投入}这两项\textbf{达标}（接收覆盖$100\%$、退出0；总投入30工时在预算内）；\item 但\textbf{“错误改善”这一项\emph{没}达标}——要求“点估计至少+2个百分点\textbf{\emph{且}区间下界$>0$}”，实际点估计$+4\%$过了第一关、但\textbf{区间下界$-0.45<0$，第二关没过}。\end{itemize}
所以基于事先规则，\textbf{本次决定是“\emph{暂缓发布}K18”}——\textbf{不是因为它变差了，而是因为“\emph{改善\emph{是否真实存在}}”的证据\emph{还不够硬}}（不能排除零改善）。投入要算清且\textbf{不强行折成同一个金额}：复核$24$工时（120条、每条12分钟）+开发$4$工时+维护$2$工时$=30$工时；计算资源另算，\textbf{“人工小时”和“机时”量纲不同，没有业务给出的换算系数就不相加（第6章预算原则）}。还有一个诚实的记账：\textbf{这一轮\emph{还没有}重训下游模型}，所以\textbf{重训成本、用途收益\emph{都还没发生}}——\textbf{\emph{不能}把“标签错误率从$8\%$降到$4\%$”当成“\emph{已经取得的模型收益}}”（模型压根没动）。

那\textbf{若要继续推进，正确姿势是什么}？两条，且都\textbf{\emph{预先}设计}：其一，\textbf{预先设计一个\emph{新的、更大的审计批次}}（或改进合并/加权规则），用更足样本把区间压窄到下界$>0$；其二，若要改策略，\textbf{重新开发后\emph{另作一次独立验收}}（旧开发样本不能再当裁判）。\textbf{绝不能“\emph{不断加看样本、直到区间\emph{碰巧}过线}”——那是数据窥探，过线也是假的}（7.2“反复看区间”陷阱）。

\subsection{知识版本与持续质量的维护}
\label{sec:knowledge-quality-assessment-correction-version-maintenance}

先理解“\textbf{一个结论\emph{绑定}在什么上}”：\textbf{“K18\emph{暂缓发布}”这个结论，是\emph{这次}“固定配置 + 那次审计”的产物——它\textbf{\emph{不能自动转移}到“\emph{后来}修好词典之后的新版本}”}。修好词典后要重新走流程、重新评估，旧的“暂缓”不代表新的也不能发，旧的“通过”也不自动顺延。

为了让每个版本\textbf{事后\emph{可复查、可复现结论}}，一个\term{可复查版本}（reviewable version）不能只存“结果标签”，而要\textbf{\emph{绑定}并保存以下\emph{全套}上下文}：
\begin{itemize}\item \textbf{样本范围}（这版覆盖哪些对象）；\item \textbf{样本标注、或目标构造配置}（规则/教师/变换版本）；\item \textbf{来源依赖}（用了哪些模型、规则、词典及其版本）；\item \textbf{修改理由}（为什么改、根因是什么）；\item \textbf{审计设计}（怎么抽样、查什么）、\textbf{参考依据与\emph{验证结果}}。\end{itemize}
两个“\textbf{光靠名字不够}”的提醒：其一，\textbf{递增的版本号（K17/K18）\emph{只能}区分“这是两个不同名字”，\emph{不能}帮你\emph{恢复}“当时是\emph{怎么生成}的”}——恢复生成过程靠的是上面那份绑定清单，不是编号；其二，\textbf{\emph{样本说明文档}（datasheet 等）也要随\emph{每次维护}更新}，因为\textbf{结论是在“当时那些条件”下成立的，条件变了、说明不更新，旧结论就失去了成立依据}\scite{gebru2021datasheets}。

不同类型数据要\textbf{保存不同的“变更内容”}：
\begin{itemize}\item \textbf{静态标签、关系}：要能\textbf{区分并保存“\emph{新增、删除、替换、未决}}”四种\emph{状态差异}（不只是存最新值——要知道相对上版多了什么、少了什么、改了什么、哪些还悬着）。\item \textbf{\emph{动态生成}的表示目标}（蒸馏/自监督/模型生成的向量或目标）：这类\textbf{没有手写来源、且依赖随机过程}，除了配置，还要\textbf{额外保存“\emph{用的模型}、\emph{预测输入的变换}、\emph{随机配置（种子）}，以及“\emph{能否\emph{重现}数值}”}——没有种子/版本，同一个版本号也复现不出同样的向量（第6章Progress\&Compress也强调过这点）。\end{itemize}为了给“\textbf{哪个版本、由\emph{什么过程}、用\emph{什么来源}产生}”这种溯源关系一个\textbf{通用、可交换的描述方式}，可以借鉴\term{W3C PROV}溯源模型：它把世界抽象成几类节点——\textbf{\emph{实体}（entity：产物，如一个标注版本/数据集）、\emph{活动}（activity：生成它的过程，如审计/训练）、\emph{责任主体}（agent：人/系统/模型）、以及\emph{派生/关联关系}（谁用了谁、谁生成谁）}\scite{knowledge-quality-assessment-correction-w3c2013provdm}。\textbf{本章\emph{不要求}你采用某一特定存储系统或软件}（数据库、文档、专门的溯源系统都行），但有一条硬要求：\textbf{这些“版本—过程—来源—责任人”的关联\emph{不能}只存在于“\emph{临时日志/某个人脑子里}}”——临时日志会被清、人会忘，溯源必须是持久、可查的记录。

图~\ref{fig:knowledge-quality-assessment-correction-version-feedback}用\textbf{两种箭头}展示“\emph{共享来源一旦更新}，影响怎么\emph{追踪}和\emph{闭环}”：
\begin{itemize}\item \textbf{实线（生成/依赖）}：表示“\emph{原始样本、来源版本、构造配置}\emph{产出}了标注版本、验证报告”的\textbf{强依赖}——它让\textbf{每份验证报告都能\emph{绑定回}它所评价的那个\emph{具体标注版本}}（不会出现“报告说的是哪一版”说不清）。\item \textbf{虚线（反馈/触发）}：表示“来源更新后，\textbf{\emph{触发}一次\emph{复查}、\emph{重生成}、或\emph{回退}}”的动作——它是流程信号，不是产物依赖。\end{itemize}
还要注意\textbf{版本共存}：\textbf{新版本在\emph{通过验证之前}，\emph{旧版本仍保留为“可恢复基线}}”——验证没过时可以回退，且新旧结果可对比；\textbf{不要}在新版本还没验证时就把旧版覆盖删除（否则验证失败无路可退、也无从比较）。

\begin{figure*}[htbp]
  \centering
  % 图7-4：来源、构造、样本标注和验证报告之间的可追溯版本依赖。
\begin{tikzpicture}[
  x=1mm,
  y=1mm,
  font=\figurefont\footnotesize,
  text=BookInk,
  data/.style={draw=FigureInput,fill=FigureInput!8,line width=0.8pt,text width=20mm,minimum height=12mm,align=center,inner sep=1.4mm},
  source/.style={draw=FigureModel,fill=FigureModel!8,line width=0.9pt,text width=20mm,minimum height=12mm,align=center,inner sep=1.4mm},
  config/.style={draw=BookMuted,fill=BookPaper,line width=0.8pt,text width=20mm,minimum height=12mm,align=center,inner sep=1.4mm},
  knowledge/.style={draw=BookTeal,fill=BookTeal!7,line width=1.0pt,text width=23mm,minimum height=12mm,align=center,inner sep=1.4mm},
  validation/.style={draw=BookAmber,fill=BookAmber!8,line width=0.9pt,text width=24mm,minimum height=12mm,align=center,inner sep=1.4mm},
  release/.style={draw=FigureOutput,fill=FigureOutput!6,line width=0.9pt,text width=19mm,minimum height=12mm,align=center,inner sep=1.4mm},
  dependency/.style={knowledge arrow,draw=BookMuted,line width=0.9pt,rounded corners=1mm},
  action/.style={knowledge arrow,draw=BookAmber,line width=0.9pt,dashed,rounded corners=1mm},
  publish/.style={knowledge arrow,draw=BookTeal,line width=1.0pt},
  edge label/.style={fill=white,inner sep=0.8pt,font=\figurefont\scriptsize,text=BookMuted,align=center}
]
  \node[data] (raw) at (10,20) {原始观测\\版本D5};
  \node[source] (source-old) at (40,42) {原来源配置\\版本M3};
  \node[source,draw=BookAmber] (source-new) at (40,-1) {隔离失效规则\\配置M4};
  \node[config] (config) at (75,20) {对象与预测请求\\配置C2};

  \node[knowledge,draw=BookAmber] (knowledge-old) at (75,42) {标注版本K17\\保留旧路径};
  \node[knowledge,draw=BookAmber] (knowledge-new) at (75,-1) {标注版本K18\\重生成候选};

  \node[validation] (validation-old) at (111,42) {验证报告R17\\绑定K17};
  \node[validation] (validation-new) at (111,-1) {新审计报告R18\\绑定K18};
  \node[release] (release-old) at (148,42) {已发布K17\\可回退};
  \node[release] (release-new) at (148,-1) {候选K18\\本例暂缓发布};

  \draw[dependency] (raw.east) -- ++(4,0) |- (source-old.west);
  \draw[dependency] (raw.east) -- ++(4,0) |- (source-new.west);
  \draw[dependency] (source-old) -- (knowledge-old);
  \draw[dependency] (config.north) -- (knowledge-old.south);
  \draw[dependency] (source-new) -- (knowledge-new);
  \draw[dependency] (config.south) -- (knowledge-new.north);
  \draw[dependency] (knowledge-old) -- (validation-old);
  \draw[dependency] (knowledge-new) -- (validation-new);
  \draw[publish] (validation-old) -- (release-old);
  \draw[publish] (validation-new) -- (release-new);

  \draw[action] (source-old.east) -- ++(5,0) |- (source-new.east);
  \draw[action] (knowledge-old.east) -- ++(5,0) |- (knowledge-new.east);
  \draw[action] (release-new.east) -- ++(3,0) |- (release-old.east);
  \node[edge label,text=BookAmber,anchor=east] at (54,20.5) {来源更新};
  \node[edge label,text=BookAmber,anchor=west] at (95,20.5) {重生成};
  \node[edge label,text=BookAmber,anchor=east] at (160,20.5) {验证失败回退};
  \node[font=\figurefont\scriptsize,text=BookAmber,anchor=south,align=center] at (75,51) {来源配置改变后，保留K17，并重算受影响标注形成K18};

  \node[anchor=north,font=\figurefont\scriptsize,text=BookMuted,align=center] at (82,-13) {实线：执行输入、生成、绑定或决策依赖；虚线：反馈动作。};
\end{tikzpicture}

  \caption{标注版本、来源依赖与质量反馈。实线表示原始样本、来源版本和构造配置对标注版本及验证报告的生成依赖；虚线表示来源更新后触发复查、重生成或回退。旧版本与新版本均保留，验证报告绑定具体标注版本；复查反馈用于形成新证据，不能把同一审计结果无限重复当作独立验证。}
  \label{fig:knowledge-quality-assessment-correction-version-feedback}
\end{figure*}

当\textbf{来源（模型/规则/词典）或任务规范\emph{更新}}时，维护流程有一个\textbf{固定的“受影响—对照—决策}”顺序，和7.3的根因定位一脉相承：
\begin{enumerate}\item \textbf{\emph{沿调用依赖清单}确定“\emph{哪些对象}受影响}”（谁引用了这个来源/规则）；\item \textbf{对\emph{冻结}样本\emph{分别}运行新旧配置}（只切来源、其他不变，做受控对照）；\item \textbf{执行\emph{局部}概率审计}（在受影响范围抽样估计变化，而非全量）；\item \textbf{比较版本差异}（配对待错转移、覆盖变化）；\item \textbf{决策：\emph{发布}、\emph{限制范围}、还是\emph{回退}}。\end{enumerate}
触发方式有两种、用途不同：\textbf{定期抽查}是\textbf{用来\emph{发现缓慢退化}}（词典悄悄过期、模型慢慢漂移，没人通知你）；\textbf{事件触发检查}是\textbf{用来\emph{响应}“来源/定义\emph{发生变化}}”（已知换了，主动查影响）。最后重复那条监测纪律：\textbf{连续监测\emph{仍然只产生线索}}（报警不等于错误）；而\textbf{一旦“\emph{监测用的样本}”被拿去\emph{调阈值、选方案}，它们就参与了方案选择，\textbf{要发布结论时必须\emph{另备新的独立证据}}验证}（监测样本不能再当裁判，全章一致）。

用$x_2$这条主线把7.4收口：\textbf{章首的\emph{开发排查}}确认了“$x_2$是\emph{共享来源（旧词典）错误}”，于是K18采取“\emph{隔离}那两条失效规则、用\emph{固定配置}\emph{重算总体}”的修正——它对症地处理了根因，但还没修词典本身。随后\textbf{独立验收得到“净改善4个百分点”的\emph{点估计}}，但因为\emph{区间跨0、证据不足}而\textbf{\emph{暂缓发布}}。

这里\textbf{版本清单该\emph{保留}什么}很能体现成熟度：要\textbf{\emph{完整}保留“这个暂缓决定”以及“\emph{尚未解决}的错误（漏修的3条、误伤的3条、词典本体未修）”}，\textbf{\emph{不能}只把“\emph{更好看}的点估计$+4\%$”存下来}——只存好数字，后人会误以为K18已验证有效。同时要防止反向误用：\textbf{“暂缓K18”\emph{不等于}“重新证明了K17\emph{可靠}”}——K17\textbf{已确认的旧问题（8\%错误率、共享词典失效）照样存在}，仍要\textbf{保留对它的\emph{风险处置}和\emph{复查安排}}，不能因为新版暂缓就让旧版继续带病运行而无人跟进。

维护不是线性收尾，而是一个\textbf{按“\emph{已确认根因}}”\emph{把问题\emph{送回正确环节}}的\emph{反馈环}（呼应7.3结尾）：
\begin{itemize}\item 若决定\textbf{彻底\emph{修好词典}}（而不只是隔离）到 产生\textbf{新的来源配置、新的待验收版本}，走新版本完整流程；\item 若发现是\textbf{照片与植株\emph{绑定错误}}（对象问题）到 \textbf{返回“\emph{标注对齐}}”环节（第6章6.3）；\item 若问题出在\textbf{模型\emph{预测输出}、或\emph{裁剪/构造条件}}到 \textbf{返回“\emph{相应的获取环节}}”（第4/5章），从源头改生成条件。\end{itemize}
\textbf{反馈\emph{退回}哪一步，严格由“已确认的原因”决定，而不是哪步方便改哪步}。两条贯穿全章的数据纪律在此收口：其一，\textbf{新的开发过程中“\emph{已经用过}的证据”要\emph{继续标明用途}（开发/诊断用）}，不能在后续冒充独立证据；其二，\textbf{后续的\emph{发布判断}\emph{另需独立评价依据}}（新留出集）——开发证据和发布证据始终分开，质量闭环才能一圈圈转下去而不自我证明。

最后一个\textbf{版本管理最容易制造“假进步”的地方——\emph{覆盖/总体组成}变化}。下面三种情况都会让“平均错误率”改变，\textbf{但改变\emph{未必}来自质量真的变好，而可能只是“被统计的对象\emph{换了一拨}}”：
\begin{itemize}\item \textbf{修正时\emph{删除}了难例}（把总判不对的样本剔出数据集）；\item \textbf{规范\emph{缩小}了适用范围}（把难覆盖的区域划到范围外/不适用）；\item \textbf{新来源\emph{补入}了一批少数类}（总体组成变了）。\end{itemize}
这些情况下，平均错误率\textbf{可能因为“\emph{总体组成改变}”}而下降，与“同样对象标得更准”毫无关系（第6章阈值、7.2范围口径都警告过）。所以\textbf{发布说明必须\emph{同时}报告四样东西}：\textbf{其一，\emph{共同对象}上的\emph{配对}变化（同一批对象前后）、其二，\emph{新增}对象、其三，\emph{退出}对象、其四，\emph{目标总体定义}有没有变、以及其五，\emph{仍未解决}的问题}。\textbf{把“共同对象的净变化”和“对象进出/组成变化”分开报告}，质量反馈才能\textbf{准确地\emph{返回}到“监督信号获取、来源选择、语义对齐、融合”中真正出问题的那一环}，而\textbf{\emph{不会}把“每一次版本\emph{组成}变化”都\emph{误写成}“质量\emph{进步}”}。

\section{本章小结}
\label{sec:knowledge-quality-assessment-correction-summary}

把全章收回章首那一万条标签。\textbf{第一组要分清的是“三种数字/三个问题”}：
\begin{itemize}\item \textbf{重点排查在入选的$100$条里确认了$40$条错}——它\textbf{只回答“\emph{入选区域}里查到了什么”}（检测器命中率），\textbf{回答不了整体质量}。\item \textbf{要估计整体质量}，必须\textbf{先明确“总体是谁、参考标准是什么、按什么\emph{概率抽样设计}选的样”}，再用\textbf{总体权重}（分层$W_h$、不等概率$1/\pi_i$）\textbf{纠正“样本组成不均衡}”，才能从样本无偏推总体。\item 即便如此，\textbf{“\emph{设计方差}}”只描述“\emph{选样造成的波动}”（区间多宽）；而\textbf{“\emph{无法裁决}带来的端点范围（识别范围$L$到$U_{\mathrm{err}}$）}”，\textbf{保留的是“\emph{证据本身还不能确定}”的那部分}——这两种不确定来源完全不同，普查能消前者、消不了后者。\end{itemize}
还要记住\textbf{不同监督信号要用不同的核验属性，\emph{不能}统一糊成一个“类别正确率}”：\textbf{任务答案}比对错、\textbf{关系}核对象/方向/条件、\textbf{表示目标}查构造条件与下游用途（7.1那张表）。

\textbf{第二组：“检测线索”和“质量裁决”永远分层}。\textbf{模型高损失、来源冲突、异常日志、离群/遗忘}——这些都只是\textbf{\emph{帮助安排复核}的\emph{排序信号}}，它们\textbf{\emph{不支持}任何对错结论}；真正\textbf{支持质量裁决的，只有“\emph{原始对象}、\emph{规范}、\emph{独立证据}}”（照片、采集记录、盲化复标）。

\textbf{第三组：“标签对错”和“训练贡献”是两个维度}。\textbf{Data Shapley}\textbf{在训练子集中公平地平均“边际效用}”（回答“平均而言这条样本对任务有多大功劳”，要跨大量重训、且依赖固定的学习算法+评价规则）；\textbf{TRAK}\textbf{用低维投影+局部牛顿，\emph{近似}“指定预测对训练样本的响应}”（快但是局部近似、不等于Shapley全局分配）。两者\textbf{都依赖学习算法和评价范围、都\emph{不能}代替“标签真值核验}”——高贡献不等于标签对、有影响不等于内容真。

\textbf{第四组：改完之后的“两种验证”分开做}：\textbf{对\emph{已查对象}，独立复核能确认“这些改动改对没有”}（结果复核）；\textbf{评价\emph{策略}}时则必须“\emph{先冻结规则、再用\emph{没参与开发}的证据}”（留出验证）。本章教学案例K17到K18正是范例：同口径配对得到$+4$个百分点净改善的\emph{点估计}，但近似区间仍包含恶化（$[-0.45,8.45]$跨0），证据不足，所以\textbf{\emph{暂缓发布}}——\textbf{绝不把“\emph{改过了}”当成“\emph{已经有效}}”。

最后把“\textbf{能不能发}”的完整判据列一次，它是\textbf{多目标、缺一不可}的，不能用某一项好看替代其他：\textbf{其一，新增错误（对到错）、其二，覆盖变化（有没有靠删难例/缩范围刷分）、其三，关键分组（各类/区域有没有被平均掩盖恶化）、其四，下游用途（标签变了模型到底变好没——且要控制变量）、其五，投入（人工/算力/维护全周期成本）}。本章K18\textbf{尚未重训下游模型}，所以\textbf{标签错误率的变化\emph{不能}算作“已取得的模型收益}”。

两条\textbf{任何统计技巧都补不回来}的底线要永远记住：\textbf{“参考依据本身错/系统遗漏”\emph{不由抽样区间自动补偿}}（区间只管抽样波动，管不了证据错和抽样框外对象）；\textbf{“缩小使用范围”\emph{不能}冒充“原总体质量改善}”（把难题排除后正确率上升是组成变化，不是质量进步）。

全章的\textbf{落脚点是“\emph{留痕}让质量环能转起来}”：把\textbf{来源、配置、版本、验证报告、未决问题}\textbf{持久地保留下来}，才能\textbf{按“\emph{已确认的根因}}”把问题\textbf{准确\emph{返回}到“获取、对齐、融合”等相应环节}，并让\textbf{新一轮的开发与验收\emph{继续保持证据边界}}（开发证据不当裁判、留出证据不进开发）。这样，“知识的质量”才不是一次性的验收，而是一个\textbf{可持续、可追溯、不自我证明的闭环}。

